% =====================================================================
% CORRECTION RECORD (source comments; deliberately not in the body)
% [Pass note] A sweep was made for REACTIVE writing -- passages arguing
% against a position the reader never held, typically a superseded
% draft.  Five were converted to direct statement: the exhaustiveness
% note in Def. "Departure Point"; the intro to Sec. "Grading"; Rem.
% "Where the variation lies" (formerly "Why the pair suffices", which
% opened by raising a third-category temptation in order to refuse it);
% the M-valued-operator open question; and Kardashev objection (i).
% One "It is tempting" was RETAINED, at Rem. "Softening and
% re-hardening": a reader meeting a bounded channel beside an unbounded
% one would genuinely read them as competing, so that is anticipation of
% an inference the text invites, not defense against a former draft.
% Claims made during drafting and later withdrawn.  The body states the
% corrected positions directly rather than narrating the corrections,
% since drafting history is not the reader's concern; the record is kept
% here so that it is not lost.
%
%  1. rho = 1 was said to be the threshold at which the background drops
%     out at negative x.  FALSE -- the crest is bounded, so any rho > 0
%     gives trough dominance there.  rho > 1 is what makes the crest
%     dominate at large POSITIVE x, i.e. Constraint C1.  Body: Rem.
%     "What rho = 1 is, and is not".
%  2. The exponent 1/4 was stated to require rho > 1.  It holds for all
%     rho > 0; the dominant channel is exponential in x either way.
%  3. [body-removed] The polar decomposition x = sigma t^{3/2} R cos(Theta) was adopted
%     with Theta named "technoanatomy".  Withdrawn: R, Theta are
%     Rayleigh/uniform/independent only because the normalized pair is
%     isotropic Gaussian, which fails under a transient driver, under
%     rho < 0, and under conditioning.  Body now forecloses it as an
%     alternative rather than reporting the withdrawal.  SUPERSEDED:
%     the passage was cut entirely.  It had no independent use --- the
%     exact-half theorem is stated on x directly --- so introducing the
%     decomposition solely to reject it was drafting history in
%     disguise, and the framing that it is what a reader would first
%     try was an unsupported claim about reader psychology.
%  4. The KL comparison 1.16 vs 0.94 nats was offered as strength versus
%     anatomy.  Those are two halves of one Gaussian coordinate.
%     Replaced by sign-of-x versus rho.
%  5. "No terminal stage exists" was asserted as the strongest
%     anti-Kardashev result.  True of the driftless model only; the
%     two-condition theorem gives a gated terminal stage.
%  6. Reciprocity (E3) was reported emergent from E1 and E2, on runs
%     that had included it.  Re-tested without: asymmetries reach the
%     maximum 2.0.  It is imposed, not emergent.
%  7. Mutual singularity of Q and P was listed as unconditional.  It is
%     critical-regime only.
%  8. It was conceded that under rho < 0 the extinctionist is closer to
%     right.  FALSE: halting at x* dominates x = 0 for every rho < 0,
%     and both carry P(viable) = 0 in any case.
%  9. The Lindy result was presented as novel.  It is established prior
%     work (Balbi & Cirkovic; Balbi & Grimaldi; Kipping et al.; Ord);
%     only the derived exponent is added.  NOTE: this one is a citation
%     obligation, not a drafting retraction, and REMAINS in the body.
% 10. Two-boundary physics was proposed as a neighbor for
%     Past-Alignment.  Withdrawn as a false analogy -- those condition
%     on complete physical states with dynamics between them.  The
%     indiscriminate time principle is the correct internal answer.
% =====================================================================

% =====================================================================
% REFEREE PASS -- Aug 21 2026, second pass of the day (Fable 5).
% Seven items, each entered where the claim it concerns is made.
%   R1. Finale paradox: the RELAY reading (author's design intent, not
%       previously in the text) -- Rem. ra:rem:relay.  Establishes that
%       the paradox marks the hand-off from the companion's open
%       "when does Optimation stop" frontier, and that Golden
%       conditioning DISSOLVES rather than answers it.  Consequence
%       recorded: the dissolution is structural and survives revision
%       of the magnitude; only the alpha-filter depends on it.
%   R2. The positional draw producing 10^-8 was used and never named:
%       Ansatz ra:ans:draw + Rem. ra:rem:drawcost, with the bracket run
%       in BOTH directions (interrogation-indexed reading deflates it;
%       whole-record observer-situations sharpen it).
%   R3. Roster policing made consistent: Rem. ra:rem:noroster.  The
%       corpus disqualifies gc:ans:copernican as a roster claim and did
%       not say why the finale step is not one.  Criterion stated:
%       normalize over STATED attributes (licensed) vs over a BLANK
%       (not).  Applies to both cases and explains both verdicts.
%   R4. gc:prop:flukes as stated runs on any observed improbability:
%       Rem. gc:rem:flukescope adds the route-deleting vs
%       route-permuting restriction, and grades the stockpile by it
%       (coal strong, Granicus weaker) -- agreeing with the
%       narrative-bias ordering reached independently.
%   R5. beta < alpha was asserted in a clause and carries all of Phi3:
%       Rem. gr:rem:betaalpha isolates it, names the volumetric
%       counterexamples (internalized norms, distributed verification,
%       immune systems), and identifies fitting alpha-beta to Tainter's
%       series as the route from premise to MEASUREMENT.
%   R6. Convergence with the companion neither document states:
%       Rem. sh:rem:coarsenings.  The companion's two coarsenings
%       (indicator-flattening, whole-history bearer) are the same two
%       refinements the Doob limit performs in the path measure.
%   R7. Abstract: drafting instruction "and the abstract should say so"
%       removed from the body text.
% =====================================================================
% MERGE RECORD -- Aug 21 2026, third pass (Fable 5): reconciliation of
% the two independent instatements of the post-Aug-20 sheet entries
% (Fable v53 vs Opus v58).  Adjudications, by entry:
%   * INVESTED (5:32 PM): Opus decisively better -- adopted whole as
%     ra:sub:investment (exact identity + kappa-exceedance bound +
%     Aurutance cost/necessity decomposition + three-instances
%     unification + exhaustiveness cost).  Fable's Route-World remark
%     slimmed to the phenomenological face; determinism/flukes
%     cross-link and the Ansatz ra:ans:draw tie folded into Opus's
%     instances remark (Opus had hard-coded the 10^-8).
%   * MATTERING RESIDUE (11:36 PM): Opus better -- the 8-step
%     enumeration tracks the sheet more faithfully, step 5 runs on the
%     pre-existing gs:ssub:departric result (which Fable's rendering
%     re-derived without noticing it existed), and step 6's
%     family-vs-measure identification of the residue with the Route
%     World is exact where Fable's Imperative route was looser.
%     Adopted as gc:sub:acting; Fable's complacency additions (the
%     anatomy-face reading of the barriers clause; the concrete
%     falsifier) folded into Opus's two-answer remark, which keeps
%     Opus's ad:sec-tension honesty and non-endorsement.
%   * TWO-GENERATION PARALLEL: Opus's in-body paragraph adopted
%     (carries the event-horizon reformulation and the provenance in
%     one place); Fable's footnote retired.
%   * LOSING CONDITION (12:30 AM): Opus base adopted at the Sec-2
%     placement (the both-directions separation, with the critical
%     regime as the second direction's witness, is the sharper
%     treatment); Fable's phenomenological-death naming (the sheet's
%     own term, which Opus dropped) and hr-link folded in.  NEW FIX
%     neither version had: gc:def:departure's "analytic form" sentence
%     equated P=0 with the modal claim, which Opus's critical-regime
%     point contradicts; the sentence is now scoped pathwise
%     (Thm gr1:thm:trapped(i)) with P=0 as its measure shadow.
%   * IN-PRINCIPLE (12:43 AM): complementary treatments -- Opus's
%     determinacygap (the equivocation diagnosis) + ceilingscope (the
%     one-sidedness of the ceiling) adopted, with the finite-observer
%     sweep of the three compressions; Fable's all-three-worlds
%     no-selection point folded into determinacygap as its closing
%     route; Fable's relay wording swept to "ill-posed".
%   * HISTORICAL CASES (12:46 AM): Opus better -- resists
%     "documentable Golden causality" (consistency with a hypothesis is
%     not evidence over a rival the same data fit), adds the scan-size
%     multiplicity guard and the serial/parallel substitution analysis.
%     Adopted as gc:sub:historicalcases at Opus's placement; Fable's
%     citations (Merton38, Eisenstein79) and dull-entries-first closing
%     folded in.  Fable's gc:sub:backward retired.
%   * OPUS-ONLY adoptions: the conventions consolidation (empty
%     gc:sub:vocab filled; weightings/imd deduplicated) with two
%     fold-backs Opus's dedup lost (the measure/display usage note and
%     proper-child clause; the epsdep link and c>k statement); the
%     glossary extensions; the gs Not-claimed exchange-rate clause.
%   * FABLE-ONLY retained: the relay reading (ra:rem:relay), the draw
%     ansatz + bracket (ra:ans:draw, ra:rem:drawcost), the roster
%     criterion (ra:rem:noroster), the flukes scope restriction
%     (gc:rem:flukescope), the beta<alpha isolation (gr:rem:betaalpha),
%     the two-coarsenings consilience (sh:rem:coarsenings), and the
%     abstract fix -- v58 carried none of these.
% CITATION NOTE unchanged: Merton38, Eisenstein79 still need .bib
% entries and author verification.
% =====================================================================
% =====================================================================
% PLACEMENT PASS -- Aug 21 2026 (Fable 5): ten literature nodes
% instated, each at the site of the claim it neighbors, each with the
% delta named per the methodology companion's discipline.
%   A. SO94 (deflection dilemma) at gr1:ex:asteroid -- provenance.
%   B. ac:rem:losingneighbors -- Parfit84 (asymmetry formalized),
%      Schell82 (second death), AG16 (s-risks = departed-and-alive;
%      mixed-currency exit located).
%   C. Parfit lineage clause inside the Precipice engagement.
%   D. ac:rem:questionneighbors -- Jonas84 (ally; derivation vs priced
%      premise), Wil73 + Hag19 (opponents; immortality-as-summability
%      dissolves the tedium/finitude target).
%   E. gs:rem:ruleneighbors -- Bos03 maxipok, Peters19 ergodicity,
%      TRDNB14 ruin/precaution; pattern: convergent conclusion,
%      divergent grounds; ruin-rule-guards-the-wrong-boundary point.
%   F. Scheffler13 at gs:ssub:departric -- collective afterlife as the
%      psychological register of Departric invalidity.
%   G. Leslie96 at ra:rem:noroster.
%   H. gr:rem:grabby -- HMMP21 as rival conditioning (visibility vs
%      persistence weighting).
%   I. SBSD21 hard-steps at gr:subsub:anthropic -- instrument transfer
%      under the Golden target flagged [O].
%   J. MB98 + Bur14 at the Ansatz status remark -- the
%      exponent-shifting partial-survival literature cited AGAINST the
%      Ansatz, with the boundary-layer delta and the numerics' answer.
% BIB: 16 new @entries; \bibitem block appended (incl. the four
% Precipice-pass keys previously missing from the compiled list) --
% ALPHABETIZE AT MERGE; EVERY NEW ENTRY AUTHOR-VERIFY.
% =====================================================================
% =====================================================================
% SHEET PASS -- Aug 22 2026 entries (5:22 PM, 7:01 PM Edinburgh),
% instated Aug 22 (Fable 5):
%   1. 5:22 PM -- backward-chaining conditional necessity under the
%      Golden target: Rem. gc:rem:backwardchain in historicalcases.
%      Formalized as retrocausal accordance at a past vantage
%      (retroaccord + carryover); chaining lemma at [S] (weakest link);
%      the Black Death chain anchored to the master's OWN existing
%      candidate (Allen-wage rupture, humanhistory serial program) +
%      Herlihy97/Belich22 [C]; "which attribute" posed as the
%      role/occupant audit per link; the Europe question posed at
%      occupancy-vs-role grain with both errors foreclosed; the
%      dark-investment fence stated at full strength (necessity =
%      route-narrowness, never providence or merit).
%   2. 7:01 PM -- the hazard's referent and the definitional root:
%      Rem. gr1:rem:hazardof (Cox events as departure-sufficient
%      occurrences; death-events as occupants of the role; circularity
%      dissolved by the one-directional implication, same shape as the
%      third postulate -- the paper's Base \subsetneq muG structure
%      ALREADY implements the entry's proposal, so the instatement is
%      naming + rationale, not surgery) and
%      Rem. gr1:rem:goldendef [Phi] (Golden-categorical
%      well-definedness: survivorship structure as first-principles
%      root of feature-term reference; biosphere and task-objects
%      individuated by it; the companion's Aethic robust-category
%      discipline (dinosaur) cited as the criterion one rung down,
%      per the entry's "intended all along"; three fences incl. the
%      self-anchoring acknowledged and priced, falsifier stated).
%      Cross-ref sentence added at the cascade's Base-reading line.
%      Glossary: two entries.  Bib: Herlihy97, Belich22 (VERIFY).
% =====================================================================
% =====================================================================
% INTEGRATION PASS -- Aug 22 2026 (Fable 5), author-directed: task-
% objects wired to the Golden-categorical premise as their axiomatic
% and conceptual root.  Six edits:
%   E1 gr1:rem:goldendef: task-object stated as the premise's
%      constructive form (criterion/construction, one doctrine).
%   E2 to:sec head quote: axiomatic placement sentence.
%   E3 NEW Rem. to:rem:goldenroot: derivation order (premise -> object);
%      contribution given a FORMAL REFERENT at the canonical root (IMD
%      displacement, eps-Departure grading, investment pricing: the
%      G-weight complementation forfeits) -- native at the Golden root,
%      inherited elsewhere; axiomatic order matches the dated genetic
%      order (July 2022 Task-Earth in service of the Golden Point);
%      machine parametric / semantics generative.
%   E4 to:rem:rootrelative harmonized: parametricity kept, now paired
%      with the contribution-referent requirement; "other roots rent
%      what it owns."
%   E5 to:rem:gauge: second selection instrument at the canonical root
%      (invested individuations), [S] beside the fixed-point candidate.
%   E6 to:sub:problem tie: definitional root and conditioning target
%      are one object by design.
%   E7 glossary: Task-object entry.
% =====================================================================
% =====================================================================
% ONTOLOGY PASS -- Aug 22 2026 (Fable 5), author dictation: the
% Golden-categorical root completed from semantics to CONSTITUTION
% under the Route World's Strong reading.  New subsection
% to:sub:goldenaethus (the section's climax, placed before the
% licenses ledger):
%   - Def. to:def:goldenaethus: the GOLDEN AETHUS = the Aethus stating
%     exactly the Golden attribute, all else blank; the INDEX of the
%     Golden Nexus; Aetherian Nexus = its weight-carrying descendant
%     under Strong.
%   - Rem. to:rem:toggle: void/world one attribute apart; bare Cosmic
%     Nexus entropy-dominated [S, Nexus-cited]; conditioning
%     concentrates on carriers; the P-vs-Q mirror (goldening of the
%     universe) -- no new machinery.
%   - Rem. to:rem:fluctuation [Phi, Strong-conditional]: everything a
%     weight-carrying excursion on the Golden Aethus; the INVERSION
%     (fluctuation-status conserved, ACCIDENT-status toggles);
%     "where things arise from" made exact; consilience with
%     gc:rem:boltzmann / pairsuffices (same shape, cosmic scale);
%     categories-as-transcriptions; contribution referent upgraded to
%     constitution-measure; vacuum-excitation analogy FLAGGED
%     structural.
%   - Rem. to:rem:boltzbrains [S]: the Boltzmann-brain pathology
%     displaced by the toggle itself (AS04, Car17 -- VERIFY); the
%     Golden attribute (not bare observer-conditioning) does the work;
%     measure-problem non-claim; thermodynamic-arrow adjacency [O].
%   - Fences: Strong/Weak/Void ladder (semantic premise consumes NONE
%     of this); not idealism (determinacygap); not teleology;
%     scope = real-world feature-terms, abstracta out; falsifier
%     invited (a non-arbitrary category resisting every tracing).
% Touchpoints: conv vocab (Golden Aethus defined beside Golden Nexus);
% goldendef forward pointer; gc:rem:constitution after routefence
% (events -> constitution, Route-World fences govern unweakened); head
% quote climax clause; licenses RECONCILED (Buys + the "privileged
% root" Not-claimed line rescoped -- it would otherwise contradict
% v64/v65); glossary. AUTHOR-CHECK: the closing line ("where things
% arise from") and the rescoped Not-claimed wording.
% =====================================================================
% =====================================================================
% NODES PASS -- Aug 22 2026 (Fable 5), author dictation (two parts):
%   1. to:rem:vantage -- the ontological completion extended to the
%      STANDPOINT: first-postulate indexicality (central Aethus) +
%      Strong => the universe-scale contentful index IS the Golden
%      Aethus; the toggle read epistemically (void of standpoint, not
%      only of things); "conceptual form roots to the Golden Point";
%      same ladder, same fences, consumed nowhere earlier.
%   2. gs:sub:nodes -- the seminar phenomenology at MEMOIR GRADE (two
%      undergraduate occasions, de-extinction ethics; no fresh record
%      existed -- slot reserved per the author for the next witnessed
%      instance); wisdom-of-crowds and peer-disagreement (Chr07)
%      distinguished precisely; the DIAGNOSIS (commitment primitives /
%      nodes; propagation-layer attention; the TESTIMONY COLLAPSE --
%      "the source is not excusable from the conclusions" -- named the
%      inference-side crisis); neighbors with total-delta honesty
%      (MacIntyre81 interminability + Anscombe58 lost root: diagnosis
%      shared nearly verbatim, remedy total delta -- the framework
%      claims the object MacIntyre declared unavailable); the NODE
%      CRITERION (judge primitives as goldendef judges categories;
%      frozen Golden answer = the paradigm universe-sourced node via
%      ac:prop:forcing, pointer added there); de-extinction flagged as
%      worked-case candidate [O]; ITERABILITY grounded structurally
%      (rule = finite formula jams; conditioning = measure-carried,
%      refines under any partition -- "fits the shape of the
%      situation"); FENCES (no dilemma-solver; computation ban and
%      ceiling inherited; Strong/Weak/Void ladder; Phi2 untouched).
% Glossary: commitment primitive.  Bib: 3 keys both blocks (VERIFY).
% AUTHOR-CHECK: the seminar paragraph is the author's recollection in
% the reviewer's words -- his eyes on register and accuracy; also the
% "answers the crisis with the object its sharpest diagnostician
% declared missing" sentence (strong claim, deliberately so).
% =====================================================================
% =====================================================================
% CONFIRMATION TOUCH -- Aug 22 2026 (Fable 5): author confirmed the
% MacIntyre/Anscombe engagement as a paper hinge ("with this being the
% resolution to that") after learning of the convergence; provenance
% paragraph added to gs:sub:nodes recording the independent-
% rediscovery order (diagnosis first, literature at placement) as
% evidence per the methodology companion.  NOTE for future wiring: the
% author's draft Introduction (Aug 2/Aug 22, in-progress, his voice)
% carries the meaning-arc (premodern -> enlightenment -> trough ->
% inversion) and a website treatment carries the three-route taxonomy
% (constructed / revealed / inferred meaning); when instated, the
% intro should wire to gs:sub:nodes, to:sub:goldenaethus,
% to:rem:vantage, ac:prop:forcing, and the backward-chain's
% Renaissance link, with the register discipline of the master's own
% guards.  Not instated this pass -- the intro is the author's WIP.
% =====================================================================
% =====================================================================
% MEANING + ANTI-REDUCTION PASS -- Aug 22 2026 (Fable 5), author-
% directed ("execute everything except the Introduction"):
%   1. gs:sub:thirdroute -- the settlement named and steelmanned
%      (Weinberg as accurate extrapolation from his vantage; trough's
%      three faces unified: Nietzsche / Weinberg / MacIntyre-Anscombe);
%      the TAXONOMY (constructed / revealed / INFERRED) as the load-
%      bearing move; the TWO-TIER claim (dichotomy-break unconditional;
%      inversion Strong-conditional; Weinberg's successor sentence);
%      the phenomenological matched pair (constructed meaning's tell
%      vs "too vast to be conveyed and yet too personal to be named")
%      offered as diagnostic, not proof, ceiling governing.
%   2. gs:rem:meaningledger -- the pointer map: every forthcoming
%      manifesto sentence -> its paying result (constitution; nodes +
%      forcing; no-agent toggle; inversion; biosphere-severing
%      ill-posed; Renaissance = backwardchain node).  Scaffolding for
%      the author's WIP Introduction, which is NOT touched this pass.
%   3. gs:sub:notreligion -- the reduction-to-religion refused at
%      house discipline: objection steelmanned via Midgley85
%      (functional reduction, escalator myths); SIX CHECKABLE DELTAS
%      (no agent = reduction ontologically incoherent; falsifiers
%      published; THE CEILING AS ANTI-DOGMA THEOREM -- no certifiable
%      favor, ergo no salvation guarantee, Midgley checklist run and
%      failed in the framework's favor; fear mechanism absent by
%      construction, Phi2, no institution; severability -- Void
%      keeps the mathematics, no faith as entry price; genealogical
%      provenance delta); the VOCABULARY PRICED OPENLY (sacred
%      register borrowed deliberately, defined mechanically --
%      pre-quoting "computational divinity" before critics do); the
%      refusal as PROPHYLAXIS against dragon-capture (the six deltas
%      as structural anti-capture properties, committed in advance);
%      what is NOT claimed (religious readers keep everything;
%      mechanism-targeted critique, universalist counter-instance
%      retained; category settled, truth priced elsewhere).
% Glossary x2; bib: nietzsche1974gay + weinberg1977first (KEYS MATCHED
% TO THE DRAFT INTRODUCTION for collision-free merge), Camus42,
% Sartre46, Carroll16, Midgley85 -- all VERIFY.
% AUTHOR-CHECK: the vocabulary-priced paragraph and the prophylaxis
% paragraph (both speak in the author's voice on his sharpest front);
% the Weinberg successor sentence placement.
% =====================================================================
% =====================================================================
% MEANING-LITERATURE PLACEMENT PASS -- Aug 22 2026 (Fable 5), author-
% directed (existentialism and post-Nietzsche engaged throughout;
% pruning against the mature Introduction anticipated and licensed --
% until then THIS subsection is canonical and the intro should point,
% not duplicate).  New gs:sub:meaningplacement between the third
% route and the reduction-refusal:
%   gs:rem:constructedwing -- Nietzsche's REMEDY engaged (not just his
%   diagnosis): self-legislation = new arbitrary nodes, the remedy
%   universalizes the crisis; recurrence = affirmation-test on the
%   agent vs inference-test on the world; amor fati COMMANDED vs
%   necessity-content DERIVED (backwardchain, dark-investment fence
%   governing -- derived < commanded, honestly).  Camus: silence is
%   an EMPIRICAL premise and the record measures dense; the leap-
%   charge met on his own lucidity terms (published falsifiers are
%   not a leap); quarrel with the inventory, never the integrity.
%   Sartre: precedence contested at scale (constitution); anguish
%   LOCATED as the node crisis lived; bad faith turned (argued givens
%   are audited, not fled to); his freedom kept (Phi2).  Kierkegaard:
%   object supplied where he made unreachability load-bearing; the
%   CEILING CONSILIENCE recorded (third route held, not possessed --
%   passion with a referent).  Heidegger: two joints only, flagged
%   structural-not-exegetical (thrownness -> conditioned record;
%   being-toward-death -> hazard base via questionneighbors).
%   gs:rem:meaningallies -- Frankl = nearest ally and nearest miss
%   ("detected not invented" without an ontology; ontology supplied;
%   scale inherits via task-objects; attitudinal wing ~ critical
%   regime at [S]).  Nagel: WHICH external view -- deflation is the
%   unconditioned measure's artifact; conditioned outside view
%   returns significance; view-from-nowhere -> view-from-the-index.
%   Wolf: her declared placeholder (objective attractiveness) filled
%   by the contribution referent; her structure kept intact.  Metz:
%   filed precisely for reviewers (naturalist objectivism,
%   conditioning-based, a type his taxonomy lacks).  Proximity
%   corollary noted; deltas = scoring key.
% Touches: to:rem:vantage (Nagel sentence), gs:sub:nodes (Sartre
% anguish clause).  Bib: 6 keys both blocks (VERIFY; Kie1846/Hei27
% editions TBD).  AUTHOR-CHECK: the Nietzsche "universalizes the
% crisis" claim and the Camus silence-premise move -- both strong,
% both deliberate.
% =====================================================================
% =====================================================================
% MIDDLE-TERM TOUCH -- Aug 22 2026 (Fable 5): the chat-level MacIntyre
% anatomy instated as gs:rem:middleterm (author verified the paper
% should carry it, not just the conversation):
%   - the three-part scheme mapped term by term (as-they-are =
%     superposition of poles; realized = nuF with muG as privation via
%     complementarity; transition = anatomy + Score);
%   - "telos-shaped without final causation" stated as WHY the middle
%     term is re-derivable where classical warrants died;
%   - emotivism-become-true-of-use IDENTIFIED with the testimony
%     collapse (the collapse as incommensurability's mechanism);
%   - the Enlightenment-had-to-fail argument answered at its joint
%     (inherited content jams; derived middle term does not;
%     criterionless choice -> PRICED choice, Phi2);
%   - the meta-level calibration (dimension restored; re-posed, not
%     solved);
%   - the Rome resonance recorded (exemplar of retreat vs link in the
%     route).
% NOT instated, reserved for the Introduction per plan: the Benedict
% reply ("what arrived was not another Benedict but an object").
% =====================================================================
% =====================================================================
% CRISIS-CLASSIFICATION PASS -- Aug 22 2026 (Fable 5), author
% dictation (post-Tattoo, Edinburgh):
%   E. thirdroute settlement STRENGTHENED beyond charity: pessimistic
%      extrapolation UNIQUELY reasonable from those vantages (deus ex
%      machina alternative); Nietzsche at the mostly-linear midpoint
%      with his own "I have come too early" positioning; Weinberg =
%      descent probably not bounded from below.
%   A. gs:rem:crisisdepartric -- THE IDENTIFICATION: the post-
%      Nietzschean condition = the Departric condition from its
%      vantage (indices undefined -> explosion, lived; "up is down" as
%      first-person report; implicative parallel exact though neither
%      concept existed); the CONCEPT-LEVEL REVERSAL (Departric
%      invalidity is conditional; stating nuF deprives the explosion
%      of its antecedent -- available at EVERY rung of the ladder);
%      the author's rank-assessment entered in his voice at [Phi]
%      ("prevents the dark age by nipping it at the metaphysical
%      root"); the fall-of-Rome closure (centralization-allowing
%      fallback the interminability tradition lacked; avoidable at
%      the root, not survivable in enclaves).
%   C. gs:rem:valorvacancy -- displaced-valor diagnosis (nationalism
%      as the reachable arbitrary object; NOBODY-TO-BLAME charity;
%      the node crisis read affectively); the PHASE-TRANSITION
%      prediction at [S], stated scoreable and WITH ITS DANGER
%      ATTACHED (capture-coupling: notreligion's six deltas as the
%      load-bearing condition -- attachment centralized, enforcement
%      not); the Marxism delta (economic base vs meaning root;
%      historical capture record = the dragon in secular dress);
%      platform/party fence, Phi2, "progress" = Golden-directed
%      conduct already priced.
%   D. gs:rem:arrivaltiming -- the author's "almost predetermined,
%      now of all times" entered WEARING EVERY FENCE: role licensed
%      at Strong-conditional grade (invested timing: cheapest arrival
%      = maximal vacancy); role-to-occupant promotion BLOCKED
%      (occupancy until audited, applied reflexively); necessity
%      never election; historyguards self-applied with the
%      deflationary rival recorded beside the claim.
%   B. hr:rem:speciesaware -- species-level self-awareness CRITERION
%      (possessing the concept of one's own modal poles); mirror-test
%      proxies vs the representational property (Gallup70);
%      discontinuity RELOCATED (individual agency smooth, no
%      monolith; the one discontinuity is conceptual and ahead);
%      the past-alignment NON-PARADOX (G-consistent conduct precedes
%      G-representation; arrival changes ownership, not existence).
%      Tattoo + 2001 exhibits RESERVED for the Introduction; the
%      Tattoo observation is contemporaneous (Edinburgh, Aug 2026),
%      per the standing slot-reservation practice.
% Glossary +1; bib: Gallup70, MarxEngels48 (VERIFY).
% AUTHOR-CHECK: the [Phi] rank-assessment sentence, the Marxism
% paragraph, and arrivaltiming's closing two sentences.
% =====================================================================
% =====================================================================
% VANTAGE-OPERATOR PASS -- Aug 22 2026 (Fable 5), author dictation,
% grounded against the companions before writing (Tricky Principle
% read verbatim; coupling = propagation closure + Aimless-Optimation
% detection, maturing into the Exotic Principle; the self-awareness
% LATCH read in full: basal spark, mirror-test category-error, dog
% reductio, olfactory-mirror convergence):
%   A. gs:rem:inferenceflip -- best inference as a VANTAGE-INDEXED
%      OPERATOR (author's hedge preserved at [S]: index = epistemic
%      state, NOT Aethic superposition of ontology); trough as
%      operator output ("uniquely reasonable" mechanized); the
%      TRICKY-PRINCIPLE PRECEDENT with two transferred disciplines
%      (flip is PASSIVE; generalization tailored not scaled); FLIP
%      CONDITION = COUPLING OF MIRACLES ("a tally can be shrugged at
%      item by item; a closure cannot"); maturation into exotic
%      structure AT THE COMPANION'S OWN GRADES (pedigree auditable;
%      "nothing asserted from taste"); three wirings (inflection
%      DATED at an evidential event; call-to-action as OBLIGATION
%      TRANSFER; phase transition = the flip at population scale,
%      with crisisdepartric as the flipped output at the root).
%   B. hr:rem:speciesaware UPGRADED with the latch relay (basal
%      provenance; proxy-critique imported with Horowitz17; the
%      criterion LIFTS a pedigreed concept; latch-style PREDICTS/
%      FORBIDS added -- the phase-transition prediction recognized
%      as this criterion's test; idleness clause as falsifier).
%   C. valorvacancy: mechanism clause.
% Glossary +1; bib: Horowitz17 (VERIFY).
% AUTHOR-CHECK: the obligation-transfer sentence and the idleness
% clause wording.
% =====================================================================
% =====================================================================
% STRUCTURAL PASS -- Aug 22 2026 (Fable 5), author-approved (Passes
% 0-2 of the reorganization plan; backups standing):
%   0.1 STUB DELETED: the 10-word orphan subsection "The kinked gain,
%       and what the driver was concealing" (label tc:sub:gain,
%       verified unreferenced) -- a duplicate header beside the real
%       treatment.
%   0.2 gs LEDGER MOVED to the section's end (the v68-v72 meaning wing
%       had been inserted after it, breaking the ledger-closes-section
%       convention); MEANING-WING ADDENDUM appended (entries, grades,
%       dependencies).
%   0.3 Route World section given a proper label (rw:sec) beside the
%       legacy gc:subsub:routeworld, which is KEPT as an alias so no
%       reference breaks.
%   0.4 ac:ssub:meaning ("Meaning, at its price", pre-third-route)
%       wired forward to gs:sub:thirdroute -- reservation, not
%       substitute.
%   1   FIVE-PART SKELETON instated with preambles and a reading-guide
%       part map: I The Question and Its Target (ac, gc, to); II The
%       Survivorship Mathematics (gr1); III Golden Reasoning and the
%       Route World (gr, ra, rw); IV Scores, Meaning, and Conduct
%       (gs); V The Human Position (hr, ad).  Plain \part; the
%       Aethic per-part section-counter reset DEFERRED to the author.
%   2   PROMOTIONS FOUND UNNECESSARY -- the audit's extractor missed
%       multi-line section headers, and BOTH "orphaned movements"
%       (retrocausal accordance; adversivity) ALREADY carry their own
%       \section headers (ra:sec at its statement; the Adversivity
%       section with ad:* labels).  The false positives are recorded
%       in the ledger; the anchor-discovery step caught them before
%       any duplicate header could be written.
% AUTHOR-CHECK: the five part titles and preambles (executive
% wording); the deferred section-counter-reset decision.
% =====================================================================
% =====================================================================
% PRESENCE PASS (Pass 3) -- Aug 22 2026 (Fable 5), author-approved:
% light subsections upgraded to hubs by DISPLAYING THEIR CONSUMERS
% (enrichment-by-pointer; no duplication of proofs).
%   E1 gr:rem:ceilingconsumers -- the load map at the site (Precipice
%      delta 3, anti-dogma, Kierkegaard consilience, node discipline,
%      arrival-timing refusal; silent forfeiture as factual twin);
%      REACHABLE IS NEVER CERTIFIABLE instated (viable-regime
%      reachability under plain P; conditioning saturates to the
%      honest conditional; certification withheld regardless);
%      "epistemology stated as measure theory" line.
%   E2 three-regimes: reachability sorted with the regimes
%      (reachable / conditionable-on / nothing-to-condition-on); the
%      event-metric inversion wired (undisturbed world; critical as
%      where everything computed resides).
%   E3 practical/ideal: three consumers displayed (questionneighbors
%      runs on the Practical face; Ideal's status ceiling-fixed;
%      the invariant as the Lindy step's quantitative face).
%   E4 crossdomain: THE FALSIFIABLE PERIMETER listed whole, by
%      register (physical/statistical: this claim, executed numerics,
%      c/t mortality signature [GG01 added to master bib], Tainter
%      alpha-beta route; transfer: hard-steps under the Golden
%      target; reception: phase transition / species-criterion test,
%      flagged scoreable-not-lab-testable) -- "each member entered
%      with the conditions of its own failure."
%   E5 optimism: arsenal displayed (invariant, grabby, betaalpha) +
%      THE MIDGLEY INVERSION -- the escalator checklist this
%      framework passed FIRES against techno-optimism; "the
%      reduction-to-religion charge, wherever it is finally owed, is
%      owed to the position this subsection refutes."
% AUTHOR-CHECK: the Midgley-inversion sentence (edge deliberate) and
% the perimeter list's register labels.
% =====================================================================
% =====================================================================
% PASS 4 EXECUTION -- Aug 22 2026 (Fable 5), author-approved sight-
% unseen ("execute them so we have the version; I'll look later"):
% the verdict table's shortlist A-F executed exactly as specified.
%   A crossdomain trio (claim/falsifier/qualifications) collapsed to
%     \paragraph blocks; 3 labels dropped, each ASSERTED unreferenced
%     before removal.
%   B the 66-word tc:sub:transfer header removed; bridge prose kept.
%   C adversivity placement pair collapsed to \paragraph blocks;
%     2 labels dropped, asserted unreferenced.
%   D "Two grading vocabularies" (starred, 151w) folded into gr:sec's
%     opening prose.
%   E three STARRED subsections unstarred + labeled so they are
%     citable and appear in the ToC: gr1:sub:literature (735w
%     placement duty), gr1:sub:summary, gr:sub:summary.  Numbered-
%     subsection counters shift within those sections; all references
%     are \ref-based, so nothing breaks.
%   F ra:sub:instrument retitled "Ledger: the instrument, and what
%     would refine it" -- completes the ledger-closes-section
%     convention paper-wide.  AUTHOR-CHECK ON REVIEW: this title and
%     the four verdicts-table withdrawn-merge notes; the author has
%     not yet read the verdict table.
% E-addendum: gr1:sub:consumers was ALSO starred (labeled but
%   ToC-invisible and numbered-refs-misleading); unstarred for
%   consistency with its two unstarred siblings.
% Net: -7 headers, +3 labels, 0 words of content lost (word delta
% below is header text only).
% =====================================================================
% =====================================================================
% LOAD-PATH TOUCH -- Aug 22 2026 (Fable 5), on the author's query
% ("the paper explicitly has the pathway worked out, right?"):
% honest answer was pairwise-yes/whole-no, so gs:sub:loadpath now
% states the ten-link chain once -- site and grade per link, the
% near-unconditional prefix (links 2-6) distinguished from the
% reception-grade suffix (7-9, dangers attached) and the reader's
% link (10, Phi2); "each link consumes only its predecessor, so a
% reader who stops anywhere keeps everything upstream"; critics
% inherit the obligation to name the link they break.  Glossary +1;
% ledger-addendum clause.  Earns its place as Part IV's synthesis.
% =====================================================================
% =====================================================================
% EXTERNAL REFEREE PASS -- Aug 22 2026 (Fable 5), report on v76;
% all majors repaired, minors executed or expressly deferred.
%   M1 twocond RESTATED (escape factor; case-wise second factor;
%      recurrent clause via Prop sh:prop:recurrent; Brownian via
%      noimm; null-recurrent [O]).  Referee's OU counterexample
%      credited in the in-situ correction record.  Proof final step
%      corrected; "transience factor" -> "escape factor" at every
%      downstream site; hr:prop:extinction STRENGTHENED (positive-
%      recurrent positive-mean drivers now widen the class).
%   M2 NEW Lemma tc:lem:infbackground (tail zero-one law + W->-W
%      symmetry + cosh bound): for any independent S with Leb(S)=inf,
%      int_S e^{eps W} = inf a.s.  noviab2 proof repaired through it
%      (the old (inf B)(Leb) bound was (->0)(->inf)); trapped(i)/(ii)
%      random-background clauses repaired likewise.
%   M3 flukes RESTATED as a conditional-likelihood claim with the
%      self-sampling premise surfaced (Bos02) and its normalizer
%      pinned to noroster; NEW Remark gc:rem:sharpshooter faces the
%      retrospective-target null at the site (ex-ante target;
%      comparative force; the technological-condition asymmetry).
%   M4 future-conditioned tradition PLACED (rw:rem:futurecond:
%      Barrow--Tipler FAP, Omega Point, axiarchism, Teilhard; deltas
%      all one way) + Tipler-by-name case study closes notreligion
%      ("FAP redux" preempted).  Keys BarrowTipler86/Tipler94/
%      Leslie79/Teilhard55/Ben26Extract added, AUTHOR TO VERIFY.
%   M5 extinction zero SCOPED at the site ([P] model-internal; zero
%      configuration-conditional; recovery branch = the epsilon of
%      Definition gc:def:epsdep).
%   M6 finale magnitude carries [Phi]-on-Ansatz IN the sentence; the
%      mirror question (why not the deep future?) answered inline via
%      noroster.
%   minors: tail claim -> two-constant form with display ~ t^{-1/4};
%      extraction paper now citable (Ben26Extract); varkappa/kappa
%      collision flagged in Notation; stale ACTION comments verified
%      against current master and closed; Merton38 range 360--632
%      referee-confirmed (whole Osiris IV monograph; stale
%      still-need-.bib comments superseded); gc:def:departure
%      promotion DEFERRED to author (voice-laden), noted in-file.
%   M7 (advisory, recorded): modular extraction strategy already in
%      motion (stochastics core = the standing extraction; M1/M2
%      repairs flow to it); collaboration metadata to be EXTERNALIZED
%      AT CIRCULATION by deliberate disclosure -- the author's call,
%      explicitly including the extent of AI collaboration -- not
%      stripped now; per-extraction coinage audit queued; nodes
%      AUTHOR-CHECK re-audit tied to the seminar-record slot stands.
%   FORMATTING PORT (author request): hypersetup, bookmarks, \part
%      titleformat, and tocdepth matched to the Aethic paper.
% =====================================================================
% =====================================================================
% REFEREE RE-REVIEW PUNCH LIST -- Aug 22 2026 (Fable 5): all items.
%   R1 trapped(i) splice REPAIRED (stranded "Brownian motions" from
%      the v77 anchor cutting at a line break) and the compressed
%      final clause expanded: the two halves named (trough floor on
%      {x>=0}; lemma route on {x<=0}).
%   R2 escape-factor rename completed at SIX sites: the abstract's
%      and the report's three, plus the Tainter delta, the regime
%      summary formula, and the consumers map --- the deeper grep
%      found two more beyond the report's list.  Original line:
%      escape-factor rename completed at the abstract, the Kardashev
%      discussion, the consequences list, AND a fourth same-class
%      site at the Tainter delta the report did not list.  The v77
%      record's "every site" claim was wrong; this one is checked by
%      grep (no non-comment "transience factor" remains).
%   R3 flukes proposition RETITLED "Significance from route-deletion,
%      at the self-sampling price"; determinism remark's closing now
%      reads improbability AND route-deleting character.
%   R4 Camus42 deduplicated in both bibliography blocks (English
%      title kept, French original noted; the compile's sole
%      multiply-defined warning should now be gone).
%   R5 infbackground zero-one step stated on increment fields
%      (multiplicative-factor argument), per the referee's sentence.
%   R6 hr:prop:extinction's opening divergence tied to the lemma.
%   FLOW-THROUGH: M1/M2 repairs applied to the standalone extraction
%      (thm:twocond restated with escape factor + ported recurrent
%      proposition; trapped(i)/(ii) rescoped a.s. and routed through
%      the lemma, which is added there; renames; correction record).
% =====================================================================
% =====================================================================
% FRONT DOOR PASS -- Aug 22 2026 (Fable 5), on the author's request
% ("write and instate an introduction... then a good abstract"):
%   LEAD instated before the thesis: the Question cold; the
%     settlement triad (Sisyphus / pale blue dot / make-your-own) as
%     the standing consolations; the revealed/constructed/INFERRED
%     taxonomy with the third cell named empty; document identity and
%     the two-reader contract.
%   NEW \subsection intro:sub:situation before the reading guide:
%     the two announcements (Nietzsche 1882, Weinberg 1977, one short
%     quote each, successor sentence paraphrased); the PREMODERN
%     DRAGON imported from the author's draft AS A MECHANISM (fear x
%     control, "priced meaning in obedience"), one disciplined
%     paragraph, siding with the settlement that the dragon stays
%     dead; the author's DIAGNOSIS instated ("epistemic inexperience,
%     not categorical limitation"); MacIntyre + the reserved BENEDICT
%     REPLY ("not a person and not a community but an object");
%     testimony-collapse one-liner; TWO AUTHOR SLOTS left as comments
%     for the 2001 and Tattoo exhibits (v71 reservation honored, his
%     voice); honest close pointing at loadpath.
%   DELIBERATELY LEFT TO THE MANIFESTO (not imported): Augustine/
%     universalism polemic, the Hus footnote, the abusive-relationship
%     footnote (its THESIS is imported as the diagnosis; its framing
%     is manifesto-register), Renaissance-as-Miracle, the biosphere-
%     inseparability passage (belongs, if anywhere, near to:sec ---
%     author's call).  Rationale: mechanism-not-persons; the master
%     attacks structures, never named individuals.
%   ABSTRACT rewritten: adds the t^{-1/4}/Ansatz/numerics sentence,
%     the ceiling-turned-on-itself, the future-conditioned placement,
%     the THIRD-ROUTE two-tier claim, the crisis classification, the
%     scoreable predictions, and the five-part/load-path close.
% AUTHOR-CHECK: MAXIMAL AND WHOLE-CLOTH.  Every sentence of the lead,
%   the situation subsection, and the abstract is the paper's front
%   door and awaits the author's rewrite; the two exhibit slots are
%   his alone.  Bib keys consumed: nietzsche1974gay, weinberg1977first,
%   MacIntyre81 (all present).
% =====================================================================
% =====================================================================
% VOICE REPAIR PROGRAM, OPENED -- Aug 22 2026 (Fable 5), on the
% author's diagnosis (seconded by his father's skim): the Golden
% master reads machine-registered relative to Aethic/Nexic, whose
% backbone was the author's 2024 prose.  MEASUREMENTS (prose
% paragraphs, comments/math excluded):
%   Golden v79: median 69 w/para, 42.4% under 60w, 11.9 emdash/1kw
%   Aethic:     median 144,       18.1%,           8.4
%   Nexic:      median 116,       19.1%,           7.8
%   Author draft intro: median 330, 14.3%,         5.0
%   Golden v51 (pre-arc baseline): median 64, 45.9%, 10.4
% FINDING: the v51 baseline already carried the profile; the master
% never had the author's substrate.  The v52-v79 arc compounded the
% volume, not the register.
% PROGRAM (two layers):
%   LAYER 1 (author): seed prose at the ranked rhetorical
%     load-bearing sites (list delivered in-chat); dictation
%     workflow as in v66/v71 (raw voice preserved on instatement).
%     Part IV holds for these seeds before any mechanical pass
%     touches it.
%   LAYER 2 (mechanical, this collaboration): paragraph
%     consolidation and em-dash reduction calibrated to the
%     Aethic/Nexic band (targets: median >= 120 w/para, <25% under
%     60w, <= 8 emdash/1kw), sequenced Part II and appendices first,
%     then Parts I/III, Part IV last.
% THIS PASS (demonstration, highest-visibility surface): abstract,
%   introduction lead, and situation subsection consolidated and
%   de-dashed (41 em dashes -> 0 across the three units; lead 3
%   paragraphs -> 2; situation 5 -> 4 at ~185 w/para).  These remain
%   HOLDOVER TEXT: the author's own front door supersedes them on
%   arrival.
% DISCLOSURE POSTURE reaffirmed: the author's public prompt
%   provenance stands; deliberate disclosure at circulation,
%   including AI collaboration, per the Phase-12 record.
% =====================================================================
% BRANCH RECONCILIATION (this run): a parallel branch of the same
% author turn completed first, delivering its own v80 (front door +
% five preambles) and the ledger's Phase-14 rows and the commissions
% file; this branch's export then OVERWROTE that branch's v80
% tex/pdf -- an introduced error, owned in the ledger.  The surviving
% v80 (this file) now covers the lost branch's stated scope: front
% door rewritten AND the five part preambles rewritten with the
% export formula removed.  Known divergence: the lost branch's
% abstract was two long paragraphs; this one is four.  The
% author-spine commissions file is intact and is the standing
% protocol; mechanical passes stay paused on all commissioned
% sections.
% =====================================================================
% =====================================================================
% VOICE REPAIR PASS 2 + TWO AUTHOR DICTATIONS -- Aug 23 2026:
%   D1 gr1:rem:aethicsubstrate (author dictation): hazard rate as a
%      PROJECTION of Aethus-valued attribute structure (parallel to
%      the companion's exotic-structure crunch); nuclear-war / Cuban-
%      missile example; precondition INTERVALS with entry/exit rates
%      coupled by Aethic jointness weights, not causal links; non-
%      reductionism via the fundamental theorem of empirical Aethae;
%      [O] program question: golden-causality legibility at attribute
%      level via accordance.
%   D2 gr1:rem:forwardforcing (author dictation, sharpened in-model):
%      P(viable | survival to t) -> 1 at a soft threshold set by the
%      slowest non-viable clock; on {T = infinity} outright, since
%      T=inf entails Lambda<inf a.s.  Forward mirror of the
%      Optimation climb; ceiling fence explicit.
%   VOICE: intro's older middle subsections (thesis, notclaimed)
%      rewritten to the measured register; model-in-three-lines left
%      as is (functional list genre).  Introduction now register-
%      complete end to end.  NOT instated: the author's reflexive
%      observation that this collaboration itself may be golden
%      causality -- blocked by the paper's own occupancy fences
%      (arrivaltiming); recorded here only.
% AUTHOR-CHECK: both dictation remarks await his read; the anchor-
%   segment program (his hand) remains the long game.
% =====================================================================
% =====================================================================
% AUTHOR DICTATION D3 -- Aug 23 2026: sch:rem:notexhausted.
%   The schedule channel declared LIVE under continuity (correcting
%   the prior merge-into-other-channels practice); IDENTIFICATION:
%   accordance-legible modern golden causality IS the third term
%   operating (sharpens aethicsubstrate's [O] to [S]); scoreable
%   texture added (prerequisite-shaped clustering vs uniform spread);
%   sequential-achievement class with Newton as citable exemplar
%   (Ben25Seq, verify-flagged), kept separate per the author's
%   instruction; the framework's own production entered ONCE as a
%   hypothetical candidate member, expressly non-load-bearing, with
%   the arrivaltiming fence cited in-text; retention clause wired to
%   forwardforcing (asymptotic conditioning nails remaining
%   prerequisites, stationary-action character).
% Reflexive golden causality: author confirms the PROVENANCE RECORD
%   is its home; this ledger and his public commentary carry it.
% =====================================================================
% STILL OPEN after this pass, in priority order: (a) derive the
% positional weighting from the Copernican draw (closes ra:ans:draw);
% (b) DONE Aug 21: app:numerics executed -- falsifiers did not fire;
% results recorded in the appendix and in the extraction paper; (c) the identification map (gs:sub:limit),
% which the Route World and the Score both wait on; (d) rho_t dynamics
% (gr1:subsub:fluct), the paper's stated principal exposure.
% =====================================================================

% =====================================================================
% PASS RECORD -- sheet entries 5:32 PM Aug 20 through 12:46 AM Aug 21
% (Edinburgh), instated this pass (Fable 5, Aug 21 2026):
%   1. 5:32 PM Aug 20 -- invested Golden causality + the currency
%      register: Rem. gc:rem:invested, closing the Golden Optimation
%      subsection.
%   2. 11:36 PM Aug 20 -- 12:17 AM Aug 21 -- the mattering residue /
%      Departric invalidity derivation: new subsection
%      gc:sub:matteringresidue (placed before the Route-World fences,
%      which bound it), with the complacency objection and its
%      two-part reply as Rem. gc:rem:complacency; the population-
%      robustness half of the reply carried as a premise with its own
%      falsifier.  Rendering choice, per the sheet's own query: stepped
%      derivation in the finale-paradox style, not prose, with the
%      conclusion's exact strength isolated in Step 4.  The entry's
%      second half (define the Departure Point by the third postulate's
%      two-generation statement) was found ALREADY instated at
%      Rem. gc:rem:thirdpostulate; only the provenance footnote (the
%      discipline against naive transfer; the late, deliberate
%      convergence) was added there.
%   3. 12:30 AM Aug 21 -- phenomenological death; existential first
%      principles; the continuous threat against the single-crossroads
%      picture: Rem. gc:rem:phenomdeath, directly after
%      Def. gc:def:departure.
%   4. 12:43 AM Aug 21 -- modal placement of "unavailable in
%      principle" (the ceiling binds all three worlds, so supposing
%      availability refutes the model, not the Route World; no reductio
%      is assemblable from it): paragraph added inside
%      Rem. gc:rem:routefence.
%   5. 12:46 AM Aug 21 -- the hypothesis read backward; the
%      Enlightenment prerequisite fan at candidate grade: new
%      subsection gc:sub:backward (between the empirical-signature
%      subsection and the Lindy step), fenced by the look-elsewhere bar
%      and required to pass the narrative-bias asymmetry test.
% CITATION NOTE: keys Merton38 and Eisenstein79, introduced in
% gc:sub:backward, need .bib entries (Merton, "Science, Technology and
% Society in Seventeenth-Century England", 1938; Eisenstein, "The
% Printing Press as an Agent of Change", 1979); author to verify
% editions before merging.
% =====================================================================

% =====================================================================
%  GOLDEN REASONING -- MASTER DOCUMENT
%  Sections 1-2 assembled 2026-08-01.
%
%  SOURCES.  This file is the concatenation of two standalone section
%  files, each of which carries its own full provenance and correction
%  ledger in its header:
%      section-soft-horizons.tex   -> Section 1
%      section-two-channel.tex     -> Section 2 (revision 2)
%  Those ledgers are NOT reproduced here.  Keep the standalone files:
%  they are the authoritative record of who originated what, of the
%  correction log, and of citation-verification status.  Consult them
%  before posting.
%
%  BIBLIOGRAPHY.  The two sections shared four references (AS15, GJW99,
%  VMS79, VY85); the merged list is their union, 35 entries, deduped.
%  The file is SELF-SWITCHING: if biblatex is installed it is used,
%  with the .bib written at compile time by filecontents; otherwise a
%  manual thebibliography is used instead.  Both are generated from the
%  same 35-key list and their key sets were checked equal.  To force
%  one path, edit the \IfFileExists test below.
%
%  NOTE ON VERIFICATION.  The fallback (thebibliography) path is
%  compile-verified.  The biblatex path is NOT -- biblatex and biber
%  were unavailable in the environment where this file was assembled.
%  If biber errors, the usual causes are (a) needing a biber run
%  between pdflatex runs, or (b) backend=biber vs bibtex mismatch.
%
%  CROSS-REFERENCES.  Section 2's references to "Section One" have been
%  converted to \ref{sh:sec} in this merged file.  Label prefixes are
%  disjoint: sh: for Section 1, tc: for Section 2.
% =====================================================================

\documentclass[11pt]{article}

\usepackage[margin=1.1in]{geometry}
\usepackage{amsmath,amssymb,amsthm,mathtools}
\usepackage{booktabs}
\usepackage{array}
\usepackage{enumitem}
\usepackage{xcolor}
\definecolor{grlink}{rgb}{0.10,0.20,0.55}
\usepackage{hyperref}
% >>> FORMATTING PORT (Aug 22 2026, author request): hyperlink
% appearance, PDF sidebar bookmarks, and \part styling matched to the
% Aethic paper.  grlink color retained above for any future opt-in.
\hypersetup{
colorlinks,
citecolor=black,
filecolor=black,
linkcolor=black,
urlcolor=black,
bookmarksnumbered=true,
bookmarksopen=true,
bookmarksopenlevel=1,
}
\usepackage{titlesec}
% Aethic \part format
\titleformat{\part}[block]
  {\centering\Huge\bfseries}
  {Part \thepart:}
  {5pt}
  {}
% Aethic-style per-part section reset, available but NOT applied
% (structural-pass note: numbering-stability conservatism):
\newcommand{\GRPart}[1]{\part{#1}\setcounter{section}{0}}

% ---- bibliography database (written at compile time) ----------------
\begin{filecontents*}[overwrite]{golden-reasoning.bib}
@article{AD13, author={Aurzada, Frank and Dereich, Steffen}, title={Universality of the asymptotics of the one-sided exit problem for integrated processes}, journal={Ann. Inst. Henri Poincar\'e Probab. Stat.}, volume={49}, year={2013}, pages={236--251}}
@incollection{AS15, author={Aurzada, Frank and Simon, Thomas}, title={Persistence probabilities and exponents}, booktitle={L\'evy Matters V}, series={Lecture Notes in Mathematics}, volume={2149}, publisher={Springer}, address={Cham}, year={2015}, pages={183--224}}
@article{BBM05, author={Ben Arous, G\'erard and Bogachev, Leonid V. and Molchanov, Stanislav A.}, title={Limit theorems for sums of random exponentials}, journal={Probab. Theory Relat. Fields}, volume={132}, number={4}, year={2005}, pages={579--612}}
@article{BMS13, author={Bray, Alan J. and Majumdar, Satya N. and Schehr, Gr\'egory}, title={Persistence and first-passage properties in nonequilibrium systems}, journal={Adv. Phys.}, volume={62}, number={3}, year={2013}, pages={225--361}}
@article{Bha82, author={Bhattacharya, Rabi N.}, title={On the functional central limit theorem and the law of the iterated logarithm for Markov processes}, journal={Z. Wahrscheinlichkeitstheorie verw. Gebiete}, volume={60}, year={1982}, pages={185--201}}
@book{Bov06, author={Bovier, Anton}, title={Statistical Mechanics of Disordered Systems: A Mathematical Perspective}, publisher={Cambridge University Press}, year={2006}}
@article{DW15, author={Denisov, Denis and Wachtel, Vitali}, title={Exit times for integrated random walks}, journal={Ann. Inst. Henri Poincar\'e Probab. Stat.}, volume={51}, number={1}, year={2015}, pages={167--193}}
@article{Der80, author={Derrida, Bernard}, title={Random-energy model: limit of a family of disordered models}, journal={Phys. Rev. Lett.}, volume={45}, year={1980}, pages={79--82}}
@article{Der81, author={Derrida, Bernard}, title={Random-energy model: an exactly solvable model of disordered systems}, journal={Phys. Rev. B}, volume={24}, number={5}, year={1981}, pages={2613--2626}}
@book{Dur19, author={Durrett, Rick}, title={Probability: Theory and Examples}, edition={5}, publisher={Cambridge University Press}, year={2019}}
@book{FW12, author={Freidlin, Mark I. and Wentzell, Alexander D.}, title={Random Perturbations of Dynamical Systems}, edition={3}, publisher={Springer}, year={2012}}
@article{GJW99, author={Groeneboom, Piet and Jongbloed, Geurt and Wellner, Jon A.}, title={Integrated Brownian motion, conditioned to be positive}, journal={Ann. Probab.}, volume={27}, number={3}, year={1999}, pages={1283--1303}}
@article{Gol71, author={Goldman, Malcolm}, title={On the first passage of the integrated Wiener process}, journal={Ann. Math. Statist.}, volume={42}, year={1971}, pages={2150--2155}}
@article{Gom25, author={Gompertz, Benjamin}, title={On the nature of the function expressive of the law of human mortality, and on a new mode of determining the value of life contingencies}, journal={Philos. Trans. R. Soc. Lond.}, volume={115}, year={1825}, pages={513--583}}
@article{Hor67, author={H\"ormander, Lars}, title={Hypoelliptic second order differential equations}, journal={Acta Math.}, volume={119}, year={1967}, pages={147--171}}
@article{IW94, author={Isozaki, Yasuki and Watanabe, Shinzo}, title={An asymptotic formula for the Kolmogorov diffusion and a refinement of Sinai's estimates for the integral of Brownian motion}, journal={Proc. Japan Acad. Ser. A}, volume={70}, number={9}, year={1994}, pages={271--276}}
@book{KT81, author={Karlin, Samuel and Taylor, Howard M.}, title={A Second Course in Stochastic Processes}, publisher={Academic Press}, year={1981}}
@article{Kol34, author={Kolmogorov, Andrei N.}, title={Zuf\"allige Bewegungen (zur Theorie der Brownschen Bewegung)}, journal={Ann. of Math. (2)}, volume={35}, year={1934}, pages={116--117}}
@article{Lac91, author={Lachal, Aim\'e}, title={Sur le premier instant de passage de l'int\'egrale du mouvement brownien}, journal={Ann. Inst. H. Poincar\'e Probab. Statist.}, volume={27}, year={1991}, pages={385--405}}
@book{Lan04, author={Lando, David}, title={Credit Risk Modeling: Theory and Applications}, publisher={Princeton University Press}, year={2004}}
@article{MK04, author={Molchan, George and Khokhlov, Andrei}, title={Small values of the maximum for the integral of fractional Brownian motion}, journal={J. Stat. Phys.}, volume={114}, year={2004}, pages={923--945}}
@article{MP01, author={Milevsky, Moshe A. and Promislow, S. David}, title={Mortality derivatives and the option to annuitise}, journal={Insurance Math. Econom.}, volume={29}, number={3}, year={2001}, pages={299--318}}
@article{MV12, author={M\'el\'eard, Sylvie and Villemonais, Denis}, title={Quasi-stationary distributions and population processes}, journal={Probab. Surv.}, volume={9}, year={2012}, pages={340--410}}
@article{Mar49, author={Maruyama, Gisiro}, title={The harmonic analysis of stationary stochastic processes}, journal={Mem. Fac. Sci. Ky\=ush\=u Univ. Ser. A}, volume={4}, year={1949}, pages={45--106}}
@article{McK63, author={McKean, Jr., Henry P.}, title={A winding problem for a resonator driven by a white noise}, journal={J. Math. Kyoto Univ.}, volume={2}, year={1963}, pages={227--235}}
@article{Mol99, author={Molchan, George M.}, title={Maximum of a fractional Brownian motion: probabilities of small values}, journal={Comm. Math. Phys.}, volume={205}, year={1999}, pages={97--111}}
@article{RVY06, author={Roynette, Bernard and Vallois, Pierre and Yor, Marc}, title={Some penalisations of the Wiener measure}, journal={Jpn. J. Math.}, volume={1}, year={2006}, pages={263--290}}
@book{RY99, author={Revuz, Daniel and Yor, Marc}, title={Continuous Martingales and Brownian Motion}, edition={3}, publisher={Springer}, year={1999}}
@article{Sin92, author={Sinai, Yakov G.}, title={Distribution of some functionals of the integral of a random walk}, journal={Theoret. and Math. Phys.}, volume={90}, year={1992}, pages={219--241}, note={translated from Teoret. Mat. Fiz. \textbf{90} (1992), no. 3, 323--353}}
@article{VMS79, author={Vaupel, James W. and Manton, Kenneth G. and Stallard, Eric}, title={The impact of heterogeneity in individual frailty on the dynamics of mortality}, journal={Demography}, volume={16}, year={1979}, pages={439--454}}
@article{VY85, author={Vaupel, James W. and Yashin, Anatoli I.}, title={Heterogeneity's ruses: some surprising effects of selection on population dynamics}, journal={Amer. Statist.}, volume={39}, year={1985}, pages={176--185}}
@article{Wah78, author={Wahba, Grace}, title={Improper priors, spline smoothing and the problem of guarding against model errors in regression}, journal={J. Roy. Statist. Soc. Ser. B}, volume={40}, year={1978}, pages={364--372}}
@book{Bos02, author={Bostrom, Nick}, title={Anthropic Bias: Observation Selection Effects in Science and Philosophy}, publisher={Routledge}, address={New York}, year={2002}}
@article{Car83, author={Carter, Brandon}, title={The anthropic principle and its implications for biological evolution}, journal={Philos. Trans. R. Soc. Lond. A}, volume={310}, year={1983}, pages={347--363}}
@article{Dys60, author={Dyson, Freeman J.}, title={Search for artificial stellar sources of infrared radiation}, journal={Science}, volume={131}, number={3414}, year={1960}, pages={1667--1668}}
@article{Got93, author={Gott III, J. Richard}, title={Implications of the Copernican principle for our future prospects}, journal={Nature}, volume={363}, number={6427}, year={1993}, pages={315--319}}
@misc{Han98, author={Hanson, Robin}, title={The Great Filter --- Are We Almost Past It?}, howpublished={manuscript}, year={1998}}
@article{Kar64, author={Kardashev, Nikolai S.}, title={Transmission of information by extraterrestrial civilizations}, journal={Soviet Astronomy}, volume={8}, year={1964}, pages={217--221}, note={Russian original, Astron. Zh. \textbf{41} (1964), 282--287}}
@misc{SDO18, author={Sandberg, Anders and Drexler, Eric and Ord, Toby}, title={Dissolving the Fermi paradox}, howpublished={preprint}, year={2018}}
@book{Tal12, author={Taleb, Nassim Nicholas}, title={Antifragile: Things That Gain from Disorder}, publisher={Random House}, address={New York}, year={2012}}
@book{WB00, author={Ward, Peter D. and Brownlee, Donald}, title={Rare Earth: Why Complex Life Is Uncommon in the Universe}, publisher={Copernicus}, address={New York}, year={2000}}
@article{Nexus, author={Benander, Ajax}, title={Nexic Reasoning: Defining a Generalized Calculus Over Anthropic Parameters}, year={2024}, month={December}, doi={10.31219/osf.io/jvkcp}, url={https://philpapers.org/rec/BENNRD}, note={Also at osf.io/jvkcp}}
@misc{Ben26Found, author={Benander, Ajax}, title={The Copernican Method Against the Copernican Principle: Anthropic Typicality over Partial Information States}, howpublished={Preprint}, year={2026}}
@misc{Ben26Optimal, author={Benander, Ajax}, title={The Optimal Earth: A Paradigmatic Re-Rendition of Natural History from Coupled Improbability}, howpublished={Working draft}, year={2026}}
@article{BC21, author={Balbi, Amedeo and {\'C}irkovi{\'c}, Milan M.}, title={Longevity is the key factor in the search for technosignatures}, journal={Astron. J.}, volume={161}, year={2021}, pages={222}, doi={10.3847/1538-3881/abec48}}
@article{BG24, author={Balbi, Amedeo and Grimaldi, Claudio}, title={Technosignatures longevity and Lindy's law}, journal={Astron. J.}, volume={167}, number={3}, year={2024}, pages={119}, doi={10.3847/1538-3881/ad217d}}
@misc{Kip20, author={Kipping, David and Frank, Charles and Scarlata, Claudia}, title={Contact inequality: first contact will likely be with an older civilization}, howpublished={preprint}, year={2020}, note={arXiv:2010.12358}}
@misc{Ord23, author={Ord, Toby}, title={The Lindy Effect}, howpublished={preprint}, year={2023}, note={arXiv:2308.09045}}
@article{Crary23, author={Crary, Alice}, title={The toxic ideology of longtermism}, journal={Radical Philosophy}, volume={2}, number={14}, year={2023}, pages={49--57}}
@misc{GM21, author={Greaves, Hilary and MacAskill, William}, title={The case for strong longtermism}, howpublished={Global Priorities Institute Working Paper 5-2021, University of Oxford}, year={2021}}
@book{Ord20, author={Ord, Toby}, title={The Precipice: Existential Risk and the Future of Humanity}, publisher={Bloomsbury}, address={London}, year={2020}}
@article{CSB10, author={{\'C}irkovi{\'c}, Milan M. and Sandberg, Anders and Bostrom, Nick}, title={Anthropic shadow: observation selection effects and human extinction risks}, journal={Risk Analysis}, volume={30}, number={10}, year={2010}, pages={1495--1506}}
@incollection{MacAskill22, author={MacAskill, William}, title={Are we living at the hinge of history?}, booktitle={Ethics and Existence: The Legacy of Derek Parfit}, publisher={Oxford University Press}, year={2022}, note={AUTHOR TO VERIFY edition and pages at merge}}
@misc{Ben26Method, author={Benander, Ajax}, title={[Methodology companion: discovery-side genealogy, shades of similarity, and linguistic masking]}, howpublished={Working draft}, year={2026}, note={PLACEHOLDER TITLE --- author to supply exact title at merge}}
@article{Merton38, author={Merton, Robert K.}, title={Science, technology and society in seventeenth century {E}ngland}, journal={Osiris}, volume={4}, year={1938}, pages={360--632}, note={AUTHOR TO VERIFY edition used}}
@book{Eisenstein79, author={Eisenstein, Elizabeth L.}, title={The Printing Press as an Agent of Change}, publisher={Cambridge University Press}, year={1979}, note={AUTHOR TO VERIFY edition used}}
@book{Herlihy97, author={Herlihy, David}, title={The Black Death and the Transformation of the West}, publisher={Harvard University Press}, year={1997}, note={AUTHOR TO VERIFY}}
@book{Belich22, author={Belich, James}, title={The World the Plague Made: The Black Death and the Rise of Europe}, publisher={Princeton University Press}, year={2022}, note={AUTHOR TO VERIFY}}
@book{Parfit84, author={Parfit, Derek}, title={Reasons and Persons}, publisher={Oxford University Press}, address={Oxford}, year={1984}}
@book{Schell82, author={Schell, Jonathan}, title={The Fate of the Earth}, publisher={Knopf}, address={New York}, year={1982}}
@book{Jonas84, author={Jonas, Hans}, title={The Imperative of Responsibility: In Search of an Ethics for the Technological Age}, publisher={University of Chicago Press}, address={Chicago}, year={1984}, note={German original 1979}}
@misc{AG16, author={Althaus, David and Gloor, Lukas}, title={Reducing Risks of Astronomical Suffering: A Neglected Priority}, howpublished={Center on Long-Term Risk, report}, year={2016}}
@book{Scheffler13, author={Scheffler, Samuel}, title={Death and the Afterlife}, publisher={Oxford University Press}, address={New York}, year={2013}}
@article{Bos03, author={Bostrom, Nick}, title={Astronomical waste: the opportunity cost of delayed technological development}, journal={Utilitas}, volume={15}, number={3}, year={2003}, pages={308--314}}
@misc{TRDNB14, author={Taleb, Nassim Nicholas and Read, Rupert and Douady, Rapha{\"e}l and Norman, Joseph and Bar-Yam, Yaneer}, title={The precautionary principle (with application to the genetic modification of organisms)}, howpublished={preprint}, year={2014}, note={arXiv:1410.5787}}
@article{Peters19, author={Peters, Ole}, title={The ergodicity problem in economics}, journal={Nature Physics}, volume={15}, year={2019}, pages={1216--1221}}
@book{Leslie96, author={Leslie, John}, title={The End of the World: The Science and Ethics of Human Extinction}, publisher={Routledge}, address={London}, year={1996}}
@article{HMMP21, author={Hanson, Robin and Martin, Daniel and McCarter, Calvin and Paulson, Jonathan}, title={If loud aliens explain human earliness, quiet aliens are also rare}, journal={Astrophys. J.}, volume={922}, year={2021}, pages={182}}
@article{SBSD21, author={Snyder-Beattie, Andrew E. and Sandberg, Anders and Drexler, K. Eric and Bonsall, Michael B.}, title={The timing of evolutionary transitions suggests intelligent life is rare}, journal={Astrobiology}, volume={21}, number={3}, year={2021}, pages={265--278}}
@article{MB98, author={Majumdar, Satya N. and Bray, Alan J.}, title={Persistence with partial survival}, journal={Phys. Rev. Lett.}, volume={81}, year={1998}, pages={2626--2629}}
@incollection{Bur14, author={Burkhardt, Theodore W.}, title={First passage of a randomly accelerated particle}, booktitle={First-Passage Phenomena and Their Applications}, publisher={World Scientific}, year={2014}, pages={21--44}}
@article{GG01, author={Gavrilov, Leonid A. and Gavrilova, Natalia S.}, title={The reliability theory of aging and longevity}, journal={J. Theoret. Biol.}, volume={213}, year={2001}, pages={527--545}}
% [Referee pass Aug 22: four future-conditioned-tradition keys + the
%  extraction key, AUTHOR TO VERIFY editions/translations.]
@book{BarrowTipler86, author={Barrow, John D. and Tipler, Frank J.}, title={The Anthropic Cosmological Principle}, publisher={Oxford University Press}, address={Oxford}, year={1986}}
@book{Tipler94, author={Tipler, Frank J.}, title={The Physics of Immortality: Modern Cosmology, God and the Resurrection of the Dead}, publisher={Doubleday}, address={New York}, year={1994}}
@book{Leslie79, author={Leslie, John}, title={Value and Existence}, publisher={Basil Blackwell}, address={Oxford}, year={1979}}
@book{Teilhard55, author={Teilhard de Chardin, Pierre}, title={The Phenomenon of Man}, publisher={Harper}, address={New York}, year={1959}, note={French original 1955; trans.\ B.~Wall}}
@misc{Ben26Extract, author={Benander, Ajax K.}, title={Survivorship structure of the two-channel muffled hazard}, howpublished={Companion extraction manuscript}, year={2026}}
@misc{Ben25Seq, author={Benander, Ajax K.}, title={Sequential achievement}, howpublished={Companion manuscript; AUTHOR TO VERIFY title and year}, year={2025}}
@article{SO94, author={Sagan, Carl and Ostro, Steven J.}, title={Dangers of asteroid deflection}, journal={Nature}, volume={368}, year={1994}, pages={501}}
@book{Hag19, author={H{\"a}gglund, Martin}, title={This Life: Secular Faith and Spiritual Freedom}, publisher={Pantheon}, address={New York}, year={2019}}
@incollection{Wil73, author={Williams, Bernard}, title={The Makropulos case: reflections on the tedium of immortality}, booktitle={Problems of the Self}, publisher={Cambridge University Press}, year={1973}, pages={82--100}}
@article{AS04, author={Albrecht, Andreas and Sorbo, Lorenzo}, title={Can the universe afford inflation?}, journal={Phys. Rev. D}, volume={70}, year={2004}, pages={063528}, note={AUTHOR TO VERIFY}}
@misc{Car17, author={Carroll, Sean M.}, title={Why {B}oltzmann brains are bad}, howpublished={preprint}, year={2017}, note={arXiv:1702.00850; AUTHOR TO VERIFY venue}}
@book{MacIntyre81, author={MacIntyre, Alasdair}, title={After Virtue: A Study in Moral Theory}, publisher={University of Notre Dame Press}, year={1981}, note={AUTHOR TO VERIFY edition}}
@article{Anscombe58, author={Anscombe, G. E. M.}, title={Modern moral philosophy}, journal={Philosophy}, volume={33}, number={124}, year={1958}, pages={1--19}}
@article{Chr07, author={Christensen, David}, title={Epistemology of disagreement: the good news}, journal={Philos. Rev.}, volume={116}, number={2}, year={2007}, pages={187--217}}
@book{nietzsche1974gay, author={Nietzsche, Friedrich}, title={The Gay Science}, publisher={Vintage}, address={New York}, year={1974}, note={trans. W. Kaufmann; orig. 1882. AUTHOR TO VERIFY edition (key matched to draft Introduction)}}
@book{weinberg1977first, author={Weinberg, Steven}, title={The First Three Minutes}, publisher={Basic Books}, address={New York}, year={1977}, note={key matched to draft Introduction}}
@book{Camus42, author={Camus, Albert}, title={The Myth of Sisyphus}, publisher={Gallimard}, address={Paris}, year={1942}, note={French original \emph{Le Mythe de Sisyphe}, 1942; AUTHOR TO VERIFY translation edition}}
@misc{Sartre46, author={Sartre, Jean-Paul}, title={Existentialism Is a Humanism}, howpublished={lecture, pub. 1946}, note={AUTHOR TO VERIFY edition}}
@book{Carroll16, author={Carroll, Sean}, title={The Big Picture: On the Origins of Life, Meaning, and the Universe Itself}, publisher={Dutton}, year={2016}}
@book{Midgley85, author={Midgley, Mary}, title={Evolution as a Religion: Strange Hopes and Stranger Fears}, publisher={Methuen}, address={London}, year={1985}}
@book{Kie1846, author={Kierkegaard, S{\o}ren}, title={Concluding Unscientific Postscript}, year={1846}, publisher={[edition TBD]}, note={AUTHOR TO VERIFY edition/translation}}
@book{Hei27, author={Heidegger, Martin}, title={Being and Time}, year={1927}, publisher={[edition TBD]}, note={AUTHOR TO VERIFY edition/translation}}
@book{Frankl59, author={Frankl, Viktor E.}, title={Man's Search for Meaning}, publisher={Beacon}, address={Boston}, year={1959}, note={orig. 1946; AUTHOR TO VERIFY edition}}
@article{Nagel71, author={Nagel, Thomas}, title={The absurd}, journal={J. Philos.}, volume={68}, number={20}, year={1971}, pages={716--727}}
@book{Wolf10, author={Wolf, Susan}, title={Meaning in Life and Why It Matters}, publisher={Princeton University Press}, year={2010}}
@book{Metz13, author={Metz, Thaddeus}, title={Meaning in Life: An Analytic Study}, publisher={Oxford University Press}, year={2013}}
@article{Gallup70, author={Gallup, Gordon G.}, title={Chimpanzees: self-recognition}, journal={Science}, volume={167}, year={1970}, pages={86--87}}
@misc{MarxEngels48, author={Marx, Karl and Engels, Friedrich}, title={Manifesto of the Communist Party}, year={1848}, note={AUTHOR TO VERIFY edition}}
@article{Horowitz17, author={Horowitz, Alexandra}, title={Smelling themselves: dogs investigate their own odours longer when modified in an ``olfactory mirror'' test}, journal={Behavioural Processes}, volume={143}, year={2017}, pages={17--24}, note={AUTHOR TO VERIFY}}
@misc{Torres21, author={Torres, {\'E}mile P.}, title={Against longtermism}, howpublished={Aeon}, year={2021}}
@article{Bostrom19, author={Bostrom, Nick}, title={The vulnerable world hypothesis}, journal={Global Policy}, volume={10}, number={4}, year={2019}, pages={455--476}, doi={10.1111/1758-5899.12718}}
@book{Tainter88, author={Tainter, Joseph A.}, title={The Collapse of Complex Societies}, publisher={Cambridge University Press}, year={1988}}
@article{EF24, author={Engler, John-Oliver and Fischer, Jan Niklas}, title={Drawing blanks and winning: quantifying global catastrophic risk associated with human ingenuity}, journal={Sustainable Futures}, volume={7}, year={2024}, pages={100165}}
@book{Cooper19, author={Cooper, Keith}, title={The Contact Paradox: Challenging Our Assumptions in the Search for Extraterrestrial Intelligence}, publisher={Bloomsbury Sigma}, address={London}, year={2019}}
@incollection{Plutarch, author={Plutarch}, title={Whether Land or Sea Animals Are Cleverer ({D}e sollertia animalium)}, booktitle={Moralia}, note={\S7; the saying is attributed there to Bion of Borysthenes}}
@article{Hei46, author={Heider, Fritz}, title={Attitudes and cognitive organization}, journal={J. Psychol.}, volume={21}, year={1946}, pages={107--112}}
@article{CH56, author={Cartwright, Dorwin and Harary, Frank}, title={Structural balance: a generalization of Heider's theory}, journal={Psychol. Rev.}, volume={63}, year={1956}, pages={277--293}}
@article{MKKS11, author={Marvel, Seth A. and Kleinberg, Jon and Kleinberg, Robert D. and Strogatz, Steven H.}, title={Continuous-time model of structural balance}, journal={Proc. Natl. Acad. Sci. USA}, volume={108}, year={2011}, pages={1771--1776}}
@article{AKR05, author={Antal, Tibor and Krapivsky, Paul L. and Redner, Sidney}, title={Dynamics of social balance on networks}, journal={Phys. Rev. E}, volume={72}, year={2005}, pages={036121}}
@book{Wal79, author={Waltz, Kenneth N.}, title={Theory of International Politics}, publisher={McGraw-Hill}, year={1979}}
@book{Wal87, author={Walt, Stephen M.}, title={The Origins of Alliances}, publisher={Cornell University Press}, year={1987}}
@article{Aethus, author={Benander, Ajax}, title={Aethic Reasoning: A Comprehensive Solution to the Quantum Measurement Problem}, year={2024}, month={November}, url={https://philpapers.org/rec/BENARA-5}, note={PhilSci-Archive eprint 25713}}
@misc{WeightAlg, author={Benander, Ajax}, title={The Weight Algebra of Information States: A Two-Layer Commutative Semiring for Superposition and Pruning}, howpublished={Working draft}, year={2026}, month={July}}
@misc{AethusBrief, author={Benander, Ajax}, title={Aethic Reasoning in Brief: From a Counterfactual Semantics to the Quantum Measurement Problem}, howpublished={Preprint}, year={2026}, url={https://philpapers.org/rec/BENARI-3}, note={The compressed companion; PhilPapers record BENARI-3}}
@misc{GoldenScience22, author={Benander, Ajax}, title={Golden Science}, howpublished={Manuscript}, year={2022}, url={https://aethus.org/aethic-code/files/GoldenScience_Nov2022.pdf}, note={Dated August 2022; archived copy November 2022}}
@book{kripke1980naming, title={Naming and Necessity}, author={Kripke, Saul A.}, year={1980}, publisher={Harvard University Press}, address={Cambridge, Massachusetts}}
@book{Pomeranz00, author={Pomeranz, Kenneth}, title={The Great Divergence: China, Europe, and the Making of the Modern World Economy}, publisher={Princeton University Press}, year={2000}}
@book{Allen09, author={Allen, Robert C.}, title={The British Industrial Revolution in Global Perspective}, publisher={Cambridge University Press}, year={2009}}
@book{Diamond97, author={Diamond, Jared}, title={Guns, Germs, and Steel: The Fates of Human Societies}, publisher={W. W. Norton}, year={1997}}
@book{Arrian, author={Arrian}, title={Anabasis of Alexander}, note={1.15.8}}
@incollection{PlutarchAlex, author={Plutarch}, title={Life of Alexander}, booktitle={Parallel Lives}, note={\S16}}
@book{Wittgenstein69, author={Wittgenstein, Ludwig}, title={On Certainty}, editor={Anscombe, G. E. M. and von Wright, G. H.}, publisher={Blackwell}, address={Oxford}, year={1969}}
@book{Sextus, author={Sextus Empiricus}, title={Outlines of Pyrrhonism}, note={The five modes of Agrippa, {PH} I.164--177}, year={c.\,200}}
@book{Mackie77, author={Mackie, J. L.}, title={Ethics: Inventing Right and Wrong}, publisher={Penguin}, address={Harmondsworth}, year={1977}}
@misc{AethusDecision, author={Benander, Ajax}, title={Newcomb's Problem Without a Fixed Past: A Decision Rule From Observer-Indexed Determinacy}, howpublished={Working draft}, year={2026}}
@article{Barbour1995, author={Barbour, Julian B.}, title={General Relativity as a Theory of Two Fields on Absolute Space}, journal={General Relativity and Gravitation}, volume={27}, pages={1-24}, year={1995}}
@misc{NexicFoundation, author={Benander, Ajax}, title={The Accordance Principle: Foundations of Nexic Reasoning}, howpublished={Working draft}, year={2026}}
@book{Benatar06, author={Benatar, David}, title={Better Never to Have Been: The Harm of Coming into Existence}, publisher={Oxford University Press}, address={Oxford}, year={2006}}
@misc{VHEMT, author={{Voluntary Human Extinction Movement}}, title={The Voluntary Human Extinction Movement}, howpublished={\url{https://www.vhemt.org}}, note={Accessed 2026}}
@article{ALICE2012Temperature, author={{ALICE Collaboration}}, title={Enhanced production of multi-strange hadrons in high-multiplicity proton--proton collisions}, journal={Nature Physics}, volume={13}, pages={535--539}, year={2017}, note={The LHC's quark--gluon plasma at $\sim5$ trillion K, the hottest temperature yet produced}}
@article{Kovachy2015Picokelvin, author={Kovachy, Tim and Hogan, Jason M. and Sugarbaker, Alex and Dickerson, Susannah M. and Donnelly, Christine A. and Overstreet, Chris and Kasevich, Mark A.}, title={Matter wave lensing to picokelvin temperatures}, journal={Physical Review Letters}, volume={114}, pages={143004}, year={2015}}
@book{burgess2007animals, author={Burgess, Colin and Dubbs, Chris}, title={Animals in Space: From Research Rockets to the Space Shuttle}, publisher={Springer--Praxis}, address={Chichester}, year={2007}}
@misc{Pattali2026_CRISPR_epigenome_mod, author={Pattali, Rijul and others}, title={Programmable epigenome editing for reactivation of silenced genes}, howpublished={Preprint}, year={2026}, note={Cited for the general capability; the specific reference is to be confirmed at deposit}}
@inproceedings{Brannen2006TheLM,
  title={The Lepton Masses},
  author={Carl Andrew Brannen},
  year={2006},
  url={https://api.semanticscholar.org/CorpusID:17711106}
}
@article{koide1983new,
  author    = {Koide, Yoshio},
  title     = {New view of quark and lepton mass hierarchy},
  journal   = {Physical Review D},
  volume    = {28},
  number    = {1},
  pages     = {252--254},
  year      = {1983},
  month     = {Jul},
  publisher = {American Physical Society},
  doi       = {10.1103/PhysRevD.28.252}
}
\end{filecontents*}
\begin{filecontents*}[overwrite]{gr-manualbib.tex}
\begin{thebibliography}{99}\small
\bibitem{Allen09} R.~C.~Allen, \emph{The British Industrial
Revolution in Global Perspective}, Cambridge University Press, 2009.

\bibitem{AKR05} T.~Antal, P.~L.~Krapivsky, and S.~Redner, \emph{Dynamics of social balance on networks}, Phys.\ Rev.\ E \textbf{72} (2005), 036121.

\bibitem{Arrian} Arrian, \emph{Anabasis of Alexander}, 1.15.8.

\bibitem{AD13} F.~Aurzada and S.~Dereich,
\emph{Universality of the asymptotics of the one-sided exit problem for
integrated processes}, Ann.\ Inst.\ Henri Poincar\'e Probab.\ Stat.\
\textbf{49} (2013), 236--251.

\bibitem{AS15} F.~Aurzada and T.~Simon,
\emph{Persistence probabilities and exponents}, in: L\'evy Matters V,
Lecture Notes in Mathematics 2149, Springer, Cham, 2015, pp.~183--224.

\bibitem{BC21} A.~Balbi and M.~M.~\'Cirkovi\'c, \emph{Longevity is
the key factor in the search for technosignatures}, Astron.\ J.\
\textbf{161} (2021), 222.

\bibitem{BG24} A.~Balbi and C.~Grimaldi, \emph{Technosignatures
longevity and Lindy's law}, Astron.\ J.\ \textbf{167} (2024), no.~3,
119.

\bibitem{Barbour1995} J.~B.~Barbour, \emph{General Relativity as a
Theory of Two Fields on Absolute Space}, Gen.\ Rel.\ Grav.\ \textbf{27}
(1995), 1--24.

\bibitem{BBM05} G.~Ben~Arous, L.~V.~Bogachev, and
S.~A.~Molchanov, \emph{Limit theorems for sums of random exponentials},
Probab.\ Theory Relat.\ Fields \textbf{132} (2005), no.~4, 579--612.

\bibitem{Aethus} A.~Benander, \emph{Aethic Reasoning: A
Comprehensive Solution to the Quantum Measurement Problem}, November 2024;
PhilPapers BENARA-5, PhilSci-Archive eprint 25713.

\bibitem{AethusBrief} A.~Benander, \emph{Aethic Reasoning in
Brief: From a Counterfactual Semantics to the Quantum Measurement
Problem}, preprint, 2026, PhilSci-Archive eprint 30439.

\bibitem{AethusDecision} A.~Benander, \emph{Newcomb's Problem Without a
Fixed Past: A Decision Rule From Observer-Indexed Determinacy}, working
draft, 2026.

\bibitem{Ben26Found} A.~Benander, \emph{The Copernican Method
Against the Copernican Principle: Anthropic Typicality over Partial
Information States}, preprint, 2026.

\bibitem{Ben26Optimal} A.~Benander, \emph{The Optimal Earth: A
Paradigmatic Re-Rendition of Natural History from Coupled
Improbability}, working draft, 2026.

\bibitem{GoldenScience22} A.~Benander, \emph{Golden Science},
manuscript, August 2022; archived copy dated November 2022.
\newblock \url{https://aethus.org/aethic-code/files/GoldenScience_Nov2022.pdf}

\bibitem{NexicFoundation} A.~Benander, \emph{The Accordance Principle:
Foundations of Nexic Reasoning}, working draft, 2026.

\bibitem{Nexus} A.~Benander, \emph{Nexic Reasoning: Defining
a Generalized Calculus Over Anthropic Parameters}, 2024,
\texttt{doi:10.31219/osf.io/jvkcp}.

\bibitem{WeightAlg} A.~Benander, \emph{The Weight Algebra of
Information States: A Two-Layer Commutative Semiring for Superposition
and Pruning}, working draft, July 2026.

\bibitem{Benatar06} D.~Benatar, \emph{Better Never to Have Been:
The Harm of Coming into Existence}, Oxford University Press, 2006.

\bibitem{Bha82} R.~N.~Bhattacharya,
\emph{On the functional central limit theorem and the law of the iterated
logarithm for Markov processes},
Z.~Wahrscheinlichkeitstheorie verw.\ Gebiete \textbf{60} (1982), 185--201.

\bibitem{Bos02} N.~Bostrom, \emph{Anthropic Bias: Observation
Selection Effects in Science and Philosophy}, Routledge, New York, 2002.

\bibitem{Bostrom19} N.~Bostrom, \emph{The vulnerable world
hypothesis}, Global Policy \textbf{10} (2019), no.~4, 455--476.

\bibitem{Bov06} A.~Bovier,
\emph{Statistical Mechanics of Disordered Systems: A Mathematical
Perspective}, Cambridge University Press, 2006.

\bibitem{BMS13} A.~J.~Bray, S.~N.~Majumdar, and G.~Schehr,
\emph{Persistence and first-passage properties in nonequilibrium
systems}, Adv.\ Phys.\ \textbf{62} (2013), 225--361.

\bibitem{burgess2007animals} C.~Burgess and C.~Dubbs, \emph{Animals in
Space: From Research Rockets to the Space Shuttle}, Springer--Praxis,
Chichester, 2007.

\bibitem{Car83} B.~Carter, \emph{The anthropic principle and its
implications for biological evolution}, Philos.\ Trans.\ R.\ Soc.\ Lond.\
A \textbf{310} (1983), 347--363.

\bibitem{CH56} D.~Cartwright and F.~Harary, \emph{Structural balance: a generalization of Heider's theory}, Psychol.\ Rev.\ \textbf{63} (1956), 277--293.

\bibitem{ALICE2012Temperature} ALICE Collaboration, \emph{Enhanced
production of multi-strange hadrons in high-multiplicity proton--proton
collisions}, Nature Phys.\ \textbf{13} (2017), 535--539.

\bibitem{Cooper19} K.~Cooper, \emph{The Contact Paradox:
Challenging Our Assumptions in the Search for Extraterrestrial
Intelligence}, Bloomsbury Sigma, London, 2019.

\bibitem{Crary23} A.~Crary, \emph{The toxic ideology of
longtermism}, Radical Philosophy \textbf{2} (2023), no.~14, 49--57.

\bibitem{DW15} D.~Denisov and V.~Wachtel,
\emph{Exit times for integrated random walks}, Ann.\ Inst.\ Henri
Poincar\'e Probab.\ Stat.\ \textbf{51} (2015), 167--193.

\bibitem{Der80} B.~Derrida,
\emph{Random-energy model: limit of a family of disordered models},
Phys.\ Rev.\ Lett.\ \textbf{45} (1980), 79--82.

\bibitem{Der81} B.~Derrida,
\emph{Random-energy model: an exactly solvable model of disordered
systems}, Phys.\ Rev.\ B \textbf{24} (1981), no.~5, 2613--2626.

\bibitem{Diamond97} J.~Diamond, \emph{Guns, Germs, and Steel: The
Fates of Human Societies}, W.~W.~Norton, 1997.

\bibitem{Dur19} R.~Durrett,
\emph{Probability: Theory and Examples}, 5th ed., Cambridge University
Press, 2019.

\bibitem{Dys60} F.~J.~Dyson, \emph{Search for artificial stellar
sources of infrared radiation}, Science \textbf{131} (1960), no.~3414,
1667--1668.

\bibitem{Sextus} Sextus Empiricus, \emph{Outlines of Pyrrhonism};
the five modes of Agrippa, \emph{PH} I.164--177.

\bibitem{EF24} J.-O.~Engler and J.~N.~Fischer, \emph{Drawing
blanks and winning: quantifying global catastrophic risk associated with
human ingenuity}, Sustainable Futures \textbf{7} (2024), 100165.

\bibitem{FW12} M.~I.~Freidlin and A.~D.~Wentzell,
\emph{Random Perturbations of Dynamical Systems}, 3rd ed., Springer, 2012.

\bibitem{Gol71} M.~Goldman,
\emph{On the first passage of the integrated Wiener process},
Ann.\ Math.\ Statist.\ \textbf{42} (1971), 2150--2155.

\bibitem{Gom25} B.~Gompertz,
\emph{On the nature of the function expressive of the law of human
mortality, and on a new mode of determining the value of life
contingencies}, Philos.\ Trans.\ R.\ Soc.\ Lond.\ \textbf{115} (1825),
513--583.

\bibitem{Got93} J.~R.~Gott~III, \emph{Implications of the
Copernican principle for our future prospects}, Nature \textbf{363}
(1993), no.~6427, 315--319.

\bibitem{GM21} H.~Greaves and W.~MacAskill, \emph{The case for
strong longtermism}, Global Priorities Institute Working Paper 5-2021,
University of Oxford, 2021.

\bibitem{GJW99} P.~Groeneboom, G.~Jongbloed, and J.~A.~Wellner,
\emph{Integrated Brownian motion, conditioned to be positive},
Ann.\ Probab.\ \textbf{27} (1999), 1283--1303.

\bibitem{Han98} R.~Hanson, \emph{The Great Filter --- Are We
Almost Past It?}, manuscript, 1998.

\bibitem{Hei46} F.~Heider, \emph{Attitudes and cognitive organization}, J.~Psychol.\ \textbf{21} (1946), 107--112.

\bibitem{Hor67} L.~H\"ormander,
\emph{Hypoelliptic second order differential equations},
Acta Math.\ \textbf{119} (1967), 147--171.

\bibitem{IW94} Y.~Isozaki and S.~Watanabe,
\emph{An asymptotic formula for the Kolmogorov diffusion and a refinement
of Sinai's estimates for the integral of Brownian motion},
Proc.\ Japan Acad.\ Ser.\ A \textbf{70} (1994), 271--276.

\bibitem{Kar64} N.~S.~Kardashev, \emph{Transmission of
information by extraterrestrial civilizations}, Soviet Astronomy
\textbf{8} (1964), 217--221; Russian original, Astron.\ Zh.\
\textbf{41} (1964), 282--287.

\bibitem{KT81} S.~Karlin and H.~M.~Taylor,
\emph{A Second Course in Stochastic Processes}, Academic Press, 1981.

\bibitem{Kip20} D.~Kipping, C.~Frank, and C.~Scarlata,
\emph{Contact inequality: first contact will likely be with an older
civilization}, preprint, 2020, \texttt{arXiv:2010.12358}.

\bibitem{Kol34} A.~N.~Kolmogorov,
\emph{Zuf\"allige Bewegungen (zur Theorie der Brownschen Bewegung)},
Ann.\ of Math.\ (2) \textbf{35} (1934), 116--117.

\bibitem{Kovachy2015Picokelvin} T.~Kovachy et al., \emph{Matter wave
lensing to picokelvin temperatures}, Phys.\ Rev.\ Lett.\ \textbf{114}
(2015), 143004.

\bibitem{kripke1980naming} S.~A.~Kripke, \emph{Naming and
Necessity}, Harvard University Press, Cambridge, Massachusetts, 1980.

\bibitem{Lac91} A.~Lachal,
\emph{Sur le premier instant de passage de l'int\'egrale du mouvement
brownien}, Ann.\ Inst.\ H.\ Poincar\'e Probab.\ Statist.\ \textbf{27}
(1991), 385--405.

\bibitem{Lan04} D.~Lando,
\emph{Credit Risk Modeling: Theory and Applications}, Princeton
University Press, 2004.

\bibitem{Mackie77} J.~L.~Mackie, \emph{Ethics: Inventing Right and
Wrong}, Penguin, Harmondsworth, 1977.

\bibitem{Mar49} G.~Maruyama,
\emph{The harmonic analysis of stationary stochastic processes},
Mem.\ Fac.\ Sci.\ Ky\=ush\=u Univ.\ Ser.\ A \textbf{4} (1949), 45--106.

\bibitem{MKKS11} S.~A.~Marvel, J.~Kleinberg, R.~D.~Kleinberg, and S.~H.~Strogatz, \emph{Continuous-time model of structural balance}, Proc.\ Natl.\ Acad.\ Sci.\ USA \textbf{108} (2011), 1771--1776.

\bibitem{McK63} H.~P.~McKean, Jr.,
\emph{A winding problem for a resonator driven by a white noise},
J.\ Math.\ Kyoto Univ.\ \textbf{2} (1963), 227--235.

\bibitem{MV12} S.~M\'el\'eard and D.~Villemonais,
\emph{Quasi-stationary distributions and population processes},
Probab.\ Surv.\ \textbf{9} (2012), 340--410.

\bibitem{MP01} M.~A.~Milevsky and S.~D.~Promislow,
\emph{Mortality derivatives and the option to annuitise},
Insurance Math.\ Econom.\ \textbf{29} (2001), 299--318.

\bibitem{MK04} G.~Molchan and A.~Khokhlov,
\emph{Small values of the maximum for the integral of fractional
Brownian motion}, J.\ Stat.\ Phys.\ \textbf{114} (2004), 923--945.

\bibitem{Mol99} G.~M.~Molchan,
\emph{Maximum of a fractional Brownian motion: probabilities of small
values}, Comm.\ Math.\ Phys.\ \textbf{205} (1999), 97--111.

\bibitem{VHEMT} L.~U.~Knight, \emph{The Voluntary Human
Extinction Movement}, founded 1991; \texttt{vhemt.org}.

\bibitem{Ord23} T.~Ord, \emph{The Lindy effect}, preprint, 2023,
\texttt{arXiv:2308.09045}.

\bibitem{Pattali2026_CRISPR_epigenome_mod} R.~Pattali et al.,
\emph{Programmable epigenome editing for reactivation of silenced genes},
preprint, 2026 (specific reference to be confirmed at deposit).

\bibitem{Plutarch} Plutarch, \emph{Whether Land or Sea Animals Are
Cleverer} (\emph{De sollertia animalium}), \S7, in \emph{Moralia};
the saying is attributed there to Bion of Borysthenes.

\bibitem{PlutarchAlex} Plutarch, \emph{Life of Alexander},
\S16, in \emph{Parallel Lives}.

\bibitem{Pomeranz00} K.~Pomeranz, \emph{The Great Divergence:
China, Europe, and the Making of the Modern World Economy}, Princeton
University Press, 2000.

\bibitem{RY99} D.~Revuz and M.~Yor,
\emph{Continuous Martingales and Brownian Motion}, 3rd ed., Springer,
1999.

\bibitem{RVY06} B.~Roynette, P.~Vallois, and M.~Yor,
\emph{Some penalisations of the Wiener measure}, Jpn.\ J.\ Math.\
\textbf{1} (2006), 263--290.

\bibitem{SDO18} A.~Sandberg, E.~Drexler, and T.~Ord,
\emph{Dissolving the Fermi paradox}, preprint, 2018.

\bibitem{Sin92} Ya.~G.~Sinai,
\emph{Distribution of some functionals of the integral of a random walk},
Theoret.\ and Math.\ Phys.\ \textbf{90} (1992), 219--241; translated
from Teoret.\ Mat.\ Fiz.\ \textbf{90} (1992), no.~3, 323--353.

\bibitem{Tainter88} J.~A.~Tainter, \emph{The Collapse of Complex
Societies}, Cambridge University Press, 1988.

\bibitem{Tal12} N.~N.~Taleb, \emph{Antifragile: Things That Gain
from Disorder}, Random House, New York, 2012.

\bibitem{Torres21} \'E.~P.~Torres, \emph{Against longtermism},
Aeon, 2021.

\bibitem{VMS79} J.~W.~Vaupel, K.~G.~Manton, and E.~Stallard,
\emph{The impact of heterogeneity in individual frailty on the dynamics
of mortality}, Demography \textbf{16} (1979), 439--454.

\bibitem{VY85} J.~W.~Vaupel and A.~I.~Yashin,
\emph{Heterogeneity's ruses: some surprising effects of selection on
population dynamics}, Amer.\ Statist.\ \textbf{39} (1985), 176--185.

\bibitem{Wah78} G.~Wahba,
\emph{Improper priors, spline smoothing and the problem of guarding
against model errors in regression},
J.\ Roy.\ Statist.\ Soc.\ Ser.\ B \textbf{40} (1978), 364--372.

\bibitem{Wal87} S.~M.~Walt, \emph{The Origins of Alliances}, Cornell University Press, 1987.

\bibitem{Wal79} K.~N.~Waltz, \emph{Theory of International Politics}, McGraw-Hill, 1979.

\bibitem{WB00} P.~D.~Ward and D.~Brownlee, \emph{Rare Earth: Why
Complex Life Is Uncommon in the Universe}, Copernicus, New York, 2000.

\bibitem{Wittgenstein69} L.~Wittgenstein, \emph{On Certainty},
G.E.M.~Anscombe and G.H.~von Wright (eds.), Blackwell, Oxford, 1969.

\bibitem{koide1983new} Y.~Koide, \emph{New view of quark and lepton mass
hierarchy}, Phys.\ Rev.\ D \textbf{28} (1983), 252--254.

\bibitem{Brannen2006TheLM} C.~A.~Brannen, \emph{The lepton masses},
preprint, 2006.


% --- PLACEMENT PASS ADDITIONS (Aug 21 2026): entered as a block; TO BE
% ALPHABETIZED INTO THE LIST AT MERGE.  Every entry AUTHOR-VERIFY.
% (Includes the four keys from the Precipice pass, previously present
% only in the @-entry block.)
\bibitem{Ord20} T.~Ord, \emph{The Precipice: Existential Risk and the
Future of Humanity}, Bloomsbury, London, 2020.
\bibitem{CSB10} M.~M.~{\'C}irkovi{\'c}, A.~Sandberg, and N.~Bostrom,
\emph{Anthropic shadow: observation selection effects and human
extinction risks}, Risk Analysis \textbf{30} (2010), 1495--1506.
\bibitem{MacAskill22} W.~MacAskill, \emph{Are we living at the hinge of
history?}, in Ethics and Existence: The Legacy of Derek Parfit, Oxford
University Press, 2022.
\bibitem{Ben26Method} A.~Benander, methodology companion (title to be
supplied by the author at merge), working draft, 2026.
\bibitem{Herlihy97} D.~Herlihy, \emph{The Black Death and the
Transformation of the West}, Harvard University Press, 1997.
\bibitem{Belich22} J.~Belich, \emph{The World the Plague Made: The
Black Death and the Rise of Europe}, Princeton University Press, 2022.
\bibitem{Merton38} R.~K.~Merton, \emph{Science, technology and society
in seventeenth century England}, Osiris \textbf{4} (1938), 360--632.
\bibitem{Eisenstein79} E.~L.~Eisenstein, \emph{The Printing Press as an
Agent of Change}, Cambridge University Press, 1979.
\bibitem{Parfit84} D.~Parfit, \emph{Reasons and Persons}, Oxford
University Press, Oxford, 1984.
\bibitem{Schell82} J.~Schell, \emph{The Fate of the Earth}, Knopf, New
York, 1982.
\bibitem{Jonas84} H.~Jonas, \emph{The Imperative of Responsibility},
University of Chicago Press, Chicago, 1984.
\bibitem{AG16} D.~Althaus and L.~Gloor, \emph{Reducing Risks of
Astronomical Suffering: A Neglected Priority}, Center on Long-Term
Risk report, 2016.
\bibitem{Scheffler13} S.~Scheffler, \emph{Death and the Afterlife},
Oxford University Press, New York, 2013.
\bibitem{Bos03} N.~Bostrom, \emph{Astronomical waste: the opportunity
cost of delayed technological development}, Utilitas \textbf{15}
(2003), 308--314.
\bibitem{TRDNB14} N.~N.~Taleb, R.~Read, R.~Douady, J.~Norman, and
Y.~Bar-Yam, \emph{The precautionary principle (with application to the
genetic modification of organisms)}, preprint arXiv:1410.5787, 2014.
\bibitem{Peters19} O.~Peters, \emph{The ergodicity problem in
economics}, Nature Physics \textbf{15} (2019), 1216--1221.
\bibitem{Leslie96} J.~Leslie, \emph{The End of the World: The Science
and Ethics of Human Extinction}, Routledge, London, 1996.
\bibitem{HMMP21} R.~Hanson, D.~Martin, C.~McCarter, and J.~Paulson,
\emph{If loud aliens explain human earliness, quiet aliens are also
rare}, Astrophys.\ J. \textbf{922} (2021), 182.
\bibitem{SBSD21} A.~E.~Snyder-Beattie, A.~Sandberg, K.~E.~Drexler, and
M.~B.~Bonsall, \emph{The timing of evolutionary transitions suggests
intelligent life is rare}, Astrobiology \textbf{21} (2021), 265--278.
\bibitem{MB98} S.~N.~Majumdar and A.~J.~Bray, \emph{Persistence with
partial survival}, Phys.\ Rev.\ Lett. \textbf{81} (1998), 2626--2629.
\bibitem{Bur14} T.~W.~Burkhardt, \emph{First passage of a randomly
accelerated particle}, in First-Passage Phenomena and Their
Applications, World Scientific, 2014, 21--44.
\bibitem{GG01} L.~A.~Gavrilov and N.~S.~Gavrilova, \emph{The
reliability theory of aging and longevity}, J.\ Theoret.\ Biol.\
\textbf{213} (2001), 527--545.
\bibitem{BarrowTipler86} J.~D.~Barrow and F.~J.~Tipler, \emph{The
Anthropic Cosmological Principle}, Oxford University Press, Oxford,
1986.
\bibitem{Tipler94} F.~J.~Tipler, \emph{The Physics of Immortality},
Doubleday, New York, 1994.
\bibitem{Leslie79} J.~Leslie, \emph{Value and Existence}, Basil
Blackwell, Oxford, 1979.
\bibitem{Teilhard55} P.~Teilhard de Chardin, \emph{The Phenomenon of
Man}, Harper, New York, 1959 (French original 1955).
\bibitem{Ben26Extract} A.~K.~Benander, \emph{Survivorship structure
of the two-channel muffled hazard}, companion extraction manuscript,
2026.
\bibitem{Ben25Seq} A.~K.~Benander, \emph{Sequential achievement},
companion manuscript (AUTHOR TO VERIFY title and year).
\bibitem{SO94} C.~Sagan and S.~J.~Ostro, \emph{Dangers of asteroid
deflection}, Nature \textbf{368} (1994), 501.
\bibitem{Hag19} M.~H\"agglund, \emph{This Life: Secular Faith and
Spiritual Freedom}, Pantheon, New York, 2019.
\bibitem{Wil73} B.~Williams, \emph{The Makropulos case: reflections on
the tedium of immortality}, in Problems of the Self, Cambridge
University Press, 1973, 82--100.
\bibitem{AS04} A.~Albrecht and L.~Sorbo, \emph{Can the universe
afford inflation?}, Phys.\ Rev.\ D \textbf{70} (2004), 063528.
\bibitem{Car17} S.~M.~Carroll, \emph{Why Boltzmann brains are bad},
preprint arXiv:1702.00850, 2017.
\bibitem{MacIntyre81} A.~MacIntyre, \emph{After Virtue: A Study in
Moral Theory}, University of Notre Dame Press, 1981.
\bibitem{Anscombe58} G.~E.~M.~Anscombe, \emph{Modern moral
philosophy}, Philosophy \textbf{33} (1958), 1--19.
\bibitem{Chr07} D.~Christensen, \emph{Epistemology of disagreement:
the good news}, Philos.\ Rev. \textbf{116} (2007), 187--217.
\bibitem{nietzsche1974gay} F.~Nietzsche, \emph{The Gay Science},
trans.\ W.~Kaufmann, Vintage, New York, 1974 (orig.\ 1882).
\bibitem{weinberg1977first} S.~Weinberg, \emph{The First Three
Minutes}, Basic Books, New York, 1977.
\bibitem{Camus42} A.~Camus, \emph{The Myth of Sisyphus}, Gallimard,
Paris, 1942 (French original \emph{Le Mythe de Sisyphe}).
\bibitem{Sartre46} J.-P.~Sartre, \emph{Existentialism Is a
Humanism}, lecture, published 1946.
\bibitem{Carroll16} S.~Carroll, \emph{The Big Picture}, Dutton,
2016.
\bibitem{Midgley85} M.~Midgley, \emph{Evolution as a Religion:
Strange Hopes and Stranger Fears}, Methuen, London, 1985.
\bibitem{Kie1846} S.~Kierkegaard, \emph{Concluding Unscientific
Postscript}, 1846.
\bibitem{Hei27} M.~Heidegger, \emph{Being and Time}, 1927.
\bibitem{Frankl59} V.~E.~Frankl, \emph{Man's Search for Meaning},
Beacon, Boston, 1959 (orig.\ 1946).
\bibitem{Nagel71} T.~Nagel, \emph{The absurd}, J.\ Philos.\
\textbf{68} (1971), 716--727.
\bibitem{Wolf10} S.~Wolf, \emph{Meaning in Life and Why It Matters},
Princeton University Press, 2010.
\bibitem{Metz13} T.~Metz, \emph{Meaning in Life: An Analytic Study},
Oxford University Press, 2013.
\bibitem{Gallup70} G.~G.~Gallup, \emph{Chimpanzees: self-recognition},
Science \textbf{167} (1970), 86--87.
\bibitem{MarxEngels48} K.~Marx and F.~Engels, \emph{Manifesto of the
Communist Party}, 1848.
\bibitem{Horowitz17} A.~Horowitz, \emph{Smelling themselves: dogs
investigate their own odours longer when modified in an ``olfactory
mirror'' test}, Behavioural Processes \textbf{143} (2017), 17--24.
\end{thebibliography}
\end{filecontents*}
\IfFileExists{biblatex.sty}{%
  \usepackage[backend=biber,style=numeric,sorting=none,maxbibnames=99,giveninits=true]{biblatex}%
  \addbibresource{golden-reasoning.bib}%
  \newcommand{\PrintRefs}{\printbibliography[title={References}]}%
}{%
  \newcommand{\PrintRefs}{\input{gr-manualbib.tex}}%
}

\numberwithin{equation}{section}

\theoremstyle{plain}
\newtheorem{theorem}{Theorem}[section]
\newtheorem{proposition}[theorem]{Proposition}
\newtheorem{lemma}[theorem]{Lemma}
\newtheorem{corollary}[theorem]{Corollary}
\newtheorem{claim}[theorem]{Claim}

\theoremstyle{definition}
\newtheorem{definition}[theorem]{Definition}
\newtheorem{ansatz}[theorem]{Ansatz}
\newtheorem{convention}[theorem]{Convention}
\newtheorem{constraint}{Constraint}
\newtheorem{principle}{Principle}

\theoremstyle{remark}
\newtheorem{remark}[theorem]{Remark}
\newtheorem{example}[theorem]{Example}

\newcommand{\E}{\mathbb{E}}
\newcommand{\PP}{\mathbb{P}}
\newcommand{\R}{\mathbb{R}}
\DeclareMathOperator{\Var}{Var}
\DeclareMathOperator{\Cov}{Cov}
\DeclareMathOperator{\Leb}{Leb}
\newcommand{\eqd}{\overset{d}{=}}
\newcommand{\one}{\mathbf{1}}
\newcommand{\Imm}{\mathcal{I}}
\newcommand{\Fc}{\mathcal{F}}
\newcommand{\Lc}{\mathcal{L}}
\newcommand{\hm}{\mathfrak{h}}
\newcommand{\Ac}{\mathcal{A}}
\newcommand{\Sc}{\mathsf{S}}

\allowdisplaybreaks

\newtheorem{premise}{Premise}
\renewcommand{\thepremise}{$\Phi$\arabic{premise}}
\newtheorem{objection}[theorem]{Objection}

\title{Golden Reasoning\\[4pt]
\large Soft Horizons, Survivorship, and the Two-Channel Muffled Hazard}
\author{Ajax Benander\thanks{Email: \href{mailto:akb9408@rit.edu}{akb9408@rit.edu}.}}
\date{August 2026}

\begin{document}
\maketitle

\begin{abstract}\small
\noindent A biosphere that survives indefinitely must satisfy two
conditions, not one. We model total hazard as a two-channel process,
an external background muffled by accumulated capacity plus an
internal term that the same capacity feeds, and we place the noise
on the negative logarithmic derivative of the hazard, one
integration further out than is standard. Survival then requires a
\emph{sign} condition on how hazard responds to capacity, which we
call \emph{technoanatomy}: where the sign is adverse, viability is
identically zero for every parameter value and every driver, so that
no amount of capacity helps, while where it is favourable, viability
is the product of that indicator with the driver's escape
probability, and the two conditions are independent. In the critical
case the survival tail is $t^{-1/4}$, the persistence exponent of
integrated Brownian motion transferred to the soft-killing
functional by an explicitly flagged Ansatz whose pre-registered
falsifiers are tested numerically. Conditioning on survival
converges to a Doob $h$-transform, locally absolutely continuous
with respect to the unconditioned law but mutually singular from it
at infinity, and from the singularity follows an epistemic ceiling:
no finite observation certifies which measure one is under. The
ceiling is proved once and then turned on this work's own
conclusions as firmly as on its targets.

Two familiar positions fall to the same structure. Programmes that
measure civilisational progress by capacity alone are indexing the
multiplicand and not the indicator, and programmes that would
withdraw capacity leave the background unmuffled, reaching viability
zero by the other route. Both conclusions are derived rather than
argued separately.

The model is one of three movements. The second is a conditioning
argument: that the record we hold reads as the past-light-cone
restriction of a conditioning on an event in our \emph{future} light
cone, the indefinite survival of the biosphere, with the companion
program's observer-conditioning recovered as its pre-technological
restriction. The future-conditioned tradition of Barrow and Tipler,
the Omega Point, and axiarchism is placed against the argument delta
by delta, and the strong forward reading is developed with a
phenomenology and a discriminating quantity but, by the ceiling, no
decisive test.

The third movement is a set of positions, each conditional on a
single adoptable premise and each fenced: a grounding for the
normative vocabulary, a decision rule whose limit form is a ratio of
persistence-harmonic values, an ethics, and, for the reader the
abstract's first sentence does not reach, a third route to meaning,
\emph{inferred} rather than revealed or constructed. The route is
stated at two tiers, an unconditional claim that the
constructed-versus-discovered dichotomy is broken and a
Strong-conditional claim that the settlement's inference inverts;
the post-Nietzschean crisis is classified as a formal invalidity and
reversed at its root; and the framework's reception-scale
predictions are entered scoreable, with the conditions of their
failure attached. The dependency between the movements runs one way,
so that the results and the argument consume nothing from the
positions, and a reader unpersuaded by the positions holds the first
two intact. The whole is organised in five parts, with the chain
from evidence to conduct stated once, link by link, at its grades.
\end{abstract}

\setcounter{tocdepth}{2} % Aethic depth (referee minor 8: trims the ToC)
\tableofcontents
\bigskip

% ---------------------------------------------------------------------
\section*{Notation}\label{sec:notation}
\addcontentsline{toc}{section}{Notation}

\noindent Symbols are listed by the object they name. Where a letter does
double duty the two uses are separated by section, and the collisions
that remain are flagged. One is flagged here rather than in a table
note: $\varkappa$, the schedule channel's mean offspring
(\S\ref{sch:sub:schedule}), is not $\kappa$, the exceedance factor of
the imported Accordance bound (\S\ref{gc:sub:imported}, used again in
\S\ref{ra:sec}); the glyphs are close and the objects are not.

\smallskip
\noindent\textbf{The hazard model} (\S\ref{gr1:sec}).
\begin{center}\small
\begin{tabular}{@{}ll@{}}
\toprule
$h^{\mathrm{tot}}_t$ & total hazard rate, \eqref{tc:eq:hazard}\\
$B_t$ & external background; $B_t=\sum_i e^{\epsilon W^{(i)}_t}$, random\\
$\varsigma_b$ & muffling factor, logistic, $\varsigma_b(0)=1$ exactly\\
$\Gamma$ & internal coefficient (the trough at zero capacity)\\
$\rho$ & \emph{technoanatomy}: hazard's response to capacity,
Definition~\ref{gr1:def:anatomy}\\
$x_t$ & muffling exponent (\emph{technostrength}); $x=\int\mu$,
\eqref{tc:eq:x}\\
$\mu_t$ & driver, $\mu=-\tfrac{d}{dt}\log h$; Brownian post-onset,
\eqref{tc:eq:driver}\\
$\sigma$ & driver volatility \qquad $b$ --- muffling offset \qquad
$k$ --- $\log(B/\Gamma)$\\
$\Lambda_t$ & cumulative hazard $\int_0^t h\,ds$\\
$T$ & departure time (Cox construction, Definition~\ref{sh:def:cox})\\
$\tau$ & hard-wall functional \qquad $\Imm=\{\Lambda_\infty<\infty\}$ ---
viability, Definition~\ref{sh:def:viab}\\
$h_\ast$ & hazard floor under adverse anatomy \qquad $x_\ast$ --- its
argument\\
\bottomrule
\end{tabular}
\end{center}

\smallskip
\noindent\textbf{Measures and the survivorship limit}
(\S\S\ref{gr1:sub:cond}, \ref{gr:sec}).
\begin{center}\small
\begin{tabular}{@{}ll@{}}
\toprule
$\PP$ & unconditioned law --- the \emph{General} biosphere\\
$Q$ & survivorship limit, a Doob $\hm$-transform --- the \emph{Golden}
biosphere, Definition~\ref{gr:def:gg}\\
$\PP^{(t)}$ & $\PP(\cdot\mid T>t)$, the \emph{family} interpolating them;
not a measure, Remark~\ref{gr:rem:familyvsmeasure}\\
$\hm$ & persistence-harmonic function; $dQ/d\PP|_{\Fc_t}=\hm(X_t)/\hm(X_0)$,
\eqref{sh:eq:Qdensity}\\
$s(\cdot)$ & scale function \qquad $\Fc_t$ --- filtration\\
\bottomrule
\end{tabular}
\end{center}

\smallskip
\noindent\textbf{The schedule channel} (\S\ref{sch:sub:schedule}).
\begin{center}\small
\begin{tabular}{@{}ll@{}}
\toprule
$n$ & gates in the schedule \qquad $m$ --- options per gate\\
$\lambda_{k,j}$ & rate of option $j$ at gate $k$; $\lambda_k=\sum_j$,
$\pi_k=\lambda_{k,1}/\lambda_k$\\
$Z_t$ & outstanding debt \qquad $r$ --- gates owed per wrong gate\\
$\varkappa$ & mean offspring $r(1-\pi)$; criticality at
$\E[\log\varkappa]=0$\\
$u$ & completion weight, least root of $u=s\,f(u)$\\
$c_{\mathrm{eff}}$ & effective contrast, $r\log_{10}(1+(B+\Gamma)/\lambda)$\\
$\gamma_k,\nu_k$ & staged leak on credit and driver; both zero at $k=n$\\
\bottomrule
\end{tabular}
\end{center}

\smallskip
\noindent\textbf{Decision, ethics, and the Route World}
(\S\S\ref{gs:sec}, \ref{gc:subsub:routeworld}).
\begin{center}\small
\begin{tabular}{@{}ll@{}}
\toprule
$A,B$ & options (\S\ref{gs:sec}); $P(A)$ --- Aethic weight, read as
attainability\\
$\operatorname{Aur}$ & Aurutance, $\log$ of the decision ratio\\
$\mathrm{AM},\mathrm{RM}$ & automorality (act-indexed), relative morality
(maxim-indexed)\\
$\mathcal G$ & the Golden Point as a conditioning event\\
$\Lc_t,\Lc_\ast$ & accumulated and total log-evidence,
Definition~\ref{gc:def:evidence}\\
$\kappa$ & exceedance parameter of the Accordance Principle,
\eqref{gc:eq:retroexceed}\\
\bottomrule
\end{tabular}
\end{center}

\smallskip
\noindent\textbf{Collisions the reader should know about.} $s$ is the
scale function in \S\ref{gr1:sec} and the per-gate survival factor in
\S\ref{sch:sub:schedule}; $T$ is the departure time throughout except in
\S\ref{gr:subsub:measured}, where $T=4$ is a conditioning horizon; $n$ is
the background count in \eqref{tc:eq:hazard}, the gate count in
\S\ref{sch:sub:schedule}, and the institution count in \S\ref{ad:sec};
$b$ is the muffling offset and, in \S\ref{ad:sec}, a Heider weight;
$\kappa$ is the exceedance parameter and, in Appendix~\ref{app:horizon},
a gain minimum; $\varkappa$ (curly) is the schedule's mean offspring and
is a distinct symbol. $\Phi$ with a numeral indexes premises; $\Phi$
alone is the normal CDF.

% ---------------------------------------------------------------------
\section*{Glossary of coined terms}\label{sec:glossary}
\addcontentsline{toc}{section}{Glossary of coined terms}

\noindent Terms introduced by this work or by its companions, with the
place each is defined. Terms marked \textsuperscript{$\dagger$} are
imported from the companion program and are stated at
\S\ref{gc:sub:imported}.

\begin{description}\small\setlength\itemsep{1pt}
\item[Technostrength] The muffling exponent $x$: accumulated capacity to
lower hazard. \eqref{tc:eq:x}
\item[Technoanatomy] $\rho$: the \emph{sign} of hazard's response to
capacity. Definition~\ref{gr1:def:anatomy}
\item[Crest / trough channel] The external and internal terms of
\eqref{tc:eq:hazard}.
\item[Soft wall] The killing functional replacing a hard boundary;
Ansatz~\ref{sh:ansatz} relates the two.
\item[Golden Point] The greatest fixed point $\nu F$ of the safety
operator: a state admitting no foreclosing continuation.
\item[Departure Point] Its dual $\mu G$: a state all of whose
continuations foreclose --- the biosphere's losing condition, which
precedes physical death and is separable from it in both directions.
Remark~\ref{ac:rem:losingcondition}
\item[Departure-sufficient event] The referent of the hazard: an
occurrence whose realization forecloses, with death-events its typical
occupants and nothing biological presupposed.
Remark~\ref{gr1:rem:hazardof}
\item[Golden-categorical robustness] The definitional-root premise:
feature-terms are well-defined to the degree their reference is robust
relative to the asymptotic survivorship structure; the companion's
Aethic-categorical criterion at full depth.
Remark~\ref{gr1:rem:goldendef}
\item[Task-object] An object individuated by its contribution to a
designated root Aethus; the constructive form of Golden-categorical
robustness, with the Golden Point its generative root and the referent
of contribution. Definition~\ref{to:def:taskobject},
Remark~\ref{to:rem:goldenroot}
\item[Golden Aethus] The Aethus stating exactly the Golden attribute,
all else blank; the index of the Golden Nexus, and --- under the
Strong reading --- the substrate on which every thing is a
weight-carrying fluctuation. Definition~\ref{to:def:goldenaethus}
\item[Commitment primitive (node)] An undefended normative placement
from which practical argument propagates; judged non-arbitrary to the
degree it traces to the Golden root. \S\ref{gs:sub:nodes}
\item[Inferred meaning (the third route)] Meaning discovered in the
record and its conditioning structure --- neither constructed by the
observer nor revealed through a non-natural channel; unconditional at
the dichotomy-breaking tier, Strong-conditional at the inversion
tier. \S\ref{gs:sub:thirdroute}
\item[Nietzschean settlement] The modern operating frame holding that
constructed meaning is the only honest option; engaged and superseded
at \S\ref{gs:sub:thirdroute}, its diagnosis retained.
\item[Species-level self-awareness] The criterion: a species is
self-aware when it possesses the concept of its own modal poles
($\mu G$, $\nu F$) as such; on this test the discontinuity is
conceptual and, as of writing, ahead.
Remark~\ref{hr:rem:speciesaware}
\item[Inference flip] The vantage-indexed best-inference operator's
change of output at a named evidential arrival; the meaning-question's
flip condition is the coupling of Miracles, on the Tricky-Principle
precedent. Remark~\ref{gs:rem:inferenceflip}
\item[Load path] Part~IV's ten links stated as one auditable chain,
coupling through owned conduct, graded prefix-to-suffix.
\S\ref{gs:sub:loadpath}
\item[General / Golden Biosphere] A biosphere under $\PP$ / under $Q$.
Definition~\ref{gr:def:gg}
\item[Goldening] The interpolation $\PP^{(t)}$ between them.
\item[Golden Filter] The reweighting survivorship performs on an
ensemble. Proposition~\ref{ac:prop:filter}
\item[Golden Tendency] The ensemble consequence of the Filter.
Principle~\ref{ac:prin:tendency}
\item[Golden Question / Answer] \emph{Should our biosphere live
forever?}, and what the Filter does to frozen answers.
\S\ref{ac:sub:question}
\item[Golden Imperative] Passage through the Golden Point as the only
route avoiding $\mu G$. Principle~\ref{ac:prin:imperative}
\item[Golden Extrapolation] Self-placement at finite $t$ on the
goldening. Definition~\ref{ac:def:extrapolation}
\item[Golden Initiative] Persistence-directed practice, named from
inside. Definition~\ref{gs:def:initiative}
\item[Golden Score / Aurutance] The decision ratio $\hm(A)/\hm(B)$ and
its logarithm. \S\ref{gs:sub:limit}
\item[Departric] Prefix: automorality negative. \textbf{Biotic}:
positive. \textbf{Neffective}: zero. Table~\ref{gs:tab:moralities}
\item[Schedule / gate / debt / leak] The third hazard channel and its
parts. \S\ref{sch:sub:schedule}
\item[Route World] The world of the Strong extrapolation hypothesis.
Definition~\ref{gc:def:routeworld}
\item[Golden Optimation] Optimation under a future-cone conditioning
event. Definition~\ref{gc:def:goldenopt}
\item[Golden causality] The pattern in which Golden Optimation accounts
for a present event. Definition~\ref{gc:def:goldencausality}
\item[Investment / Invested Golden causality] The weight-versus-necessity
bound: a realized route's improbability lower-bounds how much more it
accomplishes, at exceedance confidence. Proposition~\ref{ra:prop:investment};
phenomenological face at Remark~\ref{gc:rem:invested}
\item[Departric invalidity / Mattering residue] That past $\mu G$ the
adopted ordering is undefined rather than small --- a claim about an
index --- and the pruning whose residue the dominance argument runs on.
\S\ref{gs:ssub:departric}, \S\ref{gc:sub:acting}
\item[Retrocausal (conditioning)] Conditioned on an attribute in one's
future light cone; nothing further. Definition~\ref{gc:def:retrocausal}
\item[Retrocausal Accordance] The Accordance Principle run forward: one's
present cannot be arbitrary to one's future.
Principle~\ref{gc:prin:retroaccord}, and \eqref{ra:eq:bound} for the
bound form.
\item[Finale paradox] Why one occupies the present rather than deep time,
on a reading with no resource to discharge the fluke; the relay point at
which the companion's ``when does Optimation stop'' frontier is
dissolved rather than answered. \S\ref{ra:sub:finale},
Remark~\ref{ra:rem:relay}
\item[Adversivity] Net antagonism in an institutional balance.
\S\ref{ad:sec}
\item[Aethus / Aethae]\textsuperscript{$\dagger$} An information state
and its refinements.
\item[Nexus / Cosmic Nexus]\textsuperscript{$\dagger$} A weighting scheme
over Aethae; the unconditioned one.
\item[Optimation]\textsuperscript{$\dagger$} The fluke content a
conditioned record displays.
\item[Accordance Principle]\textsuperscript{$\dagger$} The typicality
constraint binding attribute by attribute.
\item[Imperfection Magnitude Drop]\textsuperscript{$\dagger$} A
weight-suppressing perturbation short of foreclosure.
\end{description}

% ---------------------------------------------------------------------

\section{Introduction}\label{intro:sec}

Will the biosphere continue indefinitely? This work treats that
question as one with enough structure to carry theorems at one end
and an account of meaning at the other, and it makes the crossing in
one direction only, so that the mathematics never borrows from the
philosophy and every claim along the way is carried at a stated
grade. Concretely, the document contains a stochastic model of
survivorship together with its theorems (Part~II); a conditioning
argument over the historical record (Parts~I and~III); and a set of
normative, ethical, and existential positions that consume the first
two while being consumed by nothing (Parts~IV and~V). A reader who
wants only the probability may therefore read Part~II and stop,
holding everything this work proves, while a reader who came for the
older question underneath will find it answered, at its grade, in
Part~IV.

\emph{One must imagine Sisyphus happy. A pale blue dot. You make
your own meaning.} The reader knows these consolations, because they
are the standing settlement of the educated present: the shared
conclusion that the universe supplies no meaning, that construction
is what remains, and that maturity consists in bearing this well.
This paper engages the settlement on its own terms, steelmanning it
in \S\ref{gs:sub:thirdroute} as the uniquely reasonable position
given what its authors could see, and then argues that its inventory
was incomplete. Meaning has been offered as \emph{revealed}, and
after the revelation failed it has been offered as
\emph{constructed}; but there is a third cell in that table, meaning
\emph{inferred}, read off an object under stated inference
discipline with published falsifiers and a proven ceiling on
certainty, and that cell has stood empty less because anyone
examined and rejected it than because nothing had ever seemed to
occupy it. This document is an attempt to occupy it at a stated
price, and the price is displayed before anything is purchased.

\subsection{The thesis}\label{intro:sub:thesis}

Capacity to act on the world is usually treated as the quantity that
decides whether a civilisation lasts, and this work argues that it is
the wrong quantity, because the right one is a sign rather than a
magnitude. Consider a system facing hazard from two directions, an
external background it can learn to muffle and an internal channel
that its own capacity feeds. Accumulating capacity lowers the first
and does something to the second, and what it does to the second is a
structural property of the system rather than a matter of how much
capacity it holds. Call that property technoanatomy and write it
$\rho$. The central result is that its sign is decisive: adverse
$\rho$ makes indefinite survival impossible outright, not unlikely
but impossible, for every parameter value and every history, while
favourable $\rho$ makes survival possible without making it likely,
since a second and independent condition must also hold.

Two consequences follow immediately, and both cut against widely held
positions. A programme that measures progress by energy captured or
capability acquired is measuring a quantity that cannot distinguish
the two cases, while a programme that would withdraw capacity leaves
the background unmuffled and arrives at the same zero from the other
side. These are not two objections but one, and the paper derives
them together.

\subsection{The model in three lines}\label{intro:sub:model}

A reader who wants the mathematics before the vocabulary needs only the
following, which \S\ref{gr1:sec} develops properly.

\emph{The hazard.} $h^{\mathrm{tot}}_t=B_t\varsigma_b(x_t)+\Gamma
e^{-\rho x_t}$, \eqref{tc:eq:hazard}: an external background $B_t$
attenuated by a muffling factor $\varsigma_b$ that capacity $x$ improves,
plus an internal term that capacity feeds or starves according to the
sign of $\rho$. The normalisation is $\varsigma_b(0)=1$, so at zero
capacity the crest is the bare background.

\emph{The driver.} $x_t=\int_0^t\mu_s\,ds$ with $\mu$ Brownian: the noise
sits on $\mu=-\tfrac{d}{dt}\log h$, one integration beyond where standard
practice puts it. That placement is the paper's principal modelling
claim, and it is what moves the persistence exponent from $\tfrac12$ to
$\tfrac14$.

\emph{Departure.} $T$ is the first time the cumulative hazard exhausts an
independent exponential (Definition~\ref{sh:def:cox}), and viability is
the event $\{\Lambda_\infty<\infty\}$ --- that the total hazard ever
accrued stays finite. Departure with probability one and infinite life
expectancy are compatible, and both hold in the critical case.

\subsection{What this work does not claim}\label{intro:sub:notclaimed}

One reconciliation is owed at the outset, because it is the objection
a reader will arrive holding. This paper builds a decision rule and an
ethics on a conditional distribution, and a conditional distribution
entails no obligation, so the two must be squared before anything else
is said. The reconciliation is a division of labour between two
premises. $\Phi2$ (Premise~\ref{gr:phi2}) governs what the apparatus
outputs, which is distributions, from which no duty follows. $\Phi6$
(Premise~\ref{ac:phi6}) governs what an agent does with a choice the
apparatus never makes, namely which grounding to occupy, given that
some grounding is occupied whether or not anything is said. The two
premises operate at different levels and do not meet, with the
consequence that everything normative in this work is conditional on
$\Phi6$, which is offered and priced and may be declined, and an
agent who declines it contradicts nothing proved here.
\S\ref{ac:sub:countermeasure} states the split in four checkable
steps, and \S\ref{gs:sub:fence} restates it where the decision rule
might otherwise be over-read.

One structural fact should be stated alongside it, because it is the
protection the mathematics needs from the rest of the document. This
work contains three kinds of thing: results, meaning the hazard model
of \S\ref{gr1:sec} with its theorems and the survivorship transform
of \S\ref{gr:sec}; a graded argument, meaning the conditioning
principle of \S\ref{gc:sec}, carried at \textbf{[S]} with its
extrapolation hypotheses at \textbf{[O]}; and positions, meaning the
grounding of \S\ref{ac:sec}, the decision rule of \S\ref{gs:sec},
the ethics of \S\ref{gs:sub:ethics}, and the placement of
\S\ref{hr:sec} against extinctionism and techno-optimism, all
carried at \textbf{[$\Phi$]} and all conditional on the one
adoptable premise. The dependency between the three runs one way. The
results consume nothing from the positions, the graded argument
consumes nothing from the positions, and the positions consume both,
so a reader unpersuaded by the positions is left holding the results
and the argument intact, with nothing in the positions usable against
them. A reader who finds the positions objectionable, reading a paper
on survivorship conditioning that takes a stance on longtermism as
advocacy, should read the results and the argument first and the
positions never, and will have read everything this work proves.

Beyond that, three standing limits. The epistemic ceiling
(Theorem~\ref{gr:thm:ceiling}) forbids certifying any system's position
under survivorship from finite observation --- including one's own, and
including this work's. No actor, institution, or programme is scored
anywhere in this paper. And the model's practical relevance rests on
technoanatomy being a process rather than a fixed parameter, which
\S\ref{gr1:subsub:fluct} does not yet supply; three sections record that
as their inherited dependency, and it is the paper's principal exposure.

\subsection{The situation: two announcements and an empty cell}
\label{intro:sub:situation}

Two sentences, ninety-five years apart, frame the condition this
work addresses. The first is Nietzsche's, from 1882: ``God is dead.
God remains dead. And we have killed him.'' \cite{nietzsche1974gay}
Read structurally, what the sentence announced is that the
\emph{revealed} cell of the meaning inventory had emptied, and with
it, by an inference that went unexamined because no alternative was
visible, that meaning must henceforth be made rather than found. The
second is Weinberg's, from 1977, delivered from inside the
enterprise that had done the emptying: ``The more the universe seems
comprehensible, the more it also seems pointless.''
\cite{weinberg1977first} His successor sentence offered
comprehension itself as the consolation, research as one of the few
things lifting human life above the level of farce, and that is the
settlement in its noblest form: lucidity as the remaining good, held
by a mind that saw no further inventory to check.

Where the revealed cell's contents came from deserves one
disciplined paragraph, because it fixes what any successor must not
be. The premodern supply of meaning ran on a mechanism that this
program's companion writings name \emph{the premodern dragon}: the
coupling of ancestral fear, the inheritance of a lineage that spent
its formative epochs as prey, with the human appetite for control
over other humans. The gods that domesticated the terror also
collected the rent, and the arrangement priced meaning in obedience.
What the scientific enlightenment did, on this reading, was not so
much refute a doctrine as remove a mechanism, and the whiplash of
the modern condition follows from the fact that the operation which
removed a false floor removed the only floor then in inventory. The
settlement is the sound of a civilisation deciding, reasonably
enough on its evidence, that floors were never load-bearing to begin
with. On the essential point the present work sides with the
settlement without reservation: the dragon must stay dead, and
\S\ref{gs:sub:notreligion} prices in six checkable deltas the claim
that nothing here restores it. But the open question was never
whether to go back. It is whether a dead dragon and an empty cell
exhaust the options.

The diagnosis this framework offers of Nietzsche's announcement is
therefore epistemic inexperience rather than categorical limitation:
humanity had never yet held an object from which meaning could be
inferred, and mistook never having held one for the impossibility of
its existing. Nobody in 1882 was positioned to detect the mistake,
which is why the settlement was, and is steelmanned here as, the
uniquely reasonable reading of its evidence. The deadlock that
followed, described most exactly by MacIntyre as interminable and
incommensurable assertion \cite{MacIntyre81}, is classified in this
work as a crisis on the inference side: every tradition rightly
discounts every other tradition's root commitments, and its own,
because commitments admit no inference in either direction, a
condition to which \S\ref{gs:sub:nodes} gives its criterion.
MacIntyre closed by waiting for another, doubtless very different,
Benedict. The fallback proposed here is not a person and not a
community but an \emph{object}: a conditioning target with a
mathematics, a ladder of grades, published falsifiers, and a proven
ceiling on what anyone may claim to know about it. Part~IV states
what inferring from it costs, in the near-unconditional tier and the
Strong tier separately, and enters its reception-scale predictions
scoreable.

% [AUTHOR SLOT — the 2001 exhibit: the species-self-awareness scene,
%  in the author's voice, one paragraph. Reserved per the v71
%  dictation; wire to Remark hr:rem:speciesaware.]
% [AUTHOR SLOT — the Tattoo exhibit: the valor observation, in the
%  author's voice, one paragraph. Reserved per the v71 dictation;
%  wire to Remark gs:rem:valorvacancy.]

None of the mathematics depends on this subsection, since the
dependency runs the other way, as \S\ref{intro:sub:notclaimed} has
already priced, and a reader who wants the theorems alone loses
nothing by proceeding directly to Part~II. A reader for whom this
subsection named their situation should know that the paper was
written in the conviction that such readers exist at scale, and that
\S\ref{gs:sub:loadpath} states, link by link and at its grades, what
this work believes it can offer them.

\subsection{How to read this}\label{intro:sub:reading}

The document is long and its parts are separable. Four entry points.

\emph{For the mathematics}, start at \S\ref{gr1:sec} and read to
\S\ref{gr:sec}. That stretch is self-contained: a hazard process, its
persistence asymptotics, the anatomy trichotomy, the two-condition
theorem, and the survivorship transform. It consumes nothing from the
companion program and can be assessed on its own terms.

\emph{For the argument about capacity and risk}, read
\S\ref{gr1:sub:cond} for the two-condition theorem and then
\S\ref{hr:sec}, which applies it to extinctionism and to
techno-optimism in four pages.

\emph{For the philosophical grounding}, read \S\ref{ac:sec}. It is the
most demanding section and depends on results stated later; a reader who
wants it first should take \S\ref{intro:sub:model} as the model and
return.

\emph{For the anthropic material}, read \S\ref{gc:sec}, beginning with
\S\ref{gc:sub:imported}, which states the companion apparatus this
section consumes; its human-history argument
(\S\ref{gc:sub:humanhistory}) is the one to read if reading only one.

\emph{For the argument that the Golden Point is the right target}, read
\S\ref{ra:sec}, which develops the retrocausal accordance principle,
runs it on a worked case at human scale, and derives from it the
\emph{shape} the target must have --- and then identifies the Golden
Point as what has that shape. It is the primary argument for the target;
\S\ref{gc:sec}'s reductio clears the ground for it.

\emph{For the framework's most distinctive prediction}, read
\S\ref{gc:subsub:routeworld}, the Route World, which develops the strong
forward reading of the conditioning argument and states what would ---
and by theorem what would not --- distinguish it from its rivals. It is
carried at \textbf{[O]} throughout --- save the conduct argument of
\S\ref{gc:sub:acting}, at \textbf{[$\Phi$]} on an
adopted valuation --- and can be read after \S\ref{gc:sec}
without the intervening sections.

Two conventions run throughout. Every substantive claim carries a grade
--- \textbf{[P]} proved, \textbf{[C]} cited, \textbf{[A]} conditional on a
stated Ansatz, \textbf{[S]} argued at scaling or schematic level,
\textbf{[$\Phi$]} premise carrying its own falsifier, \textbf{[O]} open --- and the grades are load-bearing:
\S\ref{gc:sec}'s central claims are \textbf{[$\Phi$]} and \textbf{[S]},
and what it proves is proved conditional on them, which it says. And each
section closes with a ledger stating what it costs, what it buys, and
what it does not claim. A reader in a hurry can read the ledgers.


% ============================================================
% §2 — AGENCY'S CURSE AND THE GOLDEN QUESTION
% Compiled working copy, REBUILD 5 — post complete-section
% referee round (B1–B4, C discharged). Order:
%   2.1 The problem                (+ B4 selectability falsifier)
%   2.2 The countermeasure
%   2.3 The Golden Tendency
%   2.4 The Golden Question
%   2.5 Golden Extrapolation       (rev 2: B1 monotonicity
%                                   repair; B2 gradient named)
%   2.6 The Golden Imperative      (+ gc:rem:pairsuffices)
%   2.7 Placement
%   2.8 Ledger                     (rev 2: B3 heading split; B4
%                                   Costs line; C fixes)
% followed by the consolidated owed-master-edits file.
%
% Compilation notes:
%   - Numbering PROVISIONAL; cross-references symbolic.
%     FILENAMES RETAIN WRITING-ORDER NUMBERS — block order below
%     is authoritative, not the filenames.
%   - Encoding clean (no literal Φ/§ on non-comment lines).
%   - Counter guard {5} in the countermeasure; {0} pair at §6's
%     head in the owed-edits block, NOT yet in the master. §2
%     adds exactly ONE premise (Φ6).
%   - Owed-edits: 1.5 DRAFTED; remaining ⛔: 1.1, 1.4, 1.6
%     (apply), 1.7 (intro), 1.8–1.10 (mechanical), 1.11
%     (Newcomb companion).
%   - Presumes the master's preamble and labels; a review copy,
%     not standalone-compilable.
% ============================================================

As of the structural pass the paper is organized in five parts, and
the part map is the fastest orientation: Part~I poses the Question
and fixes its objects; Part~II is the survivorship mathematics,
self-contained; Part~III is the inference discipline --- Golden
reasoning, accordance, the Route World --- with the ceiling as its
governing limit; Part~IV is the normative and existential engine,
the Score through the third route; Part~V is the human position.
A reader wanting only the probability reads Part~II; only the
philosophy of meaning, Part~IV after Part~I's vocabulary; the
whole argument, in order.

% =====================================================================
% STRUCTURAL PASS (Aug 22 2026): five-part organization instated.
% Plain \part used (no per-part section-counter reset); the Aethic
% convention \newcommand{\Part}[1]{\part{#1}\setcounter{section}{0}}
% is available by a one-line change and is DEFERRED to the author
% (numbering-stability conservatism).
% =====================================================================
\part{The Question and Its Target}

\noindent Part~I fixes the question and its objects, and proves
nothing quantitative. Agency's curse is posed and the Golden Question
is stated at its price (\S\ref{ac:sec}); the conditioning target is
then chosen, with the Departure Point built as a least fixed point
together with its cascade, continuity defended, and the rival
periodization engaged (\S\ref{gc:sec}); and task-objects supply the
individuation that the perturbation arguments consume, an arc that
closes on the Golden-categorical premise and, under the Strong
reading only, on the Golden Aethus (\S\ref{to:sec}). A reader leaves
this part holding a target, an ontology of what is conditioned on,
and the vocabulary in which every later part speaks.

\section{Agency's curse and the Golden Question}\label{ac:sec}

\begin{quote}\small
An agent cannot stop acting, and cannot see far enough to know what its
acting does. This section names the resulting problem --- \emph{agency's
curse} --- argues that no agent escapes it, and proposes the one
countermeasure the predicament leaves available. Since some question
must be held ungrounded, prefer one whose answer the agent's own
existence constrains. That is $\Phi6$, this work's single doorway
through which obligation can enter.

The candidate offered is the Golden Question --- \emph{should our
biosphere live forever?} --- and what recommends it is not that its
answer can be proved but that survivorship selects it: a biosphere
organizing its practice around the negative does less to persist, and
the Filter thins it accordingly. The affirmative is therefore selected
rather than deduced, which is a positional claim and not a truth claim,
and the difference is kept in view throughout.

Two things scope the whole. Whether a biosphere \emph{can} reach what
the answer affirms is settled by $\rho$ and known to nobody
(\S\ref{gr:sub:ceiling}); and the section's practical relevance rests on technoanatomy being a
process rather than a fixed parameter --- which
Remark~\ref{hr:rem:forcesprocess} argues the framework itself requires,
and which \S\ref{gr1:subsub:fluct} has yet to supply.
\end{quote}
% ------------------------------------------------------------
% SECTION-HEAD QUOTE BLOCK — OWED (register 4.1). House pattern
% not visible from outside the master; AUTHOR TO SELECT. The
% child's why-game (2.1) and the burning building are the
% section's two native images if an epigraph in that register
% is wanted instead of a sourced quotation.
% ------------------------------------------------------------

%############################################################
%%##  FILE: sec2_1_the_problem.tex
%############################################################

% ============================================================
% §2.1 — The problem (agency's curse)
% REVISION 2, against referee report (§F scope).
% Grade [Φ] throughout.
%
% MERGE NOTES — disposition of referee items:
%   ACCEPTED IN FULL: B1–B4, C1–C5, D1–D7. C6 (offered) taken:
%     spectrum image placed in Condition 2.
%   MODIFIED, WITH REASON (per report §F):
%     E4 asks that the falsifier statement be kept verbatim, but
%     C1's repair changes what the first horn claims: chains now
%     terminate in something (a non-reason), so a floor that is
%     merely accessible no longer refutes the proposition — the
%     constructivist terminus would meet the old falsifier without
%     touching the claim. Amended by five words: the floor must be
%     accessible AND itself a reason. Everything else verbatim.
%   FIGURE COUNT (brief §1.1: one per passage): burning building
%     (Condition 1), spectrum (Condition 2), driving-blind (the
%     conjunction paragraph — the single weight moment B2 requests).
%     Three figures, three passages, none stacked.
%   CROSS-REF SCOPE: gc:prop:recursive and gc:rem:twoprivilege
%     added per C5. These exceed brief 2.1's cross-ref list
%     (Aethus only); flagged for the merge rather than silently
%     widened.
%   - D5: enumitem CONFIRMED LOADED (register §7); the earlier
%     hand-label fallback caveat is discharged.
%   - B4 (complete-section report): ONE sentence added to
%     ac:rem:threeprops — the selectability falsifier — despite
%     that report's §E stating 2.1–2.4 "should not be touched."
%     B4 explicitly places the falsifier "at ac:rem:threeprops,"
%     which I read as the override; if §E was meant to win, the
%     sentence lifts out cleanly and the falsifier survives in
%     the ledger's Costs entry alone.
%   Label change: ac:claim:nonempty → ac:prop:nonempty (D1).
%   ac:rem:admissible split (D6): admissibility stays under the
%     old label; three properties now ac:rem:threeprops.
%   - \setcounter{premise}{5} guard still belongs before Φ6 in
%     2.3, not here (no premise environment in 2.1).
%   - \ref{ac:sub:countermeasure}, \ref{ac:sub:placement} remain
%     forward placeholders (2.3, 2.2).
% ============================================================

\subsection{The problem}\label{ac:sub:problem}

Aethic reasoning demotes several classical metaphysical problems --- some to irrelevance, some to attributions of constructs \cite{Aethus}. The problem of this subsection is comparatively undemoted: the demotions pass over it, and what they leave standing is the business of the present section. We call it \emph{agency's curse}. It is the conjunction of two conditions on any agent.

\begin{enumerate}[label=\emph{Condition \arabic*.}]
    \item \emph{Forced action.} An agent is at every moment committed to acts of influence on the universe; inaction, sleep, and self-destruction are themselves selections among alternatives with divergent consequences. The position is that of someone crossing a burning building toward a child: the tempo is set by the fire, not by the crosser, and every hesitation is itself a move. But the comparison flatters. The ordinary case ends; the house has walls; the parent knows what they are trying to do. Here the stakes of each selection are unbounded relative to the time allotted for it, and the situation admits neither exit nor completion --- declining to decide is itself a decision, so the condition cannot be suspended, and it attaches to agency as such rather than to human psychology.
    \item \emph{No accessible grounding.} The agent has no access to an objective morality, meaning, or interpretation of the universe against which these forced selections could be weighed --- if any such thing exists at all. Picture a spectrum from inert matter to omniscience. A rock is not cursed, because it does not act; an omniscient agent is not cursed, because it can see; agents occupy the one altitude at which the curse can attach, high enough to be bound by Condition 1 and too low to weigh what it binds them to. The altitude is not a human peculiarity: it is a side effect of the kind of thing an agent is.\footnote{The refusal of human privilege here is a standing commitment of the paper, made again in the choice of conditioning target, at recursive anchoring, and --- at its full width --- at continuity (\ref{gc:prop:recursive}, \ref{gc:rem:twoprivilege}, \S\ref{gc:sub:continuity}).}
\end{enumerate}

The curse is the conjunction. Bound to selections of unbounded consequence and given nothing to weigh them against, an agent is driving blind --- and blind not at the periphery, where instruments can be added, but in the most intimate of its outlooks, at the point where any instrument would have to be read.

Stating this exactly requires three terms.

\begin{definition}[Questions, answers, and frozen questions]\label{ac:def:frozen}
A \emph{question} is a state an agent takes to be held or not held; an \emph{answer} is the state the agent assigns to it. An agent typically assigns an answer by borrowing from the answers to other questions, recursively --- \emph{where are the keys} borrows from \emph{where were they last held}, and so on, until the borrowing terminates. A \emph{frozen question} is a question whose answer the agent firmly holds but which borrows from no question, answer, or other information the agent can access; its answer is a \emph{frozen answer}.
\end{definition}

The ostension is the child's why-game: each answer is met with another \emph{why} until the adult can only say \emph{I don't know}. The question immediately preceding that admission is frozen --- it is the foundation of every question that traced to it, while itself resting on nothing the agent can produce.

\begin{remark}[What the definition excludes]\label{ac:rem:admissible}
A frozen question must be neither paradoxical nor inconsequential to the agent's other commitments; a question whose answer changed nothing downstream, or which admitted no coherent answer, would be frozen only vacuously. The restriction is substantive: a later candidate for the role will have to be shown to satisfy it, and the demonstration will be owed, not assumed.
\end{remark}

\begin{remark}[Three properties, and which one carries the section]\label{ac:rem:threeprops}
The frozen questions the curse concerns have three further properties. They are held by an agent under forced action (Condition 1); they concern what the agent is to \emph{do}, not merely what is the case; and they are \emph{selectable} --- the predicament fixes that some frozen question is held, not which one. The third property carries the remainder of the section: everything built after this point turns on the difference between an unavoidable set and an unavoidable member. What would refute it, in its own vicinity: a demonstration that an agent's frozen questions are fixed by their circumstances rather than chosen within them --- Proposition~\ref{ac:prop:nonempty} would stand, and everything built after it would fall. Nearby notions in the literature share some of these properties and not others; the placement is deferred to \S\ref{ac:sub:placement}, which will trade on exactly these three.
\end{remark}

With the vocabulary in place, the second condition sharpens into a claim.

\begin{proposition}[Non-emptiness; \textbf{[$\Phi$]}]\label{ac:prop:nonempty}
An agent under Conditions 1 and 2 holds at least one frozen question at every moment.
\end{proposition}

The proposition is argued, not proved, and the argument is a dilemma over an exhaustive disjunction --- there is an objective morality or there is not --- so that the two horns below are all there is. What would refute it is either an agent under forced action all of whose answers terminate in something the agent can access that is itself a reason --- a regress with an accessible floor --- or a transmission protocol of the kind ruled out on the second horn below. Neither has been exhibited; the claim is graded accordingly.

\emph{First horn: there is no objective morality.} Then for any forced selection $x$, the chain descending from \emph{why is $x$ to be done} terminates --- not in nothing, but in something that is not itself a reason: a preference, a convention, a disposition. A constructivist is free to call the terminus a fact, and it is one. What it is not is a reason of the kind the chain was descending through, and the question why \emph{that} should bear the weight of everything above it borrows from nothing the agent can access. The frozen question arrives at once.

\emph{Second horn: there is an objective morality.} This is the horn that makes the curse survive the existence of the very thing whose absence seemed to generate it, and it turns on a communication problem. Suppose the objective moral truth exists, and suppose --- granting the most favorable case --- that some agent is in full possession of it and undertakes to transmit it, by any channel whatsoever: written across the sky, planted in dreams, implanted beneath consciousness.

Consider the recipient's information state after the transmission. It is compatible with exactly three hypotheses: the content is true; the content is a deception; the content is irrelevant, applying in some circumstances but not the recipient's. The claim is not that deception is probable, or that the recipient will in fact feel doubt --- it is that \emph{no evidence available to the recipient separates the three hypotheses}. Any certificate of the transmission's veracity would itself have to be an objective moral truth, and the only source of such truths is the channel already under assessment. The request for the discharging evidence re-invokes the very source it was meant to check; the regress is not long, it is a loop of length one. The trichotomy --- true, deceptive, or irrelevant --- is therefore itself a question whose answer the recipient can hold but cannot derive from anything accessible. It is a frozen question, and the horn closes.\footnote{Aethically restated: the recipient holds the transmitted content as an attribute whose warrant is blank, and the blankness cannot be filled from within --- any filling would itself be an attribute with a blank warrant from the same source. The vocabulary is \S\ref{gc:sub:vocab}'s, and the horn is a special case of it.}% CONVERGENCE RESTORATION (verification report C2): reconstructed from the report's description (routing through blankness, forward to gc:sub:vocab); substitute the converged footnote's wording if it differs.

\begin{remark}[Why the second horn is structural, and where it could fail]\label{ac:rem:structural}
The argument is about the recipient's information state, not the recipient's psychology. A recipient of perfect equanimity, untroubled by any felt doubt, is in the same position: the three hypotheses remain undistinguished by anything in their possession, and their equanimity is a frozen answer to the trichotomy, not an escape from it. The horn would fail if a certificate of veracity could be produced whose warrant did not draw on the transmitting source --- an independent channel to the same objective truth. But an independent channel is just a second transmission, and the trichotomy attaches to it identically. What the horn cannot rule out, and does not need to, is the truth of the content; it rules out only the recipient's ever standing in an accessible derivational relation to it.
\end{remark}

Both horns close. On either, the agent holds at least one frozen question; and Condition 1 makes the holding load-bearing, since every forced selection is weighed --- when it is weighed at all --- against a frozen answer or against nothing.

\begin{remark}[What the curse removes, and what it does not]\label{ac:rem:authority}
The curse does not imply that any frozen answer is probably wrong. A frozen answer may be excellently calibrated; nothing above bears on that, and no claim about the accuracy of anyone's morality has been made. What the curse removes is authority of a specific kind: any appeal available \emph{between} agents whose frozen answers conflict. Each can trace their chain to where it stops; neither chain reaches anything the other is obliged to accept; and there the exchange ends --- not because either party is unreasonable, but because the resource an adjudication would spend does not exist. Collectives inherit the position with the stakes multiplied --- collectively held frozen answers have historically been catastrophic --- and the removal binds even that sentence: the condemnation is itself issued from a frozen answer and inherits its standing. It is an instance of the structure under discussion, not an exception to it.
\end{remark}

Four neighbors resemble the curse; none is it. The simulation problem is defanged by Aethic reasoning directly: an Aethus is not demoted in its reality by the scenario, and the felt metaphysical blow is an artifact \cite{Aethus}. The problem of induction concerns whether what sits behind the whole of human inference operates as extrapolation from within would have it. Epistemological skepticism, and its Aethic counterpart, concern whether the information held by the Aethus can be attained by an agent with any surety; Aethic accidentalism carries the same burden into metaphysics \cite{Aethus}.

\begin{remark}[The differentia, and why quietism is unavailable]\label{ac:rem:differentia}
Each of the neighbors admits a quietist response: an agent may register the problem, shrug, reinterpret their place in reality, and continue, at no cost the problem itself can name. This is not a defect in the neighbors; it is what they are like, because each concerns the metaphysical \emph{outside} --- what stands behind induction, whether the world is as presented, whether the Aethus can be reached. Agency's curse concerns the metaphysical \emph{us}, and it binds at the point of forced action: the shrug is itself a selection, the continuation is itself a selection, and each is weighed against a frozen answer or against nothing. Quietism toward the curse is therefore not a response to it but an instance of it. The difference is in kind, not degree, and no ranking of the problems is claimed or needed.
\end{remark}

One thing has gone unexamined. The argument fixes that some frozen question is held; it is silent on which. Whether the silence is an opening is the question of \S\ref{ac:sub:countermeasure}.

% ============================================================
% DELIVERABLE LISTS (brief §0) — updated for revision 2
% ============================================================
%
% 1. CLAIMS REGISTER
%    - Agency's curse is the conjunction of forced action and no
%      accessible grounding. [Φ]
%    - Both conditions are structural: forced action cannot be
%      suspended (declining is deciding); groundlessness attaches
%      to agency as such, at the one "altitude" between inert
%      matter and omniscience where the curse can attach. [Φ]
%      (spectrum image per C6)
%    - Non-anthropocentrism footnote ties the structural claim to
%      the paper's standing commitment (per C5). [Φ]
%    - Admissibility: a frozen question is neither paradoxical nor
%      inconsequential; demonstration owed by 2.4 for the Golden
%      Question. [Φ]
%    - The curse-relevant frozen questions are (i) held under
%      forced action, (ii) action-directed, (iii) selectable;
%      selectability carries the section. [Φ] (2.2 will cite)
%    - Non-emptiness (ac:prop:nonempty): ≥1 frozen question at
%      every moment. [Φ] — argued via dilemma; falsifier in-place
%      (accessible floor that is itself a reason, or a source-
%      independent veracity certificate).
%    - First horn now closes the constructivist objection: chains
%      terminate in a non-reason (preference, convention,
%      disposition); the frozen question is why THAT bears the
%      weight. [Φ] (per C1)
%    - Second horn: no evidence available to the recipient
%      separates true/deceptive/irrelevant; the request re-invokes
%      the source; loop of length one. [Φ] — structural; failure
%      condition in ac:rem:structural.
%    - The curse removes inter-agent authority, not correctness;
%      no claim that frozen answers are probably wrong. [Φ]
%      (per C3)
%    - Collectively held frozen answers have historically been
%      catastrophic, AND the condemnation itself issues from a
%      frozen answer and inherits its standing — self-application.
%      [Φ] (per C4)
%    - Neighbors admit quietism; the curse does not (quietism is
%      an instance, not a response). Difference in kind; no
%      ranking. [Φ]
%    - The curse is comparatively undemoted by the Aethic
%      demotions. [Φ]
%
% 2. CROSS-REFERENCES USED
%    \cite{Aethus} — three times: demotions (opening), simulation
%                    defanging, skepticism/accidentalism.
%    \ref{gc:prop:recursive}, \ref{gc:rem:twoprivilege} — C5
%                    footnote; EXCEEDS brief 2.1's cross-ref list,
%                    added at referee request, flagged for merge.
%    \ref{ac:sub:countermeasure} — forward, ONCE, final sentence
%                    (opening gesture cut per D3/A6). RESOLVED by
%                    the 2.3 draft (sec2_3_the_countermeasure.tex).
%    \ref{ac:sub:placement} — forward to 2.2. RESOLVED by the 2.2
%                    draft (sec2_2_placement.tex).
%    No \cite{Nexus}.
%
% 3. NEW LABELS CREATED
%    ac:sub:problem
%    ac:def:frozen
%    ac:rem:admissible      (admissibility only, post-split)
%    ac:rem:threeprops      (new in rev 2, from D6 split)
%    ac:prop:nonempty       (renamed from ac:claim:nonempty, D1)
%    ac:rem:structural
%    ac:rem:authority       (new in rev 2, C3+C4)
%    ac:rem:differentia
%    (referenced here, created elsewhere: ac:sub:countermeasure
%     [2.3, drafted], ac:sub:placement [2.2, drafted])
% ============================================================

%############################################################
%%##  FILE: sec2_3_the_countermeasure.tex
%############################################################

% ============================================================
% §2.3 — The countermeasure (Φ6)
% REVISION 2, against referee report (§F scope).
%
% MERGE NOTES — disposition of referee items:
%   ACCEPTED IN FULL: B (constraint unpacked + conditioning-event
%     connection), C1 (the conjunction, with the referee's own
%     caution honored — see below), C2 (declining has a shape),
%     C3 (positional, not beneficiary-based), C4 (enumerated
%     check), D1 (broken pointer fixed), D2 (grade tag dropped
%     from premise bracket), D3 (falsifier and Φ2 reconciliation
%     now named remarks), D4 (close's forward ref kept, the
%     falsifier's folded to a textual pointer).
%   C1 CAUTION — AGREED, with the reason stated: the inverse
%     relation ("constraint tends to fall as consequence rises")
%     is an unpriced empirical-ish heuristic that the conjunction
%     point does not need; adding it would create a new claim
%     needing its own falsifier for no structural gain. A first
%     draft of the scarcity remark drifted toward asserting it
%     and the sentence was cut. Left out.
%   D1: the cost paragraph now points at the RULE via placeholder
%     \ref{gs:sub:rule} (§7.2–7.3); gs:sub:fence remains pointed
%     at once, from the Φ2 remark, correctly. Both are forward
%     placeholders to be supplied when §7 is drafted.
%   LENGTH: ~1,400 words of body against the brief's 900–1,200.
%     The overage is entirely referee-requested additions (B +
%     C1–C4 ≈ 430 words). If the merge must reclaim, the C4
%     enumeration compresses to prose most cheaply — but it is
%     also the most checkable passage in the subsection, so cut
%     it last.
%   THE COUNTER GUARD IS LOAD-BEARING. \setcounter{premise}{5}
%     below makes this premise print Φ6. The paired
%     \setcounter{premise}{0} at the START OF §6 is owed (brief
%     §1.3; owed-edits item 1) and is NOT in this file. Without
%     the pair, Φ3 silently renumbers and every textual $\Phi3$
%     in §§6–7 breaks. Do not "clean up" either guard.
%   Owed edit #3 (gr:phi2 patch) drafted at the end of this file;
%     it travels with the section, alongside the hr:rem:notmistake
%     patch that 2.4 will carry.
%   enumitem CONFIRMED LOADED (register §7); dependency
%     discharged.
%   TRANSFER NOTE: a stale copy of this file (rev 1) reached the
%     referee and was reported byte-identical; the disk copy has
%     been rev 2 throughout (verifiable: header line 3, remarks
%     ac:rem:phi6falsifier / ac:rem:constraint / ac:rem:scarce /
%     ac:rem:phi2split present). Referee-report items B, C1–C4,
%     D1–D4 are all addressed herein; re-referee against THIS
%     copy.
% ============================================================

\subsection{The countermeasure}\label{ac:sub:countermeasure}

\S\ref{ac:sub:problem} closed on a silence: the argument fixes that some frozen question is held, and says nothing about which. The third property of Remark~\ref{ac:rem:threeprops} --- selectability --- is where the whole of \S2's positive content enters, and it enters as one premise. An agent cannot decline to stand somewhere. What can differ is the footing.

\setcounter{premise}{5} % GUARD — do not remove; pairs with \setcounter{premise}{0} at the start of \S6. See merge notes.
\begin{premise}[Minimal arbitrariness; \emph{meta-normative selection rule}]\label{ac:phi6}
Among the frozen questions an agent cannot avoid holding, prefer one whose answer is constrained by the agent's own existence over one whose answer is not.
\end{premise}

\begin{remark}[The falsifier]\label{ac:rem:phi6falsifier}
What would refute it, in its own vicinity, is one of two exhibits. Exhibit a frozen question whose answer is compelled by the agent's existence to a degree comparable with the candidate the section goes on to supply, and which delivers a \emph{different} answer to the Golden Question --- then $\Phi6$ selects a family rather than a point, and its use downstream fails even though the premise stands. Or show that ``constrained by the agent's own existence'' admits no ordering at all --- then the premise selects nothing and is empty rather than false. The first exhibit attacks the use; the second attacks the premise. Neither has been produced, and the section's downstream claims are priced against both.
\end{remark}

\begin{remark}[What the constraint is, and what it is not]\label{ac:rem:constraint}
The premise's whole content sits in one phrase, and the phrase should be said rather than leant on. The constraint is positional, not truth-conferring; nothing here reaches truth. Some frozen questions are such that the answer an agent holds is bound up with the agent's being in a position to hold it --- the holding is not a free selection among alternatives but something the position has already narrowed. Others are not: an agent could hold either answer without anything about its situation differing. The first kind carries information about the agent's circumstances; the second is a coin already flipped, where remembering how it landed adds nothing. In the paper's own terms: $\Phi6$ asks for the frozen question that behaves most like a conditioning event. The whole apparatus is built on conditioning, and the premise is the demand that one's grounding be a conditioned quantity rather than a free one --- which is what makes $\Phi6$ a resident of this framework rather than an import to it.
\end{remark}

\begin{remark}[Why the selection is scarce]\label{ac:rem:scarce}
Constraint alone is cheap, and degenerately so: \emph{am I asking something?} is fully constrained by the asking and entirely inconsequential. What $\Phi6$ asks for is a conjunction --- an answer heavily constrained by position \emph{and} consequential enough to carry what Condition 1 puts on frozen answers. The conjunction is what makes the premise non-trivial: candidates are not everywhere. And it is where a promise is cashed. Remark~\ref{ac:rem:admissible} held that a later candidate for the role would owe a demonstration of admissibility rather than an assumption of it; the conjunction is that debt restated as a criterion, and the demonstration follows once the machinery it consumes is in place.
\end{remark}

Two further things about $\Phi6$, and then the reconciliation a reader of \S\ref{gr:sec} will arrive already holding.

\emph{First, it is a selection rule over an unavoidable set, not a derivation of a value.} $\Phi6$ ranges over groundings, not over actions. It does not say the agent ought to have a grounding; Proposition~\ref{ac:prop:nonempty} established that the agent has one whether or not anything is said, and that its content is underdetermined by everything the agent can access. What $\Phi6$ adds is only this: arbitrariness admits of degrees, and among foundations the agent must occupy anyway, less is preferable to more. The ordering is ordinal and nothing stronger --- some frozen questions are more constrained by the agent's existence than others, and no figure will be put on the compulsion, here or later. And a preference for the least arbitrary footing must not be inflated into the claim that the least arbitrary footing is \emph{correct}. It is less arbitrary. Nothing in the paper can make it correct, and \S\ref{gr:sub:ceiling} is the standing reason: correctness of a grounding is exactly the kind of thing no finite observation certifies.

\emph{Second, no ought is derived from an is.} The objection every reader will have is that $\Phi6$ smuggles a norm out of a fact --- that ``your existence constrains this answer'' is an is, and ``therefore hold it'' is an ought, and the gap between them is the oldest gap there is. The answer is that the ought is not derived because it is not new. Condition 1 of \S\ref{ac:sub:problem} already places the agent under forced action; the practical necessity of standing on \emph{some} foundation is a structural fact about agency, in place before $\Phi6$ says anything, supplied by the predicament and not by the premise. $\Phi6$ operates entirely downstream of that necessity, on the one question the predicament leaves open --- which arbitrary foundation to occupy. A rule for choosing among footings an agent will occupy regardless derives no obligation to stand; the standing was never optional. That is what \emph{meta-normative} means in the premise's kind, and it is the whole of the premise's ambition.

\begin{remark}[What $\Phi2$ denies, and what $\Phi6$ claims]\label{ac:rem:phi2split}
Premise~\ref{gr:phi2} declares Golden reasoning inferential: it licenses conditional expectations, not duties. $Q$ is a conditional distribution; a distribution entails no obligation; Remark~\ref{gr:rem:notlongtermism} leans on exactly this, and the fence of \S\ref{gs:sub:fence} will lean on it again. If $\Phi6$ contradicted that, \S\ref{ac:sec} would be dismantling the paper's own load-bearing wall. The split: $\Phi2$ is a claim about what the mathematics \emph{outputs} --- and the mathematics outputs distributions, from which no duty follows, exactly as $\Phi2$ says. $\Phi6$ is a claim about what an agent does with a choice the mathematics never made for them --- the choice of grounding, which Proposition~\ref{ac:prop:nonempty} shows is forced on the agent by their situation and left open by everything the apparatus proves.

That the two premises do not meet can be checked by enumeration rather than taken from prose. The section's normative chain has four steps:
\begin{enumerate}[label=(\roman*)]
    \item an agent must hold some frozen answer (Proposition~\ref{ac:prop:nonempty});
    \item among the frozen questions the agent cannot avoid holding, $\Phi6$ prefers the constrained (Premise~\ref{ac:phi6});
    \item one candidate is maximally constrained --- discharged once the machinery it consumes is built;
    \item if $\Phi6$ is adopted and the candidate answered affirmatively, certain comparisons follow (\S\ref{gs:sub:rule}).
\end{enumerate}
At no step does a proposition of the form ``$Q$ implies that one ought to $x$'' appear. Step (i) is structural, step (ii) is adopted, step (iii) is a constraint claim, and step (iv) is conditional on both adoptions. The apparatus outputs distributions at every step and duties at none.

One further protection, stronger than the level-distinction. The grounding runs through a fact about the situation of \emph{whoever is asking}, not through the interests of anyone who will exist in the futures $Q$ describes. Remark~\ref{gr:rem:notlongtermism} answers an objection that requires such beings --- beneficiaries whose weight the apparatus stands accused of smuggling. Here there are none in the argument at all: no future being's interests appear at any of the four steps, and an agent alone in the last hour of a biosphere could run them unchanged.

An agent who declines $\Phi6$, then, is in conflict with nothing the paper proves, and the paper is, in exactly $\Phi2$'s sense, neutral about the declining. But declining has a shape, and \S\ref{ac:sub:problem} already named it. Declining $\Phi6$ in favor of a competing selection criterion is coherent, and the section holds no argument against it. Declining it and thereby standing outside the problem is not a position: indifference among frozen questions is not abstention from the choice but the choice made without a criterion --- a frozen answer of the least examined kind. The structure is Remark~\ref{ac:rem:differentia}'s, met one level up: quietism toward the selection is an instance of the selection.
\end{remark}

What the paper is not neutral about is the price of adopting. An agent who adopts $\Phi6$ \emph{and} answers the Golden Question affirmatively does thereby acquire commitments --- and the decision rule of \S\ref{gs:sub:rule} is those commitments made computable, which is a sentence that should be read as a cost rather than an advertisement. The accounting is this: $\Phi2$ keeps the apparatus obligation-free; $\Phi6$ is the single doorway through which obligation can enter; and everything on the far side of the doorway is the agent's adopted grounding, carried at \textbf{[$\Phi$]}, revocable by the agent and certified by nothing. Whether the doorway leads anywhere --- whether there exists a frozen question whose answer an agent's existence does compel --- is the question \S\ref{ac:sub:question} exists to answer.

% ------------------------------------------------------------
% OWED EDIT #3 (draft) — one-sentence patch to gr:phi2's
% surrounding text, to be applied at the Φ2 site in §6:
%
%   "Φ2 governs what the apparatus outputs; it is silent on the
%    separate question of how an agent grounds duties the
%    apparatus never issued, for which see the selection rule of
%    \S\ref{ac:sub:countermeasure}."
%
% Without this pointer, §2 and §7 read as contradicting Φ2
% (brief §1.5, Collision 1; owed-edits list, item 3).
% ------------------------------------------------------------

% ============================================================
% DELIVERABLE LISTS (brief §0) — revision 2
% ============================================================
%
% 1. CLAIMS REGISTER
%    - Φ6 (ac:phi6): among unavoidable frozen questions, prefer
%      one whose answer is constrained by the agent's own
%      existence. [Φ], kind: meta-normative selection rule.
%      Falsifiers in-place (ac:rem:phi6falsifier): (a) a
%      comparably-compelled rival delivering a different Golden
%      Answer (kills the use), (b) no ordering on the constraint
%      (kills the premise).
%    - The constraint is positional, not truth-conferring; the
%      constrained kind carries information about the agent's
%      circumstances, the unconstrained kind is a flipped coin.
%      [Φ] (new, per B)
%    - Φ6 asks for the frozen question that behaves most like a
%      conditioning event; the grounding is demanded to be a
%      conditioned quantity. [Φ] (new, per B — houses Φ6 in the
%      framework)
%    - The selection is a CONJUNCTION: heavy positional constraint
%      AND consequence sufficient for Condition 1's load; cheap
%      constraint is degenerate; candidates are not everywhere.
%      [Φ] (new, per C1; cashes ac:rem:admissible's promise)
%    - Inverse relation between constraint and consequence: NOT
%      claimed (referee's own caution, agreed; unpriced heuristic).
%    - Φ6 ranges over groundings, not actions; ordinal only; no
%      cardinality anywhere. [Φ]
%    - "Less arbitrary" is not "correct" (gr:sub:ceiling). [Φ]
%    - Not is–ought: the ought is pre-existing (Condition 1); Φ6
%      operates downstream, deriving nothing. [Φ]
%    - Φ2 split: Φ2 governs outputs, Φ6 governs occupancy;
%      checkable by the four-step enumeration; no step has the
%      form "Q implies one ought to x". [Φ] (enumeration new,
%      per C4)
%    - The grounding is positional, not beneficiary-based: no
%      future being's interests appear at any step; this protects
%      gr:rem:notlongtermism independently of the level
%      distinction. [Φ] (new, per C3)
%    - Declining Φ6 for a competing criterion is coherent and
%      unopposed; declining it to stand outside the problem is
%      not a position — a frozen answer of the least examined
%      kind, ac:rem:differentia one level up. [Φ] (new, per C2)
%    - Cost stated: Φ6 + affirmative Golden Answer ⇒ acquired
%      commitments; §7's rule is those commitments made
%      computable; neutrality is about adoption, not price. [Φ]
%
% 2. CROSS-REFERENCES USED
%    \S\ref{ac:sub:problem} ×2, Remark~\ref{ac:rem:threeprops},
%    Remark~\ref{ac:rem:admissible},
%    Proposition~\ref{ac:prop:nonempty} ×3,
%    Remark~\ref{ac:rem:differentia},
%    Premise~\ref{gr:phi2}, Remark~\ref{gr:rem:notlongtermism}
%    ×2, \S\ref{gr:sub:ceiling}
%    \S\ref{ac:sub:question} — forward, ONCE, final sentence
%      (D4: falsifier's instance folded to "the candidate the
%      section goes on to supply"; ac:rem:scarce and step (iii)
%      folded to "once the machinery it consumes is built/in
%      place" — all order-safe after the Tendency/Question swap).
%      RESOLVED by the 2.4 draft
%      (sec2_4_the_golden_question.tex), which also cashes
%      ac:rem:scarce's stated debt (both halves, explicitly).
%    \S\ref{gs:sub:fence} — forward to §7.5, once (Φ2 remark)
%    \S\ref{gs:sub:rule} — forward to §7.2–7.3, twice (step (iv)
%      + cost paragraph); NEW placeholder, per D1
%
% 3. NEW LABELS CREATED
%    ac:sub:countermeasure   (resolves 2.1's forward reference)
%    ac:phi6
%    ac:rem:phi6falsifier    (new in rev 2, per D3)
%    ac:rem:constraint       (new in rev 2, per B)
%    ac:rem:scarce           (new in rev 2, per C1)
%    ac:rem:phi2split        (new in rev 2, per D3)
%    (referenced here, created elsewhere: ac:sub:question [2.4,
%     drafted]; still open: gs:sub:fence [§7.5], gs:sub:rule
%     [§7.2–7.3])
%
% 4. COUNTER GUARD: \setcounter{premise}{5} present; paired
%    \setcounter{premise}{0} at §6's start remains an owed edit.
% ============================================================

%############################################################
%%##  FILE: sec2_5_the_golden_tendency.tex
%############################################################

% ============================================================
% §2.5 — The Golden Tendency
% REVISION 3 — adopts the five convergence corrections from the
% consolidated register (§2 items 2.1–2.5), on top of rev 2
% (referee §G) and rev 2.1 (post-check fixes).
%
% MERGE NOTES — convergence adoptions this revision:
%   2.1  ac:prop:filter hypothesis strengthened: the ordering is
%        now on HAZARDS, not survival functions. Monotone-in-t
%        dominance needs the likelihood ratio S_θ(t)/S_θ'(t)
%        nondecreasing in t, which is exactly hazard ordering;
%        survival-function ordering at each t does not deliver
%        it. Statement and one-line justification amended.
%   2.2  ac:prin:tendency antecedent made load-bearing: the
%        conditioned-from region is now named — the region on
%        which total hazard tends to zero, existing exactly when
%        ρ > 0 by gr1:prop:trichotomy(i).
%   2.3  ac:rem:regimes trapped case: "false" (which contradicted
%        the principle two paragraphs above) replaced — the
%        antecedent FAILS in trapped, no vanishing-hazard region
%        exists, the principle does not fire rather than misfires.
%   2.4  C3 SETTLED (my rev-2.1 narrowing superseded by fact):
%        the background channel \eqref{tc:eq:B} IS a sum of many
%        independent random exponentials (Appendix app:background,
%        partition function with freezing). The 2024 aggregation
%        argument described a component and took it for the
%        whole. "Licensed nothing" replaced accordingly; the
%        rev-2.1 TODO is discharged and removed.
%   2.5  Hardcoded \S5 → \S\ref{gr1:sec}.
% Standing from earlier revisions: derivation cited not rebuilt;
%   2024 equations not transferred; Lindy priority conceded;
%   Tardigrade cut (gr:sub:crossdomain if wanted);
%   gr1:thm:twocond declined (2.4.2 owns should/can, referee A6
%   accepted); ensemble-not-trajectory fence per report B;
%   \E per house macro; \eqref{} for equations per watch item;
%   A1 opening line is original, substitutable at merge.
% ============================================================

\subsection{The Golden Tendency}\label{ac:sub:tendency}

\S\ref{ac:sub:countermeasure} closed on whether the doorway leads anywhere; the candidate that answers is named in \S\ref{ac:sub:question}, and its credentials will consume machinery this subsection must build first.

A late observer does not see what was made; they see what is left. Survivorship is not a force --- it selects nothing, prefers nothing, and does not act; it is what remains when a population is observed at a later time than it was generated. This subsection states what remains, in two parts, and then says how much of it survives contact with the scoping the paper already imposes on itself.

The first part is arithmetic.

\begin{proposition}[The Golden Filter; \textbf{[P]}]\label{ac:prop:filter}
Let a cohort of biospheres originate together with types $\theta \sim \nu$, each type carrying hazard rate $\lambda_\theta(\cdot)$ and survival function $S_\theta(t) = \exp\!\big({-\int_0^t \lambda_\theta(\tau)\,d\tau}\big)$. Order types by efficiency so that $\lambda_\theta$ is pointwise nonincreasing in $\theta$. Then the composition of the cohort still standing at $t$ is the originating composition reweighted by survival,
\[
    dN_t(\theta) \;\propto\; S_\theta(t) \, d\nu(\theta),
\]
and it stochastically dominates $\nu$, with the dominance nondecreasing in $t$.
\end{proposition}

The proof is the statement: the number of type-$\theta$ members still present is the number that started times the probability one lasted, and normalizing gives the reweighting; the dominance is monotone in $t$ because the likelihood ratio between a more and a less efficient type, $S_\theta(t)/S_{\theta'}(t) = \exp\!\int_0^t (\lambda_{\theta'} - \lambda_\theta)\,d\tau$, is nondecreasing in $t$ exactly when the hazards are ordered --- which is why the hypothesis orders hazards rather than survival functions. Nothing is hypothesized beyond the ordering. All of the content of the Filter is in what it gets applied to, and the deflation is worth insisting on, because the Filter is routinely mistaken for a substantive posit when it is a bookkeeping identity about two different distributions that were never the same object.

The galaxy is what the identity looks like when it is allowed to run. Relax the Rare-Earth hypothesis to galactic scale, let biospheres undergo abiogenesis independently and far apart, each drawing a hazard rate at random. The fragile fizzle; the durable persist and fizzle later. Watch the populations rather than the biospheres: the durable accumulate, because arrivals outpace departures, while the fragile turn over in waves and hold roughly constant. Then the observation recurses --- among the durable, the most durable stand in exactly this relation to the rest --- and the skew propagates into the tail as far as the tail extends.\footnote{Strictly, the galaxy is the renewal rather than the cohort version: origination is ongoing, so the standing population is a convolution of the origination rate with survival rather than a simple reweighting of a single $\nu$. Proposition~\ref{ac:prop:filter} is exact for the cohort; the galaxy inherits the same concentration conclusion asymptotically by the same mechanism, and nothing below turns on the difference.}

The Filter is best named by what it sits between. On one side is the distribution of \emph{generation} --- what abiogenesis produces, the Cosmic Nexus. On the other is the distribution of \emph{abundance} --- what a late observer finds. The Filter is the map from the first to the second, and the map is nothing but conditioning on having lasted. In the paper's vocabulary this is the pair $\PP$ and $Q$ of Definition~\ref{gr:def:gg}: General and Golden, generation and abundance, the two objects the galaxy separates by hand.

\begin{remark}[The Filter accumulates; it does not aim]\label{ac:rem:noaim}
The galaxy contains no striving. No biosphere in it is trying to be durable, no selection pressure is applied to any individual, and no biosphere's prospects improve by so much as an instant on account of the Filter. The durable accumulate because they were durable, and that is the whole mechanism. The paper's anti-teleological commitment (Remark~\ref{gc:rem:teleology}, and the calibration attached to Optimation's name in \cite{Nexus}) is not a caution appended to this passage but a description of it: \emph{is populated by}, \emph{accumulates}, \emph{is selected} are exact; \emph{aims}, \emph{seeks}, \emph{progresses toward} are not, and every teleological reading of what follows is an error introduced by the reader rather than a claim inherited from the model.
\end{remark}

The apparatus that makes this precise is already in the paper.\footnote{Provenance. A 2024 draft derived the Filter by hand. It defined survival efficiency as a quantile $x$ of the survival probability $s(t,x)$ at horizon $t$, introduced a ``Survival Nexus'' $N_S \equiv \{S > t\}$, and computed the efficiency expected under it as $\E[X_t] = \int_0^1 s\,x\,dx \big/ \int_0^1 s\,dx$, exhibiting toy hazard families under which $\E[X_\infty]$ rose to $3/4$ and $5/6$ from the unconditioned $1/2$. What that construction anticipated: the object (a survival-conditioned reweighting of efficiency), the direction (the conditioned expectation exceeds the unconditioned one, and increases with the horizon), and the two-Nexus framing (generation against abundance), which survives here intact. What it did not anticipate, and could not have: that the limit exists at all rather than merely trending; that the limiting object is a Doob transform (Theorem~\ref{sh:thm:Q}) rather than a family of reweightings indexed by $t$; that the approach is clockless; and that the result is scoped by anatomy in the way Remark~\ref{ac:rem:regimes} records. Its measured descendant is the anatomy selection of \S\ref{gr:subsub:anatsel}, where the crude $3/4$ becomes $0.75 \to 0.84 \to 0.91$ across horizons. The 2024 equations are superseded and are not reproduced.}
Three identifications carry it. The Filter \emph{is} the survivorship interpolation $\PP^{(t)} \to Q$ of \S\ref{sh:sub:interp}, whose limit Theorem~\ref{sh:thm:Q} establishes. That high efficiency accumulates \emph{is} Lemma~\ref{sh:lem:monotone}: every bias field is a nondecreasing function of the state, so conditioning on survival to any horizon is an upward tilt, at every horizon, with no further assumption. And the rate at which abundance departs from generation is the goldening rate $s/4t$ of \eqref{gr:eq:rate}, with the persistence exponent entering as the linear-response coefficient.

So far the ensemble has been a population of biospheres; the Tendency reweights a different one --- the continuations of a single history, held in disagreeing superposition. The substitution is free, because survivorship reweighting does not consult what indexes its ensemble: Proposition~\ref{ac:prop:filter} multiplies each member's weight by that member's survival function and renormalizes, indifferent to whether the members are separated in space or in superposition. The one assumption the substitution would cost --- that continuations carry no term coupling them --- is discharged by construction at \S\ref{sh:sub:interp}, where the object is a conditioning of a single path measure, so no such term arises to be assumed away.

\begin{principle}[The Golden Tendency; \textbf{[S]}]\label{ac:prin:tendency}
Let a biosphere have favorable technoanatomy, $\rho > 0$. Then the state space contains a region on which the total hazard tends to zero (Proposition~\ref{gr1:prop:trichotomy}(i)), and as $t$ grows, the futures in which the biosphere has descendants at horizon $t$ are drawn increasingly from that region.
\end{principle}

\begin{remark}[The claim is about an ensemble, not a path]\label{ac:rem:notrajectory}
This is the section's central fence, and it is stated here in the form a reader arriving from \S\ref{gr:subsub:violence} will already expect. The Tendency does not say that a biosphere improves. Nothing here moves any system along any path; $\rho$ does not change in the fixed model, and what conditioning does is reweight an ensemble. The selection is Darwinian in form rather than developmental --- which is exactly what the folk version of the technological-adolescence argument gets wrong, and what Cooper \cite{Cooper19} is right to object to in objecting that no evidence supports a trajectory of moral improvement. The ensemble form is not a weaker version of the trajectory form. It is the version that survives that rebuttal, and it is the reason the section can afford the claim at all.

Where, then, would a trajectory reading live? Not nowhere, and not forbidden --- unbuilt. It becomes available only where technoanatomy carries its own dynamics, so that anatomy may escape upward and fail to return (\S\ref{gr1:subsub:fluct}, Corollary~\ref{gr1:cor:recurrentanatomy}); those dynamics are \textbf{[O]}. A reader who wants improvement is requesting a piece of the framework that does not yet exist --- and which Remark~\ref{hr:rem:forcesprocess} argues is required for independent reasons anyway.
\end{remark}

\begin{remark}[What the antecedent costs, and who is not entitled to it]\label{ac:rem:regimes}
The antecedent is anatomy, not regime, because the Tendency holds in two of the three cells that \S\ref{gr:subsub:regimes} distinguishes and they are exactly the cells with $\rho > 0$ --- the cells in which the vanishing-hazard region the statement names exists at all. In the \emph{critical} regime it holds in the form stated, and the machinery above is the machinery of it: the tail is polynomial, the limit exists, and the reweighting is the whole content. In the \emph{viable} regime it holds and stops being interesting, because survival has positive probability outright and conditioning on it is an ordinary conditional rather than a limit statement. In the \emph{trapped} regime the antecedent fails outright: $\rho \le 0$ by Theorem~\ref{gr1:thm:trapped}, and by Proposition~\ref{gr1:prop:trichotomy} no region of vanishing hazard exists to be drawn from --- survivors concentrate on a single path, and the principle does not misfire there; it does not fire. So the Tendency is not a theorem, and calling it one would be the section's most attractive error. What is proved is that conditioning tilts upward (Lemma~\ref{sh:lem:monotone}) and that the limit exists where the tail is polynomial (Theorem~\ref{sh:thm:Q}). That descendants' futures are \emph{in fact} so drawn requires in addition that this biosphere has $\rho > 0$ --- and by \S\ref{gr:sub:ceiling} no one is entitled to that antecedent from finite observation. The Tendency is therefore stated by an author who cannot certify its condition. That is the normal case rather than a defect: the conditional is what the mathematics supports, and the ceiling is what forbids upgrading it.
\end{remark}

Two things remain before the subsection closes.

\emph{The Lindy relation is prior work, and the priority is not the paper's.} That civilizational longevity follows a Lindy relation is established in \cite{BC21}, \cite{BG24}, \cite{Kip20}, and \cite{Ord23}; \S\ref{gr:sub:consequences}(c) states this and states that the paper's only addition is the derived exponent. The Tendency inherits that concession without amendment. The 2024 material reached the Lindy relation by a different route --- aggregate many independent hazards and the aggregate flattens, so near-constant hazard is what many hazards become rather than an assumption imposed on them --- and that argument is sound about the object it addresses. The model, moreover, contains that object: the background channel of \eqref{tc:eq:B} is exactly a sum of many independent random exponentials, treated in Appendix~\ref{app:background} as a partition function with freezing. Where the 2024 material erred is that a component was described and taken for the whole. This paper does not diffuse $h$; it diffuses $-\tfrac{d}{dt}\log h$, one integration further out, which is what makes $\log h$ a $C^1$ object, the state degenerate and $3/2$-self-similar, and the survival tail polynomial rather than exponential. The noise placement is the locating novelty flagged in the relation-to-literature passage of \S\ref{gr1:sec}. The tail is a property of the evolution, and the aggregation argument --- correct about the background channel --- never reaches it.

\emph{Contingency is not a competitor.} The standing reply to the Tendency is that it makes descendants sound determined --- that a claim of this shape cannot hold of beings who choose. The gas law is the familiar answer: a statement about a volume is not falsified by the fact that no molecule obeys it, and the order at the aggregate is not a constraint imposed on any constituent from outside. But the framework's own answer is sharper, and it is not an analogy. The Tendency is a statement about a measure over futures, not a prediction about any one of them. Contingency is not what the Tendency competes with; contingency is what the measure is defined over. An objector who establishes that the future is radically open has established the existence of the object the Tendency reweights, and has said nothing whatever about the reweighting. The two claims cannot conflict, because they are not claims about the same thing.

What the Tendency does not supply is an agent-side counterpart. The mathematics describes a reweighting; whether a biosphere's conscious effort has an analogous structure, and what a decision rule built on it would compute, is taken up in \S\ref{gs:sub:initiative} and not here. A related caution: the observation that a biosphere might generate a sub-component targeting its own extinction pressures --- and that we are, as it happens, such a sub-component --- reads as a coincidence too large to be an accident. It is not an exclamation but an inference, and the form it takes is the contrapositive of the Theorem of Aimless Optimation \cite{Nexus}: unequal weights under two Nexae entail non-independence from $H_N$. Stated that way it licenses a recorded correlation and nothing more.

% ============================================================
% DELIVERABLE LISTS (brief §0) — revision 3
% ============================================================
%
% 1. CLAIMS REGISTER
%    - Golden Filter (ac:prop:filter): cohort with HAZARD-ordered
%      types; standing composition = survival-reweighted ν;
%      stochastic dominance nondecreasing in t via monotone
%      likelihood ratio. [P] — hypothesis strengthened per
%      convergence 2.1; proof-is-statement retained with the
%      one-line MLR justification.
%    - Galaxy is the renewal version; same conclusion
%      asymptotically. [P] (footnote)
%    - Filter = map from generation (ℙ) to abundance (Q);
%      conditioning on having lasted (gr:def:gg). [S]
%    - The Filter accumulates and does not aim. [Φ]
%    - Three identifications: interpolation (sh:sub:interp,
%      sh:thm:Q); upward tilt (sh:lem:monotone); rate s/4t
%      (\eqref{gr:eq:rate}). Inherited tags.
%    - Index substitution free; exchangeability discharged at
%      sh:sub:interp. [P] (body prose)
%    - Golden Tendency (ac:prin:tendency): given ρ > 0, a
%      vanishing-total-hazard region exists
%      (gr1:prop:trichotomy(i)) and descendant-futures are drawn
%      increasingly from it. [S] — antecedent now LOAD-BEARING
%      per convergence 2.2.
%    - Ensemble-not-path fence; trajectory reading unbuilt, [O]
%      (gr1:subsub:fluct, gr1:cor:recurrentanatomy).
%    - Trichotomy: critical (holds, contentful), viable (holds,
%      uninteresting), trapped (ANTECEDENT FAILS — no
%      vanishing-hazard region; principle does not fire).
%      Per convergence 2.3; "false" removed.
%    - Nobody entitled to ρ > 0 from finite observation. [Φ]
%    - Lindy is prior work; paper adds the exponent only.
%    - C3 SETTLED per convergence 2.4: the 2024 aggregation
%      argument correctly describes the background channel
%      \eqref{tc:eq:B} (app:background, partition function with
%      freezing) and took the component for the whole; the tail
%      is a property of the evolution of −d/dt log h. [S]
%    - Contingency is what the measure is defined over. [Φ]
%    - Brain/sub-component: recorded correlation via Aimless
%      Optimation contrapositive. [S]
%    - 2024 provenance recorded (unchanged).
%
% 2. CROSS-REFERENCES USED
%    Definition~\ref{gr:def:gg}, \S\ref{sh:sub:interp},
%    Theorem~\ref{sh:thm:Q}, Lemma~\ref{sh:lem:monotone},
%    \eqref{gr:eq:rate}, \S\ref{gr:subsub:regimes},
%    Proposition~\ref{gr1:prop:trichotomy} ×2,
%    Theorem~\ref{gr1:thm:trapped}, \S\ref{gr:sub:ceiling},
%    \S\ref{gr:sub:consequences}, \S\ref{gr:subsub:anatsel},
%    Remark~\ref{gc:rem:teleology}, \S\ref{gr:subsub:violence},
%    \S\ref{gr1:subsub:fluct},
%    Corollary~\ref{gr1:cor:recurrentanatomy},
%    Remark~\ref{hr:rem:forcesprocess}, \S\ref{gr1:sec},
%    \eqref{tc:eq:B}, Appendix~\ref{app:background}
%    \cite{Nexus} ×2; \cite{BC21}, \cite{BG24}, \cite{Kip20},
%    \cite{Ord23}
%    \S\ref{gs:sub:initiative} — forward to §7.1 (label per the
%      fixed §7 map, register item 3.3)
%    NOT USED, with reason: gr1:thm:twocond — 2.4.2 owns
%      should/can (referee A6 accepted).
%
% 3. NEW LABELS CREATED
%    ac:sub:tendency
%    ac:prop:filter
%    ac:rem:noaim
%    ac:prin:tendency
%    ac:rem:notrajectory
%    ac:rem:regimes
% ============================================================

%############################################################
%%##  FILE: sec2_4_the_golden_question.tex
%############################################################

% ============================================================
% §2.4 — The Golden Question
% REVISION 2, against the verification report (§B priority, plus
% §D encoding and §E items touching this file).
%
% MERGE NOTES — disposition of report items:
%   B ACCEPTED IN FULL, and the referee's diagnosis was right in
%     a way my own propagation note concealed: the ensemble
%     discipline had been applied to the VOCABULARY and not to
%     the argument's SUBJECT. The nihilist counterexample is
%     decisive against the rev-1 form — bare asking fails
%     ac:rem:constraint's own test. Relocated to the ensemble:
%     the asymmetry now comes from the answer being ACTED ON at
%     biosphere scale, through ac:prop:filter.
%   CONVERGED with the referee's 2.4: label ac:prop:forcing
%     adopted; promoted to a proposition environment per the
%     recommendation. The referee's own flagged weak joint —
%     whether acting on the negative NECESSARILY lowers
%     persistence — is carried here as a near-analytic middle
%     step with its falsifier stated in the fence, not as a
%     necessity claim. That is the honest strength: near-analytic
%     is not analytic, and the exhibit that would break it is
%     named.
%   KEPT from rev 1 per the report: the conditioning setup and
%     "into every system in which the asking is placed, the
%     asker is placed as well" — explicitly re-purposed as setup
%     that establishes the question's subject and no more.
%   PAYOFF adopted: the selection leans on ac:prop:filter (every
%     cell), not ac:prin:tendency (needs ρ>0) — so the Golden
%     Answer's selection does not inherit the Tendency's
%     antecedent, and 2.4.2's should/can separation now falls
%     out of the mechanism. One sentence added in 2.4.2 to
%     receive it.
%   E items: the scientific-method argument CUT from 2.4.4 (it
%     ran near-verbatim twice; the pointer remains and 2.2 keeps
%     the argument); $x^*$ → $x_\ast$ per master notation;
%     encoding per §D throughout.
%   OWED EDIT #2 (hr:rem:notmistake patch) retained at the foot;
%     still to be fitted to the remark's actual text.
% ============================================================

\subsection{The Golden Question}\label{ac:sub:question}

\subsubsection{The statement}\label{ac:ssub:statement}

$\Phi6$ selects nothing unless a candidate exists. The candidate:

\begin{quote}
    \emph{Should our biosphere live forever?}
\end{quote}

The contrast that locates it is with the cogito, and one sentence carries it: the cogito secures an existence claim and grounds nothing further, while the Golden Question is chosen for its downstream reach --- its interest is entirely in what it grounds. (The historical placement beside Descartes is a separate and stronger claim, and it is \S\ref{ac:sub:placement}'s to make at its price.)

Two debts fall due here, and they are the two halves of Remark~\ref{ac:rem:scarce}'s conjunction. The \emph{consequence} half is the admissibility demonstration that Remark~\ref{ac:rem:admissible} put on account: the question is not paradoxical, since either answer is coherent and neither dissolves the asking; and it is not inconsequential, since its answer bears on every forced selection Condition 1 names --- a biosphere that should persist forever and one that need not are different backdrops for every decision an agent under forced action makes. That reach is the whole of the not-Descartes contrast restated as a property. The \emph{constraint} half is the forcing argument.

Begin with what the asking itself supplies, which is less than it appears. By some future point the biosphere persists or it does not; the asking is an act of a persisting biosphere's sub-component, so wherever the asking occurs, persistence to date has already been conditioned on. The asker does not survey the dichotomy from above it: into every system in which the asking is placed, the asker is placed as well. This establishes what the question is about --- the asker's own situation rather than a detached one --- and it establishes no more. The bare fact of asking does not narrow the answer: an asker holding the negative sits on the same branch, with nothing about their situation differing, and by Remark~\ref{ac:rem:constraint}'s own test a position of that kind is a coin already flipped. The narrowing needs a mechanism, and the mechanism is the one the section already owns.

\begin{proposition}[Forcing; \textbf{[$\Phi$]}]\label{ac:prop:forcing}
In the ensemble of biospheres in which a frozen answer to the Golden Question is held and acted upon at biosphere scale, survivorship reweights toward the affirmative: the weight carried by continuations in which the negative is held falls with the persistence of its host, and the affirmative incurs no corresponding penalty.
\end{proposition}

\noindent (One placement for later: in the vocabulary of
\S\ref{gs:sub:nodes}, the frozen answer is the paradigm
\emph{universe-sourced commitment node} --- the existence proof
consumed there.)

The argument is three steps, each already established elsewhere. \emph{Consequence.} By the admissibility just discharged, the answer bears on practice: an answer with no downstream difference would be inconsequential and inadmissible, so a biosphere-scale frozen answer organizes what its holder does. \emph{The middle step.} A biosphere whose practices are organized around \emph{this biosphere should not persist} does less to persist. The step is near-analytic --- persistence-directed practice is what the affirmative organizes, and its absence is what the negative licenses --- but near-analytic is not analytic, and the fence below states what would break it. \emph{Reweighting.} Proposition~\ref{ac:prop:filter} does the rest: the standing ensemble thins in hosts of the acted-on negative and does not thin in hosts of the affirmative, so the answer found held in persisting biospheres skews affirmative, increasingly in $t$. The asymmetry is real because it turns on what happens when the answer is \emph{acted on}, not on the bare fact of asking.

\begin{remark}[A stronger sense in which the question is the asker's]
\label{ac:rem:indecomposable}
One consideration bears on \emph{which question} $\Phi6$ selects rather
than on which answer survivorship leaves, and it raises the price of the
first falsifier of Remark~\ref{ac:rem:phi6falsifier}: a rival frozen
question, comparably constrained, delivering a different answer.

The companion argues that a reasoner cannot detach from its biosphere
even notionally. The dependence runs past life support to
contextual-epistemic grounding, so the attempted separation dissolves the
Aethic truth conditions the reasoner's own statements rested on, in the
way imagining masslessness dissolves the frame under Mach's principle
\cite{Nexus, Barbour1995}; the chain of inferences by which anything is
fixed terminates in the reasoner's being a component of the biosphere.

If that holds, the Golden Question's subject is not merely correlated
with the asker's existence --- it is what the asker's capacity to frame
\emph{any} question is constituted by. That is a stronger reading of
``constrained by the agent's own existence'' than the positional one
argued above, and a rival question would have to match it.

Two limits, since the argument is easy to over-draw. It establishes a
dependence and not an asymmetry: the negative answer is framed with the
same constituted apparatus, so nothing here discriminates between
\emph{answers}, and the forcing argument's reliance on survivorship is
not relieved. And it is the companion's argument at the companion's
grade, imported here only to make the falsifier harder to satisfy rather
than to close it.
\end{remark}

\begin{remark}[What agency is, in this framework, and why it is not a
primitive]\label{ac:rem:agencyisnot}
One background commitment should be surfaced here, since it governs how
every reference to an agent in this section is to be read. Agency is not
a primitive of the framework: the companion's act-robustness doctrine
forbids it \cite{Aethus}. What stands in for it is \emph{contingency that
continues to move} --- an attribute whose weight in the Aethic
progression keeps shifting across successive Aethae without settling to
stated. That is the difference between a choice still open and a fact
already in the books: the object in the sky, once its distance is
resolved, is stated and one-and-done; whether to attempt the difficult
thing is a weight that the progression keeps revising until it does not.
An agent, on this reading, is a system some of whose relevant attributes
are of the second kind, and ``the agent chooses'' is shorthand for the
weight settling. Nothing in this section requires more than that, and
Remark~\ref{ac:rem:bridge} below is what keeps the shorthand from being read
as causal purchase.
\end{remark}

\begin{remark}[Position, not causal purchase]\label{ac:rem:bridge}
$\Phi6$ is a rule for an agent; Proposition~\ref{ac:prop:forcing} is about an ensemble, and the bridge must be stated rather than implied. An individual's holding does not move their biosphere's hazard --- one agent among a biosphere's many is not a lever on it. What the individual has is position: they are asking from inside a biosphere that has persisted, and by Proposition~\ref{ac:prop:filter} that is evidence about which answer the biospheres one can be asking from tend to hold. The relation is evidential, not causal --- the relation the record bears to the Accordance Principle throughout, and the one Remark~\ref{gc:rem:teleology} insists on when it refuses to let a conditioning event act. The constraint $\Phi6$ asks after is positional in exactly this sense: what the asker's existence narrows is not which answer they can hold --- the nihilist holds the negative unimpeded --- but which answer their existence is evidence for at the scale the question itself concerns. No situational feature is invoked beyond existence within a persister, and the narrowing runs through the same survivorship that structures the whole paper.
\end{remark}

\begin{objection}[Why call the Golden Point good?]\label{ac:obj:good}
Grant everything. Survivorship selects the affirmative; biospheres
organized around the negative thin out. Why does that make the Golden
Point \emph{good}? And if this work will not supply an objective
morality, why is the conclusion not simply that persisting forever is bad
and a biosphere ought to stop?
\end{objection}

\noindent The objection is met by conceding it, and the concession costs
nothing that was ever claimed.

\emph{Nothing here says the Golden Point is good.}
Remark~\ref{ac:rem:forcing} states the limit already: what is
established is positional stability, not truth, and reading a selection
effect as a demonstration is the error Remark~\ref{gc:rem:teleology} guards
against. $\Phi6$ asks for the \emph{least arbitrary} grounding among
those an agent must occupy anyway, not the correct one --- and
\S\ref{gr:sub:ceiling} is the standing reason no argument could supply
the latter, since correctness of a grounding is exactly what no finite
observation certifies. A reader who concludes that the Golden Point is
not good has contradicted nothing in this work.

\emph{The second half has an answer that asks the objector to concede
nothing about value.} One may hold that a biosphere ought to stop.
Proposition~\ref{ac:prop:forcing} does not refute the position; it
locates it. Held and acted on at biosphere scale, that answer removes its
own holder, so the ensemble still holding it thins with time. The
position is not false --- it is \emph{unstable}, and the difference is
the whole of what this subsection claims. An answer may be perfectly
coherent and still not be one the biospheres one can be asking from tend
to hold.

\begin{remark}[What the forcing argument shows, and what it cannot]\label{ac:rem:forcing}
The proposition is a conditioning claim of exactly the paper's usual shape, and it inherits the paper's usual discipline. It fixes what survivorship selects among persisting biospheres, and hence what an asker is positioned to find; it does not fix what is the case. Selection effects do not act (Remark~\ref{ac:rem:noaim}), and being positioned inside the selection is not evidence that the selected answer is true --- the constraint is positional, not truth-conferring, which is $\Phi6$'s own scoping (Remark~\ref{ac:rem:constraint}) applied to $\Phi6$'s best candidate. What would refute the argument, in its own vicinity: exhibit a biosphere organizing its practice around the negative while doing no less to persist --- the middle step is near-analytic, not analytic, and this is its exhibit; or exhibit a coupling by which affirmative-organized practice lowers persistence comparably, restoring the symmetry the argument denies. And the standing rival exhibit of Remark~\ref{ac:rem:phi6falsifier} remains open: a frozen question comparably compelled that delivers a different answer would leave this argument intact and its use void.
\end{remark}

One consequence is worth its own paragraph. The selection just argued leans on Proposition~\ref{ac:prop:filter}, which reweights in every cell of the trichotomy, and not on Principle~\ref{ac:prin:tendency}, which needs $\rho > 0$ to name a destination. The Golden Answer's selection therefore does not inherit the Tendency's antecedent: the answer is selected wherever there is survivorship at all, and $\rho > 0$ governs only whether what it affirms is \emph{available}. That is the separation of \S\ref{ac:ssub:shouldcan} arriving from the mechanism rather than by stipulation.

\subsubsection{Should and can}\label{ac:ssub:shouldcan}

The question asks whether the biosphere \emph{should} live forever, and the affirmative is what survivorship selects wherever an asker can stand. Whether it \emph{can} is a separate matter, and Theorem~\ref{gr1:thm:twocond} is the separation --- a separation the mechanism has already exhibited, since the selection ran through Proposition~\ref{ac:prop:filter}, which operates in every cell, while availability is $\rho$'s business alone. In the trapped regime the purchase is exactly nothing: $P(\text{viable}) = 0$ identically (Theorem~\ref{gr1:thm:trapped}), so an agent there is positioned inside the selection of the Golden Answer and inhabits a biosphere that cannot act on it. The answer is yes and the thing is unavailable, and both halves of that sentence stand at full strength.

This is the paragraph that keeps the section from reading as optimism, and it should be taken at the calibration the paper has already set: we are not practically Golden to any interesting tolerance, and Golden reasoning does not claim we are (Remark~\ref{gr:rem:weakpremise}). Nor does the forced answer sort the options it condemns: Remark~\ref{hr:rem:concession} shows the extinctionist future and the settled-at-$x_\ast$ future carrying viability zero alike, so affirming the Golden Answer while settling is not a middle path but the same zero by a longer route. What follows from a forced \emph{should} and an uncertified \emph{can} is not despair and not reassurance; it is Remark~\ref{hr:rem:forcesprocess}'s conclusion that anatomy must be a process --- the one response that is neither.

\subsubsection{Meaning, at its price}\label{ac:ssub:meaning}

The meaning claim, stated at exactly its strength: for an agent who adopts $\Phi6$, the slot that philosophy labels \emph{meaning} is filled, and filled non-arbitrarily. That is the whole of it. The claim is downstream of an adopted premise, carried at \textbf{[$\Phi$]}, revocable with the adoption, and certified by nothing --- the far side of \S\ref{ac:sub:countermeasure}'s doorway, described rather than promoted.

Stating the claim's compatibility with Remark~\ref{hr:rem:notmistake} is what earns the paragraph. That remark disclaims a verdict on human meaning-making issuing \emph{from a source outside human self-interpretation} --- and resists, at the paper's grade, the temptation to read the framework as supplying one. Nothing here supplies one. The grounding runs through $\Phi6$, which is adopted rather than output; the answer is positional rather than true; and an agent who declines the premise receives no verdict of any kind, favorable or otherwise. What the disclaimed reading wanted was meaning underwritten from outside; what \S\ref{ac:sec} delivers is an agent's unavoidable slot filled from inside, by a criterion the agent holds. Those are different kinds of claim, and only the first is disclaimed. The unconditional form --- that this is what meaning \emph{is}, search concluded --- would collide with Remark~\ref{hr:rem:notmistake} and $\Phi2$ simultaneously, and it is not made.


This subsection was written before the apparatus that pays it existed;
the reader should know the debt is now discharged. The meaning here
purchased at its stated price is developed in full as the third route
--- inferred meaning, with its two-tier claim, its literature, and its
defenses --- at \S\ref{gs:sub:thirdroute}, and this early statement
should be read as that development's reservation, not its substitute.

\subsubsection{An alternative phrasing}\label{ac:ssub:phrasing}

One rephrasing is worth naming and not running: that Golden reasoning stands to agency's curse as the scientific method stands to skepticism, with \emph{least circular} and \emph{minimally arbitrary} as the same selection in two vocabularies. The argument that earns the parallel, and the burden the parallel carries, are \S\ref{ac:sub:placement}'s in full.

% ------------------------------------------------------------
% OWED EDIT #2 (draft) — sentence to be APPENDED to
% hr:rem:notmistake, fitted to its actual text at merge (I have
% only the brief's quoted fragments):
%
%   "The disclaimer concerns an unconditional verdict issuing
%    from outside human self-interpretation; the Φ6-conditional
%    claim of \S\ref{ac:ssub:meaning} is of a different kind ---
%    an adopted grounding, revocable and certified by nothing ---
%    and the two are compatible for the reasons given there."
%
% Without this pointer, 2.4.3 and hr:rem:notmistake read as
% contradicting (brief §1.5, Collision 2; owed-edits item 2).
% 2.4.3 is not to be merged without it.
% ------------------------------------------------------------

% ============================================================
% DELIVERABLE LISTS (brief §0) — revision 2
% ============================================================
%
% 1. CLAIMS REGISTER
%    - The Golden Question stated (quote display). [Φ]
%    - Not-Descartes contrast (one sentence); full placement at
%      2.2. [Φ]
%    - Admissibility discharged: not paradoxical, not
%      inconsequential (bears on every forced selection). [Φ] —
%      consequence half of ac:rem:scarce's conjunction; ALSO
%      supplies step one of the forcing argument.
%    - Setup: every asking conditions on persistence to date;
%      establishes the question's subject and NO MORE — bare
%      asking fails ac:rem:constraint's test (nihilist on the
%      same branch). [Φ] (rev 2: demoted from argument to setup)
%    - FORCING (ac:prop:forcing, proposition): in the ensemble of
%      biospheres holding and ACTING ON a frozen answer at
%      biosphere scale, survivorship reweights toward the
%      affirmative; negative-holding hosts thin (ac:prop:filter),
%      affirmative unpenalized. [Φ] — three steps: consequence
%      (admissibility), near-analytic middle, reweighting.
%    - Middle step priced as NEAR-ANALYTIC with falsifier: a
%      biosphere organizing practice around the negative while
%      doing no less to persist. (Referee's weak joint, carried
%      at honest strength rather than as necessity.)
%    - Bridge (ac:rem:bridge): individual holding is not a lever
%      on biosphere hazard; position is EVIDENTIAL not causal
%      (Accordance-record relation; gc:rem:teleology); what
%      existence narrows is which answer it is evidence for at
%      biosphere scale, not which answer an individual can hold.
%      [Φ]
%    - Fence (ac:rem:forcing): fixes what survivorship selects /
%      what an asker is positioned to find, not what is the
%      case; positional not truth-conferring; falsifiers
%      in-place (middle-step exhibit; symmetric-penalty
%      coupling; rival question still open). [Φ]
%    - PAYOFF: selection leans on ac:prop:filter (every cell),
%      not ac:prin:tendency (ρ>0) — the Golden Answer's
%      selection does not inherit the Tendency's antecedent;
%      should/can falls out of the mechanism. [Φ]
%    - Should/can: gr1:thm:twocond; trapped P(viable)=0;
%      "the answer is yes and the thing is unavailable."
%      Extinctionist and settled-at-x_ast viability zero alike
%      (hr:rem:concession); anatomy-as-process
%      (hr:rem:forcesprocess). Inherited.
%    - Meaning claim, Φ6-conditional ONLY; compatibility with
%      hr:rem:notmistake stated as difference in kind. [Φ]
%    - Scientific-method parallel NAMED, not run (dedup per §E);
%      argument and burden at 2.2. [Φ]
%
% 2. CROSS-REFERENCES USED
%    \S\ref{ac:sub:placement} ×2, Remark~\ref{ac:rem:scarce},
%    Remark~\ref{ac:rem:admissible},
%    Remark~\ref{ac:rem:constraint} ×2,
%    Remark~\ref{ac:rem:phi6falsifier}, Remark~\ref{ac:rem:noaim},
%    \S\ref{ac:sub:countermeasure},
%    Proposition~\ref{ac:prop:filter} ×4,
%    Principle~\ref{ac:prin:tendency},
%    Remark~\ref{gc:rem:teleology},
%    Theorem~\ref{gr1:thm:twocond}, Theorem~\ref{gr1:thm:trapped},
%    Remark~\ref{gr:rem:weakpremise},
%    Remark~\ref{hr:rem:concession},
%    Remark~\ref{hr:rem:forcesprocess},
%    Remark~\ref{hr:rem:notmistake}
%    NOT USED, with reason: gr:sub:ceiling — enters via twocond
%      and weakpremise; a third restatement of "not correct"
%      declined (unchanged from rev 1).
%
% 3. NEW LABELS CREATED
%    ac:sub:question
%    ac:ssub:statement
%    ac:prop:forcing         (rev 2: converged label with the
%                             referee's draft; was body prose)
%    ac:rem:bridge           (new in rev 2)
%    ac:rem:forcing          (retitled fence; falsifiers updated)
%    ac:ssub:shouldcan
%    ac:ssub:meaning
%    ac:ssub:phrasing
% ============================================================

%############################################################
%%##  FILE: sec2_6_extrapolation.tex
%############################################################

% ============================================================
% §2.5 (compiled order) — Golden Extrapolation
% REVISION 2, against the complete-section report (B1, B2, A4).
% Filename retains WRITING-order numbering; compiled block order
% is authoritative.
%
% MERGE NOTES:
%   B1 ACCEPTED IN FULL. The generation paragraph was doubly
%     broken: "each contributed the same thing" is denied by
%     hr:rem:monotone (contributions increase — the scrapbin's
%     equality intuition carried through unrepaired), and
%     gc:rem:twoprivilege was misused to forbid a stakes ranking
%     when it forbids anchor privilege — a misuse my own closing
%     distinction refuted three sentences later. Rewritten on
%     monotonicity alone: a monotone sequence has no interior
%     maximum, so the superlative is always ahead. The
%     biosphere/lineage distinction dropped per the report;
%     human-history comparative retained as the report's
%     defensible statement.
%   B2 ACCEPTED IN FULL. The gradient is now NAMED: the
%     goldening interpolation ℙ^(t) → Q of gr:def:gg, indexed by
%     age per gr:def:practical, with the number attached
%     (ε-Goldenness over windows 4εt, gr:prop:ratio +
%     gr:rem:invariant). This also discharges the A4 fence
%     concern by construction: ℙ^(t) is a family of measures,
%     not a path.
%   NOTATION: \mathcal{F}_s used for the window field per the
%     report's ℱ_s; verify against the master's macro.
%   Standing from rev 1: "Dimensional Null Aethus" not used
%     pending verification; positive form over three denials;
%     restatement register.
% ============================================================

% [Placement pass, Aug 21: the Question's published neighbors.]
\begin{remark}[The Question's neighbors: an ally and two opponents]
\label{ac:rem:questionneighbors}
The Golden Question has one close published ally and two standing
opponents, and naming all three fixes what the affirmative does and
does not claim.

\emph{The ally.} Jonas's first imperative --- ``that there be a
mankind'' --- is the affirmative answer held as the ground of an
ethics for the technological age \cite{Jonas84}, and it is the
nearest published neighbor of $\Phi6$ read through the Question. The
delta is the direction of support: Jonas \emph{derives} the
imperative, from an ontology of purposes this framework neither needs
nor asserts, whereas here the affirmative is priced as a premise with
its falsifier stated and its adoption left to the reader
(\S\ref{ac:sub:countermeasure}) --- and the forcing proposition
supplies what Jonas's derivation cannot, a survivorship mechanism by
which biospheres that persist come to hold the affirmative without
anyone's having proved it.

\emph{The opponents.} Williams's tedium argument and H\"agglund's
finitude argument hold that endless life would be, respectively,
unlivable and valueless --- meaning requiring mortal stakes
\cite{Wil73,Hag19}. Both run at the individual scale, and the
category shift to the biosphere already blunts them; but the sharper
reply is supplied by the hazard base itself, and it is worth stating
because it dissolves rather than resists the objection. The Golden
Point is not the removal of danger: under the Cox construction the
hazard is strictly positive at every instant forever, and unbounded
lifespan means only that the accumulated hazard converges ---
immortality as a convergent series of mortal moments, never as
invulnerability. So the condition both opponents claim meaning
requires --- live stakes, real fragility, something genuinely at risk
--- is satisfied at \emph{every} time under the affirmative, and what
the objections target is a picture of immortality (stakes ended,
danger over) that the formalism does not contain. The finitude
tradition's quarrel is with a Golden Point no one here has defined.
\end{remark}

\subsection{Golden Extrapolation}\label{ac:sub:extrapolation}

The conditioning of \S\ref{sh:sub:interp} may be anchored at any Aethus in a lineage; Proposition~\ref{gc:prop:recursive} establishes that the anchoring is recursive, and Remark~\ref{gc:rem:twoprivilege} draws the two consequences --- no spacetime event is privileged, no point in a lineage is privileged, and the appearance of privilege is an artifact of which cone one occupies. Golden Extrapolation is what those results look like read from the inside.

\begin{definition}[Golden Extrapolation; \textbf{[$\Phi$]}]\label{ac:def:extrapolation}
To place oneself by Golden Extrapolation is to take one's position as an output of the conditioning rather than as its anchor: a point on the goldening interpolation $\PP^{(t)} \to Q$ of Definition~\ref{gr:def:gg}, indexed by one's own age, occupied for the same reason any point occupies its own.
\end{definition}

The reading is available because the mathematics already refuses the alternative. An anchor is a stipulation, and Proposition~\ref{gc:prop:recursive} shows the conditioning does not require one --- so a lineage point that took itself for the anchor would be claiming a privilege the construction declines to grant anywhere. What replaces the privilege is not diminishment: every point on the interpolation is an output in the same sense, and a position's being unprivileged says nothing about its stakes.

Nor is the gradient a figure. It is the family of measures the paper already owns, and placement on it carries a number. Definition~\ref{gr:def:practical} indexes the interpolation by age --- a biosphere of age $t$ is a Practical Golden Biosphere to tolerance $\varepsilon$ on windows of length $s$ where $\PP^{(t)}$ and $Q$ agree to within $\varepsilon$ on $\mathcal{F}_s$ --- and by Proposition~\ref{gr:prop:ratio} with Remark~\ref{gr:rem:invariant}, a position at age $t$ confers $\varepsilon$-Goldenness over windows of length $4\varepsilon t$: a constant fraction of one's own history, at every age. Self-placement by Extrapolation is therefore measured, not merely unprivileged. And $\PP^{(t)}$ is a family of measures, not a path, so the reading sits on the right side of Remark~\ref{ac:rem:notrajectory} by construction rather than by phrasing.

One comparative remains, and the result that governs it is Remark~\ref{hr:rem:monotone}: stakes rise monotonically along the chain, because a step late in it carries the improbability of every step behind it. A monotone sequence has no interior maximum. So the present may well be the most consequential moment in human history, and cannot be the most consequential in the biosphere's --- under monotonicity that superlative is always ahead. The claim ranks no positions; what rises is what is at issue at them, and that is \S\ref{hr:subsub:stakes}'s business.

% ============================================================
% DELIVERABLE LISTS (brief §0) — revision 2
% ============================================================
%
% 1. CLAIMS REGISTER
%    - Golden Extrapolation (ac:def:extrapolation): position as
%      OUTPUT of the conditioning; a point on the goldening
%      interpolation ℙ^(t) → Q, indexed by age. [Φ] — the
%      gradient now NAMED (per B2); no independent argument.
%    - Self-anchoring claims a privilege granted nowhere. [Φ]
%    - Placement carries a number: ε-Goldenness over windows
%      4εt, a constant fraction of one's history at every age
%      (gr:def:practical, gr:prop:ratio, gr:rem:invariant). [Φ]
%      over cited results.
%    - ℙ^(t) is a family of measures, not a path — fence met by
%      construction (ac:rem:notrajectory). [Φ]
%    - Generation comparative (per B1, rewritten): stakes rise
%      monotonically (hr:rem:monotone); a monotone sequence has
%      no interior maximum; possibly the most consequential
%      moment in HUMAN history, cannot be in the biosphere's —
%      the superlative is always ahead. No equality claim; no
%      use of gc:rem:twoprivilege for ranking; no
%      biosphere/lineage distinction. [Φ]
%    - Stakes themselves are hr:subsub:stakes's. [Φ]
%
% 2. CROSS-REFERENCES USED
%    \S\ref{sh:sub:interp}, Proposition~\ref{gc:prop:recursive}
%    ×2, Remark~\ref{gc:rem:twoprivilege} (positioning only, NOT
%    ranking), Definition~\ref{gr:def:gg},
%    Definition~\ref{gr:def:practical},
%    Proposition~\ref{gr:prop:ratio},
%    Remark~\ref{gr:rem:invariant},
%    Remark~\ref{ac:rem:notrajectory},
%    Remark~\ref{hr:rem:monotone}, \S\ref{hr:subsub:stakes}
%    NOT USED, with reason: "Dimensional Null Aethus" — term
%      unverified; restore at merge if current.
%
% 3. NEW LABELS CREATED
%    ac:sub:extrapolation
%    ac:def:extrapolation
% ============================================================

%############################################################
%%##  FILE: sec2_7_imperative.tex
%############################################################

% ============================================================
% §2.6 (compiled order) — The Golden Imperative
% Drafted against brief 2.7, at the established register.
% Filename retains WRITING-order numbering; compiled block order
% is authoritative.
%
% MERGE NOTES:
%   - The Departric Dichotomy Theorem's IVT/MVT proof (scrapbin
%     226–251) is NOT transferred in any form (brief §1.4(b),
%     confirmed by the register's §6 disposition). The dichotomy
%     is INHERITED from gc:subsub:cascade, where it is
%     unconditional.
%   - DE-DUP with hr:sub:environment per brief pitfall and owed
%     item 1.5: this subsection states the MODAL form only; the
%     expectancy form stays at §8, with the in-body pointer here
%     and the §8-side pointer drafted in the owed-edits file
%     (1.5 now DRAFTED).
%   - Register §6 dispositions adopted: Broad Landscape's strong
%     material lands here (life as disagreeing superposition
%     between Departure and Golden; abiotic-vs-Golden as the
%     fundamental dichotomy). The generation line went to
%     Extrapolation. The fish/algae/RNA-world figure stack is
%     cut — the intermediate-stage claim is stated plainly.
%     "Immortal struggle," "immunity device," and the unpriced
%     Rare-Earth "wins almost every time" are cut per brief.
%   - Grades kept apart per brief: complementarity [P]
%     (inherited); the imperative reading [Φ], Φ6-conditional.
% ============================================================

\subsection{The Golden Imperative}\label{ac:sub:imperative}

The mathematics first, because it is already done. By \S\ref{gc:subsub:cascade}, $G = \neg F \neg$, so $\nu F$ and $\mu G$ are complements outright: the dichotomy between reaching the Golden Point and reaching the Departure Point is inherited rather than argued, holds for arbitrary branching, and is exhaustive on definite states. \textbf{[P]}, and nothing in this subsection adds to it. What the subsection adds is a modal reading and a reframing, both carried at \textbf{[$\Phi$]}.

\begin{principle}[The Golden Imperative; \textbf{[$\Phi$]}]\label{ac:prin:imperative}
Reaching the Golden Point is the only route by which the Departure Point is avoided.
\end{principle}

The modal form is exactly as strong as the complementarity beneath it and no stronger: with exactly one of the two points given to occur, avoiding one \emph{is} reaching the other, and the principle is a near-immediate consequence rather than a separate result --- which is the accurate presentation. Calling it an \emph{imperative} is the overlay, and the overlay imports $\Phi6$ and the Golden Answer: an agent who has adopted neither reads the principle as a true disjunction with no instruction in it, and the paper gives them no argument that they are wrong to. The expectancy form of the same claim --- that every alternative to the Golden Point carries a finite expectancy --- is \S\ref{hr:sub:environment}'s, stated once there; the two forms point at each other and neither is repeated.

% [Merge: Opus placement+content adopted; phenomenological-death naming
% and hr-link folded in from the Fable instatement.]
\begin{remark}[Departure is the losing condition, not death]
\label{ac:rem:losingcondition}
One consequence of the modal definition should be stated before the
outlook it licenses, because it is where this framework's primitives
part company with those of the fields it borders. Astrobiology and
futurism take \emph{extinction} as the quantity to estimate. This work
does not, and the reason is that departure and death come apart in both
directions.

\emph{A biosphere can be departed and alive.} Past $\mu G$ every
continuation fails to advance, so the Golden Point is unreachable ---
modally, not merely improbably --- and viability is zero. Nothing has yet
happened to it. Its hazard has fired no more than before; its members
are unaware; the physical extinction lies ahead and may lie far ahead.
What has been lost is the \emph{possibility} of the thing worth
reaching, and that loss is complete at the moment of departure and not
at the moment of death. In the framework's own terms, a biosphere
departs when its fate is sealed and not when the sealing is executed,
and the interval between the two can be long. The framework's name for the
first event is the \emph{phenomenological death} of the biosphere: the
death this work prices arrives exclusively with the Departure Point,
even where the biological end is yet to come, and the gap between the
two is exactly the bounded remainder the trapped tail measures. This is
the reading \S\ref{hr:sub:extinction} consumes when it finds the
undisturbed world already departed --- entry into $\mu G$ by
configuration is the death, and the hazards that later fire are its
consequences.

\emph{And a biosphere can be undeparted and certain to die.} The
critical regime is the case: departure occurs with probability one, and
life expectancy is nonetheless infinite (\S\ref{gr1:sub:verdict}),
because advancing continuations exist and merely carry measure zero.
Such a biosphere has not departed --- $\mu G$ is a modal condition and
this one fails it --- and the difference is not a technicality, since
everything the survivorship transform computes is computed there.

So neither implication holds. Extinction probability is therefore not
the quantity this framework estimates, and a reader who translates
$\PP(\mathcal G)$ into ``probability of not going extinct'' has
substituted a measure claim for a modal one and will get the trapped and
critical cases exactly backwards.

\emph{What replaces it.} The existential dichotomy for a biosphere is
Departure Point against Golden Point, not death against life. Physical
death is downstream of the first and uninformative about it: where
anatomy is adverse it is certain and its expectancy is finite
(Theorem~\ref{gr1:thm:trapped}(ii)); where anatomy is favourable but
recurrence holds it is certain and its expectancy is infinite; and in
neither case does the timing of death report the state that decided the
outcome. The losing condition is the closure of possibility, and it is
the only thing in this framework that a biosphere can lose \emph{once}.
\end{remark}


\begin{remark}[The outlook, in one sentence]\label{ac:rem:race}
The dichotomy admits a plain statement of what it makes of life, and the
statement should be given once at full width before the reframing below.
Life, from its earliest origins on a planet, undergoes a race to reach
its Golden Point before its Departure Point claims it. That is not a
metaphor laid over the mathematics but the mathematics read in prose:
the two points are complementary fixed points, every actual biosphere
carries weight on both, and the weight moves. The reading corrects the
character the Departure Point is usually given. It is not a dark finish
line recently come into view, to be crossed or avoided by human ingenuity
alone --- as if it appeared with us. It has been the other pole of every
biosphere's superposition since that biosphere began, and what is recent
is only the capacity to move the weight deliberately.
\end{remark}

That correction has a consequence for how the record should be read, and
it is a continuity argument rather than a rhetorical one. The naive
picture places the Departure Point ahead of us as one possible
destination among several --- a crossroads at the modern day, at which
humanity may or may not turn the wrong way. The framework's picture is
that at \emph{every} moment of a biosphere's history there stood the
same possibility of diverging onto it, and that our ancestors, in
surviving improbably, were not preparing for a contest we would later
fight but fighting the same one. By
Principle~\ref{gc:prin:continuity} the ancestral struggle and the
agency-mediated one are one activity under different mediation; what
Remark~\ref{ac:rem:losingcondition} adds is that they have one
\emph{losing condition} as well, and it was available throughout. The
Miracles of \S\ref{gc:sub:naturalhistory} are accordingly not
lucky preliminaries to a modern crisis. They are earlier rounds of it,
and the record we hold is the record of the rounds that were won.

Two things follow for reading the record. The stakes are higher now, by
Remark~\ref{hr:rem:monotone}, and only by that --- the difference is
magnitude, not kind. And no moment in the record was ever safe: a
biosphere's not having departed at any given point is a fact about that
point and confers nothing on the next, which is why
\S\ref{ac:sub:extrapolation} places one on a gradient rather than past a
threshold.

% [Placement pass, Aug 21: the losing condition's nearest published
% neighbors, deltas named per the methodology companion's discipline.]
\begin{remark}[The losing condition's neighbors, and the deltas]
\label{ac:rem:losingneighbors}
Three bodies of writing stand near Remark~\ref{ac:rem:losingcondition},
and each should be located exactly, because each contains part of the
correction and none contains the object.

\emph{Parfit's asymmetry} is the source intuition of the field: that
the difference between near-total and total destruction exceeds the
difference between none and near-total, because the second difference
destroys the future rather than more of the present \cite{Parfit84}.
That is the modal claim in embryo --- what total destruction uniquely
removes is not more people but the \emph{possibility} --- and this
framework reads its own fixed point as the formalization Parfit's
intuition was waiting for: $\mu G$ is what ``destroying the future''
denotes, the asymmetry becomes a theorem about the poles rather than a
judgment about magnitudes, and the separation results say where the
intuition's usual event-probability rendering fails it.

\emph{Schell's second death} is the nearest literary precursor: the
death of humankind as an object distinct from, and worse than, the sum
of first deaths --- ``the death of death,'' the cancellation of the
unborn \cite{Schell82}. The kinship with phenomenological death is
close enough that the delta must be said plainly: Schell names the
object and mourns it; the fixed point defines it, separates it from
biological death in both directions, and prices what follows. A named
object supports rhetoric; a defined one supports theorems.

\emph{The s-risk literature} has, under the banner ``fates worse than
extinction,'' been circling the departed-and-alive quadrant --- futures
in which a biosphere persists at scale inside a configuration from
which nothing advances \cite{AG16}. The framework supplies that
literature the category it has lacked: such futures are not
\emph{worse than} the losing condition, they are the losing condition,
occupied at length --- departure with the trapped tail still running.
And the exchange runs the other way at a stated price: a
suffering-focused reader is precisely the mixed-currency reader of
\S\ref{gc:sub:acting}, for whom the residue argument honestly breaks;
the framework locates their exit rather than disputing their values.
\end{remark}

The reframing is what Definition~\ref{gc:def:departure} already carries, read at full width: the Departure Point is not the destruction of life, it is the inevitability of the destruction --- so a biosphere past its Departure Point differs from the abiotic in tense, not in outcome. Two consequences. \emph{The fundamental dichotomy moves.} The division that matters is not extant against extinct, and not abiotic against biotic --- the latter boundary is not an extremum by any fundamental measure --- but abiotic-or-Departric against Golden, since the complementarity above is exactly the statement that everything ends up on one side of that line or the other. \emph{Life as it stands is an intermediate.} A biosphere still inside the disagreeing superposition between the two points (Remark~\ref{gc:rem:pairsuffices}) is neither of the things the dichotomy names; its standing is not a category but a pendency, resolved in one direction or the other by the same complementarity. Neither consequence is a demotion --- pendency is where all the section's stakes live --- but both are $\Phi6$-conditional readings of a \textbf{[P]} result, and the tags do not mix.

% ============================================================
% DELIVERABLE LISTS (brief §0)
% ============================================================
%
% 1. CLAIMS REGISTER
%    - Complementarity: νF and μG complements outright via
%      G = ¬F¬; inherited from gc:subsub:cascade; arbitrary
%      branching; exhaustive on definite states. [P] — cited,
%      not re-proved; the scrapbin's IVT/MVT proof not
%      transferred.
%    - Golden Imperative (ac:prin:imperative): reaching the
%      Golden Point is the only route by which the Departure
%      Point is avoided. [Φ] — modal form; exactly as strong as
%      the complementarity; near-immediate consequence.
%    - The imperative READING imports Φ6 + the Golden Answer; a
%      non-adopter reads a true disjunction with no instruction,
%      and is given no argument otherwise. [Φ]
%    - De-dup discharged: modal form here, expectancy form at
%      hr:sub:environment, mutual pointers (§8 side in owed
%      file, item 1.5).
%    - Departure Point = inevitability of destruction, not
%      destruction (gc:def:departure); a post-Departure
%      biosphere differs from the abiotic in tense, not outcome.
%      [Φ] reading of inherited framing.
%    - Fundamental dichotomy: abiotic-or-Departric vs Golden,
%      not extant/extinct or abiotic/biotic (the latter not an
%      extremum by any fundamental measure). [Φ]
%    - Life inside the superposition is a pendency, not a
%      category; not a demotion. [Φ]
%
% 2. CROSS-REFERENCES USED
%    \S\ref{gc:subsub:cascade}, Definition~\ref{gc:def:departure},
%    Remark~\ref{gc:rem:pairsuffices} (added per complete-section
%    report §C — the w : 1−w result the superposition sentence
%    invokes), \S\ref{hr:sub:environment}
%    NOT USED, with reason: gc:eq:departurefix (brief's list) —
%      the complementarity citation carries the whole load; the
%      fixed-point equation adds machinery the modal statement
%      does not consume. Add at merge if the master's usage
%      pairs them.
%
% 3. NEW LABELS CREATED
%    ac:sub:imperative
%    ac:prin:imperative
% ============================================================

%############################################################
%%##  FILE: sec2_2_placement.tex
%############################################################

% ============================================================
% §2.2 — Placement
% Drafted against brief §5 (deferred subsection, now written,
% with the argument-first material of 2.1/2.3/2.4/2.5 fixed).
%
% MERGE NOTES:
%   - POSITION (verification report F, recommendation ADOPTED
%     provisionally, reversible at the author's ruling):
%     Placement now sits at the END of the argument, before the
%     ledger — final slot after 2.6/2.7 when they exist. The
%     reasoning that deferred its WRITING applies to where it
%     SITS: every reference in this file (Φ6, the Golden
%     Question, ac:rem:phi6falsifier) is now backward, and the
%     hinge delta's punchline ("Φ6 is such an operation") lands
%     on a reader who has met Φ6. All references are symbolic,
%     so the move costs nothing but renumbering. Filename kept
%     for continuity; the compiled doc carries the new order.
%   - CITATION KEYS ARE PLACEHOLDERS: Wittgenstein69,
%     Mackie77, Camus42, and the trilemma key (Sextus or
%     Albert1968 — merge to pick per the master's bibliography
%     style). None verified against the master's .bib.
%   - The brief names phil_sci.tex as the model for stating a
%     delta against a near neighbor. I do not have that file;
%     the deltas below follow the structure the brief attributes
%     to it (state the similarity honestly, name the delta as
%     one property, price the addition) and should be checked
%     against the actual model at merge.
%   - Spends the material 2.1 and 2.4 were instructed to leave:
%     the Bacon/Descartes historical placement, the
%     skepticism↔scientific-method analogy in full, and the
%     not-Descartes contrast beyond one sentence.
%   - Grade [Φ] throughout except the fecundity expectation,
%     which is [O] and priced as such.
% ============================================================

\subsection{Placement}\label{ac:sub:placement}

The problem and its countermeasure now stand, and four neighbors sit close enough that silence about them would read as either ignorance or identity. This subsection states each delta once, at one property each, and then makes the historical claim the section has so far declined to make --- at its price.

\emph{Hinge propositions.} The closest neighbor by a distance is Wittgenstein's hinges \cite{Wittgenstein69}, and a reader who knows them will have assumed identity since Definition~\ref{ac:def:frozen}. The resemblance is real and should be conceded in full: a hinge is exempt from grounding, functions as the foundation of a practice, and answers the child's \emph{why} with \emph{it just is} --- all of which is the frozen question's phenomenology exactly. The delta is Remark~\ref{ac:rem:threeprops}, property for property. Hinges are world-directed: \emph{here is a hand}, \emph{the world existed before my birth} --- claims about what is the case, held by believers whether or not any action is forced. Frozen questions of the curse's kind are action-directed and bind at forced action. And --- the property that carries everything --- hinges are not selectable: they are inherited with a form of life, held involuntarily, and the hinge framework contains no operation by which an agent chooses among candidate hinges, because on that framework the choosing is not a thing an agent does. $\Phi6$ \emph{is} such an operation. The premise the countermeasure turns on is thus not merely absent from the hinge framework but inexpressible in it, and that is the difference doing work. Neither framework is thereby refuted; they answer different questions, and the frozen question is not a hinge renamed.

\emph{The Agrippan trilemma.} Regress, circle, or dogma \cite{Sextus}: every justification ends in one of the three, and the frozen question is the dogmatic horn --- encountered from inside, by an agent under forced action, rather than catalogued from outside by an epistemologist. That much is placement, not addition. The addition is one sentence: the trilemma is symmetric over stopping points --- it has no resources to prefer one dogma to another, since its whole content is that justification ends --- and $\Phi6$ is precisely a preference over stopping points. The section's claim against the trilemma is therefore exactly as strong as $\Phi6$ and no stronger: a meta-normative selection among the dogmatic horn's occupants, priced at \textbf{[$\Phi$]}, falsified where $\Phi6$ is (Remark~\ref{ac:rem:phi6falsifier}).

\emph{Error theory.} Mackie's argument from queerness \cite{Mackie77} shares the first horn's terminus --- no objective values --- and draws from it a semantic thesis: moral claims are systematically false. The curse is indifferent to the semantics; it operates identically whether moral claims are false, truth-valueless, or expressive, because its engine is the agent's information state rather than the status of moral language. The sharper difference is that the curse survives the negation of Mackie's conclusion. Error theory assumes the absence of objective values; the second horn of \S\ref{ac:sub:problem} is what happens in their \emph{presence}, and the communication argument has no counterpart in Mackie because his framework never needs one. The curse is not error theory with different vocabulary; it is what remains standing on both sides of the question error theory answers.

\emph{Absurdism.} The nearest neighbor in temperament rather than structure \cite{Camus42}. The absurd is a confrontation --- the demand for meaning against a silent universe --- and the canonical response is a stance: revolt, lucidity, the refusal of both hope and suicide. The curse is narrower and colder: a structural fact about an information state, with no confrontation in it and no stance demanded. And the tradition sits inside the problem rather than beside it: the absurdist's assurance of unsureness is itself a frozen answer, held with the same groundlessness as any conviction it replaces. What the curse's framework adds is not a better stance but a different kind of object --- $\Phi6$ is a criterion, and a criterion is answerable to falsifiers in a way a stance is not.

\emph{The historical claim.} What remains is the analogy the section has been holding back since its first sentence, and it is stated here in full so that its price can be attached in the same breath. Skepticism was never dissolved, and could not be; it was defanged, by a procedure that created a demarcation --- the scientific against the ascientific --- and the demarcation proved fecund. The procedure's own defense is circular and knows it: to deny the method's standing is to assert a framework in which knowledge and empirical observation are uncorrelated, which defeats the asserting, so the defense is circular and, by its own measure, the least circular available --- no rival argument even bounds its own circularity. The parallel now claimed: agency's curse stands to philosophy as skepticism stood to inquiry; $\Phi6$ creates a demarcation among groundings --- constrained by the agent's existence, or not --- as the method created one among claims; and \emph{least circular} and \emph{minimally arbitrary} are the same selection expressed in the two vocabularies, which is why the analogy is structural rather than incidental.

The price. Everything in the analogy up to fecundity is a restatement of results already graded --- the demarcation exists if $\Phi6$ selects, and Remark~\ref{ac:rem:phi6falsifier} says what would show it selects nothing. The placement of the present moment beside Bacon and Descartes is a different kind of claim: it asserts that the demarcation will produce work, and that is a forecast, evidentially empty today, carried at \textbf{[O]} and at nothing less. The scientific demarcation earned its history in retrospect, by centuries of what it enabled; this one has enabled \S\ref{ac:sub:question} and \S\ref{gs:sec}, and whether that is the beginning of a comparable retrospect or the whole of a small one is exactly what no one now writing can know. The analogy is offered with its falsifier and its forecast separated, and only the first is asserted.

% ============================================================
% DELIVERABLE LISTS (brief §0)
% ============================================================
%
% 1. CLAIMS REGISTER
%    - Hinge delta: resemblance conceded in full; difference =
%      the three properties of ac:rem:threeprops, with
%      selectability decisive — Φ6 is inexpressible in the hinge
%      framework. Neither framework refuted. [Φ]
%    - Trilemma delta: frozen questions are the dogmatic horn
%      met from inside; the trilemma is symmetric over stopping
%      points, Φ6 is a preference over them; the addition is
%      exactly Φ6-strong. [Φ]
%    - Error-theory delta: the curse is indifferent to moral
%      semantics and survives the negation of Mackie's
%      conclusion; the communication argument (second horn) has
%      no Mackian counterpart. [Φ]
%    - Absurdism delta: structural fact vs confrontation;
%      criterion vs stance; the absurdist's assurance of
%      unsureness is itself a frozen answer. [Φ]
%    - Historical analogy, structural part: skepticism defanged
%      by a demarcation; method's defense least-circular by its
%      own measure; Φ6's demarcation parallel; least-circular =
%      minimally-arbitrary identification. [Φ] — priced by
%      ac:rem:phi6falsifier.
%    - Historical analogy, fecundity part: placement beside
%      Bacon/Descartes is a forecast, [O], asserted as a
%      forecast only; falsifier and forecast explicitly
%      separated.
%
% 2. CROSS-REFERENCES USED
%    Definition~\ref{ac:def:frozen}, Remark~\ref{ac:rem:threeprops},
%    Remark~\ref{ac:rem:phi6falsifier} ×2, \S\ref{ac:sub:problem},
%    \S\ref{ac:sub:question}
%    \cite{Wittgenstein69}, \cite{Sextus}, \cite{Mackie77},
%    \cite{Camus42} — ALL PLACEHOLDER KEYS, verify/supply at
%    merge; trilemma cite may be Albert1968 instead per house
%    preference.
%    \S\ref{gs:sec} — forward placeholder \ref per the fixed §7
%     label map (register 3.3); resolves when §7 is drafted.
%
% 3. NEW LABELS CREATED
%    ac:sub:placement        (resolves forward refs in 2.1, 2.3
%                             rev 2 note, and 2.4)
% ============================================================

%############################################################
%%##  FILE: sec2_8_ledger.tex
%############################################################

% ============================================================
% §2.8 (compiled order) — Ledger
% REVISION 2, against the complete-section report (B3, B4, C).
%
% MERGE NOTES:
%   B3 ACCEPTED: "Inherited dependency" restored to the master's
%     sense (structural exposure, per gc:sub:ledger's usage);
%     the bookkeeping inventory relocated under a new "Merge
%     conditions" heading, kept per A5.
%   B4 ACCEPTED: selectability now priced — falsifier supplied
%     at ac:rem:threeprops (the ONE touch to stable 2.1, made
%     because B4 places it there explicitly; flagged in 2.1's
%     notes for confirmation against §E's no-reopening
%     preference) and Costs entry added below.
%   C: drafting-history clause cut from Buys (the observation
%     survives in this file's audit-delta note (i), where it
%     belongs); \ref* → \ref; Costs itemize → enumerate per
%     gr:sub:ledger.
%   AUDIT DELTAS from plan (unchanged from rev 1):
%     (i) the forcing repair added a claim not in the brief —
%         the Golden Answer's selection is independent of the
%         Tendency's antecedent. Listed under Buys and priced
%         under Costs via ac:prop:forcing's falsifiers.
%     (ii) the trapped-regime disposition changed twice; the
%         ledger records the final form only.
%     (iii) forcing's falsifiers are the middle-step exhibit and
%         the symmetric-penalty coupling; the brief's "asker off
%         the survival branch" prices only the setup.
%   FORMAT: enumerate for Costs (matching gr:sub:ledger);
%     conform the remainder to the master's ledgers at merge.
% ============================================================

\begin{remark}[The register of this section, stated for the reader who
noticed]\label{ac:rem:register}
This section is philosophical and does not pretend otherwise. It is not
armchair, in the specific sense that it links to the mathematics at every
load-bearing step --- the two-condition theorem, the Filter, the ceiling
--- and the relation runs one way: \emph{the mathematics limits the
possibility space, and the philosophy is what stands out in what remains}.
Where the section crosses from one register to the other it does so by
naming a premise, carried at \textbf{[$\Phi$]}, rather than by letting a
philosophical claim inherit a mathematical grade. So the section merges
the practices of two disciplines deliberately, and the merger has a
discipline of its own: nothing here revokes the grounding of
\S\ref{gr1:sec} or \S\ref{gr:sec}, and everything here that is not a
theorem says so.
\end{remark}

\subsection{Ledger}\label{ac:sub:ledger}

\paragraph{Costs.}
\begin{enumerate}
    \item \emph{Non-emptiness} (Proposition~\ref{ac:prop:nonempty}), \textbf{[$\Phi$]}. Falsifier: an agent under forced action all of whose answers terminate in something accessible that is itself a reason; or, on the second horn, a veracity certificate whose warrant does not draw on the transmitting source.
    \item \emph{Selectability} (Remark~\ref{ac:rem:threeprops}, third property), \textbf{[$\Phi$]}. Falsifier: a demonstration that an agent's frozen questions are fixed by their circumstances rather than chosen within them --- Proposition~\ref{ac:prop:nonempty} would stand, and everything built after it would fall.
    \item \emph{Minimal arbitrariness} ($\Phi6$, Premise~\ref{ac:phi6}), \textbf{[$\Phi$]}, kind: meta-normative selection rule. Falsifiers: a frozen question comparably compelled by existence that delivers a different Golden Answer (kills the use); no ordering on ``constrained by the agent's own existence'' (kills the premise).
    \item \emph{Forcing} (Proposition~\ref{ac:prop:forcing}), \textbf{[$\Phi$]}. Falsifiers: a biosphere organizing its practice around the negative answer while doing no less to persist (the middle step is near-analytic, not analytic, and this is its exhibit); a coupling by which affirmative-organized practice lowers persistence comparably. The setup separately voided by an asking that does not condition on persistence.
    \item \emph{The Golden Filter} (Proposition~\ref{ac:prop:filter}), \textbf{[P]}, under the hazard-ordering hypothesis; exact for a cohort, asymptotic for the renewal case.
    \item \emph{The Golden Tendency} (Principle~\ref{ac:prin:tendency}), \textbf{[S]}, antecedent $\rho > 0$ inside the statement, the conditioned-from region existing by Proposition~\ref{gr1:prop:trichotomy}(i); regime scoping inherited from \S\ref{gr:subsub:regimes} unchanged; ensemble form only.
    \item \emph{The meaning claim} (\S\ref{ac:ssub:meaning}), \textbf{[$\Phi$]}, $\Phi6$-conditional throughout; falls with $\Phi6$'s falsifiers.
    \item \emph{Golden Extrapolation} (Definition~\ref{ac:def:extrapolation}), \textbf{[$\Phi$]}, a reading of cited results (Proposition~\ref{gc:prop:recursive}, Remark~\ref{gc:rem:twoprivilege}, Proposition~\ref{gr:prop:ratio}); the generation comparative governed by Remark~\ref{hr:rem:monotone} alone --- possibly maximal in human history, never in the biosphere's, the superlative always ahead.
    \item \emph{The Golden Imperative} (Principle~\ref{ac:prin:imperative}), \textbf{[$\Phi$]} overlay on the inherited \textbf{[P]} complementarity of \S\ref{gc:subsub:cascade}; instructive force $\Phi6$-conditional.
    \item \emph{Placement's four deltas}, \textbf{[$\Phi$]}, each priced at $\Phi6$'s strength where they consume it; \emph{the fecundity forecast}, \textbf{[O]}, asserted as forecast only, its falsifier and forecast explicitly separated.
\end{enumerate}

\paragraph{Buys.} A grounding for the paper's normative vocabulary that is stated rather than assumed --- one premise, one doorway, and a four-step chain a hostile reader can check for the absence of any ``$Q$ implies one ought'' step. A meaning claim explicitly conditional on an adopted premise and therefore compatible with Remark~\ref{hr:rem:notmistake}. A should/can separation that arrives from the mechanism --- the Golden Answer's selection runs through the Filter, which reweights in every cell, while availability is $\rho$'s alone --- rather than by stipulation; in consequence, the selection does not inherit the Tendency's antecedent. A self-placement that is measured, not merely unprivileged. A provenance record that prices the 2024 material honestly, including where it took a component for the whole. And a placement that locates the section's one premise as inexpressible in its nearest neighbor's framework rather than merely absent from it.

\paragraph{Not claimed.} That the Golden Answer is true. That any agent ought to adopt $\Phi6$. That the biosphere will reach the Golden Point, or that $\rho > 0$ obtains --- no one is entitled to the antecedent (\S\ref{gr:sub:ceiling}). That agency's curse is dissolved; it is countered, and the counter is adopted. That anything in the section generates a duty --- $\Phi2$ stands, and the enumeration of \S\ref{ac:sub:countermeasure} exhibits it standing. That any biosphere improves --- every claim is ensemble-form, and the trajectory reading remains unbuilt at \textbf{[O]}. That an individual's holding moves a biosphere's hazard --- position, not causal purchase. That constraint confers truth, that compulsion admits a number, or that constraint and consequence trade off by any stated relation. That any generation's contribution equals any other's --- contributions rise with the stakes. That the historical placement beside Bacon and Descartes is anything but a forecast.

\paragraph{Inherited dependency.} \S\ref{ac:sec} inherits the exposure that \S\ref{gc:sub:ledger} records, and more acutely. Everything after the countermeasure consumes survivorship over futures, and under a fixed adverse $\rho$ there is nothing of the kind to consume: \S\ref{ac:ssub:shouldcan} has nothing to separate, Principle~\ref{ac:prin:imperative} is satisfied vacuously, and the section's claims are then true and useless. The exposure is reduced only where technoanatomy is a process rather than a parameter, and those dynamics are \textbf{[O]} (\S\ref{gr1:subsub:fluct}); Remark~\ref{hr:rem:forcesprocess} is the standing argument that the process is required regardless. The exposure is real and is not reduced by the arguments' internal soundness.

% ============================================================
% DELIVERABLE LISTS (brief §0) — revision 2
% ============================================================
%
% 1. CLAIMS REGISTER
%    - No new claims; the ledger prices the section's. New Costs
%      line: selectability (per B4). Inherited-dependency
%      exposure restored to the master's sense (per B3):
%      fixed adverse ρ makes the section vacuously true and
%      useless; reduced only by [O] anatomy dynamics.
%    - New Not-claimed line: generational equality of
%      contribution (disclaims the withdrawn B1 claim, matching
%      the A3-praised pattern of fencing withdrawn claims).
%
% 2. CROSS-REFERENCES USED
%    Internal: ac:prop:nonempty ×2, ac:rem:threeprops, ac:phi6,
%    ac:prop:forcing, ac:prop:filter, ac:prin:tendency,
%    ac:ssub:meaning, ac:def:extrapolation, ac:prin:imperative,
%    ac:sub:countermeasure, ac:ssub:shouldcan.
%    Master: gr1:prop:trichotomy, gr:subsub:regimes,
%    gc:prop:recursive, gc:rem:twoprivilege, hr:rem:monotone,
%    gc:subsub:cascade, hr:rem:notmistake, gr:sub:ceiling,
%    sh:thm:Q, sh:lem:monotone, gr1:thm:trapped,
%    gr1:thm:twocond, gc:rem:teleology, gr:prop:ratio,
%    gc:sub:ledger, gr1:subsub:fluct, hr:rem:forcesprocess.
%    Forward: gs:sec, gs:sub:initiative, gs:sub:rule,
%    gs:sub:fence.
%
% 3. NEW LABELS CREATED
%    ac:sub:ledger
% ============================================================

%############################################################
%%##  FILE: sec2_owed_master_edits.tex
%############################################################

% ============================================================
% OWED EDITS TO THE MASTER — consolidated patch file for §2
% (register §1, "the highest-risk category"). These edits live
% OUTSIDE the new subsection files and are the ones that get
% lost. Apply at merge; check off in place.
%
% Register item numbers retained. Status: [DRAFTED] = patch text
% below, apply verbatim-ish after fitting to surrounding text;
% [ACTION] = mechanical, no drafting needed; [PENDING] = blocked
% on an unwritten subsection; [NOTE] = cannot be drafted from
% outside the master, instructions only.
% ============================================================

% ---- 1.1 [DRAFTED] — counter guard, head of §6 --------------
% Insert immediately before §6's first premise. Pairs with the
% \setcounter{premise}{5} guard in 2.3 (sec2_3, in place).
% WITHOUT BOTH, Φ3 renumbers to Φ4 and every textual $\Phi3$ in
% §§6–7 breaks.
%
%   \setcounter{premise}{0} % GUARD — do not remove; pairs with
%   % \setcounter{premise}{5} in §2.3. §2's premise prints Φ6;
%   % §6's print Φ1–Φ5 unchanged. See §2.3 merge notes.
%
% Guard arithmetic watch (register §7): currently §2 adds exactly
% ONE premise (Φ6). Neither 2.4, 2.2, nor the drafted §7 plan
% adds another. If any later draft does, the {5} shifts — recheck
% both guards before merge.

% ---- 1.2 [DRAFTED] — patch at gr:phi2 -----------------------
% One sentence at the Φ2 site in §6 (also drafted at the foot of
% sec2_3; kept in both places deliberately):
%
%   "Φ2 governs what the apparatus outputs; it is silent on the
%    separate question of how an agent grounds duties the
%    apparatus never issued, for which see the selection rule of
%    \S\ref{ac:sub:countermeasure}."

% ---- 1.3 [DRAFTED] — patch at hr:rem:notmistake -------------
% Sentence to APPEND to the remark (also at the foot of sec2_4).
% MUST be fitted to the remark's actual text, which has not been
% available to the drafting side:
%
%   "The disclaimer concerns an unconditional verdict issuing
%    from outside human self-interpretation; the Φ6-conditional
%    claim of \S\ref{ac:ssub:meaning} is of a different kind ---
%    an adopted grounding, revocable and certified by nothing ---
%    and the two are compatible for the reasons given there."

% ---- 1.4 [DRAFTED] — sentence at gr1:thm:twocond ------------
% At the point where the theorem appears:
%
%   "The theorem also scopes the Golden Tendency of
%    \S\ref{ac:sub:tendency}: the forced answer of
%    \S\ref{ac:ssub:shouldcan} purchases nothing toward the
%    capacity this theorem separates from it."
%
% (Drafted to serve double duty with 2.4.2's citation; trim if
% the theorem's site already gestures forward.)

% ---- 1.5 [DRAFTED] — de-duplicate hr:sub:environment --------
% The Imperative subsection now exists and carries the MODAL
% form with its in-body pointer to §8. The §8-side pointer, to
% be inserted at hr:sub:environment where the only-route claim
% currently appears:
%
%   "The modal form of this claim --- that the Golden Point is
%    the only route by which the Departure Point is avoided ---
%    is stated once, at \S\ref{ac:sub:imperative}; the present
%    subsection keeps the expectancy form, that every
%    alternative carries a finite one."
%
% Fit to the passage; if §8's sentence already IS the modal
% form, replace it with the expectancy form and the pointer.

% ---- 1.6 [DRAFTED] — pointer at hr:subsub:stakes ------------
% One sentence where purpose/stakes are discussed:
%
%   "Purpose and its grounding are treated once, in
%    \S\ref{ac:sub:question}; the present subsection confines
%    itself to stakes."
%
% Fit to the actual passage; the point is a pointer, not a
% second treatment accreting in either direction.

% ---- 1.7 [NOTE] — §1 Introduction ---------------------------
% Still empty; must be written knowing §2 exists. The scrapbin's
% "To begin this paper…" opening presumed this material came
% first and is unusable as-is. Not a §2 deliverable; scheduling
% decision owed.

% ---- 1.8 [DONE, verified Aug 22] — stray rule ----------------
% Current master: no uncommented rule line exists (grep ^==== empty).

% ---- 1.9 [DONE, verified Aug 22] — duplicate preamble --------
% Current master: exactly one \title/\author/\date and one
% \allowdisplaybreaks outside comments.  The nine-fold blocks
% belonged to the pre-merge master this note was written against.

% ---- 1.10 [VERIFIED Aug 22, retained] — uncited bib entries --
% Britannica_coal, Magoon_coal, CL67, Ric44, Ric45 remain uncited
% in the master body; retained pending the coal/recurrence
% material's return, and flagged for every extraction's cut list.

% ---- 1.11 [NOTE] — Newcomb companion ------------------------
% decision.tex exists in the corpus, keyed AethusDecision.
% Needed: (a) bib entry; (b) replace the live "VERIFY AT MERGE"
% footnote at gc:rem:teleology; (c) upgrade that paragraph from
% "structural parallel alone" to a citation. Requires the
% companion's actual title/venue for the bib entry — not
% draftable from here.

% ============================================================
% LABEL RULINGS ADOPTED (register §3) — for the record
%   3.1  Φ6 label: ac:phi6 (house pattern gr:phi1…gr:phi5).
%   3.2  ac:sub:placement: MOOT — 2.2 is now drafted
%        (sec2_2_placement.tex) and creates the label; no stub
%        needed, nothing dangles.
%   3.3  §7 map FIXED and already consistent with all §2 drafts:
%        gs:sec, gs:sub:initiative (7.1), gs:sub:rule (7.2),
%        gs:sub:limit (7.3), gs:sub:regimes (7.4),
%        gs:sub:fence (7.5), gs:sub:ledger (7.6).
%   3.4  gs:sub:fence mispointer: fixed in 2.3 rev 2 (D1); the
%        cost paragraph and enumeration step (iv) point at
%        gs:sub:rule.
% ============================================================

% ---- E.4 [ACTION] — bib entries for Placement's cites --------
% Four philosophy citations used in Placement (house-style keys
% per verification report §E): Wittgenstein69 (On Certainty),
% Sextus (Outlines of Pyrrhonism), Mackie77 (Ethics: Inventing
% Right and Wrong), Camus42 (Le Mythe de Sisyphe). None yet in
% the master's bibliography; entries owed at merge.

% ---- F [RECORD] — Placement moved to end of argument ---------
% Verification report F recommendation adopted provisionally:
% Placement sits before the ledger (after 2.6/2.7 when written);
% everything between shifts up one. All ac: references are
% symbolic; nothing dangles. REVERSIBLE at the author's ruling —
% if reversed, restore 2.2's slot and re-check that "Φ6 is such
% an operation" is survivable for a reader who has not met Φ6.

\section{Continuing from Nexic reasoning}\label{gc:sec}

\begin{quote}\small
The companion development derives a typicality constraint --- the
Accordance Principle --- and applies it to the natural-historical record
under a weighting it calls Optimation: the cosmic weighting conditioned
on observer-existence \cite{Nexus, Ben26Found, Ben26Optimal}. That
weighting is there a construction, adopted because observer-existence is
the obvious thing to condition on. This section argues that it is the
wrong thing, that the right thing is the \emph{Golden Point}, and that
Optimation is recovered exactly as a consequence rather than assumed as
a premise.

Two claims carry the argument. The \emph{Past-Alignment Principle}
(\S\ref{gc:sub:pastalign}) states that conditioning the cosmic weighting
on the Golden Point and then discarding every attribute outside one's
own past light cone returns precisely the Nexus of Optimation --- so
that the record we observe is, attribute for attribute, the past-light-cone
restriction of the Golden Nexus. The \emph{Golden-Conditioning
Hypothesis} (\S\ref{gc:sub:conditioning}) argues by contradiction that
this is no coincidence: observer-existence is too weak a condition to
have produced the record, since it survives perturbations that the
record does not. What conditioning on observers cannot explain,
conditioning on the Golden Point does.

The consequence is a reversal of order. Optimation ceases to be the
framework's starting point and becomes its first derived result; the
Golden Biosphere ceases to be an object defined by survival asymptotics
and becomes the conditioning target from which the rest follows. What
the argument does \emph{not} settle is the future: whether the alignment
that holds behind us extends ahead of us is left to the three
hypotheses of \S\ref{gc:sub:extrapolation}, which are stated and not
adjudicated.
\end{quote}

% ---------------------------------------------------------------------
\subsection{Vocabulary and standing conventions}\label{gc:sub:vocab}

\begin{convention}[Terms carried from the companion]\label{gc:conv:vocab}
The argument runs in the companion's vocabulary
\cite{Nexus, Ben26Found}, and the six terms it uses are fixed here.
An \emph{Aethus} (plural \emph{Aethae}) is a partial information state:
some attributes stated, all others genuinely blank, individuated by its
stated content and closed under entailment from it. The \emph{Cosmic
Nexus} is the occurrence-weighted measure over Aethae, with no
conditioning absorbed into it. \emph{Optimation} is that measure
conditioned on the biosphere-level observer attribute, and the
\emph{Nexus of Optimation} the resulting weighting. The
\emph{Aetherian Nexus} is the Aethus of our own natural history ---
the one whose stated attributes we read off the record. \emph{Centric
unfolding} is the accumulation of stated attributes along a
particular reasoner's history. And the \emph{Golden Point} is the event
of a biosphere's passage into the Golden regime, in the sense
\S\ref{gr:sec} makes precise; the \emph{Golden Nexus} is the cosmic
weighting conditioned on that event. The \emph{Golden Aethus} is the
Aethus stating exactly that event as a future-referring attribute, all
else blank --- the index of that conditioning
(\S\ref{to:sub:goldenaethus}).
Two usage notes are carried from the companion's fuller statement.
\emph{Optimation} names the conditioned measure; the fluke content a
record so conditioned displays against the Cosmic Nexus's own rates is
what the measure \emph{shows}, and this work uses the word for the
measure and speaks of its fluke content when the display is meant. And
the Aetherian Nexus is not a weighting but an Aethus --- the one of
which our own Aethus is a proper child --- and is used in that sense
throughout.
To \emph{complement} a stated attribute is to replace it by a
strictly mutually exclusive alternative: the perturbed Aethus states
something incompatible with what ours states at that place, rather than
lacking what ours has or carrying it alongside further content. Aethae
being individuated by stated content, every perturbation in this
framework is such a substitution, and perturbed Aethae are accordingly
disjoint from ours rather than related to it by inclusion. An
\emph{Imperfection Magnitude Drop} is what results from complementing a
stated attribute correlated with the remainder of an Aethus: the
success-intersection of the perturbed history is displaced by many
orders of magnitude, becoming drastically minute without reaching zero
\cite[\S IMD]{Nexus}. The \emph{High-Contrast Earth Hypothesis} asserts
the size of that displacement and not its sign, the sign being supplied
separately by the Accordance Principle. The Drop is the graded
counterpart of departure and the object Definition~\ref{gc:def:epsdep}
refines, and the \emph{Strong Accordance Principle} supplies the
direction as the statement $c>k$, with $k$ the prior-odds toll
\cite{Ben26Optimal}.
\end{convention}

\begin{convention}[Grades]\label{gc:conv:grades}
As elsewhere: \textbf{[P]} proved, \textbf{[C]} cited, \textbf{[S]}
argued at scaling or schematic level, \textbf{[O]} open, and
\textbf{[$\Phi$]} a philosophical premise carrying its own falsifier.
This section's two central claims --- the conditioning principle and the
Golden Tendency's antecedent --- are \textbf{[$\Phi$]} and \textbf{[S]}
respectively, and the extrapolation hypotheses are \textbf{[O]} by
construction. What is \textbf{[P]} here is arithmetic downstream of a
premise: the weight-ratio bound, the disjointness lemma it rests on, and
the credence and Lindy laws of \S\ref{gc:subsub:routeworld}, each proved
conditional on an input the section states and does not prove. Nothing
here is asserted at a grade the argument does not reach.
\end{convention}

% ---------------------------------------------------------------------
\subsection{Imported apparatus}\label{gc:sub:imported}

This section consumes results from the companion program as premises,
and a reader without those texts cannot check the arguments that follow.
The apparatus is therefore stated here, in the companion's own terms and
at its own grades, so that what is borrowed is visible and what is
argued here is separable from it. Nothing in this subsection is claimed
as new; nothing outside it is assumed.

\begin{principle}[Accordance; imported at \textbf{[$\Phi$]}]
\label{gc:prin:accordance}
For a natural-historical attribute $A$ with canonical alternative $A'$
its complement, and $H$ the production of an observer,
\begin{equation}\label{gc:eq:accordance}
\frac{\PP(H\mid A')}{\PP(H\mid A)}\;\lesssim\;\frac{\PP(A)}{\PP(A')},
\end{equation}
read attribute by attribute over a partial state and never as a global
optimisation over complete histories: \emph{improbability on the path to
an observer must be bought by observer-relevance}
\cite{Nexus, NexicFoundation}.
\end{principle}

The relation $\lesssim$ is not a tolerance. It is \emph{defined} by an
exceedance form: for any $\kappa\ge1$,
\begin{equation}\label{gc:eq:accordexceed}
\Pr\left[\frac{\PP(H\mid A')}{\PP(H\mid A)}>
\kappa\,\frac{\PP(A)}{\PP(A')}\right]\;\le\;\frac{1}{1+\kappa},
\end{equation}
the probability taken over the realisation, so violations are priced
rather than forbidden and the whole tail profile is fixed with no free
parameter. \eqref{gc:eq:accordance} is the $\kappa=1$ reading. Two
consequences travel with it and are used in this work: the alternative
is \emph{canonical} --- the complement, so nothing is analyst-supplied
--- and the exceedance form is its own multiplicity guard, bounding the
expected number of $\kappa$-fold violations among $n$ interrogated pairs
by $n/(1+\kappa)$, with the standing caution that $n$ must count
attributes \emph{scanned} and not merely those interrogated.

\begin{convention}[The sum-collapse]\label{gc:conv:collapse}
Over a finite product of $n$ interrogable slots of $b$ values each,
alternatives at perturbation depth $\Delta$ are counted by
$\binom{n}{\Delta}(b-1)^{\Delta}$ and suppressed by
$10^{-c_{\mathrm{eff}}\Delta}$, so the total closes by the binomial
theorem at $\bigl(1+(b-1)10^{-c_{\mathrm{eff}}}\bigr)^{n}$. The realised
configuration carries a share of order unity exactly when
$c_{\mathrm{eff}}>\log_{10}\bigl(n(b-1)\bigr)$ --- contrast per slot
outrunning the logarithm of combinatorial supply \cite{Ben26Optimal}.
\end{convention}

\begin{remark}[What is imported, and what this section owes]
\label{gc:rem:importscope}
Three notes on the standing of the above.

\emph{They are premises here, not results.} Each is argued in the
companion and none is re-argued in this work, so every claim in
\S\ref{gc:sec} that consumes one inherits its grade. A reader who
rejects Principle~\ref{gc:prin:accordance} should read this section's
conclusions as conditionals on it.

\emph{The companion's own limits transfer with them.} In particular the
Accordance Principle is explicitly \emph{not} a licence for global
position inferences: a position in a completed roster has no
attribute-wise translation, and the retort that the position could be
made an attribute fails because its answer space presupposes the roster
it was meant to replace \cite{NexicFoundation}. Where this work makes a
position inference --- Ansatz~\ref{gc:ans:copernican} --- it does so on
its own account and says so.

\emph{The vocabulary this section runs on.} \S\ref{gc:sub:humanhistory}
and \S\ref{gc:sub:naturalhistory} state their arguments in the 2022 Laws
--- Second Law, Third Law, Golden Linkage, the $L$-threshold --- while
this subsection states the companion's constraint in the companion's
current notation. The two are the same apparatus under older and newer
letters; the companion's own notational-provenance section supplies the
translation and records that the Laws are its constraint as first
written, and a reader moving between the papers should use it rather than
translate by hand \cite{Nexus}. Where this work refines a companion
operation --- the topological criterion on Imperfection Magnitude Drops
refines the companion's counterfactual operation, and the coupling
analysis of \S\ref{gc:subsub:cascade} consumes its closure formalisation
--- the companion now carries canonical anchors for both, and the
dependency is checkable there.

\emph{And some companion machinery is used more lightly than the above.}
Where a further companion result is consumed at a single site, it is
stated where it is used rather than here; the Accordance Principle and the
sum-collapse are the two this section leans on repeatedly, and the
vocabulary they run in is fixed at \S\ref{gc:sub:vocab}.
\end{remark}



% ---------------------------------------------------------------------
\subsection{The problem: what does Optimation condition on?}
\label{gc:sub:problem}

Optimation is introduced in the companion as the cosmic weighting
conditioned on observer-existence, and the choice is natural: an
observer is what is doing the reasoning, so an observer is what the
reasoning conditions on. The choice is nevertheless a choice, and it
is load-bearing --- every rarity verdict the companion issues is
relative to it.

The difficulty is that observer-existence is \emph{robust} in a way the
record is not. Strip from our history the coal that powered
industrialization, or the metallurgical sequence that made rocketry
possible, or any of a dozen similar contingencies, and what remains is
still a world with observers in it. Observerhood survives these
perturbations; the specific character of our record does not. A
conditioning target that is indifferent to the difference cannot
explain why we find the difference realized.

That gap is what the rest of this section fills. The proposal is that
the operative condition was never observer-existence but the Golden
Point, and that Optimation is what Golden-conditioning looks like when
restricted to the region of spacetime one can actually see.


\noindent Two questions follow, and the rest of this section answers them
in order. \emph{Which} target should replace observer-existence, and on
what grounds --- \S\ref{gc:sub:continuity}. And what reason is there to
think the replacement is not merely better motivated but actually the
condition under which the record was produced ---
\S\S\ref{gc:sub:pastalign}--\ref{gc:sub:humanhistory}, in two
independent legs.

% ---------------------------------------------------------------------
\subsection{Choosing the target: continuity and non-privilege}\label{gc:sub:continuity}

Before the argument for Golden-conditioning, a prior question: why the
Golden Point rather than some more familiar target? ``Technologically
advanced civilization'' is the obvious candidate and is the wrong one,
for two reasons that compound.

It is not well-defined. Any threshold one names --- writing, metallurgy,
industry, spaceflight --- is a stipulation, and the arbitrariness is not
cosmetic: every rarity verdict issued relative to such a target inherits
it.

And it is \emph{discontinuous with natural history}. A conditioning
target must be applicable across the whole record if the record is to be
read as one object. ``Civilization'' applies to the last few millennia
and to nothing before, so a framework built on it cannot run a single
conditioning across both the end-Cretaceous impact and the Industrial
Revolution. It must instead splice two analyses at a joint it has no
principled way to locate.

% [Referee pass Aug 22, minor 6, DEFERRED TO AUTHOR: the referee
%  suggests promoting the definition's internal/external division and
%  the provenance footnote into their own remark so the environment
%  states the definition alone.  The restructure is voice-laden
%  (which paragraphs are the definition is an authorial call), so it
%  is recorded here rather than executed.]
\begin{definition}[Departure Point]\label{gc:def:departure}
The \emph{Departure Point} is the event past which \emph{every}
continuation fails to advance --- the first moment at which all
scenarios lead to a failure to advance past that point. The modal
quantifier is the definition and not an embellishment of it: what is
being named is the closure of possibility, not the observation of
non-progress.

\emph{This is what makes it the dual of the Golden Point.} That was
defined as the point past which a biosphere is too capable at self
preservation to regress past it --- which, stated with the quantifiers
exposed, is the point past which \emph{some} available course has
\emph{every} continuation avoiding foreclosure. The asymmetry between
the two is deliberate: departure is $\forall\exists$ --- every course
admits a foreclosing continuation --- and Goldenness is
$\exists\forall$ --- some course admits none.

Two $\forall$-statements could not be complements of one another;
$\exists\forall$ and $\forall\exists$ can be, and
\S\ref{gc:subsub:cascade} shows that these two are, by exhibiting them
as the complementary fixed points of De Morgan-dual operators. The
quantifier asymmetry is therefore not a wrinkle in the duality but its
mechanism.

The pair is exhaustive on definite states, and an Aethus carries weight
across both. \S\ref{gc:subsub:grading} makes that precise, since it is
best stated once the grading is in hand.

The reading this forces bears stating, since a looser gloss
would lose it. Goldenness is not the guarantee that nothing can go
wrong; it is the possession of a \emph{course} along which nothing
does, whatever the environment answers. It is a claim about available
strategy, not about the absence of hazard --- which is why a Golden
biosphere still faces a universe that would destroy it and is Golden
anyway.

Both are \emph{points} for the same reason: irreversibility. A state
one may still leave is not a point but a phase, and neither concept is
about phases.

Two consequences of stating it this way.

\emph{Biological death is a mechanism, not the definition.} Death is one
way a trajectory is foreclosed, and it happens to be the most visible
one; irreversible loss of the capacity to advance is another, and the
definition is indifferent between them. What is extracted is the
\emph{phenomenology of foreclosure}, taken apart from the particular
mechanism of biosphere-death in which it is most legible --- and once
extracted it applies wherever there is a trajectory to foreclose,
which is why the notion carries across the scales
Principle~\ref{gc:prin:continuity} requires.

\emph{Non-advancement alone does not suffice.} A biosphere that fails
to progress for an age has not thereby reached its Departure Point; it
has had a lull, and a lull is reversible. Only when no continuation
remains that advances has the point been passed. This is the
distinction that the phrase ``failure to advance'' elides and that the
quantifier restores.

Two consequences for how the model of \S\ref{gr1:sec} should be read.
The hazard $h^{\mathrm{tot}}_t$ of \eqref{tc:eq:hazard} is the hazard
\emph{of foreclosure}, not of death; the muffling exponent $x$ is
accumulated against that; and every Imperfection Magnitude Drop is a
movement toward it. And the analytic form of the modal claim is \emph{pathwise}
divergence --- $\Lambda_\infty=\infty$ along every continuation, which
the trapped regime satisfies outright
(Theorem~\ref{gr1:thm:trapped}(i)) --- of which
$\PP(\Lambda_\infty<\infty)=0$ is the measure shadow; the two come
apart exactly where Remark~\ref{ac:rem:losingcondition} separates
them, since the critical regime carries the shadow while advancing
continuations survive at measure zero. This is why the
trapped regime of Proposition~\ref{gr1:prop:trichotomy}(iii) has
departed: not because a biosphere halting at $x_\ast$ has stopped, but
because with $\rho<0$ the viability probability is identically zero, so
no continuation available to it advances. Halting is not the departure;
having nowhere left to go is.

The term is used rather than ``extinction'' deliberately, and the
reason is the same one that governs the choice of the Golden Point
itself. ``Extinction'' is a biological category, applying to taxa and
carrying the vocabulary of one discipline; the Departure Point must
apply at every scale the record contains --- to a Mesozoic lineage as
readily as to a biosphere tomorrow --- and must name an event in the
Aethic block universe rather than a fact about organisms. It is
therefore defined at the same primitive level as the Golden Point,
which is what allows Principle~\ref{gc:prin:continuity} to be stated at
all.

\emph{Internal and external.} The Departure Point divides by the
locus of its cause. An \emph{internal} Departure Point is brought about
by factors within the biosphere's own homeworld; an \emph{external} one
is inflicted from outside, and is roughly constant in magnitude over
long timescales because the cosmos is roughly constant in makeup. That
division is the ontological origin of the two channels of
\eqref{tc:eq:hazard}: the crest channel carries the external hazard and
is bounded, the trough channel carries the internal and is not.
Technoanatomy (Definition~\ref{gr1:def:anatomy}) is then the statement
of how the second responds to capacity acquired against the
first.\footnote{\emph{Provenance.} The Departure Point, its
internal/external division, the constancy of the external threat, and
the observation that a biosphere may come to draw energy out of itself
are all set out in our \emph{Golden Science} of August 2022
\cite{GoldenScience22}, some years before the present formalism and
before the drafting collaboration. That document also states the Golden Point as the point
past which a biosphere cannot regress, derives a positive immortality
probability from an exponentially decaying hazard, argues for
``biosphere'' over ``civilization'' as the unit, and names the Trace
Connection as an open problem --- which \S\ref{gc:sub:pastalign} is the
attempt to close. The formalism here should be read as supplying
machinery for commitments already made there, not as generating
them.}
\end{definition}


\subsubsection{The implication cascade: departure as a least fixed point}
\label{gc:subsub:cascade}

Definition~\ref{gc:def:departure} carries a recursion that is worth
making explicit, because it is already available in the companion
framework in finished form and because it changes what the model's
hazard is a hazard \emph{of}.

The internal and external hazards of \eqref{tc:eq:hazard} \emph{imply}
departure; they are not exhaustive of it. A biosphere may reach a
configuration from which no route avoids those hazards, and by the
modal definition that configuration \emph{is} a Departure Point,
attained before either hazard has fired. The statement ``no matter what
you do, you reach a Departure Point'' is itself a Departure Point ---
and since that statement can again be true one step earlier, the
implication propagates.

\begin{remark}[This is the third postulate's operator, and it is proved]
\label{gc:rem:thirdpostulate}
The companion's third postulate supplies exactly this recursion
\cite{Aethus}. Over states ordered by descent, with $\mathrm{Base}$
collecting the non-recursive modes, it defines a safety operator and
its dual
\begin{align*}
F(Y)&=\bigl\{A\notin\mathrm{Base}:\ \exists B\sqsubseteq A,\
\forall C\sqsubseteq B,\ C\in Y\bigr\},\\
G(X)&=\mathrm{Base}\cup\bigl\{A:\ \forall B\sqsubseteq A,\
\exists C\sqsubseteq B,\ C\in X\bigr\},
\end{align*}
both monotone, so that by Knaster--Tarski the greatest fixed point
$\nu F$ and least fixed point $\mu G$ exist for arbitrary branching and
are complements of one another; for finitely-branching trees the
closure ordinal is at most $\omega$, so the cascade terminates.

Read $\mathrm{Base}$ as departure by internal or external hazard
--- departure-sufficient occurrences in the sense of
Remark~\ref{gr1:rem:hazardof}, with nothing about death presupposed.
Then $G$ is the doom operator in the sense above --- every route has a
doomed continuation --- and the Departure Point is
\begin{equation}\label{gc:eq:departurefix}
\mathcal{D}\;=\;\mu G\;\supsetneq\;\mathrm{Base}.
\end{equation}
Three consequences, each already carried by the companion's
machinery rather than assumed here.

\emph{The implication is one-directional by design.} The companion
states the third postulate as an implication and not a biconditional,
which is precisely the containment in \eqref{gc:eq:departurefix}: the
hazards suffice for departure and are not necessary to it.

\emph{The two-generation form is the companion's, and was not imported.}
The quantifiers above are the third postulate's own: $F$ asks for a child
all of whose grandchildren are safe, $G$ for a state every child of which
has a doomed grandchild, and it is the \emph{pair} of generations that
makes the operators monotone and the fixed points complementary. Read
that way the Golden Point is an event horizon in a precise sense --- the
state past which the capacity to trigger departure is closed off, which
is $\nu F$'s $\exists\forall$ said in the other direction --- and the
Departure Point is its dual, the state at which every remaining
possibility retains that capacity. This parallel is recorded rather than
used, and its provenance is worth a line: the companion's derivation was
conducted under a discipline of \emph{not} reaching for parallels to the
Nexic and Golden material, on the ground that the surprise of agreeing
superpositions gave no warrant to expect every construction to transfer.
The parallel therefore holds without having been arranged, which is a
weaker thing than a proof and a stronger thing than a resemblance.

\emph{The Golden/Departure duality is a theorem, not a construction.}
The Golden Point is $\nu F$ --- there exists a route all of whose
continuations remain safe --- and the Departure Point is $\mu G$. Since
$G=\neg F\neg$, the two fixed points are complements outright, so the
duality asserted in Definition~\ref{gc:def:departure} need not be
stipulated: it is inherited, and the pair is exhaustive.

\emph{The complementarity is a fact about definite states.} The
fixed-point analysis lives on the Boolean lattice of definite states,
where the result above is exact. An Aethus is a weighted superposition
of such states, and what that implies for the categories is taken up in
\S\ref{gc:subsub:grading}.

\emph{Human extinction sits in the difference
$\mu G\setminus\mathrm{Base}$.} It is not itself a biosphere-level
departure by internal or external hazard,
since the biosphere persists; but with no muffling agent remaining,
every continuation reaches $\mathrm{Base}$
(Proposition~\ref{hr:prop:extinction}). It therefore enters at the
first step of the cascade. The classical reading --- that human
extinction is a Departure Point because it is only a matter of time ---
is the $n=1$ case of the recursion, and \eqref{gc:eq:departurefix} is
what licenses it.
\end{remark}

\begin{remark}[The model's hazard is a lower bound]
\label{gc:rem:hazardlower}
The consequence for \S\ref{gr1:sec} should be stated plainly, since it
qualifies every rate in this work. The hazard $h^{\mathrm{tot}}_t$ of
\eqref{tc:eq:hazard} is the hazard of $\mathrm{Base}$: it prices the
internal and external mechanisms and nothing else. The hazard of
\emph{departure} is the hazard of $\mu G$, which by
\eqref{gc:eq:departurefix} is at least as large and in general strictly
larger --- because a biosphere departs when its fate is sealed, not
when the sealing is executed, and the interval between those can be
long.

So the model's timescales are conservative in a specific direction:
where it reports a departure time $T$, the Departure Point proper may
have been passed considerably earlier, and the survival curves of
\S\ref{gr1:sub:verdict} are upper bounds on the probability of not
having departed. Nothing in the asymptotic analysis is disturbed by
this --- a cascade that fires earlier by a bounded lead time leaves
exponents untouched --- but the quantitative readings inherit it, and
the anatomy urgency of \S\ref{gr1:subsub:fluct} is sharpened by it: the
window in which one may still act closes at the cascade's rate, not at
the hazard's. Computing that rate is open \textbf{[O]};
\S\ref{sch:sub:schedule} is a first attempt at it, exhibiting the excess
as a third hazard channel whose completion weight is a least fixed point
\eqref{sch:eq:fixedpoint}.

A second approach, complementary rather than competing, is available in
the companion and is worth recording here because it constrains what the
rate can be. A sealing-hazard is not a hazard on a trajectory but a
\emph{settlement} rate on the books: what it prices is the acquisition of
statedness by the attribute that the fate is sealed, and that is an
event of the present index whatever its referent's date --- statedness
being a property of the indexing Aethus and never of a clock reading, so
that blanks lack tense \cite{Aethus, AethusDecision}. Three consequences follow, and the first is
restrictive. The weight the current books place on a fixed
future-referring attribute is, by the framework's law of total
probability, a \emph{bounded martingale} along the progression --- the
current state prices its own next resolutions --- so its expected rate of
increase vanishes identically and \emph{the sealing-hazard cannot be a
drift}. What survives as a per-unit-time object is the settlement flux,
the weight mass absorbed into statedness per unit proper time, together
with the direction-blind quadratic-variation rate, which a refined
account could consult alongside it. And since a bounded martingale
converges, the current weight audits its own hazard: it decomposes into
the price of eventual settlement plus the expected never-settled residue,
so the number available at an instant is already an integral over the
flux.

That last observation has a consequence for how the open problem should
be posed, and it is the one worth carrying. The flux and the level are
not rival candidates for \emph{the} sealing-hazard; they are the two ends
of a one-parameter family, obtained by pricing settlement at a discount,
\begin{equation}\label{gc:eq:sealprice}
h_\gamma(B)\;=\;\E_B\bigl[e^{-\gamma(\rho-t)}W_\rho\bigr],
\qquad \gamma\ge0,
\end{equation}
with $\rho$ the earlier of settlement and the path's terminus. At
$\gamma\to0$ optional stopping returns the level itself; at
$\gamma\to\infty$ the exponential concentrates and $\gamma h_\gamma$
returns the instantaneous rate. And the kernel's shape is \emph{forced}
rather than chosen: requiring the pricing to commute with refinement of
the progression --- the price of the children's prices returning the
parent's --- makes the delay factor multiplicative across concatenated
intervals, whose bounded monotone solutions are exactly the exponentials
\cite{Aethus}.

Two things follow for this work. The exponential of the Cox construction
(Definition~\ref{sh:def:cox}) is therefore derivable from partition consistency
rather than adopted with it. And this paper already uses both ends of the
family without having named them as ends: \eqref{tc:eq:hazard}'s
$h^{\mathrm{tot}}$ is a rate, the $\gamma\to\infty$ face; the schedule
channel's completion weight \eqref{sch:eq:fixedpoint} is a level, the
$\gamma\to0$ face. Reading them as one family says what the open rate of
this remark actually is --- not a single number awaiting computation, but
a curve whose urgency parameter is itself a modelling commitment, and
whose mixtures over urgencies are, by Bernstein's theorem, exactly the
completely monotone kernels.

The construction is developed at individual scale in \cite{Aethus}, where
the settling attribute is lineage termination; what transfers here is the
form, by Principle~\ref{gc:prin:continuity}, and what does not transfer
is any of its quantities.
\end{remark}

\begin{proposition}[One permanent floor suffices; \textbf{P} given the
model]\label{gc:prop:onefloor}
Suppose some channel $h_j$ of \eqref{tc:eq:hazard} acquires a
\emph{permanent floor}: there are $c>0$ and $t_0$ with
$h_j(t)\ge c$ for all $t\ge t_0$, no route remaining by which it falls
arbitrarily again. Then $\Lambda_\infty=\infty$, so
$\PP(\Lambda_\infty<\infty)=0$ and the state lies in $\mu G$.
\end{proposition}

\begin{proof}
The channels are non-negative, so $h^{\mathrm{tot}}\ge h_j$ pointwise,
and
\[
\Lambda_\infty=\int_0^{\infty}h^{\mathrm{tot}}_t\,dt
\;\ge\;\int_{t_0}^{\infty}h_j(t)\,dt
\;\ge\;\int_{t_0}^{\infty}c\,dt=\infty. \qedhere
\]
\end{proof}

\noindent Two things the proof does \emph{not} use are worth marking,
since both are tempting and one is false. It uses no property of the
other channels: in particular the internal channel need not
participate, and a permanent floor in the external channel alone
places the state in $\mu G$ --- which is the claim's whole point, that
one does not need the base hazard entire. And it does not pass through
``$\Lambda_\infty<\infty$ requires $h^{\mathrm{tot}}_t\to0$'', which
would be the natural route and is not available: a non-negative
function may have finite integral without tending to zero, as tall
narrow spikes of shrinking width show. The floor hypothesis makes the
detour unnecessary, since it bounds the integrand below on a set of
infinite measure directly.

\begin{example}[Human extinction as an Imperfection Magnitude Drop]
\label{gc:ex:humanIMD}
Human extinction is an Imperfection Magnitude Drop, and for reasons
that are structural rather than sentimental. The grades differ across
the parts and are marked where they fall: (A) and (B) below are claims
about historical individuation \textbf{[$\Phi$]}, their combination
with Proposition~\ref{gc:prop:onefloor} is \textbf{[P]} given the
model, and the smallness of $\varepsilon$ is imported at the
companion's grade (Remark~\ref{gc:rem:whatcarries}). Two reasons, kept
separate because they are referred to separately below:

\begin{description}\itemsep2pt
\item[\normalfont(A) \emph{The lineage is spent.}] The cascade to human
evolution ran through the human--chimpanzee last common ancestor, a
particular historical individual which no longer exists to be run
through again.
\item[\normalfont(B) \emph{The setting is spent.}] The specific sequence
of climatic and geographical events that brought that lineage down to
the grasslands has been cashed out, the world now standing at a
different stage.
\end{description}

\noindent The pathway is therefore not merely improbable on repetition
--- it is \emph{unrepeatable}, its preconditions spent.

Two claims must be kept apart here, since the argument moves between
them and they carry different values. That \emph{the human route
recurs} is not improbable but unavailable: (A) and (B) name stated
attributes that are gone, so this contributes nothing at all to what
follows. That \emph{the Golden Point is reached without} (A) and (B) is
a different proposition, and it is minute rather than zero --- an
Imperfection Magnitude Drop, in the exact sense the next paragraphs
derive. Everything positive that remains after human extinction belongs
to the second; the first is simply closed.

\emph{The derivation, run on (A) and (B).} An Imperfection Magnitude
Drop is a targeted perturbation away from one's own Aethus, yielding a
natural history whose success-intersection is drastically minute but
never quite zero \cite[\S IMD]{Nexus}. Human extinction complements (A)
and (B): the last common ancestor is no longer a stated attribute
available to be run through, and the climatic-geographical setting has
been cashed out. Any Aethus continuing from that state therefore
carries those two attributes complemented.

Two established results then apply, in an order that matters. The
High-Contrast Earth Hypothesis fixes the \emph{magnitude}: complementing
an attribute correlated with the remainder of the Aethus displaces
$P(H\mid\cdot)$ by many orders of magnitude, to an effectively random
order --- the hypothesis asserting the size of the gap and not its sign
\cite{Nexus}. The Accordance Principle then fixes the \emph{direction},
placing ours on the greater side. The perturbed Aethae are accordingly
drastically minute, which is what it is to be an Imperfection Magnitude
Drop.

\emph{Why the deep-time reply does not reach it.} ``Re-evolving humans
is unlikely'' invites the answer that deep time makes unlikely things
happen, and a post-extinction biosphere has deep time in abundance. That
answer would land against a claim about rates. It does not land here,
because the displacement is a property of the \emph{perturbation} rather
than of a waiting period: complementing a correlated attribute
re-randomizes the success-intersection wherever in time the continuation
is read, and elapsed time supplies no mechanism returning it to the
unperturbed value. The claim is thus made within probability and not
outside it, as the Disaster-Difficulty Unification Principle requires
--- the two sit on the severity spectrum wherever their probability puts
them, and what distinguishes them is not their kind but that their
probability does not accumulate.

\subsubsection{Alternative lineages, by disjointness}
\label{gc:subsub:disjointness}

The objection that remains is the one about successors: humans will not
return, but some other lineage --- another primate, or the corvid or
cetacean lineages in which complex cognition has arisen independently
--- might find its own way to the Golden Point. It is answered by the
derivation already given, and by a single observation about what such a
lineage is.

\begin{proposition}[Every alternative lineage is an Imperfection
Magnitude Drop relative to (A) and (B)]\label{gc:prop:disjoint}
Let $L$ be any extant lineage other than our own, and let $A_L$ be the
Aethus of a natural history in which $L$ reaches the Golden Point.
The attributes of $L$ mediating its route are \emph{disjoint} from
those mediating (A) and (B): $L$ did not descend through the
human--chimpanzee last common ancestor, and was not the lineage driven
into the African grassland setting. That disjointness is not an added
premise --- it is what makes $L$ an alternative lineage rather than
ours. Hence $A_L$ carries (A) and (B) complemented, and by the
displacement above $A_L$ is an Imperfection Magnitude Drop relative to
our own Aethus.
\end{proposition}

Three features make the proposition economical. The disjointness is
definitional, so it need not be established lineage by lineage. The
displacement is already in hand, so no second mechanism is introduced.
And the conclusion attaches to \emph{each} alternative individually, so
nothing turns on how many there are. Each is a drop, for the same
reason, by the same result.

Two clarifications, since the proposition is easy to over-read.

\emph{It does not say a successor is impossible.} An Imperfection
Magnitude Drop is minute and never zero, so $A_L$ retains positive
weight for every $L$. What is claimed is that the weight is displaced
by the contrast, not that it is annihilated --- and the residual is
precisely the $\varepsilon$ of \S\ref{gc:subsub:grading}.

\emph{It does not depend on which lineages are named.} The primates,
corvids and cetaceans are referents rather than premises: they say what
the disjoint attributes belong to. The proposition would hold for a
lineage nobody has thought to nominate, since disjointness from (A) and
(B) is a consequence of being a different lineage and not a fact
discovered about any particular one.

\subsubsection{The conclusion}\label{gc:subsub:imdconclusion}

The pieces now assemble, and the last step is worth running in full
rather than by reference, since it is where a claim about \emph{lineages}
becomes a claim about the \emph{Departure Point} and the two are not
obviously the same kind of thing.

(A) and (B) remove the human route;
Proposition~\ref{gc:prop:disjoint} makes every alternative route an
Imperfection Magnitude Drop. Consider then any branch carrying no new
muffling agent. On such a branch no route remains by which one
re-arises --- that is what the two previous results jointly establish
--- and a channel whose reduction depends on an agent that no longer
exists and cannot re-arise is a channel that will never again fall.
Fix any $c$ below its prevailing value and any $t_0$ after the loss:
$h^{\mathrm{ext}}(t)\ge c$ for all $t\ge t_0$, which is precisely a
\emph{permanent floor} in the sense of
Proposition~\ref{gc:prop:onefloor}.

That proposition then closes the argument in two lines. Channels are
non-negative, so $h^{\mathrm{tot}}\ge h^{\mathrm{ext}}$ pointwise, and
\[
\Lambda_\infty\;\ge\;\int_{t_0}^{\infty}h^{\mathrm{ext}}(t)\,dt
\;\ge\;\int_{t_0}^{\infty}c\,dt\;=\;\infty ,
\]
so $\PP(\Lambda_\infty<\infty)=0$. By
Definition~\ref{gc:def:departure} that is the analytic form of the
modal claim --- no continuation advances --- and the state therefore
lies in $\mu G$. One such channel suffices, and no property of the
other channel, of $\rho$, or of the driver is used anywhere in the
step: a single floor anywhere in the total hazard is enough, whatever
the rest of the model is doing.

It is worth marking what this does \emph{not} require. It does not
require that the floor be high, only that it be positive and permanent;
a biosphere left at a very low but irreducible external hazard departs
just as surely as one left at a high one, merely later. And it does not
require extinction to be the mechanism --- what carries the step is the
loss of the \emph{agent}, and any event removing the capacity to muffle
without removing the biosphere would carry it identically.
On every such branch, therefore, human extinction is a Departure Point
\emph{outright} --- in the modal sense of
Definition~\ref{gc:def:departure}, and not merely a step toward one.

Across the full Aethic weight the statement is graded rather than
outright. Positive weight survives on the branches
Proposition~\ref{gc:prop:disjoint} prices --- those on which some
alternative lineage reaches the Golden Point --- and $\varepsilon$ is
that weight. It is positive because an Imperfection Magnitude Drop is
minute and never zero, and small because each such branch takes the
displacement of the High-Contrast Earth Hypothesis
(Remark~\ref{gc:rem:whatcarries}). Its formal statement is
\S\ref{gc:subsub:grading}.

\begin{remark}[What carries the argument, and what would break it]
\label{gc:rem:whatcarries}
The conclusion closes an objection that would otherwise stand against
\S\ref{hr:sec} --- that a successor lineage might re-evolve the muffling
capacity, so the accumulation is inherited rather than forfeited --- and
it closes it with the same machinery that establishes the drop, not
with an additional argument. What is contributed here is only the
identification of (A) and (B) as the attributes perturbed, together
with the observation that any alternative lineage is disjoint from them
(Proposition~\ref{gc:prop:disjoint}); the displacement is then supplied
by the High-Contrast Earth Hypothesis and its direction by the
Accordance Principle, both imported at the companion's grade
\cite{Nexus, Ben26Optimal}.

Three things would break it, and they attack different components.
That (A) or (B) is misidentified --- that the last common ancestor or
the climatic sequence is not in fact a stated attribute of the Aethus in
the relevant sense --- is an attack on this section's own contribution.
That the contrast is not vast, or that the interrogated attributes fall
outside the correlated class the Hypothesis is scoped to, is an attack
on the imported magnitude. And that the direction runs the other way is
an attack on the imported Accordance step. The first is where a critic
should aim if the target is the present argument rather than the
framework it runs on.
\end{remark}
\end{example}

\subsubsection{Grading, and why two states suffice}
\label{gc:subsub:grading}

The cascade settles what a Departure Point \emph{is}. What remains is
that no realistic state satisfies it exactly, since some weight always
survives on branches that advance --- so the operative statements are
graded, and the grading attaches to weight rather than to kind.

\begin{remark}[The recovery branch: a genuine Aethic superposition,
and the grading it supports]\label{gc:rem:boltzmann}
An Imperfection Magnitude Drop does not drive the probability of later
recovery to zero, and this is definitional rather than incidental: a
drop is a perturbation whose success-intersection is drastically minute
\emph{but never quite zero} \cite{Nexus}. The weight remaining after
human extinction is therefore exactly the drop's residual weight ---
that of the perturbed Aethae in which the Golden Point is attained
regardless. Nothing further is needed to say why it is positive, and
nothing further is needed to say why it is small: the first is what an
Imperfection Magnitude Drop is, and the second is the displacement of
Example~\ref{gc:ex:humanIMD}.

That small positive weight sits in apparent tension with a Departure
Point defined by a strict modal quantifier: if \emph{some} continuation
advances, however faintly, then not every continuation fails, and the
label should not apply.

\emph{The post-drop state is an Aethic superposition, and that is the
whole of the resolution.} An Imperfection Magnitude Drop does not
replace one state by another; it produces a state carrying both a
departed component and a recovery component, each genuinely stated, in
whatever proportion the physics dictates \cite{WeightAlg}. Neither
component is an artifact of description, and neither is deleted by the
other being large. The recovery branch is therefore \emph{present in
the Aethus}, not absent from it and merely improbable --- which is
precisely the condition a strict universal quantifier is sensitive to.

\emph{So the label does not attach, and grading is forced rather than
chosen.} If a component that advances is genuinely stated, then it is
false that every continuation fails, and
Definition~\ref{gc:def:departure} does not apply however small that
component's weight. What the situation calls for is not a verdict about
whether the state is departed but a measurement of how much of it is:
the superposition is graded, so the description of it must be graded
too. Nothing is rounded and nothing is discarded to reach that
conclusion --- the grading is read off the state rather than imposed on
it.

\begin{remark}[The graded object at individual scale, and a result that
transfers]\label{gc:rem:deathgrade}
The definition below has an exact individual-scale instance in the
companion, and the companion proves something about it this work does
not: for a child $B$ of a state with survival weight $s_A$, define its
\emph{death grade} $g(B):=1-s(B)/s_A$, the fraction of survival weight
the child destroys. Then the level-price is \emph{affine} in the grade,
and a graded slide of grade $g$ is equivalent at zero urgency to a
mixture placing weight $q=g$ on total settlement and $1-q$ on no change
\cite{Aethus}. Read across, an $\varepsilon$-Departure Point is a child of
death grade $1-\varepsilon$ at biosphere scale, and the affinity and the
mixture equivalence transfer unchanged --- which is worth recording
because it says the grading is linear in exactly the coordinate the
survivor's own accounting uses (Remark~\ref{sh:rem:accumulator}), and
not in its logarithm.
\end{remark}

\begin{definition}[$\varepsilon$-Departure Point]\label{gc:def:epsdep}
A state is an \emph{$\varepsilon$-Departure Point} when the total
Aethic weight of continuations that advance is at most $\varepsilon$. A
Departure Point in the strict sense of
Definition~\ref{gc:def:departure} is the case $\varepsilon=0$.
\end{definition}

Human extinction is then an $\varepsilon$-Departure Point with
$\varepsilon$ the weight of the recovery branch --- which is to say,
in the phrase that states the situation exactly, \emph{almost precisely
a Departure Point, probabilistically}. That is a stronger claim than
the binary label would have been, not a weaker one: it carries a
number where the label carried only a verdict, and it is true where the
label was not.

This settles the question the definition deferred.

\begin{remark}[Where the variation lies]\label{gc:rem:pairsuffices}
Write $w(A)$ for the total Aethic weight that state $A$ places on
continuations which advance. The categories are $\nu F$ and $\mu G$,
complementary and exhaustive on definite states
(\S\ref{gc:subsub:cascade}). What varies across biospheres is therefore
not \emph{which} category one occupies but \emph{how weight is
distributed} across the two: a state at $0<w<1$ is a superposition of
the pair in the proportion $w:1-w$, exactly as
Remark~\ref{gc:rem:boltzmann} describes the post-drop state.

Three consequences.

\emph{The weights carry what a finer taxonomy would.} The framework
supplies Aethic weights already, and $w$ is a number where a category
would be a verdict; the finer information is in the weight and is
available without further vocabulary.

\emph{The grading attaches to the state, not to the kind.}
Definition~\ref{gc:def:epsdep} grades by weight, which is the correct
place for grading to sit: an $\varepsilon$-Departure Point is a state
whose Departure-tending weight is $1-\varepsilon$.

\emph{Every actual biosphere is superposed, and this is the subject
matter.} None is secured and none has $w=0$ while advancing
continuations remain, so every case of interest carries weight on both
poles. That is also why \S\ref{gr:sec}'s interpolation $\PP^{(t)}$ is
the right object: it tracks the redistribution of weight between the
two, rather than arrival at either.
\end{remark}

\begin{remark}[Where the strictness survives, and the two layers]
\label{gc:rem:twolayers}
Two things should be kept apart, because the grading does not reach
everywhere.

\emph{The model's verdicts remain strict.} Within \S\ref{gr1:sec}, the
configuration $x\equiv0$ gives $\PP(\Lambda_\infty<\infty)=0$
\emph{exactly}, and adverse anatomy gives it exactly for every
$\rho\le0$: these are identities, not small numbers. The $\varepsilon$
of Definition~\ref{gc:def:epsdep} does not live inside the model at
all. It lives at the Aethic level, where the weighting distributes over
\emph{which configuration obtains} --- and the recovery branch is
precisely the branch on which a new muffling agent arises, which the
model of a fixed biosphere does not represent. So the model says what
happens given a configuration, strictly; the Aethic weight says how
much weight each configuration carries.

\emph{This is what the two-layer algebra is for.} The companion's
construction is an unweighted layer carrying the structural facts and a
weighted layer laid over it without disturbing them
\cite{WeightAlg}. The cascade of \S\ref{gc:subsub:cascade} is a Boolean
recursion and belongs to the first: $\mu G$ is a set, computed from
implication alone. The grading of
Definition~\ref{gc:def:epsdep} belongs to the second. Neither
displaces the other. Making the cascade itself graded would require a
fixed point of an $M$-valued doom operator, which is open \textbf{[O]}
and is the natural next question this remark raises.
Remark~\ref{sch:rem:mvalued} exhibits one concrete instance of such a
fixed point --- single-type and carrying one number rather than ranging
over the lattice, but computed rather than assigned.
\end{remark}
\end{remark}


\begin{principle}[Continuity; \textbf{[$\Phi$]}]\label{gc:prin:continuity}
The struggle our ancestors underwent against the Departure Point in
natural history is continuous, philosophically and ontologically, with a
society's agency-mediated struggle against the Departure Point. They are
not analogs but the same activity under different mediation.
\end{principle}

The Golden Point is the target that respects this, and it does so
because of what kind of definition it has: \emph{survival conditioning
of its type}. That is primitive enough to apply without alteration to a
Mesozoic mammal evading theropod predation and to a modern society
avoiding its own weapons. Both lower a hazard; both prevent an
Imperfection Magnitude Drop; neither requires the vocabulary of the
other. The substance is continuous where the societal construct is not.

\begin{objection}[Agency]\label{gc:obj:agency}
A Mesozoic mammal does not know it is fighting the Departure Point.
There is no agency, no goal, no strategy --- only an animal avoiding a
predator. Calling that continuous with deliberate societal
risk-management equivocates on the word ``fighting''.
\end{objection}

\noindent The objection assumes agency is what individuates the
activity. It is not. By \S\ref{to:sec}, what individuates a task-object
is its \emph{contribution to the root}, and agency is nowhere in that
criterion --- which is precisely why task-objects were required before
this principle could be stated.

What replaces agency is the mediating mechanism, and Optimation
conditions on the mechanism rather than around it. Where a society
lowers its hazard through deliberation, a lineage lowers its through
evolutionary selection; the conditioning does not care which, because
what it fixes is that the hazard be lowered. It therefore reaches
\emph{through} the mechanism to whatever contingencies the mechanism
requires in order to operate.

\begin{example}[Conditioning through the medium: the
$\alpha$-gal case; \textbf{[S]}]\label{gc:ex:medium}
Our $\alpha$-gal case exhibits the structure exactly, and its
details matter because they show the conditioning having no alternative
route \cite{Nexus}.

Some thirty million years ago the Catarrhine lineage came under
pathogen pressure of a particular kind. The pathogen presented
$\alpha$-gal, a carbohydrate our own cells also produced, so that no
immune response was available which did not attack the host. What
survived carried a mutation disabling \textit{GGTA1}, the gene
responsible for producing it --- a gene retained by every other
mammalian lineage, and lost in ours alone.\footnote{The inactivation of
\textit{GGTA1} in Old World primates is established, as is
pathogen-mediated selection as its leading explanation. The
\emph{severity} of the episode --- whether it approached extinction of
the lineage --- is an inference and not a finding, and the argument does
not rest on it: pressure sufficient to remove a gene universal among
mammals is all that the reasoning requires, and a milder episode
delivering that pressure would serve the example unchanged.}

Now read it under Optimation. The gene was ancient and useful, which is
why it is universal among mammals; and evolutionary selection removes a
gene if and only if there is pressure to remove it. That is the
mechanism's own condition of operation, and it is not negotiable. So
under ordinary conditions the loss was not merely improbable but close
to unavailable: selection had nothing to act on, because the gene was
not costing anything.

A conditioning requiring the loss therefore cannot act on the gene. It
must act on whatever \emph{supplies the pressure} --- and what supplied
it was the contingent existence of a pathogen presenting the very
carbohydrate in question, which converts the gene from an asset into the
one thing standing between the lineage and extinction. Optimation
reweights the pathogen's occurrence to prominence, and the loss follows
\emph{through} selection rather than around it.

This is the general shape. Optimation does not reach past evolutionary
selection to install outcomes; it reweights the contingencies that make
selection deliver them. Agency, where it exists, is one such mechanism
among others --- the deliberative rather than the Darwinian medium ---
and the conditioning is indifferent between them.
\end{example}

\begin{remark}[What the principle buys, and its cost]
\label{gc:rem:continuitycost}
It buys the object the rest of this work analyzes. Without continuity,
\S\ref{gr1:sec}'s hazard process would model two different things
spliced at an arbitrary joint; with it, one process runs from
abiogenesis to the present under a single conditioning, and the
technological present is a \emph{position} along it rather than a change
of subject.

The cost is that the principle is a philosophical premise and not a
result, and its scope should be stated at full width rather than left to
be inferred, because it is larger than the conditioning question that
motivated it. Continuity licenses treating a Mesozoic mammal's evasion
of predation, a body's suppression of its own cancers, and an
institution's not turning its capacity against itself as values of
\emph{one variable}. That is what permits \S\ref{ad:sec}'s institutional
machinery to bear on \S\ref{gr1:sec}'s hazard dynamics at all: without
it, the adversivity model would describe the politics of societies and
the survival model would describe the fate of lineages, and no argument
would carry between them.

So the premise is not merely buying a conditioning target. \emph{It is
what makes this one work rather than several.} Should it fail, the
sections do not degrade gracefully --- they separate, and each retains
whatever it had independently while the connections lapse.

\emph{Falsifier:} a demonstration that deliberative and evolutionary
hazard-reduction differ in a respect the conditioning must track ---
that some feature of the record is explicable under one and not the
other. The agency objection is the natural candidate and is answered
above; a better one would consist in exhibiting such a feature rather
than in restating the intuition. A second and more targeted falsifier
follows from the width just claimed: exhibit a hazard-lowering mechanism
at one scale whose governing variable provably fails to correspond to
anything at another --- a property of institutional self-restraint with
no analog in organismal or lineage-level hazard control, or the
reverse.
\end{remark}

\begin{remark}[The principle transfers form, and not only vocabulary]
\label{gc:rem:continuitydeductive}
One further thing the principle does is easily missed, because the
statement above is defensive in shape --- it says what continuity
licenses one to treat as a single variable, which sounds like permission
to use a word. It also licenses \emph{inference}, and an instance is now
available.

\S\ref{sch:sub:schedule} represents pre-onset hazard as a sequence of
gates. Continuity says the pre-onset struggle and the post-onset one are
the same activity under different mediation. It therefore predicts that
post-onset internal hazard is representable the same way --- that the
trough channel of \eqref{tc:eq:hazard} should decompose into gates, its
smoothness being a coarse-graining rather than a fact.
Remark~\ref{sch:rem:internalgates} carries out the decomposition and
finds two things the prediction did not have to deliver: the exponential
form falls out of geometric compounding rather than being assumed, and
the anatomy trichotomy turns out to be the schedule's own criticality
test with capacity substituted for gate count. The construction also
agrees with the boundary-scaling derivation of
\S\ref{gr:subsub:boundary}, which was arrived at independently.

That is the principle doing deductive work. It did not merely permit the
two hazards to be discussed in one vocabulary; it said what form one
should take given the other, and the form checked out against a
derivation that did not consume it. A reader who regards continuity as a
philosophical convenience should notice that a convenience does not
generate a correct structural prediction about a channel it was not
formulated to describe.

The falsifier stated above is unchanged and now has a sharper instance
available: exhibit an internal hazard whose dependence on capacity
provably admits no gate decomposition, and the transfer fails at exactly
the place this remark claims it succeeds.
\end{remark}

\subsubsection{Why continuity is wanted: the non-privilege of humans}
\label{gc:subsub:nonprivilege}

The principle has now been stated and defended, but not motivated, and
its motivation is not tidiness. Continuity is wanted because it is what
makes the framework non-anthropocentric, and the non-anthropocentrism is essential rather than incidental.

Begin from what the target is. \emph{The Golden Point is a property of a
biosphere} --- not of a civilization, not of a species. That is why it
applies to an alien biosphere as readily as to ours, and applies without
translation: nothing in ``survival conditioning of its type'' mentions
humans, technology, or society, so nothing must be reinterpreted when
the biosphere in question is not ours.

\begin{proposition}[Recursive anchoring; \textbf{[S]}]
\label{gc:prop:recursive}
Because the target is biosphere-level and the principle is continuous
across mediation, the conditioning may be anchored at \emph{any} Aethus
in the lineage consistent with ours. Take a Mesozoic mammal, its Aethus
a parent of ours rather than an Imperfection Magnitude Drop, and ask
what the continuation of Aethic attributes within \emph{its} future
light cone must look like for its biosphere to reach the Golden Point.
The answer runs through the remainder of the lineage --- and so through
us --- without us having been named anywhere in the question.
\end{proposition}

\noindent Our position is therefore an \emph{output} of the conditioning
rather than a stipulated anchor for it. Humans are not a privileged
split point; the technological present is where the recursion happens to
pass, and its role is emergent from a primitive that mentions no
species. The same question asked of a mammal in the Maastrichtian, of a
lineage that has not yet speciated, or of a biosphere elsewhere entirely,
is one question.

\begin{remark}[Three things a human-centered target would cost]
\label{gc:rem:privilegecost}
The alternative --- conditioning on ``technological civilization'' or on
observers --- is not merely less elegant. It forfeits three things, in
increasing order of seriousness.

\emph{Robustness across biospheres.} An alien case would require
separate treatment, since the target would be defined by features of our
own history. A framework whose central conditioning cannot be applied to
the systems it is meant to reason about has not achieved generality; it
has assumed its instance.

\emph{Robustness across perturbations.} \S\ref{gc:subsub:pastalign-argument}
evaluates histories in which humans do not arise. Under a human-centered
target those histories cannot be evaluated at all --- the target simply
fails to apply --- and precisely the perturbations that carry the
argument become unaskable.

\emph{Non-circularity.} A human-centered target takes \emph{us} as
primitive, so the framework would build its base unit out of the thing
it is supposed to explain. The Continuity Principle is what prevents
this, and that is the deepest reason to want it.
\end{remark}

\begin{remark}[A distinction to keep: definition versus conditioning]
\label{gc:rem:defvscond}
Two claims are easily run together and must not be, since the argument's
validity depends on their separation.

The \emph{definition} of the Golden Point --- a survival property of a
biosphere, in the sense \S\ref{gr1:sec} makes precise --- is
stipulative. Someone who rejects everything else in this section may
adopt it, and the hazard analysis of \S\ref{gr1:sec} stands under it
regardless.

The \emph{conditioning claim} --- that this property is what the
Accordance Principle has been conditioning on --- is not stipulative and
is what \S\S\ref{gc:sub:conditioning}--\ref{gc:sub:humanhistory} argue
for, by reductio on natural history and again on human history. Running
the two together would make the argument circular, defining the
conditioning target into its role and then observing that it occupies
it. Keeping them apart is what leaves the argument something it could
have failed at.
\end{remark}

\begin{remark}[Two non-privilege claims, and their symmetry]
\label{gc:rem:twoprivilege}
The section now makes two claims of this shape, and they are the same
claim along different axes.
Principle~\ref{gc:prin:pastalign} indexes its light-cone restriction to
one's own present \emph{because that is where the observable cone is
anchored, not because the event is distinguished}; the construction at
any other event returns that event's cone.
Proposition~\ref{gc:prop:recursive} says the same of position in a
lineage: the conditioning may be anchored anywhere consistent, and
returns that anchor's continuation. No privileged spacetime event, no
privileged lineage point --- and in both cases the appearance of
privilege is an artifact of which cone the reasoner happens to occupy.
\end{remark}

\begin{remark}[A third axis: no privileged tense]\label{gc:rem:notense}
The two axes above concern \emph{where} a reasoner stands. A third
concerns \emph{when}, and it follows from the same two results rather
than needing anything new.

We read the ancestral record as a settled chain in which, at juncture
after juncture, the beneficial thing happened --- which is what
conditioning predicts of any cone read from its own end. By
Proposition~\ref{gc:prop:recursive} the construction is available at any
anchor, so an observer anchored far enough downstream reads \emph{our}
actions the same way: a settled chain, cropped to their past cone, and by
Principle~\ref{gc:prin:pastalign} aligned with their own Nexus of
Optimation. We are, to that anchor, what the Cretaceous is to us. The
distinction between the static record one reads and the dynamic outcomes
one produces is therefore not a distinction in the objects. It is an
artifact of which cone the reasoner occupies --- the sentence the
previous remark ends on, now applied to tense.

Two guards, both standing elsewhere in this work and both binding here.

\emph{Ensemble, not trajectory.} Nothing in this reweights anyone's
actions. Conditioning reweights the ensemble a reasoner is drawn from,
and the downstream anchor's reading says which ensemble rather than what
any biosphere will do; Remark~\ref{ac:rem:notrajectory} is disqualifying about
the other reading.

\emph{And determinacy is indexed to information, not to date.} The
downstream cone holds our actions as determinate entries because that
anchor has read them; ours does not, and nothing about the reading
reaches back. That is exactly the footing on which the companion treats
Newcomb's problem --- determinacy indexed to an information state rather
than to a position in time, so that an unread outcome is a blank entry
and blanks lack tense \cite{AethusDecision} --- and
Remark~\ref{gc:rem:teleology}'s refusal to let a conditioning event
\emph{act} is unamended by it.

The relation to Principle~\ref{gc:prin:continuity} is close and is not
identity, and the two are easily merged.
Continuity is a claim about the \emph{object}: the ancestral struggle and
the agency-mediated one are one activity under different mediation. This
is a claim about the \emph{reading}: the epistemic relation to that
object is the same from any anchor. Neither entails the other --- a
single activity could be read asymmetrically across anchors, and
readings could coincide for activities that were not one --- so they are
neighbors, and the argument here consumes only the second.
\end{remark}

% ---------------------------------------------------------------------
% [Instated Aug 21 2026 (Fable 5), author-commissioned: the extended
%  engagement with the Precipice periodization.  House discipline
%  applied: the position stated at its strongest before any delta; what
%  survives stated at the end; the whole fenced from the Route World.
%  CITATION NOTE: keys Ord20, CSB10, MacAskill22, Ben26Method are NEW
%  and need .bib verification (Ben26Method's title especially --- it is
%  the methodology companion and the author should supply the exact
%  title at merge).]
\subsubsection{The rival periodization: the Precipice, engaged}
\label{gc:subsub:precipice}

The continuity principle has a standing rival, and it is the most
serious one the field possesses: the periodization of Ord's
\emph{Precipice} \cite{Ord20}. This work owes it a direct engagement
--- not because the rivalry is total, but for the opposite reason. Of
everything in the existential-risk literature, the Precipice's central
definition is the nearest approach to the losing condition as this
framework draws it, and the methodology companion's discipline applies
with full force to nearest neighbors \cite{Ben26Method}: where a
position is close, the deltas must be named exactly, because a dense
neighborhood is precisely where a distinct position gets misfiled as a
variant of what each reader already holds. The deltas here are three
--- the quantity, the shape of time, and the endpoint --- and each is
stated below after the position itself is stated at its strongest.

\paragraph{The position, at its strongest.} The Precipice defines an
existential catastrophe not as extinction but as \emph{the destruction
of humanity's long-term potential} --- explicitly including
unrecoverable collapse and unrecoverable dystopia alongside extinction
--- and argues that humanity entered, in 1945, a period different in
kind from all prior history: capability has outrun wisdom, anthropogenic
risk now exceeds the natural background by orders of magnitude (the
book's considered figures: something like one in six per century
against roughly one in ten thousand), the period is short by
civilizational standards, and the task it sets is \emph{existential
security} --- risk brought low and locked low --- after which a Long
Reflection can settle what the future is for. It is the most careful
and most quantified statement of its picture in the literature, its
treatment of its own uncertainties is honest, and its definition is
modal in spirit: ``destruction of potential'' is closure-of-possibility
language, and the inclusion of non-extinction catastrophes is exactly
the concession that death and loss come apart. The definition
descends from Parfit's asymmetry \cite{Parfit84}
(Remark~\ref{ac:rem:losingneighbors}), and the descent matters below:
what the Precipice inherits from Parfit is the intuition, and what it
does not inherit --- because nothing in the lineage supplies it --- is
a formal object for the intuition to be \emph{about}. All of that should be
said before any of what follows, because what follows is not a
dismissal. It is the claim that the Precipice's definition is right,
that its own operationalization does not honor it, and that the
periodization built on that operationalization inverts in the cases
that decide.

\paragraph{Delta one: the quantity. The ledger reverts to events.}
The Precipice's definition points at membership in $\mu G$ --- a
configuration property, the closure of possibility
(Definition~\ref{gc:def:departure},
Remark~\ref{ac:rem:losingcondition}) --- but its arithmetic is
risk-per-century of catastrophe \emph{events}. The two come apart in
both directions, and the separations are not corner cases; they are
this framework's central theorems. A biosphere can be departed with no
catastrophe on any ledger: past $\mu G$ nothing has yet happened ---
hazards fire later, members are unaware
(Remark~\ref{hr:rem:silentforfeit}) --- and the sharpest instance is
the undisturbed world of \S\ref{hr:sub:extinction}, which minimizes
per-century catastrophe probability on the event bookkeeping and is
already departed on the modal one: the safest world the ledger can
describe is a lost world. And a biosphere can be undeparted while the
event bookkeeping condemns it: in the critical regime departure is
certain, expectancy is nonetheless infinite, advancing continuations
exist, and everything the survivorship theory computes is computed
there --- a world an event metric writes off with probability one is
the world in which the entire structure of hope resides. So a ranking
of worlds by existential-catastrophe probability misorders them
relative to the quantity the Precipice's own definition names, and
this framework's position is not that Ord's definition is wrong but
that it deserves a formalization his ledger cannot give it: the
framework sides with the definition against the arithmetic.

\paragraph{Delta two: the shape of time. Gradient, not age.} The
Precipice's picture is a threshold crossed --- one short age, different
in kind, bracketed by Trinity on one side and by security or ruin on
the other. Against this the framework sets three results it did not
build for the purpose. First, continuity
(Principle~\ref{gc:prin:continuity}, and the cascade of
\S\ref{gc:subsub:cascade}): the Departure Point has been available at
every moment of the biosphere's history, the record's Miracles are
earlier rounds of the same engagement rather than preliminaries to a
modern one, and what 1945 changed is mediation and magnitude ---
never the arrival of a losing condition that was absent before.
Second, and this is the part that preserves everything the Precipice
is right to insist on: the stakes claim is not disputed but
\emph{re-derived}, and strengthened. Remark~\ref{hr:rem:monotone}
obtains by combinatorial construction what the Precipice obtains by
periodization --- each moment carries the highest stakes the lineage
has yet faced, because a failure late in the chain forfeits the whole
accumulated product --- and a monotone sequence distinguishes every
later point over every earlier one without distinguishing any point
categorically (Remark~\ref{gc:rem:periodization},
\S\ref{ac:sub:extrapolation}). The gradient picture therefore keeps
the urgency and discards only the step. Third, the quantitative heart
of the kind-claim --- the ratio of anthropogenic to natural risk ---
divides a numerator estimated by judgment by a denominator estimated
from the track record, and the track record is precisely what
survivorship conditioning prices: a record read under conditioning on
its reader shows rare natural catastrophe whether or not the
underlying rates were kind. The literature knows this as anthropic
shadow \cite{CSB10}, and the Precipice acknowledges it; the delta is
that this framework supplies the machinery
(\S\ref{gr1:sub:cond}) in which the caveat stops being a footnote and
becomes the subject --- the bias field is explicit, its magnitude is
computable in-model, and ``natural risk is small'' and ``we are among
survivors'' are hypotheses the finite record cannot separate. A
periodization whose different-in-kind claim rests on that ratio rests
on the one number the framework's own theorems say a survivor cannot
read off their record.

\paragraph{Delta three: the endpoint. Security decomposes.} The
Precipice's terminal state --- existential security --- is risk
brought low \emph{and known to be low}: a far side, reached and
certified, on which the Long Reflection can convene. The framework
decomposes this into one component it preserves and one it proves
unavailable. Risk brought low is anatomy work --- $\rho$ held
positive, patching capacity maintained against acquired capability ---
and it is continuous, never completed, with the losing condition
permanently adjacent; on this the two pictures can trade freely.
Known to be low is a certification, and the ceiling
(Theorem~\ref{gr:thm:ceiling}, Remark~\ref{gr:rem:ceilingscope})
forbids it to any finite observer: no bounded record places a
biosphere on the favorable side, forfeiture itself is silent
(Remark~\ref{hr:rem:silentforfeit}), and the unavailability points in
exactly the direction that would have been consoling. So ``the
Precipice, safely passed'' is not a state anyone could know themselves
to occupy, and a research program whose success condition is
unobservable in principle has misdrawn its own finish line --- the
metaphor is doing load-bearing work the theory cannot support, since a
precipice is the kind of thing that has a far side, and a
superposition one inhabits is not. What replaces the finish line is
not despair but the correct object: maintenance of the anatomy
condition, graded and continuous, with the Golden Point as
$\nu F$ --- a modal fixed point one may be inside without ever
certifying it --- rather than security as a date.

\paragraph{A corollary: the hinge dissolves rather than resolves.}
The debate downstream of the Precipice --- whether the present is the
``hinge of history'' \cite{MacAskill22, GM21} --- is conducted on a
shared assumption this framework denies: that hingeness is
categorical, so that the present either is or is not the special
moment. The framework's answer is neither side's. Every present is the
hingiest yet, by the monotone-stakes construction --- so the
influence-skeptic's base-rate argument attacks a claim the framework
never makes --- and no present is categorically special, by
continuity and by the tense-relativity of the conditioning's labels
(Remark~\ref{gc:rem:carryover}) --- so the hinge-proponent's
conclusion is an artifact of reading a label that moves with the
observer as a property of the moment. A graded hinge is had for free;
a categorical one is unavailable; and the binary framing, not either
answer to it, is what the formalism removes.

\paragraph{What survives, and the fences.} Much survives, and it
should be listed rather than implied: the urgency (re-derived, by
construction rather than by period); the dominance of the
anthropogenic channel (the trough channel is the unbounded one, and
the framework gives that asymmetry a functional form rather than an
estimate); the priority of risk-work over expansion (the
anti-Kardashev results agree with the Precipice against the
accelerationists); and most of the book's concrete policy content,
which re-grounds without loss on the gradient picture. What does not
survive is the quantity, the age, and the finish line. Three fences
bound the engagement. Nothing here consumes the Route World --- the
counter runs from the losing-condition analysis, the monotone-stakes
construction, and the ceiling, so a reader who rejects every
conditioning argument in this section keeps deltas one and three
whole and loses only the third prong of delta two. Nothing here
disputes the Precipice's numbers by offering rivals --- the
framework's claim is that the decisive ratio is unreadable from a
survivor's record, which is a reason to compute less, not
differently. And nothing here diminishes the practical convergence:
a reader persuaded by the Precipice to work on existential risk has
been persuaded of the right conduct, on this framework's own showing
(\S\ref{gc:sub:acting}), by an argument whose bookkeeping this
subsection corrects and whose urgency it keeps.

\paragraph{The provenance pattern, recorded once.} The engagement
instantiates the methodology companion's criterion exactly
\cite{Ben26Method}: the nearest neighbor's \emph{definition} already
contains the correction --- destruction of potential \emph{is}
closure of possibility --- while the neighborhood's density is what
kept the modal formalization unoccupied, each camp filing the vacancy
as territory already held. The prediction that discipline makes about
reception is stated here so it can be scored: this engagement will be
misfiled, confidently and in both directions, as ``a variant of Ord
with extra mathematics'' and as ``a rejection of existential-risk
studies.'' It is neither, and the deltas above are the scoring key.


% ---------------------------------------------------------------------
With the target fixed and its motivation stated, the question becomes
evidential: is there reason to believe the record was in fact produced
under it? The remainder of the section argues that there is, first by
saying what alignment with such a target would look like
(\S\ref{gc:sub:pastalign}), then by two independent reductios --- one on
natural history (\S\ref{gc:sub:conditioning}), one on human history
(\S\ref{gc:sub:humanhistory}) --- of which the second does not require
the premise the first leans on.

% ---------------------------------------------------------------------
\subsection{What alignment with the target would mean}\label{gc:sub:pastalign}

\subsubsection{Statement}\label{gc:subsub:pastalign-statement}

\begin{principle}[Past-Alignment; \textbf{[$\Phi$]}]
\label{gc:prin:pastalign}
Construct the \emph{Golden Nexus} by conditioning the Cosmic Nexus on
the Golden Point. From it, discard every stated attribute whose
referential event lies outside one's own past light cone --- in the
stricter reading, discarding the spacelike attributes as well as the
future-directed ones. What remains is the Nexus of Optimation.

Equivalently: the past light cone of Aethic attributes we actually
observe is the same one as the past light cone through our position in
the Aethic block universe of the Golden Nexus.\footnote{More precisely
than ``our actual light cone'': owing to random contingency in centric
unfolding, our own Aethus is a proper child of the Aetherian Nexus
rather than the Aetherian Nexus entire, so the object being compared is
the past light cone of the Aetherian Nexus's own Aethic block universe.}
\end{principle}

Two features of the statement are deliberate and should not be read past.

\emph{The restriction is to the past cone only.} The principle takes no
position on whether attributes in one's future light cone agree with the
near-optimal part of the Golden Nexus. That silence is not modesty but
scope: one cannot reliably infer future-cone attributes prior to
engaging in Golden reasoning, so a principle that assumed agreement
there would be assuming its own conclusion.
\S\ref{gc:sub:extrapolation} treats the forward question separately and
does not settle it.

\emph{The spacetime point is not privileged.} The principle is indexed
to one's own present because that is where the observable cone is
anchored, not because that event has any distinguished status. The same
construction at any other event returns that event's cone.

\begin{remark}[What the principle claims in the companion's own terms]
\label{gc:rem:oec}
Read against the Optimal-Earth Conjecture \cite{Ben26Optimal}, the
principle sharpens what ``optimal'' is optimal \emph{for}. The
conjecture there holds our record observer-optimal on effectively any
interrogable attribute. Past-Alignment says the optimand is not
vaguely observer-production but specifically the Golden Point: one asks
for the optimal Aethic history \emph{to the Golden Point}, and then
surgically reopens every attribute whose occurrence lies in the future
cone of the indexing event. If the principle holds, cropping the Golden
Nexus to one's past light cone reproduces Optimation exactly --- and it
is the exactness, not the plausibility, that would make this a
first-principles derivation of Optimation rather than a restatement.
\end{remark}

\subsubsection{The argument from prerequisite attributes}
\label{gc:subsub:pastalign-argument}

The supporting argument has two steps.

\emph{Step one: identify the attributes that position us for the Golden
Point.} Our record contains features which, whatever else they do, make
this biosphere unusually able to reach the Golden regime. The capacity
to build rocket vehicles is the cleanest example, since it is what would
allow the biosphere to persist past the eventual death of its star.

\emph{Step two: perturb the record and watch them fail.} Consider
Aethic perturbations in which the dinosaurs never go extinct, in which
animal life never evolves at all, or in which Africa does not exist or
develops very differently. Each comparison presupposes that the
perturbed history and ours share a common term against which a missing
capacity registers as missing, which ordinary object-individuation
cannot supply and \S\ref{to:sec} does: what is compared across these
histories is a Task-Earth and a task-lineage
(Definition~\ref{to:def:taskobject}), not an Earth and a lineage. With
that in place, each perturbation removes the capacity:

\begin{itemize}\itemsep2pt
\item A world in which the dinosaurs persist is one in which, by the
Accordance Principle's own verdict, the lineage does not reach our
level of intelligence, and so does not build rocket vehicles.
\item A world of fungi lacks the animal capacities the construction
requires.
\item A world without the specific African climatic sequence leaves
great apes arboreal, without the bipedal orientation and strategic
cognition the construction requires.\footnote{The climatic sequence in
Africa is what drove our great-ape ancestors to bipedality.}
\end{itemize}

The examples are illustrative rather than exhaustive, and the
generalization is what matters. By the substitution property of the
Accordance Principle \cite{Nexus}, the argument runs at \emph{any}
rung: a trait of ours may itself be placed in the conditioning seat, and
alternative natural histories are then deduced to fail its occurrence
almost surely. The maneuver requires the trait to be identifiable in
histories where it does not occur, which is a task-object
identification rather than an object one
(\S\ref{to:sub:licenses}). So for any Golden-relevant trait one names,
perturbations
of natural history systematically undermine its occurrence probability.

\begin{corollary}[Local optimality of the record; \textbf{[S]}]
\label{gc:cor:localopt}
If for every Golden-relevant trait --- individuated as a task-object,
per \S\ref{to:sec} --- the perturbed histories systematically
depress its probability, then our Aethus locally maximizes the
probability of the Golden Point among its neighbors --- and our natural
history is accordingly the past-light-cone restriction of the optimal
Aethus for the Golden Nexus.
\end{corollary}

\begin{remark}[The grade, stated honestly]\label{gc:rem:grade}
Corollary~\ref{gc:cor:localopt} is \textbf{[S]} and not
\textbf{[P]}. It inherits the substitution property's own license ---
rung-screening, flagged as such at the source \cite{Ben26Optimal} ---
and it establishes \emph{local} optimality among interrogated
perturbations, not global optimality over the space of histories. What
would strengthen it is a perturbation exhibited that leaves a
Golden-relevant trait's probability \emph{undiminished}; what would
break it is a systematic class of such perturbations.
\end{remark}


% ---------------------------------------------------------------------
Past-Alignment as stated records an agreement. It does not yet show the
agreement is more than coincidence, and the two arguments that follow are the
argument that it is not.

% ---------------------------------------------------------------------
\subsection{The conditioning principle, and the two arguments for it}\label{gc:sub:conditioning}

Past-Alignment as stated is an observation about agreement. The
following argues that the agreement is not coincidental --- that
Nexic natural history was never conditioned on Optimation in the first
place, but on the Golden Nexus throughout.

\begin{principle}[Golden-Conditioning; \textbf{[S]}, by reductio]
\label{gc:prin:conditioning}
Suppose, for contradiction, that the Accordance Principle conditions on
personal observer-existence rather than on the Golden Point.

Observerhood is robust under Aethic perturbation. One may strip
probabilistic prerequisites of the Golden Point --- rocket capability,
the presence of industrialization-driving coal --- and replace them by
their complements without disturbing one's standing as an observer. A
world in which Earth simply lacked the resources to begin an industrial
revolution is a world whose inhabitants are no less observers.

Now select a handful of such attributes: necessary for the Golden Point,
unlikely to survive reopening in the Aetherian Nexus, and admissible in
the sense of Remark~\ref{gc:rem:boundrequires} --- reopening them must
leave the dependent structure intact, so that the resulting Aethae
remain individuated by the same task-objects as ours
(\S\ref{to:sub:licenses}). The totality of
the counterfactual Aethae so obtained carries \emph{substantially more}
Aethic weight than the Aetherian Nexus, while carrying a comparatively
minute probability of the Golden Point.

If Accordance were conditioned on observer-existence alone, it would
therefore place us in that high-weight counterfactual sum rather than in
the Aetherian Nexus. But the Aetherian Nexus is in fact our parent
Aethus. Contradiction. Hence the Accordance Principle is not conditioned
on observers, so long as ``observer'' is robust enough to survive these
perturbations.
\end{principle}

\subsection{The undeniable case: human history}\label{gc:sub:humanhistory}

Principle~\ref{gc:prin:conditioning} admits two arguments, and they
differ in one respect that decides their order. Both need observerhood
to be robust under the perturbations considered --- that one can strip
prerequisites of the Golden Point without disturbing anyone's standing
as an observer. On natural history a sufficiently demanding notion of
``observer'' can contest that premise (Remark~\ref{gc:rem:circularity}).
On human history it cannot: no one holds that human observers required
Alexander's survival, or Newcastle's coal, or Newton. Observers existed
before each and would have continued after. So the human-history
argument is the one on which the robustness premise is simply not
deniable, and it is put first for that reason; the natural-history
argument follows as the same reasoning in a domain where the premise is
contestable and the conclusion correspondingly weaker.

The argument also does something the natural-history one cannot. What
these contingencies deliver is not observers but specifically the
technological condition --- and observer-conditioning is indifferent to
every one of them while Golden-conditioning is not. That isolates what
the conditioning must be sensitive to, and it says something about the
companion that the companion cannot say about itself: its
conditioning target, an eventual observer in one's biosphere, is exactly
right for the pre-technological record and provably insufficient beyond
it. Not a criticism --- the companion reads the pre-human record and
stops, by design --- but the reason for stopping becomes visible only
here.

\subsubsection{Flukes establish their own significance}
\label{gc:subsub:flukes}

\begin{proposition}[Significance from route-deletion, at the
self-sampling price; \textbf{[S]}]
\label{gc:prop:flukes}
Let $F$ be an event established as a \emph{fluke} --- of low outlook
probability $p$ --- with no claim yet made about its significance.
Suppose $F$ were statistically independent of the Golden Point. By the
Second Law its Golden Probability equals its outlook probability, $g=p$,
and we should expect with probability $1-p$ to occupy the counterfactual
$\neg F$. We do not. Hence $F$ is not independent of the Golden Point,
and by the Third Law its Golden Linkage exceeds the L-threshold
\cite{Nexus, GoldenScience22} --- in the notation of
\S\ref{gc:sub:imported}, the Accordance bound \eqref{gc:eq:accordance}
with $H$ the Golden Point, $A:=F$, $A':=\neg F$, so that
$\PP(G\mid\neg F)/\PP(G\mid F)\lesssim p/(1-p)$ with $\kappa$-fold
violations priced by \eqref{gc:eq:accordexceed}
(Remark~\ref{gc:rem:importscope} records the translation).

Two premises deserve their names before the conclusion is stated at
its licensed strength. The step ``we should expect with probability
$1-p$ to occupy $\neg F$'' is a \emph{self-sampling} premise ---
location weighted within a reference class spanning the $F$- and
$\neg F$-worlds --- of the lineage this corpus polices elsewhere
\cite{Bos02}; it is adopted exactly once, at the root, and
Remark~\ref{ra:rem:noroster} governs its normalizer. Granting it,
what the proposition delivers is a \emph{conditional-likelihood}
claim: \emph{for route-deleting $F$}
(Remark~\ref{gc:rem:flukescope}), observing $F$ raises the Golden
linkage above the L-threshold. Significance is deduced given the
counterfactual's route-deleting character --- and establishing that
character is case work, as the scope remark says, not a free rider on
the flukiness. Aggregating over a set of route-deleting flukes, the
union of their counterfactual Aethae carries a far greater Aethic
weight than our own, and the principle refuses to book our location
within the lesser one as brute luck.
\end{proposition}

% [Referee pass, Aug 21: as stated the proposition's template runs on
% any observed fluke whatever, including incidental ones.  The
% restriction that makes it sound is stated here.]
\begin{remark}[What the proposition quantifies over, and the
counterfactuals it excludes]\label{gc:rem:flukescope}
The template above --- $F$ is improbable, we observe it, therefore $F$
is linked to the Golden Point --- runs on \emph{any} observed
improbability if nothing restricts the range of $F$, and would then
license absurdities: the precise configuration of a Cretaceous
raindrop is astronomically improbable and is not thereby significant.
The restriction that blocks this is already implicit in the
proposition's use and is stated here, because a reader is entitled to
it before the stockpile rather than after.

\emph{The counterfactual must delete a route, not permute one.} Write
$\neg F$ for the alternative. The inference is licensed when
$\PP(G\mid\neg F)$ is genuinely lowered by the substitution --- when
the counterfactual removes a prerequisite of the Golden Point --- and
is void when $\neg F$ merely reshuffles which particular events fill
the same structural role. The raindrop case is a permutation: some
configuration obtains, and every one of them leaves the downstream
dependency structure untouched, so $\PP(G\mid\neg F)=\PP(G\mid F)$ and
the bound \eqref{gc:eq:accordance} is satisfied with room to spare
rather than violated. Improbability alone therefore establishes
nothing; improbability \emph{of a route-deleting attribute}
establishes the linkage, and the proposition should be read with that
restriction throughout.

\emph{The restriction is not free, and it is what the stockpile is
selected for.} Deciding whether a counterfactual deletes or permutes
is a claim about dependency structure and is not delivered by the
attribute's weight --- which is
Remark~\ref{gc:rem:flattening}'s point arriving early: a marginal test
on weights recovers the hubs of the dependency graph, and the
route-deleting attributes are exactly the hubs. So the criterion is
principled and its application is \textbf{[S]}, case by case. This is
also the honest register in which the stockpile's items differ. The
coal has the route-deleting character on the economic historians' own
analysis --- Allen's price ratios make mechanization pay in one place
and nowhere else, so the counterfactual removes the route rather than
relocating it \cite{Allen09} --- while the Granicus is weaker on
exactly this axis, since the claim that no comparable transmission
survives Alexander's death at the Granicus is a strong fragility
premise about the Hellenistic world and not a computed dependency.
That ordering is the same one Remark~\ref{gc:rem:narrativebias} reaches
from the evidential side, and the agreement of two independent
criteria on which items are strongest is worth more than either.
\end{remark}

\begin{remark}[The retrospective-target null, faced at the site]
\label{gc:rem:sharpshooter}
The strongest rival to the proposition is not the determinist's
substitutability but the \emph{retrospective-target} null: any
realized history, read backward, contains improbable events that were
prerequisites of whatever that history's distinctive features happen
to be, so a stockpile of flukes pointing at the realized outcome is
expected under no conditioning whatever --- the sharpshooter who
paints the target around the shots. The paper's answer exists but
lived two sections away; it belongs here. First, the target is
\emph{ex ante}: the Golden Point is picked out by the fixed-point
structure of \S\ref{gc:subsub:cascade} --- a future event defined
before and independently of which record happens to obtain --- so the
target is drawn before the shots in the only sense the null requires.
Second, the evidential force claimed is \emph{comparative}, and the
comparison is stated where the historical cases are heard
(\S\ref{gc:sub:historicalcases}): consistency with a hypothesis is
not evidence for it over a rival the same data also fit. Third, what
survives that discipline is the human-history leg's genuinely
distinctive content: the stockpile's flukes are prerequisites of the
\emph{technological} condition specifically, to which
observer-conditioning is provably indifferent --- and that asymmetry
between the rivals, not raw improbability, is what the proposition's
linkage tracks.
\end{remark}

\begin{remark}[This settles the contingency-determinism debate, in one
direction]\label{gc:rem:determinism}
Historians divide over whether the path to modernity was contingent or
overdetermined, and the dispute is usually conducted by weighing
narratives. Proposition~\ref{gc:prop:flukes} makes it an inference.

The determinist holds that the apparent flukes were substitutable: no
Alexander, then someone else; no British coal, then some other route.
Substitutability is exactly high $b$ --- a good chance of the Golden
Point without $F$ --- hence low Golden Linkage, hence $g \approx p$,
hence the expectation that we occupy $\neg F$. The Third Law --- the
Accordance bound \eqref{gc:eq:accordance} with $H$ the Golden Point ---
bounds it quantitatively: $b/a < p/(1-p)$, so the flukier the event, the more
severely substitutability is constrained by our observing it. The
determinist is not refuted by a better narrative; they are constrained
by an inequality whose inputs are the event's improbability and its
route-deleting character (Remark~\ref{gc:rem:flukescope}) --- nothing
else.
\end{remark}

\subsubsection{The stockpile}\label{gc:subsub:stockpile}

Two cases, deliberately ordered so that the weaker rhetorically comes
first.

\emph{A single sword-stroke.} At the Granicus in 334 BCE, two months
into an eleven-year campaign, the Persian noble Spithridates brought a
battle-axe down on Alexander's helmet and split its crest; Cleitus the
Black severed his arm before the second blow \cite{Arrian,
PlutarchAlex}. Had the stroke landed, the expedition would have ended
within days of beginning, and with it the Hellenistic world, its
Alexandrias, and the transmission through which much of what followed
ran.

\emph{Where the coal happened to be.} Pomeranz argues that England was
a ``fortunate freak'' whose coal lay close to water and ports, making
steam economically feasible, while China's coal sat in the northwest,
far from the Jiangnan textile economy that might have used it ---
a geological contingency placing coal and the Americas nearer the
western than the eastern end of Eurasia \cite{Pomeranz00}. Allen adds
the labor side: Newcastle carried the highest ratio of labor to energy
costs in the world in the early eighteenth century, and only under that
ratio did mechanization pay \cite{Allen09}. His summary judgment is
an Imperfection Magnitude Drop stated by an economic historian who has
never encountered one: there was a single path to the twentieth century,
and it ran through northern Britain.

A third case is available but requires reframing. Diamond's account of
Eurasian advantage \cite{Diamond97} is \emph{determinist} at the level
of human history --- geography made the outcome near-inevitable --- and
is criticized on that ground. What the present argument takes from it
lies one level down: the distribution of domesticable species and the
continent's east--west axis are contingent \emph{across possible
Earths} while determining \emph{within} ours. That is precisely the
Accordance profile, low outlook probability with high Golden Linkage,
and Diamond's determinism is at the wrong level rather than mistaken.

\subsubsection{Why this leg is the stronger one}
\label{gc:subsub:strongerleg}

Two reasons, and the second is the more important.

\emph{The robustness premise becomes uncontroversial.} The natural-history
reductio needs it granted that observers survive the removal of the
end-Cretaceous impact, which can be resisted. Nothing comparable is
available here. No one holds that human observers required Alexander's
survival, or Newcastle's coal seams, or Newton. Humans existed for
millennia before each and would have continued after their removal. The
premise that carried the weight and could be denied is here simply not
deniable.

\emph{The conclusion sharpens.} What these contingencies deliver is not
observers but specifically the \emph{technological} condition ---
skyscrapers, the internet, the capacity to leave a dying star.
Observer-conditioning is indifferent to every one of them;
Golden-conditioning is not. So the human-history argument does not merely repeat
the first at lower risk; it isolates what the conditioning must be
sensitive to.

\begin{remark}[The narrative-bias objection, and the test that answers
it]\label{gc:rem:narrativebias}
The serious objection is selection in the record rather than in the
world. Dramatic near-misses make better stories, are more often written
down, and are more often repeated; a history dense with them may report
a fact about historiography.

The framework predicts an asymmetry that narrative bias does not, and
the prediction is already on record. Attributes with high Golden Linkage
magnitude but \emph{low connectivity} to other attributes should have
low Golden Probability --- lucky avoidance of adversity outranks lucky
survival of it, so isolated close calls should be \emph{rare} in the
record while densely connected contingencies should be
\emph{numerous} \cite{GoldenScience22}. Narrative bias inflates both
alike; the framework inflates one and suppresses the other.

The discrimination favors the framework for a plain reason: \emph{the strongest cases are the dull ones}. Nobody narrates
the Newcastle labor-to-energy price ratio, and no popular account turns
on Jiangnan's distance from northwestern coal seams. Yet that is where
the peer-reviewed economic history converges \cite{Pomeranz00,
Allen09}. Were narrative selection generating the pattern, the
undramatic structural contingencies would be missing --- and they are
the best-attested items in the file. The Granicus is the vivid
illustration; the coal is the evidence, and the ordering of
\S\ref{gc:subsub:stockpile} inverts the rhetorical one deliberately.
\end{remark}


% ---------------------------------------------------------------------
Both legs concern the past cone, and neither speaks to the future. That
silence is deliberate --- \S\ref{gc:subsub:pastalign-statement} recorded
it as scope rather than modesty --- and the three positions available on
the forward question are set out next without adjudication.

% ---------------------------------------------------------------------
\subsection{The contestable case: natural history}\label{gc:sub:naturalhistory}

The same reductio, run on the natural-historical record. Its premise
--- that observerhood survives the perturbations considered --- is
contestable here in a way it was not above, and the argument is
correspondingly the weaker of the two; it is included because it reaches
further back, and because it is where the conditioning claim was first
posed.

\begin{remark}[The Trace Connection, named in 2022]
\label{gc:rem:trace}
The claim this subsection advances was posed as an open problem, under
its own name, four years before the present argument. Our
\emph{Golden Science} of August 2022 \cite{GoldenScience22} sets out the
Golden Principle in four numbered points, of which the third states that
our observations are far closer to representing those of an average
Golden-Point biosphere than those of an average position in the universe
--- Past-Alignment in prose --- and the fourth states that some
connection, causal or causal-mimicking, holds between our existence at
this time and place and the Golden Point itself. That paper then names
the connection and declines to supply it: it is to be called \emph{the
Trace Connection} until it is officially derived.

\S\ref{gc:sub:conditioning} is offered as that derivation. What the 2022
statement supplies, and what no later argument could manufacture for
itself, is that the gap was identified, named, and dated before any
candidate filling existed --- so the reductio cannot be read as having
been reverse-engineered to fit a conclusion already in hand. The same
document also carries, in prose and without the stochastic apparatus of
\S\ref{gr1:sec}: the hazard-function framing of biosphere survival; the
observation that an exponentially decaying hazard yields a strictly
positive limiting survival probability; the division of the failure mode
into internal and external channels, together with the insistence that a
Golden biosphere must lessen \emph{both} rather than either; the thesis
that domination is inefficient because it feeds the internal channel,
which is $\Phi3$; the Kardashev critique including the
self-consumption argument; and the definition of the Golden Point as the
point past which a biosphere cannot regress, drawn there against the
Hubble horizon. The formal content of
\S\S\ref{gr1:sec}--\ref{ad:sec} is later; the apparatus it formalizes is
not.
\end{remark}

\begin{corollary}[What remains; \textbf{[S]}]\label{gc:cor:remains}
What the reductio leaves standing is the empirical artifact of
Past-Alignment itself: conditioning on the existence of the Golden Point
somewhere in the Aethic block universe reproduces exactly the
inferential form of the Accordance Principle that yields the
stated-attributes of Nexic reasoning. Among the contingent natural
histories one might find oneself at the end of --- a fungal world, a
dinosaurs-survive world, a wholly lemurine one --- the one that must be
conditioned on to reproduce the Nexic state of affairs we actually
observe is the one in which the Golden Point exists.
\end{corollary}

\begin{remark}[A second robustness, and why the target's modal form was
the right one]\label{gc:rem:finitehorizon}
Principle~\ref{gc:prin:conditioning} turns on a robustness: observerhood
survives perturbations the record does not, so a conditioning
indifferent to the difference cannot have produced the difference. The
schedule channel of \S\ref{sch:sub:schedule} exhibits a second target
with the same defect, and it is nearer to hand than observerhood.

\emph{Survival to a finite horizon is robust in exactly that way.} A
biosphere may survive any stretch one names without having climbed its
schedule at all --- by not being struck, at bare background, on the
lucky branch that Remark~\ref{sch:rem:asymptote} shows never leaves the
support of a finite-horizon conditioning. So conditioning on
having-lasted leaves the record unexplained for precisely the reason
conditioning on observers did: it is satisfied by histories in which
nothing was climbed.

What survives the objection is the asymptotic form, and the point worth
extracting is that the target was already stated in it.
Definition~\ref{gc:def:departure} gives Goldenness as $\exists\forall$
--- some course admits \emph{no} foreclosing continuation --- which is a
condition on the whole of a future and not on any initial segment of
one. A reading of the Golden Point as a passage completed at some finite
time would have inherited the robustness above and failed here. The
modal form is therefore not a technical nicety of the definition; it is
what makes the target capable of doing the work
Principle~\ref{gc:prin:conditioning} asks of it.

The cost is recorded rather than absorbed: this raises what the
conditioning must be, and by \S\ref{gr:sub:ceiling} no one is entitled
to know that it obtains.
\end{remark}

\subsubsection{Bounding the counterfactual weight}
\label{gc:subsub:weightbound}

The reductio of \ref{gc:prin:conditioning} turns on a single inequality:
that the counterfactual Aethae obtained by reopening Golden-Point
prerequisites carry substantially more weight than the Aetherian Nexus.
Stated bare, that is an assertion, and an objector may simply deny it ---
perhaps coal is common given a Carboniferous, perhaps the metallurgical
sequence is near-forced given hands and fire. This subsection replaces
the assertion with a construction, a bound, and a statement of exactly
what empirical input remains outstanding.

\paragraph{The construction.} Fix stated-attributes
$a_1,\dots,a_k$, each a prerequisite of the Golden Point and none a
prerequisite of observerhood. Let $\mathcal A$ denote the Aetherian
Nexus, in which every $a_j$ obtains, and let $\mathcal A_j$ denote the
Aethus obtained by reopening $a_j$ alone to its logical complement,
holding the remaining $k-1$ at their realized values.

\begin{lemma}[Disjointness; \textbf{[P]}]\label{gc:lem:disjoint}
The Aethae $\mathcal A_1,\dots,\mathcal A_k$ are pairwise disjoint, and
each is disjoint from $\mathcal A$. For $j\neq j'$, the pair
$\mathcal A_j$ and $\mathcal A_{j'}$ disagree on \emph{both} slot $j$
and slot $j'$; and $\mathcal A_j$ disagrees with $\mathcal A$ on slot
$j$.
\end{lemma}

\begin{proposition}[Weight ratio; \textbf{[P]}]\label{gc:prop:ratio}
Write $\alpha_j$ for the probability that slot $j$ takes its realized
value, conditional on the remaining slots taking theirs. Then
\begin{equation}\label{gc:eq:ratio}
\frac{W(\mathcal A_j)}{W(\mathcal A)}=\frac{1-\alpha_j}{\alpha_j},
\end{equation}
and by Lemma~\ref{gc:lem:disjoint} the single-flip Aethae may be summed:
\begin{equation}\label{gc:eq:bound}
R \;:=\; \frac{W_\times}{W(\mathcal A)}
\;\ge\; R_1 \;=\; \sum_{j=1}^{k}\frac{1-\alpha_j}{\alpha_j}.
\end{equation}
\end{proposition}

\begin{remark}[What the bound does and does not require]
\label{gc:rem:boundrequires}
Three features of \eqref{gc:eq:bound} are worth isolating, because they
are what make it usable.

\emph{Independence is not required.} Only disjointness is, and that
holds by construction. Dependence among the $a_j$ changes what the
individual ratios \emph{are} --- the $\alpha_j$ become conditional
rather than marginal, as stated --- but not whether the terms may be
added. This matters because the attributes at issue are plainly not
independent, and an argument needing them to be would be unusable.

\emph{It is a lower bound, and deliberately a loose one.} The
single-flip Aethae are a vanishing fraction of the counterfactual set:
the full sum over all Aethae differing in at least one slot is
$\prod_j\alpha_j^{-1}-1$, larger by many orders. But the full sum
requires the joint structure, and the reductio needs only that $R\gg1$.
An undercount that suffices is preferable to an exact figure that is
contestable.

\emph{Admissibility is governed by the Imperfection Magnitude Drop
criterion.} A perturbation enters the sum only if it leaves the
dependent structure intact --- if reopening $a_j$ does not re-randomise
what depends on it. A perturbation that triggers a drop does not yield a
clean single-slot flip, and Lemma~\ref{gc:lem:disjoint} fails for it.
Coal in an inaccessible location is admissible on this criterion: it may
be probabilistically catastrophic without re-randomizing the structure
downstream. A relocation of the lineage's speciation to another
continent is not, since the material that later mediates the outcome is
not present there to be held fixed.
\end{remark}

\paragraph{How large must $R$ be, and how improbable must the
attributes be?} The reductio requires only $R\gg1$. Evaluating
\eqref{gc:eq:bound} for $k$ attributes each at probability $\alpha$:

\begin{center}\small
\begin{tabular}{@{}rrrrr@{}}
\toprule
$\alpha$ & $k=1$ & $k=3$ & $k=5$ & $k=8$\\
\midrule
$10^{-1}$ & $9$ & $27$ & $45$ & $72$\\
$10^{-2}$ & $99$ & $297$ & $495$ & $792$\\
$10^{-3}$ & $999$ & $3.0\times10^{3}$ & $5.0\times10^{3}$ & $8.0\times10^{3}$\\
$10^{-4}$ & $1.0\times10^{4}$ & $3.0\times10^{4}$ & $5.0\times10^{4}$ & $8.0\times10^{4}$\\
\bottomrule
\end{tabular}
\end{center}

\noindent The bar is far lower than the informal argument suggested. A
\emph{single} attribute at $\alpha=10^{-3}$ delivers $R\sim10^{3}$
unaided; no portfolio is needed, and the burden reduces to establishing
one prerequisite as genuinely improbable. Correspondingly, the objection
that the counterfactual sum might be comparable requires \emph{every}
candidate prerequisite to have $\alpha$ of order one --- that coal
placement, metallurgical sequence, and the rest are each more likely
than not given the rest of the record --- which is a considerably
stronger claim than the objection is usually asked to bear.

\paragraph{Estimating $\alpha$: the sister-case method.} The remaining
empirical input is supplied by an instrument the companion already
develops, and it should be used rather than reconstructed
\cite{Nexus}.

By the Theorem of Aimless Optimation, Optimation is the Cosmic Nexus
conditioned on the attribute $H_N$ that an Aetherian exists in one's
biosphere; and if an attribute $X$ is statistically independent of
$H_N$, the two Nexae assign it the same weight. Independent attributes
are therefore \emph{controls}: their observed frequency estimates the
Cosmic Nexus probability directly, unlifted by conditioning. The
companion's worked case is the Catarrhini node --- Old World monkeys
furnish over a hundred extant species across twenty-five million years
with no attainment of ape-level intelligence, and their fate being
independent of $H_N$, that frequency transfers as the baseline for our
own lineage at the same node.

This is exactly the quantity \eqref{gc:eq:bound} requires. For an
attribute $a_j$ admitting such a control, $\alpha_j$ is read off the
control's frequency, and its contribution to $R_1$ follows.

\begin{remark}[A corollary on once-only attributes, and the independence
check that governs it]\label{gc:rem:onceonly}
A corollary of the same theorem bears on attributes observed exactly
once, and it should be read as an inference rather than a
consistency claim. Once a first occurrence has secured what the Golden
Point requires, a \emph{second} occurrence is statistically independent
of $H_N$ --- it changes nothing about whether an Aetherian exists ---
and by the theorem it therefore retains its Cosmic Nexus weight and is
not lifted. Optimation supplies what is required and does not
overshoot.

The consequence is a discriminating observation. Write $\lambda$ for
the unconditioned expected count across all opportunities. Under the
hypothesis that $\lambda\ll1$, Optimation must lift the attribute for
$H_N$ to obtain, and lifts it to exactly one; \emph{exactly one
occurrence is what that hypothesis predicts}. Under the hypothesis that
no lift was needed, occurrences are Poisson at rate $\lambda$ and
exactly one is a coincidence of probability $\lambda e^{-\lambda}$. The
likelihood ratio between them is $\left(\lambda
e^{-\lambda}\right)^{-1}$:
\begin{center}\small
\begin{tabular}{@{}rrrrrr@{}}
\toprule
$\lambda$ & $1$ & $3$ & $5$ & $10$ & $20$\\
\midrule
ratio favoring lift & $2.7$ & $6.7$ & $30$ & $2.2\times10^{3}$ &
$2.4\times10^{7}$\\
\bottomrule
\end{tabular}
\end{center}
So a single occurrence tells strongly against the configuration having
been \emph{common} --- $\lambda\ge5$ is disfavored thirtyfold or worse
--- while discriminating only weakly between a genuinely minute
$\lambda$ and one of order unity. Its evidential role is accordingly to
exclude the objection that the prerequisites were unremarkable, not to
supply the magnitude. The magnitude comes from the independent
per-opportunity estimate, and where that gives $\lambda\ll1$ the
attribute contributes approximately $1/\lambda$ to $R_1$.

The corollary is only as good as its individuation, and this is where
the difficulty sits. Rather than proceed directly to an estimate, we
calibrate the instrument on two attributes for which the data are good.
Both return nothing, and the pattern in \emph{how} they return nothing
is what fixes the strategy.

\paragraph{Calibration I: coal.} Britain's coal is the obvious first
candidate, and it fails at the endowment grain. Carboniferous and
younger basins are known on virtually every continent --- the
Appalachians, the Ruhr, the Donets, the Kuznetsk, the Nord-Pas de
Calais--Saar--Lorraine field, and others --- giving eight or more on any
reasonable individuation, and Britain's is not the most favorable, the
Appalachian seams being unusually large and unusually exposed. With
$\lambda\approx8$ we obtain $\alpha=1-e^{-\lambda}\approx0.9997$ and a
contribution to $R_1$ of about $3\times10^{-4}$: nil.

The verdict is confirmed by the author whose thesis is most at stake.
Allen observes that the Ruhr went unexploited not for want of coal but
because Dutch peat held fuel prices down, removing the return on
improving Rhine transport; and he draws the counterfactual explicitly,
that had German coal been developed three centuries earlier the
industrial revolution might have been a Dutch--German achievement rather
than a British one \cite{Allen09}. He identifies Belgium, around
Li\`ege and Mons, as possessing a comparable ratio of wages to energy
cost. So the claim that the path to the twentieth century ran through
northern Britain was always about the \emph{conjunction} of cheap energy
with dear labor and the institutions to exploit both --- never about
the seams, which were widely available. A reader who recalls that claim
should read it that way, and the estimator's verdict agrees with Allen
rather than contradicting him, reached independently and from geology.

\paragraph{Calibration II: domesticable megafauna.} A second candidate,
drawn from a different literature, fails in the same manner. Diamond's
count gives fourteen large mammal species domesticated of some hundred
and forty-eight candidates, with thirteen of the fourteen wild ancestors
confined to Eurasia \cite{Diamond97} --- a striking asymmetry, and the
basis of a well-known thesis. But the relevant question for the
estimator is not how the successes distributed; it is how many
\emph{continents} possessed a workable endowment. Diamond's own
candidate counts are seventy-two for Eurasia against fifty-one for
sub-Saharan Africa, which is not the profile of a rare attribute. Taking
continents as the opportunities gives $\lambda\approx2$ and a
contribution of about $0.16$: nil again.

Diamond's \emph{conjunction} does better --- megafauna together with an
east--west principal axis and the crop package on one landmass gives
$\lambda\approx0.16$ and a contribution near $5.8$ --- which is a real
figure, and still four orders below what a single attribute would need
to carry the reductio alone.

\begin{proposition}[The bound on \emph{parallel} attributes;
\textbf{[S]}]\label{gc:prop:squeeze}
The two failures share a cause, and it is structural rather than
accidental. Call an attribute \emph{parallel} when its opportunities are
a set of co-existing sites --- landmasses, basins, continents --- so
that $\lambda=Np$ with $N$ the site count. For an attribute individuated
at continental scale, $N\approx5$. To obtain
$(1-\alpha)/\alpha\sim10^{3}$ then requires $\lambda\sim10^{-3}$, hence
a per-site probability of order $2\times10^{-4}$ --- a claim that cannot
responsibly be defended for any attribute of the form ``this landmass
possesses feature $X$''. Deep time can deliver minute $\lambda$ for
parallel attributes because its site count is astronomical; recent
history cannot, because its site count is five.

Coal and megafauna are both parallel attributes. That, and not their
recency, is why they failed.
\end{proposition}

\paragraph{Serial attributes escape the bound.} Proposition
\ref{gc:prop:squeeze} is a claim about parallel attributes only, and the
restriction is not a technicality --- it identifies the wrong turn that
coal and megafauna represent. An attribute may instead be
\emph{serial}: a conjunction of successive branch points within a single
trajectory, each of which had to go one way for the next to be reached.

\begin{proposition}[Serial attributes; \textbf{[S]}]
\label{gc:prop:cascade}
Let an attribute consist of $m$ successive branch points, the $i$th
resolved favorably with probability $p_i$ conditional on its
predecessors. Then $\alpha=\prod_{i\le m}p_i$ and the contribution to
$R_1$ is $\alpha^{-1}-1$, which grows \emph{exponentially in $m$}. At
$p=0.75$ throughout: $m=8$ contributes about $9$, $m=16$ about $99$,
and $m=25$ about $1.3\times10^{3}$.
\end{proposition}

\noindent The opportunity structure is what differs. A parallel
attribute is bounded by how many sites exist, and there are five
continents. A serial attribute is bounded by how many branch points the
trajectory contains, and a well-articulated historical episode contains
many. Recent history is therefore not weak evidence; it was being
individuated in the form that makes it weak.

This is how the candidates should be stated. Not ``Alexander existed'',
which is a parallel attribute of no force, but \emph{that he survived
and succeeded through the entire sequence} --- the wound at the
Granicus, Issus, Gaugamela, the near-fatal injury among the Malli, the
absence of a disabling illness before consolidation, and the
Hellenization surviving the immediate fragmentation of the empire ---
each a branch point with its own $p_i$. Likewise the Roman survival of
the Second Punic War is a sequence, not a fact: the recovery after
Cannae, Hannibal's failure to march, the holding of the Italian allies,
Scipio's African campaign. Likewise the passage of Christianity from a
provincial sect to an imperial religion, and its part in what followed.
Each such episode is one attribute in \eqref{gc:eq:bound}, and each
carries a contribution set by the length of its cascade rather than by
the number of continents.

\begin{remark}[What serial attributes cost]\label{gc:rem:serialcost}
The gain in magnitude is paid for in the independence check, and the
price should be stated plainly, since it is where such an estimate would
fail.

For each branch point one must hold that its failure had no adequate
\emph{substitute} --- that had Alexander died at the Granicus, the
Hellenization of the East does not simply proceed under another
commander. History is rich in substitutes, and each $p_i$ that admits
one contributes nothing, since the attribute is then not load-bearing
for $H_N$ at that step. The estimator's discipline accordingly transfers
from grain, where it sat for parallel attributes, to substitutability.

Two consequences follow. Cascade length must be counted in
\emph{substitution-free} branch points, not in dramatic incidents, which
will reduce $m$ considerably below the narrative count. And the
per-step probabilities must be conditional on the predecessors, since
the steps are plainly not independent --- a commander who has survived
four battles is not a random draw at the fifth. Both corrections push
$\alpha$ upward, and a careful estimate should expect the eventual
contribution to fall well short of the naive product.
\end{remark}

\begin{remark}[Individuate the branch point by its role, not its
occupant]\label{gc:rem:roleindividuation}
The substitutability audit has a corollary that governs how each branch
point must be \emph{stated}, and getting it wrong is the most likely
route to an inflated estimate.

The counterfactual at each step must be taken over the \emph{role}, not
over the occupant. The question at the relevant branch point is not
whether Christianity in particular arose and spread, but whether
\emph{anything} arose that discharged the role it discharged --- in the
window, with the reach, and with the consequences for what followed. So
individuated, the competing cults of the period are not rivals to the
hypothesis but candidate fillers of the same role, and their availability
raises $\alpha$ at that step rather than lowering it: where a role admits
$\kappa$ independent candidates each of probability $q$, the
role-probability is $1-(1-q)^{\kappa}$, not $q$.

Role individuation is therefore the \emph{conservative} choice, and that
is precisely its merit. It concedes the substitution objection in
advance rather than meeting it afterward: a critic who replies that some
other movement would have served is not refuting the estimate but
restating the quantity already being estimated. What survives such an
audit is a residue of steps at which no filler was available with
comparable probability --- and the claim, where a cascade is offered, is
about that residue.

It also relocates where improbability may legitimately be found. It is
not to be sought in the striking identity of the occupant, which is
where narrative attention naturally goes, but in the narrowness of the
window in which \emph{any} occupant could have appeared. The
transitional periods are the interesting ones for that reason, and the
dramatic incidents largely are not.
\end{remark}

\begin{remark}[The limitation is present ignorance, not the structure]
\label{gc:rem:measurement}
One further point about the status of these estimates, and it
cuts in an unusual direction.

The obstacle to evaluating $\prod_i p_i$ for any actual cascade is not
that the quantity is ill-defined or that the method is unsound. It is
that counterfactual history is not yet in a state to supply the
$p_i$ with defensible error bars --- the substitutability audits and the
window widths that Remark~\ref{gc:rem:roleindividuation} requires are
matters on which the discipline has views but not measurements. The
uncertainty is \emph{ours}, and contemporary; it is not a property of
the cascades.

Two things follow. The paper should not, and does not, offer a numerical
cascade estimate: what it offers is the structure
\eqref{gc:eq:bound}, the serial--parallel distinction, and the
individuation discipline, with the arithmetic left for inputs that do
not yet exist. And the expectation that such inputs will exist is a
substantive bet rather than a certainty --- one that quantitative
counterfactual history will mature to the point of supplying audited
branch-point probabilities. If it does, the enumeration becomes
routine and the bound follows. If it does not, the weight claim remains
structurally established and empirically unfilled, which is a weaker
position than we would like and a stronger one than assertion
\textbf{[O]}.
\end{remark}

\begin{remark}[An unnamed significance is a prediction, and it has a
falsification test]\label{gc:rem:unnamed}
Remark~\ref{gc:rem:measurement} leaves the framework in a position that
looks, at first, like an unfalsifiable auxiliary. We do not know what
the establishment of a universalizing cult in the late-antique
Mediterranean forced through downstream; the analysis nonetheless
concludes that something Golden-relevant was forced through, on the
strength of \ref{gc:prop:flukes} alone. Stated carelessly this is the
canonical bad move --- every fluke acquires an unnamed significance,
nothing counts against the framework, and the Accordance analysis
becomes a device for dignifying whatever happened.

It is not that, and the reason is that the unnamed significance is a
\emph{prediction} with a definite content and a definite refutation.

\emph{Content.} For a cascade $C$ satisfying the two admissibility
conditions --- genuinely improbable, and load-bearing in the sense of
Remark~\ref{gc:rem:roleindividuation} --- the claim is that there exists
a consequence of $C$ on which the probability of the Golden Point
depends, and that it will be identified. The claim is made before the
consequence is known, which is what makes it a prediction rather than a
description.

\emph{Refutation.} Exhibit a route to the Golden Point that does not
require $C$. If the terminus is reachable without the cascade, then
$C$ was not load-bearing for it, and the predicted consequence does not
exist to be found. This is the substitutability audit of
Remark~\ref{gc:rem:serialcost} run at the \emph{terminus} rather than at
an intermediate --- which is the only place it can be run when the
intermediate is unnamed, and which is a harder audit rather than an
easier one, since the chain to be traced is longer.

\emph{And the criterion is diachronic.} A single unnamed significance
is a debt; a growing ledger of them is a symptom. The framework is
progressive if the ratio of identified to unidentified consequences
improves as counterfactual history matures, and degenerating if the
unnamed list only lengthens while the named one does not. We state this
as the test we intend the program to be held to: \emph{if the
historiography improves and the significances remain unnamed, that
counts against the framework and is not to be absorbed by it.} The
passage of time is thereby made a test rather than an excuse, which is
what the position in Remark~\ref{gc:rem:measurement} would otherwise
lack.

One asymmetry should be noted, since it limits what confirmation is
available. Failure to exhibit a $C$-free route is weak evidence at best,
because such routes are counterfactual claims subject to the same
immaturity that motivated the remark. The test is therefore one-sided in
the ordinary way: it can refute a particular significance claim and
cannot establish one, and the framework's standing rests on
whether the identifications accumulate.
\end{remark}

\paragraph{The consequence: aggregate, but of serial attributes.} The
reductio requires $R\gg1$, not $R>10^{3}$, and $R_1$ is a \emph{sum}.
Even under the corrections of Remark~\ref{gc:rem:serialcost}, a modest
stockpile suffices: twenty attributes contributing about ten apiece
give $R_1\approx200$, and a single well-attested cascade may supply as
much on its own.

\begin{center}\small
\begin{tabular}{@{}rrr@{}}
\toprule
attributes & each $\approx5$ & each $\approx10$\\
\midrule
$10$ & $50$ & $100$\\
$20$ & $100$ & $200$\\
$40$ & $200$ & $400$\\
\bottomrule
\end{tabular}
\end{center}

\noindent So the program is an enumeration of serial contingencies,
each articulated as a cascade and each audited for substitutability. The
record supplies candidates readily --- the military sequences above; the
demographic rupture of the Black Death and its effect on the very wage
structure Allen's account requires; and the theses of \cite{Diamond97}
beyond megafauna, several of which admit serial rather than parallel
statement. What the two calibrations established is not that recent
history is barren but that it must be interrogated along the trajectory
rather than across the map.

\begin{remark}[Why the negative calibrations are retained]
\label{gc:rem:whyretain}
Two attributes were examined and both returned nothing. It would be
tidier to omit them, and it would be a mistake.

They demonstrate that the instrument has teeth. A method that never
returns nothing cannot discriminate, and a reader entitled to suspect
that the estimator was constructed to deliver a congenial answer is
answered by two cases where it did not. They also fix the grain
discipline \emph{before} any positive estimate is offered, so that the
eventual stockpile is credible rather than convenient. And they preempt
the two objections a reader arrives with --- that coal was surely the
scarce thing, and that Diamond's asymmetry surely settles the matter ---
with the estimator's verdict rather than with assertion.

Most importantly, the pattern in the two failures is itself the
finding. Had one failed and the other succeeded, the first would read as
a red herring. Both failing, and failing for the same reason, is what
establishes Proposition~\ref{gc:prop:squeeze} and thereby the
aggregation strategy. The negative results carry the argument.
\end{remark}

This is where the argument's remaining empirical debt sits, and it is
now specific and bounded. What \eqref{gc:eq:bound} establishes is the
structure: given the $\alpha_j$, the bound follows. What
Proposition~\ref{gc:prop:squeeze} establishes is that no single recent
attribute will supply it, so the debt is a stockpile of order twenty
contingencies, each estimated by the sister-case method and each passing
the independence check. The weight claim is accordingly no longer
asserted but \emph{conditionally computed} --- reduced from a
plausibility judgment about an unexamined sum to an enumeration of
bounded terms, each obtained by an instrument calibrated in advance and
each with its required magnitude stated \textbf{[O]}. What remains owed
is the enumeration, and \emph{that} is a research task rather than an
argumentative gap: the terms are individually modest, the sources are
identified, and the standard each must meet is fixed before any is
computed.
\end{remark}

\begin{remark}[The conditioning is not teleological]
\label{gc:rem:teleology}
A reader's first objection will be that conditioning on a future event
to account for present observations is teleology: the Golden Point
reaching backward to arrange its own prerequisites. The framework
answers this at its foundations rather than here, and the answer should
be cited rather than improvised.

The relevant result is the \emph{indiscriminate time principle}, a
corollary of the extrusion principle \cite{Aethus}, whose content is a
pair of mutual non-informations: knowing where the events an attribute's
blanks refer to sit in space and time tells one nothing about when those
blanks resolve for oneself, and knowing when they resolve for oneself
tells one nothing about where they sit. Blankness is a property of the
indexing Aethus, never of a clock reading --- \emph{blanks lack tense},
and an observer's unlearned past is exactly as open as their unformed
future \cite{AethusBrief}.

The objection therefore rests on a premise the framework denies. It
assumes that the Golden Point's futurity makes it a different kind of
object from an unresolved attribute of the past --- that the not-yet-real
is a category. On the Aethic account there is no such category: the
Golden Point is a blank attribute relative to our Aethus, in precisely
the sense that an unread outcome of the deep past is blank, and
conditioning on it is the same operation in both cases. The asymmetry
one feels is real but \emph{entropic} rather than ontological: the
future blocks discriminating information more stringently because it
sits downstream of the entropic arrow, while the past merely happens to
withhold it --- and informational uncertainty is the same thing under
either obstruction \cite{Aethus}.

The mirror-image case matters, because it is a structure the
reader can check without granting anything particular to this work.
Newcomb's problem conditions on a \emph{past} event --- the predictor's
sealed placement --- which is blank relative to the agent's information
state; and any treatment on which the agent's decision bears on that
placement must decline the assumption at issue here, that a past
occurrence is thereby a settled fact. Past-Alignment conditions on a
\emph{future} event that is blank relative to ours. Same operation,
opposite direction, and by the indiscriminate time principle the
direction is not a difference. A reader who accepts the one and rejects
the other owes an account of which of the two mutual non-informations
they mean to keep.%
\footnote{The companion treatment is \cite{AethusDecision}, which develops
the Newcomb case on exactly this footing: determinacy indexed to an
information state rather than to a position in time, so that an unread past
outcome is a blank entry and blanks lack tense.}

What the Golden Point does here is therefore what any conditioning event
does: it selects, and does not act. The Aetherian Nexus is not
\emph{caused} to contain rocketry by a later passage; the weighting
conditioned on that passage assigns overwhelming weight to histories
containing it, and we read the record we are in. Whether anything
stronger than selection holds --- whether the future cone is constrained
rather than merely uninformative --- is exactly what
\S\ref{gc:sub:extrapolation} declines to settle, and the Void hypothesis
is the position on which nothing stronger holds at all.
\end{remark}

\begin{remark}[Why this is not circular, and where it could fail]
\label{gc:rem:circularity}
The obvious objection is circularity: conditioning on the Golden Point
and then finding its prerequisites is no achievement, since the
prerequisites are what the conditioning selects for. The objection would
land against a bare appeal to Past-Alignment, and it is worth being
explicit that it does not land here, because the argument of
\ref{gc:prin:conditioning} is \emph{comparative} rather than
confirmatory. Two candidate conditioning targets are placed side by
side; they issue \emph{different} predictions about the record --- one
indifferent to industrial prerequisites, one not --- and the record
discriminates between them. A comparison that could have gone the other
way is not a circle.

Three ways it could fail, in decreasing order of seriousness. The weight claim remains the decisive one, though
\S\ref{gc:subsub:weightbound} has reduced it from an assertion to a
single measurement: the bound \eqref{gc:eq:bound} is proved, and one
prerequisite at $\alpha\sim10^{-3}$ suffices, so what would remove the
contradiction is now the specific demonstration that \emph{every}
candidate prerequisite has $\alpha$ of order one. Second, the
argument requires ``observer'' to be robust under exactly these
perturbations; a sufficiently demanding notion of observerhood --- one
entailing industrial capacity, say --- would collapse the two
conditioning targets together and dissolve the reductio, though at the
cost of making ``observer'' do work it is not usually asked to do.
Third, the reductio establishes that the target is \emph{not}
observer-existence; that it is specifically the Golden Point is
supported by Past-Alignment rather than forced by the contradiction, and
some third target could in principle serve.
\end{remark}


% ---------------------------------------------------------------------
Both arguments are now in hand. The natural-history one carries the
weakness Remark~\ref{gc:rem:circularity} flags --- its robustness
premise is contestable --- and the human-history one of
\S\ref{gc:sub:humanhistory} does not; a reader who accepts only the
second has accepted the principle, and the first is then the same
inference extended into deeper time.

% ---------------------------------------------------------------------
\subsection{Extrapolating forward: three hypotheses}
\label{gc:sub:extrapolation}

Past-Alignment is a claim about the past cone and is silent ahead. Three
positions on the forward question are available, and this section states
them without adjudicating.

\begin{definition}[Extrapolation hypotheses; \textbf{[O]}]
\label{gc:def:extrapolation}
\begin{description}\itemsep3pt
\item[Void.] The alignment of our past cone with the Golden Nexus's is
a brute and uninfluential coincidence. Whether or not some connection
obtains, no law or inferential rule constrains our future cone to align
with the Golden Nexus's at the same spacetime position.

\item[Weak.] There is an explicit distribution over futures: the Nexus
describing our future cone assigns each proper child Aethus a weight
proportional to some function of its probability of the Golden Point ---
in the simplest case, that probability itself. Comparing two possible
futures is then possible by an Accordance-like metric, their relative
likelihood tracking their Golden-Point probabilities. Futures in which
the Golden Point does not occur are either incomparable under the
method or, in the strong reading, of probability zero relative to one's
Aethus.

\item[Strong.] The Golden Aethus is already stated in one's own Aethus,
and has been since the beginning of centric unfolding. On this reading
the stated attributes of Nexic reasoning were conditioning on the Golden
Point throughout; we occupy the Aethic block universe with a physical
moving present some centuries before that event, while the conditioning
that it occurs is already fully resolved and supplies exactly the
natural-historical content readable off our past cone. Contingency
remains in \emph{which} route is taken --- which stated attributes we
come to gather, conditioned on their consisting with the Golden Point
--- but outright Aethic disjointness from the Golden Point is
forbidden.
\end{description}
\end{definition}

\begin{remark}[Three questions the Strong hypothesis raises]
\label{gc:rem:strongquestions}
If the Strong hypothesis holds, three metaphysical questions become
pressing and none is answered here \textbf{[O]}: why future-cone Aethic
information should be available at all; how and why that information is
so intimately bound up with the initiation of centric unfolding, being
perhaps its earliest and leading piece of Aethic content; and what the
implications are for the entropic principle of Aethic reasoning and its
account of how the past is distributed between past and future Aethic
cones. Deriving either of the second or third from the other would be a
natural line of attack.
\end{remark}

\begin{remark}[One consideration bearing on the forward question, which
does not settle it]\label{gc:rem:asymptoteextrap}
The three hypotheses are stated and not adjudicated, and this remark
adjudicates nothing. It records a consideration that bears on the first
of them and raises its price.

By Remark~\ref{sch:rem:asymptote}, the climbing of a schedule is
selected for by the asymptotic condition and by no finite-horizon one.
If the record was produced under conditioning at all --- which
\S\ref{gc:sub:conditioning} argues and does not prove --- then the
conditioning event is an infinite-horizon one: a property of the entire
future cone rather than of any initial segment of it.

The Void hypothesis holds that no law or inferential rule constrains our
future cone. That sits awkwardly beside a conditioning event which
\emph{is} a property of that cone entire. The awkwardness is not a
refutation, and the distinction that saves Void should be stated for it:
a selection is not a law, and a reader may hold that the record is what
asymptotic conditioning produces while denying that anything about the
asymptote constrains what follows. What the consideration does is make
that a conjunction rather than a single claim --- Void must now carry
both halves, where as originally stated it carried one.
Remark~\ref{gc:rem:voidontology} sketches a form in which it might carry
them.

Nothing here supports Weak or Strong, which face their own difficulties
(Remark~\ref{gc:rem:strongquestions}); nothing here claims the
conditioning obtains; and \S\ref{gr:sub:ceiling} forbids anyone from
knowing that it does.
\end{remark}

\begin{remark}[How Void could look: a proposal, not an adjudication;
\textbf{[O]}]\label{gc:rem:voidontology}
Void is stated as a denial --- that no law or inferential rule
constrains our forward cone --- and a denial has no ontology. Since
Remark~\ref{gc:rem:asymptoteextrap} raises its price, it is only fair to
record that a positive form is available which pays that price rather
than evading it. What follows is a proposal about how Void \emph{could}
be realized, offered without endorsement and without prejudice to the
other two.

\emph{The construction.} Let the conditioning event be region-indexed in
its own statement: attributes whose referential events lie in the past
cone are constrained as one conditioning would have them, attributes
elsewhere as another. This is not two operations --- a conjunction of
region-indexed constraints is a single event in the product space, and
conditioning on it is one act. Read as a gate: \emph{if past-cone,
condition on $A$; otherwise on $X$}. A smooth interpolation between the
two regimes serves as well as a step and is the more natural object.

\emph{The weakening that makes it work.} $A$ need not \emph{be} the
infinite-survival condition. It need only approximate it closely enough
that Principle~\ref{gc:prin:pastalign} is observed --- and
Past-Alignment is an empirical claim about what the record looks like,
which more than one event can produce. So $A$ may be categorically
distinct from infinite survival while agreeing with it to within
observational reach, and Void's forward denial becomes the claim that
$X$ departs from $A$ where nothing has yet checked.

\emph{What that buys.} The natural objection is over-determination:
every region is eventually some observer's past cone, so a later
observer treats as past-conditioned what an earlier one treated as
future-conditioned, and consistency would force $X\equiv A$ everywhere
--- collapsing the construction into a single conditioning. Under the
approximation reading it does not bite. The two need only agree
\emph{approximately} on regions that are somebody's past, and
approximate agreement is compatible with categorical distinctness. The
objection bites only against the strict identity, which the construction
does not assert.

\emph{The constraint that remains.} Whatever indexes the gate cannot be
a distinguished spacetime position: Proposition~\ref{gc:prop:recursive}
lets the conditioning be anchored at any Aethus in the lineage and
returns that anchor's continuation, and
Remark~\ref{gc:rem:twoprivilege} forbids privilege outright. The index
must therefore be schematic --- computed by every observer from their
own anchor and yielding the same partition of their own cone. That is
not an obstacle so much as a specification: an anchor-invariant index
exists in this work already, since the goldening interpolation of
\S\ref{gr:subsub:practical} is age-relative and its windows are a fixed
fraction of a biosphere's own history.

\emph{And the coincidence need not be explained.} If the transition is
quasi-abrupt and falls near our position, Void owes no account of why.
That is not evasion within Void's own terms: its thesis is that such
alignments are brute and uninfluential, and a hypothesis whose content
is bruteness pays in brute facts by design. The fair scoring is that
this is a cost by ordinary lights --- three brute facts where the
hypothesis as stated carried one --- and not a cost by Void's, and a
reader will weigh those differently depending on what they think
explanation is for.

\emph{Where it risks collapsing into Weak.} The smooth version's danger
is that a graded conditioning strength over futures is a weighting of
futures, which is Weak's position and not Void's. Whether daylight
remains turns on whether the interpolation index is a function of
Golden-Point probability. This work keeps those apart ---
Remark~\ref{gr1:rem:whatTis} insists $\PP(\Lambda_\infty<\infty)$ is the
probability of advancing rather than of not dying --- so the gap is real
but narrow, and a proposal of this kind should be checked against it
before being relied on.

None of this adjudicates. It says only that Void has a form available in
which it is a positive hypothesis about what the conditioning event is,
rather than a refusal to say; and the form leaves open exactly what the
rest of this subsection leaves open, which is what the conditioning
actually is.
\end{remark}

\begin{remark}[Weak as the intermediate, and what each refusal commits
it to; \textbf{[O]}]\label{gc:rem:weakposition}
Symmetry requires the same treatment of the remaining hypothesis, and
Weak turns out to be the most structurally constrained of the three ---
its position is fixed by two refusals, each carrying a commitment it
cannot decline.

\emph{Refusing Void commits Weak to privileging the Golden Point.} If
some rule constrains the forward cone, Weak's is indexed to
Golden-Point probability and to nothing else --- so the Golden Point is
not one outcome among many but the quantity that sets the measure. This
commitment is shared with Strong, and it is what separates both from
Void: the forward question's first division is whether the Golden Point
indexes anything ahead.

\emph{Refusing Strong commits Weak to our contingent position doing
work.} Under Strong the conditioning is already fully resolved and every
anchor returns the same settled object, so where one stands is
inferentially idle. Under Weak the distribution is over \emph{one's own}
proper children, so the anchor does real work.

That second commitment needs care, because stated carelessly it collides
with Remark~\ref{gc:rem:twoprivilege}. It does not. What non-privilege
forbids is a metaphysically distinguished position at which the
conditioning is anchored and nowhere else. Weak's anchor is
\emph{schematic}: every Aethus indexes its own distribution over its own
children, by the same rule, and Proposition~\ref{gc:prop:recursive} says
precisely that the conditioning may be anchored anywhere and returns
that anchor's continuation. Weak is therefore not in tension with
recursive anchoring but is the hypothesis that takes it forward.

Which admits a sharper statement of Weak's position than the definition
gives. \emph{Weak is the only one of the three under which recursive
anchoring has forward content.} Under Void there is nothing ahead to
anchor; under Strong anchoring makes no difference, since all anchors
return the same resolved object; under Weak each anchor returns a
different distribution over its own children, and the differences are
the hypothesis's whole content.

\emph{One consequence for the Aetherian Nexus, which requires a
distinction the section has not yet drawn.}
Principle~\ref{gc:prin:conditioning} demotes observer-conditioning as an
\emph{explanation} of the record: observerhood is too robust under
perturbation to have produced it. That leaves untouched a second role.
Something must specify where ``we'' are, and
Principle~\ref{gc:prin:pastalign}'s own footnote already leans on the
Aetherian Nexus for exactly this, noting that our Aethus is a proper
child of it rather than the Nexus entire. \emph{Target} and \emph{index}
are different offices, and failing at the first says nothing about the
second. Under Weak the index does forward work, since the distribution
is over the children of a position the index picks out --- so the
Aetherian Nexus recovers a standing under Weak that it has under neither
of the others, without any part of the reductio being withdrawn
\cite{Nexus}.

\emph{And the costs, since the other two have had theirs stated.} Weak
owes an argument for its functional form: the definition says the weight
is proportional to \emph{some function} of Golden-Point probability, and
any increasing function would serve, with nothing yet selecting one ---
so the hypothesis as stated is a family rather than a position. It owes
an account of how the index is computed, which the schematic reading
requires be available to every anchor alike. And its weighting is
\emph{relative} --- a share of a total the hypothesis does not bound ---
so it inherits the caution of Remark~\ref{sch:rem:absolute} without
amendment: that a future carries the largest share of the forward
measure is not that it is likely, and comparison by an Accordance-like
metric is comparative in exactly the sense that supplies no absolute
reassurance. None of the three is supplied here.
\end{remark}

\subsection{Relation to the rest of this work}\label{gc:sub:relation}

Two compatibility questions arise immediately, and one structural
parallel is worth flagging as an open lead.

\emph{Compatibility with the vantage restriction.}
\S\ref{gr:sec} insists that Golden reasoning is a vantage and not an
obligation --- inferential, licensing conditional expectations rather
than duties, partly on the ground that $Q$ is a measure nothing
occupies. The Void and Weak hypotheses sit comfortably with that
restriction. The Strong hypothesis does not: if the Golden Aethus is
already stated and disjointness is forbidden, then the future-directed
reasoning has a purchase the vantage restriction disclaims. A reader
adopting Strong should understand that it buys something
$\Phi2$ deliberately gave away, and should re-price the argument
accordingly.

\emph{Compatibility with the epistemic ceiling.}
\S\ref{gr:sub:ceiling} holds that no finite observation can certify a
biosphere as Golden, since $Q$ and $\PP$ are locally equivalent and
separate only at infinite horizon. Past-Alignment appears to sit close
to this, since it asserts an agreement between observable content and
the Golden Nexus. The two are reconcilable, and the reconciliation is
that they concern different cones: the ceiling forbids distinguishing
$Q$ from $\PP$ over \emph{future} trajectories, while Past-Alignment
compares \emph{past} attribute content. Nothing in Past-Alignment
certifies that we shall reach the Golden Point; it claims that our
record is what conditioning on that event would have produced, which is
a claim about behind us.

\emph{An open structural parallel.} The trichotomy of
Definition~\ref{gc:def:extrapolation} is suggestively parallel to the
regime trichotomy of \S\ref{gr:subsub:regimes} --- trapped, critical,
viable. In particular the Strong hypothesis resembles the viable
regime, where $\PP(T=\infty)>0$ and $Q\ll\PP$ makes the Golden Biosphere
an ordinary sub-population one may belong to, while Void resembles the
trapped case in which no passage occurs. Should the parallel be exact,
the extrapolation question would not be independent of the anatomy
question: which hypothesis obtains would be settled by $\rho$ and by the
driver's scale function rather than by metaphysics, and
Theorem~\ref{gr1:thm:twocond} would decide it. We flag this as a lead
and do not claim it \textbf{[O]}.

\begin{remark}[The reversal of order]\label{gc:rem:reversal}
If the argument of this section is accepted, the framework's
architecture inverts. Optimation is no longer the starting weighting but
the first derived result --- what Golden-conditioning looks like cropped
to the observable cone. The Golden Biosphere is no longer characterized
by survival asymptotics and then found interesting; it is the
conditioning target, and the survival asymptotics of
\S\ref{gr1:sec} are the account of what such a target consists in. The
sections that follow are then not a sequence of separate results but the
development of a single object from different sides: \S\ref{gr1:sec}
gives its dynamics, \S\ref{gr:sec} its ensemble structure and what may
be inferred from occupying it, and \S\ref{ad:sec} the institutional
machinery that would determine whether a given biosphere reaches it.
\end{remark}

% ---------------------------------------------------------------------
\subsection{Ledger}\label{gc:sub:ledger}

\paragraph{Costs.} \emph{Past-Alignment} (\ref{gc:prin:pastalign}) ---
\textbf{[$\Phi$]}, the section's principal commitment;
\emph{falsifier:} an attribute of the record whose presence
Golden-conditioning does not predict, or a Golden-relevant trait whose
probability is undiminished under perturbation
(Remark~\ref{gc:rem:grade}). \emph{The weight claim} within
\ref{gc:prin:conditioning} --- \textbf{[S]}, asserted rather than
computed, and the reductio's load-bearing step;
\emph{falsifier:} a demonstration that the counterfactual sum does not
outweigh the Aetherian Nexus. \emph{Robustness of observerhood} ---
\textbf{[$\Phi$]}; \emph{falsifier:} a defensible notion of observer
entailing the industrial prerequisites, which would collapse the two
conditioning targets. \emph{Rung-screening}, inherited with the
substitution property \cite{Ben26Optimal}, and carrying its license
unchanged.

\paragraph{Buys.} Optimation derived rather than posited, and derived
from a comparison that could have failed. A conditioning target that
explains what observer-conditioning cannot --- why the record contains
its industrial and technological prerequisites at all. An explicit
statement of what is \emph{not} claimed about the future, together with
the three hypotheses that exhaust the forward question. And a
reorganization of the surrounding work into the development of one
object, stated in Remark~\ref{gc:rem:reversal}.

\paragraph{Inherited dependency.} Both arguments now depend on
technoanatomy being a \emph{process} rather than a parameter, for the
reason given in Remark~\ref{hr:rem:forcesprocess}: under a fixed
adverse $\rho$ the cascade pronounces the question closed before it is
asked, and the section's claims are then vacuously true and useless.
The dynamics of $\rho_t$ are open \textbf{[O]}
(\S\ref{gr1:subsub:fluct}), so this section's arguments are conditional
on a piece of the framework that is not yet built. That is a real
exposure and is not reduced by the arguments' internal soundness.

\paragraph{Not claimed.} That the future cone aligns with the Golden
Nexus; that the Golden Point will occur; that any of the three
extrapolation hypotheses is correct; and that the parallel with the
regime trichotomy of \S\ref{gr:subsub:regimes} is more than suggestive.

\section{Task-objects: individuation for a perturbed history}
\label{to:sec}

\begin{quote}\small
Every argument of \S\ref{gc:sec} quantifies over perturbed natural
histories. The dinosaurs never evolve; Africa develops otherwise; a
Golden Trek contains no Earth. Each such argument asks whether some
attribute survives the perturbation, and each therefore presupposes an
answer to a question it never asked: what makes a perturbed Earth still
an Earth, or a perturbed Aristotle still an Aristotle. Under an ontology
of possible worlds the question has a standard answer. Under an ontology
of Aethae it does not, and the perturbation arguments are ill-posed
without a replacement.

This section supplies the replacement, which our earliest dated
work already contains: the \emph{task-object}, an object individuated
not by perceived independence from other objects but by the contribution
it makes to a designated root. The concept was arrived at
independently and coincides with Kripke's distinction between rigid and
contingent designators \cite{kripke1980naming}; the departure from Kripke is
not in the distinction but in what is done with it, and in a denial of
the metaphysics his rigid designators presuppose
(\S\ref{to:sub:kripke}). Axiomatically the section develops
Remark~\ref{gr1:rem:goldendef}: the Golden-categorical premise
supplies the root, and this section builds the objects the premise
individuates --- criterion there, construction here, one doctrine ---
and, under the Route World's Strong reading only, the section closes
on the framework's ontological climax, the Golden Aethus
(\S\ref{to:sub:goldenaethus}).
Two liabilities are carried openly: the
correspondence between objects and task-objects is many-to-many in both
directions, so a selection principle is owed and only sketched
(\S\ref{to:sub:selection}); and the reply to the standard identity
objection runs through a commitment about consciousness whose cost is
priced in \S\ref{to:sub:identity}.
\end{quote}

% ---------------------------------------------------------------------
\subsection{The problem: perturbation requires individuation}
\label{to:sub:problem}

Consider the argument of \S\ref{gc:subsub:pastalign-argument} in the
form it was actually run. One perturbs the record --- the dinosaurs
persist --- and observes that the capacity to build rocket vehicles does
not survive. For this to be an argument rather than a stipulation, the
perturbed history must remain comparable to ours: it must still contain
a planet, a biosphere, a lineage, against which the missing capacity can
be registered as missing.

That comparability is what ordinary object-talk cannot supply here.
Perturb the record enough that no informed person would call the
resulting planet Earth, and the question ``does Earth still have
rocketry?'' does not become false; it becomes unasked. Yet the argument
needs it asked, because the Imperfection Magnitude Drop it is trying to
exhibit is precisely a drop \emph{in the substituted planet's}
contribution to the Golden Point.

The difficulty is sharpened by the framework's own commitments. Under an
Aethic ontology there are no completed determinate histories across
which an object could retain a categorical substance, so nothing plays
the role that a rigid designator's referent plays in the possible-worlds
picture (\S\ref{to:sub:kripke}). The individuation must therefore come
from somewhere else, and the proposal is that it comes from the root the
analysis is already conditioned on --- which is
Remark~\ref{gr1:rem:goldendef} arriving as this section's solution:
the definitional root of well-formed feature-terms and the
conditioning target of the perturbation arguments are one object, and
the coincidence is the design.

% ---------------------------------------------------------------------
\subsection{The definition}\label{to:sub:definition}

\begin{definition}[Task-object; \textbf{[$\Phi$]}]\label{to:def:taskobject}
Fix a root Aethus $\mathcal R$. A \emph{task-object} relative to
$\mathcal R$ is an object individuated not by its perceived independence
from other objects but by the contribution it makes to
$\mathcal R$. Where the root is the Golden Point, a \emph{Task-Earth} is
--- for instance --- the planet on which the makers of the Golden Point
speciated, and the corresponding \emph{task-humans} and
\emph{task-events} are individuated likewise.\footnote{%
\emph{Provenance.} The concept and the Task-Earth example are
ours, recorded 1 July 2022 and reproduced in the natural-historical
development of August 2022 \cite{GoldenScience22}, which makes this the
earliest dated item in the present corpus --- predating the extrusion
principle by some weeks and the Golden Point material it was written to
serve by a month. The formalization below, the placement against Kripke,
and the pricing of the identity objection are later.}
\end{definition}

% [Instates the task-object/Golden-categorical integration: the premise
% as the concept's axiomatic and conceptual root, per the author's
% direction of Aug 22.]
\begin{remark}[The axiomatic root, and what contribution is]
\label{to:rem:goldenroot}
Two claims fix this definition's place in the framework, and both run
in the same direction: from Remark~\ref{gr1:rem:goldendef} to here.

\emph{Task-objects arise from the Golden-categorical premise.} The
root parameter $\mathcal R$ is, at its origin and at its canonical
value, the Golden Point, and the derivation order is premise first:
Golden-categorical robustness is the criterion of well-defined
reference, and the task-object is what the criterion \emph{builds}
--- the object-form of Golden-categorical individuation. The
parametric definition above is the abstraction from that case, not a
neutral genus of which the Golden root is one species, and the
distinction matters below.

\emph{At the canonical root, contribution is a defined quantity.}
``The contribution it makes to $\mathcal R$'' is, for an arbitrary
root, a phrase owed a semantics. At the Golden root the framework
already possesses one, assembled from its own operations: complement
the individuating attributes and read the displacement --- the
Imperfection Magnitude Drop of \S\ref{gc:sub:vocab}, graded by the
$\varepsilon$-Departure apparatus (Definition~\ref{gc:def:epsdep})
and priced by the investment bound
(Proposition~\ref{ra:prop:investment}) --- so that an object's
contribution to the Golden Point is the $G$-weight its complementation
forfeits. Nothing analogous exists for an arbitrary root unless its
inquiry supplies it. So the concept's well-posedness is native at the
Golden root and \emph{inherited} elsewhere, exactly where an analog
referent for contribution is supplied --- which is the precise sense
in which the premise is the concept's conceptual root and not merely
its first customer.

\emph{The axiomatic order matches the genetic order.} The provenance
footnote above records that Task-Earth was built in July 2022 in
service of the Golden Point material it predates on paper; the premise
converts the service into the source. That the concept's stated
axiomatics now runs in the direction its construction actually ran is
recorded deliberately: it is the difference between a root assigned in
retrospect and a root that was generative, and the corpus's dated
record is what makes the claim checkable rather than honorific.

None of this retracts the root-relativity stated next; it locates it.
The \emph{machine} is parametric in $\mathcal R$; the
\emph{semantics} is generative at the Golden Point.
\end{remark}

\begin{remark}[Root-relativity is by design, not by accident]
\label{to:rem:rootrelative}
The root is a parameter. Task-object individuation is root-relative by
construction: the machinery requires no particular root, and a
different inquiry with a different root would individuate differently
and legitimately --- provided it supplies what
Remark~\ref{to:rem:goldenroot} shows the Golden root supplies
natively, a formal referent for contribution. This matters for how the
concept should be read outside this paper: nothing in the
\emph{machinery} is committed to prosperity, survival, or any other
particular optimand. What is committed is a pair: that some root must
be designated before individuation is well-defined, and that at the
designated root contribution must mean something. The Golden Point is
where this program obtains both at once, and where the concept was
generated; other roots rent what it owns.
\end{remark}

\begin{remark}[Why this is not a redescription of ordinary objects]
\label{to:rem:notredescription}
Ordinary objects and task-objects come apart in both directions, and the
examples are ours. Earth $3.9$ billion years ago is called
Earth, despite differing from the present planet in continent
configuration, atmospheric composition, and much else. A perfect
duplicate in another star system would not be called Earth, despite
differing in none of them. Neither verdict tracks contribution to a
root, and both are stable features of ordinary usage --- which is
precisely why ordinary usage cannot serve where the analysis needs
comparability across perturbation.
\end{remark}

% ---------------------------------------------------------------------
\subsection{Relation to Kripke, and the departure}\label{to:sub:kripke}

The distinction between a designator that picks out its object robustly
under counterfactual variation and one that picks it out by a
description which the object might not have satisfied is Kripke's
\cite{kripke1980naming}. The present concept was arrived at independently, by a
different route --- asking what continuity there is between Earth's
persistence in the past and the prospects of a biosphere's near future
--- and the coincidence is worth recording as a convergence rather than
claimed as a discovery.

The departure is twofold.

\emph{First, in what is done with the distinction.} Kripke's conclusion
is prohibitive: do not use contingent designators where rigid ones are
required. The present move is instead to adopt the awkward category as
both ontologically apt and operationally useful --- the pattern by which
physics adopted Newton's first law over the intuition of motion
requiring a mover. That pattern is not general license, and should not
be read as one: counterintuitiveness is usually a warning, conspicuously
so in moral matters, and the claim is only that this is one of the
uncommon cases in which the received intuition is out of step with the
ontology, plausibly because what served human societies through
evolution was never selected for tracking it.

\emph{Second, and more radically, in a denial of the metaphysics.} The
Aethic framework does not accept the objects to which rigid designators
are supposed to point. With Aethae rather than possible worlds as the
base unit, there is no categorical substance persisting across worlds
for such a designator to track: by the Aethic reversal principle, the
underlying determinate reality required to constitute such an object
cannot be well-defined in the first place \cite{Aethus}. The scope of
this denial should be stated carefully --- \emph{what is denied is
Kripke's metaphysical presupposition, not his linguistics.} The modal
argument survives untouched as a claim about how names function in
language; the present claim lives a level below, about what that
functioning tracks.

\begin{remark}[Robustness recovered without substance]
\label{to:rem:robustness}
What made rigid designation attractive was robustness under
perturbation, and that is recoverable without the metaphysics. Placing
several description classes of an object in Aethic superposition yields
a designation that survives perturbation by landing in another class
when one fails --- robustness as a property of the superposition rather
than of a substance beneath it. The strong form of the resulting
position is that an Aethic ontology must make task-objects fundamental
and admit rigid-designator talk only as non-arbitrary shorthand for
them.
\end{remark}

% ---------------------------------------------------------------------
\subsection{The two-way fan, and the selection problem}
\label{to:sub:selection}

\begin{proposition}[Many-to-many in both directions; \textbf{[S]}]
\label{to:prop:fan}
The correspondence between objects and task-objects is many-to-many.
One object admits indefinitely many task-objects: Earth may be
individuated as a spherical body composed largely of iron, or as a
region within a star's habitable zone, and so on without limit. One
task-object admits indefinitely many objects: fix Task-Earth as the
speciation site of the Golden Point's makers, and every planet
satisfying it qualifies --- including one on which Hawaii never formed,
which is not our Earth and is a Task-Earth.
\end{proposition}

The fan is the concept's principal technical liability, since it means a
task-object is not determined by the object it individuates, and some
selection is required. Our own assessment is that the
selection carries irreducible speculation, but speculation that might be
linguistically systematized; the criterion offered is to prefer the
task-object yielding the simplest or most informative interplay with the
wider question the root encodes.

\begin{remark}[The criterion is a gauge choice, and the paper already
has one]\label{to:rem:gauge}
That criterion is recognizably a gauge fixing: one probes the freedom
until the description aligns most economically, exactly as one selects
coordinates in which a physical problem is cleanest. The comparison is
ours and is substantive here, because
\S\ref{gr1:sub:coords} performs the same operation on the survival
model --- initial data there form a gauge orbit invisible to every
asymptotic exponent, and the analysis proceeds by identifying what the
freedom cannot touch. The two are the same methodological move at
different levels: find the quantities the arbitrary choice leaves
invariant, and state the theory in those.

This suggests where a systematization would come from, and the framework
supplies a candidate it does not yet apply here. The construction for
degree-of-existence over description classes \cite{Aethus} selects
canonical classes as \emph{fixed points} --- those whose further
refinements no longer improve description quality --- which is exactly
the shape of criterion the task-object selection problem needs. Whether
that construction transfers is open \textbf{[O]}, and it is the first
thing we would attempt.

At the canonical root a second instrument becomes available, and it is
this work's own. Among the task-objects an object admits, the
conditioning is not indifferent: it concentrates on the individuations
that are load-bearing --- those whose defining contribution the record
shows \emph{invested}, in the sense
Proposition~\ref{ra:prop:investment} prices --- so ``prefer the
task-object whose contribution carries the weight'' is a selection
criterion the Golden root supplies for free, where an arbitrary root
supplies none. It is offered at \textbf{[S]} beside the fixed-point
candidate, not in place of it; the two would be expected to agree
where both apply, and testing that agreement is part of the same open
item.
\end{remark}

% ---------------------------------------------------------------------
\subsection{The identity objection, and what answering it costs}
\label{to:sub:identity}

The standard objection is Kripke's own. Aristotle might not have taught
Alexander; he would have been Aristotle regardless; so the person is
what the name tracks, and no task-human individuation is required.
Under the present proposal the corresponding task-human --- the man who
taught Alexander in the sense of equipping him for some specified
subsequent achievement --- would simply not have been instantiated, and
the objection says this gets the identity facts wrong.

The reply is not to deny the objection's premise but to run it out.

\begin{remark}[The reply, as a reductio of the objection's own criterion]
\label{to:rem:reductio}
Grant that identity is settled by continuity of consciousness through a
common point of divergence, which is what the objection's force rests
on. Under centric unfolding, any other person's self-describing Aethus
is reachable by ascending to a sufficiently early ancestor Aethus ---
plausibly in infancy or before, wherever a substantial portion of the
universe enters disagreeing superposition --- and perturbing there. The
objection's own criterion therefore identifies not merely the untutored
Aristotle with Aristotle, but \emph{every} person with every other: the
set of human-pointing rigid designators collapses to a single
equivalence class.

The structure of the reply should be read carefully, since its force
depends on it. It is not a positive thesis that all persons are
identical, offered for acceptance. It is a demonstration that the
criterion the objection relies on, applied consistently under centric
unfolding, cannot deliver the fine-grained identity facts the objection
needs --- so the objection cannot be used against task-object
individuation without first supplying a criterion that survives its own
generalization.

Two costs are recorded. The reply commits to reading continuity of
consciousness through the framework's own treatment of centric unfolding
\cite{Aethus}, which is a substantive commitment and not a neutral one.
And a reader may take the collapse as a reductio of the framework rather
than of the objection; the reply is that the collapse
follows from the objection's criterion together with centric unfolding
as stated, so contesting it while retaining Aethic reasoning would
require an extra-centric unfolding proposal, which nothing here
supplies \textbf{[O]}.
\end{remark}

% ---------------------------------------------------------------------
% [Instates the author's ontological-completion dictation (Aug 22): the
% Golden-categorical root extended from semantics to constitution under
% the Route World's Strong reading; the Golden Aethus named; the toggle
% argument; the fluctuation thesis; the Boltzmann displacement; fences.]
\subsection{The ontological completion: the Golden Aethus}
\label{to:sub:goldenaethus}

Everything in this section so far has run at the level of reference:
a criterion (Remark~\ref{gr1:rem:goldendef}), its constructions
(Definition~\ref{to:def:taskobject}), and a formal referent for
contribution (Remark~\ref{to:rem:goldenroot}). This subsection states
what the arc closes on when the Route World's Strong reading is added,
and it should be said plainly at the outset both that this is the
framework's categorical and ontological climax and that it is
\emph{conditional}: everything here consumes the Strong reading, and
nothing earlier in the section does. A reader who declines the Route
World keeps the semantics whole and loses only what follows.

\begin{definition}[Golden Aethus]\label{to:def:goldenaethus}
The \emph{Golden Aethus} is the Aethus that states exactly the Golden
attribute --- a Golden biosphere in the future light cone, the Golden
Point event as a future-referring stated attribute --- with every
other attribute genuinely blank. It is the \emph{index} of the Golden
Nexus: conditioning the Cosmic Nexus on the Golden Aethus's stated
content is what the Golden Nexus \emph{is} (\S\ref{gc:sub:vocab}).
It is distinct from the Aetherian Nexus, which is the Aethus of our
whole natural history: under the Strong reading the Aetherian Nexus is
a weight-carrying descendant of the Golden Aethus --- the record is
stated content accreted within the Golden Aethus's blanks.
\end{definition}

\begin{remark}[The toggle: void and world]\label{to:rem:toggle}
Two measures stand one attribute apart, and the distance between them
is the argument. The bare Cosmic Nexus absorbs no conditioning
(\S\ref{gc:sub:vocab}), and its typical configurations are the
high-entropy ones: unconditioned occurrence-weight is
entropy-dominated, ordered macrostructure is exponentially atypical
within it, and the companion's own reading of the bare measure's
default tendencies is the same fact in Nexic vocabulary
\cite{Nexus} \textbf{[S]}. Call that face what it is --- the void.
Now state the one attribute the Golden Aethus carries, and the weight
redistributes onto exactly the histories that can \emph{carry} it: a
Golden biosphere is not statable over vacuum, so the conditioning
concentrates on deep ordered records, biospheres, lineages, the whole
furniture of a world. The existence-in-the-weighting of every ordered
thing therefore toggles with a single stated attribute --- void on one
side of the statement, world on the other. And the toggle has a
biosphere-scale mirror already on the books: it is $\PP$ against $Q$
(Definition~\ref{gr:def:gg}) run at cosmic register --- the goldening
of the universe rather than of a biosphere --- which is why nothing in
this subsection requires machinery the paper does not already own.
\end{remark}

\begin{remark}[Everything as a fluctuation on the Golden Aethus;
\textbf{[$\Phi$]}, Strong-conditional]\label{to:rem:fluctuation}
Under the Strong reading, the toggle licenses an ontological
statement, and it is the strongest sentence in this work: each thing,
object, and circumstance in the universe is a \emph{fluctuation on the
Golden Aethus} --- a weight-carrying excursion within its blanks,
superposition content that the Golden Nexus renders load-bearing.

The word ``fluctuation'' is doing exact work, and the inversion inside
it should be stated with care. Under the bare Cosmic Nexus, ordered
structure is a fluctuation in the dismissive sense: improbable,
incidental, owed nothing. Under the conditioning, the same structures
are fluctuations still --- they remain the variance-structure of a
measure, excursions rather than substances --- but they are no longer
accidents: the conditioning concentrates on them because the stated
attribute cannot be carried without them. What toggles is not
fluctuation-status but \emph{accident}-status. That is the precise
sense of ``where things arise from'': things arise in the weighting,
from the Golden Aethus, as its resolution-structure --- and the
framework has described this shape before. A biosphere-scale state is
already, on the books, a weighted superposition of the poles
(Remark~\ref{gc:rem:boltzmann}, Remark~\ref{gc:rem:pairsuffices});
the present thesis is that description run on the universe's contents
--- the same shape at cosmic scale, a consilience rather than an
addition.

Two consequences discharge obligations incurred earlier. First, the
categorical claim: if things themselves arise from the Golden Aethus,
then categories drawn non-arbitrarily should trace to the Golden Point
--- not because a semantics was chosen well but because the
individuation being read off is the world's own. The
Golden-categorical premise and the task-object machinery are, under
Strong, not conventions but transcriptions, and the ``intended all
along'' of Remark~\ref{to:rem:goldenroot} acquires its full meaning:
the criterion was tracking constitution. Second, the contribution
referent becomes a constitution-measure: the $G$-weight a
complementation forfeits (Remark~\ref{to:rem:goldenroot}) is, under
Strong, how much of the thing there \emph{is} --- the analogy to read
it by, flagged as structural and consuming no physics, being
excitations over a ground state: objects stand to the Golden Aethus as
excitations to a vacuum, the geometry being that of a conditioned
measure's typical set and nothing more.
\end{remark}

\begin{remark}[The vantage itself: non-indeterminacy roots to the same
point]\label{to:rem:vantage}
The fluctuation thesis concerns what things \emph{are}. A further
extension, the author's, concerns the standpoint from which anything is
determinately anything at all --- and it is worth stating separately
because it is what carries the ontological completion into registers as
abstract as morality (\S\ref{gs:sub:nodes}). The companion's first
postulate already makes determinate reasoning \emph{indexical}: without
a central Aethus there is no validity, no pruning, no context --- the
procedural statement that all logic runs relative to an index
\cite{Aethus}. The question the postulate leaves open is which index a
\emph{universe's} worth of determinate content has. Under the Strong
reading the answer is the one this subsection has already built: the
contentful universe-scale index is the Golden Aethus, of which every
determinate vantage in the record --- the Aetherian Nexus, and our own
Aethus within it --- is a weight-carrying descendant. So the toggle of
Remark~\ref{to:rem:toggle} reads epistemically as well as
ontologically: remove the conditioning and what remains is not merely
the void of things but the void of \emph{standpoint} --- indeterminacy
with no structure in any particular direction, in which anything goes
because nothing is stated anywhere. What separates vantage and context
from that abyss is the same object that separates world from void. (The
view-from-nowhere tradition meets the framework exactly here, and is
engaged at Remark~\ref{gs:rem:meaningallies}: a literal nowhere is
the abyss, so the honest question was never whether to view from an
index but which index the viewing consumes \cite{Nagel71}.)
Conceptual form, not only thing-ness, roots to the Golden Point ---
graded by the same ladder, priced by the same fences, and consumed
nowhere earlier than this remark.
\end{remark}

\begin{remark}[A payoff: the fluctuation-observer pathology displaced]
\label{to:rem:boltzbrains}
One consequence is substantial enough to state separately, at
\textbf{[S]}. Under a bare high-entropy measure conditioned only on
the existence of an observer-moment, the typical observer is a
\emph{minimal} fluctuation --- the Boltzmann-brain pathology of the
cosmological literature \cite{AS04,Car17}: records deceptive,
environments absent, induction self-undermining. Under conditioning on
the Golden attribute the pathology is displaced by the toggle itself,
with no added postulate: a minimal deluded fluctuation contributes
nothing to carrying a Golden biosphere, so the conditioned weight
concentrates on observers embedded in the deep, real, ordered
structure the attribute requires --- observers whose records are
records. Two scope notes keep this honest. Observer-conditioning alone
(Optimation in the bare sense) does not obviously achieve the
displacement --- it is the Golden attribute's structural demands that
do the work, which is an argument \emph{for} the Golden target from an
unexpected quarter. And nothing here claims a solution to the
cosmological measure problem: the claim is internal to the framework's
weighting, and is priced as such. One further adjacency is recorded at
\textbf{[O]} and not claimed: whether the record's thermodynamic
asymmetry is itself induced by the future-referring conditioning --- a
future-boundary condition that tense-freedom renders legitimate --- is
open.
\end{remark}

\emph{The price, paid in full.} Four fences. \emph{Conditionality,
as a ladder}: under Strong, the thesis is the framework's ontology;
under Weak, it holds of the record --- the past cone is
Golden-conditioned by Past-Alignment --- and is silent forward; under
Void it is unavailable, and the semantic premise of
Remark~\ref{gr1:rem:goldendef} stands alone on its comparative
defense, having consumed none of this. \emph{Not idealism}: the
Golden Aethus is nobody's mind; statedness is a property of an index
and not of anyone's access
(Remark~\ref{gc:rem:determinacygap}), and every verdict here is
realist in exactly that sense. \emph{Not teleology}: arising-from is
arising-in-the-weighting; the conditioning does not act, choose, or
design, and the standing fences of \S\ref{gc:sec} govern unweakened.
\emph{Scope}: the quantifier ``essentially all'' ranges over
real-world feature-terms per Remark~\ref{gr1:rem:goldendef};
abstracta are outside it by that remark's own scoping, and the
categorical claim carries its own check --- a non-arbitrary real-world
category that resists every tracing to the root would break the
quantifier, and exhibiting one is the invited falsification.

The arc of this work's categorical machinery --- criterion,
construction, referent, constitution --- terminates here. This is,
conditionally and at its stated price, where things arise from.

\subsection{What this licenses}\label{to:sub:licenses}

With individuation in hand, three things earlier in the paper become
well-formed rather than merely persuasive.

\emph{The perturbation arguments.} \S\ref{gc:subsub:pastalign-argument}
may compare our record with histories in which the dinosaurs persist,
because both contain a Task-Earth and a task-lineage against which a
missing capacity registers as missing. Without task-objects the
comparison has no common term and the argument does not run.

\emph{The substitution property.} The generalization across rungs
\cite{Nexus} places a trait of ours in the conditioning seat and
deduces that alternative histories fail its occurrence. That maneuver
requires the trait to be identifiable in histories where it does not
occur, which is a task-object identification and not an object one.

\emph{Imperfection Magnitude Drops.} A drop is a fall in contribution to
the root under substitution \cite{Ben26Optimal}. Since a task-object
just is an object individuated by that contribution, the drop is a
change in a task-object's value rather than a comparison between
incommensurable objects --- which is what makes it a quantity at all.

\begin{remark}[Order of development]\label{to:rem:order}
The dating is worth one line, because it inverts the order of
presentation. Task-objects are dated 1 July 2022; the Golden Point
material follows in August \cite{GoldenScience22}; the extrusion
principle and with it tense-free blankness follow weeks after that. The
individuation scheme therefore came first and the ontology that
vindicates it came second --- the concept was constructed to serve a
question about a biosphere's prospects, and only later found itself
entailed by an ontology built for unrelated reasons. We record this as
provenance rather than as evidence, but it does bear on the charge that
task-objects are a device fitted to the argument they support.
\end{remark}

% ---------------------------------------------------------------------
\subsection{Ledger}\label{to:sub:ledger}

\paragraph{Costs.} \emph{Task-object individuation}
(\ref{to:def:taskobject}) --- \textbf{[$\Phi$]};
\emph{falsifier:} a perturbation argument of the kind
\S\ref{gc:sec} requires, shown to be well-formed under ordinary object
individuation, which would make the apparatus idle. \emph{The denial of
rigid designators' referents} --- \textbf{[$\Phi$]}, inherited from the
Aethic ontology \cite{Aethus} and not argued here;
\emph{falsifier:} the ontology's. \emph{The selection principle} ---
\textbf{[O]}, owed and not supplied; the fan of
Proposition~\ref{to:prop:fan} is real and a criterion beyond
``simplest or most informative'' is not given. \emph{The identity reply}
--- \textbf{[$\Phi$]}, resting on continuity of consciousness read
through centric unfolding, with the singleton collapse acknowledged as
a cost a reader may decline to pay.

\paragraph{Buys.} Well-formedness for the perturbation arguments, the
substitution property, and Imperfection Magnitude Drops, none of which
survives ordinary object individuation intact. A principled account of
why ordinary usage fails here rather than an appeal to its being
inconvenient. And a concept whose earliest dated statement precedes both
the material it serves and the ontology that vindicates it. And, under
the Strong reading only, the constitution reading of
\S\ref{to:sub:goldenaethus}, with the fluctuation-observer
displacement it purchases.

\paragraph{Not claimed.} That Kripke's modal argument is mistaken; that
all persons are identical; that the Golden Point is a privileged root \emph{of the
machinery} --- its generative standing for the semantics
(Remark~\ref{to:rem:goldenroot}) and, under Strong only, its
constitutive standing for the ontology (\S\ref{to:sub:goldenaethus})
\emph{are} claimed, and priced where stated; or that the selection
problem is solved.


\part{The Survivorship Mathematics}

\noindent Part~II is the mathematics alone, and nothing normative
enters it. The two-channel hazard is built on the Kolmogorov pair,
the two-condition theorem separates the regimes, the schedule channel
is developed as the third hazard, and the survivorship limit is taken
to the Doob transform with its three modes (\S\ref{gr1:sec}).
Everything that the later parts consume is proved here at
\textbf{[P]} or carried at a stated grade, so a reader who wants only
the probability can stop when this part ends, and it is for that
reader that the standalone extraction exists.

\section{The model and its survivorship structure}\label{gr1:sec}

\begin{quote}\small
This section is mathematics. It establishes the object the philosophical
section reasons about, and it is organized by what that reasoning
consumes rather than by order of discovery: the object
(\S\ref{gr1:sub:object}), its coordinates (\S\ref{gr1:sub:coords}), the
verdict on survival (\S\ref{gr1:sub:verdict}), and the structure of
conditioning on survival (\S\ref{gr1:sub:cond}). Material that maps the
terrain around the model rather than establishing it --- horizon theory,
the free-energy background, the waiting room and gain profile, the
transfer computations, and the numerics --- is placed in appendices,
where it is available without standing in the way.

Claims are tagged \textbf{[P]} proved here, \textbf{[C]} classical or
cited, \textbf{[A]} conditional on the soft-wall Ansatz
(Ansatz~\ref{sh:ansatz}), \textbf{[S]} scaling-level, \textbf{[O]}
open. The tag \textbf{[C]} is load-bearing and frequent: the persistence
exponent, the harmonic function of the conditioned process, the
free-energy asymptotics of the background and the boundary
classification of one-dimensional diffusions are all imported, and this
section claims none of them. What it claims is the model, the
coordinates, and the assembly.
\end{quote}

\subsection{Relation to the literature}\label{gr1:sub:literature}

Four bodies of work border this section, and the reader deserves a
precise map of what is classical, what is adjacent, and what is claimed
as new.

\emph{The classical layer.} The state process is Kolmogorov's 1934
diffusion \cite{Kol34}, the primordial hypoelliptic example
\cite{Hor67}. Its persistence theory is a mature subject with a clean
lineage --- McKean \cite{McK63}, Goldman \cite{Gol71}, Lachal
\cite{Lac91}, Sinai \cite{Sin92}, Isozaki--Watanabe \cite{IW94},
Denisov--Wachtel \cite{DW15} --- surveyed in \cite{AS15} and, as the
\emph{random acceleration process}, possessing an extensive physics
literature reviewed in \cite{BMS13}. The conditioned (never-hitting)
process, our limit object in \S\ref{sh:sub:interp}, is constructed by
Groeneboom, Jongbloed and Wellner \cite{GJW99}. None of this is ours;
we import it.

\emph{The adjacent frameworks.} Our survivorship interpolation belongs
to two established genres. First, Theorem~\ref{sh:thm:Q} is a
\emph{penalization} statement in the sense of Roynette--Vallois--Yor
\cite{RVY06}: a limit of Wiener-type measures reweighted by a
time-dependent functional, identified as a martingale change of measure
--- here the weight is survival itself, and the polynomial tail places
it in the clockless class. Second, the mode classification of
Theorem~\ref{sh:thm:modes} sits against the theory of quasi-stationary
distributions and $Q$-processes surveyed in \cite{MV12}: classical QSD
theory is at home in the exponential-tail class (our modes (I) and
(III)), while the critical polynomial mode admits no proper QSD in fixed
coordinates, its Yaglom object being self-similar. We use both
frameworks as positioning, not as inputs; the proofs here are
self-contained given the cited persistence theory.

\emph{The applied layer: what this is a model of.} Stochastic hazard
rates are not new, and the wrapper $S_t=\exp(-\int_0^th)$ with $h$ a
diffusion is entirely standard: it is the reduced-form, or
intensity-based, framework of credit risk \cite{Lan04} and of stochastic
mortality, where the force of mortality is routinely modeled by a
diffusion, the mean-reverting Brownian Gompertz specification of Milevsky
and Promislow \cite{MP01} being a canonical instance. What distinguishes
the present model is not the wrapper but \emph{where the noise is
placed}: standard practice diffuses $h$ or $\log h$, whereas we diffuse
$-\tfrac{d}{dt}\log h$, one integration further out. The consequences are
structural rather than cosmetic --- $\log h$ becomes $C^1$, the state
becomes degenerate and $3/2$-self-similar, and the governing persistence
exponent moves from $\tfrac12$ to $\tfrac14$ (Remark~\ref{sh:rem:iwp},
\S\ref{sh:subsub:tail}). We stress the point because it locates the
novelty precisely, and disclaims the rest.

\emph{The conceptual ancestor.} The thesis of \S\ref{sh:sub:interp} ---
that selection on survival makes the observed ensemble diverge
systematically from the unbiased one --- is not new in substance. It is
the frailty phenomenon of Vaupel, Manton and Stallard \cite{VMS79} and
Vaupel and Yashin \cite{VY85}: when the frail are removed first the
population hazard falls below the individual hazard, the gap widening
with age, producing mortality deceleration at the population level with
no corresponding deceleration for any individual. We claim no priority
over that insight, and regard this section as exhibiting its
\emph{critical dynamical} analog. The differences are structural and
are the point. Frailty is a static random variable fixed at birth and the
analysis is a mixture computation over a population; here the randomness
is a driver evolving in time, and the object is a limit of path measures
singular with respect to the unbiased law
(Theorem~\ref{sh:thm:singular}). Where the frailty literature computes
how an observed hazard curve bends, the interpolation of
\S\ref{sh:sub:interp} identifies the limiting law of the survivor's
entire trajectory, together with the drift \eqref{sh:eq:G} that
conditioning manufactures.

\emph{What is claimed as new.} To our knowledge, the following are new
at least in assembly, and we flag them as this section's contribution:
the log-derivative placement of the noise, with its gauge analysis
(\S\ref{sh:subsub:dimensional}) and the two-sided soft-wall analysis
whose remaining gap is located exactly (Ansatz~\ref{sh:ansatz}); the
scale-function horizon trichotomy with mortal critical case
(\S\ref{sh:subsub:scale}); the adverse-drift survival rate
$3\nu^2/8\sigma^2$ with its front-load-and-coast optimal profile
(Proposition~\ref{sh:prop:advdrift}); and the survivorship
interpolation read as a monotone flow with the tail index as its
linear-response coefficient, classified across drivers into saturation,
transformation and concentration (\S\ref{sh:sub:interp}). Where a
statement merely re-expresses known results in this frame, the tags and
citations say so.

\begin{convention}\label{sh:conv}
For positive functions, $f\sim g$ means $f/g\to1$; $f\asymp g$ means
$0<\liminf f/g\le\limsup f/g<\infty$; $f\lesssim g$ means $f\le Cg$.
Hatted quantities are nondimensionalized by the intrinsic timescale of
Definition~\ref{sh:def:units}. $\Phi$ is the standard normal distribution
function; $\one_A$ is an indicator; $\Leb$ is Lebesgue measure. All processes
are defined on a filtered probability space satisfying the usual conditions,
and $W$ denotes a standard one-dimensional Brownian motion.
\end{convention}

% ---------------------------------------------------------------------
\subsection{The object}\label{gr1:sub:object}\label{tc:sub:model}\label{sh:sub:model}

% ---------------------------------------------------------------------

\subsubsection{Setup: hazard, credit, and the Kolmogorov diffusion}
\label{sh:subsub:setup}

\begin{definition}[Model]\label{sh:def:model}
Fix $\sigma>0$ and initial data $(h_0,\mu_0,D_0)\in(0,\infty)\times\R\times\R$.
Define the \emph{improvement rate}, \emph{credit}, \emph{hazard}, and
\emph{cumulative hazard} by
\begin{equation}\label{sh:eq:model}
\mu_t=\mu_0+\sigma W_t,\qquad
D_t=D_0+\int_0^t\mu_s\,ds,\qquad
h_t=h_0e^{-(D_t-D_0)},\qquad
\Lambda_t=\int_0^t h_u\,du,
\end{equation}
and write $x_t=\log h_t$, so that $\mu_t=-\dot x_t$. We also set
$I_t=\int_0^t W_s\,ds$, so that $D_t=D_0+\mu_0t+\sigma I_t$.
\end{definition}

\begin{remark}[Which quantity accumulates, and why it is not the
logarithm]\label{sh:rem:accumulator}
One choice in the construction below is easy to make silently and worth
making out loud, because this work returns to its consequence three
times. Along any single path the natural additive ledger is the
\emph{cumulative} hazard $\Lambda$, and survival is $e^{-\Lambda}$ on
that path. Across a superposition the two do not commute: $\Lambda$ is
per-path and does not average --- $\E[e^{-\Lambda}]\neq e^{-\E[\Lambda]}$
by Jensen, with the gap in the survivor's favour --- while the
\emph{weight} does average, linearly, because averaging weights is what
the framework's law of total probability is.

So the canonical accumulator is the weight and not its logarithm, and
$-\log$ survival is a per-path reparametrisation rather than a quantity
one may average over Aethae \cite{Aethus}. The Cox construction below
respects this by keeping $\Lambda$ inside the expectation, which is why
$\PP(T>t)=\E[e^{-\Lambda_t}]$ and not $e^{-\E[\Lambda_t]}$.

The same gap is what Remark~\ref{sch:rem:gap2} records for completion
weight, what Remark~\ref{sch:rem:bpre} records for the criticality
criterion, and what Remark~\ref{sch:rem:n2} records for the companion's
contrast. Four instances of one fact: a geometric mean is what a path
experiences and an arithmetic mean is what an ensemble reports, and
everything in this work that looks like optimism under averaging is that
gap read from the wrong side.
\end{remark}

\begin{definition}[Departure time; Cox construction]\label{sh:def:cox}
Let $E\sim\mathrm{Exp}(1)$ be independent of $W$ and define
$T:=\inf\{t\ge0:\Lambda_t\ge E\}\in(0,\infty]$. Then
$\PP(T>t\mid\Fc^W_\infty)=e^{-\Lambda_t}$ and
$\PP(T=\infty\mid\Fc^W_\infty)=e^{-\Lambda_\infty}$, with the convention
$e^{-\infty}=0$. The \emph{survival function} is
$S_t=\PP(T>t\mid\Fc^W_\infty)=e^{-\Lambda_t}$.
\end{definition}

\begin{definition}[Viability]\label{sh:def:viab}
The \emph{viability event} is $\Imm:=\{\Lambda_\infty<\infty\}$. On $\Imm$
the conditional survival probability $S_\infty=e^{-\Lambda_\infty}$ is
strictly positive; off $\Imm$ it vanishes. Thus
\begin{equation}\label{sh:eq:twotier}
\PP(T=\infty)\;=\;\E\!\left[e^{-\Lambda_\infty}\one_\Imm\right]
\;\le\;\PP(\Imm),
\end{equation}
and we distinguish throughout between \emph{viability} ($\PP(\Imm)$, the
probability that immortality is achievable) and \emph{immortality proper}
($\PP(T=\infty)$).
\end{definition}

% [Instates the 7:01 PM Aug 22 sheet entry, first half: the referent of
% the hazard, the anti-smuggling rationale, and the non-biconditional
% dissolution of the circularity worry.]
\begin{remark}[What the hazard is a hazard of]\label{gr1:rem:hazardof}
A reading question should be settled at the definition, because an
accusation hangs on it. If the Cox event is glossed ``biosphere
death,'' then the framework has installed at its root exactly the
primitives --- organism, death, biosphere --- it elsewhere claims to
individuate from survivorship structure, and the accusation of
smuggling would be fair. The gloss this work intends is different and
should be stated plainly: the Cox event is a \emph{departure-sufficient
occurrence} --- an event whose realization forecloses --- and $h$ is
the arrival rate of sufficiency, with nothing about cells, organisms,
or death presupposed. Death-events are then the typical
\emph{occupants} of the departure-sufficiency \emph{role}, which is
where their operational anchoring (extinction events are what a record
shows) survives intact --- but the role, not the occupant, is what the
model consumes, and the role is stated in the framework's own
primitive, the closure of possibility.

The circularity this invites is apparent and dissolves by direction.
The worry: the hazard is defined by reference to the Departure Point,
and the Departure Point (\S\ref{gc:subsub:cascade}) is built over the
hazard's own events. The dissolution: the implication runs one way,
and only one, by design. $\mathrm{Base}$ is defined by the analytic
model; $\mu G$ is its closure under inevitability; the containment
$\mu G\supsetneq\mathrm{Base}$ \eqref{gc:eq:departurefix} says the
Cox events \emph{suffice} for departure and are not necessary to it
--- so the hazard is an implies-departure object, never a
departure-defining one, and the unfavorable modal cases are added by
$G$'s recursion, never by re-writing the hazard. This is the same
non-biconditional shape as the companion's third postulate, and
deliberately so (Remark~\ref{gc:rem:thirdpostulate}): where a
biconditional would close a circle, a one-way implication that
recurses to a wider net leaves the definitional order a straight line
--- events at the root, closure above them, and the arrow never
reversing. The mathematics is untouched by any of this; what the
remark fixes is the interpretation layer, and with it the answer the
smuggling objection receives.
\end{remark}

% [Instates the 7:01 PM Aug 22 sheet entry, second half: the
% definitional-root premise.]
\begin{remark}[Golden-categorical well-definedness; \textbf{[$\Phi$]}]
\label{gr1:rem:goldendef}
The preceding remark removes death and biosphere from the root; this
one states what stands there instead, and it is entered as a premise
at its own grade because it is a semantic thesis and not a theorem.

\emph{The premise.} The well-definedness of real-world feature-terms
takes the Golden Point --- the asymptotic survivorship structure ---
as its first-principles root: a term is well-defined, in the sense
this framework can certify, to the degree its reference is robust
relative to that structure. The \emph{biosphere} itself is the first
application: it is individuated as the carrier of the survivorship
structure --- the unit whose continuations the fixed points of
\S\ref{gc:subsub:cascade} classify --- rather than presupposed as a
biological given. And the \emph{task-object} is not a second
application but the premise's constructive form: \S\ref{to:sec} is
where this premise builds its objects, the root Aethus of
Definition~\ref{to:def:taskobject} instantiated at the Golden Point
being the case from which the concept arose and the case at which
``contribution'' has a formal referent
(Remark~\ref{to:rem:goldenroot}). The premise and the task-object are
one doctrine at two registers --- criterion and construction --- and
the section develops the second. The premise generalizes a discipline the corpus
already practices: the companion requires categorical definitions to
survive transitions between neighboring Aethae --- its natural-kind
terms are rebuilt to that standard, ``dinosaur'' by anatomical
superclade and mechanism-role rather than by roster or genome
\cite{Nexus} --- and the present premise is that criterion's intended
completion: robustness across Aethae matures into robustness under
the Golden conditioning, Aethic-categorical into
\emph{Golden-categorical}, one criterion at two depths.

\emph{The price, paid openly.} Three fences. The model's mathematics
nowhere consumes the premise --- the theorems run on whatever unit is
plugged into them --- so a reader who declines it loses only the
answer to the smuggling objection at the level of reference, and
nothing downstream. The premise is semantic, not idealist:
well-definedness is not existence, and no feature waits on the
framework to be. And the self-anchoring is acknowledged rather than
hidden: rooting the framework's semantics in its own central object
is a circle of exactly the kind every foundational scheme carries
somewhere, and the claim made for this placement of it is
comparative --- the survivorship structure is a root the framework
can actually individuate, where pretheoretic biology hands it the
individuation problems of another science as axioms. The
non-biconditional shape guards here too: features may admit other
adequate definitions; the premise says where \emph{this} framework's
certificates come from, and only that, at \textbf{[$\Phi$]} with the
comparative claim its falsifier --- exhibit a unit-individuation the
framework needs that the survivorship root cannot supply and the
premise fails where it stands. One forward pointer completes the
placement: under the Route World's Strong reading the premise acquires
an ontological completion --- the categories trace to the root because
the \emph{things} do --- developed and priced at
\S\ref{to:sub:goldenaethus}, which this remark does not consume.
\end{remark}

The pair $(\mu,D)$ is the \emph{Kolmogorov diffusion} \cite{Kol34}:
\begin{equation}\label{sh:eq:kolm}
d\mu_t=\sigma\,dW_t,\qquad dD_t=\mu_t\,dt,\qquad
\Lc=\tfrac{\sigma^2}{2}\partial_\mu^2+\mu\,\partial_D .
\end{equation}

\begin{lemma}[Regularity]\label{sh:lem:hypo}
The free process $(\mu_t,D_t)$ is Gaussian with mean
$(\mu_0,\,D_0+\mu_0t)$ and covariance
$\bigl(\begin{smallmatrix}\sigma^2t&\sigma^2t^2/2\\
\sigma^2t^2/2&\sigma^2t^3/3\end{smallmatrix}\bigr)$;
in particular it admits the explicit smooth transition density written down
by Kolmogorov \cite{Kol34}. More generally, $\Lc$ is hypoelliptic: with
$X_1=\sigma\partial_\mu$ and $X_0=\mu\partial_D$ one has
$[X_1,X_0]=\sigma\partial_D$, so H\"ormander's bracket condition holds at
first order \cite{Hor67}. Consequently the process killed at rate $e^{x}$
retains smooth densities, which is the regularity relevant to the killed
problem of \S\ref{sh:subsub:reduction}.
\end{lemma}

\begin{proof}
Gaussianity and the moments follow from
Lemma~\ref{sh:lem:secondorder} below and linearity of
\eqref{sh:eq:model} in $W$; the bracket computation is immediate.
\end{proof}

Three structural remarks frame everything that follows.

\begin{remark}[No boundary in $h$]\label{sh:rem:noboundary}
The logarithmic parametrization makes $h>0$ automatic: there is no absorbing
or reflecting boundary in the hazard itself, and the entire problem concerns
the growth of the integrated driver.
\end{remark}

\begin{remark}[Credit is spent only through negative rate]\label{sh:rem:monot}
$D$ is nondecreasing exactly while $\mu\ge0$. Accumulated credit cannot be
destroyed except by first driving the improvement rate negative; this
monotone coupling is the source of all horizon intuition, and its
quantitative teeth are extracted in \S\ref{sh:subsub:softhorizon}.
\end{remark}

\begin{remark}[Nonparametric-Bayes reading of the model]\label{sh:rem:iwp}
Definition~\ref{sh:def:model} places an integrated Wiener process prior on
$-\log h$ --- equivalently, the classical cubic smoothing-spline prior of
Wahba \cite{Wah78}, since a prior whose second derivative is white noise
is precisely one whose first derivative is Brownian. The model is
therefore the canonical smoothness prior on the log-hazard rather than an
exotic construction, and the results below may be read as the survival
asymptotics that prior implies.
\end{remark}

\begin{lemma}[Exact gauge role of the initial hazard]\label{sh:lem:hgauge}
Write $\Lambda^{(h_0)}$ for the cumulative hazard with initial value $h_0$.
Then $\Lambda^{(h_0)}_t=h_0\Lambda^{(1)}_t$ pathwise; hence $\Imm$ does not
depend on $h_0$, while
$\PP(T=\infty)=\E[e^{-h_0\Lambda^{(1)}_\infty}\one_\Imm]$ is continuous and
nonincreasing in $h_0$, strictly decreasing whenever $\PP(\Imm)>0$, with
$\lim_{h_0\downarrow0}\PP(T=\infty)=\PP(\Imm)$ and
$\lim_{h_0\uparrow\infty}\PP(T=\infty)=0$.
\end{lemma}

\begin{proof}
The identity is \eqref{sh:eq:model}; the limits follow by monotone and
dominated convergence.
\end{proof}

Thus $h_0$ can never affect a dichotomy, only a value: it is the first and
most complete gauge degree of freedom. The same will be shown for
$(\mu_0,D_0)$ at the level of asymptotic exponents.

\subsubsection{Primitives and the hazard}\label{tc:subsub:primitives}

\begin{definition}[Two-channel model]\label{tc:def:model}
Fix $\sigma>0$, $\epsilon>0$, $n\in\mathbb N$, $\Gamma>0$ and
$\rho>1$. Define
\begin{align}
\text{driver:}&\quad \mu_t=\mu_0+\sigma W_t,
&&[\mu]=\mathsf T^{-1},\ [\sigma]=\mathsf T^{-3/2},
\label{tc:eq:driver}\\
\text{muffling exponent:}&\quad x_t=x_0+\int_0^t\mu_s\,ds,
&&[x]=1,
\label{tc:eq:x}\\
\text{background:}&\quad B_t=\sum_{i=1}^{n}e^{\epsilon W^{(i)}_t},
&&[B]=\mathsf T^{-1},\ [\epsilon]=\mathsf T^{-1/2},
\label{tc:eq:B}
\end{align}
and let the total hazard be the sum of two channels,
\begin{equation}\label{tc:eq:hazard}
h^{\mathrm{tot}}_t
=\underbrace{B_t\,\varsigma_b(x_t)}_{\text{crest}}
+\underbrace{\Gamma\,e^{-\rho x_t}}_{\text{trough}},
\qquad
\varsigma_b(x):=\frac{1+e^{-b}}{1+e^{x-b}}
=\frac{1+e^{-b}}{2}\Bigl(1-\tanh\tfrac{x-b}{2}\Bigr),
\end{equation}
with \emph{offset} $b>0$. The normalizing factor $1+e^{-b}$ is chosen so
that $\varsigma_b(0)=1$ exactly, whence $h^{\mathrm{crest}}|_{x=0}=B_t$:
the void limit is the background itself, for every $b$. Setting $b=0$
recovers the unshifted form $B_t(1-\tanh\frac{x}{2})$
Survival is $\bar F_t=e^{-\Lambda_t}$ with
$\Lambda_t=\int_0^t h^{\mathrm{tot}}_u\,du$, and the departure time $T$
is given by the Cox construction of the driftless case. There is \emph{no}
capacity parameter and \emph{no} leak: $x$ is the undamped integral of
$\mu$, exactly as in the driftless case --- a stipulation about the
\emph{completed} regime, which Definition~\ref{sch:def:leak} exhibits as
the $k=n$ boundary case of a staged leak. The same applies one derivative
up: the Brownian driver of \eqref{tc:eq:driver} is likewise a post-onset
object, since an improvement rate is something an agent supplies and
before the onset state there is none.
\end{definition}

\begin{remark}[Regimes]\label{tc:rem:regimes}
The three regimes of \eqref{tc:eq:hazard}:
\begin{center}\small
\begin{tabular}{@{}lll@{}}
\toprule
regime & total hazard & reading\\
\midrule
$x\gg b$ & $\approx B_te^{\,b-x}$ & shielding; background passes through
multiplicatively\\
$-k/\rho\ll x\ll b$ & $\approx B_t$ & \emph{void}: the system
contributes nothing\\
$x\ll0$ & $\approx\Gamma e^{\rho|x|}$ & internal damage; $B$ has dropped
out entirely\\
\bottomrule
\end{tabular}
\end{center}
The middle line is exact at $x=0$: the normalization gives
$h^{\mathrm{crest}}|_{x=0}=B_t$, so the void limit is the background
itself. The offset $b$ places $x=0$ \emph{inside} the void rather than
at its upper edge --- see Remark~\ref{tc:rem:whyoffset}.
\end{remark}

\subsubsection{The four constraints}\label{tc:subsub:constraints}

The model is answerable to four structural requirements. We state them
as they were posed and record where each is met.

\begin{constraint}[Background reaches the total at positive $x$]
\label{tc:C1}
Met for $\rho>1$, and \emph{only} then: for $x\gg b$ the crest channel
is $B_te^{\,b-x}$, multiplicative in $B$, and by
Proposition~\ref{tc:prop:dropout}(ii) it dominates the trough at all
large positive $x$ precisely when $\rho>1$. This --- not
Constraint~\ref{tc:C2} --- is what the brittleness condition buys.
\end{constraint}

\begin{constraint}[Background drops out at negative $x$]\label{tc:C2}
Met by Proposition~\ref{tc:prop:dropout}(i), for \emph{every} $\rho>0$.
Since $\Gamma$ is a primitive constant the trough contains no $B$
whatsoever, so at $x\ll0$ the background disappears from the total
entirely, not merely in its fluctuations.
\end{constraint}

\begin{constraint}[the driftless case undisturbed]\label{tc:C3}
Met as a theorem, \S\ref{tc:sub:asymptotics}: both walls are
$o(t^{3/2})$ and hence gauge.
\end{constraint}

\begin{constraint}[No smuggled primitives]\label{tc:C4}
Met by Theorem~\ref{tc:thm:gain}: with $B$ constant the total hazard is
$h^{\mathrm{tot}}(0)e^{-\Psi(x)}$ for a single soft-kink gain $\Psi$, so
the model is the driftless case's with the credit passed through a kink.
\end{constraint}

\begin{proposition}[Which channel dominates, and where; \textbf{P}]
\label{tc:prop:dropout}
\emph{(i) Negative $x$.} For every $\rho>0$,
$\Gamma e^{-\rho x}/\bigl(B\varsigma_b(x)\bigr)\to\infty$ as
$x\to-\infty$, since $\varsigma_b$ is bounded while $e^{\rho|x|}$ is
not. Hence $h^{\mathrm{tot}}\sim\Gamma e^{\rho|x|}$, free of $B$:
Constraint~\ref{tc:C2} needs no condition on $\rho$ beyond positivity.

\emph{(ii) Positive $x$.} As $x\to+\infty$ the crest is
$\sim Be^{\,b-x}$ and the trough is $\Gamma e^{-\rho x}$, so the crest
dominates for all large $x$ if and only if $\rho>1$. For $\rho<1$ the
two cross at
\begin{equation}\label{tc:eq:crossover}
x_{\times}=\frac{k+b}{1-\rho},\qquad k:=\log\frac{B}{\Gamma},
\end{equation}
beyond which the trough dominates and the background drops out
\emph{even for survivors}.
\end{proposition}

\begin{proof}
(i) $\varsigma_b\le1+e^{-b}$ by Proposition~\ref{tc:prop:bounded}, while
$\Gamma e^{\rho|x|}\to\infty$ for any $\rho>0$. (ii) Compare exponents
$b-x$ and $-k-\rho x$; they are equal at \eqref{tc:eq:crossover}, and
for $\rho>1$ the crest exponent is the larger for all $x$ exceeding a
finite threshold.
\end{proof}

\begin{remark}[What $\rho=1$ is, and is not]\label{tc:rem:rhorole}
Three separate claims must be kept apart, and the first two are easily
run together.
\emph{(a)} $\rho=1$ is \emph{not} the threshold for the background to
drop out at negative $x$; that happens for every $\rho>0$, by
(i).
\emph{(b)} $\rho=1$ \emph{is} the threshold at which the background
ceases to reach the survivor's hazard: for $\rho<1$ internal damage
overtakes the external channel beyond $x_{\times}$ and the system is
effectively closed to the external world, violating
Constraint~\ref{tc:C1}. This is a sharp threshold in a structural
parameter, and it concerns \emph{openness}, not survival.
\emph{(c)} $\rho$ does \emph{not} affect the survival exponent. For any
$\rho>0$ the dominant channel is an exponential in $x$, so the killing
wall is a level of $x$ and the persistence problem is the driftless case's up
to a rescaling of the credit by a constant. The exponent is $\tfrac14$
throughout, which strengthens Constraint~\ref{tc:C3} rather than
qualifying it.
\end{remark}

\subsubsection{Boundedness and the causal asymmetry}
\label{tc:subsub:asymmetry}

\begin{proposition}[The crest channel is bounded; \textbf{P}]
\label{tc:prop:bounded}
$0<\varsigma_b(x)<1+e^{-b}$ for all $x\in\R$, the supremum being
approached only as $x\to-\infty$. Hence the external channel can
\emph{reduce} the background hazard without bound, but can
\emph{increase} it by at most the factor $1+e^{-b}$, which tends to $1$
as $b$ grows.
\end{proposition}

\begin{proof}
$(1+e^{-b})/(1+e^{x-b})$ is decreasing in $x$ with limits $1+e^{-b}$ and
$0$.
\end{proof}

\begin{remark}[Why this matters]\label{tc:rem:asymmetry}
The slogan is ours: \emph{strength can shield from the
background but never amplify it}, and it is now enforced structurally
rather than assumed. The earlier multiplicative form $B_te^{-x}$
violated it --- at $x<0$ it multiplied the background without bound,
which asserts that a failure of the system's own muffling magnifies the
external world, a causal claim the ontology does not license.
Amplification now arises only through the trough channel, which is
internal damage: a distinct mechanism, correctly separated from the
external one. Proposition~\ref{tc:prop:dropout} is the formal expression
of that separation. Note how much the offset strengthens the slogan: at
$b=0$ the background could still be doubled, whereas for $b\gtrsim6$ the
permitted amplification is under a quarter of one percent. Shielding is
unbounded; amplification is, for practical purposes, forbidden.
\end{remark}

% ---------------------------------------------------------------------
\subsubsection{The construction as an iterated operator}
\label{gr1:subsub:operator}

The placement of the noise --- on $\mu$, the negative logarithmic
derivative of the hazard, rather than on the hazard or its logarithm ---
is the model's single most consequential choice, and this subsection
states the sense in which it is not a choice at all.

\begin{definition}[The hazard operator]\label{gr1:def:operator}
For a positive function $f$ of time write
\begin{equation}\label{gr1:eq:operator}
\mathcal D[f]\;:=\;-\frac{d}{dt}\log f .
\end{equation}
\end{definition}

\noindent Applied to the survival function it returns the hazard, and
applied to the hazard it returns this model's driver:
\begin{equation}\label{gr1:eq:ladder}
\mathcal D[\bar F]=h,\qquad \mathcal D[h]=\mu .
\end{equation}
The driver is therefore the \emph{second iterate of the operation that
defines the hazard in the first place}, and the model's noise sits at
level two of a recursion whose level one is the hazard itself. This is a
different statement from ``we adopt a smoothness prior on the log-hazard
and choose the integrated-Wiener member of the family,'' which is how
the same object appears from the nonparametric-Bayes side, and it is a
stronger one: no member of the family is being selected, because the
ladder \eqref{gr1:eq:ladder} has rungs and the model occupies one of
them.\footnote{%
\emph{Provenance, with its stages kept apart.} The operator
\eqref{gr1:eq:operator} is ours, derived in May 2021 by a route
worth recording because it is the physical one rather than the textbook
one: examine $\frac{d}{dt}(1-\bar F)$ at $t=0$, read the resulting
quantity as an instantaneous \emph{felt danger} propagating through the
system, rewrite the survival function with that quantity as a parameter,
replace the product $ht$ in the exponent by the time integral
$\int_0^t h$, and solve for $h$ --- which returns
\eqref{gr1:eq:operator}. The hazard thus arrives as a rate one can point
at, not as a definition one is handed, which is the same difference that
\S\ref{gr1:subsub:operator} turns on. The \emph{second} application is
later and its stages should not be run together: our record
places the present stochastic use of $\mu$ on the hazard in autumn 2025,
with an earlier and as yet unlocated use of a doubled operator on the
survival function, in connection with the Golden Point rather than with
hazard rates, somewhere in 2021--2023 and under a name of its own. That
earlier use is author-attested and not verified here; if the journal
entries are recovered, this footnote should be revised to cite them.
The logistic deformation of \S\ref{tc:subsub:asymmetry} is later still,
arising during the preparation of the present work from the causal
constraint recorded there.}

\begin{remark}[Where the ladder stops]\label{gr1:rem:ladderstops}
The rungs are not equally hospitable, and the model sits on the last one
that is analytically open. Noise at level one --- $\log h$ Brownian ---
makes the credit Brownian and the persistence exponent $\tfrac12$. Noise
at level two --- $\mu$ Brownian, the present model --- makes the credit
integrated Brownian and the exponent $\tfrac14$
(Theorem~\ref{sh:thm:persist}). At level three the credit is a twice
integrated Brownian motion, and no closed form for the exponent is
known: existence and monotonicity are available for fractionally
integrated processes \cite{AD13,AS15}, and the persistence of iterated
partial sums has been studied, but the value is open. Level one is close
to a definition and level three is close to an open problem; level two
is the rung on which the object is both nontrivial and solvable, and it
is the rung the hazard construction reaches by being applied twice.
\end{remark}

\begin{remark}[The two channels, in operator terms, and why they differ]
\label{gr1:rem:channelsoperator}
Applying $\mathcal D$ to the two channels of \eqref{tc:eq:hazard}
separately, at fixed background, gives
\begin{equation}\label{gr1:eq:channelops}
\mathcal D\bigl[h^{\mathrm{trough}}\bigr]=\rho\,\mu,
\qquad
\mathcal D\bigl[h^{\mathrm{crest}}\bigr]=\varsigma(x-b)\,\mu,
\qquad
\varsigma(u)=\frac{e^{u}}{1+e^{u}} .
\end{equation}
The internal channel therefore preserves the operator relation
\emph{exactly}, with constant gain $\rho$; the external channel
preserves it only up to a logistic weight running from $0$ to $1$.

The asymmetry is the causal constraint of
\S\ref{tc:subsub:asymmetry} in its sharpest form, and it is stated as such. An exact relation $h=h_0e^{-x}$ is unbounded above as
$x\to-\infty$: it says that a failure of the system's own muffling
multiplies the background hazard without limit, which is a causal claim
the ontology declines --- a biosphere may fail to shield itself from the
external world, but it cannot thereby make that world more dangerous
than it is. No such constraint applies inward. \emph{One may harm
oneself without bound; one may not amplify the world.} Hence the trade
recorded in \eqref{gr1:eq:channelops}: the internal channel keeps the
exact relation because nothing forbids it, and the external channel must
be deformed because something does. The logistic is the minimal
deformation purchasing boundedness, and the offset $b$ is how much of it
is bought.

Three consequences are already in the text and are named here rather
than derived again. The technoanatomy $\rho$ of
Definition~\ref{gr1:def:anatomy} \emph{is} the internal channel's
operator gain, $\mathcal D[h^{\mathrm{int}}]/\mu$, which is a third and
perhaps the cleanest answer to why that exponent deserves the name: it
is the rate at which effort converts into change in self-inflicted
hazard, protective when positive and wounding when negative. The void
of \S\ref{gr1:sub:anatomy} is the region where $\varsigma(x-b)\approx0$,
so that $\mathcal D[h^{\mathrm{crest}}]\approx0$ even for $\mu\neq0$ ---
the operator is applied and nothing comes out, which is what
Proposition~\ref{tc:prop:mingain} measures from the other side. And the
concealed gain $\Psi'$ of Corollary~\ref{tc:cor:concealed} is exactly
the hazard-weighted blend of the two relations in
\eqref{gr1:eq:channelops}, which is why its range is $(0,\rho)$ and why
it tends to $1$ under shielding and to $\rho$ under damage.
\end{remark}


\subsubsection{Reduction to a persistence problem}\label{sh:subsub:reduction}

\begin{proposition}[Feynman--Kac]\label{sh:prop:FK}
Let $u(t,x,\mu):=\PP_{x,\mu}(T>t)=\E_{x,\mu}\bigl[e^{-\Lambda_t}\bigr]$.
Then $u$ is the (classical, by Lemma~\ref{sh:lem:hypo}) solution of
\begin{equation}\label{sh:eq:FK}
\partial_tu=\frac{\sigma^2}{2}\partial_\mu^2u-\mu\,\partial_xu-e^{x}u,
\qquad u(0,\cdot,\cdot)\equiv1 .
\end{equation}
\end{proposition}

\begin{proof}
It\^o's formula applied to $s\mapsto
u(t-s,x_s,\mu_s)\exp(-\int_0^se^{x_r}dr)$ exhibits the right side as a
martingale if and only if \eqref{sh:eq:FK} holds; regularity is supplied by
hypoellipticity of the killed generator.
\end{proof}

The killing rate $e^x$ vanishes rapidly for $x\ll0$ and diverges for
$x\gg0$: at exponent level it is an absorbing wall at $x=0$, of thickness
$O(1)$ in $x$ --- a factor $e$ in the hazard --- which is negligible against
excursions of size $t^{3/2}$. We elevate this to a labeled hypothesis, since
it is the single non-rigorous reduction on which the tail asymptotics of
\S\ref{sh:subsub:tail} rest.

\begin{ansatz}[Soft-wall reduction]\label{sh:ansatz}
Let $\tau:=\inf\{t\ge0:D_t\le0\}$ denote the hard-wall passage time of the
Kolmogorov diffusion. Then for every admissible initial datum there exist
constants $0<c\le C<\infty$ and $t_0<\infty$ with
\[
c\,\PP(\tau>t)\ \le\ \PP(T>t)\ \le\ C\,\PP(\tau>t),\qquad t\ge t_0,
\]
after an $O(1)$ adjustment of the initial credit. Equivalently,
$\log\PP(T>t)=\log\PP(\tau>t)+O(1)$.
\end{ansatz}

\begin{remark}[Status of the Ansatz]\label{sh:rem:ansatzstatus}
One half is rigorous modulo a cited input. \emph{Lower bound:} on the event
$\{D_u\ge f(u)\ \forall u\le t\}$ with $f(u)=2\log\frac{2+u}{2}$ one has
$\Lambda_t\le h_0e^{-D_0}\int_0^\infty\bigl(\tfrac{2+u}{2}\bigr)^{-2}du<\infty$,
whence $\PP(T>t)\ge e^{-C}\,\PP(D\ge f\text{ on }[0,t])$; that a logarithmic
moving boundary does not change the persistence exponent of an integrated
process is the standard moving-boundary phenomenon (see the survey
\cite{AS15} and references therein; cf.\ also the robustness expressed
by the universality results \cite{AD13,DW15}). \emph{Upper bound:} from
$\PP(T>t)\le e^{-K}+\PP(\Lambda_t\le K)$ and
$\Lambda_t\ge h_0e^{-D_0}\Leb\{u\le t:D_u\le0\}$, the event
$\{\Lambda_t\le K\}$ forces $D$ to remain positive off a set of bounded
Lebesgue measure. What is missing is the (plausible, unproved here)
statement that \emph{persistence with a bounded occupation tolerance has the
same exponent as strict persistence}. Everything tagged \textbf{[A]} below
depends on exactly this gap and nothing else; \S\ref{sh:sub:numerics}
describes its direct numerical test.

The standing literature nearest the gap should be cited against it,
because a referee will and should. Persistence with \emph{partial
survival} --- each zero crossing killing with probability $p$ ---
carries exponents that vary continuously with $p$
\cite{MB98, Bur14}, so ``a softened wall leaves the exponent alone''
is false in general, and the Ansatz cannot rest on softness as such.
What it rests on is the structure of this softening: the killing here
is rate-based and confined to a boundary layer of width $O(1)$ in the
credit, against excursions of scale $t^{3/2}$ --- a vanishing fraction
of the process's range, where partial survival taxes every crossing
at fixed price no matter the scale. The distinction is exactly what
the numerics interrogate, and the measured constancy of the
soft-to-hard ratio is evidence that the boundary-layer case sits in
the exponent-preserving class the moving-boundary results
\cite{AS15, DW15} delimit from the rigorous side.
\end{remark}

\subsubsection{The transfer, stated}\label{tc:subsub:transferstated}

\begin{theorem}[Transfer of the driftless case; \textbf{A}]
\label{tc:thm:transfer}
In the two-channel model, for every admissible parameter set,
\begin{equation}\label{tc:eq:transfer}
\PP(T>t)\;\asymp\;\underbrace{\hm\bigl(\mu_{0}-\tfrac{\epsilon^{2}}{2},
\;D_{0}-\log n-b\bigr)}_{\text{shifted initial data}}\;\cdot\;
\underbrace{e^{-\Theta\left(n(b/\sigma)^{2/3}\right)}}_{\text{room
passage}}\;\cdot\;t^{-1/4},
\end{equation}
with $\hm$ the persistence-harmonic function of the driftless case. In
consequence: $\E[T]=\infty$ while $T<\infty$ a.s.; residual life is
asymptotically Pareto$(\tfrac14)$; quantiles grow as the fourth power of
rarity; the effective population hazard is $1/4t$; and the survivorship
interpolation, its pure Doob limit, its $s/4t$ linear response and its
critical mode all hold verbatim, the limit being the same $Q$-process
computed at the shifted data.
\end{theorem}

\begin{proof}
Above the room the hazard is the driftless case's up to bounded constants
(Lemma~\ref{tc:lem:aboveroom}); the wall is linear and absorbed exactly
into the shifted data (Theorem~\ref{tc:thm:absorb}, with the offset $b$
contributing a further constant); the band collapses
(Theorem~\ref{tc:thm:collapse}); the room contributes a bounded additive
term to $\Lambda$, hence a multiplicative prefactor
(Proposition~\ref{tc:prop:roomcost}). The listed consequences are those
of the driftless case, transported through
\eqref{tc:eq:transfer}.
\end{proof}

\begin{remark}[How close, in one sentence]\label{tc:rem:howclose}
Exactly as close as it is possible to be: every difference between the
two models is a shift of initial data or a bounded multiplicative
constant, and we proved both to be gauge. The two-channel model
buys its ontology --- a positive fluctuating background, a bounded
external channel obeying the causal asymmetry, an unbounded internal
channel, an ordered pair of onsets and a waiting room --- at a cost of
precisely zero asymptotic exponents, and it repays part of the price by
shrinking the one gap the earlier section had to assume.
\end{remark}

% ---------------------------------------------------------------------
\subsection{Coordinates: gauge, strength, anatomy}\label{gr1:sub:coords}\label{tc:sub:polar}

\subsubsection{Dimensional structure, self-similarity, and the gauge orbit}
\label{sh:subsub:dimensional}

\begin{definition}[Units and intrinsic scales]\label{sh:def:units}
Assigning $[t]=\mathsf T$ gives $[\mu]=\mathsf T^{-1}$,
$[\sigma]=\mathsf T^{-3/2}$, $[D]=1$, $[h]=\mathsf T^{-1}$. The unique
timescale constructible from $\sigma$ alone is $\sigma^{-2/3}$; the
time-equivalents of the initial data are
\begin{equation}\label{sh:eq:timescales}
t_\mu:=\frac{\mu_0^2}{\sigma^2},\qquad
t_D:=\Bigl(\frac{D_0}{\sigma}\Bigr)^{2/3}\quad(D_0>0);
\end{equation}
and the unique dimensionless state coordinate on the quadrant
$\{\mu>0,D>0\}$ is the \emph{horizon coordinate}
\begin{equation}\label{sh:eq:kappa}
\kappa:=\frac{\sigma^2D}{\mu^3},\qquad\text{satisfying}\qquad
\frac{t_D}{t_\mu}=\kappa_0^{2/3}.
\end{equation}
Hatted variables are $\hat\mu=\mu\sigma^{-2/3}$, $\hat D=D$,
$\hat h=h\sigma^{-2/3}$, $\hat t=t\sigma^{2/3}$.
\end{definition}

\begin{lemma}[Second-order structure]\label{sh:lem:secondorder}
$\Var(I_t)=t^3/3$, $\Cov(W_t,I_t)=t^2/2$, and for $u\ge1$,
$\E[I_1I_u]=\tfrac u2-\tfrac16$.
\end{lemma}

\begin{proof}
By Fubini, $\E[I_aI_b]=\int_0^a\!\int_0^b(\alpha\wedge\beta)\,d\beta\,d\alpha$.
For $a=b=t$ this is $2\int_0^t\!\int_0^\beta\alpha\,d\alpha\,d\beta=t^3/3$.
Next, $\E[W_tI_t]=\int_0^t(\alpha\wedge t)\,d\alpha=t^2/2$. Finally, for
$u\ge1$, $\int_0^u(\alpha\wedge\beta)\,d\beta
=\alpha u-\alpha^2/2$ for $\alpha\le1\le u$, and integrating over
$\alpha\in[0,1]$ gives $u/2-1/6$.
\end{proof}

\begin{proposition}[Exact self-similarity]\label{sh:prop:selfsim}
Fix $\sigma$ and $c>0$. If $(\mu,D)$ starts from $(\mu_0,D_0)$, then
\[
\bigl(\mu_{c^2t},\,D_{c^2t}\bigr)_{t\ge0}\ \eqd\
\bigl(c\,\mu'_t,\,c^3D'_t\bigr)_{t\ge0},
\]
where $(\mu',D')$ starts from $(\mu_0/c,\,D_0/c^3)$. Equivalently,
$(I_{ct})_{t\ge0}\eqd(c^{3/2}I_t)_{t\ge0}$.
\end{proposition}

\begin{proof}
Brownian scaling $(W_{c^2t})_t\eqd(cW_t)_t$, integrated once in time.
\end{proof}

\begin{proposition}[The gauge orbit]\label{sh:prop:gauge}
The initial data $(h_0,\mu_0,D_0)$ affect asymptotics only through
prefactors and timescales, never through exponents. Precisely:
\emph{(i)} $h_0$ is exactly gauge in the sense of
Lemma~\ref{sh:lem:hgauge};
\emph{(ii)} the deterministic contribution
$x_0-\mu_0t$ to $x_t$ is affine in $t$, hence
$(x_0-\mu_0t)/t^{3/2}\to0$, while $\sigma I_t/t^{3/2}\eqd\sigma I_1$ is
nondegenerate for every $t$: the initial data live in the subspace that the
$t^{3/2}$ fluctuations cannot resolve;
\emph{(iii)} under the scaling of Proposition~\ref{sh:prop:selfsim} the
initial data transform as $(\mu_0,D_0)\mapsto(c\mu_0,c^3D_0)$, so any
scale-invariant asymptotic quantity is constant on the two-parameter orbit
$\{(c\mu_0,c^3D_0):c>0\}$ and, when finite and positive, is a function of
$\kappa_0$ alone.
\end{proposition}

\begin{proof}
(i) is Lemma~\ref{sh:lem:hgauge}; (ii) is immediate from
Proposition~\ref{sh:prop:selfsim} and Lemma~\ref{sh:lem:secondorder};
(iii) holds because the orbits of the scaling action on the open quadrant
are exactly the level curves of $\kappa$.
\end{proof}

The sharp form of the slogan --- that even the \emph{prefactors} of survival
asymptotics depend on the data only through a single equivalent timescale ---
is Proposition~\ref{sh:prop:degrees} below.

% ---------------------------------------------------------------------
\subsubsection{The sign of the credit, and an exact half}
\label{gr1:subsub:sign}

One fact about the driftless model is used repeatedly below and is worth
isolating, because it is the engine of a proof and not merely a
description.

\begin{theorem}[Adverse configuration occupies exactly half of
logarithmic time; \textbf{P}]\label{tc:thm:half}
For the driftless driver, almost surely
\begin{equation}\label{tc:eq:half}
\lim_{S\to\infty}\frac1S\,\Leb\bigl\{s\in[0,S]:x_{e^{s}}<0\bigr\}
=\frac12 ,
\end{equation}
so the set $\{t:x_t<0\}$ --- the times at which the system's own
muffling sits \emph{above} rather than below its background hazard ---
has logarithmic density exactly $\tfrac12$, for every parameter value
and every initial condition.
\end{theorem}

\begin{proof}
$Z_s:=x_{e^{s}}/(\sigma e^{3s/2})$ is stationary, centered Gaussian and
ergodic (the Lamperti argument of Theorem~\ref{sh:thm:noimm}), so
Birkhoff gives density $\PP(Z_0<0)=\tfrac12$ by symmetry.
\end{proof}

\begin{remark}[What the exact half is for, and what it is not]
\label{gr1:rem:halfuse}
Two uses and one disclaimer.

\emph{Use.} It is the engine of Theorem~\ref{tc:thm:noviab2}: where the
driftless argument needed an unbounded integrand, the two-channel model
needs only that the system spends a positive fraction of logarithmic
time at hazard above background --- and the fraction is exactly one
half, for every parameter value. Boundedness of the hazard does not buy
immortality; it changes which argument rules it out.

\emph{Sharpening.} The no-viability proof of
\S\ref{gr1:sub:verdict} used only the qualitative input
$\PP(Z_0<-\varepsilon)>0$; here the constant is exact.

\emph{Disclaimer, and it is the important one.} This is a statement
about a \emph{state coordinate}, not about technoanatomy. Every
biosphere, however well constituted, spends half of logarithmic time
with $x<0$; the sign of the credit is a condition the system passes
through and re-enters, not a property it has. Technoanatomy in the sense
of Definition~\ref{gr1:def:anatomy} is the exponent $\rho$, which governs
what capacity \emph{does} when acquired and does not fluctuate at all.
Nothing about the frequency with which a biosphere is adversely
positioned bears on whether its capacity is adversely directed, and the
theorem should not be read as saying otherwise.
\end{remark}


% ---------------------------------------------------------------------
\subsection{The verdict: certain departure, infinite expectation}
\label{gr1:sub:verdict}\label{sh:sub:verdict}

\begin{remark}[What $T$ is the time of]\label{gr1:rem:whatTis}
Throughout this section $T$ is the time of the \emph{Departure Point}
(Definition~\ref{gc:def:departure}) --- the moment past which no
continuation advances --- not of biological death, and
$h^{\mathrm{tot}}$ is that event's hazard. Death is one mechanism of
foreclosure and irreversible loss of the capacity to advance is
another; the analysis does not distinguish them, because for
advancement they are the same event. It does, however, distinguish both
from a mere lull: a biosphere that has stopped progressing but could
still resume has not reached $T$. Where the classical results imported below speak of survival
and first passage, the vocabulary is theirs and the referent here is
departure. Read that way, $\PP(\Lambda_\infty<\infty)$ is the
probability of \emph{advancing}, not of not dying --- which is what the
Golden Point was always a claim about.
\end{remark}

\subsubsection{Impossibility of immortality}\label{sh:subsub:noimm}

\begin{theorem}[No viability; \textbf{P}]\label{sh:thm:noimm}
For every $(h_0,\mu_0,D_0,\sigma)$ in the driftless model
\eqref{sh:eq:model}, $\Lambda_\infty=\infty$ almost surely. Consequently
$S_\infty=0$ a.s., $\PP(\Imm)=0$, $\PP(T=\infty)=0$, and $T<\infty$ a.s.
\end{theorem}

\begin{proof}
Since $\Lambda_\infty=h_0e^{-D_0}\int_0^\infty e^{-\mu_0u-\sigma I_u}du$
and the prefactor is a fixed positive constant, it suffices to prove that
the integral diverges a.s.

\emph{Step 1 (Lamperti transform).} Define $Z_s:=e^{-3s/2}I_{e^s}$,
$s\ge0$. As a linear functional of $W$, $Z$ is a centered Gaussian process.
By Proposition~\ref{sh:prop:selfsim}, for each $r\ge0$,
$(Z_{s+r})_{s\ge0}=(e^{-3(s+r)/2}I_{e^re^s})_{s\ge0}
\eqd(e^{-3s/2}I_{e^s})_{s\ge0}=(Z_s)_{s\ge0}$: $Z$ is stationary. By
Lemma~\ref{sh:lem:secondorder}, for $s\ge0$,
\begin{equation}\label{sh:eq:lampcov}
\E[Z_0Z_s]=e^{-3s/2}\Bigl(\frac{e^{s}}{2}-\frac16\Bigr)
=\frac12e^{-s/2}-\frac16e^{-3s/2}\ \longrightarrow\ 0 .
\end{equation}

\emph{Step 2 (ergodicity).} A stationary Gaussian process whose covariance
vanishes at infinity is mixing \cite{Mar49}, hence ergodic.

\emph{Step 3 (Birkhoff).} Fix $\varepsilon>0$. Since
$Z_0\sim N(0,\tfrac13)$, $p_\varepsilon:=\PP(Z_0\le-\varepsilon)>0$. By the
continuous-time Birkhoff theorem,
$\tfrac1S\Leb\{s\in[0,S]:Z_s\le-\varepsilon\}\to p_\varepsilon$ a.s., so
$\Leb\{s\ge0:Z_s\le-\varepsilon\}=\infty$ a.s.

\emph{Step 4 (change of variables).} Let
$A:=\{u\ge1:I_u\le-\varepsilon u^{3/2}\}$. Under $u=e^s$ the set in Step~3
maps into $A$ with
$\Leb\{s:Z_s\le-\varepsilon\}=\int_A\frac{du}{u}\le\Leb(A)$,
the inequality because $u\ge1$. Hence $\Leb(A)=\infty$ a.s.

\emph{Step 5 (divergence).} Choose $u_1=u_1(\varepsilon,\mu_0,\sigma)$ so
that $\sigma\varepsilon u^{3/2}-\mu_0u\ge0$ for $u\ge u_1$. Then
\[
\int_0^\infty e^{-\mu_0u-\sigma I_u}\,du
\ \ge\ \int_{A\cap[u_1,\infty)}e^{\sigma\varepsilon u^{3/2}-\mu_0u}\,du
\ \ge\ \Leb\bigl(A\cap[u_1,\infty)\bigr)\ =\ \infty
\qquad\text{a.s.}\qedhere
\]
\end{proof}

\begin{remark}[The asymmetry of exponential response]\label{sh:rem:asym}
The driver is symmetric; the response is not. A favorable excursion
suppresses the integrand $e^{-D}$ multiplicatively toward zero ---
\emph{bounded} benefit for $\Lambda$ --- while an unfavorable excursion of
relative size $\varepsilon$ inflates it without bound. Exponential response
to a symmetric driver is structurally hostile to survival; the mechanism is
the sign-oscillation that makes low-dimensional random walks recurrent,
aggravated by this asymmetry. Identifying exactly what a horizon must break
--- the symmetry of the driver, not the response --- is the burden of
\S\ref{sh:subsub:scale}.
\end{remark}

\subsubsection{No viability, by a different route}\label{tc:subsub:noviab}

the driftless case's no-viability theorem was proved by exhibiting an integrand
that \emph{blows up}: on a set of infinite measure the credit satisfies
$D_u\le-\varepsilon u^{3/2}$, whence $e^{-D_u}\ge
e^{\sigma\varepsilon u^{3/2}}\to\infty$. That argument is unavailable
here whenever the crest channel acts alone, since $\varsigma_b$ is
\emph{bounded}. The conclusion nevertheless survives, by a weaker and
more robust mechanism: a positive floor on an infinite set. The
mechanism needs one lemma, and the lemma is worth stating in the
generality it naturally has.

\begin{lemma}[Infinite sets defeat the driftless background;
\textbf{P}]\label{tc:lem:infbackground}
Let $W$ be a standard Brownian motion, $\epsilon>0$, and let
$S\subset[0,\infty)$ be a measurable set, possibly random but
independent of $W$, with $\mathrm{Leb}(S)=\infty$ almost surely. Then
\[
\int_S e^{\epsilon W_u}\,du\;=\;\infty\qquad\text{almost surely,}
\]
and consequently the same holds with the background
$B_u=\sum_{i=1}^n e^{\epsilon W^{(i)}_u}\ge e^{\epsilon W^{(1)}_u}$
of \eqref{tc:eq:B} in place of $e^{\epsilon W_u}$.
\end{lemma}

\begin{proof}
Conditioning on $S$, it suffices to treat deterministic $S$ of
infinite measure. For each $T$, finiteness of
$\int_{S\cap[T,\infty)}e^{\epsilon W_u}du$ is unchanged by the
multiplicative factor $e^{\epsilon W_T}$ --- it is decided by
$\int_{S\cap[T,\infty)}e^{\epsilon(W_u-W_T)}du$ --- and the integral
over $S\cap[0,T]$ is a.s.\ finite, so the event
$\{\int_S e^{\epsilon W_u}du=\infty\}$ lies in
$\sigma(W_u-W_T:u\ge T)$ for every $T$, hence in the tail
$\sigma$-field of the increments, and has probability $0$ or $1$ by
Kolmogorov's zero--one law. Suppose the probability were $0$. Then
$\int_S e^{\epsilon W_u}du<\infty$ a.s., and applying the same
conclusion to the Brownian motion $-W$ gives
$\int_S e^{-\epsilon W_u}du<\infty$ a.s.\ as well; adding,
\[
2\,\mathrm{Leb}(S)\;\le\;\int_S\bigl(e^{\epsilon W_u}
+e^{-\epsilon W_u}\bigr)du\;<\;\infty,
\]
contradicting $\mathrm{Leb}(S)=\infty$.
\end{proof}


\begin{theorem}[No viability in the two-channel model; \textbf{P}]
\label{tc:thm:noviab2}
For every $\sigma,\epsilon>0$, $n\ge1$, $b\ge0$, $\rho>0$ and every
$\Gamma\ge0$ --- including $\Gamma=0$, where the total hazard is
bounded --- one has $\Lambda_{\infty}=\infty$ almost surely. Hence
$\bar F_{\infty}=0$ a.s.\ and $T<\infty$ a.s.
\end{theorem}

\begin{proof}
Discard the trough, which is nonnegative. On $\{x_u\le0\}$ we have
$\varsigma_b(x_u)\ge\varsigma_b(0)=1$, so
\[
\Lambda_{t}\;\ge\;\int_{0}^{t}B_u\varsigma_b(x_u)\,du
\;\ge\;\int_{S\cap[0,t]}B_u\,du,
\qquad S:=\{u:x_u\le0\}.
\]
By Theorem~\ref{tc:thm:half} the set $S$ has logarithmic density
$\tfrac12$; writing $u=e^{s}$,
$\mathrm{Leb}(S\cap[0,t])=\int_{0}^{\log t}\one\{Z_s\le0\}e^{s}ds
\to\infty$ a.s., since the integrand is positive on a set of density
$\tfrac12$ and $e^{s}\to\infty$. So $S$ is an $x$-measurable set of
infinite measure, independent of the background's driving motions, and
Lemma~\ref{tc:lem:infbackground} gives
$\int_S B_u\,du=\infty$ a.s., whence $\Lambda_{\infty}=\infty$.
\end{proof}

% CORRECTION RECORD (referee pass, Aug 22 2026): the previous proof
% concluded from (inf_{u<=t} B_u) x Leb{x<=0}, an indeterminate
% (->0)x(->inf) product, since inf_t B_t = 0 a.s. -- a fact the paper
% itself records at Remark gr1:rem:trappedscope.  The gap bit exactly
% at the Gamma=0 clause the statement advertises.  The lemma above
% (tail zero-one law + W -> -W symmetry) closes it in the generality
% the model needs; the theorem's conclusion was true as stated.



\subsubsection{The survival tail and its consequences}\label{sh:subsub:tail}

\begin{theorem}[Persistence of integrated Brownian motion; \textbf{C}]
\label{sh:thm:persist}
Let $(\mu,D)$ start from $(m,0)$ with $m>0$ and let
$\tau=\inf\{t:D_t\le0\}$. There is a constant $c_1>0$ such that
\begin{equation}\label{sh:eq:persist}
\PP_{(m,0)}(\tau>t)\ \sim\ c_1\Bigl(\frac{m^2}{\sigma^2t}\Bigr)^{1/4},
\qquad t\to\infty .
\end{equation}
The exponent $\tfrac14$ --- half the Brownian value, halved by one
integration --- is the persistence exponent of integrated Brownian
motion, and its provenance is a chain worth recording precisely. McKean
\cite{McK63} computed the first-passage (``winding'') law of the
Kolmogorov diffusion; Goldman \cite{Gol71} derived the passage-rate
asymptotic from McKean's law; Lachal \cite{Lac91} obtained the exact
first-passage distributions; Sinai \cite{Sin92} initiated and settled,
up to constants, the discrete counterpart $n^{-1/4}$ for the integrated
random walk; Isozaki--Watanabe \cite{IW94} supplied the sharp constants;
and Denisov--Wachtel \cite{DW15} proved the exact $n^{-1/4}$ asymptotics
universally for centered integrated random walks with $2+\delta$
moments, via a discrete potential theory (see also the universality
result of Aurzada--Dereich \cite{AD13} for integrated processes).
Starting states in the open quadrant obey the same power with a
prefactor proportional to the harmonic function of
Proposition~\ref{sh:prop:harmonic}; this general-start form is
established in \cite{GJW99}, which recovers \cite{Gol71,IW94} as
special cases. For the surrounding theory see the survey \cite{AS15}
and, on the physics side --- where the model is the \emph{random
acceleration process} and has an extensive literature of its own ---
the review \cite{BMS13}.
\end{theorem}

\begin{remark}[The $m^{1/2}$ law is exact given the $m=1$ asymptotic]
By Proposition~\ref{sh:prop:selfsim} with $c=1/m$ (and $\sigma=1$),
$\PP_{(m,0)}(\tau>t)=\PP_{(1,0)}(\tau>t/m^2)$, so once
\eqref{sh:eq:persist} is known for one $m$ it holds for all, prefactor
scaling as $m^{1/2}$ exactly. The initial rate enters through its
time-equivalent $t_\mu$ and nothing else.
\end{remark}

\begin{claim}[Survival tail; \textbf{A}]\label{sh:claim:tail}
Under Ansatz~\ref{sh:ansatz}, for every initial datum there are
constants $0<c_-\le c_+<\infty$, depending on
$(\hat h_0,\hat\mu_0,\hat D_0)$, with
$c_-t^{-1/4}\le\PP(T>t)\le c_+t^{-1/4}$ for all large $t$;
equivalently
\begin{equation}\label{sh:eq:tail}
\PP(T>t)\ \asymp\ t^{-1/4},\qquad t\to\infty .
\end{equation}
\end{claim}

All of the following consequences are elementary integrations of
\eqref{sh:eq:tail}; each inherits the tag \textbf{[A]}. Write
$\alpha=\tfrac14$.

\begin{proposition}[Consequences of the tail; \textbf{A}]
\label{sh:prop:tailconseq}\label{sh:cor:moments}\label{sh:cor:mean}%
\label{sh:cor:trunc}\label{sh:cor:quant}\label{sh:cor:residual}%
\label{sh:cor:lindy}
All of the following are elementary consequences of \eqref{sh:eq:tail}
by regular variation, with $\alpha=\tfrac14$; we state them together
rather than prove them severally, since each is a one-line integration
and none is original.

\begin{center}\small
\begin{tabular}{@{}lll@{}}
\toprule
quantity & value & reading\\
\midrule
$\E[T^{p}]$ & finite iff $p<\tfrac14$ &
even $\E[\sqrt T]=\infty$; heavier-tailed than Cauchy\\
$\E[T]$ & $\infty$, every initial condition &
null recurrence: certain departure, infinite mean\\
$\E[T\wedge\mathcal T]$ & $\asymp\tfrac{4C}{3}\mathcal T^{3/4}$ &
an observed mean is an artifact of the window\\
$t_q$ with $\PP(T>t_q)=q$ & $\asymp(C/q)^{4}$ &
quantiles grow as the fourth power of rarity\\
$\PP(T>\lambda t\mid T>t)$ & $\to\lambda^{-1/4}$ &
residual life Pareto$(\tfrac14)$, mean infinite at every age\\
$r(t)=-\tfrac{d}{dt}\log\PP(T>t)$ & $\approx 1/4t$ &
a Lindy law; $\int^\infty r=\infty$, so departure is certain\\
\bottomrule
\end{tabular}
\end{center}

\noindent Two remarks are worth more than the derivations. The truncated
mean shows that any finite dataset reports a finite mean lifetime
growing as the $3/4$ power of its own duration, so the quantity is not
estimable. And the effective hazard $1/4t$ sits far from the mortality
boundary --- any positive exponent gives certain departure --- but deep
inside the infinite-mean regime, whose boundary is exponent $1$.
\end{proposition}

% ---------------------------------------------------------------------
\subsection{Technoanatomy: the sign of the trough exponent}
\label{gr1:sub:anatomy}

Section~\ref{gr1:sub:object} introduced the trough channel
$\Gamma e^{-\rho x}$ under the standing restriction $\rho>1$, imposed so
that the background would reach the survivor's hazard
(Constraint~\ref{tc:C1}). That restriction has never been examined, and
lifting it turns out to supply the model's missing structural
parameter. We lift it now: $\rho$ ranges over $\R$.

\begin{definition}[Technoanatomy]\label{gr1:def:anatomy}
The exponent $\rho$ is the \emph{technoanatomy}\footnote{%
\emph{Provenance.} The notion, its square--cube framing, and the
exponential-against-linear scaling that
\S\ref{gr:subsub:boundary} later derives as $\rho=-(\alpha-\beta)$ all appear in
our journal three years before the present model, and the entry
is reproduced here in full because it states the mechanism more
compactly than the formal development does. Dated 23 July 2023,
1:42~pm, under the heading \emph{Golden Biosphere Analog to the Square
Cube Law}:
\begin{quote}\small
``Just thought of the craziest idea! You know the quote about how
advanced aliens might attack us like how boys throw stones at frogs? The
implication I get from this is that such an alien biosphere is portrayed
as just a scaled up version of our own so to speak in every aspect of
society, such that we just get a biosphere which is far more advanced
technologically, but interacts about itself in similar ways, such as
they are precisely are to humans as humans are to ants.

So, here's why I think this won't work:

The idea is basically an analog of the square cube law to
Astrobiology, working in a similar way to how you can't just scale an
ant up to the size of an elephant. There has to be a transformation in
the process instead.

The idea, essentially, is that the capacity for us to destroy ourselves
grows exponentially with technology, whereas the defense mechanisms put
in place with our current society, given the tragedy of the commons and
all, grows only linearly with technology. If we assigned each human
comparatively limitless power, such that the advanced analog of
`boys' might attack humans with the same mindset as human boys attack
frogs, then such an advanced civilization will implode as soon as it
spawns. Why? Because the linear defense mechanisms of a human society
would be immediately overtaken by the vast exponential issues propagated
by such a technology.

Anyways, so this demonstrates essentially two things: A) we will not
remain constant in societal mindset with time, and B) we will be
essentially forced through the keyhole of the Golden Point mindset long
before we reach the capacity of supposedly treating humans like frogs to
boys.''
\end{quote}
The entry must be read as its author wrote it, which is
\emph{conditionally on survivorship} --- the conditioning having been the
paradigmatic heart of the program since at least 2022, on our own
record, and therefore predating this entry as well. Read that way, every
one of its claims is recovered below; read unconditionally, two of them
would appear to fail. The distinction bears directly on how much the
entry anticipates, and read correctly the answer is all of it.

Six things are worth marking, since each became a separate part of the
apparatus. The exponential-against-linear contrast is
$\rho<0$: destruction capacity growing as $e^{\alpha x}$ against
enforcement growing more slowly is exactly
\eqref{gr:eq:rhoextract}, and the entry states the mechanism without the
surface-versus-volume derivation that \S\ref{gr:subsub:boundary}
supplies. The square--cube analogy is the source of the names
\emph{technostrength} and \emph{technoanatomy}, so
Remark~\ref{gr:rem:namesclosure}'s observation that the joke turned out
to be the mechanism is a retrospective reading of this entry rather than
a later discovery. ``Implode as soon as it spawns'' is
Theorem~\ref{gr1:thm:trapped}. And ``Golden Point'' is the earlier name
for what the present work calls the Golden Biosphere.

The entry's two numbered conclusions are the remaining pair, and both
are theorems here once the conditioning is respected. Item~(B) --- that
a biosphere is forced through the keyhole \emph{before} attaining the
capacity in question --- is
Proposition~\ref{gr1:prop:tolerance}: the affordable duration of adverse
anatomy is $\delta\lesssim\Gamma^{-1}e^{-|\rho|x}$, so conditional on
being alive at large technostrength, adverse anatomy is essentially
absent, and a biosphere cannot have arrived there without having already
passed through. The keyhole is not a force acting on biospheres --- 
unconditionally, Corollary~\ref{gr1:cor:recurrentanatomy} allows a
biosphere simply to fail --- but a property of the survivors, which is
what the entry claims. Item~(A) --- that societal arrangement will not
remain constant --- is the anatomy selection measured in
\S\ref{gr:subsub:anatsel}, where the conditioned distribution of $\rho$
shifts away from its prior. Neither is a prediction the model
contradicts; both are among the results it was later built to state.}
of the system: the
degree to which increasing muffling capacity fails to be turned back
against the system that acquired it. Reading the trough channel
$\Gamma e^{-\rho x}$ as internal hazard:
\begin{center}\small
\begin{tabular}{@{}lll@{}}
\toprule
& internal hazard as $x$ grows & reading\\
\midrule
$\rho>0$ & falls, $\propto e^{-\rho x}$ & capacity is patched as it is
built\\
$\rho=0$ & constant at $\Gamma$ & capacity neither patched nor
weaponized\\
$\rho<0$ & rises, $\propto e^{|\rho|x}$ & \emph{the system eats itself}\\
\bottomrule
\end{tabular}
\end{center}
\end{definition}

\begin{remark}[Why this is the right home for the notion]
\label{gr1:rem:anatomyhome}
Two things recommend $\rho$ over the phase coordinate $\Theta$ of
\S\ref{gr1:sub:coords}, with which it is easily conflated.
First, $\Theta$ is a \emph{state coordinate}: it fluctuates, is uniform
under $\PP$, and every system occupies every value of it half of
logarithmic time (Theorem~\ref{tc:thm:half}). Anatomy in the intended
sense is not something a system passes through but something it
\emph{is}. Second, and decisively, the intended notion is a claim about
what \emph{strength does} --- whether acquiring capacity raises or
lowers one's own hazard --- and that is a statement about a derivative,
$\partial_x\log h^{\mathrm{int}} = -\rho$, not about a position. The
phase coordinate is retained under its own name; technoanatomy is
$\rho$.
\end{remark}

\begin{example}[The asteroid, and why the channels must be two]
\label{gr1:ex:asteroid}
The reason internal and external hazard cannot be folded into one term
is best seen in a case where the same capability serves both. Consider a
biosphere that becomes progressively better at deflecting asteroids.
Every increment of that capability lowers the external hazard: the
crest channel falls, and by \eqref{tc:eq:hazard} it falls
multiplicatively in the background, exactly as it should.

But an apparatus that can redirect a body large enough to end a
biosphere is, by construction, an apparatus that can \emph{aim} such a
body. The same increment therefore also creates a capability that did
not previously exist and that is directed at the biosphere itself. What
happens to the total hazard depends on nothing about asteroids and
everything about whether the biosphere is so arranged that the new
capability can be turned inward --- which is precisely $\rho$.

With $\rho>0$ the deflection apparatus is built into an arrangement that
patches its own misuse faster than it creates the possibility of it, and
the net effect of the capability is protective. With $\rho<0$ it is not,
and the biosphere has purchased a reduction in external hazard at the
price of a larger increase in internal hazard --- the trade being worse
the more capable the apparatus. The example is not incidental to the
model: it is why \eqref{tc:eq:hazard} carries two channels rather than
one, and why the second is unbounded above while the first is not
(Proposition~\ref{tc:prop:bounded}). A single-channel model cannot
represent a capability that cuts both ways, and every capability worth
having does.

The dilemma is not this work's discovery, and the priority should be
recorded where the example lives: Sagan and Ostro stated it for this
very case --- deflection capability is aiming capability, so the
apparatus meant to remove the hazard becomes the larger hazard
\cite{SO94}. What the model adds is a functional form and a verdict:
the dilemma is the two-channel structure itself, its resolution is a
property of the biosphere's arrangement rather than of the apparatus,
and the two-condition theorem is what their dilemma looks like once
$\rho$ is a parameter rather than a worry.
\end{example}

\begin{remark}[On not measuring $\rho$ yet: a deliberate reservation]
\label{gr1:rem:noproxy}
The obvious next question is how one would measure $\rho$ for a real
system. No procedure is offered here, and the omission is deliberate
rather than a gap awaiting routine work. The reasoning is given here, partly because it is the framework's own reasoning applied to itself,
and partly so that a later proposal can be judged against it.

\emph{What is open is narrower than it may appear.} The exponent is
already fully defined: $\rho=-\partial_x\log h^{\mathrm{int}}$, the
operator gain of the internal channel \eqref{gr1:eq:channelops}. What is
absent is the \emph{identification map} --- which real configurations
carry which values --- and that is a different absence from a vague
concept. It is the difference between having no definition of
temperature and having no thermometer, and only the second is a thing
one might reasonably decline to supply in haste.

\emph{A premature proxy would be adverse anatomy applied to
measurement.} A scored index for technoanatomy would be a capability
acquired for the purpose of lowering hazard which, once acquired,
becomes the target of optimization and thereby raises it: institutions
would come to manage the index rather than the quantity, and the
quantity would be free to worsen behind it. That is Goodhart's law, and
in this framework's own vocabulary it is $\rho<0$ with the measurement
apparatus as the capability. Proposing an anatomy metric carelessly is
therefore not merely a methodological risk; it is an instance of the
failure the metric exists to detect.

\emph{Two structural features make the delay affordable.} First,
$\rho$ is a \emph{derivative and not a level} --- a rate of response of
self-inflicted hazard to acquired capability, not a stock. It is
correspondingly resistant to proxying by any static score, since to
game it one would have to control how hazard responds to capability,
which requires already possessing the property. Second, and more
usefully, \emph{the trichotomy of Proposition~\ref{gr1:prop:trichotomy}
turns on sign alone}. Trapped, critical and viable are separated by
whether $\rho$ is negative, zero or positive, with no magnitude
entering anywhere. An investigator may therefore work at sign level
today --- asking, of a particular capability in a particular system,
whether its acquisition raises or lowers self-inflicted hazard --- and
obtain the regime classification, which is most of what the model is
for, without committing to any index.

\emph{And the right formalization is downstream, not absent.} If $\rho$
is eventually obtained as in \S\ref{ad:sub:toward} --- computed from
institutional structure rather than scored onto it --- then there is no
index to manage, because the quantity is derived from what the system
\emph{is} rather than assigned to it by an assessor. That is the
formalization this reservation is holding the place for. Reaching for a
proxy in the interim would foreclose it, which is to say: it would
trade a route that advances for one that does not \textbf{[O]}.
\end{remark}

\begin{proposition}[The anatomy trichotomy; \textbf{P}]
\label{gr1:prop:trichotomy}
Write $h(x)=B\varsigma_b(x)+\Gamma e^{-\rho x}$ with $B,\Gamma>0$.
\begin{enumerate}[label=\emph{(\roman*)},itemsep=1pt]
\item If $\rho>0$, $h$ is strictly decreasing with $h(x)\to0$ as
$x\to+\infty$: capacity is unboundedly protective.
\item If $\rho=0$, $h$ is strictly decreasing with
$h(x)\to\Gamma>0$: capacity is protective but against a floor.
\item If $\rho<0$, $h\to B(1+e^{-b})$ as $x\to-\infty$ and
$h\to\infty$ as $x\to+\infty$, with a unique interior minimum. Writing
$A:=Be^{b}$ and $|\rho|=-\rho$, the minimum sits at
\begin{equation}\label{gr1:eq:xstar}
x_\ast=\frac{1}{1+|\rho|}\log\frac{A}{|\rho|\Gamma},
\qquad
h_\ast=A^{\frac{|\rho|}{1+|\rho|}}\,\Gamma^{\frac{1}{1+|\rho|}}
\Bigl(|\rho|^{\frac{1}{1+|\rho|}}+|\rho|^{-\frac{|\rho|}{1+|\rho|}}\Bigr),
\end{equation}
in the regime $x_\ast\gtrsim b$ where the crest is exponential.
\end{enumerate}
\end{proposition}

\begin{proof}
$\varsigma_b$ is strictly decreasing with limits $1+e^{-b}$ and $0$;
$e^{-\rho x}$ is decreasing for $\rho>0$, constant for $\rho=0$, and
increasing for $\rho<0$. (i) and (ii) are immediate. For (iii),
approximating the crest by $Ae^{-x}$ above the offset, $h'=0$ gives
$Ae^{-x}=|\rho|\Gamma e^{|\rho|x}$, whence $x_\ast$; substituting back
and collecting the common factor
$A^{|\rho|/(1+|\rho|)}\Gamma^{1/(1+|\rho|)}$ gives $h_\ast$. Uniqueness
follows since $h''>0$ wherever $h'=0$. (Verified numerically to three
significant figures at $b=3$, $k=6$, $\rho=-0.2$: predicted
$x_\ast=8.84$, $h_\ast=1.744\times10^{-2}$; measured $8.88$ and
$1.757\times10^{-2}$.)
\end{proof}

\begin{theorem}[Non-positive anatomy is fatal, and fast; \textbf{P}]
\label{gr1:thm:trapped}
Let $\rho\le0$. Then:

\emph{(i) Viability is zero, for the random background and pathwise.}
$\Lambda_\infty=\infty$ almost surely, so $\PP(\Lambda_\infty<\infty)=0$.

\emph{(ii) The tail is exponential at frozen background.} If $B_t\equiv
B>0$, then $h_\ast:=\inf_x h(x)>0$ and
\begin{equation}\label{gr1:eq:trappedtail}
\PP(T>t)\;\le\;e^{-h_\ast t}\qquad\text{for all }t,
\end{equation}
whence $\E[T]<\infty$. For the random background no such floor
exists --- $\inf_tB_t=0$ a.s.\ (Remark~\ref{gr1:rem:trappedscope})
--- and no exponential bound is claimed there; part (i) is what holds
pathwise, its recurrent-driver case routed through
Lemma~\ref{tc:lem:infbackground}.

\emph{(iii) In either case no power-law structure survives}: neither the
$t^{-1/4}$ tail, nor infinite life expectancy, nor Pareto residual life,
nor the Lindy reading of \S\ref{gr1:sub:verdict}.
\end{theorem}

\begin{proof}
\emph{(i)} If $x_t\to+\infty$ the trough $\Gamma e^{|\rho|x_t}\to\infty$;
if $x_t\to-\infty$ the crest tends to $B_t(1+e^{-b})$ and
$\int_0^\infty B_t\,dt=\infty$ a.s.\ by
Lemma~\ref{tc:lem:infbackground} with $S=[0,\infty)$ (the earlier
appeal to recurrence of the driving Brownian motions was informal).
On any other path at least one channel is bounded below along a set of
infinite measure, and the two halves are worth naming: either
$\mathrm{Leb}\{t:x_t\ge0\}=\infty$, where for $\rho\le0$ the trough
alone floors the hazard at $\Gamma e^{-\rho x_t}\ge\Gamma$ there ---
a deterministic floor on an infinite set --- or
$\mathrm{Leb}\{t:x_t\le0\}=\infty$, where the crest with
$S=\{t:x_t\le0\}$ routes through Lemma~\ref{tc:lem:infbackground}
exactly as in Theorem~\ref{tc:thm:noviab2}. In every case
$\Lambda_\infty=\infty$.

\emph{(ii)} For $\rho=0$, $h\ge\Gamma>0$ independently of the background.
For $\rho<0$ at frozen $B$, $h_\ast$ is the interior minimum of
Proposition~\ref{gr1:prop:trichotomy}(iii), strictly positive since both
channels are; then $\Lambda_t\ge h_\ast t$ and
$\PP(T>t)=\E[e^{-\Lambda_t}]\le e^{-h_\ast t}$.
\end{proof}

\begin{remark}[Why the exponential rate is not pathwise for the random
background]\label{gr1:rem:trappedscope}
The scope of \eqref{gr1:eq:trappedtail} deserves stating, since an
earlier form of this theorem claimed the bound held pathwise and used no
property of the driver.

For $\rho=0$ that is exact: the trough alone floors the hazard at
$\Gamma$, whatever the background does. For $\rho<0$ it is not, and the
reason is that $h_\ast$ is a function of $B$. The two-channel background
$B_t=\sum_i e^{\epsilon W^{(i)}_t}$ of \eqref{tc:eq:hazard} is random with
$\inf_t B_t=0$ almost surely, so there is no $B_0>0$ below which the
background never falls, and no single $h_\ast$ floors the hazard along
every path. What survives unconditionally is part~(i) --- the cumulative
hazard still diverges, since a vanishing crest requires $x_t\to-\infty$
which is exactly where the crest is at its largest --- so the
\emph{verdict} is unchanged and only the \emph{rate} is scoped.

Read correctly, then, \eqref{gr1:eq:trappedtail} is an annealed or
frozen-background statement, and the correct reading of $\E[T]<\infty$ is
the same. Nothing downstream consumes the rate: \S\ref{gr1:sub:cond}
consumes the trichotomy and the two-condition theorem, and
\S\ref{gs:sub:regimes}'s trapped row consumes the qualitative shape,
which part~(i) supplies. Whether an unconditional exponential rate holds
for the random background --- plausibly yes, at a smaller constant, by an
ergodic argument on $B$ --- is open \textbf{[O]}.
\end{remark}

\begin{remark}[What the trapped regime looks like]
\label{gr1:rem:trapped}
The reading of $\rho<0$ should be said plainly, because it inverts the
intuition the rest of the model builds. With negative technoanatomy the
\emph{safest available configuration is not the strongest but the
weakest}: $h\to B(1+e^{-b})$ as $x\to-\infty$, so a system with bad
anatomy is safest at background hazard with no capacity at all, and
every increment of capacity beyond $x_\ast$ is net harmful. The system
is not merely at risk; it is trapped between an external hazard it
cannot shed without capacity and an internal hazard that capacity
creates. Equation~\eqref{gr1:eq:xstar} prices the trap: the best
achievable hazard $h_\ast$ is a geometric interpolation between the
external scale $A$ and the internal scale $\Gamma$, weighted by anatomy.
\end{remark}


% ---------------------------------------------------------------------
\subsection{Transience, and the two-condition theorem}
\label{gr1:sub:twocond}

Anatomy governs whether capacity protects. It does not by itself deliver
survival, because protection must also be \emph{acquired}, and whether
it is depends on the driver. This subsection assembles the two into the
model's central structural result.

\subsubsection{The feedback condition}\label{gr1:subsub:feedback}

By Proposition~\ref{gr1:prop:trichotomy}, $\rho>0$ makes $h\to0$ as
$x\to+\infty$; but $\Lambda_\infty<\infty$ additionally requires
$x_t\to+\infty$ fast enough for $\int^\infty e^{-\min(1,\rho)x_t}dt$ to
converge, and $x_t=\int_0^t\mu_s\,ds$ grows without bound only if the
driver does not return. That is a transience condition on $\mu$, and
the boundary classification of one-dimensional diffusions settles it
exactly.

\begin{remark}[What the theorem computes]\label{gr1:rem:twocondreading}
By Remark~\ref{gr1:rem:whatTis}, $\PP(\Lambda_\infty<\infty)$ is the
probability that the biosphere \emph{advances} --- that some
continuation available to it reaches the Golden Point. The
two-condition theorem below therefore does not say what it takes to
avoid dying; it says what it takes for a route to exist at all, and the
two factors are the two things a route requires. That reading is what
makes the trapped regime intelligible: with $\rho<0$ the probability is
identically zero, so a biosphere there has departed in the strict
sense --- not because it has stopped, but because nothing it can do
advances.
\end{remark}

\begin{remark}[Transience is a feedback loop]\label{gr1:rem:feedback}
The canonical transient driver is
\begin{equation}\label{gr1:eq:uou}
d\mu_t=\nu\,\mu_t\,dt+\sigma\,dW_t,\qquad \nu>0,
\end{equation}
whose drift is proportional to the state: \emph{the rate of improvement
grows with the improvement already achieved}. This is the
self-reinforcing dynamic that the technological-singularity literature
posits, written as a diffusion, and it is not an addition to the model
but the member of the driver family whose scale function is bounded.
The connection made in \S\ref{gr1:subsub:twocondthm} is therefore not an
analogy: the condition under which the model admits immortality
\emph{is} the condition under which the driver exhibits feedback.
\end{remark}

\subsubsection{The theorem}\label{gr1:subsub:twocondthm}

\begin{theorem}[Two conditions for viability; \textbf{P}]
\label{gr1:thm:twocond}
Let the driver be a regular diffusion on $\R$ with scale function $s$,
and let $\rho\in\R$ be the technoanatomy. Then
\begin{equation}\label{gr1:eq:twocond}
\PP\bigl(\Lambda_\infty<\infty\bigr)
\;=\;
\underbrace{\one\{\rho>0\}}_{\text{anatomy}}
\;\cdot\;
\underbrace{\frac{s(\mu_0)-s(-\infty)}{s(+\infty)-s(-\infty)}}
_{\text{escape}},
\end{equation}
the second factor --- the driver's \emph{escape probability} --- read
as follows. When at least one scale limit is finite, it is the
displayed formula, with the conventions that $s(+\infty)=\infty$
(and $s(-\infty)$ finite) gives $0$, and $s$ bounded above with
$s(-\infty)=-\infty$ gives $1$. When \emph{both} scale limits are
infinite the driver is recurrent and the factor is not determined by
$s$ at all (Theorem~\ref{sh:thm:trichotomy}(iv)): it equals
$\one\{\bar m>0\}$ for positive recurrent drivers with integrable
stationary mean $\bar m$ (Proposition~\ref{sh:prop:recurrent}), it
equals $0$ for the Brownian driver (Theorem~\ref{sh:thm:noimm}), and
for null-recurrent drivers beyond Brownian it is open \textbf{[O]}.
Both factors are necessary: neither good anatomy without feedback, nor
feedback without good anatomy, admits viability.
\end{theorem}

% CORRECTION RECORD (referee pass, Aug 22 2026): the previous
% statement read the second factor as 0 WHENEVER s(+inf)=inf, which
% contradicted the appendix's own trichotomy(iv)/recurrent(a) and is
% false: an Ornstein--Uhlenbeck driver mean-reverting to a positive
% level has both scale limits infinite yet gives viability 1 (ergodic
% theorem: x_t/t -> mean > 0).  The counterexample is the referee's;
% the restatement above routes the both-infinite case through the
% appendix, which had it right all along.  Downstream wording
% ("transience factor" -> "escape factor") updated at every site.

\begin{remark}[The forward forcing: survival as asymptotically
decisive evidence; sketch at \textbf{[P]}, reading at the ladder's
grades]\label{gr1:rem:forwardforcing}
A configuration-level consequence of the theorem deserves its own
statement, because it is the forward-time mirror of the conditioning
argument's past-direction climb. Let the configuration, the sign
$\rho$ together with the driver class, be drawn from any prior that
gives positive mass to the viable set, the configurations for which
the theorem's product is positive. Non-viable configurations have
survival probabilities decaying to zero, exponentially in the trapped
regime and polynomially in the critical one, while viable
configurations hold a positive limit, so the conditional probability
of viability given survival to time $t$ tends to one: there is a soft
threshold in $t$, set by the slowest non-viable clock, past which a
surviving system is almost certainly in the configuration the
two-condition theorem describes. On the event of survival forever the
statement is no longer asymptotic but outright, since $T=\infty$
entails $\Lambda_\infty<\infty$ almost surely and hence the viable
configuration itself. In the dictation's words: each attribute
intersection of the biosphere surviving through proper time $t$ tilts
the disagreeing superposition of poles, the tilt becomes
asymptotically decisive, and under full conditioning on indefinite
survival the competence the Golden Point marks is not made likely but
implied, there being a threshold beyond which surviving with
non-Golden dangers of any kind still standing becomes arbitrarily
unlikely. This is the to-the-future extension of the past-direction
result that the same conditioning forces the evolutionary climb by
which the Golden Point is reachable at all. Two fences hold
unchanged. The mixture sketch is stated for the model's configuration
space, with the identification of the viable configuration as the
Golden Point's marker carried at the framework's usual grades and not
at \textbf{[P]}; and the epistemic ceiling governs exactly as before,
because this remark concerns the conditional law and not any
observer's epistemic state, so survival so far, at any finite time,
certifies nothing for the one surviving.
\end{remark}


\noindent A third factor joins these once the schedule channel of
\S\ref{sch:sub:schedule} is carried: \eqref{sch:eq:threecond} prefixes
the product with the probability that the schedule completes at all, which
attenuates where the anatomy indicator annihilates.

\noindent The theorem also scopes the Golden Tendency of
\S\ref{ac:sub:tendency}: the answer survivorship selects at
\S\ref{ac:ssub:shouldcan} purchases nothing toward the capacity this
theorem separates from it.

\begin{proof}
On $\{\mu_t\to+\infty\}$: eventually $\mu_t\ge1$, so $x_t\ge x_{t_0}+
(t-t_0)$ and, for $\rho>0$, $h(x_t)\lesssim e^{-\min(1,\rho)(t-t_0)}$,
whence $\Lambda_\infty<\infty$. For $\rho\le0$,
Theorem~\ref{gr1:thm:trapped}(i) gives $\Lambda_\infty=\infty$ almost
surely --- for the random background and regardless of the driver --- so
the first factor is necessary. On
$\{\mu_t\to-\infty\}$, $x_t\to-\infty$ and $h\to B(1+e^{-b})>0$ for
every $\rho$, so $\Lambda_\infty=\infty$; on the recurrent event the
outcome is governed by Proposition~\ref{sh:prop:recurrent} ---
$\bar m>0$ gives $D_t/t\to\bar m$ and $\Lambda_\infty<\infty$, while
$\bar m\le0$ (nondegenerate) gives $\Lambda_\infty=\infty$ by its
parts (b)--(c) --- with the Brownian driver settled by
Theorem~\ref{sh:thm:noimm} and general null recurrence left open,
which is the statement's scope clause. The hitting
probabilities of $\pm\infty$ are the classical scale-function formula
\cite{KT81,RY99}.
\end{proof}

\begin{corollary}[The solvable case; \textbf{P}]\label{gr1:cor:solvable}
For the feedback driver \eqref{gr1:eq:uou},
$s'(v)=\exp(-\nu v^{2}/\sigma^{2})$ is integrable at both ends, and
\begin{equation}\label{gr1:eq:sigmoid}
\PP\bigl(\Lambda_\infty<\infty\bigr)
=\one\{\rho>0\}\cdot\Phi\!\Bigl(\frac{\mu_0\sqrt{2\nu}}{\sigma}\Bigr).
\end{equation}
The viability probability is thus a \emph{product of two thresholds of
different character}: a sharp one in the technoanatomy, at $\rho=0$,
which is a discontinuity; and a smooth Gaussian sigmoid in the initial
improvement rate, of width $\sigma/\sqrt{2\nu}$, which sharpens to a
step only in the small-noise limit.
\end{corollary}

\begin{remark}[Against the singularity literature]
\label{gr1:rem:againstsing}
Equation \eqref{gr1:eq:twocond} states, in the model's own terms, an
objection to the singularity and techno-optimist literatures which is
sharper than the usual one. Those literatures supply
the second factor and are silent on the first: feedback is argued for at
length, and whether capacity is turned back on its holder is not a
parameter they carry. The model's verdict on that omission is not that
it leaves the conclusion unsupported but that it inverts it. With
$\rho<0$, feedback does not merely fail to secure survival; it
\emph{accelerates the failure}, since $\mu_t\to+\infty$ drives
$x_t\to+\infty$ and hence, by
Proposition~\ref{gr1:prop:trichotomy}(iii), $h\to\infty$
superexponentially. Supplying feedback alone is strictly worse than
supplying neither: a system with bad anatomy and no feedback sits near
$h_\ast$, while one with bad anatomy and feedback runs to unbounded
hazard as fast as its improvement compounds.
\end{remark}

\subsubsection{If anatomy itself fluctuates}\label{gr1:subsub:fluct}

Technoanatomy has been treated as a constant, which is an idealization
and the model's most obvious place to give. Suppose instead
$\rho=\rho_t$ is a process. Two consequences follow at scaling grade,
and together they sharpen rather than blunt the event-horizon reading.

\begin{proposition}[The tolerance for adverse anatomy collapses with
strength; \textbf{S}]\label{gr1:prop:tolerance}
Let anatomy dip to $\rho_t=-|\rho|<0$ for a duration $\delta$ at a time
when technostrength is $x$. The contribution to the cumulative hazard
is $\asymp\delta\,\Gamma e^{|\rho|x}$, so that such an excursion is
survivable only if
\begin{equation}\label{gr1:eq:tolerance}
\delta\;\lesssim\;\Gamma^{-1}e^{-|\rho|x}.
\end{equation}
The affordable duration of adverse anatomy therefore decays
\emph{exponentially in technostrength}.
\end{proposition}

\begin{corollary}[Recurrent anatomy is fatal; \textbf{S}]
\label{gr1:cor:recurrentanatomy}
If $\rho_t$ is recurrent on $\R$ --- if anatomy returns to negative
values infinitely often --- then $\Lambda_\infty=\infty$ almost surely,
\emph{whatever the driver}. Each negative excursion occurring at
technostrength $x_t$ contributes $e^{|\rho|x_t}$, and along any path on
which strength accumulates this diverges superexponentially. Viability
therefore requires $\rho_t$ to be eventually positive almost surely ---
transient upward, or absorbed above some $\rho_c>0$ --- and the
two-condition theorem becomes a \emph{two-transience} condition, with
$\PP(\text{viable})$ a product of two scale functions rather than an
indicator times one.
\end{corollary}

\begin{remark}[This is the event horizon the fixed model could not
state]\label{gr1:rem:horizonproper}
Three readings, and the third is the one to keep.

\emph{The threshold becomes an escape rather than a level.} With $\rho$
fixed, the anatomy condition is a sign test one either passes or fails
from the outset. With $\rho$ fluctuating it is an event --- anatomy
escaping upward and not returning --- which is what an event horizon
is: not a place in the state space but a passage after which the past
is unreachable.

\emph{``Hard to regress past'' becomes quantitative.} Regression is
never impossible; what \eqref{gr1:eq:tolerance} says is that the
affordable \emph{duration} of it collapses exponentially in strength. A
biosphere may be badly configured for a long while when weak and for
almost no time at all when strong.

\emph{An ordering constraint falls out.} Since the cost of anatomical
failure grows with strength while the benefit of capacity does not
grow with anatomy, anatomy must be secured \emph{before} strength
accumulates. A biosphere that acquires capacity first and attends to
its own arrangement afterwards is running the two in the fatal order,
and \eqref{gr1:eq:tolerance} prices exactly how fatal: by the time it
turns to the problem, its margin for error has shrunk by
$e^{|\rho|x}$.
\end{remark}

\noindent The dynamics of $\rho_t$ are not supplied here, and supplying
them is the model's principal open direction: an anatomy that degrades
when capacity outruns the patching of it, $d\rho=(\text{patching})-
\lambda\,\dot x\,dt$ or similar, would make technoanatomy emergent
rather than parametric. Whether the asymptotics of the resulting
two-process system remain tractable is \textbf{[O]}.

\begin{remark}[Where the earlier results now sit]\label{gr1:rem:wherenow}
The trichotomy locates the rest of this section. The driftless model of
\S\S\ref{gr1:sub:verdict}--\ref{gr1:sub:cond} --- certain departure,
infinite life expectancy, the $t^{-1/4}$ tail, the soft horizon and the
whole survivorship interpolation --- is the case $\rho>0$ with a
\emph{recurrent} driver: good anatomy, no feedback. It is the boundary
between the trapped regime, where viability is zero and life expectancy
finite at frozen background (Theorem~\ref{gr1:thm:trapped}), and the
viable regime, where
immortality has positive probability
(Corollary~\ref{gr1:cor:solvable}). That is the same structural position
the driftless case occupies within the scale-function classification of
Appendix~\ref{app:horizon}, and it is not a coincidence: both
trichotomies are the same statement resolved along different axes ---
one along anatomy, one along the driver.
\end{remark}


% ---------------------------------------------------------------------
\subsubsection{The internal channel as a perturbation}
\label{gr1:subsub:perturb}

The two channels admit a perturbative treatment in the internal
coupling $\Gamma$, and it is worth recording, both because it supplies
the model's one series expansion and because the series behaves in an
instructive way at the anatomy threshold.

\begin{proposition}[First-order internal correction; \textbf{P} to
first order, \textbf{S} for the convergence claim]
\label{gr1:prop:perturb}
Write $\Lambda_t=\Lambda^{\mathrm{crest}}_t+\Gamma\int_0^te^{-\rho
x_s}ds$. Expanding the survival functional in $\Gamma$,
\begin{equation}\label{gr1:eq:perturb}
\PP(T>t)\;=\;\PP_{\mathrm{crest}}(T>t)\left[\,1-\Gamma\,
\E_{Q_{\mathrm{crest}}}\!\Bigl[\int_0^t e^{-\rho x_s}\,ds\Bigr]
+O(\Gamma^{2})\,\right],
\end{equation}
where $Q_{\mathrm{crest}}$ is the crest-only survivorship measure of
\S\ref{gr1:sub:cond}. The leading correction is therefore $\Gamma$ times
the \emph{survivor ensemble's} expected exposure to the internal
channel, not the unconditioned ensemble's --- the conditioning enters
the correction rather than sitting outside it.
\end{proposition}

\begin{proof}[Proof of the first-order term]
$e^{-\Lambda}=e^{-\Lambda^{\mathrm{crest}}}(1-\Gamma\int
e^{-\rho x}+O(\Gamma^{2}))$ pointwise; take expectations and divide by
$\E[e^{-\Lambda^{\mathrm{crest}}}]$, which converts the expectation of
the bracket into a $Q_{\mathrm{crest}}$-expectation by the definition of
that measure.
\end{proof}

\begin{remark}[The expansion fails exactly at the anatomy threshold]
\label{gr1:rem:perturbfail}
Under $Q_{\mathrm{crest}}$ the credit grows as $x_s\asymp\sigma
s^{3/2}$, so the correction term behaves as
$\int^\infty e^{-\rho\sigma s^{3/2}}ds$, which converges if and only if
$\rho>0$. Numerically at $\sigma=1$: the integral over $[0,60]$ is
$0.90$ at $\rho=1$ and $6.7$ at $\rho=0.05$, but $2.2\times10^{10}$ at
$\rho=-0.05$ and $1.0\times10^{60}$ at $\rho=-0.3$, with essentially all
the mass in the far tail once $\rho<0$.

The reading is given because it is a check on the model rather
than a decoration. A perturbation series whose radius of convergence
collapses at a parameter value is the standard signature of a genuine
phase transition there, as opposed to a modeling artifact or a
definitional boundary: the expansion fails because the object it is
expanding around ceases to exist, not because the algebra becomes
awkward. That $\Gamma$-perturbation theory breaks down precisely at
$\rho=0$ is therefore independent evidence that the trichotomy of
\S\ref{gr1:sub:anatomy} is a real division of the model and not a
convenient way of describing it. It also explains why no uniform
treatment of the two channels is available: for $\rho\le0$ the internal
channel is not a perturbation of anything, being the dominant term at
every large $x$.
\end{remark}

\begin{remark}[What is available in closed form, and what is not]
\label{gr1:rem:closedform}
A brief inventory, since the model is now large enough that the question
arises. \emph{Exact:} the viability product \eqref{gr1:eq:twocond}; the
Gaussian sigmoid \eqref{gr1:eq:sigmoid}; the trapped-regime optimum
$x_\ast$ and floor $h_\ast$ of \eqref{gr1:eq:xstar}; the gain minimum
$\kappa(\rho)e^{-(\rho b+k)/(1+\rho)}$; the first-order term of
\eqref{gr1:eq:perturb}. \emph{Exact as an exponent, unknown as a
constant:} the survival tail $Ct^{-1/4}$, whose prefactor involves the
Groeneboom--Jongbloed--Wellner harmonic function and is not evaluated
here. \emph{Series:} the bias field expansion
$(1-s/t)^{-1/4}=1+s/4t+O(s^2/t^2)$ of \eqref{gr:eq:rate}, in the
interpolation parameter $s/t$; and \eqref{gr1:eq:perturb}, in $\Gamma$.
\emph{Variational rather than closed:} the adverse-drift rate
$3\nu^{2}/8\sigma^{2}$ and its optimal profile, obtained from the
Cameron--Martin action by Cauchy--Schwarz --- the model's analog of an
Euler--Lagrange computation, and the natural processing step whenever a
conditioned path law is wanted at large deviation scale.
\emph{Not available:} the trapped regime's survivor tail beyond the
exponential upper bound, and the asymptotics of any model in which
$\rho$ carries its own dynamics.
\end{remark}


% ---------------------------------------------------------------------
% ---------------------------------------------------------------------
\subsection{The third channel: schedule}\label{sch:sub:schedule}

\eqref{tc:eq:hazard} carries two channels, and
Remark~\ref{gc:rem:hazardlower} records that they do not exhaust the
hazard of departure: $h^{\mathrm{tot}}$ prices $\mathrm{Base}$ and
nothing else, while the hazard of departure is the hazard of $\mu G$,
which by \eqref{gc:eq:departurefix} is at least as large and in general
strictly larger. That remark calls the larger rate open \textbf{[O]}.

This subsection is a first attempt at it. The claim is modest: the rate
has a computable shape, the shape is a third channel, and the
distribution it needs is one the companion already supplies.

\subsubsection{A collision of vocabulary, settled first}
\label{sch:subsub:vocab}

\begin{convention}[Endogenous, internal, schedule]\label{sch:conv:internal}
The companion divides departure into \emph{external} --- inflicted from
outside the biosphere --- and \emph{internal}, defined there as the
complement: the remaining cases, all occurring within the biosphere's
scope \cite{Nexus}. That category is a residue and is correspondingly
wide; the companion's own instance of it is a biosphere too
underdeveloped to reach a milestone given its resources and timing,
which is a scheduling failure and not a self-inflicted one.

The internal channel of \eqref{tc:eq:hazard} is narrower: hazard the
biosphere inflicts on itself with capacity it has acquired, whose
response to that capacity is the technoanatomy $\rho$ of
Definition~\ref{gr1:def:anatomy}.

The two senses are a category and one of its members. This work uses
\emph{endogenous} for the companion's residue and reserves
\emph{internal} for the channel:
\begin{equation}\label{sch:eq:partition}
\text{endogenous}\;=\;\underbrace{\text{internal}}_{\text{capacity
turned inward}}\;+\;\underbrace{\text{schedule}}_{\text{derailment
before capacity can be held}}\;+\;\cdots
\end{equation}
The ellipsis is substantive. A category defined by exclusion is not
shown exhausted by exhibiting two members, and nothing here claims these
are the last.\footnote{The rename is a convention of this work and is
offered rather than imposed; \emph{non-external} and \emph{endemic}
would serve. One point of order regardless of the choice: the companion
is deposited, so a reader arriving from the deposit will meet
\emph{internal doomsday} in the residue sense, and the convention is
owed to bridge the two states of that text whether or not the term is
changed there.}
\end{convention}

\subsubsection{The schedule}\label{sch:subsub:model}

\begin{definition}[Schedule; \textbf{[S]}]\label{sch:def:schedule}
A \emph{schedule} is a sequence of $n$ gates to be passed in order. At
gate $k$ there are $m$ competing exponential clocks with rates
$\lambda_{k,1},\dots,\lambda_{k,m}$, the first corresponding to the
realized option. Write
\begin{equation}\label{sch:eq:clocks}
\lambda_k=\sum_{j=1}^{m}\lambda_{k,j},
\qquad
\pi_k=\frac{\lambda_{k,1}}{\lambda_k},
\end{equation}
so that gate $k$ fires at rate $\lambda_k$ and resolves to the realized
option with probability $\pi_k$. Write $k_t$ for the number of gates
passed by time $t$. The \emph{onset state} is $k=n$: the first state at
which the muffling exponent $x$ of \eqref{tc:eq:x} can be held at all.
\end{definition}

The rates are the primitives and $\lambda_k,\pi_k$ are read off them,
which matters because it is the per-option rates a natural-historical
account would estimate --- not a tempo and a fallibility chosen
separately.

Stage dependence is retained deliberately. The reason the object is a
\emph{sequence} rather than a set is that each event's effect depends on
the state its predecessor left: the climatic and ecological
baton-passing that carried a lineage down from the trees is not a set of
independent draws, and a constant $\pi$ would model it as one.
Stage-dependent $\pi_k$ is the minimum that preserves the structure;
state-dependent $\pi(\cdot)$ is what the picture actually asks for, and
is not supplied here \textbf{[O]}.

One point of scope, since it is a natural worry and the construction
already absorbs it. An event later in the sequence arriving before its
predecessors needs no transition of its own: it is one of the $m$
options competing at the gate it arrives at, and resolving to it is a
wrong resolution like any other. The $m-1$ non-realized options are not
a homogeneous class --- each carries its own rate by
Definition~\ref{sch:def:rates} and its own consequence --- so
heterogeneity of failure is present without being labeled by type. What
the chain does not carry is partial progress within a gate: it is Markov
between them, which is a simplification and is chosen rather than
derived.

The discreteness is a limitation of the same kind as calling Earth a
sphere. Exotic structure in the ordering is not recoverable from a chain
over $m^n$ articulable realizations; the shape of the dependence is, and
that is what is wanted here.

\subsubsection{Where the rates come from}\label{sch:subsub:rates}

Definition~\ref{sch:def:schedule} takes the per-option rates as given.
One specification of how they are drawn is worth carrying, because it
determines more than it looks like it should.

\begin{definition}[Log-spread rates; \textbf{[S]}]\label{sch:def:rates}
At each gate the $m$ rates are drawn afresh and independently with
$\log\lambda_{k,j}$ \emph{logistic} --- linear across the bulk with both
tails capped, since a shortest feasible time bounds the fast end and a
longest tolerable one bounds the slow --- and \emph{which option is the realized one is
independent of its rate}. The correct option may be the fastest of the
$m$, or the slowest, and nothing about the draw favors either. The
lognormal reading is the default and carries its own justification: a
log-rate that is a sum of many independent multiplicative contributions
is normal by the central limit theorem, so the law is inherited rather
than chosen, and the capped variants are refinements for the extreme
tails.
\end{definition}

Three consequences, in ascending order of how much they change.

\emph{The waiting time is set by the correct option's own rate.} A gate
fires at $\lambda_k=\sum_j\lambda_{k,j}$ and resolves correctly with
probability $\pi_k=\lambda_{k,1}/\lambda_k$, so the expected number of
attempts before a correct resolution is $1/\pi_k$ and each costs
$1/\lambda_k$. The product is $1/\lambda_{k,1}$: however many wrong
options crowd the gate, what a biosphere waits on is the realized
option's own clock. Exposure per gate is therefore log-spread along with
the rate, and the difference between a lucky gate and an unlucky one is
orders of magnitude of waiting rather than a factor.

\emph{The contrast between options is a random order of magnitude, with
the sign withheld.} $\log(\lambda_{k,1}/\lambda_{k,j})$ is a difference
of two log-uniform draws: symmetric about zero, spread across the range,
and committing to no direction. That is the shape of the High-Contrast
Earth Hypothesis, and the correspondence is exact rather than
suggestive --- the Hypothesis asserts a random order of magnitude between
independently constructed alternatives and \emph{deliberately withholds
the sign}, with the Strong Accordance Principle supplying direction
separately. \eqref{sch:def:rates} supplies the same shape from a
mechanical premise about competing clocks, and the direction arrives the
same way it does there: from conditioning, below. So the Hypothesis
acquires a microfoundation at the gate level, and does not have to be
imported as a posit at this scale --- though what replaces it is another
\textbf{[S]}, and the trade is one posit for a more checkable one rather
than a derivation from nothing.

The correspondence should be stated at the right strength, since it is
not an identity of distributions. By \eqref{sch:eq:cefff} the contrast a
perturbation carries is $r\log_{10}(1+(B+\Gamma)/\lambda)$, which in the
regime supplying many orders (Remark~\ref{sch:rem:tension}) is $r$ times
a \emph{log-rate gap} --- and under Definition~\ref{sch:def:rates} a
log-rate gap is a difference of normals, hence normal, hence symmetric
and spread across orders with no committed sign. The companion states
the Hypothesis as uniform rather than normal. What the collapse
consumes, however, is not the shape but two functionals of it --- the
typical contrast $\E[c]$ and the averaged one
$-\log_{10}\E[10^{-c}]$ --- and both are defined for either law. So the
construction is compatible with the Hypothesis in the sense the argument
uses it, without reproducing its stated distribution.

The scale of the spread is a further matter, and the accurate thing to say is not that it awaits measurement. Nothing available to an observer
inside a biosphere fixes what \emph{sets} the range over which competing
rates are distributed: one meets the distribution already running, with
no vantage from which its parameters could be explained rather than
merely estimated. That is \S\ref{gr:sub:ceiling}'s structure applied one
level down, and the scale is scoped as uncertifiable for the reason
$\rho$ is, rather than listed as a parameter someone has yet to
supply.

\emph{And conditioning selects the fast-correct branch, which is why the
realized path looks optimal.} Exposure at gate $k$ is $1/\lambda_{k,1}$,
paid at hazard $B+\Gamma$ by Proposition~\ref{sch:prop:gating}, so a
branch's survival weight carries $\exp\{-(B+\Gamma)/\lambda_{k,1}\}$ at
every gate. Branches in which the realized option happened to be slow
pay exponentially for it. Conditioning on survival therefore selects,
gate by gate, for histories in which the correct option was among the
fastest available --- not because anything chose them, but because the
alternatives were charged for the wait.

That is the mechanism, and should be stated as one. Under the Cosmic Nexus the
realized option is a random draw and is usually slow; under conditioning
one is overwhelmingly likely to find oneself on a branch where it was
fast at gate after gate. The appearance of an optimal path threaded
through a combinatorial mess is what survivorship does to a log-spread
rate landscape, and requires no optimization to produce it. This is the
Optimal-Earth phenomenology arriving from the schedule's own machinery
rather than being assumed alongside it \cite{Nexus, Ben26Optimal}.

\subsubsection{Incompleteness decays capacity}\label{sch:subsub:leak}

The first commitment is that until the schedule completes, the muffling
exponent cannot be held. This departs from
Definition~\ref{tc:def:model}, which states that there is no capacity
parameter and no leak; the departure is deliberate and is confined.

\begin{definition}[Staged leak; \textbf{[S]}]\label{sch:def:leak}
Replace \eqref{tc:eq:x} and \eqref{tc:eq:driver} by
\begin{equation}\label{sch:eq:leak}
dx_t=\mu_t\,dt-\gamma_{k_t}x_t\,dt,
\qquad
d\mu_t=-\nu_{k_t}\mu_t\,dt+\sigma\,dW_t,
\end{equation}
with $\gamma_k,\nu_k$ nonincreasing in $k$ and $\gamma_n=\nu_n=0$. The
schedule of coefficients is \emph{geometric}: each closed gate
multiplies the leak by a constant factor,
\begin{equation}\label{sch:eq:geomleak}
\gamma_k=\gamma_0\,\varphi^{\,k},\qquad
\varphi=(\gamma_{n-1}/\gamma_0)^{1/(n-1)}<1,
\end{equation}
with $\gamma_0$ enormous --- so that before any gate is closed the credit
is pinned at the origin to a degree with no practical upper bound ---
and $\gamma_n=0$ at completion.
\end{definition}

The solution of the first equation is
$x_t=\int_0^te^{-\gamma_k(t-s)}\mu_s\,ds$: the credit is still the
integral of the driver, accumulated additively, with each increment
buffered by a multiplicative decay from the moment it lands. The leak is
not an alternative to \eqref{tc:eq:x} but a discounting of it.

The second equation is the same commitment one derivative up, and it is
what the picture requires rather than an extra assumption. The driver
$\mu=-\tfrac{d}{dt}\log h$ is an \emph{improvement rate}, and a rate of
improvement is something an agent supplies. Before the onset state there
is none, so $\mu$ is not a free random walk either: neither the stock nor
the rate is sustained, and the two dampings are one absence expressed at
two levels. A biosphere's ancestors do not hold a small amount of
technostrength --- they hold none, and what the pair of equations says is
that excursions away from zero are transient rather than accumulated.

At stage $k$ the deterministic part obeys $\tfrac{d}{dt}\log x=-\gamma_k$:
capacity acquired before completion bleeds away at a rate each completed
gate reduces. The leak is by hand, and making it emergent --- $\gamma$
falling out of the schedule's own structure rather than being assigned
to it --- is the model's natural next target \textbf{[O]}.

\begin{remark}[What a closed gate buys, and why almost none of it is
bought early]\label{sch:rem:backloaded}
\eqref{sch:eq:geomleak} makes gates equal in $\log\gamma$ rather than in
$\gamma$, and the consequence for the payoff is severe enough to state
on its own.

The credit a biosphere can sustain at stage $k$ scales as
$1/\gamma_k=\varphi^{-k}/\gamma_0$, so it grows \emph{geometrically} in
gates closed. But the hazard depends on the credit \emph{exponentially}
--- both channels of \eqref{tc:eq:hazard} do, by construction --- so the
relief a closed gate delivers goes as
$\exp\{-\rho C\varphi^{-k}/\gamma_0\}$, doubly exponential in $k$.
Closing the last gate is worth more than closing all the others
together, by a margin that is not close.

With $\gamma_0$ enormous the reading sharpens. The early gates take a
biosphere from a credit that is unimaginably pinned to one merely very
pinned, which is no relief at all in the only currency the hazard reads.
Practical relief arrives in the final stretch, and a schedule
half-climbed is nearer in hazard to one not begun than to one completed.

Set this beside Remark~\ref{hr:rem:monotone}, which has stakes rising
monotonically for combinatorial reasons --- a late step carries the
improbability of every step behind it. The present observation is a
second and independent reason for the same monotonicity, arriving from
the leak's geometry rather than the chain's combinatorics. Two arguments
of different provenance agreeing on the profile is a better position
than either alone.
\end{remark}

\begin{remark}[The existing model is the completed case]
\label{sch:rem:limit}
$\gamma_n=0$ is not a convenience. It makes
Definition~\ref{tc:def:model} the $k=n$ boundary case of
\eqref{sch:eq:leak} exactly, so nothing proved in
\S\S\ref{gr1:sub:verdict}--\ref{gr1:sub:cond} is disturbed: the
persistence exponent $\tfrac14$, the trichotomy, the two-condition
theorem and the survivorship interpolation are all statements about the
completed regime, and this channel prefixes them rather than amending
them.

The price is that the incomplete regime has its own asymptotics --- and
they can be given, which is the next proposition's business. The
persistence problem is not the one \S\ref{sh:subsub:tail} solves, and
its exponent does not carry.
\end{remark}

\begin{proposition}[The incomplete regime has no polynomial tail;
\textbf{[S]}]\label{sch:prop:exponent}
Under \eqref{sch:eq:leak} with $\gamma_k,\nu_k>0$ held fixed, the pair
$(\mu,x)$ is a two-dimensional linear system with both eigenvalues
negative, driven by white noise; it is therefore ergodic with a
stationary Gaussian law centered at the origin. Persistence of $x$ is then
persistence of a stationary Gaussian process, which decays
\emph{exponentially}:
\begin{equation}\label{sch:eq:incompletetail}
\PP\bigl(x\ \text{stays positive to }t\bigr)\;\asymp\;e^{-\theta(\gamma_k,\nu_k)\,t},
\qquad \theta>0 .
\end{equation}
No power law survives. The rate $\theta$ is not given in closed form
here; persistence rates for stationary Gaussian processes are generally
not, and \eqref{sch:eq:incompletetail} identifies the class rather than
the constant.\footnote{Simulation at $\sigma=1$, $8\times10^{4}$ paths.
With $\gamma=1$ and the driver \emph{undamped} ($\nu=0$) the decay is a
power law --- fitted slope $-0.476$, log--log correlation $0.9995$ ---
which is the damped randomly accelerated particle and is not the regime
this definition describes. With the driver damped, the fit inverts
decisively: at $(\gamma,\nu)=(1,0.5)$ and $(1,1)$ the correlations
against $t$ are $1.0000$ and $0.9999$ with rates $0.222$ and $0.312$,
against log--log correlations of $0.975$ and $0.985$.}
\end{proposition}

\begin{remark}[What incompleteness costs, and what completion buys]
\label{sch:rem:suspends}
The contrast with the completed regime is not a change of exponent but a
change of kind, and it is the sharpest thing the channel says about why
the schedule is worth climbing.

\S\ref{gr1:sub:verdict} establishes that the completed model departs with
probability one and yet has infinite life expectancy: the tail is
$t^{-1/4}$, residual life is Pareto, quantiles grow as the fourth power
of rarity, and the effective hazard is $1/4t$ --- falling at every age.
By \eqref{sch:eq:incompletetail} the incomplete regime has none of it.
An exponential tail has finite expectation, memoryless residual life,
quantiles growing logarithmically, and a hazard that does not fall at
all. \emph{Every consequence of Proposition~\ref{sh:prop:tailconseq} is a
consequence of having completed}, and none of it is available before.

Two readings follow, and the second is structural.

The Lindy structure is a \emph{reward} rather than a property. A
biosphere before its onset state is not on a fat tail at low altitude;
it is on a thin one, and no amount of waiting improves its position,
since an exponential is memoryless. What completing the schedule buys is
not more of the same but a different tail entirely.

And the floor is the trapped regime's floor with one difference that
matters. Theorem~\ref{gr1:thm:trapped}(ii) gives $\PP(T>t)\le e^{-h_\ast t}$
for adverse anatomy at frozen background, and \eqref{sch:eq:incompletetail}
gives an
exponential floor for incompleteness. The two are structurally alike and
disposed of oppositely: the anatomy floor cannot be left, since $\rho$ is
a parameter, while the schedule floor is \emph{closable} --- $\gamma_k$
and $\nu_k$ vanish at completion, and the debt that stands between is
payable whenever the criterion of \S\ref{sch:subsub:open} holds. An
incomplete biosphere and a trapped one sit on the same kind of tail; only
one of them can get off it.

A biosphere may of course be both, and the two floors then compose
exactly rather than competing. If $h^{\mathrm{tot}}\ge h_\ast$ then
$\PP(T>t)=e^{-h_\ast t}\,\E\bigl[e^{-\int(h^{\mathrm{tot}}-h_\ast)}\bigr]$
identically --- a constant floor factors out of the exponential
functional --- leaving the damped-persistence problem for the recentered
hazard with its own rate $\tilde\theta$. So
$\theta_{\mathrm{joint}}=h_\ast+\tilde\theta$: neither floor governs, they
\emph{add}, and the interventions are correspondingly separable.
Completing the schedule removes $\tilde\theta$ and not $h_\ast$; moving
$\rho$ removes $h_\ast$ and not $\tilde\theta$. A biosphere that is both
trapped and incomplete must do both things, and doing either buys exactly
its own term.
\end{remark}

\subsubsection{The schedule gates the external channel}
\label{sch:subsub:gating}

Completing the chain confers no capacity. It confers only the ability to
\emph{hold} capacity --- which sounds like a small thing and is not,
because of what \eqref{tc:eq:hazard} does at $x=0$.

\begin{proposition}[An unheld $x$ leaves the crest at bare background;
\textbf{[P]}]\label{sch:prop:gating}
By the normalization of \eqref{tc:eq:hazard}, $\varsigma_b(0)=1$
exactly, so $h^{\mathrm{crest}}\big|_{x=0}=B_t$; and the trough does not
vanish there either, since $\Gamma e^{-\rho\cdot0}=\Gamma$. A biosphere
that cannot hold $x$ therefore faces the external background unshielded
\emph{and} pays the internal channel's full coefficient, at total
$h^{\mathrm{tot}}\big|_{x=0}=B_t+\Gamma$. Its cumulative hazard is
$\Lambda_t=\int_0^t(B_u+\Gamma)\,du\to\infty$, and
$\PP(\Lambda_\infty<\infty)=0$ identically --- for every parameter
value, every anatomy $\rho$, and every driver. The configuration and its
consequences are those of Proposition~\ref{hr:prop:extinction}.
\end{proposition}

\begin{remark}[The external term is small only because the schedule
succeeded]\label{sch:rem:reliance}
Read at $x>0$ the crest channel is the least of the model's worries: it
is bounded above by Proposition~\ref{tc:prop:bounded}, it is shielded
without limit as capacity accumulates, and beside an unbounded internal
channel it looks like one measly term. That reading presupposes what the
schedule supplies.

Run the dependency out. Shielding is the factor $\varsigma_b(x)$;
$\varsigma_b$ shields only for $x>0$; $x$ can be held only at $k=n$; and
$k=n$ is reached through $n$ gates each of which can go $m-1$ ways
wrong. \emph{So the survivability of the external channel is downstream
of the schedule entire.} The cosmic background does not become more
dangerous when the schedule fails; it becomes the only thing that
happens, which is the sentence Remark~\ref{hr:rem:undisturbed} writes
about a biosphere whose muffling agent has been removed.

That is not a coincidence of phrasing. \emph{An incomplete schedule and
the extinctionist's withdrawn world are the same configuration.} Both
are $x$-cannot-be-held; both sit at bare $B_t$; both are
Proposition~\ref{hr:prop:extinction}'s case; and both lie in
$\mu G\setminus\mathrm{Base}$ by \S\ref{gc:subsub:cascade} --- departed
by configuration rather than by event, with no hazard yet fired. The
channel therefore does not merely add a term. It exhibits the model's
apparently smallest channel as conditional on its largest improbability.
\end{remark}

\subsubsection{Failure incurs debt, not weight}\label{sch:subsub:debt}

What a wrong gate costs can now be said, and the natural first answer is
the wrong one. Treating a wrong gate as multiplying an onset weight
makes that weight a primitive --- a number the model is handed rather
than one it produces --- and leaves the cascade of
\S\ref{gc:subsub:cascade} inexpressible, since a single multiplicative
factor has nowhere to put the recursion by which an Imperfection
Magnitude Drop's residual carries continuations of its own.

\begin{definition}[Debt; \textbf{[S]}]\label{sch:def:debt}
A wrong gate does not suppress a weight. It \emph{incurs debt}: $r$
further gate-successes are owed, from the state now occupied, before
completion is reached. Write $Z_t$ for the outstanding debt, with
$Z_0=n$. Each firing consumes one gate and, with probability $1-\pi$,
adds $r$:
\begin{equation}\label{sch:eq:debt}
Z\;\longmapsto\;
\begin{cases}
Z-1 & \text{with probability }\pi,\\
Z-1+r & \text{with probability }1-\pi .
\end{cases}
\end{equation}
\emph{Completion} is the event $Z=0$.
\end{definition}

The natural-historical reading is the one that motivates it. Human
extinction does not leave the biosphere holding a fractional prospect;
it leaves it owing a further sequence --- some other lineage must run
its own improbable gymnastics from wherever the world now stands --- and
the onset state recedes by exactly that sequence's length. Debt
postpones a deadline. It does not discount a prize.

$r$ need not be a primitive, and the model is better if it is not. Read
the schedule as a tree whose paths each carry a definite number of
events, differing between paths: a wrong turn then places the biosphere
on a longer path, and \emph{the excess length is the debt}. On that
reading \eqref{sch:eq:debt} is the induced dynamics on \emph{events
remaining} rather than a further posit, $r$ is a structural property of
the tree rather than a parameter, and whether it correlates with the
rates --- assumed away above --- becomes a question about how the tree
is built. Nothing downstream changes, since the induced dynamics is
\eqref{sch:eq:debt}.

Two things follow immediately and neither was available before.

\emph{The win condition is binary.} A branch either clears its debt or
does not, and any branch that clears it completes --- however deep the
debt ran, and with no residual weight to carry. The grading that a
weight formulation would posit emerges instead from the
\emph{distribution} over branches.

\emph{And the recursion is structural.} Debt spawns debt: the extra
gates are gates, they can go wrong, and going wrong incurs more. That is
the cascade's shape, and the object it makes is standard.

\begin{proposition}[The debt process is Galton--Watson; \textbf{[P]}]
\label{sch:prop:gw}
Under Definition~\ref{sch:def:debt}, $Z$ is the population of a
Galton--Watson process with offspring probability generating function
\begin{equation}\label{sch:eq:pgf}
f(u)\;=\;\pi+(1-\pi)\,u^{\,r},
\qquad\text{mean offspring}\quad \varkappa\;=\;r(1-\pi).
\end{equation}
Completion is extinction of $Z$. Hence completion occurs almost surely
if and only if $\varkappa\le1$, and for $\varkappa>1$ the probability of never
completing is $1-q^{\,n}$, where $q$ is the least root of $q=f(q)$ in
$[0,1]$.
\end{proposition}

\begin{remark}[A phase transition in the schedule's own structure]
\label{sch:rem:criticality}
$\varkappa=r(1-\pi)=1$ is a sharp threshold, and it is not a restatement of
anything already in the paper.

Below it, debt clears with probability one: however unlucky a biosphere
is, it owes a finite sequence and eventually pays it. Above it, a
positive fraction of branches accumulate debt without bound --- cycling
into ever-deeper obligation, each new gate as fallible as the last ---
and those branches never reach the state at which $x$ can be held. The
transition sits at $r=1/(1-\pi)$: a biosphere may tolerate large debt
per failure if failures are rare, or frequent failures if each is cheap,
and the product is what decides.

The threshold is structurally the anatomy threshold's sibling. $\rho=0$
divides configurations by whether acquired capacity protects; $\varkappa=1$
divides them by whether incurred debt is payable. Both are sign tests on
a structural parameter, both are sharp, and neither is visible in any
magnitude the system displays.
\end{remark}

\begin{remark}[Under log-spread rates the criterion is the geometric
mean; \textbf{[P]}]\label{sch:rem:bpre}
Definition~\ref{sch:def:rates} redraws the rates at every gate, so
$\pi_k$ --- and with it $\varkappa_k=r(1-\pi_k)$ --- is not a deterministic
sequence but an i.i.d.\ random one. The debt process is then a branching
process in a \emph{random} environment rather than a varying one, and the
criterion changes accordingly. By the Smith--Wilkinson criterion,
completion is almost sure exactly when
\begin{equation}\label{sch:eq:bpre}
\E\bigl[\log \varkappa\bigr]\;<\;0,
\end{equation}
the \emph{geometric} mean of the offspring means, and not when
$\E[\varkappa]<1$.

The two disagree, and the direction of the disagreement matters. By
Jensen $\E[\log \varkappa]\le\log\E[\varkappa]$, so a schedule may complete with
probability one while its \emph{arithmetic} mean offspring exceeds
unity. Using $\E[\varkappa]$ is therefore not a conservative simplification but
a wrong answer in the pessimistic direction: it condemns schedules that
in fact clear.\footnote{Simulation, Poisson offspring, $4000$ trials
each. With $\varkappa$ equal to $0.1$ or $5$ with equal probability,
$\E[\varkappa]=2.55$ and $\E[\log \varkappa]=-0.348$: the arithmetic criterion calls it
supercritical, the geometric criterion subcritical, and completion is
observed with frequency $1.000$. Where the two criteria agree ---
$\varkappa\in\{0.5,3\}$ and $\varkappa\in\{0.9,1.6\}$, both geometrically supercritical
--- completion frequencies are $0.738$ and $0.690$.}

One caution about what the criterion certifies. \eqref{sch:eq:bpre}
settles whether the debt clears \emph{eventually}; it does not say when.
Under a random environment the probability of not having completed by
stage $n$ decays at the tilted rate
$\gamma^\ast=-\min_{0\le s\le1}\log\E[\varkappa^{\,s}]$, which is bounded
above by $|\E\log\varkappa|$ and can fall far below it: for the
$\varkappa\in\{0.1,5\}$ environment above, $\gamma^\ast=0.0158$ at
$s^\ast=0.092$ against $|\E\log\varkappa|=0.347$, a factor of
twenty-two.\footnote{The rate identity is the standard large-deviation
form for subcritical branching in a random environment; the polynomial
correction depends on the subcriticality class.} Certain completion and
timely completion are therefore different questions, and only the first
is settled here \textbf{[O]}.

The reading is the one this work keeps arriving at. A geometric mean is
what a typical realization experiences; an arithmetic mean is what an
ensemble average reports. The quenched quantity governs, exactly as
Remark~\ref{sch:rem:gap2} says it does for the completion weight, and
here it governs the dichotomy itself rather than only the rate.
\end{remark}

\begin{remark}[Unrecoverable failures, and a doom weight from the
schedule alone]\label{sch:rem:defective}
Definition~\ref{sch:def:debt} makes every failure recoverable at a
price. Some plausibly are not --- an event arriving before its
prerequisites may foreclose rather than cost --- and the natural
encoding is $r=\infty$, which is to say a \emph{defective} offspring
law. Let a gate resolve correctly with probability $\pi$, recoverably
with probability $q_1$ at debt $r$, and unrecoverably with probability
$q_2=1-\pi-q_1$. Then
\begin{equation}\label{sch:eq:defective}
f(u)=\pi+q_1u^{\,r},\qquad f(1)=\pi+q_1<1,
\end{equation}
so $u=1$ is not a fixed point and the least root of $u=f(u)$ lies
strictly below one \emph{whatever} the recoverable branch does.
Completion probability is bounded away from one by $q_2$ alone, and
subcriticality of the recoverable branch does not rescue it.

Two consequences. The schedule now produces a positive doom weight
unaided --- a nonzero probability of never reaching the onset state,
computed from $(\pi,q_1,q_2,r)$ and nothing else --- which is a second
concrete instance of the graded object
Remark~\ref{sch:rem:mvalued} discusses. And the question whether
foreclosing options dominate acquires two answers rather than one: at
fixed mean offspring the second factorial moment enters the survival
bounds, so they may dominate the \emph{probability} of non-completion
while leaving the mean untouched.\footnote{Checked
against direct simulation of \eqref{sch:eq:debt} with a killing branch,
$4\times10^{4}$ trials each: at $(\pi,q_1,r)=(0.85,0.10,3)$ the least
root is $0.9306$ against $0.9310$ simulated; at $(0.70,0.25,2)$,
$0.9046$ against $0.9040$; at $(0.60,0.35,2)$, $0.8571$ against
$0.8557$.}
\end{remark}

\subsubsection{The offset weighting, emergent}\label{sch:subsub:emergent}

The constraint on this formulation was that it reproduce the companion's
offset weighting rather than replace it. It does, and the derivation
also says what the offset is \emph{made of}.

By Proposition~\ref{sch:prop:gating} a biosphere pays $B+\Gamma$ while
incomplete --- the background unshielded, and the internal channel at its
bare coefficient. Debt is therefore paid in hazard: clearing $Z$ gates at
rate $\lambda$ costs $D/\lambda$ in exposure, and a branch's survival
weight is $e^{-(B+\Gamma)\cdot(\text{total gates fired})/\lambda}$.

\begin{proposition}[Completion weight as a least fixed point;
\textbf{[P]}]\label{sch:prop:fixedpoint}
Let $s:=e^{-(B+\Gamma)/\lambda}$ be the survival factor per gate-time.
Where $\lambda$ is itself drawn per gate, the correct single-$\lambda$
surrogate for the mean log-weight is the \emph{harmonic} mean
$\bigl(\E[\lambda^{-1}]\bigr)^{-1}$, since $-\log u$ is to first order
linear in $1/\lambda$ --- which for a lognormal law is
$\lambda_{\mathrm{geo}}e^{-s_\lambda^{2}/2}$, smaller than the geometric
mean, so passage costs and contrasts both exceed what
geometric-mean substitution would give. The
expected completion weight from a single owed gate is the least root
$u\in[0,1]$ of
\begin{equation}\label{sch:eq:fixedpoint}
u\;=\;s\,f(u)\;=\;s\bigl(\pi+(1-\pi)u^{\,r}\bigr),
\end{equation}
and from $n$ initial gates it is $u^{\,n}$. For small $B/\lambda$,
\begin{equation}\label{sch:eq:expansion}
-\log u\;=\;\frac{B+\Gamma}{\lambda}\cdot\frac{1}{1-\varkappa}
+O\!\left(\left(\tfrac{B}{\lambda}\right)^{2}\right),
\qquad \varkappa=r(1-\pi)<1 .
\end{equation}
\end{proposition}

\begin{proof}
\eqref{sch:eq:fixedpoint} is the total-progeny functional equation for a
Galton--Watson process weighted by $s$ per individual. For
\eqref{sch:eq:expansion}, write $u=1-\varepsilon$ and expand:
$1-\varepsilon=(1-(B+\Gamma)/\lambda)(1-\varkappa\varepsilon)+O(\varepsilon^{2})$
gives $\varepsilon(1-\varkappa)=(B+\Gamma)/\lambda$. The factor $1/(1-\varkappa)$ is the expected total
progeny per individual, as it must be.
\end{proof}

\begin{remark}[What the offset is made of, and where the recursion
went]\label{sch:rem:cefff}
Two effects separate cleanly, and the separation is this subsection's
return.

\emph{Contrast per unit of depth.} A single wrong gate adds $r$ gates,
each exponential at rate $\lambda$, so the added waiting is
$\mathrm{Gamma}(r,\lambda)$ and the weight multiplies by
$\E[e^{-hT}]=(1+h/\lambda)^{-r}$ with $h:=B+\Gamma$. Matching to the
companion's $10^{-c_{\mathrm{eff}}}$,
\begin{equation}\label{sch:eq:cefff}
c_{\mathrm{eff}}\;=\;r\,\log_{10}\!\Bigl(1+\frac{B+\Gamma}{\lambda}\Bigr),
\end{equation}
exactly and with no approximation. The deterministic-exposure reading
--- $r/\lambda$ of waiting at hazard $B+\Gamma$, giving
$(B+\Gamma)r/\lambda\ln10$ --- is the $\lambda\gg B+\Gamma$ limit of it
and \emph{overstates} the contrast elsewhere, since $\log(1+x)<x$. The
direction of that error is worth naming: it makes the companion's
collapse look easier than it is, so \eqref{sch:eq:cefff} is what should
be used. The High-Contrast Earth
Hypothesis's ``extremely many orders'' is thereby not a posit but a
reading: one gets many orders when the background is high, the debt
large, or the tempo slow.

\emph{Depth accumulated.} The recursion does not change
\eqref{sch:eq:cefff}; the branching factor $1/(1-\varkappa)$ cancels out of the
per-depth contrast exactly. What it changes is how much depth there is:
expected total gates fired per owed gate is $1/(1-\varkappa)$, so a biosphere
near criticality accumulates depth without the contrast per unit of it
worsening at all. The recursion inflates the exposure, not the exchange
rate --- and inflates it without bound as $\varkappa\to1$.

The separation identifies two independent levers. A biosphere may be
badly placed because each failure is expensive ($c_{\mathrm{eff}}$
large) or because failures compound ($\varkappa$ near one), and the two are not
the same defect.
\end{remark}

\begin{remark}[Contrast wants slow gates; survival wants fast ones]
\label{sch:rem:tension}
\eqref{sch:eq:cefff} has two limits and they pull opposite ways, which
is a structural fact about the channel rather than a defect of it.

Where gates are fast against the standing hazard, $\lambda\gg B+\Gamma$,
the contrast is $c_{\mathrm{eff}}\approx r(B+\Gamma)/\lambda\ln10$ ---
small. Wrong turns cost little, passage is cheap, and alternatives are
barely suppressed: there is nothing to collapse a sum over.

Where gates are slow, $\lambda\ll B+\Gamma$, the contrast is
$c_{\mathrm{eff}}\approx r\log_{10}\bigl((B+\Gamma)/\lambda\bigr)$ --- a
\emph{log-rate gap}, hence many orders of magnitude whenever the gap is
large. But the same slowness prices survival at
$(\lambda/(B+\Gamma))^{\,r}$ per wrong turn, which is astronomically
small.

So the regime supplying the companion's contrast is the regime in which
the realized history had to be enormously lucky, and the two are one
fact seen from either end. That is not a tension to be resolved but the
Miracle phenomenology in rate coordinates: high contrast between natural
histories and low probability of the realized one are not independent
observations.
\end{remark}

\begin{remark}[The collapse is a claim about shares, not about
likelihood]\label{sch:rem:absolute}
A companion point, and the one most easily misread. The sum-collapse
establishes that the realized ordering carries a share of order unity
\emph{of the total}. It says nothing whatever about the size of that
total, and under a log-scale randomization the total is minute.

Two facts about extremes make this precise. The best of $m$ draws from a
spread log-rate law is not far above the field: the maximum of $m$
normals sits $\sqrt{2\ln m}$ deviations above the mean asymptotically,
and only $1.54$ deviations at $m=10$ --- so the optimum is not
qualitatively better than a typical option, it is the same law read at
its edge. And the optimum still faces the standing hazard: with
$\lambda\ll B+\Gamma$ the weight of even the best option at a gate is
$\approx\lambda/(B+\Gamma)$, orders below one, and $n$ gates multiply
those orders.

A worked instance makes the two magnitudes visible together. Take
$\log_{10}\lambda$ normal with mean $15$ decades below $B+\Gamma$ and
deviation $3$, with $m=10$, $n=20$, $r=3$. The best option at a gate
sits $4.6$ decades above typical; the weight of the all-best path is
about $10^{-208}$; a single wrong turn costs $c_{\mathrm{eff}}\approx31$
decades; and the collapse condition $c_{\mathrm{eff}}>\log_{10}(n(m-1))
=2.26$ is met roughly fourteen times over.\footnote{Direct simulation,
$2\times10^{5}$ draws; the extreme-value theory agrees to two decimals,
$\mu+1.5388s=-10.38$ against $-10.39$ observed.}

So relative dominance is overwhelming and absolute weight is
$10^{-208}$, and the two are not in tension --- both follow from the
same spread, which is what makes one alternative dominate its rivals and
what makes all of them improbable. The reading to avoid is therefore
explicit: \emph{the highest-weight ordering is not the likely ordering}.
Nothing in the companion's apparatus, and nothing here, converts a share
of a small number into a large one. This is the same discipline
\S\ref{gr:sub:consequences} applies to the survivor's own position, and
it applies with more force here because the numbers are worse.
\end{remark}

\begin{remark}[Which contrast the collapse consumes]\label{sch:rem:n2}
A random contrast does not enter the companion's sum-collapse the way a
deterministic one does, and the difference has to be recorded.

An alternative at depth $\Delta$ is suppressed by
$10^{-\sum_{i\le\Delta}c_i}$, so with independent draws the per-tier
factor is $(\E[10^{-c}])^{\Delta}$ and the binomial closes at
$\bigl(1+(b-1)\E[10^{-c}]\bigr)^{n}$. The operative quantity is
therefore $-\log_{10}\E[10^{-c}]$, not $\E[c]$, and by Jensen the first
is strictly smaller. The gap is severe rather than technical: for $c$
uniform on $(0,2M)$ the averaged contrast is $\log_{10}(2M\ln10)$,
growing only \emph{logarithmically} in the spread; for $c$ normal with
mean $\mu$ and deviation $s$ it is $\mu-s^{2}\ln10/2$, negative once
$s\gtrsim\sqrt{2\mu/\ln10}$.\footnote{At $\E[c]=10$: uniform on
$(0,20)$ gives $1.66$; normal with $s=2$ gives $5.45$; normal with $s=3$
gives $-0.36$. Direct sampling understates the collapse --- four million
draws return $1.25$ for the last --- because $\E[10^{-c}]$ is carried by
a tail sampling does not reach, which is itself the point.}

The resolution is the one this work keeps reaching. The averaged
contrast is dominated by alternatives no one encounters; the claim the
collapse is used for concerns the share carried by the \emph{realized}
configuration, which is quenched. One refinement to ``runs on $\E[c]$'':
the quenched tier weight is a sum of order $\bigl(n(b-1)\bigr)^{\Delta}$
lognormal terms and is extreme-value dominated at small $\Delta$, so the
binding constraint sits at the shallowest tier and carries a fluctuation
allowance,
\begin{equation}\label{sch:eq:quenchedcrit}
\bar c\;>\;\log_{10}\bigl(n(b-1)\bigr)
\;+\;\frac{s\sqrt{2\ln\bigl(n(b-1)\bigr)}}{\ln10},
\end{equation}
relaxing toward the $\E[c]$ form as $\Delta$ grows, since the fluctuation
scales as $\sqrt\Delta$ against a mean linear in $\Delta$. At $s=3$,
$n=20$, $b=10$ the threshold moves from $2.26$ to $6.45$, which
$\E[c]=10$ still clears. The companion's
own hedge --- that its deterministic use is the union-bounded envelope
of an almost-sure claim --- points the same way. So the collapse
survives randomization in the quenched reading and fails in the annealed
one, and a reader wanting the annealed statement needs a contrast far
larger than $\E[c]$ suggests \textbf{[O]}.
\end{remark}

\begin{remark}[This is the graded cascade, in one concrete case]
\label{sch:rem:mvalued}
Remark~\ref{gc:rem:twolayers} keeps the cascade's Boolean recursion and the
$\varepsilon$-grading of Definition~\ref{gc:def:epsdep} in separate
layers, and records what would join them: a fixed point of an
$M$-valued doom operator, called open \textbf{[O]} and named there as
the natural next question.

\eqref{sch:eq:fixedpoint} is an instance of exactly that object. The map
$u\mapsto s\,f(u)$ is monotone on $[0,1]$, its least fixed point is the
graded analog of $\mu G$'s least fixed point on the Boolean lattice,
and the grading it carries is a weight rather than a label. The
completion probability is therefore \emph{computed} as a fixed point
rather than assigned, which is what the two-layer remark asks for.

The claim should be taken at its size. This is a single-type process
carrying one number, not the $M$-valued operator over the full lattice
of states the general construction would need; what is exhibited is the
first concrete case and the shape it confirms.

What the case does supply is the general construction's
\emph{definition}, and the choice is not the obvious one. The Boolean
cascade characterizes $\mu G$ modally: every course admits a foreclosing
continuation, which is a fact about the state space. A graded version
cannot be built on that, because a modal fact carries no measure. It can
be built on the dynamical reading --- \emph{all paths lead to} the
Departure Point, rather than all courses make it exist --- since that is
a hitting statement, and the measure of paths hitting a set is
canonically the \emph{least} fixed point of the associated operator.
\eqref{sch:eq:fixedpoint} is exactly such a hitting measure, which is why
it arrived as a least fixed point without anyone imposing it.

Two things follow for the general construction. The operator to build is
the one whose fixed point is the hitting measure of the departure set,
and Kleene iteration from $\bot$ reaches it wherever the weight algebra's
multiplication is Scott-continuous. And the projection back to the
Boolean layer is a \emph{threshold} rather than a linear map --- Boolean
$\mu G$ is where the surviving weight is zero, not where the graded value
takes a particular magnitude --- so the consistency requirement is
$\mathrm{eval}(u)=\one\{u=0\}$, and it should be imposed at definition
time rather than proved afterwards. Whether the construction so specified
generalizes to the lattice remains open \textbf{[O]}.
\end{remark}

\begin{remark}[Which rate an observer should use]\label{sch:rem:gap2}
The completion weight $u^{\,n}$ is an expectation over branches, and
near criticality it is dominated by the rare branch that incurred little
debt: by \eqref{sch:eq:expansion} the \emph{mean} total progeny is
$1/(1-\varkappa)$, while the median sits far below it. An observer reasoning
about their own history should therefore not use the mean. Expected
completion weight is carried by branches that got lucky early, and a
biosphere is not entitled to assume it is on one --- the same
quenched--annealed discipline the reference sheet records for the
survivorship transform.\footnote{The same gap is recorded for the survivorship transform in Appendix~\ref{app:sheet}; the two are instances of one phenomenon rather than one result applied twice.}
\end{remark}

\subsubsection{The composite}\label{sch:subsub:composite}

The schedule does not win Golden-Point probability. It wins the
probability of reaching the one state at which $x$ stops decaying, and
everything the model already proves takes over from there. That gives
the total a two-stage shape:
\begin{equation}\label{sch:eq:composite}
\PP(\text{advance})
=\underbrace{\PP\bigl(Z\ \text{clears}\bigr)}_{\text{schedule:
\eqref{sch:eq:fixedpoint}}}
\;\times\;
\underbrace{\E\Bigl[\PP\bigl(\Lambda_\infty<\infty\mid X_{\tau}\bigr)
\Bigr]}_{\text{the model, from the completion state}},
\end{equation}
with $\tau$ the completion time and $X_\tau=(\mu_\tau,D_\tau)$ the
\emph{model's} state at that moment --- the Kolmogorov pair of
Definition~\ref{tc:def:model}, and not the schedule's debt $Z$, which is
zero at $\tau$ by construction. Factoring the second stage by
Theorem~\ref{gr1:thm:twocond},
\begin{equation}\label{sch:eq:threecond}
\PP\bigl(\Lambda_\infty<\infty\bigr)
=\underbrace{u^{\,n}}_{\text{schedule}}
\;\cdot\;\underbrace{\one\{\rho>0\}}_{\text{anatomy}}
\;\cdot\;\underbrace{\frac{s(\mu_0)-s(-\infty)}{s(+\infty)-s(-\infty)}}
_{\text{escape}} .
\end{equation}

\begin{remark}[The three factors are not of one kind]
\label{sch:rem:threefactor}
Anatomy \emph{annihilates}: $\rho\le0$ gives zero, not a small number,
by Theorem~\ref{gr1:thm:trapped}. The schedule \emph{attenuates}: a
wrong gate costs debt, and debt cleared is completion, so the factor is
drastically minute and never quite zero --- which is what it is to be an
Imperfection Magnitude Drop in the companion's exact sense \cite{Nexus}.
So the schedule factor is the residual $\varepsilon$ of
Definition~\ref{gc:def:epsdep} computed rather than named, and the
channel's connection to the cascade is an identification rather than an
analogy.
\end{remark}

\begin{remark}[The decomposition is hierarchical, not a product]
\label{sch:rem:hierarchical}
\eqref{sch:eq:composite} should be read exactly. It is not an
independence factorization --- $X_\tau$ depends on how long the debt
took to clear, so the two factors are not independent --- but a
hierarchical decomposition, and what makes it usable is that the
influence runs one way. The schedule's law is autonomous: $\lambda_k$
and $\pi_k$ consult the gate count and nothing else, so the schedule may
be solved first, without reference to the driver. The leak then carries
the schedule's verdict into the driver through $\gamma_{k_t}$, and the
model runs from the state that leaves. Schedule, then driver, then
model, with no arrow back.

That one-directionality is what a state-dependent schedule threatens,
and it threatens it in only one of the two ways such dependence can be
read. Let $\pi$ depend on the schedule's own \emph{configuration} ---
which gates are closed, in what order, with what debt outstanding ---
and the schedule remains autonomous, since configuration is generated by
the schedule itself. It becomes a multitype branching process, with
types the configurations and criticality governed by the top Lyapunov
exponent of products of mean matrices rather than by the scalar
criterion of \S\ref{sch:subsub:open}, and the hierarchy survives intact.
Let $\pi$ instead depend on the \emph{driver} --- on $x_t$ or $\mu_t$
--- and the schedule consults what it was supposed to precede. The
hierarchy collapses, \eqref{sch:eq:action} becomes a genuine control
problem, and the optimal history must trade action spent early against
exposure shortened.

The baton-passing picture that motivates state dependence is of the
first kind: what an event does depends on the state its predecessor
left, which is configuration. So the version the natural-historical
reading requires costs a multitype upgrade and not the control problem.
Whether the driver \emph{also} influences gate outcomes is a modeling
question about the world rather than a mathematical one, and the answer
adopted here is that it does not. Gate outcomes are stochastic in the
direct sense: what makes an option the realized one is a fact about the
onward schedule it opens, not about the credit held when it fires. The
hierarchy therefore holds, \eqref{sch:eq:action} does not become a
control problem, and the multitype upgrade is the whole cost of state
dependence.

One refinement that adoption invites is recorded rather than pursued. If
an option is realized in virtue of the ease of the gates it opens, then
correctness is not a label attached to one of the $m$ branches but a
\emph{derived} property of the subtree below it --- which makes the
schedule a tree carrying a value function rather than a sequence with a
marked option, and makes the realized path something survivorship picks
out rather than something given in advance. That is consistent with
Definition~\ref{sch:def:rates}, whose independence claim concerns an
option's \emph{own} rate and not its descendants'. Whether the tree
formulation reproduces the sequence one at the level of
\eqref{sch:eq:fixedpoint} is open \textbf{[O]}.
\end{remark}

\begin{remark}[Only the asymptote selects the schedule]
\label{sch:rem:asymptote}
The channel prices what climbing costs. Whether anything \emph{selects}
for climbing is a separate question, and the answer separates the two
conditioning events this work uses.

\emph{Finite horizons do not select it.} Condition on $\{T>t\}$ for any
finite $t$. A biosphere that never completes its schedule sits at bare
background by Proposition~\ref{sch:prop:gating} and survives to $t$ with
probability $e^{-Bt}>0$. The un-climbed branches therefore carry
positive weight under every finite-horizon conditioning, and the lucky
route --- survive by not being struck, and never pay the schedule at all
--- remains in the support throughout.

\emph{The asymptote does.} Condition instead on the viability event
$\Imm=\{\Lambda_\infty<\infty\}$ of Definition~\ref{sh:def:viab}. By
Proposition~\ref{sch:prop:gating},
$\PP\bigl(\Imm\mid\text{schedule incomplete}\bigr)=0$ \emph{exactly} ---
not small. The exclusion is at the level of support rather than of
weight, and it is the only conditioning in this work that achieves it.

So the schedule is selected for by the asymptotic condition and by
nothing weaker. The crossover is worth seeing quantitatively, because it
is where the balance actually sits. Survival at bare background is
$e^{-Bt}$; survival once completed and critical is $\asymp Ct^{-1/4}$ by
Claim~\ref{sh:claim:tail}. A polynomial beats an exponential eventually
and always, so completion wins in the end for every $C>0$; but the
crossover $T_\ast$ solving $e^{-BT}=CT^{-1/4}$ recedes as $C$ falls, and
$C$ is set by the state completion leaves
(Remark~\ref{sch:rem:initialdata}). A biosphere completing into a poor
state finds the schedule worth having climbed only at horizons that run
away toward the asymptote.

The reading is the one that matters for what this channel is for.
\emph{Nothing about surviving a finite stretch gives a reason to climb.}
Only conditioning on there being no last moment makes the climb pay ---
which is why the schedule appears in the record of a biosphere whose
history was conditioned, and would not appear in the record of one that
had merely lasted.
\end{remark}

\begin{remark}[The schedule supplies the initial data the model assumes]
\label{sch:rem:initialdata}
\eqref{sch:eq:composite} makes an identification that should be stated plainly:
\emph{the completion state is the initial datum} that
Definition~\ref{sh:def:model} takes as given. The model begins where the
schedule ends, and the schedule delivers not a point but a distribution
over points --- shaped by how long the debt took to clear and how far
$x$ decayed while it did.

By Proposition~\ref{sh:prop:gauge} the initial data are gauge: they act
through prefactors and timescales, never through exponents. So the
schedule cannot move the persistence exponent $\tfrac14$, cannot move
the trichotomy, and cannot move the two-condition theorem. What it sets
is the prefactor --- and by Theorem~\ref{sh:thm:persist} the prefactor
is proportional to the persistence-harmonic function $\hm$ of the
starting state.

That is the same $\hm$ the Golden Score computes
(Proposition~\ref{gs:prop:score}), and \S\ref{gs:sub:limit} records as
its principal cost that no map exists from a decision to a point in
$(\mu,D)$. The schedule produces exactly such points. Whether decisions
bearing on a biosphere's schedule --- on $\lambda$, on $\pi$, on debt
incurred --- thereby acquire computable harmonic values is the first
place a candidate identification map might be found. We flag it as a
lead and claim nothing \textbf{[O]}.
\end{remark}

\begin{remark}[The variational form, stated and not solved]
\label{sch:rem:variational}
The second factor of \eqref{sch:eq:composite} is a conditioned path
problem, and Remark~\ref{gr1:rem:closedform} already names the tool: the
adverse-drift rate and its front-load-then-coast profile come from the
Cameron--Martin action by Cauchy--Schwarz, the model's analog of an
Euler--Lagrange computation and the natural processing step wherever a
conditioned path law is wanted at large-deviation scale.

The composite asks for that computation across both stages. One seeks
the path minimizing
\begin{equation}\label{sch:eq:action}
S[\mu]\;=\;\int_0^{\infty}\left[\frac{\dot\mu_t^{2}}{2\sigma^{2}}
\;+\;h\bigl(x_t,k_t\bigr)\right]dt
\end{equation}
subject to \eqref{sch:eq:leak}, to the debt dynamics
\eqref{sch:eq:debt}, and to $\Lambda_\infty<\infty$: a least-action path
through a landscape whose running cost is the hazard and whose
constraint is the schedule. The pathfinding reading is exact rather than
figurative --- the action is a length and the optimal history a geodesic
in it.

Two features of the answer are predictable without solving it, and both
are checks on a solution's behavior. The running cost during
incompleteness is $B$ and is \emph{independent of the driver}, since $x$
is not held and $\varsigma_b(0)=1$; so the schedule stage contributes a
term the driver cannot optimize against, and the variational freedom
lives entirely after $\tau$. And because \eqref{sch:eq:composite}
factorizes, the optimal post-completion profile is the existing one,
started from the distribution Remark~\ref{sch:rem:initialdata}
describes rather than from a chosen point.

Whether the stages remain separable once $\lambda$ or $\pi$ depends on
the state --- which the baton-passing picture requires --- is the
substantive question and is open \textbf{[O]}. If they do not,
\eqref{sch:eq:action} becomes a genuine control problem rather than a
factorized one, and the optimal history would trade early action against
schedule exposure in a way the present separation forbids.
\end{remark}

\subsubsection{The internal channel is probably a schedule too}
\label{sch:subsub:internalschedule}

One structural claim precedes the open list, because it is
the channel's clearest suggestion about the rest of the model and because
it is an inference the Continuity Principle licenses rather than a
speculation laid beside it.

Principle~\ref{gc:prin:continuity} holds that hazard-lowering across
mediations is one activity and not a family of analogs. The schedule
represents pre-onset hazard as gates. If continuity is taken at the width
Remark~\ref{gc:rem:continuitycost} states it, then post-onset internal
hazard should be representable the same way, and the trough's smooth
exponential should be a coarse-graining of gates not yet identified.

\begin{remark}[The trough's exponential is geometric compounding;
\textbf{[S]}]\label{sch:rem:internalgates}
Take $h^{\mathrm{int}}$ proportional to the count of \emph{unpatched
failure modes} carried at capacity $x$. Suppose each increment of
capacity both creates modes and passes gates that patch them, and write
$q$ for the fraction left unpatched per unit of capacity. Then
$\text{unpatched}(x)=\text{unpatched}(0)\,q^{\,x}$, so
\begin{equation}\label{sch:eq:internalgates}
h^{\mathrm{int}}(x)\;=\;\Gamma q^{\,x}\;=\;\Gamma e^{-\rho x},
\qquad \rho=\ln(1/q).
\end{equation}
So geometric compounding of per-gate patch rates \emph{would} produce
the exponential of \eqref{tc:eq:hazard}, with $\Gamma$ the count at
$x=0$. The construction is offered at that strength and no more: it
exhibits a mechanism yielding the observed form, not evidence that this
is the mechanism. \eqref{tc:eq:hazard} remains the paper's statement of
the internal channel, nothing downstream is rewritten in gate
coordinates, and the construction's use is heuristic --- it says what one
would look for, and becomes a derivation only if the gates are named.

Three things follow, and the second is the reason to record this.

\emph{The anatomy trichotomy becomes a compounding test.}
Proposition~\ref{gr1:prop:trichotomy} divides on the sign of $\rho$;
under \eqref{sch:eq:internalgates} that is $q$ against $1$ --- patching
outpacing creation, matching it, or falling behind.

\emph{And it is then the same test the schedule already runs.}
Remark~\ref{sch:rem:criticality} calls $\varkappa=1$ the anatomy threshold's
sibling. Under this construction they are not siblings but one test
applied twice: both ask whether a compounding factor exceeds unity, and
they differ only in what indexes the compounding --- gate count for the
schedule's debt, capacity for the trough's unpatched modes. That is what
continuity predicts, and finding it is a check on the principle rather
than a use of it.

\emph{The construction agrees with the one the paper already has.}
\S\ref{gr:subsub:boundary} derives $\rho=-(\alpha-\beta)$ from boundary
scaling --- an extractive domain growing as $e^{\alpha x}$ against
enforcement growing as $e^{\beta x}$. Under \eqref{sch:eq:internalgates}
that is $q=e^{\alpha-\beta}$, so $q>1$ exactly when $\alpha>\beta$, and
the two microfoundations return the same condition by different routes.
Neither is derived from the other.
\end{remark}

What is missing is the thing the reformulation is for: \emph{which
events}. The schedule channel names its gates as natural-historical
transitions, and the internal channel's gates would be whatever a
biosphere must get right for acquired capacity not to be turned inward
--- institutional, presumably, and \S\ref{ad:sub:toward} is where a
candidate list would come from. Until they are named, the smooth rate is the correct representation, and the reason needs precise statement: a
smooth hazard is the \emph{annealed} description of an unidentified jump
process, so by the discipline of Remark~\ref{sch:rem:gap2} it flatters
the typical history. An observer using \eqref{tc:eq:hazard}'s trough to
reason about their own biosphere is using a mean carried by branches that
patched luckily \textbf{[O]}.

\subsubsection{What is open}\label{sch:subsub:open}

\emph{Depth-dependent debt.} $r$ and $\pi$ are constant, so all failures
cost alike. Two quantities must be kept apart before that is repaired,
since they are easily run together and only one of them is ordered by
position. \emph{Stakes} --- what failing forfeits --- rise monotonically
along the chain, by Remark~\ref{hr:rem:monotone}: a late step carries
the improbability of every step behind it. \emph{Debt} --- what
recovering costs --- does not. An early gate may carry enormous $r$: a
biosphere without its moon, or with one a third the size, owes a
compensating sequence for axial stabilization and tides at the
\emph{first} gate rather than the last, and it is not obvious that any
finite sequence discharges it. Severity is not ordered by position, and
a model that assumed it was would misplace the worst gates.

Repairing this makes $Z$ a branching process in a varying environment,
$\varkappa_k=r_k(1-\pi_k)$ a sequence rather than a number, and the criterion a
condition on \emph{products} rather than on any single term: writing
$S_k=\prod_{j<k}\varkappa_j$, completion is almost sure when
$\sum_k S_k^{-1}=\infty$ and has positive failure probability when the
sum converges, under a uniform second-moment condition on the offspring
laws. For constant $\varkappa$ this recovers $\varkappa\le1$. The reading is worth
having: one catastrophic gate does not doom a schedule whose remaining
gates are benign enough, since it is the geometric mean of the $\varkappa_k$
that decides --- a single spike is survivable in a way a raised floor is
not --- with the caveat that under Definition~\ref{sch:def:rates} the
environment is random rather than merely varying, so the operative
criterion is \eqref{sch:eq:bpre} and the sum below describes the
deterministic case. The question of whether a deeper drop costs more
\emph{per unit of debt} is answered there rather than here, and answered
negatively: the rates are redrawn at each gate and correctness is
independent of them, so severity is neither constant nor monotone in
depth but resampled, and it is the distribution rather than any trend
that carries the structure. The same harmonic sum gives the rate and not
merely the dichotomy:
two-sided bounds of Agresti type put the probability of not having
completed by stage $n$ comparable to $\bigl(\sum_{k\le n}S_k^{-1}
\bigr)^{-1}$, up to the second-moment factor, so the quantitative form
comes free with the criterion and is what a bridge to
Proposition~\ref{sch:prop:exponent}'s timing statements would consume
\textbf{[O]}.

\emph{State-dependent gates.} $\lambda(\cdot)$ and $\pi(\cdot)$ are what
the baton-passing picture requires and what
Remark~\ref{sch:rem:variational} identifies as deciding whether the
composite factorizes.

\emph{The leak is assigned rather than emergent}, and the incomplete
regime's asymptotics remain unknown (Remark~\ref{sch:rem:limit}).

\emph{And the parameters are unmeasured.} $B$, $\lambda$, $\pi$, $r$ and
$n$ are what the companion's serial-attribute program
(Remark~\ref{gc:rem:serialcost}) would supply --- and \eqref{sch:eq:cefff}
makes that program's task easier rather than harder, since
$c_{\mathrm{eff}}$ need no longer be estimated directly but follows from
three quantities a historian might plausibly bound.
Remark~\ref{gr1:rem:noproxy}'s reservation applies with full force
throughout: a scored schedule index would be a measurement apparatus
optimized against, which is the failure the channel exists to price.

% ---------------------------------------------------------------------

\begin{remark}[The third term is not exhausted; author dictation,
\textbf{[S]} with imports at \textbf{[$\Phi$]}]
\label{sch:rem:notexhausted}
The working assumption until this dictation was that the schedule
channel belonged mostly to the record: the evolutionary ladder
climbed, the term's modern residue mergeable into the internal and
external channels. The continuity principle says otherwise, and
saying it plainly corrects the author's own prior practice. Nothing
in the conditioning marks the present as the moment the schedule
closed, so the honest default under continuity is that the
prerequisites for establishing the Golden Point are not yet
exhausted, and the third term is live in the modern day rather than
spent in the Pleistocene.

The dictation's central move is then an identification. What the
accordance principle under Golden conditioning renders legible as
golden causality in the modern record, the specific
accordance-yielding attributes physically observable as they occur,
is not a separate phenomenon that happens to be visible; it is the
third term operating. The standing conviction that golden causality
is more visible than the internal and external channels would
suggest is on this reading not a resolution artifact of the
attribute level (Remark~\ref{gr1:rem:aethicsubstrate}); it is
channel identity, accordance-legible concentrations being schedule
steps landing. This sharpens that remark's open question into an
identification carried at \textbf{[S]}, and it carries a scoreable
texture, entered here with the conditions of its failure: if the
identification holds, the rarity audit should find modern
concentrations clustering on prerequisite-shaped structures rather
than spreading uniformly over the record, and a uniform spread would
count against it.

The class of events in question is the sequential-achievement class
\cite{Ben25Seq}, and its citable historical exemplar keeps the
discussion honest about magnitudes: Newton's occurrence when he
occurred, the quantile of intellectual horsepower struck at that
point of the process, was an event of low probability relative to
the alternatives, individually improbable and collectively required,
which is the signature of a schedule step. The general claim leans
on the class, and on the expectation that most of its modern members
have not yet been identified; no single member is load-bearing.

One candidate modern member should be named exactly once, because
naming it is the only way to fence it. The production of the present
framework by its date is itself an event whose probability was low
relative to the alternatives, given how many things had to go right
to lead to it, and it is therefore a formally admissible candidate
instance of the class. It is entered here as hypothetical only and
expressly not load-bearing: the fluke hypothesis may simply be true
of it, and the occupancy fences of
Remark~\ref{gs:rem:arrivaltiming} bar the reflexive claim in any
case. The class carries the argument, and no member of it does.

None of the riskiness a live third term adds changes the conditional
picture, and this is the dictation's retention clause. The forward
forcing of Remark~\ref{gr1:rem:forwardforcing} reads on this channel
too: under conditioning at the asymptotic future, in character much
like a stationary-action principle, the conditional law concentrates
on schedules that complete, because reaching the Golden Point
outweighs the third term's hazard once the conditioning is
asymptotic. The prerequisites still standing are, under that
conditioning, nailed with conditional probability approaching one,
which retains in the forward direction exactly what the Optimation
climb established in the backward one.
\end{remark}

\subsection{Conditioning on survival}\label{gr1:sub:cond}\label{sh:sub:interp}

\subsubsection{The interpolation is a monotone flow of upward tilts}
\label{sh:subsub:flow}

\begin{definition}[Survivorship interpolation; bias field]
\label{sh:def:interp}
For $t>0$ set $\PP^{(t)}:=\PP(\,\cdot\mid T>t)$, and analogously
$\PP^{(t)}_\tau:=\PP(\,\cdot\mid\tau>t)$ for the hard-wall functional;
$\PP^{(0)}=\PP$ is the unbiased law. For $0\le s<t$ the \emph{bias field}
on the window $[0,s]$ is the Radon--Nikodym derivative
\begin{equation}\label{sh:eq:bias}
\beta_t(s):=\frac{d\PP^{(t)}}{d\PP}\Big|_{\Fc_s}
=\one_{\{T>s\}}\,
\frac{\PP_{X_s}(T>t-s)}{\PP_{X_0}(T>t)},
\qquad X_s=(\mu_s,D_s),
\end{equation}
the second equality by the Markov property; likewise for
$\PP^{(t)}_\tau$ with $\tau$ in place of $T$. The interpolation at finite
$t$ is thus globally absolutely continuous with respect to $\PP$, with
density $\one_{\{T>t\}}/\PP(T>t)$ on $\Fc_\infty$.
\end{definition}

\begin{lemma}[Survivorship tilts upward; \textbf{P}]\label{sh:lem:monotone}
In the driftless model, couple two initial data
$(\mu_0,D_0,h_0)\le(\mu_0',D_0',h_0)$ (componentwise in the first two
arguments) synchronously, with the same $W$. Then pathwise
$\mu'_t-\mu_t=\mu_0'-\mu_0\ge0$,
$D'_t-D_t=(D_0'-D_0)+(\mu_0'-\mu_0)t\ge0$, hence $h'\le h$,
$\Lambda'\le\Lambda$, and both $\PP_x(T>t)$ and $\PP_x(\tau>t)$ are
nondecreasing in $(\mu_0,D_0)$ (and nonincreasing in $h_0$). Consequently
every bias field \eqref{sh:eq:bias} is, on $\{T>s\}$, a nondecreasing
function of the state $X_s$: at every conditioning horizon, survivorship
bias is an upward tilt in rate and credit, and the interpolation is a
monotone flow of such tilts.
\end{lemma}

\begin{proof}
Immediate from the displayed pathwise inequalities and
Definition~\ref{sh:def:cox}.
\end{proof}

\subsubsection{The critical limit is the pure Doob transform}
\label{sh:subsub:Q}

\begin{theorem}[Existence and identification of the limit;
\textbf{C}, soft version \textbf{A}]\label{sh:thm:Q}
Assume the pointwise refined tail asymptotics
$\PP_{(\mu,d)}(\tau>t)\sim c\,\hm(\mu,d)\,t^{-1/4}$ of
Theorem~\ref{sh:thm:persist} and
Proposition~\ref{sh:prop:harmonic}, together with the martingale property
$\E_{x}[\hm(X_s)\one_{\{\tau>s\}}]=\hm(x)$ furnished by \cite{GJW99}.
Then for every $s\ge0$ and every bounded $\Fc_s$-measurable $\Psi$,
\[
\E\bigl[\Psi\mid\tau>t\bigr]\ \longrightarrow\ \E_Q[\Psi],
\qquad t\to\infty,
\]
where $Q$ is the time-homogeneous Markov law with
\begin{equation}\label{sh:eq:Qdensity}
\frac{dQ}{d\PP}\Big|_{\Fc_s}
=\one_{\{\tau>s\}}\,\frac{\hm(X_s)}{\hm(X_0)} :
\end{equation}
the \emph{pure} Doob $\hm$-transform, at eigenvalue one. Under
Ansatz~\ref{sh:ansatz} the same limit holds for
$\PP^{(t)}=\PP(\,\cdot\mid T>t)$.
\end{theorem}

\begin{proof}
Write $f_t:=\one_{\{\tau>s\}}\PP_{X_s}(\tau>t-s)/\PP_{X_0}(\tau>t)$, so
that $\E[\Psi\mid\tau>t]=\E[\Psi f_t]$. By the Markov property
$\E[f_t]=1$ for every $t$. Factorizing,
\[
f_t=\one_{\{\tau>s\}}\,
\frac{(t-s)^{1/4}\,\PP_{X_s}(\tau>t-s)}{t^{1/4}\,\PP_{X_0}(\tau>t)}
\,\Bigl(1-\frac st\Bigr)^{-1/4}
\ \longrightarrow\
f:=\one_{\{\tau>s\}}\frac{\hm(X_s)}{\hm(X_0)}
\quad\text{a.s.},
\]
using the assumed asymptotics at the (fixed, random) starting point
$X_s$ and at $X_0$, and $(1-s/t)^{-1/4}\to1$. The martingale property
gives $\E[f]=1=\lim\E[f_t]$, so by Scheff\'e's lemma $f_t\to f$ in
$L^1(\PP)$, and $\E[\Psi f_t]\to\E[\Psi f]=\E_Q[\Psi]$ for bounded
$\Psi$. That \eqref{sh:eq:Qdensity} defines a consistent, time-homogeneous
Markov family is the standard theory of Doob transforms by invariant
harmonic functions; time-homogeneity with no eigenvalue factor is
possible precisely because $\hm$ is invariant (eigenvalue one) for the
killed semigroup.
\end{proof}

\begin{remark}[Why survivorship has no clock here; \textbf{P}]
\label{sh:rem:noclock}
The mechanism deserves isolation. For \emph{regularly varying} tails the
kinematic factor $\PP_x(\tau>t-s)/\PP_x(\tau>t)\to1$ for each fixed $s$:
the age of the conditioning drops out at leading order, and the limit is
an un-tilted $h$-transform. For \emph{exponential} tails
$\PP_x(T>t)\approx C(x)e^{-\theta t}$ the same factor converges to
$e^{\theta s}$, and survivorship acquires a clock --- the eigen-tilting
of mode (III) in Theorem~\ref{sh:thm:modes}. The polynomial class is
exactly the class in which survivorship bias is, asymptotically, a
change of measure with no time-dependence. In the taxonomy of
penalizations of Wiener measure \cite{RVY06}, the survival weight
$e^{-\Lambda_t}$ thus lands in the clockless (martingale-normalizable)
class precisely because its expectation is regularly varying.
\end{remark}

\begin{proposition}[Interpolation rate; the tail index as linear-response
coefficient; \textbf{P/C}]\label{sh:prop:rate}
Under the hypotheses of Theorem~\ref{sh:thm:Q}, for fixed $s$ as
$t\to\infty$,
\begin{equation}\label{sh:eq:rate-interp}
\beta_t(s)\;=\;\one_{\{\tau>s\}}\,\frac{\hm(X_s)}{\hm(X_0)}
\Bigl(1-\frac st\Bigr)^{-1/4}\bigl(1+o(1)\bigr),
\qquad
\Bigl(1-\frac st\Bigr)^{-1/4}=1+\frac{s}{4t}+O\!\Bigl(\frac{s^2}{t^2}\Bigr).
\end{equation}
The natural interpolation parameter is therefore $s/t$, and the
persistence exponent $\tfrac14$ is the linear-response coefficient of
survivorship bias: the leading finite-$t$ correction to the survivor's
universe is $s/4t$ per window of length $s$. Approach to the limit is
polynomial --- the critical signature, contrasting with the exponential
approach of mode (I) and the exponentially growing tilt of mode (III)
below.
\end{proposition}

\begin{proof}
Insert the assumed asymptotics into \eqref{sh:eq:bias} and expand.
\end{proof}

\subsubsection{The limit object: survivorship drift and promotion to
transience}\label{sh:subsub:promotion}

\begin{theorem}[Survivorship drift; \textbf{P} given \textbf{C} inputs]
\label{sh:thm:drift}
Under $Q$ the state solves
\begin{equation}\label{sh:eq:Qsde}
d\mu_t=\sigma^2\,\partial_\mu\log\hm(\mu_t,D_t)\,dt+\sigma\,dW^Q_t,
\qquad dD_t=\mu_t\,dt,
\end{equation}
for a $Q$-Brownian motion $W^Q$: conditioning adds drift only in the
noisy coordinate. With the representation \eqref{sh:eq:harm},
\begin{equation}\label{sh:eq:G}
b_Q(\mu,D)=\frac{\sigma^2}{\mu}\,G(\kappa),
\qquad
G(\kappa):=\frac12-3\kappa\,\frac{F'(\kappa)}{F(\kappa)},
\end{equation}
and: \emph{(i)} $b_Q\ge0$ everywhere (survivorship never pushes the rate
down), by Lemma~\ref{sh:lem:monotone} applied to $\hm$;
\emph{(ii)} $G(0^+)=\tfrac12$, since
$\hm(\mu,0^+)=c_1'\mu^{1/2}$ by Theorem~\ref{sh:thm:persist};
\emph{(iii)} $G(\kappa)\to0$ as $\kappa\to\infty$, since
$\hm(0^+,d)=c_2 d^{1/6}$ is independent of $\mu$ at leading order.
Statements (ii)--(iii) use the regularity of the profile $F$ furnished by
the explicit formulas of \cite{GJW99}.
\end{theorem}

\begin{proof}
For invariant harmonic $\hm$, the transformed generator is
\[
\Lc^{\hm}f=\hm^{-1}\Lc(\hm f)
=\Lc f+\sigma^2\bigl(\partial_\mu\log\hm\bigr)\,\partial_\mu f,
\]
since only $\partial_\mu^2$ is second order in
\eqref{sh:eq:kolm}; this is \eqref{sh:eq:Qsde}. Formula \eqref{sh:eq:G}
is the logarithmic derivative of \eqref{sh:eq:harm} using
$\partial_\mu\kappa=-3\kappa/\mu$. Positivity: synchronous coupling gives
$\PP_{(\mu,d)}(\tau>t)$ nondecreasing in $\mu$, hence
$\partial_\mu\hm\ge0$ wherever the derivative exists. The boundary values
follow from the quoted prefactor laws: $F(0)=c_1'\in(0,\infty)$ and
$\kappa F'(\kappa)/F(\kappa)\to0$ as $\kappa\to0$ by smoothness of $F$ at
the origin; as $\kappa\to\infty$, $\hm\sim c_2d^{1/6}$ forces
$\mu\,\partial_\mu\log\hm\to0$.
\end{proof}

\begin{corollary}[Wall drift; promotion to transience; \textbf{P/C}]
\label{sh:cor:bessel}
Near zero credit ($\kappa\to0^+$) the survivorship drift reduces to
$\sigma^2/2\mu$, so that in this regime the conditioned rate solves
\[
d\mu_t=\frac{\sigma^2}{2\mu_t}\,dt+\sigma\,dW^Q_t .
\]
Two cautions attach to this formula.
\emph{(i)} This is the Bessel drift of index $\delta=2$, not $\delta=3$:
the exponent $\tfrac12$ in $\hm(\mu,0^+)\propto\mu^{1/2}$ is half the
exponent $1$ of the classical harmonic function $x$ for Brownian motion
killed at the origin, the halving being exactly that of the persistence
exponent from $\tfrac12$ to $\tfrac14$ under one integration.
\emph{(ii)} Under $Q$ the coordinate $\mu$ is not autonomous --- by
\eqref{sh:eq:G} its drift depends on $\kappa$, hence on both coordinates
--- so the comparison is a statement about the \emph{form} of the drift in
a limiting regime, not an identification of processes.

The correct P\'olya statement is therefore not a Bessel index but a
recurrence class, and at that level it is unconditional: under $\PP$ the
credit is null recurrent, returning to its initial level a.s.\
(Theorem~\ref{sh:thm:noimm}) with infinite expected return time
(Corollary~\ref{sh:cor:mean}), whereas under $Q$ it is transient,
$D_s\to\infty$ a.s.\ (Theorem~\ref{sh:thm:singular}). Conditioning on
survival promotes the credit from the recurrent to the transient class ---
the P\'olya $2\mathrm{D}\to3\mathrm{D}$ transition --- with no literal
change of dimension anywhere. Dimension remains, as in
\S\ref{sh:subsub:scale}, the wrong order parameter for the state dynamics;
what survivorship alters is the recurrence class of the \emph{ensemble}.
\end{corollary}

\begin{theorem}[Locally ours, globally elsewhere; \textbf{P/C}]
\label{sh:thm:singular}
\emph{(i)} $Q(\tau=\infty)=1$.
\emph{(ii)} $\hm(X_s)\to\infty$ $Q$-a.s.\ as $s\to\infty$; moreover
$D_s\to\infty$ $Q$-a.s.\ \cite{GJW99}.
\emph{(iii)} $\liminf_sD_s=-\infty$ $\PP$-a.s.
\emph{(iv)} Hence $Q|_{\Fc_s}\ll\PP|_{\Fc_s}$ for every finite $s$, with
the explicit density \eqref{sh:eq:Qdensity}, while $Q\perp\PP$ on
$\Fc_\infty$: the survivor's universe is locally absolutely continuous
with respect to ours on survivors, and globally disjoint from it. The
exit from the measure class is completed only at $t=\infty$; at every
finite conditioning the bias is the bounded tilt of
Definition~\ref{sh:def:interp}.
\end{theorem}

\begin{proof}
(i) $Q(\tau>s)=\E_\PP[\one_{\{\tau>s\}}\hm(X_s)]/\hm(X_0)=1$ for every
$s$ by the martingale property. (ii) Under $Q$, $1/\hm(X_s)$ is a
nonnegative supermartingale:
$\E_Q[1/\hm(X_{s+r})\mid\Fc_s]
=\PP_{X_s}(\tau>r)/\hm(X_s)\le1/\hm(X_s)$; it therefore converges
$Q$-a.s., and since
$\E_Q[1/\hm(X_s)]=\PP_{X_0}(\tau>s)/\hm(X_0)\to0$, Fatou forces the a.s.\
limit to vanish, i.e.\ $\hm(X_s)\to\infty$; the sharper statement
$D_s\to\infty$ is part of the description of the conditioned process in
\cite{GJW99}. (iii) is contained in the proof of
Theorem~\ref{sh:thm:noimm} (the set $A$ there is unbounded, and on it
$D_u\le D_0+\mu_0u-\sigma\varepsilon u^{3/2}\to-\infty$). (iv) Local
absolute continuity is \eqref{sh:eq:Qdensity}; mutual singularity on
$\Fc_\infty$ follows since $\{\lim_sD_s=+\infty\}$ has $Q$-probability
one and $\PP$-probability zero.
\end{proof}

\subsubsection{The three modes of survivorship}\label{sh:subsub:modes}

\begin{theorem}[Modes of the interpolation; \textbf{P/C/S}]
\label{sh:thm:modes}
As $t\to\infty$ the survivorship interpolation exhibits exactly three
asymptotic modes across the driver classes of \S\ref{sh:sub:horizon}:
\begin{enumerate}[label=\emph{(\Roman*)},itemsep=2pt]
\item \emph{Saturation} \textbf{[P]} (favorable drivers,
$\PP(T=\infty)>0$): since $\{T>t\}\downarrow\{T=\infty\}$,
\[
\bigl\|\PP^{(t)}-\PP(\,\cdot\mid T=\infty)\bigr\|_{\mathrm{TV}}
\ \le\ \frac{2\,\PP(t<T<\infty)}{\PP(T>t)}\ \longrightarrow\ 0
\quad\text{on }\Fc_\infty,
\]
and every bias is bounded by $1/\PP(T=\infty)$: survivorship saturates at
a bounded, globally absolutely continuous reweighting. No new universe.
\item \emph{Transformation} \textbf{[C/A]} (critical drivers; the
driftless model): the limit is the pure Doob transform of
Theorem~\ref{sh:thm:Q}, locally absolutely continuous and globally
singular (Theorem~\ref{sh:thm:singular}), approached at the polynomial
rate of Proposition~\ref{sh:prop:rate}. The new universe is entered, but
only at infinity.
\item \emph{Concentration} \textbf{[S]} (adverse drivers, exponential
tails; e.g.\ Proposition~\ref{sh:prop:advdrift}): on fixed windows the
limit is the eigen-tilted transform of
Remark~\ref{sh:rem:noclock} with clock
$\theta=3\nu^2/8\sigma^2$, while at macroscopic scale the conditioned
path law collapses onto the deterministic optimal profile of
Proposition~\ref{sh:prop:advdrift}(iii):
$\sup_{v\in[0,1]}\bigl|\mu_{tv}/t-\tfrac{|\nu|}4v(2-3v)\bigr|\to0$ in
$\PP^{(t)}$-probability. Survivorship becomes a law of large numbers for
the rare: the survivor ensemble is asymptotically deterministic at
leading order.
\end{enumerate}
The three modes are in bijection with the horizon trichotomy
(Theorem~\ref{sh:thm:trichotomy}, Proposition~\ref{sh:prop:recurrent}),
and they constitute a strictly finer invariant: modes (II) and (III)
share $\PP(\Imm)=0$ but are separated by the interpolation --- by whether
the survivor's universe is approached polynomially, without a clock, or
exponentially, with one.
\end{theorem}

\begin{proof}
(I) is continuity of measure from above plus the displayed
total-variation estimate, whose numerator is
$2\PP(\{t<T<\infty\})$ and whose denominator is bounded below by
$\PP(T=\infty)>0$. (II) collects Theorems~\ref{sh:thm:Q}
and~\ref{sh:thm:singular} and Proposition~\ref{sh:prop:rate}. (III) is
the large-deviation sketch of Proposition~\ref{sh:prop:advdrift}(iii)
read under the conditional measure, together with
Remark~\ref{sh:rem:noclock}; elevating it to a full proof requires the
sharp prefactor asymptotics $\PP_x(T>t)\sim C(x)e^{-\theta t}$, which we
do not establish here.
\end{proof}

\begin{remark}[The verdict]\label{sh:rem:verdict}
The section's question can now be answered in one paragraph. The
survivorship interpolation $t\mapsto\PP^{(t)}$ is a monotone flow of
upward tilts (Lemma~\ref{sh:lem:monotone}), globally absolutely
continuous at every finite $t$, whose asymptotic character is the
sharpest invariant of the system and takes exactly three forms
(Theorem~\ref{sh:thm:modes}). The driftless model realizes the critical
form: the limit exists, is the pure Doob transform by the
persistence-harmonic function, carries no clock
(Remark~\ref{sh:rem:noclock}), is approached at rate $s/4t$ with the tail
exponent as linear-response coefficient
(Proposition~\ref{sh:prop:rate}), is locally equivalent on survivors to
the unbiased law yet globally singular
(Theorem~\ref{sh:thm:singular}), and differs in generator by the
universal survivorship drift \eqref{sh:eq:G}, equal to $\sigma^2/2\mu$
at the wall and carrying the credit from the recurrent to the transient
class (Corollary~\ref{sh:cor:bessel}). The gauge structure of
\S\ref{sh:sub:model} persists to the end: the initial data cancel from
the limit density \eqref{sh:eq:Qdensity} up to the normalization
$\hm(X_0)$, and the survivorship drift is universal. The event horizon
absent from state space is therefore real, but it is a feature of the
ensemble, not of the state: it is crossed by the conditioning parameter
$t$ --- which is to say, by survival itself --- and never by the system.
\end{remark}

% [Referee pass, Aug 21: a convergence with the companion that neither
% document states.  Recorded here because it is evidence about the
% apparatus rather than about either paper.]
\begin{remark}[The same refinement the companion performs, in the path
measure]\label{sh:rem:coarsenings}
One structural agreement with the companion is worth recording, since
it was reached independently on both sides and neither development
notes it.

The companion derives classical survivorship bias as the Accordance
Principle under \emph{two coarsenings}: the observer tilt flattened
from a graded weighting to a two-valued indicator, and that indicator
read off whole histories rather than off single attributes
\cite{Nexus, NexicFoundation}. Restore the grading and relocate the
bearer, and the classical bias refines into the principle.

The interpolation of this subsection performs the same two refinements
in the path measure. The crude object is
$\one_{\{T>t\}}/\PP(T>t)$ --- an indicator, borne by the whole
trajectory --- and the limit object is $\hm(X_s)/\hm(X_0)$, a graded
weight borne by the state at each time. So the tilt is regraded from
indicator to continuum, exactly as the companion regrades it, and its
bearer moves from the completed path to the local state, exactly as the
companion moves it from the completed history to the attribute. What
Theorem~\ref{sh:thm:Q} proves is that the refinement is not a
modelling choice here but the limit the crude object converges to,
which is a stronger statement than the companion is in a position to
make for its own case.

Two things follow, and the second is the reason to record it.
Survivorship conditioning and observer conditioning are the same
operation refined in the same two directions, in the static and
dynamic registers respectively --- so the frailty ancestry of
\S\ref{gr1:sec}'s literature note and the companion's coarsening
result are two views of one fact. And an agreement of this shape is
evidence of the kind neither development can manufacture for itself:
the refinements were derived on one side from a measure-theoretic
limit and on the other from a typicality principle, with no common
step, and they arrived at the same pair \textbf{[S]}.
\end{remark}

% ---------------------------------------------------------------------
\subsection{Summary of the model section}\label{gr1:sub:summary}

% ---------------------------------------------------------------------


The system is a Kolmogorov diffusion in $(\mu,D)$ with exponential
killing; at exponent level it is the persistence problem for integrated
Brownian motion. Initial data are gauge: they act only through an
equivalent timescale $t_*$, and the one intrinsic coordinate is
$\kappa=\sigma^2D/\mu^3$ (\S\ref{sh:sub:model}). For the driftless model,
viability is impossible for every initial condition
(Theorem~\ref{sh:thm:noimm}), yet life expectancy is infinite:
$\PP(T>t)\asymp(t_*/t)^{1/4}$, with Pareto($\tfrac14$) residual life,
quantiles growing as the fourth power of rarity, and an effective Lindy
hazard $1/4t$ --- null recurrence, the two-dimensional case of the P\'olya
analogy (\S\ref{sh:sub:verdict}). No horizon exists there: the crossover
is a maximally soft quarter-power law smeared over orders of magnitude,
and recovery cost is sublinear in credit. Where genuine horizons exist,
the order parameter is the boundedness of the driver's scale function,
the horizon profile \emph{is} the normalized scale function
(Theorem~\ref{sh:thm:trichotomy}), the recurrent class is decided by the
stationary mean with mortal critical case
(Proposition~\ref{sh:prop:recurrent}), and the solvable unstable-OU
driver realizes the profile exactly as a Gaussian sigmoid
$\Phi(\mu_0\sqrt{2\nu}/\sigma)$ of width $\sigma/\sqrt{2\nu}$, a step
function only in the small-noise limit (\S\ref{sh:sub:horizon}). Under
adverse drift the survival tail is exponential with rate
$3\nu^2/8\sigma^2$, carried by a front-load-then-coast optimal
fluctuation. The destination is the survivorship interpolation
(\S\ref{sh:sub:interp}): a monotone flow of upward tilts $\PP^{(t)}$
whose critical limit is the pure Doob transform by the
persistence-harmonic function, approached without a clock at rate
$s/4t$ --- the tail exponent as linear-response coefficient --- locally
absolutely continuous on survivors yet globally singular, and generated
by the universal survivorship drift, equal to $\sigma^2/2\mu$ at the
wall: conditioning carries the credit from the recurrent to the transient
class, the P\'olya promotion effected by survival rather than by
dimension. Across
drivers the interpolation has exactly three modes --- saturation,
transformation, concentration --- in bijection with the horizon
trichotomy and strictly finer than the survival dichotomy. A horizon is
not a place. Where it exists in state space it is a ratio of scales; and
where it does not, it is real nonetheless --- located in the ensemble,
crossed not by the state but by the conditioning parameter, which is to
say by survival itself.

One channel stands outside that account and prefixes it. The schedule of
\S\ref{sch:sub:schedule} prices what must go right before the muffling
exponent can be held at all: failure incurs debt rather than discounting a
weight, the debt process is Galton--Watson with a criticality at mean
offspring one, and the completion weight is the least fixed point of
\eqref{sch:eq:fixedpoint} --- from which the companion's per-perturbation
contrast emerges as $Br/\lambda\ln10$ rather than being assumed. Its
consequence for the rest is structural: an unheld $x$ leaves the crest at
bare background, so the external channel --- the smallest term in
\eqref{tc:eq:hazard} --- is survivable only because the schedule
completed, and the state at which it completes is the initial datum
everything above takes as given.



% ---------------------------------------------------------------------
\begin{remark}[What the hazard rate is a projection of; author
dictation, \textbf{[$\Phi$]} on the companions]
\label{gr1:rem:aethicsubstrate}
The hazard rate this section models is not the ontology; it is a
projection of it. The richer object is Aethus-valued attribute
structure, and the relation is exactly the one the companion program
exhibits when exotic structure over Aethus values crunches to hazard
rates at the interface with survivorship mathematics \cite{Nexus}:
the rates are an intermediate quantity standing between the
attribute-level structure and the survival analysis, faithful to what
the analysis needs and silent about nearly everything else. On the
framework's own expectation it is the attribute level that a mature
successor engages directly, with rates retained as the projection one
computes when the question is survivorship and only survivorship.

A worked example fixes the picture. The event of nuclear war enters
the internal channel as one hazard term, yet the event itself is
heterogeneously composed of many Aethic attributes, the conditions
that jointly establish it, and an episode such as the Cuban missile
crisis reads in this vocabulary as the triggering of nearly all of
the needed preconditions with a terminal miss: followed in Aethic
time along the moving present, the total event's rate spikes as each
precondition lands and collapses at the end because the last stages
did not occur. The over-idealised model the dictation proposes treats
the preconditions as intervals, each carrying an entry rate and an
exit rate, with the rates of each interval heavily dependent on which
others are open, where the dependence is carried in the Aethic
jointness-weight sense rather than as causal linkage; that dependence
is what allows successor preconditions to fire on timescales of hours
once enough others stand open, though not in general. Nothing in this
is reductionist, and by the companion's fundamental theorem of
empirical Aethae it could not be: the Attribute Set design leaves the
decomposition chosen, so the nuclear-war event is itself the total,
one decomposition is useful to this model, others are or are not, and
no cut is more real than the event it cuts \cite{Aethus}.

One program-level question is entered here as open \textbf{[O]},
because this remark is its natural site: how much of what the rate
projection renders indiscernible is legible as golden causality when
the explicit attributes are examined. The accordance principle under
Golden conditioning already permits identifying, by the rarity of
observed events, that certain structures carry it intensely
(\S\ref{ra:sec}), and the attribute-level version of that audit is
future work with a stated instrument, and
Remark~\ref{sch:rem:notexhausted} sharpens the question into an
identification. Nothing in Part~II consumes
this remark; it prices the model's ontology and points outward, at
the companions' grades.
\end{remark}

\subsection{Where these results are used}\label{gr1:sub:consumers}

This section has been developed as a self-contained account of a hazard
process, and is readable as one. It is not, however, freestanding within
this work: each of its principal results has a consumer elsewhere, and
naming them is the quickest way to see why the section is here rather
than in a paper of its own.

\emph{The two-condition theorem} (\ref{gr1:thm:twocond}) is the paper's
most-used result. It supplies the regime trichotomy on which
\S\ref{gr:subsub:regimes} scopes the General/Golden distinction --- and
which, as recorded there, restricts the measure-theoretic claim to the
critical case. It supplies the dominance argument of
Remark~\ref{hr:rem:concession}, where the extinctionist option and the
settled-at-$x_\ast$ option are shown to carry viability probability zero
alike. And its anatomy factor is what \S\ref{ad:sec} exists to compute.

\emph{Technoanatomy} (Definition~\ref{gr1:def:anatomy}) is consumed
twice over. \S\ref{gr:sub:extraction} identifies extractive organization
with $\rho<0$, which is what converts the Extraction Hypothesis from a
vocabulary match into a claim about the sign of an exponent; and
\S\ref{ad:sub:toward} takes the derivation of $\rho$ from institutional
dynamics as its object, together with an account of why the machinery
there does not yet reach it.

\emph{The trapped-regime theorem} (\ref{gr1:thm:trapped}) does the work
in \S\ref{hr:sub:extinction}, where adverse anatomy is shown fatal
independently of the driver --- the step that makes the argument against
extinctionism turn on forfeiture rather than on comparative hazard.

\emph{The tolerance collapse} (\ref{gr1:prop:tolerance}) yields the
ordering constraint invoked in \S\ref{hr:subsub:stakes}: since the
affordable duration of adverse anatomy decays as $e^{-|\rho|x}$, anatomy
must be secured before capability accumulates, and a system that
attends to its arrangement late does so with the least margin it will
ever have.

\emph{The gauge results} of \S\ref{gr1:sub:coords} recur in a different
register in \S\ref{to:sub:selection}, where the selection of a
task-object among the many that individuate a given object is a gauge
fixing of the same shape: find the quantities the arbitrary choice
leaves invariant, and state the theory in those.

\emph{And the survivorship interpolation} of \S\ref{gr1:sub:cond} is
what \S\ref{gr:sec} is about in its entirety --- the Doob limit there
being the Golden Biosphere, and the local absolute continuity of
Theorem~\ref{sh:thm:singular} being the epistemic ceiling that
\S\ref{gr:sub:ceiling} turns against this work's own conclusions as
firmly as against its targets.

\emph{The schedule channel} (\S\ref{sch:sub:schedule}) is consumed in
three places and supplies one thing to each. It discharges the open rate
of Remark~\ref{gc:rem:hazardlower}, at the cost recorded there. It
supplies Remark~\ref{gc:rem:twolayers}'s $M$-valued fixed point in one
concrete case. And by Remark~\ref{sch:rem:initialdata} it produces the
distribution over initial data that \S\ref{gs:sub:limit} needs and does
not have --- which is a lead toward that section's principal open cost
rather than a discharge of it.

% ---------------------------------------------------------------------
\part{Golden Reasoning and the Route World}

\noindent Part~III states what conditioning on the Golden Point
licenses and what it forbids. General and Golden are distinguished
and the refutations run (\S\ref{gr:sec}); the epistemic ceiling fixes
what no finite observer certifies; retrocausal accordance is derived
as the instrument, with the finale paradox resolved as the
companion's designed relay (\S\ref{ra:sec}); and the Route World
carries the strong reading with its fences and its acting argument
(\S\ref{rw:sec}). What Part~IV consumes from here is an inference
discipline together with its limit, namely that reachable is never
certifiable.

\section{Golden reasoning}\label{gr:sec}
\setcounter{premise}{0} % GUARD --- pairs with \setcounter{premise}{5} in
% \S\ref{ac:sub:countermeasure}.  That section's premise prints $\Phi6$; these
% print $\Phi1$--$\Phi5$ unchanged.  Do not remove either guard.

\begin{quote}\small
The preceding sections were mathematics. This one is philosophy, and
says so. Its method is to let the mathematics \emph{limit the
possibility space} and then to reason philosophically inside what
remains. Two streams of load-bearing claims therefore run through it,
and are tagged accordingly: mathematical results carry the tags of
\S\ref{gr1:sec}, while philosophical premises carry \textbf{[$\Phi$]} and
are numbered $\Phi1,\dots,\Phi5$, each graded in the corpus's own
vocabulary --- stipulation, declaration, conjecture-grade, premise,
imported result --- and each priced with its falsifier in the ledger of
\S\ref{gr:sub:ledger}. Nothing is smuggled: where the argument turns on
a premise rather than a theorem, the premise is named, argued for, and
the thing whose exhibition would break it stated.

The thesis is that the fundamental categorical division among
world-scale living systems is not Pre-Civilization versus Civilization
but \emph{General} versus \emph{Golden} --- $\PP$ versus $Q$ --- and
that this division, unlike the other, is not a threshold anyone chose.
The ideal category is empty, but its practical form is not: a biosphere
of age $t$ is $\varepsilon$-Golden over windows of length $4\varepsilon
t$, so that Goldenness is bought by having lasted and by nothing else.
From this follow: the rejection of the Kardashev scale, not because its
thresholds are arbitrary (though they are) but because the model admits
a terminal stage gated by a variable the ladder cannot index; the rejection of the
violent-aliens thesis, conditional on one named premise; and a strict
epistemic ceiling which limits this section's own conclusions as much as
its targets'.
\end{quote}

\label{gr:sub:grades}

Two tag systems meet in this section and should not be confused.
\S\ref{gr1:sec} grade mathematical claims by how well established they
are: \textbf{[P]} proved, \textbf{[C]} classical or cited, \textbf{[A]}
conditional on the soft-wall Ansatz, \textbf{[S]} scaling-level,
\textbf{[O]} open. The companion program grades commitments instead by
\emph{what sort of thing they are} --- exposed postulate,
conjecture-grade, assumption, standing license, schema premise --- and
prices each with the thing whose exhibition would break it
\cite{Nexus, Ben26Found, Ben26Optimal}. The two are orthogonal
rather than rival, and this section uses both: the mathematical tags for
what it inherits from \S\ref{gr1:sec}, the companion's vocabulary for what
it adds. Each premise below accordingly carries its kind in its
statement, and \S\ref{gr:sub:ledger} assembles the whole into a ledger
in the companion's two-column format, on the companion's own principle
that a priced cost which could never be wrong is not a cost but a
decoration.

% ---------------------------------------------------------------------
\subsection{The referent problem}\label{gr:sub:referent}

\begin{premise}[The biospheric referent; \emph{methodological
stipulation}]\label{gr:phi1}
The appropriate unit of analysis for questions about the long-run
trajectory of world-scale living systems is the \emph{biosphere}: the
totality of the living and technical activity associated with a world,
considered as one system.
\end{premise}

The argument is one of referential stability across the relevant
timescale. ``Species'' fails backward: the system that became us was
not a species for most of its history, and speciation is not the joint
along which the relevant quantity --- total hazard-bearing organization
--- varies. ``Civilization'' fails backward for the same reason and
worse, since it names a phase rather than a system, so that using it as
the unit builds a phase transition into the vocabulary before any
analysis has occurred. ``Planet'' fails forward: it is a substrate, not
a system, and there is no reason the system should remain confined to
one. ``Biosphere'' is the only available term that refers to the same
thing at every point of the trajectory, which is exactly what an
asymptotic analysis requires.

This is not a terminological preference. A unit that changes referent
partway through cannot appear in a limit theorem. \S\ref{gr1:sec} concern
the $t\to\infty$ behavior of a single system; the vocabulary must
survive the limit or the theorems cannot be applied at all.

\begin{remark}[What this already costs Kardashev]\label{gr:rem:kardvocab}
Kardashev's scale \cite{Kar64} classifies \emph{civilizations} by energy
capture. Whatever its other merits, the choice of unit forecloses the
move this section makes: one cannot ask what the universe looks like
from the vantage of a long-survived biosphere if the vocabulary presumes
a civilization, since ``civilization'' carries semantic content ---
agency, society, purpose --- that the analysis is supposed to derive
rather than assume. The unit is doing work that ought to be done by
argument.
\end{remark}

% ---------------------------------------------------------------------
\subsection{General and Golden}\label{gr:sub:generalgolden}

\begin{definition}[General and Golden biospheres]\label{gr:def:gg}
A \emph{General Biosphere} is a biosphere considered under $\PP$, the
unconditioned law of \S\ref{gr1:sec}. A \emph{Golden Biosphere} is a biosphere
considered under $Q$, the survivorship limit --- the Doob $\hm$-transform
of \S\ref{gr1:sec}, the $t\to\infty$ limit of conditioning on $\{T>t\}$. The
family $\PP^{(t)}=\PP(\,\cdot\mid T>t)$ interpolating them is the
\emph{goldening} of the biosphere.
\end{definition}

\begin{remark}[The interpolation is a family; $Q$ is a measure]
\label{gr:rem:familyvsmeasure}
The distinction in that last sentence is substantive, and a specific
error follows from letting it slide. $\PP^{(t)}$ is a \emph{family}
indexed by $t$: each member is a different measure, and the index moves.
$Q$ is a single measure. So quantities defined by conditioning behave
differently along the two. Under any fixed measure the conditional
probability of a fixed event is a martingale in the filtration --- so
under $Q$ it is one. Watched \emph{along} the interpolation it is not,
because the conditioning tightens as $t$ advances and there is no single
measure for the martingale property to hold under; an extinction weight
so watched acquires a downward drift that is an artefact of the moving
index and not a fact about the process.

The framework's own reading of the difference is the cleaner one. Passing
to $Q$ is not a reweighting applied progressively as survival
accumulates; it is the conditioning made part of the base state, which in
the companion's terms is the statement of a future-referring attribute in
present books \cite{Aethus} --- available because statedness is a
property of the index rather than of the block's clock, and
Remark~\ref{gc:rem:notense} is the consequence drawn out. The Doob transform is what that operation looks like
analytically, and the reason $Q$ is a measure where $\PP^{(t)}$ is only a
family is that the former has already absorbed its condition and the
latter is still acquiring one.
\end{remark}

Three features of this definition do work that no threshold-based
division can do.

\emph{First, the split is categorical and unchosen --- in the critical
regime.} When $\PP(T=\infty)=0$, as it is under a recurrent driver with
$\rho>0$, the laws $Q$ and $\PP$ are \emph{mutually singular} on the
infinite-horizon $\sigma$-algebra: there is an event --- the credit
tending to infinity --- of $Q$-probability one and $\PP$-probability
zero. General and Golden are then not two regions of one population
separated by a line someone drew but different measures, the difference
established by a theorem rather than a stipulation. Pre-Civilization
versus Civilization has no such status, being a threshold in a
continuous variable with both the variable and the value contested.

The scoping is not a technicality and \S\ref{gr:subsub:regimes} states
it in full: outside the critical regime the singularity fails, and with
it the unchosenness. This is the section's strongest claim and it is
regime-specific.

\emph{Second, in the critical regime nothing occupies the Golden
category --- as an ideal.} By the mode classification of
\S\ref{sh:subsub:modes}, that regime is the critical one: the survivor's
law is approached but entered only at infinity. Every actual biosphere is somewhere on the interpolation
$\PP^{(t)}$, at finite $t$, and none is ideally Golden. This would make
the category inapplicable were it not for the practical form of
\S\ref{gr:subsub:practical}, which is occupiable, and which is what
Golden reasoning actually assumes of anything.

\emph{Third, the approach is quantified.} The bias field of \S\ref{gr1:sec}
gives, on a window of length $s$ at conditioning horizon $t$,
\begin{equation}\label{gr:eq:rate}
\beta_t(s)=\one_{\{T>s\}}\frac{\hm(X_s)}{\hm(X_0)}
\Bigl(1+\frac{s}{4t}+O(s^{2}/t^{2})\Bigr),
\end{equation}
so that goldening proceeds at rate $s/4t$ with the persistence exponent
$\tfrac14$ as its linear-response coefficient. ``The further out you
look, the more Golden it gets'' is not a slogan here; it is
\eqref{gr:eq:rate}.

\subsubsection{Golden reasoning in the three regimes}
\label{gr:subsub:regimes}

The two-condition theorem partitions the model, and Golden reasoning has
a different character in each cell. Since the apparatus of this section
was built inside one of them, the partition is stated before anything is
built on it.

\begin{center}\small
\begin{tabular}{@{}lll@{}}
\toprule
regime & $\PP(T=\infty)$ & what General/Golden becomes\\
\midrule
\emph{trapped}: $\rho\le0$ & $0$, $\E[T]<\infty$ &
concentration; $Q$ collapses onto one path\\
\emph{critical}: $\rho>0$, recurrent & $0$, $\E[T]=\infty$ &
\textbf{mutually singular; the split is categorical}\\
\emph{viable}: $\rho>0$, transient & $>0$ &
$Q\ll\PP$; an ordinary sub-population\\
\bottomrule
\end{tabular}
\end{center}

Three consequences, and the middle one costs this section something.

\emph{Trapped.} With adverse anatomy, viability is zero and life
expectancy is finite at frozen background
(Theorem~\ref{gr1:thm:trapped}), and survivors concentrate on the
optimal-fluctuation path: mode (III) of \S\ref{sh:subsub:modes}, where
the conditioned ensemble is asymptotically deterministic. There is a
General/Golden distinction, but the Golden ensemble has no randomness
left to study and no biosphere reaches it.

\emph{Critical.} This is where the rest of this section lives, and it is
the only cell in which the categorical claim holds. Death is certain but
life expectancy infinite; $Q$ is entered only at infinity and is
mutually singular with $\PP$; the interpolation $\PP^{(t)}$ is
non-trivial, and everything from \eqref{gr:eq:rate} through the
epistemic ceiling is a statement about it.

\emph{Viable.} With good anatomy and feedback, $\PP(T=\infty)>0$ and
conditioning on immortality is an \emph{ordinary conditional
probability}: $Q\ll\PP$ globally, total-variation convergence holds, and
the Golden Biosphere is a sub-population of the General one --- a set of
paths one could in principle be a member of, with no measure-theoretic
novelty and no unchosen categorical split. Mode (I), saturation.

The consequence for this section's central claim should be stated
without softening, because it is a real limitation. \emph{The
measure-theoretic recasting of the survivor category is available in
exactly one of the three regimes.} Where survival is genuinely
achievable, General and Golden differ by a reweighting rather than in
kind, and the categorical division reduces to the ordinary distinction
between a population and a favored part of it. The framework's most
distinctive structural claim thus holds precisely where immortality does
\emph{not}, and a reader who believes we inhabit the viable regime
should regard \S\S\ref{gr:sub:generalgolden}--\ref{gr:sub:ceiling} as
describing a case they think we are not in.

Two things blunt the loss without removing it. Which regime obtains is
decided by $\rho$ and by the driver's scale function, neither of which
is observable (\S\ref{gr:sub:ceiling}), so no one is presently entitled
to the belief that we inhabit the viable regime. And the critical regime
is not a knife-edge in the pejorative sense but the boundary the other
two are defined against: it is where the apparatus is sharpest precisely
because it is where the two failure modes meet.

Two wirings the partition now carries, added where later parts made
them load-bearing. \emph{Reachability sorts with the regimes.} In the
viable regime the Golden Point is reachable at positive probability
under $\PP$ itself --- conditioning is not what makes it reachable
but what makes it \emph{conditionable-on} where reachability
degenerates: critical, where the event is null yet modally alive and
the Doob limit is required; trapped, where there is nothing left to
condition on. \emph{And the partition prices the field's central
metric.} Ranking worlds by event-probability misorders them in both
directions across this table --- the undisturbed world minimizes
every event ledger and is already departed, while the critical
regime is condemned at probability one and is where everything the
survivorship theory computes resides
(Remark~\ref{ac:rem:losingcondition}, \S\ref{hr:sub:extinction},
\S\ref{gc:subsub:precipice}). The regime an event metric writes off
is the regime this section was built for.

\subsubsection{Practical and Ideal Golden Biospheres}
\label{gr:subsub:practical}

The ideal category is empty, but the questions one asks of a biosphere
are not questions about its infinite horizon. They concern windows ---
the next millennium, the next galactic orbit. On a window the
distinction collapses to a tolerance.

\begin{definition}[Practical and Ideal]\label{gr:def:practical}
The \emph{Ideal Golden Biosphere} is $Q$. A biosphere of age $t$ is a
\emph{Practical Golden Biosphere to tolerance $\varepsilon$ on windows
of length $s$} if $\PP^{(t)}$ and $Q$ agree to within $\varepsilon$ on
the $\sigma$-algebra $\mathcal F_s$.
\end{definition}

\begin{proposition}[Practical Goldenness is a ratio, not a level;
\textbf{S}]\label{gr:prop:ratio}
By the bias field \eqref{gr:eq:rate}, the density of $\PP^{(t)}$ with
respect to $Q$ on $\mathcal F_s$ is $(1-s/t)^{-1/4}(1+o(1))$, so the two
laws deviate at order $s/4t$. Hence
\begin{equation}\label{gr:eq:practical}
\text{$\varepsilon$-Golden on windows of length } s
\qquad\Longleftrightarrow\qquad
t\;\gtrsim\;\frac{s}{4\varepsilon},
\end{equation}
and a biosphere of age $t$ is $\varepsilon$-Golden for every question
concerning a window shorter than the \emph{Golden horizon}
$s_\varepsilon(t)=4\varepsilon t$. The tag is \textbf{[S]}: the $s/4t$
scale follows from the leading term of \eqref{gr:eq:rate}, but the exact
total-variation constant depends on the $o(1)$ structure and is not
established here.
\end{proposition}

\begin{remark}[The invariant]\label{gr:rem:invariant}
The Golden horizon is a fixed fraction of age:
$s_\varepsilon(t)/t=4\varepsilon$, independent of $t$. \emph{Every
surviving biosphere is Golden over a constant fraction of its own
history, and the fraction is four times one's tolerance.} There is no
threshold in strength, in energy, in technology, or in any state
variable whatever; the only thing that buys Goldenness is having lasted,
and what it buys is proportional. To be $1\%$-Golden over the coming
millennium requires twenty-five thousand years behind one.

This sharpens \S\ref{gr:subsub:kardashev}: the objection to Kardashev is
not merely that its thresholds are placed arbitrarily, but that the
operative threshold does not live in the space it measures.
\end{remark}

\begin{remark}[The premise of Golden reasoning is weaker than it looks]
\label{gr:rem:weakpremise}
This answers the natural suspicion that Golden reasoning assumes
something optimistic. It does not. By \eqref{gr:eq:practical}, practical
Goldenness to any fixed tolerance follows from age alone; and age is
exactly what the question already supposes. To ask what our descendants
are like across any substantial future is already to condition on there
\emph{being} descendants across that future, and that conditioning by
itself confers practical Goldenness over windows proportional to the
horizon supposed. The assumption is not that the future goes well. It is
that there is one --- and the model converts that supposition, and
nothing further, into the whole apparatus.

The corollary for the present is deflationary and should be stated
plainly: we are not practically Golden to any interesting tolerance, and
Golden reasoning does not claim we are. It describes the ensemble our
descendants would be drawn from in those futures where they exist.
\end{remark}

\begin{premise}[Golden reasoning is a vantage, not an obligation;
\emph{scope declaration}]\label{gr:phi2}
Golden reasoning is the practice of re-analyzing questions about
biospheres by taking $Q$ as the reference measure rather than $\PP$. It
is \emph{inferential}: it licenses conditional expectations, not duties.
\end{premise}

\noindent $\Phi2$ governs what the apparatus outputs; it is silent on the
separate question of how an agent grounds duties the apparatus never
issued, for which see the selection rule of
\S\ref{ac:sub:countermeasure}.

\begin{remark}[Why this is not longtermism]\label{gr:rem:notlongtermism}
The objection that any future-weighted reasoning collapses into ``act in
the interest of the future, as you happen to conceive it'' does not
apply, and the reason is structural rather than rhetorical. A normative
longtermism requires future beings whose interests ground obligations.
Golden reasoning has no such beings available: $Q$ is a measure that
nothing occupies in the critical regime
(Definition~\ref{gr:def:gg}), where moreover $\PP(T=\infty)=0$. In the
viable regime the objection would have purchase and is met differently,
by $\Phi2$ alone rather than by the vacuity of the reference class. What $Q$ supplies is a
\emph{conditional distribution}, from which one may infer what
long-survived systems are like without thereby acquiring any duty toward
them. The distinction is the ordinary one between a likelihood and a
value, and it is preserved here by the mathematics rather than by
stipulation. Whether one ought to care about the far future is a
question this section does not address and does not need.
\S\ref{gs:ssub:notlongtermism2} closes the concession made above, by
exhibiting an ordering difference that survives a normative reading.
\end{remark}

% ---------------------------------------------------------------------
Three consumers of this distinction now sit downstream and should be
visible from here. The finitude tradition's dissolution runs
entirely on the Practical face: unbounded lifespan as a convergent
series of mortal moments --- hazard strictly positive at every
instant, forever --- is what deprives the tedium and finitude
objections of their target (Remark~\ref{ac:rem:questionneighbors}).
The Ideal pole's epistemic status is fixed by the ceiling: never
certified, only approached, which is why every operative claim in
this work runs through the Practical register
(\S\ref{gr:sub:ceiling}). And the invariant above is the
quantitative face of the Lindy step --- having lasted as the only
purchase, proportionality as the price --- consumed where the Route
World's evidence method is built (\S\ref{gc:sub:lindy}).

\subsection{Technostrength and technoanatomy}\label{gr:sub:coords}

\begin{definition}[The two coordinates]\label{gr:def:coords}
\emph{Technostrength} is the muffling exponent $x_t$: the capacity the
biosphere has accumulated, positive when its own organization holds its
hazard below background and negative when it holds it above.
\emph{Technoanatomy} is the exponent $\rho$ of
Definition~\ref{gr1:def:anatomy}, the sign of
$-\partial_x\log h^{\mathrm{int}}$: whether acquiring capacity lowers or
raises the biosphere's hazard \emph{to itself}.

The two are of different kinds, and keeping them apart is the whole of
what follows. Technostrength is a \emph{state coordinate}: it
fluctuates, and by Theorem~\ref{tc:thm:half} every biosphere ---
however well constituted --- spends exactly half of logarithmic time
with $x<0$. Technoanatomy is a \emph{structural property}: it does not
fluctuate at all, it says what capacity does when acquired, and
$\rho<0$ names the biosphere whose own strength is turned against it,
which is what Definition~\ref{gr1:def:anatomy} means by eating oneself.
\emph{Adverse strength} is $x<0$; \emph{adverse anatomy} is $\rho<0$,
and only the second is a defect.
\end{definition}

Two results of \S\ref{gr1:sec} are worth restating in this vocabulary because
they are counterintuitive.

\emph{Technostrength is not cumulative --- in the critical regime.}
Under a recurrent driver $x$ oscillates on scale $t^{3/2}$ and returns
to every level infinitely often: it is not a stock that accumulates but
a magnitude that fluctuates, so a framework treating it as monotonically
increasing has the wrong kind of object. Under a \emph{transient} driver
it does accumulate, and the ladder picture is locally right --- which is
the concession \S\ref{gr:subsub:regimes} makes precise and prices.

\emph{Every biosphere is adversely configured half the time.} By
Theorem~\ref{tc:thm:half} the set $\{t:x_t<0\}$ has logarithmic density
exactly $\tfrac12$, for every parameter value and every biosphere.
Adverse \emph{strength} is therefore not a defect distinguishing some
biospheres from others but a condition all of them occupy half of
logarithmic time --- which is precisely why it cannot be what
``bad anatomy'' means, and why the two must be kept apart.

\subsubsection{What conditioning does to each coordinate}
\label{gr:subsub:measured}

\paragraph{Numerical methods.} All figures reported in this subsection,
in \S\ref{gr:subsub:anatsel}, in Remark~\ref{hr:rem:concession}, and in
the numerical check under Proposition~\ref{gr1:prop:trichotomy} were
produced by direct Monte Carlo on the two-channel model, using the exact
Gaussian update of Lemma~\ref{sh:lem:exact} for the driver-credit pair
and forward Euler for the background exponentials, at step
$\Delta=10^{-3}$. Conditioning on $\{T>t\}$ was performed by rejection
--- particles were killed by the Cox rule and the survivors reweighted
uniformly --- so the conditioned estimates are $\PP^{(T)}$ estimates and
not $Q$ estimates, exactly as the column headings below say. Standard
errors on the tabled probabilities are binomial and never exceed
$0.004$; the median-$x$ and quantile-$\rho$ entries carry bootstrap
errors of comparable size. \emph{Every $\PP(\rho>0\mid\text{survived})$
figure is prior-dependent}: $\rho\sim\mathrm{Unif}(-1,3)$ was chosen to
place a quarter of the mass on adverse anatomy and has no other
justification, and the $0.75\to0.84$ movement should be read as the
direction and rough magnitude of selection under that prior rather than
as a calibrated number. Code and seeds are available on request.

The following are simulation results for the two-channel model of
\S\ref{gr1:sec} at $\sigma=1$, $\epsilon=0.05$, $n=3$, $b=3$, $\rho=2.5$,
$k=6$, conditioning horizon $T=4$, $N=1.2\times10^{5}$ particles.

\begin{center}\small
\begin{tabular}{@{}lrr@{}}
\toprule
& General ($\PP$) & Conditioned ($\PP^{(4)}$)\\
\midrule
$\PP(x<0)$ (adverse strength) & $0.501$ & $0.0037$\\
median $x/\sigma T^{3/2}$ & $-0.001$ & $+1.085$\\
$\PP(\rho>0)$, $\rho\sim\mathrm{Unif}(-1,3)$ & $0.750$ & $0.839$\\
$5\%$ quantile of $\rho$ & $-0.80$ & $-0.44$\\
\bottomrule
\end{tabular}
\end{center}

\noindent The two $\rho$ rows depend on the prior $\mathrm{Unif}(-1,3)$
placed on anatomy, which is a choice and not a result; the $x$ rows do
not.

\begin{remark}[Truncation versus rescaling]\label{gr:rem:truncresc}
Two numbers, and they differ in kind rather than in force.
Adverse strength is suppressed by a factor of $137$ --- abruptly, and
within a few multiples of the intrinsic timescale. Technoanatomy is
selected too, but \emph{gradually}: $\PP(\rho>0)$ moves from $0.75$ to
$0.839$ at the same horizon, reaching $0.909$ only by $T=6$, with the
negative tail truncating steadily ($5\%$ quantile $-0.80\to-0.44\to
-0.14$).

The asymmetry is structural rather than incidental, and it follows from
what the two objects are. Strength is a state coordinate, so
conditioning pushes it out of a half-space it can and does re-enter ---
fast, and reversibly. Anatomy is a fixed structural parameter, so it can
be selected only through the accumulated hazard cost of never changing
--- slow, and irreversibly. Anatomy is the slower and the more decisive:
by Theorem~\ref{gr1:thm:trapped}, adverse anatomy is eventually fatal
with certainty, while adverse strength is a condition every biosphere
occupies half of logarithmic time and survives.


\end{remark}

\subsubsection{Does conditioning select technoanatomy?}
\label{gr:subsub:anatsel}

The measurement above concerns technostrength, a state coordinate; the
claim the argument needs is about technoanatomy, a structural parameter,
which therefore cannot be selected \emph{within} a trajectory. It can be
selected across an ensemble, and that is directly computable: place a
prior on $\rho$, evolve, and reweight by survival. Taking
$\rho\sim\mathrm{Uniform}(-1,3)$, so that $\PP(\rho>0)=0.75$ a priori:

\begin{center}\small
\begin{tabular}{@{}rrrr@{}}
\toprule
$T$ & $\PP(\text{survive})$ & $\PP(\rho>0\mid\text{survived})$ &
$5\%$ quantile of $\rho$\\
\midrule
$2$ & $2.8\times10^{-3}$ & $0.745$ & $-0.80$\\
$4$ & $2.2\times10^{-4}$ & $0.839$ & $-0.44$\\
$6$ & $1.9\times10^{-4}$ & $0.909$ & $-0.14$\\
\bottomrule
\end{tabular}
\end{center}

\noindent Survivorship does select for positive technoanatomy, and the
adverse tail truncates steadily. But the selection is \emph{gradual}
where the strength truncation was abrupt: adverse strength falls by a
factor of $137$ at $T=4$, while $\PP(\rho>0)$ has moved only from
$0.75$ to $0.84$ at the same horizon. Anatomy is the slower and the more
decisive: by Theorem~\ref{gr1:thm:trapped} an adverse $\rho$ is
eventually fatal with certainty, while adverse strength is a condition
every biosphere occupies half of logarithmic time and survives.

\begin{remark}[The coordinates stay separable]\label{gr:rem:separable}
Strength and anatomy cannot fuse under conditioning, and the reason is
structural rather than measured: $\rho$ is a parameter of the model and
$x$ a functional of the path, so no amount of conditioning makes one a
function of the other. They therefore remain distinct diagnostic
questions rather than two aspects of a single ``advancement'' --- which
is exactly what a one-dimensional scale cannot represent.
\end{remark}

% ---------------------------------------------------------------------
\subsection{What Golden reasoning refutes}\label{gr:sub:refutes}

\subsubsection{The Kardashev scale}\label{gr:subsub:kardashev}

The usual complaint against Kardashev's scale \cite{Kar64} is that its
thresholds --- planetary, stellar, galactic energy capture --- are
arbitrary. That complaint is correct but weak, since a scale may be
useful with arbitrary units. Golden reasoning supports three stronger
objections, of which the first is a theorem.

\emph{(i) The terminal stage exists, and is gated by a variable the
scale cannot see.} In the driftless model $\Lambda_\infty=\infty$
almost surely and no absorbing state of safety exists; in general one
does, and locating it is what makes the objection bite. By
Theorem~\ref{gr1:thm:twocond},
\[
\PP(\text{viable})=\one\{\rho>0\}\cdot
\frac{s(\mu_0)-s(-\infty)}{s(+\infty)-s(-\infty)},
\]
so there \emph{is} an absorbing state of safety --- and reaching it
requires two conditions, of which Kardashev's scale tracks at most one.
The escape factor is the accumulation of capacity, which energy
capture does index, however crudely. The anatomy factor is not on the
scale at all: $\rho$ does not appear in energy capture, is not
monotone in it, and cannot be inferred from it.

The objection is therefore not that the ladder has no top but that
\emph{climbing it is neither necessary nor sufficient for reaching the
top}. Worse, by Remark~\ref{gr1:rem:againstsing}, a biosphere with
$\rho<0$ that climbs fast does not merely fail to arrive: its hazard
diverges superexponentially in the very capacity the scale is
measuring. A scale indexed on capture alone cannot distinguish the
ascent that arrives from the ascent that is a fall. That last point is
the vulnerable-world hypothesis \cite{Bostrom19} restated with a
continuous parameter in place of a discrete draw, and
\S\ref{gr:subsub:xrisk} sets out what is and is not shared.

\emph{(ii) The indexed quantity is not monotone.} Technostrength is
stationary and recurrent in logarithmic time
(Definition~\ref{gr:def:coords}). A ladder presumes a quantity that
ratchets. Energy capture may of course ratchet while technostrength does
not --- but then energy capture is not tracking the quantity that
governs survival, which returns us to (iii).

\emph{(iii) A magnitude index cannot represent categorical selection.}
By Remark~\ref{gr:rem:truncresc}, survivorship truncates technoanatomy
while merely rescaling technostrength. A one-dimensional index of
magnitude has no coordinate in which to express a deleted half-space. It
is not that Kardashev measures the less important variable --- that
claim is false, and was corrected above --- but that it measures the one
whose selection is quantitative, and is structurally unable to express
the one whose selection is categorical.

\begin{remark}[On Dyson spheres]\label{gr:rem:dyson}
Dyson's proposal \cite{Dys60} is often read as the natural
Type~II signature. Under Golden reasoning it is better read as a
\emph{hypothesis about technoanatomy} disguised as one about
technostrength: it asserts not merely that a biosphere could capture
stellar output but that a long-lived one would be organized so as to
want to. That assertion is exactly what \S\ref{gr:sub:extraction}
disputes, and it should be recognized as an assertion rather than an
extrapolation.
\end{remark}

% [Placement pass, Aug 21: the expansionist selection story engaged.]
\begin{remark}[The expansionist mirror: grabby selection]
\label{gr:rem:grabby}
One modern argument deserves separate notice because it reaches the
Kardashev picture from selection reasoning --- this framework's own
instrument. The grabby-aliens model observes that loud, expansionist
civilizations dominate what is \emph{visible}, and runs anthropic
timing arguments under that weighting \cite{HMMP21}. Read beside this
section, that is a rival conditioning: weight by volume occupied and
the typical survivor is an expander; weight by persistence and the
two-condition theorem makes expansion-rate orthogonal to the factor
that decides. The two conditionings are not in contradiction --- one
asks what is seen, the other what survives --- but they license
opposite extrapolations from the same record, and conflating them
reproduces the Kardashev error at the level of selection effects:
visibility-weighting smuggles in the assumption that the seen and the
persisting are the same ensemble, which is exactly what the trough
channel denies. The framework's quarrel is therefore not with the
grabby model's arithmetic but with any normative or predictive reading
of it that takes loudness for viability; \S\ref{gr:sub:ceiling} bounds
what either conditioning can claim to certify.
\end{remark}

\subsubsection{The violent-aliens thesis}\label{gr:subsub:violence}

The thesis in view holds that intelligence evolves preferentially in
predators, hence that advanced biospheres are likely to be predatory or
violent. We engage the position rather than any proponent.

\begin{quote}\small
``It was the saying of Bion, that though the boys throw stones at frogs
in sport, yet the frogs do not die in sport but in earnest.''
\hfill --- Plutarch, \emph{Moralia} \cite{Plutarch}
\end{quote}

\noindent The line is Bion of Borysthenes', preserved by Plutarch, who
deploys it against hunting for sport; it has since been borrowed by
astrobiological writing as the image of what a technologically superior
biosphere might do to an inferior one. It is the right epigraph for this
subsection because it names the assumption at issue more exactly than
the conquest analogies usually offered. The serious form of the
violent-aliens thesis is not that an advanced biosphere would hate us.
It is that we would be \emph{beneath its consideration} --- that its
power could be spent across the boundary between itself and us
casually, without the expenditure ever rising to the level of a
decision. The stones are thrown in sport. That is the claim, and any
reply pitched at malice misses it.

The image also fixes what the reply must not do. It must not argue that
such a biosphere would be kind, since kindness is not what the thesis
denies; nor that it would have outgrown violence, since sport is not
violence in the relevant sense. What it must do is address the
\emph{scaling}: whether a biosphere can be arranged so that power is
spendable in that manner and still arrive at the scale where the
spending would matter. Our own route to the present framework
began from exactly this observation, three years before the model
existed, and \S\ref{gr1:sub:anatomy} records the entry.

Before the model is brought to bear, a further placement is needed, because the \emph{conclusion} here is the one most likely to be mistaken
for a new one and it is not. The conclusion --- that anything old enough to encounter must be
peaceable --- is a standing argument in SETI, sometimes called the
\emph{technological adolescence} argument, and it long predates this
paper: a civilization reaching the capacity to destroy itself must
either become socially responsible or fail to persist, so whatever
survives has already made that transition. Cooper documents the
argument and its currency, and also supplies the standing rebuttal: that
it conflates altruism with amiability, that it is typically asserted
rather than derived, and that it presumes a trajectory of moral
improvement for which there is no evidence \cite{Cooper19}. Nothing
below claims the conclusion as new. What is at issue is whether a
mechanism can be supplied that survives Cooper's objection.

The inference under dispute has the form: predatory origin
$\Rightarrow$ predatory maturity. Golden reasoning denies the
implication, on the grounds that it neglects the conditioning that must
be applied before any encounter can be contemplated. Any biosphere one could meet is old; any biosphere
that is old is drawn from $\PP^{(t)}$ for large $t$, not from $\PP$; and
$\PP^{(t)}\to Q$, under which technoanatomy is truncated. Whatever
selection operated at origin, a second and stronger selection operates
over the trajectory, and it is the second that determines what survives
to be encountered.

\setcounter{premise}{3} % GUARD --- the text names this $\Phi4$; document order
% would otherwise print $\Phi3$.  Pairs with the {2} guard at the Extraction
% Hypothesis and the {4} guard at Rarity of origins.
\begin{premise}[Origin does not transfer; \emph{premise}]
\label{gr:phi4}
Facts about the selective regime that produced intelligence in a
biosphere do not, without further argument, transfer to the selective
regime that governs that biosphere's persistence at world scale. The two
regimes act on different units, at different scales, over different
timescales.
\end{premise}

$\Phi4$ is not exotic; it is the ordinary observation that a trait
adaptive for an organism in a population need not be adaptive for a
biosphere in a universe. What the model adds is that the second regime
is not merely different but \emph{stronger} in the relevant sense: it
deletes a half-space of configurations, and it does so with a force
measured in \S\ref{gr:subsub:measured}.

What the model contributes to the frogs case in particular is a reply
that never mentions disposition. A biosphere that can spend power
casually across a self/other boundary is one that maintains such a
boundary and enforces it --- and by \S\ref{gr:subsub:boundary} the
enforcement of a boundary is interface-like while the domain it encloses
is volume-like, so the arrangement carries $\rho<0$ by
\eqref{gr:eq:rhoextract}. By Corollary~\ref{gr:cor:extractcost} its
viability probability is then exactly zero. The stone-throwers are not
ruled out because they would have grown kind; they are ruled out
because the organization that permits stone-throwing at that scale is
the organization the scaling argument says cannot reach that scale. The
frogs survive not by the boys' mercy but by the boys' having imploded
before attaining the requisite size --- which is, in a form Bion would
have recognized, the square--cube law again.

Three further features distinguish this from the folk argument, and they
are exactly the three points Cooper's rebuttal presses on. \emph{It is not
a moral claim.} What is selected is $\rho$, the sign of
$-\partial_x\log h^{\mathrm{int}}$ --- a structural property of how a
biosphere is arranged, carrying no implication that its members are
kind, and applying identically to a biosphere with no members in the
relevant sense at all. The conflation of altruism with amiability has
nothing to attach to. \emph{It is not asserted but measured.}
\S\ref{gr:subsub:anatsel} exhibits the selection on $\rho$ and finds it
\emph{gradual} --- $0.75$ to $0.91$ over the horizons computed ---
where the folk argument implicitly treats the filtering as complete.
\emph{And it is not a trajectory.} Nothing here says biospheres improve;
$\rho$ does not change at all in the fixed model, and what conditioning
does is reweight an ensemble, not move a system along a path. The
selection is Darwinian in form rather than developmental, which is
what the folk version gets wrong and Cooper is right to object to.

This much, however, only establishes that predatory maturity is not
\emph{inherited}. To conclude that it is positively \emph{selected
against}, one needs the identification of predation with negative
technoanatomy. That identification is a philosophical premise, and it is
the subject of the next subsection. Everything in
\S\ref{gr:sub:consequences} is downstream of it.

\subsubsection{Trajectory arguments for internal neglect}
\label{gr:subsub:neglect}

A third target belongs here, and it lies inside the biosphere rather
than out among the aliens. Longtermism in its strong form holds that
what most matters about our actions is their effect on the very long
run \cite{GM21}; and a recurring corollary within that literature is
that resources directed at the welfare of some populations do less for
the long-run future than the same resources directed elsewhere, so that
broad investment in the worse-off is not what a future-oriented agent
should fund. We engage the corollary as a position and name no one, as
with the violent-aliens thesis.

The position is already under sustained criticism, and it is worth being
exact about how the present objection differs, because the difference is
what makes it worth adding. The standing critiques are \emph{external}:
they fault longtermism on ethical and political grounds --- as an
ideology with damaging real-world effects and a flawed underlying moral
theory \cite{Crary23}, and as a doctrine whose cosmic accounting
licenses the discounting of present suffering \cite{Torres21}. Those
objections deny the framework's premises. The objection developed here
does the opposite: it \emph{grants} the long-run criterion entirely, and
shows the position failing by that criterion, on the framework's own
terms and in its own currency. If it succeeds it is therefore not an
alternative to the ethical critiques but a complement of a different
kind --- an internal result rather than an external one.

Stated at its strongest it is not a callous argument but an
optimization: interventions differ in long-run leverage, a
future-oriented agent should allocate to high-leverage interventions,
and if the leverage of some interventions is low then funding them is a
loss measured against the future. Nothing in that reasoning is invalid.
The question is whether the leverage calculation is complete.

Golden reasoning says it is not, and the missing term is technoanatomy.
No new premise is required, and this is the point worth making: a
position of this kind is $\Phi3$ with the boundary drawn \emph{inside}.
Writing off part of a system is the maintenance of a self/resource
distinction across an internal interface --- a region of the biosphere
to which protection is not extended --- which is structurally the same
object as the extractive boundary of \S\ref{gr:subsub:boundary}, and it
inherits the same scaling. Enforcement of an internal boundary is
interface-like; the interior it excludes is volume-like; so by
\eqref{gr:eq:rhoextract} the arrangement carries $\rho<0$, and
Corollary~\ref{gr:cor:extractcost} applies verbatim.

The consequence is that the position is self-undermining in its own
currency, and sharply rather than rhetorically. By
Corollary~\ref{gr:cor:extractcost}(ii) a biosphere with $\rho<0$ has
viability probability \emph{exactly zero} however much capacity it
acquires; by (i) its viability is zero where a $\rho>0$ biosphere's is
positive; and by
Remark~\ref{gr1:rem:againstsing} acquiring capacity faster makes its
position worse, since $x\to\infty$ with $\rho<0$ drives hazard to
infinity superexponentially. An argument that trades internal cohesion
for external capacity, offered in the name of the long-run future, is
therefore an argument for a shortened one --- not because neglect is
distasteful, which the model cannot see, but because the boundary it
requires is the structure the model identifies as fatal.

The thought experiment that makes this concrete is worth stating,
because it isolates the variable. Hold a society's technoanatomy fixed
at whatever value its present arrangements realize, and scale its
technostrength to galactic magnitude. If $\rho\ge0$, nothing goes
wrong. If $\rho<0$, the hazard does not rise gently but spikes as
$e^{|\rho|x}$: at every scale the same arrangement supplies more
micro-opportunities for the system to act against itself, and each
opens further ones. The elephant-sized ant of
Remark~\ref{gr:rem:namesclosure} is exactly this, and the square--cube
law is exactly why.

Four fences, because this is the section's most contestable
application. It refutes a \emph{class} of positions --- any that trades
internal cohesion for external capacity --- and not any author's wider
view, nor longtermism as such: a longtermism that valued internal
cohesion would be untouched, and the argument is in that sense a
constraint on longtermist reasoning rather than a refutation of it. It is downstream of $\Phi3$ and inherits its conjecture grade;
reject $\Phi3$ and this goes with it. And it establishes nothing about
what anyone ought to do, by $\Phi2$: what it establishes is that a
particular argument for neglect, evaluated by its own long-run
criterion, fails on that criterion.

% ---------------------------------------------------------------------
\subsection{A prediction outside cosmology}\label{gr:sub:crossdomain}

Most of what this section has established is a re-derivation. The
conclusions of \S\ref{gr:sub:refutes} were available before it --- the
Kardashev scale has long been thought crude, the peaceability of
long-lived civilizations is a standing SETI argument
\cite{Cooper19}, and the mechanism behind $\Phi3$ has precedents in
the collapse and existential-risk literatures
\cite{Tainter88,Bostrom19}. What the section offers is that they follow
from one premise with a named exponent, which is a unification rather
than a discovery.

Unification is a respectable contribution, but it carries an obligation
that a discovery does not. A framework which only reproduces what it was
built to reproduce is a reformulation, and the way to tell the two apart
is whether anything falls out that was not among the inputs. This
subsection states the clearest such consequence, in a domain the
framework was not designed for and where it can be checked.

\paragraph{The claim.}

The asymmetry measured in \S\ref{gr:subsub:measured} --- adverse
strength truncated abruptly, technoanatomy selected gradually --- is not
a fact about biospheres. Nothing in its derivation mentions one. It is a
consequence of conditioning a two-coordinate system on survival, where
one coordinate is a functional of the path and the other a parameter of
the model, and it should therefore appear wherever that structure does.

\begin{quote}
\emph{Cross-domain prediction.} In any population subject to
survivorship selection, carrying a fluctuating state coordinate
alongside a fixed structural parameter that governs whether the state's
growth raises or lowers hazard, conditioning on survival should
\textbf{truncate} the state coordinate --- deleting a region it can
otherwise re-enter --- and \textbf{shift} the structural parameter,
progressively and without deleting anything, the two effects being
comparable in magnitude and opposite in kind.
\end{quote}

Candidate domains are not exotic: firms with a fluctuating balance-sheet
position and a fixed governance structure; institutions with a
fluctuating resource base and a fixed constitutional arrangement;
lineages with a fluctuating population size and a fixed life-history
strategy. In each the two coordinates are separately measurable, old
members are selected from many, and the comparison the prediction asks
for is between surviving and full cohorts rather than between a model
and the sky.

\paragraph{What would count against it.}

Three outcomes would tell against the framework, and they should be stated plainly, since a prediction whose failure conditions are vague
is not doing the work this subsection is meant to do.

If the selection acted \emph{categorically on the structural parameter
and quantitatively on the state coordinate} --- the asymmetry inverted
--- the account of why the two behave differently would be wrong, since
that account turns on which of them can be re-entered.

If \emph{neither} coordinate showed selection in a population where
survivorship is known to operate, the framework would be describing a
mechanism that does not act at accessible scales, whatever its formal
standing.

And if the structural parameter showed selection \emph{as abrupt} as
the state coordinate, the distinction drawn in
\S\ref{gr:sub:coords} between a state and a structure would not be doing
the work claimed for it, and the argument of
\S\ref{gr:subsub:violence} --- which rests on that distinction to answer
the charge of positing a moral trajectory --- would lose its reply.

\paragraph{Two qualifications.}

The prediction is weaker than it may appear, in two respects that should
be recorded rather than left for a reader to find.

It is a claim about \emph{shape}, not magnitude. The framework does not
predict how large either effect should be in any particular domain,
since that depends on the hazard scale, the horizon, and the prior
spread of the structural parameter --- none of which the framework
supplies. What it predicts is that one effect is a truncation and the
other a shift, which is a qualitative signature and correspondingly
easier to satisfy by accident.

And it is not \emph{unique} to this framework. Any account on which
survival depends on a structural property will predict some selection on
that property; what is distinctive here is the predicted \emph{contrast}
in kind between the two coordinates, and only the contrast is
diagnostic. A domain exhibiting selection on both, with no difference in
character, would be consistent with much and would confirm little.

With those qualifications the prediction stands as the section's answer
to the question of whether its unification is more than a restatement.
We do not know the answer; the point of stating the prediction in this
form is that someone could find out.

% ---------------------------------------------------------------------
This prediction no longer stands alone, and the family it belongs to
is the paper's falsifiable perimeter, listed once so a reader can
audit it whole. Physical and statistical register: the present
cross-domain claim; the executed numerics, whose stated falsifiers
did not fire (\S\ref{app:numerics}); the mortality-deceleration
signature --- population hazard decaying as a universal $c/t$ under
criticality, against the plateau fingerprint of static frailty
\cite{GG01} --- carried in the standalone extraction's discussion;
and the Tainter route to measuring $\alpha-\beta$, the one posited
inequality, from collapse-era returns data
(Remark~\ref{gr:rem:betaalpha}, \cite{Tainter88}). Transfer
register: the hard-steps machinery re-run under the Golden target,
open and specified (\S\ref{gr:subsub:anthropic}, \cite{SBSD21}).
Reception register, scoreable rather than lab-testable and flagged
as such: the phase-transition prediction and the species-criterion
test it doubles as (Remark~\ref{gs:rem:valorvacancy},
Remark~\ref{hr:rem:speciesaware}). Registers differ; the discipline
does not --- each member is entered with the conditions of its own
failure.

\subsection{The Extraction Hypothesis}\label{gr:sub:extraction}

\setcounter{premise}{2} % GUARD --- the text names this $\Phi3$; see the guard
% at Origin does not transfer.
\begin{premise}[Extraction Hypothesis; \emph{conjecture-grade}]
\label{gr:phi3}
Organization built on unbounded extraction --- consuming the resources
of a locale and moving on --- has \emph{negative technoanatomy}:
$\rho<0$ in the sense of Definition~\ref{gr1:def:anatomy}, so that
internal hazard rises with acquired capacity rather than falling with
it.
\end{premise}

The premise is now a claim about the sign of an exponent rather than an
identification between a social form and a vocabulary, and this is worth
pausing on, because it is what the anatomy parameter buys. Before, one
had to be told what ``bad technoanatomy'' meant and then persuaded that
extraction was an instance. Now $\Phi3$ says something with a truth
value fixed by the model: that for extractive organization,
$-\partial_x\log h^{\mathrm{int}}<0$. The consequences follow
mechanically rather than interpretively, and they are severe.

\begin{corollary}[What $\Phi3$ costs an extractive biosphere; \textbf{P}
given $\Phi3$]\label{gr:cor:extractcost}
If extractive organization has $\rho<0$ then, by
Theorem~\ref{gr1:thm:trapped} and Theorem~\ref{gr1:thm:twocond}:
\emph{(i)} its viability is zero and, at frozen background, its life
expectancy is finite with $\PP(T>t)\le e^{-h_\ast t}$ for the floor
$h_\ast$ of \eqref{gr1:eq:xstar} (Theorem~\ref{gr1:thm:trapped}(ii));
\emph{(ii)} its viability probability is \emph{exactly zero}, whatever
feedback it achieves --- the anatomy factor in \eqref{gr1:eq:twocond}
annihilates the escape factor;
\emph{(iii)} it inhabits the trapped regime of
\S\ref{gr:subsub:regimes}, where none of the Lindy structure holds and
the survivor ensemble concentrates rather than diversifying;
and \emph{(iv)} its safest available configuration is not its strongest
but its \emph{weakest}, since $h\to B(1+e^{-b})$ as $x\to-\infty$ while
$h\to\infty$ as $x\to+\infty$.
\end{corollary}

\noindent Item (ii) is the one to hold onto. An extractive biosphere
does not merely have poor prospects; it has \emph{no} prospect of
viability, and acquiring capacity faster makes its position worse rather
than better. Item (iv) states the same thing from the other end and is
the sharpest form of the thesis: for a system that turns its strength on
itself, every increment of strength beyond $x_\ast$ is a net loss, and
the counsel of the model is to stop.

This is the section's single load-bearing philosophical premise. It is
not a theorem and is not presented as one. We argue for it.

\subsubsection{The argument from boundary cost}\label{gr:subsub:boundary}

Extraction requires a maintained distinction between the extracting
system and the extracted resource: a boundary across which value flows
one way. The boundary is not free. It must be enforced --- against
diversion, against defection, against the tendency of a system's own
subunits to treat the system as their resource, which is precisely what
a tumor does.

Now consider how the two sides of this ledger scale with technostrength,
and note that the anatomy parameter lets the scaling be written down
rather than gestured at. Let capacity be indexed by $x$. The domain over
which extraction occurs grows with reach and is volume-like,
$V\sim e^{\alpha x}$; the boundary across which enforcement is exercised
is interface-like, $S\sim e^{\beta x}$. The claim is
$\beta<\alpha$: volume outgrows surface. The unenforced fraction of the
domain is then $1-e^{-(\alpha-\beta)x}\to1$; measured against the
enforcement available to contain it, the unprotected extractive
interior grows as $V/S\sim e^{(\alpha-\beta)x}$, and internal hazard
with it:
% Derivation uses rho = -(alpha-beta): hazard scales with unenforced volume
% per unit of enforcement, matching gr1:def:anatomy's provenance note,
% sch:subsub:internalschedule (q = e^{alpha-beta}), ad:sub:toward, and
% Objection gr:obj:cheap.  All five sites agree.
\begin{equation}\label{gr:eq:rhoextract}
h^{\mathrm{int}}\;\sim\;e^{(\alpha-\beta)x}
\qquad\Longrightarrow\qquad
\rho\;=\;-(\alpha-\beta)\;<\;0 .
\end{equation}
So the argument does not merely conclude that extraction is bad
anatomy; it \emph{names the exponent}. Technoanatomy is minus the excess
of the extractive domain's growth rate over the boundary's, and the more
reach an extractive organization acquires per unit of enforcement, the
worse its anatomy is. Hence:

\begin{quote}
\emph{As a biosphere's technostrength grows, the cost of maintaining the
self/resource distinction grows faster than the capacity to maintain
it, and the shortfall is the exponent $\alpha-\beta>0$.}
\end{quote}

\noindent This mechanism is not original to the present work, and
\S\ref{gr:subsub:xrisk} records its ancestry: Tainter's account of
collapse from the declining marginal returns of complexity
\cite{Tainter88} reaches the same interior optimum from archaeological
evidence, and Bostrom's vulnerable-world hypothesis \cite{Bostrom19}
reaches a compatible conclusion from technological risk. What the
present treatment adds is the exponent \eqref{gr:eq:rhoextract} and the
hazard-theoretic consequences that follow from its sign.

Past some scale, the cheapest available target of an extractive
subsystem is the biosphere itself --- not through malice but through
gradient, exactly as a tumor exploits the body's own supply mechanisms
rather than building its own. Extraction at scale therefore turns
inward, and in the model's vocabulary this is precisely the
\emph{trough} channel of \S\ref{gr1:sec}: internal damage, unmediated by the
environment, unbounded above, and with adverse anatomy $\rho<0$ ensuring
that its costs outrun the crest channel's gains per unit.

% [Referee pass, Aug 21: the derivation names the exponent but the
% inequality carrying its sign was asserted in one clause.  Isolated
% here, with what would establish it and what would defeat it.]
\begin{remark}[The one posited inequality, and what would settle it]
\label{gr:rem:betaalpha}
Everything in \eqref{gr:eq:rhoextract} follows from the scalings once
$\beta<\alpha$ is granted, and that inequality is asserted in a clause.
It is the whole load, and it should carry its own accounting rather
than travel inside a derivation that looks computed.

\emph{What it says.} That enforcement is interface-like while the
extractive domain is volume-like --- so that reach purchased at rate
$\alpha$ can be policed only at rate $\beta<\alpha$, and the shortfall
is the exponent. The geometric intuition is the square--cube law and it
is genuinely suggestive; it is not an argument, because the objects
here are institutions rather than solids and nothing forces an
institutional cost to scale like a surface.

\emph{What would defeat it.} Enforcement mechanisms that scale with the
domain rather than with its boundary, of which there are real
candidates. Internalized norms are carried by the units themselves and
are volumetric by construction; distributed verification scales with
participants; and an immune system --- the biological analogue the
tumor image invites --- is precisely a policing mechanism distributed
through the volume it protects, which is why organisms defeat the naive
surface bound the image would impose on them. If any such mechanism
achieves $\beta\ge\alpha$ at scale, $\rho$ need not be negative for an
extractive organization and $\Phi3$ fails as a general claim. The
premise's force is therefore that no known arrangement has achieved it
\emph{at the scales in question}, which is weaker than the derivation's
presentation suggests and is the honest form.

\emph{What would establish it.} Not the geometry but the record.
Tainter's series are the right kind of evidence, since declining
marginal returns on complexity is $\beta<\alpha$ measured rather than
posited: an administrative cost growing faster than the domain it
administers is exactly the shortfall, and the archaeological cases
supply the growth rates \cite{Tainter88}. Fitting $\alpha-\beta$ to
those series --- rather than citing Tainter as a convergent conclusion,
which is all \S\ref{gr:subsub:xrisk} presently does --- would convert
$\Phi3$ from a premise into a measurement, and would supply the first
empirical estimate of a technoanatomy anywhere in this work. That is a
tractable piece of work, it is the most direct route from this
framework to data, and it is open \textbf{[O]}.

\emph{Scope of the exposure.} The two-condition theorem, the
trichotomy, and everything in \S\ref{gr1:sec} are untouched by this:
they take $\rho$'s sign as given and say what follows. What rests on
$\beta<\alpha$ is the identification of \emph{extractive organization}
as a case of $\rho<0$ --- that is, $\Phi3$ and the consequences
\S\ref{gr:sub:extraction} draws from it, and nothing else.
\end{remark}

\begin{remark}[The names were the mechanism]\label{gr:rem:namesclosure}
The argument just given is the square--cube law. That is the law from
which ``technostrength'' and ``technoanatomy'' were named --- initially
as a joke about giants. The joke turns out to be the mechanism: a
biosphere whose extractive reach has outgrown its capacity to remain one
system is an elephant-sized ant, and collapses for the same reason. The
model's brittleness exponent $\rho$ is the quantitative form of ``how
badly does the mismatch hurt''.
\end{remark}

\subsubsection{Objections}\label{gr:subsub:objections}

\begin{objection}[Bounded extraction]\label{gr:obj:bounded}
A biosphere might extract only from what is genuinely external and never
grow past the scale at which boundary cost bites.
\end{objection}

\noindent This is granted, and it is not a counterexample but a
concession of the point. A biosphere that halts its extractive growth at
a sustainable scale is not the unboundedly expanding entity whose
possibility is at issue; it is not a Kardashev trajectory. $\Phi3$
concerns \emph{unbounded} extraction, and the objection concedes that
unbounded extraction is unavailable.

\begin{objection}[Cheap boundaries]\label{gr:obj:cheap}
Sufficiently advanced technique might make boundary enforcement scale
with volume rather than surface.
\end{objection}

\noindent This is the serious objection and we do not refute it. In the
notation of \eqref{gr:eq:rhoextract} it is the claim $\beta\ge\alpha$:
enforcement scaling at least volumetrically, so that $\rho\ge0$ and the
whole corollary lapses. That is an empirical claim about the scaling of
coordination costs, it is where $\Phi3$ is falsifiable, and stating it
as an inequality between two exponents is the chief gain of the anatomy
parameter --- the objection and the premise now disagree about a
measurable number rather than about a description. We note only that
every known instance of large-scale organization --- organisms,
ecosystems, institutions --- exhibits $\beta<\alpha$, and that the
burden of the exception lies with whoever claims it.

\begin{objection}[Assuming the conclusion]\label{gr:obj:circular}
Defining bad technoanatomy as ``that which raises hazard'' and then
asserting extraction raises hazard makes the argument circular.
\end{objection}

\noindent The charge is now easy to answer precisely, and the anatomy
parameter is what makes it so. The model defines $\rho$ as
$-\partial_x\log h^{\mathrm{int}}$ and proves what follows from its
sign; that is a theorem about an exponent and mentions extraction
nowhere. $\Phi3$ claims that a particular form of organization realizes
a negative value of that exponent, and argues for it from boundary
scaling. Two distinct objects, one proved and one argued, and neither
borrowed from the other. The identification is supplied from outside,
by \S\ref{gr:subsub:boundary}, and could be false --- it would be false
if $\beta\ge\alpha$. The division of labor is exactly the one this section's method
requires: mathematics limits the possibility space, philosophy says
which region of it we occupy.

\begin{remark}[What would falsify $\Phi3$, and what would not]
\label{gr:rem:falsify}
Three falsifiers, and they are \emph{operational} in the companion's
sense --- checkable in domains that exist, rather than requiring a
capability the argument itself says no one has \cite{Ben26Optimal}.
$\Phi3$ fails if $\beta\ge\alpha$ in \eqref{gr:eq:rhoextract} --- if
enforcement is shown to scale at least volumetrically in reach; if an
extractive organization at scale is exhibited stable over timescales
long compared with its own doubling time; or if internal hazard is
observed to \emph{decrease} with technostrength in an extractive system,
which is $\rho>0$ measured directly and refutes the premise outright. The first
is the live question, and it is live in institutional, ecological and
organismal data rather than only in cosmology.

What would \emph{not} falsify it: the bare possibility that extraction
might turn out fine, or a hypothetical technology stipulated to make
boundaries cheap. The claim is about scaling, scaling is measurable
wherever large-scale organization exists, and a falsifier must therefore
exhibit the scaling rather than assert its negation.
\end{remark}

% ---------------------------------------------------------------------
\subsection{Consequences}\label{gr:sub:consequences}

The following are downstream of $\Phi3$ and inherit its status. Each is
falsifiable in principle; none is currently testable.

\emph{(a) Encountered biospheres are not extractive.} Not because
predation is impossible but because extractive organization is a
transient: it may be instantiated, but the conditioning under which any
encounter occurs selects against it. A predatory biosphere
``copy-pasted'' into the galaxy would be a configuration of the
$\PP$-ensemble, not the $Q$-ensemble, and would not persist to be met.

\emph{(b) Absence of megastructures is not evidence of absence.} If
$\Phi3$ holds, stellar-scale extraction is a signature of biospheres on
their way out, not of biospheres that lasted. Searches keyed to such
signatures \cite{Dys60} are therefore keyed to the wrong tail of the
distribution, and their null results carry less evidential weight
against the existence of old biospheres than is usually assumed.

\emph{(c) Old biospheres are Lindy --- and here the novelty is narrow,
so it is stated narrowly.} From the survival tail
$\PP(T>t)\asymp Ct^{-1/4}$ (\S\ref{gr1:sec}): residual life is
asymptotically Pareto$(\tfrac14)$ in units of attained age, and the
effective population hazard is $1/4t$. The longer a biosphere has
persisted, the longer it may be expected to persist, in exact
proportion --- while still departing with probability one.

That conclusion is \emph{not} new, and it would be a misrepresentation
to present it as such. Lindy longevity for technological civilizations
is an active subfield: Balbi and \'Cirkovi\'c model the termination rate
as a Weibull hazard and note that a decreasing rate yields a
Pareto-like fat tail of exactly Lindy character \cite{BC21}; Balbi and
Grimaldi investigate Lindy's law for technosignature longevity directly,
and show it can overturn the expectation that a first detection be
long-lived \cite{BG24}; Kipping and collaborators reach related
conclusions about the age of a first contact from lifetime distributions
\cite{Kip20}; and Ord surveys the mechanisms from which a Lindy effect
can arise at all, including one that conjures it from an uninformative
prior \cite{Ord23,Tal12}.

What this section adds to that literature is one thing only. In each of
those treatments the power law is \emph{parameterized}: the exponent is
a free shape parameter, chosen for tractability or fitted, and the
question studied is what follows from a given value of it. Here the
exponent is an \emph{output} --- $\tfrac14$, and it would have been
$\tfrac12$ had the noise entered the log-hazard rather than its
derivative (\S\ref{gr1:sec}). The claim is therefore not that biospheres are
Lindy, which is others', but that the Lindy exponent is fixed by the
smoothness of whatever drives the hazard, and is in that sense a
measurable structural fact rather than a modeling choice. Whether that
addition is worth much depends entirely on whether the driver's
smoothness is ever ascertainable, which this section does not claim.

\emph{(d) The age distribution of extant biospheres is heavy-tailed.}
Quantiles scale as the fourth power of rarity. Any sample of biospheres
found at a given epoch is dominated by the old, and moments of order
$\ge\tfrac14$ diverge: the sample mean age of such a population is an
artifact of the observation window, growing as $\mathcal T^{3/4}$.

\subsubsection{Relation to anthropic reasoning}\label{gr:subsub:anthropic}

Golden reasoning stands in an unusual relation to the observation
selection literature \cite{Car83,Got93,Bos02}, and the relation deserves precision because it is the section's clearest methodological
contribution.

Gott's Copernican argument \cite{Got93} derives a predictive
distribution for a system's remaining duration from the premise that the
observation time is uniformly distributed over the system's lifetime.
The Doomsday argument \cite{Car83,Bos02} proceeds similarly from a
reference class and an indifference principle. In both, the age
distribution is an \emph{input}, justified by a principle of
indifference.

In Golden reasoning the age distribution is an \emph{output}. The tail
$t^{-1/4}$ is not postulated; it is the persistence exponent of
integrated Brownian motion, which follows from where the noise is placed
in the hazard (\S\ref{gr1:sec}) and nothing else. The reference class is not
chosen; it is the conditioned ensemble $\PP^{(t)}$, which the model
constructs. The exponent could have been $\tfrac12$ --- it is, for a
mean-reverting driver --- and which it is depends on the smoothness of
the driver rather than on any epistemic principle.

The same contrast holds, and holds more sharply, against the
technosignature-longevity literature, which is the nearer neighbor and
is already model-based rather than indifference-based
\cite{BC21,BG24,Kip20}. Those treatments carry a hazard function with a
free shape parameter --- Weibull $k$, or a Pareto exponent $\alpha$ ---
and study the consequences of its value. The disagreement is therefore
not about method but about where a parameter comes from: there it is
chosen, here it is forced by the driver. Ord's survey is the useful
corrective in both directions \cite{Ord23}, since it shows how many
distinct mechanisms deliver a Lindy effect --- which is a warning that
observing Lindy behavior would badly underdetermine the mechanism
producing it, and hence that the exponent, not the effect, is where any
inferential weight can sit.

One further neighbor in this family should be placed because its
instrument is reusable. The hard-steps analysis --- Carter's
originally, given its modern Bayesian form by Snyder-Beattie,
Sandberg, Drexler and Bonsall --- infers from the timing of
evolutionary transitions, under observer-conditioning, that the
completed steps were individually improbable \cite{Car83, SBSD21}.
That is timing run as an instrument on a conditioned record, and
nothing in it is proprietary to the observer target: run under the
Golden target the same instrument reads transition timings as
evidence about route-narrowness --- the quantity the invested reading
prices (Proposition~\ref{ra:prop:investment}) --- and the hard-steps
literature's machinery transfers whole, with only the conditioning
attribute exchanged. Executing that transfer is a concrete open item
\textbf{[O]}, and it is the nearest thing this section has to a
ready-made empirical program built by others.

This does not vindicate the Copernican arguments; it relocates the
disagreement. Where they require a defense of indifference, Golden
reasoning requires a defense of the model. That is a better place for
the burden to sit, because a model can be wrong in identifiable ways.

\setcounter{premise}{4} % GUARD --- restores the sequence after the two guards
% above; the text names this $\Phi5$.
\begin{premise}[Rarity of origins; \emph{imported result}]
\label{gr:phi5}
Biospheres are rare: the difficulty lies in origination, not in
persistence past origination.
\end{premise}

$\Phi5$ is the rare-earth position \cite{WB00}, and it is not imported
from the general literature but from our companion program,
where it is derived rather than assumed: the Nexic framework
\cite{Nexus} argues that the astrobiological Copernican principle
is inconsistent with the mediocrity commitments that were supposed to
support it \cite{Ben26Found}, and re-reads the natural-historical record
accordingly \cite{Ben26Optimal}. Its interaction with the rest of this
section is worth making explicit, because the combination is more
coherent than either part alone. If origins are rare and survival is
Lindy, then the observed silence requires no filter ahead of us: it is
explained by the scarcity of starts, not by the destruction of the
started. Golden reasoning and rare-earth are complementary --- the first
says that what persists is not extractive and therefore not loud, the
second says there is little of it. Neither requires the Great Filter to
lie in our future \cite{Han98,SDO18}.

\subsubsection{Relation to the Nexic framework}
\label{gr:subsub:nexic}

The relation to that companion program is closer than a shared
conclusion, and stating it precisely is worth a subsection, because the
two halves turn out to partition the same problem.

\emph{The same move, in opposite temporal directions.} The Accordance
Principle constrains a realized natural-historical attribute $A$ against
its complement $A'$, with $H$ the production of observers, by
$P(H\mid A')/P(H\mid A)\lesssim P(A)/P(A')$: realized improbability must
be compensated by observer-relevance \cite{Ben26Found}. Golden reasoning
performs the same reweighting on a trajectory rather than an attribute,
the Doob density $\hm(X_s)/\hm(X_0)$ tilting the path measure by the
probability of the conditioning event downstream. Both are Bayes applied
to a selection event; they differ in what is conditioned on and in which
direction time is read. Accordance conditions on observers \emph{having
arisen} and reads backward over the record; Golden reasoning conditions
on survival \emph{to a future horizon} and reads forward over the
trajectory. The two are complementary rather than redundant:
retrodictive and predictive conditioning on the same kind of event.

\emph{The companion names the degenerate case.} This is not a
resemblance noticed from outside. The foundation paper itself identifies
survivorship reasoning as the coarsest limiting case of its own
constraint: restrict the state space to a single history's population
and degrade the graded observer-condition $H$ to a binary survival event
$S$, and the identity behind the principle reads
$P(A\mid S)/P(A'\mid S) = \bigl(P(S\mid A)/P(S\mid A')\bigr)\cdot
\bigl(P(A)/P(A')\bigr)$ --- the textbook survivorship correction
\cite{Ben26Found}. What the present section supplies is that limit
developed \emph{dynamically}: not a correction applied at one attribute
but a measure on paths, with a limit object $Q$, an exponent
$\tfrac14$, and an interpolation between the two.

\emph{Golden reasoning supplies what the companion declares itself
silent about.} The Optimal-Earth paper states as a principled limitation
that its framework is structurally silent on filters ahead: its
constraint conditions on observer-production, which our own case has
already secured, so no bound it issues bears on what becomes of
observers afterward \cite{Ben26Optimal}. That is precisely the gap this
section fills, because conditioning on $\{T>t\}$ is conditioning on a
\emph{future} event. The absence of a terminal stage, the Lindy residual
life, the $t^{-1/4}$ tail and the goldening interpolation are all
statements about what becomes of observers afterward. The two wings
therefore partition the timeline at the observer: Accordance behind,
Golden reasoning ahead.

\emph{One identification, flagged rather than assumed.} The connection
is a shared logical form plus one substantive identification, and the
identification is not free. The companion's cosmic weighting is a
primitive measure over partial information states; the $\PP$ of
\S\ref{gr1:sec} is a path measure induced by a
specific model. Reading the two as one object requires taking the
model's trajectory space as a family of partial states and its $\PP$ as
the cosmic weighting restricted to them --- a move neither paper
performs. Absent it, what is shared is the logical form, which is worth
stating and is not an entailment. Nothing in this section's results
depends on the identification being made.

\emph{An apparent clash, dissolved.} The Optimal-Earth Conjecture
\cite{Ben26Optimal} holds our record to be observer-optimal on
effectively any attribute one can interrogate, while \S\ref{gr:subsub:practical} holds
that we are not practically Golden to any interesting tolerance. These
do not conflict, and the reason is the temporal split just drawn: the
first conditions on observer-production and grades the realized past,
the second conditions on survival to a future horizon and measures
distance from $Q$. A record may be optimal in every respect Accordance
interrogates while its bearer sits at small $t$ on the goldening
interpolation. Optimality behind is not Goldenness ahead.

\subsubsection{Relation to the existential-risk and collapse
literatures}\label{gr:subsub:xrisk}

Two established positions occupy ground adjacent to $\Phi3$ and the
two-condition theorem, and both were found after those results were
written. They are recorded here because they narrow what is claimed as
new --- and because, on inspection, they support the premise rather than
threatening the paper.

\emph{Collapse from declining returns.} Tainter's account of societal
collapse \cite{Tainter88} holds that complexity is adopted as a
problem-solving strategy, that its marginal returns decline, and that
past a threshold further complexity is pure cost, at which point
collapse becomes likely. Structurally this is the U-shape of
Proposition~\ref{gr1:prop:trichotomy}(iii) arrived at from archaeology
thirty-six years earlier: a benefit-minus-cost curve in capacity with an
interior optimum, beyond which acquisition is net harmful. The
correspondence is close enough that the boundary-cost argument of
\S\ref{gr:subsub:boundary} should be read as a re-derivation of
Tainter's mechanism in hazard-theoretic terms rather than as an
independent discovery of it.

Three differences remain, and they are what this section adds. Tainter's
declining quantity is \emph{economic and energetic return}; ours is
\emph{hazard}, so the two curves are about different things even where
their shapes agree. Tainter's threshold is qualitative; ours is
$x_\ast$ of \eqref{gr1:eq:xstar}, in closed form and with a computed
floor $h_\ast$. And Tainter's account contains no analog of the
escape factor, so it cannot state the two-sided result of
Theorem~\ref{gr1:thm:twocond} --- that capacity growth is simultaneously
\emph{necessary} for viability and, absent anatomy, fatal.

\emph{The vulnerable world.} Bostrom's hypothesis \cite{Bostrom19} holds
that continued technological development will at some point attain
capabilities making civilizational devastation extremely likely unless
civilization exits what he calls the semi-anarchic default condition,
with stabilization requiring greatly amplified capacities for preventive
policing and global governance. The mapping into the present vocabulary
is direct and should be stated without hedging: \emph{the semi-anarchic
default condition is low technoanatomy, and Bostrom's stabilization
conditions are the raising of $\rho$}. Quantitative development of the
hypothesis has begun, using survival-function methods on the urn model
\cite{EF24}, which is the same class of object as
\S\ref{gr1:sub:verdict}.

What survives the comparison is narrower than the anti-acceleration
result of Remark~\ref{gr1:rem:againstsing} might suggest on first
reading, and is stated narrowly. The vulnerable-world hypothesis is an
\emph{existence claim about a discrete draw} --- that the urn contains a
black ball --- whereas $\rho$ is a \emph{continuous scaling relation}
with an exponent derived from boundary geometry. And the vulnerable-world
hypothesis is one-sided: development is dangerous. The two-condition
theorem is two-sided, requiring the same quantity that it makes lethal,
and the product structure of \eqref{gr1:eq:twocond} --- in which one
factor annihilates the other --- has no counterpart there.

\emph{Why this strengthens rather than weakens the section.} $\Phi3$ is
conjecture-grade and argued from a single line of reasoning. It now has
two independent precedents reaching a compatible conclusion from
unrelated evidence: one archaeological and economic, one
technological and policy-facing. A premise supported by convergence
across literatures that share no premises is in better standing than one
argued from scratch, and the ledger's grading of $\Phi3$ should be read
in that light --- still conjecture-grade, since none of the three
derivations is a proof, but no longer solitary.

% ---------------------------------------------------------------------
\subsection{The epistemic ceiling}\label{gr:sub:ceiling}

We close with a limitation which constrains this section's own
conclusions at least as much as its targets'.

\begin{theorem}[No finite observation certifies Goldenness]
\label{gr:thm:ceiling}
$Q$ is locally absolutely continuous with respect to $\PP$: on every
finite window the two measures are mutually absolutely continuous on the
survival event, with the explicit density of \S\ref{gr1:sec}, finite on
every finite window. They
become singular only on the infinite-horizon $\sigma$-algebra.
Consequently no observation of bounded duration can establish that a
biosphere is Golden rather than General, and no bounded evidence can
place a system on the interpolation $\PP^{(t)}$ at one $t$ rather than
another beyond the resolution the bias field permits.
\end{theorem}

\begin{remark}[What ``in principle'' may and may not mean here]
\label{gr:rem:ceilingscope}
The theorem's quantifier is \emph{bounded duration}, and prose elsewhere
in this work occasionally compresses that to \emph{in principle}. The
compression should be read narrowly, for two reasons.

\emph{It is a limit on inference, not on determinacy.} Nothing above says
there is no fact about whether a biosphere is Golden. It says no finite
record distinguishes the measures, which is compatible with the fact
being fully settled --- and in the Route World of
\S\ref{gc:subsub:routeworld} it is settled, since statedness is a
property of an Aethus and not of its bearer's access to it
(Remark~\ref{gc:rem:determinacygap}). Determinate and inaccessible is
this framework's normal condition, not a special pleading.

\emph{And the unavailability is one-sided.} The measures are singular on
the infinite-horizon $\sigma$-algebra, so the infinite record does
distinguish them --- but that record is available only to a biosphere
that persists to obtain it. A General biosphere departs at finite $T$ and
never holds it; a Golden one could, in a limit it never finishes
reaching. So the ceiling is not symmetric ignorance. \emph{No biosphere
ever learns that it is General}, and a Golden one learns it only
asymptotically --- the same shape as the fixation pole's unprovability in
\S\ref{sch:sub:schedule} and the affirmative answer's stability in
\S\ref{ac:sub:question}. This framework's uncertifiable propositions are
uncertifiable in the direction that would have been consoling.
\end{remark}

\begin{remark}[What now stands on the ceiling]\label{gr:rem:ceilingconsumers}
This subsection is short because the theorem is; its weight is
elsewhere, and the load map belongs at the site. Five results in
this work stand directly on the ceiling. The Precipice engagement's
third delta --- existential security decomposes, and
known-to-be-low is the half the ceiling forbids
(\S\ref{gc:subsub:precipice}). The anti-dogma delta of the
reduction-refusal --- dogma requires certifiable favor, and a
structure that proves no one certifies the favorable side cannot
function as a salvation guarantee (\S\ref{gs:sub:notreligion}). The
Kierkegaard consilience --- the third route is held rather than
possessed, commitment under proven uncertainty
(\S\ref{gs:sub:meaningplacement}). The node criterion's discipline
--- verdicts on commitment primitives are disciplined, never
certified (\S\ref{gs:sub:nodes}). And the arrival-timing refusal ---
no one certifies the occupant (Remark~\ref{gs:rem:arrivaltiming}).
Two companions complete the map: silent forfeiture is the ceiling's
factual twin --- what cannot be certified also does not announce
itself (Remark~\ref{hr:rem:silentforfeit}) --- and one statement is
owed here because it is consumed everywhere: \emph{reachable is
never certifiable}. In the viable regime the Golden Point is
reachable under the plain unconditioned law --- positive probability,
no conditioning required, the survivorship flow merely saturating
onto the honest conditional --- and the ceiling stands regardless: a
biosphere can be on the favorable side of the dichotomy while no
finite observation ever places it there. Reachability lives in the
model; certification is what the model withholds; and the pairing is
this framework's epistemology stated as measure theory --- local
equivalence on every window, singularity only at the horizon --- with
no translation step between the mathematics and the moral.
\end{remark}

\begin{remark}[What this costs each party]\label{gr:rem:evenhanded}
Against the violent-aliens thesis: no finite observation of a biosphere,
however extensive, could establish that it is predatory in the
asymptotically relevant sense, so the thesis cannot be confirmed by
encounter.

Against \emph{this section}: no finite observation could establish that
an encountered biosphere is \emph{ideally} Golden either. The claims of
\S\ref{gr:sub:consequences} are claims about a conditional
distribution, not predictions about any individual system, and they must
not be read as the latter.

The ceiling is, however, specific to the ideal. Practical Goldenness
(Definition~\ref{gr:def:practical}) is a finite-window property, and
Theorem~\ref{gr:thm:ceiling} does not forbid establishing it: mutual
singularity is an infinite-horizon phenomenon, and on $\mathcal F_s$ the
two laws are equivalent with bounded density. What obstructs the
practical version is not structure but sample size --- one would need an
ensemble of biospheres, and we have one. That is a contingent barrier
rather than a theorem, and the two kinds of limitation are worth keeping
apart: the ideal category is unknowable to any finite observer, the
practical one merely out of reach.

This symmetry is not a weakness of the position but a consequence of the
structure that the position is built on, and we would regard its absence
as a warning sign. An argument from survivorship that permitted
confident identification of survivors would have proved too much.
\end{remark}

% ---------------------------------------------------------------------
\subsection{What the section costs and what it buys}
\label{gr:sub:ledger}

The apparatus is priced here in the companion's format
\cite{Ben26Optimal}: costs enumerated with grade and falsifier, buys
enumerated after, and the whole kept short enough to audit.

\paragraph{Costs.} Six load-bearing commitments, no more, and their
grades are not uniform.

\begin{enumerate}[itemsep=2pt]
\item \emph{The biospheric referent} ($\Phi1$) ---
\textbf{methodological stipulation}. \emph{Falsifier:} exhibit a unit of
analysis referring to the same object at every point of the trajectory
while carrying less semantic content than ``biosphere''. The stipulation
is that none exists; one counterexample retires it, and the theorems
would then be applied in that unit instead.

\item \emph{Vantage, not obligation} ($\Phi2$) --- \textbf{scope
declaration}, and a self-imposed restriction rather than an assumption
doing work. \emph{Falsifier:} exhibit a normative conclusion derivable
from the apparatus without adding a value premise. Its failure would be
an expansion of the section's reach, not a defeat --- but it would
defeat the claim that Golden reasoning is not longtermism, which is why it is stated here rather than assumed.

\item \emph{The Extraction Hypothesis} ($\Phi3$) ---
\textbf{conjecture-grade}, argued from boundary cost but not derived,
and the section's principal exposure --- though no longer a solitary
one: \S\ref{gr:subsub:xrisk} records two independent precedents
\cite{Tainter88,Bostrom19} reaching compatible conclusions from
unrelated evidence, which is convergence rather than proof but improves
the premise's standing. In its present form it is the
claim $\rho<0$ for extractive organization
(Definition~\ref{gr1:def:anatomy}), which is why
Corollary~\ref{gr:cor:extractcost} follows from it mechanically and why
\S\ref{gr:subsub:neglect} needs no further premise. Its position in the argument is
exact and worth stating in the ledger: it is not a premise of any
theorem, but the \emph{identification} that connects the theorems to
their application, so its failure does not falsify anything in
\S\ref{gr1:sec} --- it \emph{empties}
\S\ref{gr:sub:consequences}, leaving the selection results intact and
silent about predation. \emph{Falsifiers:} the three of
Remark~\ref{gr:rem:falsify}, all operational.

\item \emph{Origin does not transfer} ($\Phi4$) --- \textbf{premise}.
\emph{Falsifier:} a mechanism by which the selective regime that
produced intelligence in a lineage transfers to the regime governing a
biosphere's persistence at world scale. Note this is the weaker of the
two anti-violence commitments and can stand alone: $\Phi4$ alone
establishes that predatory maturity is not \emph{inherited}, and only
$\Phi3$ adds that it is selected \emph{against}.

\item \emph{Rarity of origins} ($\Phi5$) --- \textbf{imported result},
carried at the companion's grade and not re-argued
\cite{Nexus, Ben26Found, Ben26Optimal}. \emph{Falsifier:} the
companion's, and nothing this section adds.

\item \emph{The inherited mathematical apparatus} --- \textbf{[A]} and
below. Every quantitative claim here --- the $t^{-1/4}$ tail, the Lindy
residual life, the $s/4\varepsilon$ criterion, the measured truncation
--- inherits the tags of \S\ref{gr1:sec}, in
particular the soft-wall Ansatz. \emph{Falsifier:} the numerical tests
specified there, chiefly a fitted exponent robustly away from
$-\tfrac14$.
\end{enumerate}

\noindent One further entry belongs in this column, though it is not a
premise: \emph{the section's novelty is narrower than its apparatus}.
The mathematics of \S\ref{gr1:sec} is classical or routine --- the
$\tfrac14$ exponent is McKean's, the background construction is an
application of known results on sums of random exponentials --- the Lindy
conclusion of \S\ref{gr:sub:consequences} is established prior work
\cite{BC21,BG24,Kip20,Ord23}; and the mechanism behind $\Phi3$ has
precedents in the collapse and existential-risk literatures
\cite{Tainter88,Bostrom19,EF24}; and the conclusion of
\S\ref{gr:subsub:violence} is the standing technological-adolescence
argument of SETI, already current and already rebutted
\cite{Cooper19}. What is claimed there is a mechanism that survives the
rebuttal, not the conclusion. What is claimed as new is the assembly
and three specific items: the measure-theoretic recasting of the
survivor category, the practical/ideal ratio criterion, and the
epistemic ceiling. A reader who finds any of those three anticipated
should discount the section accordingly, and the surrounding apparatus
provides no independent support for them.

One accounting note closes the column, and it governs what all
six costs are costs \emph{of}. The unconditional list is short, and the
ledger states it in full: that conditioning selects, that adverse
anatomy is fatal (Theorem~\ref{gr1:thm:trapped}), and that viability
requires two conditions rather than one
(Theorem~\ref{gr1:thm:twocond}). The mutual singularity of $Q$ and
$\PP$ is \emph{not} on that list; it is critical-regime only. What is priced is every sentence connecting that structure to
biospheres, aliens, or the future. No rhetorical weight anywhere in this
section is permitted to exceed the grade of $\Phi3$, and where a
consequence is stated in \S\ref{gr:sub:consequences} it is stated as
downstream of it.

\paragraph{Buys.} Against those six, the section returns the following.

The categorical division is \emph{unchosen}, in the critical regime:
General and Golden are mutually singular measures there, so the split is
established by a theorem rather than by a threshold someone placed,
which is what Pre-Civilization/Civilization and every Kardashev boundary
cannot say. \S\ref{gr:subsub:regimes} records the price --- the claim is
unavailable in the viable regime --- and the entry is worth exactly what
it is worth after that subtraction.
Practical Goldenness is \emph{computable}: $t \gtrsim s/4\varepsilon$,
with the invariant $s_\varepsilon(t)/t = 4\varepsilon$, so the question
``how Golden is this'' has an answer with quantities in it. The premise
is \emph{weaker than it looks}: age alone confers practical Goldenness,
and age is what the question already supposes, so the apparatus assumes
there is a future rather than that the future goes well. The age
Lindy \emph{exponent} is derived rather than parameterized --- a
narrower buy than it first appears, and stated as such in
\S\ref{gr:sub:consequences}, since Lindy longevity for civilizations is
established prior work \cite{BC21,BG24,Kip20,Ord23} and what is added is
that the exponent is fixed by the driver's smoothness rather than
chosen. Against the indifference-based anthropic tradition
\cite{Car83,Got93,Bos02} the contrast is larger, the age distribution
there being postulated outright; but the nearer neighbor is already
model-based, and the net gain is one parameter rather than a
paradigm.
The epistemic ceiling is \emph{symmetric}, constraining the section
exactly as tightly as its targets. And the section \emph{fills a
declared gap} in the companion: conditioning on $\{T>t\}$ is
conditioning forward, which is what \cite{Ben26Optimal} states its own
constraint cannot do.

One entry is a positive expectation rather than a dissolution, and it is
entered deliberately, since an apparatus whose entire return is relief
from other people's difficulties is answerable to nothing. The
truncation/rescaling asymmetry of \S\ref{gr:subsub:measured} is not a
fact about biospheres; it is a fact about survivorship conditioning on a
two-coordinate state. It should therefore appear in \emph{any} domain
carrying survivorship selection and an orientation-like coordinate
alongside a magnitude-like one --- firms, institutions, lineages, any
population where the old are selected from the many. The prediction is
specific: the orientation coordinate should show a deleted region among
survivors, the magnitude coordinate a shifted distribution of roughly
preserved shape, and the two divergences should be comparable in
magnitude while opposite in kind. That is checkable outside cosmology,
on data that exist, and it would embarrass the apparatus if the
selection turned out to act on magnitude categorically and orientation
quantitatively instead. It is entered here as an expectation and not a
result.

% ---------------------------------------------------------------------
\subsection{What this section does not claim}\label{gr:sub:fence}

The fence is stated hard, on the companion's discipline that an
apparatus which prices what it claims must also refuse what it does not
\cite{Ben26Optimal}.

It is not claimed that any observed system is Golden, ideally or
practically; by Theorem~\ref{gr:thm:ceiling} the ideal version is
unknowable to any finite observer (Remark~\ref{gr:rem:ceilingscope}) and
the practical version is out of reach for want of an ensemble. It is not claimed that the companion's cosmic
weighting and the path measure $\PP$ of
\S\ref{gr1:sec} are the same object; that
identification is flagged as unperformed in
\S\ref{gr:subsub:nexic} and nothing here depends on it. It is not
claimed that any individual biosphere, ours included, will behave as
$\S\ref{gr:sub:consequences}$ describes; those are statements about a
conditional distribution. No normative conclusion is claimed
($\Phi2$). And the companion's natural-historical program --- its
conjectures, its bounds, its paleontological predictions --- is not
claimed, imported, or re-argued; $\Phi5$ enters as a stated result and
nowhere else.

One thing remains on the far side of the fence, because a fence is
legible only if what it excludes can be seen. A program built on this
section would inherit three commitments, stated here as refusable rather
than as claims. It would be committed to the boundary-cost scaling of
$\Phi3$ showing up in institutional and ecological data, not merely in
argument. It would be committed to the truncation/rescaling asymmetry
appearing in survivorship-selected populations generally. And it would
be committed to the exponent's dependence on noise placement --- $\tfrac14$
for integrated drivers, $\tfrac12$ for Brownian ones --- being visible
in any domain where hazard dynamics can be measured directly. None of
the three is established here, all three could be found false without
touching a line of \S\ref{gr1:sec}, and that is the
point of naming them.

% ---------------------------------------------------------------------
\subsection{Section summary}\label{gr:sub:summary}

The unit is the biosphere ($\Phi1$), because no other term survives the
limit the theorems are taken in. The fundamental division is General
versus Golden --- $\PP$ versus $Q$ --- and it is categorical without
being chosen, since in the critical regime the two are mutually
singular; but nothing occupies
the Golden category \emph{as an ideal}, which is a measure to reason
from rather than a state to be in, and this is what distinguishes Golden
reasoning from longtermism ($\Phi2$). The practical form is occupiable
and is what the section actually assumes: a biosphere of age $t$ is
$\varepsilon$-Golden over windows of length $4\varepsilon t$, a constant
fraction of its own history, so that the premise of Golden reasoning is
not that the future goes well but merely that there is one. The two coordinates are technostrength, a
state coordinate, and technoanatomy, a structural parameter, and
conditioning treats them differently in kind: adverse strength is
truncated abruptly, $137$-fold, while anatomy is selected gradually,
$\PP(\rho>0)$ moving only from $0.75$ to $0.84$ over the same horizon.
The second is the slower and the more decisive, since adverse anatomy is
eventually fatal with certainty while adverse strength is a condition
every biosphere occupies half of logarithmic time and survives. The Kardashev scale fails not because its
thresholds are arbitrary but because the terminal stage it points at is
gated by a technoanatomy the scale cannot index, so that climbing is
neither necessary nor sufficient for arriving --- and, with adverse
anatomy, climbing fast is itself the failure. The violent-aliens thesis fails at the
inference from origin to maturity ($\Phi4$), and fails further if
extraction is bad technoanatomy ($\Phi3$) --- which we argue from
boundary cost, recovering the square--cube law from which the names
came, and which we mark as the section's one load-bearing premise, with
its falsification conditions. Downstream: encountered biospheres are not
extractive, megastructure nulls are weak evidence, and old biospheres
are Lindy with a derived exponent rather than an assumed one --- which
is where Golden reasoning improves on Copernican anthropics, by making
the age distribution an output rather than an input. With rarity of
origins ($\Phi5$), the silence needs no filter ahead. And by the
epistemic ceiling, none of this can be confirmed by encounter --- a
limitation that binds this section exactly as tightly as it binds what
the section argues against.

% ---------------------------------------------------------------------

% ============================================================
% §7 — GOLDEN SCORES: A DECISION RULE AND ITS LIMITS
% REVISION 2, against the §7 referee report (B1–B4, C).
%
% MERGE NOTES — disposition:
%   B1 ACCEPTED (shared defect, both drafts): ≍ does not deliver
%     the limit — bounded ratio only, per the master's own
%     convention. The derivation now runs on the REFINED form ∼,
%     which sh:thm:persist + sh:prop:harmonic + GJW99 supply for
%     the hard-wall functional τ and sh:thm:Q assumes; the
%     soft-wall Ansatz transfers T to τ at O(1) in the log,
%     ≍-grade, no upgrade. The recommended statement adopted
%     nearly verbatim inside gs:prop:score. The qualification is
%     in the proposition itself, not a footnote.
%   B2 ACCEPTED: the result is now Proposition gs:prop:score,
%     tagged [P] given [A] and [C] inputs, in the bracket.
%   B3 ACCEPTED, both halves: ensemble form stated first (free —
%     the Tendency at the level of practice, no [O] needed);
%     the trajectory form ROUTED to where §2 locates trajectory
%     claims (gr1:subsub:fluct, gr1:cor:recurrentanatomy, [O]
%     dynamics), which was the report's preferred repair and
%     adds the missing connection.
%   B4 ACCEPTED: trapped row finished — the rule collapses to a
%     comparison of floors h_ast via gr1:eq:xstar, and the
%     verdict is hr:rem:concession's "move ρ."
%   C: seven \S2 → \S\ref{ac:sec} (and PROPAGATED into §2's own
%     bodies, which had the same hardcoding); Claim~\ref for
%     sh:claim:tail; booktabs table; currency fence added after
%     Aurutance; §7 quote block still owed (author to select).
%   A3 RECORD: the Aurutance discrepancy flag identified an
%     error in the BRIEF, not the source — brief 7.3's gloss
%     dropped the prior terms; both drafts independently carried
%     the 2024 full form. Converged; instruction error recorded.
%   NO PREMISE ENVIRONMENT; guard arithmetic unchanged
%     (Φ6 in §2, Φ1–Φ5 in §6, Φ7 at adversivity).
% ============================================================

\section{Retrocausal accordance: the principle, the argument, and the
target}\label{ra:sec}

\begin{quote}\small
The Accordance Principle constrains one's past: a present that would be
a fluke relative to one's Nexus is evidence against that Nexus. Run
forward, it constrains one's future, and this section develops that
forward form as one cluster: the definition of \emph{retrocausal} that
keeps it clear of teleology, the principle in both its forms, the
worked example at human scale on which a reader can check it, and the
argument it powers --- that the future attribute accounting for our
occupying the present rather than deep time must be correlated with the
whole of modern development to a degree with no ordinary analogue, and
that indefinite biosphere survivorship is the candidate that fits.

The cluster is graded with more care than the rest of this work, not
less, because it sits nearer to teleology than anything else here. The
principle is \textbf{[$\Phi$]} at the companion's grade; the argument's
quantitative core is \textbf{[$\Phi$]}; the identification of its
conclusion with the Golden Point is \textbf{[S]}, an inference to the
best explanation, and says so; and the epistemic ceiling applies to all
of it.
\end{quote}

\subsection{The word, before the principle}\label{ra:sub:definition}

\begin{definition}[Retrocausal conditioning]\label{gc:def:retrocausal}
\emph{Retrocausal}, in this work, means \emph{conditioned on an Aethus
lying in one's future light cone}, and nothing further. Golden
Optimation is retrocausal in exactly this sense and in no other.
\end{definition}

\begin{remark}[What the stipulation does and does not buy]
\label{gc:rem:retrocaution}
The definition is deliberately dry, and the dryness is what keeps it
inside Remark~\ref{gc:rem:teleology}'s discipline. Nothing propagates
backward; no event acts on its own past; the verbs remain \emph{is
conditioned}, \emph{is reweighted}, \emph{is selected}, and the
forbidden ones --- \emph{aims}, \emph{seeks}, \emph{progresses toward}
--- are as forbidden here as anywhere else in this work.

What should nonetheless be said plainly is that the stipulation lands
closer to the classical notion than a purely evidential reading usually
does, and for a structural reason rather than a rhetorical one. In a
setting where causation is fundamental, conditioning on a future event
is uncontroversially inference and uncontroversially not causation,
because there is a further fact about what caused what. This framework
has no such further fact: causation is emergent from the Aethic
structure rather than primitive, and determinacy is indexed to an
information state rather than to a date, so that an unresolved future
outcome is a blank entry and blanks lack tense
(Remark~\ref{gc:rem:notense}, \cite{AethusDecision}). With no primitive
causal relation underneath, the distinction between \emph{influence} and
\emph{conditioning} has nothing to consist in beyond the conditioning
itself. The reading offered here is therefore the deflationary one, and
the proximity to the inflationary one is recorded as a feature of the
framework rather than as a claim about the world.
\end{remark}



\subsection{The principle, in two forms}\label{gc:sub:retroaccord}

The companion's
Accordance Principle constrains one's past: a present that would be a
fluke relative to one's Nexus is evidence against that Nexus. Run
forward, it constrains one's future.

\begin{principle}[Retrocausal Accordance; \textbf{[$\Phi$]}]
\label{gc:prin:retroaccord}
Let $A$ be a realized attribute of one's Aethus with canonical
alternative $A'$ its complement, and let $B$ be an attribute of one's
future cone. Then
\begin{equation}\label{gc:eq:retroaccord}
\frac{\PP(B\mid A')}{\PP(B\mid A)}\;\lesssim\;\frac{\PP(A)}{\PP(A')},
\end{equation}
weights taken under the Cosmic Nexus, with $\lesssim$ carrying the
companion's exceedance modality and no free tolerance: for any
$\kappa\ge1$,
\begin{equation}\label{gc:eq:retroexceed}
\Pr\left[\frac{\PP(B\mid A')}{\PP(B\mid A)}
>\kappa\,\frac{\PP(A)}{\PP(A')}\right]\;\le\;\frac{1}{1+\kappa},
\end{equation}
the probability taken over the realization \cite{NexicFoundation}. The
motto: \emph{one's present cannot be arbitrary with respect to one's
future.}
\end{principle}

The principle is the Accordance Principle with one substitution, and the
substitution is the whole of what is new. There, the conditioning
attribute is the production of an observer and lies behind one; here it
is an attribute of the future cone. Everything else transfers unchanged
--- the canonical alternative is the complement, so nothing is
analyst-supplied; the draw is the same Copernican one; and by Bayes
\eqref{gc:eq:retroexceed} is exactly the statement that the realized
member carries $B$-conditioned weight below $1/(1+\kappa)$, which for
$\kappa\ge1$ at most one member of a pair can do. What licenses the
substitution is that the framework indexes determinacy to an information
state rather than to a date, so a future-cone attribute is available to
condition on in the same sense a past-cone one is,
(Remark~\ref{gc:rem:notense}).

Two readings to keep apart. \emph{Not a promise:} the principle does not
say one's flukes will pay off. It says futures in which they do not are
disfavored in the stated ratio, which is an inference from the fluke and
not a claim about what the future will do --- Remark~\ref{gc:rem:teleology}
applied here. \emph{And not a tolerance:} the $\lesssim$ prices
violations rather than forbidding them, and
\eqref{gc:eq:retroexceed} fixes the entire tail profile, so a
$\kappa$-fold exceedance is not a counterexample but an event of
probability at most $1/(1+\kappa)$.

\begin{remark}[The exceedance form is the multiplicity guard]
\label{gc:rem:noisefloor}
This bears directly on the look-elsewhere worry raised above, and
supplies more of an answer than was there credited. Because
\eqref{gc:eq:retroexceed} bounds each interrogated pair's violation
probability, linearity bounds the expected number of $\kappa$-fold
violations among $n$ interrogated pairs by $n/(1+\kappa)$ --- an exact
noise floor, supplied by the principle's own tail profile rather than
imposed from outside \cite{NexicFoundation}. Scanning a hundred
attributes, one should expect up to ten spurious nine-fold flags and one
spurious ninety-nine-fold flag; a trace standing clear of that floor is
forced, and one that does not is not.

So fluke-hunting in a Route World is not the undisciplined activity
\S\ref{gc:subsub:routeworld} first suggested. It has the same
multiple-comparison discipline every inference of the kind owes, with
the correction computed rather than chosen. The concession the companion
records applies here with full force and is the operative caution:
\emph{$n$ must count the attributes implicitly scanned and not merely
those explicitly interrogated}, since a striking attribute is
interrogated because a scan flagged it. In a Route World the scan is
one's whole attention, and a correct $n$ is correspondingly enormous ---
which is a reason the identifiable traces must be conspicuous by a wide
margin before they are worth anything, and not a reason there are none.
\end{remark}



\begin{principle}[Retrocausal Accordance, bound form; \textbf{[$\Phi$]}]
\label{ra:prin:bound}
For an Aetherian whose Aethus is $H_1$, and a stated attribute
$A\supset H_1$ that is unlikely under the Cosmic Nexus, the probability
of centric unfolding to a future Aethus containing stated attribute $B$
is bounded above by the probability of $A$ given $B$ under the Cosmic
Nexus:
\begin{equation}\label{ra:eq:bound}
\PP(B\mid H_1,N_0)\;\lesssim\;\PP(A\mid B,N_0)
\qquad(H_1\subset A,\ \PP(A\mid N_0)\ll1),
\end{equation}
with $\lesssim$ carrying the exceedance modality of
\eqref{gc:eq:retroexceed}.
\end{principle}

The two forms are one principle read from either side.
\eqref{gc:eq:retroaccord} is the ratio form, and connects to the
companion; \eqref{ra:eq:bound} is the bound form, and is the one on which
the motto reads literally --- futures that would leave one's present a
fluke are capped in weight by how far they would leave it one. The
derivation of the second is the companion's own argument run forward.
If $A$ is exceptionally unlikely with respect to a future Aethus
$H_1\cap B$, then one's currently standing in $H_1$ --- for which $A$
has occurred --- is already a fluke with respect to that future; and
since one's future Aethae will eventually be one's own, and so warrant
the first postulate on one's behalf, a present that is a fluke relative
to any future one might reach violates the mediocrity constraint. The
weight of unfolding into such a future is accordingly constrained in
proportion to its dubiousness \cite{Nexus}.

\subsection{A worked case, at the scale a reader can check}
\label{ra:sub:example}

The examples elsewhere in this work are cosmic. A principle that works
only at cosmic scale is suspect, and the following runs it on a single
attribute of a single life, which is the grain at which a reader can hold
it against their own experience.

In grade school I was told, on a few occasions, that I had an
exceptionally strong presentational speaking ability. Holding that in
mind while deriving Nexic reasoning in 2021\footnote{A provenance note, so that the example is not misdated. The presentational-ability trait is what I used to \emph{formalize} the retrocausal accordance principle in 2025, when a long-standing intuition resurfaced unprompted after a presentation I gave at a summer program. The 2021 version was not presentation-based, or not solely: it ran on a blend of particularly immense flukes I had witnessed through childhood, each of them plausibly conditioned on my later being able to produce these frameworks at all, and to varying degrees of confidence. The reasoning was the same; the attribute I hung it on was not. Two further notes belong here rather than in the body. First, the principle has failed in my hands: on a day in April 2026 I ran a set of guesses through it that came out wrong, and for an instructive reason --- every inference had been conditioned on a particular result being important that turned out to be trivial. Human error in the Route World's leap of faith is as available as it ever was, and the failure is recorded because a worked failure is worth more to a reader than any success \cite{koide1983new, Brannen2006TheLM}. Second, that failure raised a question the paper's ontology makes askable and the classical one cannot: is the counterfactual in which the result was never posted one in which it was \emph{trivial all along}, or one in which the triviality is \emph{joint} with the posting in this Aethic reality? Here ``trivial'' is a trait of the initial conditions of human mathematical knowledge rather than of the mathematics itself. Aethically both are live, and the question is not idle, since it bears on whether Golden causality's demotion of an event is a fact about the event or about the branch. The journal record of the episode, with the surrounding personal circumstances, is held in the corpus's archive material and is not reproduced here.}, I had a thought that struck
me as odd and then as forced. \emph{If I go my entire life without ever
needing that skill, does my having it now not represent an immense fluke
with respect to that life?} Take the set of lives I might live in which
strong speaking ability is never called on. Living the one I am in ---
where I have it --- would be a decided fluke in that set. So, as an
effective proof by contradiction to keep this waking moment from being a
fluke, I would have to infer that somewhere in my future at least one
task requiring strong speaking ability was coming. I had constrained an
attribute of a future Aethus of mine from nothing but my own standing
fluke content. And I appear to have been right, since I now know I
needed exactly that ability to write and present the papers you are
reading.

Two things about the example, and the second is the important one.

\emph{What it displays.} It is the retrocausal accordance principle
applied to one attribute: the fluke is \emph{having the ability}, the
future attribute is \emph{a task requiring it}, and
\eqref{ra:eq:bound} says the weight of futures in which no such task
comes is capped by how much of a fluke the ability would be under them.
Every Aetherian may constrain their own future this way from their own
present, and the principle's content is exactly that this is legitimate
--- not that the constraint is tight.

\emph{What it does not test.} The confirmation is retrospective, and a
single retrospective confirmation is one draw. The principle predicts
that standing flukes will \emph{tend} to pay off, at the exceedance
grade --- $\kappa$-fold failures at probability at most $1/(1+\kappa)$
--- and one instance neither confirms nor refutes a tendency. The
example is included because it is where the principle is checkable in
kind, not because it is evidence of the principle's truth; a reader who
runs it on their own standing flukes and finds it wanting has learned
something the paper's cosmic cases could not have taught them.

\begin{remark}[Two directions, one of them deductive]\label{ra:rem:twodirs}
The retrocausal accordance principle and Golden Optimation
(\S\ref{gc:sub:goldenopt}) bear on the same fluke traits from opposite
sides, and conflating them is the error most worth guarding against here.

Golden Optimation runs \emph{from} the conditioning target \emph{to} the
trait: if the Golden Point required it, one has it. That is deduction
under the conditioning, and it explains \emph{why} one has the trait.
The retrocausal accordance principle runs \emph{from} the trait \emph{to}
the future: having it, one is constrained toward futures in which it
pays. That is inference, at exceedance grade, and it explains nothing
about why one has the trait --- only what one's having it bears on.

The converse of the deductive direction does not hold, and this work is
careful never to run it. A fluke one possesses need not have been
required by anything: to argue from ``I have this trait'' to ``the Golden
Point required it'' would be to run the deduction backward, and nothing
here does. What \emph{does} hold is the forward form only --- required,
therefore had --- which is why in a Route World the question one would
face before a notoriously difficult feat is not \emph{will I achieve
this?} but \emph{is this genuinely required, and am I the one to do it?}
That is a different question, and it remains a leap because the
requirement may simply not hold, in which case even in a Route World the
feat is not probabilistically privileged as an Aethic event.
Remark~\ref{gc:rem:routefence} says the rest: no one can determine
whether their own feat is among the requirements, and the reformulated
question is therefore not one anyone can answer in their own favour.
\end{remark}

\subsection{The finale paradox, and the argument that solves it}
\label{ra:sub:finale}

\paragraph{The paradox.} Why does one find oneself in the modern day
rather than at any earlier point in Earth's natural history? A uniform
draw over the record lands in the last century with probability of order
$10^{-8}$; and the compounding of the companion's doomsday channels
worsens this, since at net hazard $\lambda$ the survival weight is
halved by $\ln2/\lambda$ and drops toward zero thereafter, so that a
position late in a surviving history is rarer still. Being \emph{here}
is a fluke of enormous magnitude on the received reading, and the
received reading has no resource to discharge it. Call this the
\emph{finale paradox}. The retrocausal accordance principle discharges
it, and the discharge derives the shape of what one must be conditioned
on.

% [Referee pass, Aug 21: the relay reading was carried in the author's
% design intent and not in the text.  Instated here because it changes
% what the magnitude below is load-bearing for.]
\begin{remark}[The paradox's two faces, and the relay it marks]
\label{ra:rem:relay}
The paradox stated above is the Golden-side face of a question the
companion program raises on its own account, and the pair should be
seen together, because the relation between them is this work's point
of contact with the companion rather than an echo of it.

\emph{The companion's face.} Nexic reasoning reads a record conditioned
on observerhood and finds Optimation in it. That reading raises a
question the companion names and explicitly declines to answer from the
Accordance Principle alone, on \emph{hypotheses non fingo} grounds:
\emph{when does Optimation cease to apply?} \cite{Aethus}. The question
is not idle housekeeping. Observer-conditioning is conditioning on an
event in one's past cone, so the tilt it supplies is exhausted once the
event has occurred, and a conditioning that is exhausted must have a
terminus --- yet nothing in the companion's apparatus locates one, and
anything one proposed would be stipulated. The companion therefore
carries the question open, and marks it as the frontier.

\emph{The dissolution.} Golden conditioning does not answer the
question; it removes the presupposition that generates it. If the
operative conditioning event lies in the \emph{future} cone, then
Optimation is not a process running over an interval but the past-cone
restriction of a conditioning that is not exhausted at all, and there
is no terminus to find because there was never a duration to
terminate. Remark~\ref{gc:rem:carryover} is the mechanism stated
exactly: a required fluke is Golden Optimation before the observer's
present reaches it and Optimation in the companion's sense afterwards,
same event and same conditioning, with only the observer's position
changed. So ``when does Optimation stop'' asks for the date at which a
shadow ceases to be cast by an object still standing. What replaces the
question is an account of what Optimation \emph{is} from first
principles --- the past-cone shadow of a future-cone conditioning ---
which is not a smaller answer than the one sought but a different and
prior one.

\emph{What this makes the paradox.} It is the designed relay point
between the two developments: the place where the companion's open
frontier is handed across and dissolved rather than solved, and the
reason the two programs are one apparatus read in two temporal
directions (Remark~\ref{gc:rem:rungzero}). A reader tracking the corpus
should take this subsection as the join, and not as a further
application of the companion's machinery within the companion's scope.

\emph{And it fixes what the magnitude below is for.} Two things are
being done in this subsection and they carry different exposures. The
\emph{dissolution} is structural: it consumes the direction of the
conditioning and nothing quantitative, so it survives any revision of
the draw discussed at Ansatz~\ref{ra:ans:draw}. The \emph{filter} of
Proposition~\ref{gc:prop:alpha} is quantitative and does not: its force
is the smallness of a ratio, and a different weighting over positions
gives a different bound. A reader who rejects the draw therefore keeps
the relay and loses the filter, which is the same one-way dependency
this work observes elsewhere, applied here inside a single subsection.

\emph{One consequence runs back to the companion.} The companion's
ancestral-agent discussion rests part of its methodological caution on
the finale question's being open --- the frontier's openness is what
makes premature verdicts about which eras Optimation covers
illegitimate there \cite{Aethus}. Under the dissolution the caution is
not withdrawn but re-grounded, and strengthened: there is no era in
which the conditioning fails to reach, so a verdict resting on
Optimation's having powered off before some ancestral stage is
ill-posed rather than merely unsupported, its presupposition
having failed. Whether the
companion should be revised to state it that way is for the companion
\textbf{[O]}.
\end{remark}

% [Referee pass, Aug 21: the draw producing the 10^{-8} was used and
% never named.  Named here, with the bracket in both directions.]
\begin{ansatz}[The positional draw over the record; \textbf{[$\Phi$]}]
\label{ra:ans:draw}
The weighting against which the present is assessed as improbable is
\emph{duration-uniform} over the record, hazard-weighted by the
companion's doomsday channels: it is the elapsed-time measure that
delivers $\PP(A\mid N_0)/\PP(A'\mid N_0)\sim10^{-8}$, driven lower by
the compounding.
\end{ansatz}

\begin{remark}[What the draw costs, and the bracket in both directions]
\label{ra:rem:drawcost}
Ansatz~\ref{ra:ans:draw} is posited, and the reason for naming it is
that the Cosmic Nexus does not obviously deliver it. The companion's
weighting runs over Aethae --- observer-situations --- and not over
elapsed time \cite{Nexus, NexicFoundation}, and the two need not agree
even in order of magnitude. The bracket runs both ways, which is why
the ansatz is a real posit rather than a formality.

\emph{Downward.} Restrict the answer space to Aethae capable of posing
the question --- an interrogation-indexed reading of the observer tilt
--- and such situations are concentrated exactly where the questioner
is. The ratio deflates toward order one, the bound
\eqref{gc:eq:alphabound} goes slack, and the filter selects nothing.

\emph{Upward.} Weight instead by observer-situations across the whole
record, with the pre-human ones counted as the companion's own
treatment of ancestral agency declines to demote them, and the ratio
is very much \emph{worse} than $10^{-8}$: the filter tightens, and the
argument's conclusion strengthens rather than fails.

So the ansatz is not a convenience chosen to make the argument work ---
one natural alternative sharpens it and the other dissolves it --- and
the honest statement is that the filter's force is exactly as secure as
the weighting, which is presently not derived. \emph{Falsifier and
repair.} The repair is a derivation of the positional weighting from
the Copernican draw: which Aethae populate this attribute's answer
space, and at what occurrence weights. Until it exists the filter is
\textbf{[$\Phi$]} on the ansatz, and a reader who declines the ansatz
holds Remark~\ref{ra:rem:relay}'s dissolution intact \textbf{[O]}.
\end{remark}

\paragraph{The argument, in four steps.}

\emph{Step 1: the modern-day attribute set.} Let $A$ be the conjunction
of stated attributes about our biosphere that apply exclusively to
one's living in the present as opposed to any earlier stretch of natural
history --- that the hottest and coldest temperatures in the observable
universe both exist on Earth
\cite{ALICE2012Temperature,Kovachy2015Picokelvin}; that the first
vertebrate to pass the K\'arm\'an line was born within the last century
\cite{burgess2007animals}; that unlocking dormant genes by engineering is
suddenly possible \cite{Pattali2026_CRISPR_epigenome_mod}. Any parent
$A_i$ of $A$ is a subset of these developments.

\emph{Step 2: retrocausal accordance forces a future requisite.} By
\eqref{ra:eq:bound}, for each $A_i$ there is more likely than not some
future stated attribute $B_i$ for which $A_i$ is a near-necessary
prerequisite as against one's having been born far back in deep time.
This is an optimisation query. For an $A_i$ containing only a fraction
of the developments, having \emph{no} future $B_i$ that takes the rest
as necessary is unlikely by the principle; for the whole of $A$, a
handful of developments necessitated by no future is permitted by the
exceedance margin. So there is some $A_j$ for which the significant
majority of modern developments root to a single future stated attribute
$B_j$.

\emph{Step 3: the shape of $B_j$.} $B_j$ must render those developments
categorically necessary as prerequisites, against their not having
occurred across deep time. It is therefore not merely correlated with
human capacity to manipulate the environment; it is essentially the
stated attribute with the strongest possible correlation to it --- a
maximally connected graph in which one future attribute takes every
modern development as a total necessity.

\emph{Step 4: the suppression bound.} This is the quantitative core, and
it is a proposition.

\begin{proposition}[The accounting future must be almost perfectly
correlated; \textbf{[$\Phi$]}]\label{gc:prop:alpha}
Let $A$ be the conjunction of attributes distinguishing the present from
any earlier stretch of the record, $A'$ its complement --- occupying some
earlier stretch --- and let $B$ be a future attribute for which
Principle~\ref{gc:prin:retroaccord} is satisfied. Partition $B$ into a
component $C$ almost perfectly correlated with $A$ and a component $C'$
approximately independent of it, with weight $\alpha$ on the second.
Then, since $A'\cap C$ is near-invalid and $A'$ is independent of $C'$,
\begin{equation}\label{gc:eq:alpha}
\PP(A\mid N_0\cap B)\;=\;
\frac{(1-\alpha)\PP(A\mid N_0)}
{(1-\alpha)\PP(A\mid N_0)+\alpha\,\PP(A'\mid N_0)\PP(C'\mid N_0)},
\end{equation}
and requiring the present to be no less weighted than the alternative
gives
\begin{equation}\label{gc:eq:alphabound}
\alpha\;\lesssim\;\frac{\PP(A\mid N_0)}{\PP(A'\mid N_0)} .
\end{equation}
\end{proposition}

The bound's force is in its magnitude. A uniform draw over the record
places the last century at some $10^{-8}$, and the exponential
compounding of the companion's doomsday channels drives the ratio
further down --- at net hazard $\lambda$ the weight is already halved by
$\ln2/\lambda$, and ten such half-lives put the ratio near $10^{-11}$.
So $B$ must track $A$ to within one part in $10^{8}$ or far tighter:
\emph{any} independent component of $B$ larger than that opens a channel
through which the vastly more numerous earlier positions dominate, and
$\PP(A\mid N_0\cap B)$ collapses with the full rapidity the ratio
implies.

\begin{remark}[What the bound is evidence for, and what it is not]
\label{gc:rem:alphareading}
Read as a constraint the proposition is strong and the reading is
narrow: whatever future attribute accounts for one's occupying the
present rather than an earlier stretch must be correlated with the
totality of modern development to a degree with no ordinary analogue.
That is a filter on candidates, and most candidates fail it --- a
peripheral future development correlates with some of the record's
recent attributes and is independent of most, which is exactly the
$\alpha$ the bound forbids.

The step from there to the Golden Point is weaker and should be marked
as such. It is that the Golden Point satisfies the filter more plainly
than the alternatives on offer, since indefinite persistence of the
biosphere is the one candidate that takes essentially every capacity the
record has recently acquired as prerequisite rather than as incident.
That is an argument from the absence of a better candidate, and it
inherits the exhaustiveness problem Remark~\ref{gc:rem:routerobust} records:
the filter is computed, the survey of what passes it is not.

Two further cautions. The partition into $C$ and $C'$ is a modelling
choice, and a $B$ whose correlation with $A$ is graded rather than
two-component would need \eqref{gc:eq:alpha} restated over the grading.
And the whole argument is conditional on
Principle~\ref{gc:prin:retroaccord} being satisfied at all, which
\eqref{gc:eq:retroexceed} makes a probabilistic rather than a categorical
premise --- so the conclusion is that \emph{if} the present is not a
fluke with respect to one's future, the accounting future is of this
shape, and \S\ref{gr:sub:ceiling} forbids anyone from establishing the
antecedent \textbf{[O]}.
\end{remark}

% [Referee pass, Aug 21: the paper polices roster inference at
% gc:rem:copernicancost and not here.  The asymmetry is principled and
% is stated rather than left for a reader to reconstruct.]
\begin{remark}[Why this is not the inference the companion disowns]
\label{ra:rem:noroster}
An objection arrives from inside the corpus and deserves answering in
the corpus's own terms, since this work polices exactly this move
elsewhere and would be inconsistent if it did not police it here.

The companion's extraction is explicit that the global birth-rank
inference of the doomsday argument \cite{Leslie96} is \emph{not} recovered: a position
in a completed roster has no attribute-wise translation, and the retort
that the position could itself be made an attribute fails, because that
attribute's answer space presupposes the completed roster it was meant
to replace \cite{NexicFoundation}. The finale argument reasons from
one's position in the record. Why is it not the same move, and why does
Remark~\ref{gc:rem:copernicancost} disqualify
Ansatz~\ref{gc:ans:copernican} on precisely these grounds while the
present step stands?

\emph{The asymmetry is the normalizer, and it is exact.} The doomsday
inference requires the total $N$ --- the completed roster including
what has not occurred --- and $N$ is a blank. A ratio whose denominator
is blank has no answer space to draw over, which is the extraction's
objection stated mechanically. The finale ratio requires the
\emph{elapsed} record and nothing further: both $A$ and $A'$ are
stated attribute-sets of one's own Aethus, the answer space is the
partition of a stated record, and the draw is over positions already
in the books. So the finale step is an attribute-wise constraint of
exactly the licensed kind, and doomsday is not, and the discriminating
feature is not that one looks backward and the other forward but that
one normalizes over stated attributes and the other over a blank.

That also explains, rather than excuses, the treatment of
Ansatz~\ref{gc:ans:copernican}. Its normalizer is $\Lc_\ast$, the total
evidence the conditioning \emph{will} supply --- a blank, and therefore
a roster --- which is why that ansatz is disqualified from Accordance's
licence and carried as an independent posit. The two verdicts are one
criterion applied twice: normalize over what is stated and the
constraint is licensed; normalize over what is blank and it is not.
Recording the criterion in this form is what makes the corpus's
treatment of positional inference a rule rather than a pair of
judgements.

\emph{What it does not buy.} The criterion licenses the \emph{form} of
the finale ratio and says nothing about its \emph{value}: which stated
attributes populate the answer space, and at what weights, is
Ansatz~\ref{ra:ans:draw}'s question and remains open. Licensed and
underdetermined are compatible, and the argument is currently both.
\end{remark}

\paragraph{The conclusion, and its grade.} What the argument delivers is
a filter and a candidate. The filter is
Proposition~\ref{gc:prop:alpha}: whatever future attribute accounts for
one's occupying the present must be correlated with the totality of
modern development to within one part in $10^{8}$ or far tighter, and
must on sight be unmistakably optimised at a cosmic level of
significance --- since any independent component larger than that opens
a channel through which the vastly more numerous earlier positions
dominate. The candidate is the Golden Point: indefinite persistence of
the biosphere is the one future attribute on offer that takes
essentially every capacity the record has recently acquired as
prerequisite rather than as incident, and any alternative that has been
articulated does not present as ambitious enough to satisfy the
unmistakability condition. Both sentences carry the grade of their
source, \emph{in} the sentence where a reader will quote them: the
magnitude is \textbf{[$\Phi$]}-on-Ansatz~\ref{ra:ans:draw} --- under
the interrogation-indexed alternative the ratio deflates to order one
and this filter selects nothing --- so the filter's teeth are exactly
as sharp as the positional weighting is right. And the mirror question
a careful reader will pose --- under Golden conditioning, why is the
deep \emph{future} not the dominant position? --- has a one-sentence
answer that should sit here rather than be reconstructed: future
positions are blanks, not stated attributes, and
Remark~\ref{ra:rem:noroster} normalizes only over what is stated. The argument therefore both discharges the
finale paradox and, using its own momentum, argues for the Golden
Point's significance on the Aethic tree.

The two halves carry different grades and the difference is stated
rather than blurred. Steps 1--4 run on the imported principle and
standard probability, and are \textbf{[$\Phi$]} at the companion's
grade. The identification of $B_j$ with the Golden Point is
\textbf{[S]}: an inference to the best explanation, whose
``unmistakability'' criterion is doing real work and is not formal. What
would count against it is an articulated alternative $B$ that
necessitates the modern record as strongly and is not the Golden Point,
and no such alternative has been exhibited; that is the argument's
exposure, and it is the same exhaustiveness exposure
Remark~\ref{gc:rem:routerobust} records for the robustness argument.

\begin{remark}[What this argument adds to the reductio]\label{ra:rem:spine}
The conditioning argument of \S\ref{gc:sec} establishes that
observer-existence is too weak a target: it is indifferent to the
technological character of the record we hold, and a target sensitive to
that character recovers observer-conditioning as its past-cone shadow.
That clears the ground. It does not derive what the right target is,
except by naming the Golden Point and arguing that it fits.

The present argument does something the reductio cannot: it derives the
\emph{shape} of the target from the finale paradox and the retrocausal
accordance principle --- maximally necessitating of the modern record,
cosmically unmistakable --- and only then identifies the Golden Point as
the thing with that shape. That is a stronger route, and it is offered
here as the primary argument for the Golden Point as target, with the
reductio as the argument that clears the ground for it. A reader who
accepts the reductio and not this argument holds that observer-existence
is insufficient and the Golden Point is one candidate among possibly
several; a reader who accepts both holds that the Golden Point is the
candidate the record's shape selects.
\end{remark}

\begin{remark}[Golden Point Nexic Sorting, stated and deferred]
\label{ra:rem:sorting}
One further principle belongs beside the argument and is recorded here in
statement only, since its worked case is incomplete at source. Let
$A_1,\dots,A_n$ partition $A$ Aethically, let $\mathcal G$ be the Golden
Aethus and $N_0$ the Cosmic Nexus. Then the probability of an Aetherian
inhabiting $\mathcal G\cap A_i$ is proportional to the probability of
$\mathcal G$ and $A_i$ relative to $A$ under the Cosmic Nexus:
\begin{equation}\label{ra:eq:sorting}
\PP(\mathcal G\cap A_i\mid A,H_{\mathrm{in}})\;\propto\;
\PP(\mathcal G\cap A_i\mid A,N_0),
\end{equation}
so that an Aetherian's coming to inhabit a given cell is weighted
exactly by how the Golden Point and that cell co-occur from $A$ under
the unconditioned weighting --- the Golden analogue of the companion's
Optimation-weighting result. By total probability over the partition,
$\PP(\mathcal G\mid A,H_{\mathrm{in}})=\sum_i\PP(\mathcal G\mid A_i)
\PP(A_i\mid A,H_{\mathrm{in}})$, on the assumption that once a cell is
inhabited all further centric unfolding abides by that cell's own
weighting. The two-cell case around a Miracle --- $A_1$ before, $A_2$
after, with $\PP(\mathcal G\mid M)\approx\PP(\mathcal G\mid A_2)$ ---
reduces the cell ratio to
$\bigl[\PP(A_1\mid A,H_{\mathrm{in}})/\PP(A_2\mid A,H_{\mathrm{in}})\bigr]
\,\PP(M\mid A_1)$; the completion of that case is deferred
\textbf{[O]}.
\end{remark}

% [Merge: Opus's formal treatment adopted wholesale; alpha instance tied
% to Ansatz ra:ans:draw; determinism/flukes cross-link folded in.]
\subsection{Investment: what an expensive route tells you}
\label{ra:sub:investment}

The principle has a third face, and it is the one on which the Accordance
Principle is run with the Golden Point as its conditioning attribute
rather than with observer-production. Read that way it becomes
retrodictive: not a constraint on which futures one reaches, but a
reading of what a realized route's improbability \emph{purchased}.

\begin{proposition}[The investment bound]\label{ra:prop:investment}
Let $R$ be a realized route with Cosmic weight $q=\PP(R\mid N_0)$, and
$R'$ an alternative with $q'=\PP(R'\mid N_0)\gg q$.

\emph{(i) The identity, \textbf{[P]}.} Under conditioning on the Golden
Point $\mathcal G$,
\begin{equation}\label{ra:eq:invest}
\frac{Q(R)}{Q(R')}
=\frac{q}{q'}\cdot\frac{\PP(\mathcal G\mid R)}{\PP(\mathcal G\mid R')},
\qquad\text{so}\qquad
Q(R)\ge Q(R')
\iff
\frac{\PP(\mathcal G\mid R)}{\PP(\mathcal G\mid R')}\;\ge\;\frac{q'}{q}.
\end{equation}
This is exact and is an equivalence.

\emph{(ii) The inference, \textbf{[$\Phi$]}.} From \emph{observing} $R$
one may conclude that its necessity ratio is large --- but only at the
exceedance grade the Accordance Principle carries, since the step from
what is realized to what is weighted is the mediocrity step and not a
deduction. For any $\kappa\ge1$,
\begin{equation}\label{ra:eq:investexceed}
\Pr\left[\frac{\PP(\mathcal G\mid R)}{\PP(\mathcal G\mid R')}
<\frac1\kappa\cdot\frac{q'}{q}\right]\;\le\;\frac{1}{1+\kappa},
\end{equation}
the probability taken over the realization. \emph{Necessity purchased
matches weight spent, to within a $\kappa$-fold shortfall of probability
at most $1/(1+\kappa)$.}
\end{proposition}

\begin{proof}
\emph{(i)} $Q(\cdot)\propto\PP(\cdot)\PP(\mathcal G\mid\cdot)$ by
definition of the conditioning; divide. \emph{(ii)} By (i) the event
displayed is $Q(R)<\kappa^{-1}Q(R')$, which is the realized route
carrying conditioned weight a $\kappa$-fold factor below its
alternative's; \eqref{gc:eq:accordexceed} bounds exactly that at
$1/(1+\kappa)$.
\end{proof}

\begin{remark}[The exceedance is cheap, and the trade is worth stating]
\label{ra:rem:investkappa}
The two halves of the proposition should not be run together, and the
reason for separating them is that the second is where the argument's
strength actually lies.

Part~(i) is an equivalence and buys nothing on its own: it relates two
quantities neither of which is observed. Part~(ii) is what an observer
can use, and it is \emph{not} an equivalence --- one may be on a branch
the conditioning disfavours, and the principle prices that rather than
forbidding it. So the honest form of the investment reading is never
``the necessity ratio is at least $q'/q$'' but ``at least $q'/q$ divided
by whatever exceedance one is willing to tolerate.''

What makes this a qualification rather than a retreat is the shape of the
trade. Tolerating a $\kappa$-fold shortfall costs $\log_{10}\kappa$
decades of the bound and buys the failure probability down to
$1/(1+\kappa)$ --- and the first grows logarithmically while the second
collapses. At an observed cost ratio of seventeen decades: at
$\kappa=1$ the bound is seventeen decades and fails with probability at
most one-half; at $\kappa=10$, sixteen decades at one in eleven; at
$\kappa=100$, fifteen decades at one in a hundred; at $\kappa=10^{6}$,
still eleven decades, and failing with probability one in a million. Six
decades surrendered buys near-certainty.

So the investment bound is exceedance-graded and \emph{robust}, which is
a different thing from exceedance-graded and fragile. A reader who
insists on near-certainty pays a few decades out of many and keeps the
conclusion; a reader who wants the full $q'/q$ is asking for a
coin-flip's worth of confidence and should be told so.
\end{remark}

The reading this licenses is the one that gives the section its name.
Where a route of Cosmic weight $10^{-20}$ was realized and an alternative
of weight $10^{-3}$ was available, \eqref{ra:eq:investexceed} says the
realized route is better by seventeen decades at reaching the target ---
less $\log_{10}\kappa$ for whatever exceedance one tolerates, and holding
except on a branch of probability at most $1/(1+\kappa)$. Otherwise the
conditioning would weight the cheap route above it, and one would,
overwhelmingly, be on that one instead. So the observed
improbability is not merely a curiosity about how unlikely one's history
was. It is a \emph{lower bound on necessity} at a stated confidence, and the
larger it is the more the conditioning is committed to what it bought. That is what is
meant by saying Golden causality is heavily \emph{invested} in something:
the phrase names the quantity $\log(q'/q)$, and it is a measurement and
not an attitude.

\begin{remark}[The register, and what does not act]\label{ra:rem:investreg}
One sentence has to be said carefully because the natural phrasing is
forbidden. It is tempting to say that if a cheaper route had existed the
universe would have taken it. Nothing takes anything: by
Remark~\ref{gc:rem:teleology} a conditioning event does not act, and
\eqref{ra:eq:invest} is an identity about a measure. What the identity
says is that the conditioned measure \emph{weights} cheap routes above
expensive ones at equal effectiveness, so that finding oneself on an
expensive one is evidence --- at exactly the exceedance grade
\eqref{ra:eq:investexceed} makes explicit --- that the cheap ones do not
reach the target. The inference is evidential throughout, and the word
\emph{invested} is shorthand for a weight comparison rather than for an
expenditure anyone made.
\end{remark}

\begin{remark}[The two terms of Aurutance are cost and necessity]
\label{ra:rem:aurutance}
The investment bound and the decision rule of \S\ref{gs:sec} are the same
inequality read in opposite directions, and saying so vindicates a piece
of vocabulary that has looked ornamental.

Aurutance decomposes as
$\operatorname{Aur}(A;B)=[\log\PP(A)-\log\PP(B)]+[\log\hm(A)-\log\hm(B)]$
--- an attainability term and a Golden term,
\eqref{gs:eq:aurutance}. Rewrite \eqref{ra:eq:invest} in logarithms and
the investment bound is
\[
\underbrace{\bigl[\log\PP(\mathcal G\mid R)-\log\PP(\mathcal G\mid R')\bigr]}
_{\text{necessity purchased}}
\;\ge\;
\underbrace{\bigl[\log q'-\log q\bigr]}_{\text{weight spent}},
\]
which is exactly $\operatorname{Aur}(R;R')\ge0$ with the Golden term read
as $\PP(\mathcal G\mid\cdot)$ --- and, in the inferential form
\eqref{ra:eq:investexceed}, $\operatorname{Aur}(R;R')\ge-\log\kappa$
except with probability at most $1/(1+\kappa)$, so that the exceedance
tolerance is an additive allowance in Aurutance rather than a
multiplicative one. \emph{So Aurutance's two terms are the
cost and the necessity, and the sign of their sum is what decides.} The
decision rule reads the inequality forward --- which option should one
prefer, given what each costs and buys --- and the investment bound reads
it backward --- given that this one was realized, what must it have
bought. One quantity, two readings, and the second is why the vocabulary
was built on an economic root.

Two things this does not license, and the second is a fence this work has
already stated. It does not make the bound computable: $\PP(\mathcal
G\mid\cdot)$ is no more available than $\hm$ is, and
\S\ref{gs:sub:limit}'s missing map is missing here too, so the investment
reading is exact and unusable in the same breath and for the same reason.
And it does not make Aurutance a currency. The fence of
\S\ref{gs:sub:fence} stands unamended: the additivity is that of
logarithms across independent comparisons, and licenses no stock to hold,
spend, or owe. What \eqref{ra:eq:invest} supplies is an
\emph{exchange rate} within a single comparison --- how much necessity a
given improbability must buy for the comparison to come out as observed
--- which is a different object from a stock, and the distinction is
exactly the one that keeps the economic vocabulary honest.
\end{remark}

\begin{remark}[Where this structure has already appeared]
\label{ra:rem:investinstances}
The bound is not new to the paper; what is new is its statement in
general form. Three earlier results are instances.

\emph{The suppression bound.} Proposition~\ref{gc:prop:alpha}'s
$\alpha\lesssim\PP(A)/\PP(A')$ is \eqref{ra:eq:invest} with $R$ the
present and $R'$ a position in deep time: the independent component of
the accounting future is capped by exactly the cost ratio, the cap
being as small as the positional draw of Ansatz~\ref{ra:ans:draw}
makes that ratio. Remark~\ref{gc:rem:determinism}'s substitutability
bound is the same inequality met one section earlier, with $R$ a
realized historical fluke and $R'$ its allegedly substitutable
alternative --- the detector reading of
Proposition~\ref{gc:prop:flukes}, in the expenditure register.

\emph{The schedule's contrast.} \eqref{sch:eq:cefff}'s
$c_{\mathrm{eff}}=r\log_{10}(1+(B+\Gamma)/\lambda)$ is the price in
orders of magnitude that one wrong gate charges, and
Remark~\ref{sch:rem:tension} is the investment reading of it: the regime
that supplies high contrast is the regime in which the realized history
had to be enormously lucky, because contrast \emph{is} what luck was
spent on.

\emph{And the collapse.} \S\ref{gc:conv:collapse}'s condition
$c_{\mathrm{eff}}>\log_{10}(n(b-1))$ says the per-slot investment must
outrun the logarithm of the combinatorial supply --- necessity per slot
against alternatives per slot, which is \eqref{ra:eq:invest} summed over
a product space.

The generalisation worth noting is that all three are the same
comparison, and that the paper's recurring quantities --- $\alpha$,
$c_{\mathrm{eff}}$, Aurutance --- are the same ratio in different
coordinates. Where a reader meets a bound of the form \emph{improbability
must be matched by necessity}, it is this proposition.
\end{remark}

\begin{remark}[The exhaustiveness cost, again]\label{ra:rem:investcost}
The bound's use requires the alternative $R'$ to be genuinely available
and genuinely cheaper, and identifying alternatives is where every
argument in this cluster pays. A route that looks expensive against an
imagined cheap alternative that does not in fact reach the target
supports no inference at all, and the ``cheaper route'' is exactly the
kind of thing a reasoner constructs after the fact. The same
exhaustiveness exposure that Remark~\ref{gc:rem:routerobust} records, and
that Proposition~\ref{gc:prop:alpha}'s step~5 records, applies here and is
the bound's principal cost: \eqref{ra:eq:invest} is an identity, and
every use of it is only as good as the survey of alternatives it is run
against \textbf{[O]}.
\end{remark}

\subsection{What the target is, one step further out}\label{ra:sub:asymptote}

Past-Alignment identifies the conditioning target as the Golden Point.
The model's own structure says the true primitive sits one step further
out, and the point should be made here because it re-orders the relation
between this work and the companion by one level.

By Remark~\ref{sch:rem:asymptote}, the only conditioning under which the
schedule is worth climbing --- under which fighting for non-decaying
capacity is probabilistically preferred to external luck --- is
conditioning on survival to the \emph{infinite horizon}, and no
finite-horizon conditioning selects it. The Golden Point is then what
the asymptotic conditioning looks like when read as an event: the point
past which the schedule's climb is secured. Conditioning on the Golden
Point creates the schedule; conditioning on the asymptote is what makes
the schedule worth following; and Remark~\ref{gc:rem:finitehorizon} has
already shown that the modal form of the Golden Point --- $\exists\forall$,
some course admitting \emph{no} foreclosing continuation --- is what lets
it stand in for the asymptotic condition without inheriting the
finite-horizon robustness that would sink it.

On this reading observer-based Accordance is an inferential guide to the
asymptotic conditioning, and the axioms it helps one reach do not depend
on it deductively. The relation to the companion is then not ``observer-
conditioning is the shadow of Golden-conditioning'' but ``observer-
conditioning is the shadow of Golden-conditioning, which is itself the
event-form of asymptotic-survivorship conditioning'' --- the companion's
Optimation two restrictions down from the primitive. That is the
strongest form of the claim this work makes about its relation to the
companion, and it is graded \textbf{[S]}: a reading of the model, to
which the epistemic ceiling applies with full force.

\subsection{Ledger: the instrument, and what would refine it}\label{ra:sub:instrument}

The retrocausal accordance principle detects a conditioned event only
when that event is low-probability under the Cosmic Nexus from the prior
Aethus. It is a coarse instrument --- the elementary one, requiring only
a weight and a comparison, available immediately, and correspondingly
imprecise --- and \S\ref{gc:sub:credencemethod} shows what that
coarseness costs: an observer who knows to look will find a few
conditioned events, and those few will carry a large share of the
evidence, but finding all of them, or more than a small fraction of
their count, is beyond the instrument's grain.

What would refine it is not more data of the same kind but access to a
different object. The principle runs on the scalar stochastic flattening
of the Aethic structure --- a weight, a hazard, downstream a single
harmonic value --- and Remark~\ref{gc:rem:flattening} locates the
attribution loss exactly there: at the projection, not in the world. The
full attribute-dependency structure, and the inferential closure of what
is already known computed at scale over it, is what would turn the
principle's trunk and hubs into the whole graph. The paper records this
as the program's direction and claims nothing about its feasibility. It
notes only that the same holds for the companion's wing, and for a
specific reason: Remark~\ref{gc:rem:legibility} shows that the collapse
condition is also a legibility condition, so that under the conditioning
the record's near-determinism is exactly what would make such closure
informative rather than intractable.

\section{The Route World: what the strong reading predicts}\label{gc:subsub:routeworld}\label{rw:sec}

\begin{quote}\small
\S\ref{gc:sec} established that the record we hold is the past-cone
restriction of a conditioning on an event in our future cone, and left
open which of three hypotheses governs the forward cone. This section
develops the strongest of them --- the \emph{Route World}, in which the
Golden Point stands as a stated attribute of one's Aethus rather than as
a weight over futures --- and does so because it is the one hypothesis
with a phenomenology. If it holds, a biosphere approaching its Golden
Point experiences not fewer improbabilities but more, of a specific kind;
that is a claim about the texture of the near future, and unlike almost
everything else in this work, a claim about ours.

Everything here is \textbf{[O]}, with one stated exception: the
conduct argument of \S\ref{gc:sub:acting}, which is
\textbf{[$\Phi$]} on an adopted valuation and argues for the
hypothesis exactly as little as the rest. The section describes what
the hypothesis predicts, states what would distinguish it from its
rivals, disciplines the historical reading, connects the answer to the
epistemic ceiling, and derives the conduct the open hypothesis leaves
an agent; it argues for the hypothesis nowhere. What it supplies is the framework's most distinctive
prediction, and a reader should not finish this work without knowing the
strong reading has one.
\end{quote}

Void has been given a form it could take
(Remark~\ref{gc:rem:voidontology}), and Weak a statement of what its two
refusals commit it to (Remark~\ref{gc:rem:weakposition}). Strong needs
something different in kind. Void and Weak needed their commitments
drawn out, whereas Strong's difficulty is that it has no empirical handle
at all --- its three metaphysical questions are stated at
Remark~\ref{gc:rem:strongquestions}, and none of them is a question any
observation bears on. This section supplies the handle by describing the
world Strong asserts. The description is a consequence of the hypothesis
and not an argument for it.

\begin{definition}[The Route World; \textbf{[O]}]\label{gc:def:routeworld}
The \emph{Route World} is the world of the Strong extrapolation
hypothesis: the world in which the Golden Point stands as a stated
attribute of one's Aethus rather than as a weight over futures, so that
the conditioning is resolved throughout the cone and not merely behind
one. The name is meant in the sense \emph{RNA world} is meant --- a
regime with a characteristic phenomenology, described so that one can
ask whether one is in it.
\end{definition}

\begin{remark}[Whether ``stated'' is coherent for a future-cone
attribute]\label{gc:rem:statedcoherent}
The definition invites an objection that must be met before anything is
built on it. The framework rejects tensed determinacy: an unresolved
future outcome is a blank entry, and blanks lack tense. How can an
observer then hold as \emph{stated} an attribute whose referent lies in
their future cone --- is that not the very error the framework was built
to avoid?

It is not, and the reason is that the objection reads ``stated'' as
``resolved by the clock,'' which is the tensed reading the framework
denies. Statedness is a property of the \emph{index} --- of what the
observer's books carry --- and not of the referent's date
(Remark~\ref{gc:rem:notense}). An attribute is stated when it is a fixed entry
of the Aethus, whatever its referential event's coordinate; a stated
attribute with a future referent is no stranger than a stated attribute
with an unobserved past referent, and the framework has carried the
latter since its first postulate. What the Strong hypothesis asserts is
that the Golden Point is such an entry for us --- fixed in the books,
not merely weighted --- and \emph{that} is a substantive claim about
which Aethus we hold, not a category error about tense.

The analytic form of the same point is already in this work.
Remark~\ref{gr:rem:familyvsmeasure} distinguishes the interpolation
$\PP^{(t)}$, a family whose conditioning is still being acquired, from
$Q$, a single measure whose conditioning has been absorbed into the base
state --- and identifies the passage to $Q$ with the statement of a
future-referring attribute in present books. The Route World is the
world in which that passage has occurred. Weak is the world of the
family; Strong is the world of the measure; and the difference between
them is exactly the difference between an attribute weighted and an
attribute stated, which the framework can express and does not collapse.
\end{remark}

\subsection{Golden Optimation, and where it differs}\label{gc:sub:goldenopt}

The Route World's
signature is a strengthening of a phenomenon the companion already
names. Optimation reports that a record conditioned on observerhood
displays flukes against the Cosmic Nexus's own rates \cite{Nexus}; but
those flukes lie, at every moment, in the observer's \emph{past} cone,
because the conditioning event --- one's own existence --- is behind one.
The supply is bounded by what has already happened.

\begin{definition}[Golden Optimation; \textbf{[O]}]\label{gc:def:goldenopt}
\emph{Golden Optimation} is the same phenomenon under a conditioning
event in the future cone. Because the Golden Point lies ahead, the
flukes its occurrence requires need not have occurred yet: there are
moments at which a given required fluke lies in the observer's future
cone, and the observer's moving present passes over it.
\end{definition}

That is the whole differentia and it is sharp. Classical Optimation
never predicts a fluke; it accounts for flukes already in the record.
Golden Optimation predicts them, because the conditioning event has not
yet occurred and the requirements it imposes are still being met. A
Route World is therefore a world in which one keeps encountering
improbabilities \emph{prospectively}, and the rate at which one does so
is the hypothesis's observable content.

\begin{remark}[One phenomenon, and the label is tense-relative]
\label{gc:rem:carryover}
The two Optimations are not two mechanisms with a seam between them.
Take a fluke $F$ required by the Golden Point, occurring at Aethic time
$t_F$, and an observer at $t$. For $t<t_F$ the observer's present has not
reached it and $F$ is Golden Optimation; for $t>t_F$ it stands in the
record and is Optimation in the companion's sense. \emph{Same event, same
conditioning, different label} --- and what changed was the observer's
position, not anything about $F$.

So the label tracks the observer and not the phenomenon, which is
Remark~\ref{gc:rem:notense} applied to Optimation: an observer far enough
downstream reads our Golden Optimation as their past-cone record,
exactly as we read the Cretaceous. The carryover between the two is
therefore smooth by construction rather than by coincidence, and a Route
World contains no moment at which one regime hands over to another.
\end{remark}

\begin{remark}[The symmetry with the companion, and what the corpus is]
\label{gc:rem:rungzero}
One structural fact should be stated in this work rather than left for a
reader to assemble, since it is what tells the reader what the corpus
\emph{is} rather than what either paper is.

The companion's inference about the deep past rests on a single licence:
that an unread past outcome is a blank entry, and blanks lack tense, so
that a past fluke is open to inference at all
\cite{Aethus, NexicFoundation}. The companion calls that its
foundational rung, without which no inference about a past attribute can
be drawn. The present work's conditioning on a future event rests on
exactly the same licence: a future outcome is a blank entry, blanks lack
tense, so a future-cone attribute is open to conditioning in the same
sense a past-cone one is (Remark~\ref{gc:rem:notense}, and
Remark~\ref{gc:rem:statedcoherent} for the strong form).

So the two developments are one move run in opposite temporal directions
on one ontology. The companion reads backward from an observer to the
flukes that produced it; this work reads forward from a conditioning
event to the flukes it requires; and the principle that makes either
reading legitimate is the same principle in both. That is not a
coincidence of two papers written by one hand. It is a fact about the
ontology, and it is why the Golden Point is not an extension of Nexic
reasoning but its mirror image --- the same lens turned the other way.
The companion's forward pointer to this work, when it is written, should
be built on that symmetry and on the finding of
\S\ref{gc:sub:humanhistory}: observer-conditioning is exactly right for
the pre-technological record and provably insufficient beyond it, and a
target not indifferent to the technological character of the record
recovers the companion's weighting as its past-cone shadow.
\end{remark}

\begin{definition}[Golden causality]\label{gc:def:goldencausality}
\emph{Golden causality} names the pattern in which an event in the moving
present is accounted for by Golden Optimation --- by the Aethus's
conditioning on the Golden Point --- rather than by the Cosmic Nexus's
own rates.
\end{definition}

The name is admissible for a reason internal to the framework rather than
by indulgence. Causal talk here is not talk of a primitive relation of
production: this work's causal vocabulary labels patterns of conditioning
and dependence in the Aethic structure and takes no position on any
further relation underneath, because the framework posits none. Where
causation is not fundamental, a \emph{kind} of causation is a kind of
conditioning pattern, and Golden causality is one --- named, and carrying
no more than \eqref{gc:eq:retroaccord} licenses. What the name must not
smuggle is the register Remark~\ref{gc:rem:teleology} forbids: no event
pursues the Golden Point, and Golden causality describes what accounts
for something, never what something is for.

% [Merge: formal content moved to ra:sub:investment (Opus); this remark
% keeps the Route-World phenomenological face the sheet asked for.]
\begin{remark}[Invested Golden causality, as phenomenology]
\label{gc:rem:invested}
One element of the phenomenology deserves its name here, because it is
the form in which Golden causality is most naturally read off a record
from inside: the register of \emph{investment} --- ``look how much the
conditioning invested in this; it must have been necessary.'' The
register is earned rather than decorative:
Proposition~\ref{ra:prop:investment} is its exact content, an observed
route's improbability lower-bounding what it accomplished at a stated
exceedance confidence, and Remark~\ref{ra:rem:aurutance} is why the
economic vocabulary of \S\ref{gs:sec} was built to carry it ---
Aurutance's two terms are the cost and the necessity, and the
investment reading is the decision rule run backward. What the register
must never smuggle is agency: nothing takes the cheaper route or
declines it (Remark~\ref{ra:rem:investreg}), and the expenditure is the
conditioning's concentration of weight, never anyone's spending. A
Route World is, in this register, a world whose record reads as heavily
and legibly invested --- and \S\ref{gc:sub:historicalcases} is where
the temptation to read particular centuries that way is disciplined.
\end{remark}

\begin{remark}[The future-conditioned tradition, placed]
\label{rw:rem:futurecond}
The Strong reading has a published ancestry, and the placement
discipline owes it a paragraph: Barrow and Tipler's Final Anthropic
Principle \cite{BarrowTipler86}, Tipler's Omega Point program
\cite{Tipler94}, Leslie's axiarchism \cite{Leslie79}, and,
ancestrally, Teilhard \cite{Teilhard55} --- the branch of the
anthropic tradition that conditions on the \emph{future}, which is
the branch this work lives on (the companion engages the tradition's
observer-conditioning trunk; the future-conditioning branch is placed
here). The deltas are real and they all cut one way. The conditioning
of this work is a \emph{measure operation}, not a dynamics: no agent
at any grade, no attractor, no physical boundary condition --- where
the Omega program is a physical eschatology with cosmological
commitments, several since falsified. This work carries published
falsifiers and an epistemic ceiling that \emph{forbids} the
certified-salvation register the Omega program embraced
(\S\ref{gr:sub:ceiling}); and it is severable --- under Void the
mathematics and the critique stand with nothing conditioned on ---
where the tradition's packages fall whole. The delta that flatters
neither side: this framework purchases strictly less --- no
resurrection mechanics, no convergence guarantee, no collective
eschatology --- and prices what it does purchase.
\end{remark}

\subsection{The empirical signature, and its limits}\label{gc:sub:credencemethod}

Two things about the credence being computed here should be fixed before
the computation, since they determine what the computation is a
computation \emph{of}.

\emph{The credence is methodological, and it is a superposition.} A
reasoner in a Route World holds the Golden Point as a stated attribute
--- that is what the hypothesis says of their Aethus. But the reasoner's
\emph{credence} that they are in a Route World is computed from less
than their Aethus contains, since it is computed from the evidence
accessible to them and not from the stated attribute itself. So the
credence is an Aethic disagreeing superposition: the Strong hypothesis's
determinate conditioning weighed against a Weak-type alternative in which
the Golden Point is only weighted. The distinction is between
methodological credence --- what the evidence supports --- and
ontological credence --- what one's Aethus in fact carries --- and
everything below is about the first, and never asserts the second.

\emph{The procedure does not force the conclusion.} What it establishes,
if the flukes keep arriving, is that each next event acts \emph{as if}
the Golden Point were being conditioned on, to whatever precision the
accumulated evidence supports. That is the whole of the method's
ambition, and it is the correct one: a reasoner genuinely in a Route
World tends toward the conditioning by this route with arbitrarily high
credence, and never reaches it, and the ceiling of
Theorem~\ref{gr:thm:ceiling} is why the second clause is not a defect.

The Route World is testable in
principle, and the form the test takes is the subsubsection's main
result.

\begin{proposition}[Credence accumulation; \textbf{[P]} given the
premises]\label{gc:prop:credence}
Let $E$ be the conjunction of fluke attributes observed over an interval
of Aethic time, with $q=\PP(E\mid N_0)$ their unconditioned weight, and
suppose the Golden Point requires them, $\PP(E\mid G)\approx1$. Then the
Bayes factor for Strong conditioning against a non-conditioning
alternative over that interval is
\begin{equation}\label{gc:eq:bayesq}
\frac{\PP(E\mid G)}{\PP(E\mid \neg G)}\;\approx\;\frac1q,
\qquad\text{so}\qquad
\log\text{-evidence}\;\approx\;-\log q .
\end{equation}
In the stochastic model of \S\ref{gr1:sec} this specializes exactly:
the likelihood ratio between $Q$ and $\PP$ on $\Fc_t$ is the Doob
density \eqref{sh:eq:Qdensity}, so the accumulated log-evidence is
$\log\hm(X_t)-\log\hm(X_0)$, and \emph{the flukes one accumulates are
the persistence-harmonic function rising}.
\end{proposition}

Three things follow, and the second is the one that keeps the procedure
honest.

\emph{The evidence is the same object \S\ref{gs:sec} computes.} $\hm$ is
what the Golden Score is a ratio of (Proposition~\ref{gs:prop:score}),
so the Route World's credence and the decision rule consume one
quantity, and both wait on the same missing map
(\S\ref{gs:sub:limit}).

\emph{And that map is what stops the procedure being unfalsifiable.}
Absent it, ``what counts as a fluke'' is chosen after the fact, and a
sufficiently motivated observer finds flukes anywhere --- the
look-elsewhere problem, which would make \eqref{gc:eq:bayesq} an
instrument for confirming whatever it was pointed at. With it, a fluke
is whatever raises $\hm$, which the model fixes and the observer does
not. The methodology's rigor is therefore hostage to
\S\ref{gs:sub:limit}'s open problem, and should be described that way
rather than advertised as available.

\emph{The aggregate does not determine its own decomposition, and
identifiability is a matter of degree.} The likelihood ratio is a
function of the \emph{state}, $\hm(X_t)/\hm(X_0)$, and not a ledger of
which happenings supplied it: an accumulated $20$ nats --- a Bayes factor
near $5\times10^{8}$ --- is equally the signature of one unmistakable
event, of two hundred each amounting to a ten-percent tilt, or of two
thousand at one percent apiece, and the aggregate does not distinguish
the cases.

It does not follow that no individual event is identifiable, and the
qualification matters. A happening of conspicuously low weight under the
Cosmic Nexus is exactly the case
Principle~\ref{gc:prin:retroaccord} speaks to: its update factor on
futures is correspondingly large, so a future accounting for it is
favored in the stated ratio, and the principle is thereby the
\emph{detection instrument} for such events rather than merely a
constraint on them. Large single contributions can be found, and an
observer who knows to look will find some.

What resists is the rest, and the shape of the resistance is steep.
Writing $\theta$ for the weight at which a single event becomes
conspicuous and $\mu$ for the typical contribution, the share of total
evidence residing in conspicuous events falls like $e^{-\theta/\mu}$: at
$\theta$ a twenty-to-one improbability and a total of $20$ nats, three
large events put ninety-two percent of the evidence in plain sight,
twenty events put twenty percent, and two hundred put
none.\footnote{Exponential contribution sizes at fixed mean, for which
the conspicuous share is $(1+\theta/\mu)e^{-\theta/\mu}$ exactly.} And
under the log-spread laws this work uses elsewhere,
(Definition~\ref{sch:def:rates}), the two quantities that matter come
apart: with contributions lognormal at spread $3$ over two hundred
events, roughly \emph{one} carries sixty-three percent of the evidence.
A handful is findable, those few are decisive, and the proportion of
contributing events they represent is under one percent.

So the accurate statement is that a Route World should be expected to leave
a few conspicuous traces and a great many invisible ones, the conspicuous
few carrying a large share of the evidence and a negligible share of the
count. Evidential dominance and numerical proportion are different
quantities --- exactly as Remark~\ref{sch:rem:absolute} finds them for
orderings --- and an observer who located every conspicuous trace would
still not have located most of what happened. That is
\S\ref{gs:sub:fence}'s class-versus-instance discipline arriving from the
evidential side.

Three tasks are worth separating here, because the paragraph above runs
two of them together and the third is the one with a future.
\emph{Detecting} the aggregate needs no localization whatever:
\eqref{gc:eq:bayesq} reads the accumulated ratio off the state, and two
hundred invisible ten-percent tilts deliver the same decisive
$5\times10^{8}$ as one conspicuous event would. \emph{Localizing the
tail} is what Principle~\ref{gc:prin:retroaccord} does. \emph{Localizing
the bulk} has no instrument at all, and the reason is specific: the
accordance principle is a \emph{marginal} test --- it interrogates one
attribute's weight under the Cosmic Nexus and infers from that weight
alone --- so an event contributing at a magnitude no marginal test
resolves is invisible \emph{to it}, which is a fact about the instrument
and not about the event.

That distinction leaves room, and the room should be marked rather than
conceded. The accordance principle is the elementary instrument here: it
requires only a weight and a comparison, it is available immediately, and
it is correspondingly coarse. Nothing establishes that it is the best
available in principle. What a finer instrument would have to do is
compute an event's contribution to $\hm$ \emph{without routing through
that event's marginal improbability} --- and an instrument of that
description is not hypothetical, it is
\S\ref{gs:sub:limit}'s missing identification map, which is exactly a
rule assigning states to events by their own structure rather than by
their rarity. The Route World therefore acquires the same open problem
\S\ref{gs:sec} has, from the other direction, and a second reason to want
it solved: the map would convert localization from a tail-only
capability into a general one. Until then, the coarse instrument is what
there is, and the estimate that a handful of traces is findable while the
bulk is not should be read as an estimate about present instruments
rather than a claim about what a Route World conceals \textbf{[O]}.

\begin{remark}[The attribution is lost in the representation, not in the
world]\label{gc:rem:flattening}
It is worth being exact about where the localization problem comes from,
because the answer says what a better instrument would have to do rather
than merely that one is wanted.

The Aethic object is a structure: attributes standing in dependency
relations, with weights over the fan. The model of \S\ref{gr1:sec}
projects that structure onto a scalar --- a hazard, and downstream a
single harmonic value --- and the projection is many-to-one by a wide
margin. \emph{The attribution information is discarded at the
projection.} Nothing about the world hides which events contributed;
the representation this work computes in does not carry the answer,
because a number cannot. That is why the identification map is hard: it
is an inverse to a projection that discarded what the inverse needs, and
inverses of that kind are recovered by supplying structure rather than by
supplying data.

It also says what the accordance principle is finding. A marginal test on
weights locates the attributes of high dependency degree --- the ones on
which much else turns, and which are correspondingly improbable in
isolation --- so what the principle recovers is a \emph{skeleton}: the
trunk and the hubs of the dependency graph, and not its edges. Working in
the unflattened structure would be working with the graph rather than its
scalar summary, and there the bulk contributions are not small numbers
lost in an aggregate but edges with endpoints. Whether the framework's
dependency structure supports inference at that grain --- and what
volume of recorded attribute-data would make the closure computable --- is
open, and is a different open problem from the identification map rather
than the same one \textbf{[O]}.
\end{remark}

\begin{remark}[The collapse condition is also a legibility condition]
\label{gc:rem:legibility}
One consequence runs the other way and bears on the companion rather than
on this work. The companion's collapse condition,
$c_{\mathrm{eff}}>\log_{10}(n(b-1))$, is stated as a fact about
\emph{weight}: above it, the realized ordering carries a share of order
unity. It is equally a fact about \emph{inference}, and the second
reading has not been drawn.

Take the effective number of live alternatives to be the exponential of
the entropy over the tier structure. At twenty interrogable slots of ten
values the nominal space is $10^{20}$; at $c_{\mathrm{eff}}$ just above
the threshold it is a few hundred; at $c_{\mathrm{eff}}=3$ it is four;
and by $c_{\mathrm{eff}}=5$ it is one to within a percent. The collapse
is sharp, and what collapses is the space a reconstruction must search.
So the same contrast that makes the record improbable makes it
\emph{legible}: reconstruction of a high-contrast natural history
requires data scaling with the surviving handful rather than with the
nominal branching, which is the precise content of the intuition that
knowledge here should grow quickly once methods are good enough to
exploit the concentration. The threshold such methods must reach is the
collapse condition itself.

Two cautions. The concentration is \emph{conditional} --- it is what the
conditioning purchases, and without it the space is the nominal one ---
so the legibility cannot be cited as unconditional evidence for the
conditioning without circularity. And concentration is not determinism:
one ordering carrying a share of order unity leaves the others weighted
rather than excluded, and Definition~\ref{gc:def:epsdep}'s grading is exactly the
difference. What is available honestly is a prediction: if the
conditioning holds, reconstructions of deep history should converge on
less evidence than a flat-prior analysis of the same branching would
require, and that is checkable without any fluke being counted
\textbf{[O]}.
\end{remark}

One further consequence cuts against the methodology's own convenience.
How large the required flukes are depends on how much work the
conditioning has to do: a position from which the Golden Point is already comparatively
likely requires only small corrections, while deep time, from which it
was not, required the enormous ones the record displays. A biosphere
improving its position should therefore expect its Golden-Optimation
events to grow \emph{less} conspicuous, so the evidence available to
\eqref{gc:eq:bayesq} thins exactly as the situation it reports on
improves, and the method is at its most informative where the news is
worst \textbf{[O]}.

\emph{The ceiling permits the ambition and forbids the conclusion.} By
Theorem~\ref{gr:thm:ceiling} the two measures are mutually absolutely
continuous on every finite window and singular only on the
infinite-horizon $\sigma$-algebra, so the likelihood ratio is finite at
every finite time: credence may rise without bound and \emph{never
certifies}. That is not a limitation discovered here but the theorem's
own shape, and it is the exact sense in which the procedure can approach
the conclusion asymptotically while being barred from reaching it.


\subsection{The Lindy step}\label{gc:sub:lindy}

\begin{definition}[The evidence coordinate]\label{gc:def:evidence}
Write $\Lc_t$ for the log-evidence a Route World has supplied by Aethic
time $t$ --- the log-likelihood ratio of \eqref{gc:eq:bayesq}, which in
the model of \S\ref{gr1:sec} is $\log\hm(X_t)-\log\hm(X_0)$ --- and
$\Lc_\ast$ for the total the conditioning ultimately supplies. Both are
nondecreasing, and $\Lc_\ast\ge\Lc_t$ by construction.
\end{definition}

\begin{ansatz}[Copernican position in the evidence coordinate]
\label{gc:ans:copernican}
An observer's position in the accumulation is uniform:
$\Lc_t/\Lc_\ast\sim\mathrm{Unif}(0,1)$.
\end{ansatz}

\begin{proposition}[The Lindy law for accumulated evidence;
\textbf{[P]} given Ansatz~\ref{gc:ans:copernican}]\label{gc:prop:lindy}
Conditional on $\Lc_t=\ell$, the total is Pareto with unit shape and
scale $\ell$:
\begin{equation}\label{gc:eq:lindylaw}
\Pr\bigl[\Lc_\ast>x\,\bigm|\,\Lc_t=\ell\bigr]=\frac{\ell}{x},
\qquad x\ge\ell,
\end{equation}
whence the median total is $2\ell$, the geometric mean is $e\ell$, the
arithmetic mean diverges, and $\Pr[\Lc_\ast>10\ell]=1/10$. In the
weight coordinate $q=e^{-\ell}$ these read: median total weight $q^{2}$,
geometric-mean total $q^{e}$, and a one-in-ten chance the total is
$q^{10}$ or smaller.
\end{proposition}

\begin{proof}
$\Lc_\ast=\ell/U$ with $U\sim\mathrm{Unif}(0,1)$, so
$\Pr[\Lc_\ast>x]=\Pr[U<\ell/x]=\ell/x$ for $x\ge\ell$; the median solves
$\ell/x=\tfrac12$; $\E[\log(1/U)]=1$ gives the geometric mean; and
$\E[1/U]$ diverges.
\end{proof}

\begin{remark}[Reading the law, and the coordinate it must be read in]
\label{gc:rem:lindy}
Three consequences, and the third is a constraint on how the law may be
used at all.

\emph{The point estimate is conservative.} Quoting $q^{2}$ quotes a
median, and the law's mean diverges, so a Route World's total evidence
is more likely to exceed the estimate badly than to fall short of it
mildly. That asymmetry is the Pareto tail and not an artefact of the
parametrization.

\emph{Credence must be tracked in log-odds.} Posterior odds are
$O_\ast=O_0e^{\Lc_\ast}$, and since $\Lc_\ast$ is Pareto with unit
shape, $\E[O_\ast]$ diverges violently --- an expected-odds calculation
here returns infinity and reports nothing. What is well behaved is
$\log O_\ast$, whose median exceeds the present value by exactly
$\Lc_t$. Every statement in this subsubsection is accordingly a
statement about log-evidence.

\emph{And the Copernican step must be taken in the evidence coordinate,
not in time.} Applied to elapsed time the multiplier depends on the rate
at which flukes arrive --- four if the accumulation is quadratic,
$\sqrt2$ if it is square-root --- while applied to $\Lc$ it is two
regardless. Since the attributes at issue are arbitrary Aethic events
with no natural arrival rate,
Definition~\ref{gc:def:evidence}'s coordinate is the only one in which
Ansatz~\ref{gc:ans:copernican} is well posed.
\end{remark}

\begin{remark}[The Copernican premise is not licensed by Accordance, and
is the section's weakest posit]\label{gc:rem:copernicancost}
Ansatz~\ref{gc:ans:copernican} is posited rather than derived, and the reason is sharper than caution: \emph{the companion's
own foundation declines the inference it instantiates.}

The Accordance Principle binds attribute by attribute, and its extraction
is explicit that the global birth-rank inference is thereby \emph{not}
recovered --- a position in a completed roster has no attribute-wise
translation, and the retort that the position could itself be made an
attribute fails because that attribute's answer space presupposes the
completed roster it was meant to replace \cite{NexicFoundation}.
Ansatz~\ref{gc:ans:copernican} is a position-in-a-roster claim of
exactly that form, with $\Lc_\ast$ the completed roster. So it does not
inherit Accordance's licence, and Remark~\ref{gc:rem:notense}'s non-privilege
does not supply one either: non-privilege says no position is
distinguished, which is compatible with the accumulation admitting no
well-defined uniform draw at all.

What it therefore is: an independent Copernican posit, carried at
\textbf{[$\Phi$]}, load-bearing for
Proposition~\ref{gc:prop:lindy} and for nothing else in this work.
Everything upstream of it --- \eqref{gc:eq:bayesq},
\eqref{gc:eq:retroexceed}, the noise floor, the identifiability
analysis --- stands without it, since those concern evidence already
accumulated rather than evidence to come. A reader who rejects the
ansatz loses the forward extrapolation and keeps the method.

\emph{Falsifier and repair.} It fails if observer-density along
the accumulation is non-uniform --- if, for instance, observers cluster
where evidence has already accumulated, which is not obviously false and
would bias the multiplier downward. Establishing the density is the
repair, and it is open \textbf{[O]}.
\end{remark}

\subsection{What would distinguish the three worlds}
\label{gc:sub:discriminate}

A hypothesis with a phenomenology owes a statement of what would tell it
from its rivals, and this is where the Route World's status is decided.
The answer has two halves and they point in opposite directions.

\paragraph{What differs.} The three worlds make different predictions
about the \emph{rate} at which an observer encounters improbabilities
against the Cosmic Nexus's own rates.

\emph{Void.} Nothing constrains the forward cone, so improbabilities
arrive at the Cosmic rate: after one's present, flukes are as rare as the
unconditioned weighting says, and Optimation is a fact about the record
only. The evidence flux $\tfrac{d}{dt}\Lc_t$ is zero.

\emph{Weak.} The forward cone is weighted by Golden-Point probability, so
improbabilities arrive at a rate above Cosmic and below what full
conditioning gives --- and, since Weak's anchor does real work, at a rate
depending on one's position. The flux is positive and sub-maximal.

\emph{Strong.} The conditioning is resolved throughout the cone, so
required improbabilities arrive at the full conditioned rate, and by
Proposition~\ref{gc:prop:credence} the accumulated evidence
$\Lc_t=\log\hm(X_t)-\log\hm(X_0)$ rises without bound. The flux is
positive and maximal.

The discriminating quantity is therefore the flux, and it orders the
three: zero, intermediate, maximal. Two things bear on reading it.
Remark~\ref{gc:rem:carryover} says the label attaching to any given fluke
is tense-relative, so an observer cannot distinguish the worlds by
inspecting individual events; only the \emph{rate} discriminates. And by
the identifiability analysis of \S\ref{gc:sub:credencemethod}, the rate
must be read from the aggregate $\Lc_t$ rather than from a ledger of
attributed events, since the aggregate does not decompose. What one
compares is a log-likelihood ratio against a Cosmic-rate baseline: Void
predicts it stays near zero, Strong that it grows.

\paragraph{What does not.} By Theorem~\ref{gr:thm:ceiling}, $Q$ and $\PP$
are mutually absolutely continuous on every finite window and singular
only on the infinite-horizon $\sigma$-algebra. So the discriminating
quantity is finite at every finite time. Evidence for Strong can
accumulate without bound and never certify; evidence for Void can only
ever be the absence of accumulation, which is compatible with Weak at low
flux and with Strong at early time. \emph{No observation of bounded
duration decides among the three.}

That is not a limitation the Route World discovers about itself but the
ceiling's own shape, and it should be read as the point rather than as a
disappointment. What the framework predicts is a \emph{trend} --- a
log-likelihood ratio that grows under Strong, stays flat under Void, and
grows sub-maximally under Weak --- and a trend is exactly the kind of
prediction that accumulates credence without converting to certainty. A
reader who wants a decisive test is asking for something the framework
proves it cannot supply; a reader who wants a discriminating quantity has
one.

\paragraph{The scoring.} The Route World therefore has a phenomenology, a
discriminating quantity, and no decisive test --- and the last is a
theorem. That is a weaker empirical status than a falsifiable prediction
and a stronger one than a metaphysical hypothesis with no observational
content, which is where Strong stood before this section. What the
section has done is move Strong from the second category into the first,
and no further.

% [Merge: Opus's historicalcases + acting adopted as base; Fable's
% mattering-residue rendering retired; complacency second answer carries
% the anatomy-face identification and falsifier folded in.]
\subsection{Why the historical cases do not discriminate}
\label{gc:sub:historicalcases}

A reader who accepts the previous subsection's discriminating quantity
will nonetheless want to reach for something nearer to hand, and the
material is genuinely there: read a consequential century closely and
what had to go right is startling. This subsection says what such cases
do and do not support, because the temptation to over-read them is
strong and the framework's own guards are what keep the reading honest.

\paragraph{The material.} If one is in a Route World then one has been in
it throughout, and any century should show it. The scientific
enlightenment is the case most often reached for, and it repays the
reach: the transmission of algebra; the printing press and the literate
market it made \cite{Eisenstein79}, and --- less often noticed and more interesting --- a
theological temperament under which fitting exact laws to the world was
a natural thing to attempt rather than an eccentric one --- a standing
thesis of the historiography \cite{Merton38}. Counterfactuals
of the usual shape suggest themselves. Had Rome persisted in a form that
left no room for the religious configuration that followed, it is not
obvious that the temperament distinguishing Renaissance mathematics from
Archimedean mathematics would have arisen; and the material prerequisites
have their own contingencies, paper among them. Each century then sets
conditions for the next, which is the schedule structure of
\S\ref{sch:sub:schedule} in a documentary register: gates in order, with
the ordering doing work.

\paragraph{What they support.} The cases are the human-history leg's own
material (\S\ref{gc:sub:humanhistory}), and what they support is what
that leg argues: that the conditioning is sensitive to the technological
character of the record, which observer-conditioning is not. That is a
claim about \emph{which target}, and these cases bear on it.

\paragraph{What they do not.} They are past-cone cases, and the
distinction that matters is exactly the one
Definition~\ref{gc:def:goldenopt} draws. A fluke already in the record is
\emph{consistent with} Golden Optimation and is equally consistent with
ordinary Optimation --- the conditioning event behind one, the supply
bounded by what has happened. By Remark~\ref{gc:rem:carryover} the label
attaching to any past fluke is tense-relative, and a fluke that lies
behind an observer carries no information about whether the conditioning
extends ahead of them. So historical cases are \emph{Past-Alignment
consistent}, which is the milder claim, and they do not discriminate
Strong from Weak or from Void. Only the flux does, and the flux is a rate
in the observer's forward cone.

The stronger phrasing --- \emph{documentable cases of Golden causality}
--- should accordingly be resisted. By
Definition~\ref{gc:def:goldencausality} Golden causality names events
accounted for by the Golden conditioning \emph{rather than} by Cosmic
rates, and establishing that of a particular event requires establishing
the conditioning, which is what is at issue. What is documentable is
consistency, and consistency with a hypothesis is not evidence for it
over a rival the same data also fit.

\begin{remark}[Three guards, and the one that bites hardest]
\label{gc:rem:historyguards}
The historical material is where this work's own cautions apply with
most force, and they should be applied before the material is used
rather than after it is challenged.

\emph{The scan is enormous.} \eqref{gc:eq:accordexceed}'s multiplicity
guard bounds expected $\kappa$-fold violations among $n$ interrogated
pairs by $n/(1+\kappa)$, with the standing condition that $n$ counts
attributes \emph{scanned} and not merely those interrogated. For a
century read by historians, the scan is every attribute anyone has ever
found notable, and an honest $n$ is enormous. The noise floor is
correspondingly high, and a historical trace must stand clear of it by a
wide margin before it is worth anything. This is the guard that bites
hardest and it is the one most easily skipped.

\emph{Substitution is the question, not occurrence.}
\S\ref{gc:sub:humanhistory} audits by \emph{role} and not by occupant:
the question is never whether Christianity in particular arose, but
whether anything arose that discharged the role. Applied to the
temperament case this cuts both ways, and interestingly. Material
prerequisites look substitutable --- many routes lead to mass literacy,
and the printing press is one occupant of that role. A specific epistemic
disposition looks less so, since it is not obvious what else would have
discharged it; and if that holds, temperament is a \emph{serial}
attribute in the sense of Remark~\ref{gc:rem:serialcost} where the technology
is parallel, and contributes correspondingly more to $R_1$. That is a
conjecture about substitutability and not a finding, and it is exactly
the kind of claim the serial-attribute programme exists to settle
\textbf{[O]}.

\emph{And narrative bias is prior.} Remark~\ref{gc:rem:narrativebias}
states the objection and the asymmetry that answers it, and nothing in
this subsection escapes it: dramatic near-misses make better stories and
are more often recorded, so a century dense with them may be reporting a
fact about historiography. The test remains the connectivity asymmetry
that remark names, and no accumulation of vivid cases substitutes for it.
\end{remark}

\subsection{Acting under the reading, when nothing decides}
\label{gc:sub:acting}

\S\ref{gc:sub:discriminate} closes on a theorem: no observation of
bounded duration decides among the three worlds. A reader is then
entitled to ask what follows for conduct, and the answer is not

What the material is for, then, is bounded and worth stating: candidate
entries at the Past-Alignment grade, awaiting the identification map
like everything else, with the file to be assembled \emph{dull entries
first} --- print economics, paper supply, the price and reach of
literacy, before any battlefield --- exactly as the coal outweighs the
Granicus in \S\ref{gc:subsub:stockpile}. Assembling that file is the
concrete research direction the backward reading opens, and the guards
above are its methods section.

% [Instates the 5:22 PM Aug 22 sheet entry: backward-chaining
% conditional necessity; the Black Death chain; the attribute question;
% the Europe question; the dark-investment fence.]
\begin{remark}[The chain read backward: necessity propagates through
prerequisites]\label{gc:rem:backwardchain}
The material above supports one further operation, and stating it
carefully both extends the subsection and generates its next research
questions.

\emph{The operation, and its licence.} Under the Golden target,
``$A$ had to happen'' has a precise form: it is the retrocausal
accordance reading run at a vantage \emph{before} $A$ ---
Principle~\ref{gc:prin:retroaccord} with the exceedance grade it
carries, the label's tense supplied by Remark~\ref{gc:rem:carryover}.
And the form chains. If $A$ is $G$-required and $A$ presupposes $B$
--- $\PP(A\mid\neg B)\approx0$, with $\neg B$ not itself
measure-starved --- then $B$ is $G$-required at the grade of the
weaker link \textbf{[S]}: conditional necessity propagates backward
through presupposition, attribute by attribute, never period by
period.

\emph{The worked chain.} The record supplies the instance, and its
middle link is already in this work's books. The wage structure on
which Allen's account of mechanization rests is downstream of the
demographic rupture of the Black Death --- the serial-contingency
program of \S\ref{gc:sub:humanhistory} carries exactly this
candidate --- and the Renaissance's dependence on the same rupture is
a standing historiographical thesis \cite{Herlihy97,Belich22}. So two
independent routes into modernity pass through one fourteenth-century
shock, and the chain then runs backward on its own licence: the shock
presupposes the configuration it struck --- the demographic density,
the trade routes that carried transmission, the institutional
structure whose rupture did the work --- so the medieval period
\emph{up to those attributes} inherits the requirement; and that
configuration presupposes, in part, the fall-shape of Rome already
entered above as route-deletion. Under the Golden target, then, the
medieval period had to happen --- in the only sense the framework
licenses: at the grain of the attributes the chain actually consumes,
and at the grade of its weakest link.

\emph{The question each link generates.} ``Up to which attributes''
is not a rhetorical flourish but the audit of
Remark~\ref{gc:rem:historyguards} applied link by link: the
requirement attaches to the \emph{role} --- a demographic rupture of
sufficient depth to reset labor scarcity and institutional
equilibrium --- and whether \emph{Yersinia pestis} in that century
was its only available occupant is the substitution question, open
\textbf{[O]} and now precisely posed. The invested reading of
Proposition~\ref{ra:prop:investment} prices the answer's stakes: the
conditioning concentrates on the cheapest sufficient shocks, so if
cheaper occupants of the rupture-role existed and did not run, the
realized route's cost is evidence they were not in fact sufficient.

\emph{The Europe question, posed at its correct grain.} The chain as
realized runs through Europe, and the association invites two errors
that should be foreclosed in the same sentence that admits the fact.
The association is \emph{occupancy} until audited: whether comparable
rupture-to-reconfiguration routes stood open elsewhere --- the same
pandemic struck the Islamic world; China carried its own
demographic shocks --- is the role-level question, and the framework
supplies its form and no verdict. And necessity is not merit:
by Remark~\ref{gc:rem:teleology} and Remark~\ref{ra:rem:investreg}
nothing took the European route or declined another, no people is
ranked by a conditioning's concentration, and a route's being
required speaks to the narrowness of what reaches the target, never
to the worth of what occupied it.

\emph{The dark-investment fence, stated at full strength.} This
remark's least comfortable consequence should be owned rather than
softened: under the Golden target a mass-death event is conditionally
required, and the framework must say exactly what that does and does
not mean. It is a claim about route-narrowness --- that among
$G$-reaching continuations from the thirteenth century, those through
a rupture of that depth carry nearly all the weight --- and it is the
invested reading at its darkest, enormous realized cost as evidence
of how few cheaper routes there were. It licenses no providence, no
``for the best,'' no valorization of catastrophe, and no comfort:
the conditioning does not act, and an accounting is not an
absolution. A reader who feels the sentence ``the Black Death had to
happen'' as monstrous is reading its register correctly; what the
framework adds is only that the monstrousness is located in the
narrowness of the routes, not in any agency that chose them. The
whole remark is conditional on the Golden target and on
historiographical inputs carried at \textbf{[C]}/\textbf{[S]};
under Void it is history, unconditioned and unrequired.
\end{remark}

``nothing.'' It is a dominance argument, and it does not require
believing the Route World obtains.

\paragraph{The argument.} Stated in steps, so each can be refused
separately.

\begin{enumerate}\setlength\itemsep{2pt}
\item \emph{The Route World's phenomenology is categorical, not
enumerative.} What distinguishes it from the standard world is a rate
and a structure --- prospective improbability, tense-relative labelling,
evidence accumulating as $\Lc_t$ --- and not a list of instances. So it
can be adopted as a reading and applied to circumstances not on the
list.
\item \emph{If one is in the Route World, the world's phenomenology is
that one.} True by definition of the hypothesis.
\item \emph{One wants one's interpreted phenomenology to match the
world's.} This is the epistemic norm the whole work runs on and is
assumed rather than argued.
\item \emph{So if one is in the Route World, one should read and act
under Route-World phenomenology.} From 2 and 3.
\item \emph{Where the Departure Point obtains, the adopted ordering
ranks nothing.} By \S\ref{gs:ssub:departric}: past $\mu G$ the viability
event has probability zero, $Q$ is undefined as a conditional measure,
and every quantity indexed to it is undefined rather than small. This is
a claim about an index and not about the world, and it is
$\Phi6$-conditional throughout.
\item \emph{Pruning the Aethae in which departure obtains leaves exactly
the Route World.} Conditioning on non-departure is conditioning on the
viability event $\Imm$, and by
Remark~\ref{gr:rem:familyvsmeasure} that conditioning \emph{is} the
passage from the family $\PP^{(t)}$ to the measure $Q$ --- the
conditioning absorbed into the base state, which is the statement of a
future-referring attribute in present books. The Route World is the
world of the measure. So the residue is not merely Route-World-like; it
is the Route World, and the identification is exact.
\item \emph{Therefore the states in which one's conduct is ranked at all
are exactly the states in which step~4 applies.}
\item \emph{Therefore one should act under Route-World phenomenology
now}, since every alternative state is one the adopted ordering reads as
a single point.
\end{enumerate}

\paragraph{What the conclusion is.} \emph{Methodological adoption, not
ontological credence.} The argument does not say one is in the Route
World, and supplies no evidence that one is; it says that the states in
which the question of how to act has an answer are the states in which
the Route-World reading is the correct one. That distinction is the one
\S\ref{gc:sub:credencemethod} draws between what the evidence supports
and what one's Aethus in fact carries, and this argument moves only the
first.

The shape is the forced-choice argument of \S\ref{gs:ssub:departric}
generalised from a single choice to a reading. There the safe course was
scored at zero because the survival recommending it lay inside an
indifference class; here every departed world is scored at zero for the
same reason, and what remains is not a small good against a large one
but a comparison the ordering can make against comparisons it cannot.
Nothing unbounded enters, and no small probability is multiplied by a
large payoff.

\begin{remark}[The complacency objection, and the two answers it
gets]\label{gc:rem:complacency}
The objection is immediate: if one acts as though one is on a
conditioned route, does one not thereby relax --- believing oneself
saved, and so ceasing to do the work the conditioning is supposed to
require? There are two answers, and they are of very different
strengths.

\emph{The framework's own answer, which is a theorem.} The inference
from ``I am on a conditioned route'' to ``I may relax'' requires that
one's own effort be identifiable as inessential to the conditioning ---
and \S\ref{gc:sub:routefences} establishes that nobody can determine
whether their own contribution is among the requirements, since by
Theorem~\ref{gr:thm:ceiling} no finite observation places a system on
the interpolation at all. The premise the complacency
inference needs is therefore unavailable to any finite observer --- which
is a stronger statement than \emph{unverified} and a weaker one than
\emph{indeterminate}, and Remark~\ref{gc:rem:determinacygap} is why the
difference matters here. And
Remark~\ref{ac:rem:bridge} closes the other route to it: an individual
is not a lever on their biosphere's hazard, so ``my effort is
dispensable'' and ``my effort is decisive'' are equally unavailable, and
the Route World supplies no licence to either. \emph{Acting under the
Route-World reading is compatible with no relaxation of effort
whatever}, and the reading itself is what forbids the inference.

\emph{The author's empirical bet, which is not.} A second answer has
been offered and should be recorded at its own grade rather than folded
into the first: that populations do not in fact behave this way ---
that individual instability is averaged out at scale, and that a
competent society interposes barriers between any one actor and
catastrophic capacity, so that the complacent reading could not be acted
on even where it were held. The second half is at least legible in the
framework's own terms rather than an appeal to hope: an arrangement in
which no acquired capability can be turned inward by any single hand is
what $\rho>0$ \emph{is} at the grain of institutions
(Definition~\ref{gr1:def:anatomy}), so the wager is precisely the wager
that one's biosphere carries favourable anatomy at that grain --- which
makes it statable, and not thereby true. It is a claim about how human
institutions actually work, offered as a wager rather than a result,
with a concrete falsifier --- evidence that awareness of the hypothesis
degrades collective self-preservation rather than being absorbed by the
aggregation and the barriers --- and it sits uncomfortably beside
\S\ref{ad:sec}, which models institutional balance as something that
fails under exactly the adversarial loading the wager assumes away. The
paper records the bet with its falsifier and does not endorse it
\textbf{[$\Phi$]}; the first answer is the one the argument leans on,
and it leans on nothing empirical.
\end{remark}

\begin{remark}[Which step to attack]\label{gc:rem:actingcost}
The argument's steps are not equally exposed, and a reader deciding
where to push should be told which is which.

Steps 2 and 4 are definitional or immediate. Step 6 is a theorem ---
the identification of the pruned residue with the Route World is
Remark~\ref{gr:rem:familyvsmeasure}, and if it fails, the family and the
measure were never distinct and much else fails with it. Step 3 is an
epistemic norm this work assumes and does not defend; a reader who wants
their interpretation \emph{not} to track the world has other problems.

The exposed step is 5, and it is exposed in one specific way. The
indifference past $\mu G$ is an indifference \emph{of the adopted
ordering}, and an agent who holds values the ordering does not encode
--- who cares about suffering in a departed world, say --- is not
indifferent there, and for them the pruning removes states that still
matter. The argument is therefore $\Phi6$-conditional in the same sense
as everything downstream of \S\ref{ac:sub:countermeasure}, and it is
weaker than it looks for readers whose evaluative commitments outrun the
one premise this work asks them to adopt. That is not a defect to be
repaired but the scope the argument has, and step 1 is the other place to
push: if the Route World's phenomenology is only the sum of its
instances, it cannot be adopted as a reading and step 4 has nothing to
apply.
\end{remark}

\subsection{Fences}\label{gc:sub:routefences}

% [Merge: Opus's determinacygap adopted; Fable's all-three-worlds
% no-selection point folded in as its closing route.]
\begin{remark}[Determinate is not accessible, and why the distinction
cannot be relaxed here]\label{gc:rem:determinacygap}
An objection is available against the fences below, and it is good enough
that stating it is better than hoping no reader finds it. It runs as a
reductio.

\emph{The objection.} The fences say that whether one's own contribution
is among the conditioning's requirements is a thing no observer can
establish. But the Route World is precisely the hypothesis that the
conditioning is \emph{resolved} --- the Golden Point stated rather than
weighted, throughout the cone. If it is resolved, the requirements are
settled; if they are settled, there is a fact about whether any given
contribution is among them. So either that fact is available, in which
case the fences fail and the complacency inference goes through --- or it
is not available even in the resolved case, and one wonders what
``resolved'' was supposed to mean. Push the first horn and the Route
World undermines the conduct it recommends; push the second and it looks
like a hypothesis with no content. Either way, so the objection
concludes, the Route World is incoherent and its incoherence is a proof
that we are not in one.

\emph{The answer, and it is not a concession.} The objection equivocates
on ``available.'' Statedness is a property of an \emph{Aethus} --- of
what the world's books carry --- and the framework is explicit that an
Aethus is deductively closed and indexed without regard to whether its
bearer can compute over it \cite{Aethus}. So in a Route World the
requirements are determinate \emph{in the books} and remain inaccessible
\emph{to the reasoner}, and no contradiction arises, because those are
different predicates of different objects. What
Theorem~\ref{gr:thm:ceiling} forbids is not there being a fact but
\emph{certification of it from observation}: $Q$ and $\PP$ are mutually
absolutely continuous on every finite window, so no finite record
distinguishes them. That is a limit on inference and not on determinacy.

\emph{Two corrections the objection earns.} First, the fences should not
have said \emph{in principle}, and above they no longer do: the measures
are singular on the infinite-horizon $\sigma$-algebra, so an observer of
unbounded duration would distinguish them, and what the ceiling supplies
is unavailability to any \emph{finite} observer. Second, the resolved
case is not contentless. What ``resolved'' means is
Remark~\ref{gr:rem:familyvsmeasure}'s distinction --- one is under the
measure $Q$ rather than the family $\PP^{(t)}$ --- and that has
consequences the section spends its length on, including the evidence
flux that orders the three worlds. A hypothesis can be substantive and
its truth still uncertifiable; this work's own epistemic ceiling is the
standing example, and it applies to this work's conclusions as firmly as
to its targets.

\emph{And the structure is one the section already carries.}
Remark~\ref{gc:rem:flattening} makes the same distinction on the
evidential side: the attribution information is discarded \emph{at the
projection}, not in the world, so which events supplied the accumulated
evidence is determinate and not recoverable. Determinate-and-inaccessible
is the framework's normal condition rather than a special pleading
introduced to save the fences.

\emph{And a route that closes the leverage without the diagnosis.}
Even were the objector granted the availability premise arguendo, no
reductio against the Route World follows, for a structural reason
independent of the equivocation above: the ceiling is a property of the
model under every hypothesis of \S\ref{gc:sub:discriminate} --- Void,
Weak, and Strong alike --- so a premise that falsifies all three worlds
at once selects among none of them. What an exhibited certification
would refute is the mutual-absolute-continuity structure itself, at the
site where its empirical exposure lives (\S\ref{sh:sub:numerics}), and
the Route World would stand or fall with its rivals rather than before
them.
\end{remark}

\begin{remark}[The robustness argument, and why it is not relied on]
\label{gc:rem:routerobust}
A stronger conclusion suggests itself and should be marked as the
weakest thing here. If each conservative alternative --- conditioning on
some nearer, more peripheral future event rather than on the Golden
Point --- is one whose own passage will dissolve whatever objectivity it
seemed to have, then only a target at the asymptote remains stable under
the motion the argument itself describes; and one may be tempted to
conclude that nothing \emph{but} the Golden Point can be coherently
conditioned on.

Two reasons not to lean on it. It requires that the peripheral
alternatives be exhaustively surveyable, which is the same
exhaustiveness the companion's own arguments purchase carefully and this
one assumes. And it is the argument's only step that is not a
likelihood computation, so a reader who accepts \eqref{gc:eq:bayesq} and
\eqref{gc:eq:retroaccord} and rejects this paragraph loses nothing else.
The accumulation is what gives the Route World teeth; the robustness
claim is a conjecture stated beside it \textbf{[O]}.
\end{remark}

\begin{remark}[What the Route World does not license]
\label{gc:rem:routefence}
One reading has to be closed explicitly, because it is the natural one
and it is dangerous.

If required flukes occur, then --- the reading goes --- an individual
contemplating some improbable feat should ask not whether they will
succeed but whether the Golden Point requires their succeeding, and
proceed accordingly. Nothing in this work supports that, and two of its
results forbid it. By Remark~\ref{ac:rem:bridge} an individual is not a
lever on their biosphere's hazard: what they hold is a position, not a
purchase, and Golden Optimation is a statement about which histories one
finds oneself in rather than about what any agent can bring about. And
by Theorem~\ref{gr:thm:ceiling} nobody can determine whether their own
feat is among the requirements, since no finite observation places a
system on the interpolation at all. The premise of the reading is therefore
unavailable to any finite observer, which is not the same as there being
no fact of the matter (Remark~\ref{gc:rem:determinacygap}).

What survives is weaker and is the defensible form: a Route World is one in
which improbabilities are more common than the Cosmic Nexus's rates
predict. It is not one in which any particular improbability is owed to
any particular person, and an agent who acts on the stronger reading is
not applying this framework but discarding
\S\ref{gr:sub:ceiling}. Remark~\ref{gs:rem:trope} fences the neighboring
case in the same terms and for the same reason.
\end{remark}

\begin{remark}[The phenomenology extended: from events to constitution]
\label{gc:rem:constitution}
The Route World as developed in this section is a claim about
\emph{events}: which occurrences carry weight, how the record reads
under the conditioning, what the forward flux does. Its Strong reading
supports one further extension, developed where the categorical
machinery lives rather than here: from the structure of events to the
constitution of things --- the universe's contents as
weight-carrying excursions on the Aethus that states the conditioning
attribute (\S\ref{to:sub:goldenaethus}). Nothing about the Route
World's own claims strengthens or weakens by the extension, and its
fences govern there unweakened; the pointer is recorded here because a
reader of this section should know the hypothesis's full reach before
weighing it.
\end{remark}

\begin{remark}[A consequence for periodization]\label{gc:rem:periodization}
One application is worth recording because it is where the Route World
meets Principle~\ref{gc:prin:continuity} and returns something concrete.

Golden Optimation running over the ancestral record and over the present
is one process under one conditioning, so there is no categorical
transition between the environments of natural history and those of the
present --- which is continuity's own claim, arrived at from the
conditioning side. The present is nonetheless distinguished in
\emph{stakes}, by Remark~\ref{hr:rem:monotone}, and the two are
compatible for the reason \S\ref{ac:sub:extrapolation} gives: a
monotone sequence distinguishes every later point over every earlier one
without distinguishing any point categorically. The right image is a
field potential that is high here because the Golden Point is near ahead
--- and the misreading to guard against is that a high potential marks a
step. It does not: a potential rises continuously toward its source, and
what is high at the present is high because of proximity to a transition
that lies ahead, not because a transition has been crossed.

The consequence for periodization is that a boundary drawn in the past
--- an Anthropocene, on any of its proposed placements --- marks a
change in stakes rather than in kind, and continuity is what denies it
the second reading. The argument should be run from the principle and
not from the failure to locate a line, since inability to place a
boundary is a sorites and not a proof. What the Route World adds is a
positive alternative: the transition the periodization is reaching for
lies \emph{ahead}, at the Golden Point, which is the one place the
framework does mark a change in kind --- and which
\S\ref{gr:sub:ceiling} forbids anyone from dating.
\end{remark}

% ---------------------------------------------------------------------

\part{Scores, Meaning, and Conduct}

\noindent Part~IV is the normative and existential engine. The Score
is derived and fenced, Golden ethics is developed with the Departric
result at its center, the commitment-node crisis is diagnosed and
given its criterion, and the third route, inferred meaning, is stated
at its two tiers, placed against its literature, and defended against
the reduction to religion (\S\ref{gs:sec}). The part yields conduct
guidance under the computation ban and the meaning claims at their
ladder of grades.

\section{Golden scores: a decision rule and its limits}\label{gs:sec}

\begin{quote}\small
\S\ref{ac:sec} adopted a grounding and left it uncomputable. This section
computes it, and the computation is available in the model's own terms
rather than by analogy: two options' survival tails carry the same power
of $t$, the power cancels, and what remains is the ratio of their
persistence-harmonic values --- the Doob density that defines the Golden
Biosphere. The decision rule proposed informally in the earlier material
is therefore $\hm(A)/\hm(B)$, weighted by how attainable each option is.

Three things accompany it and are stated rather than deferred. The
identity is exact for the hard-wall functional and a bounded-factor
comparison for the soft-killing one, which is what the inputs' grades
purchase. The map from a decision to a point in the model's state space
does not exist, so the formula is exact and not yet usable, and
Remark~\ref{gr1:rem:noproxy}'s reasoning says why supplying one in haste would
be worse than waiting. And the rule licenses a comparison of conditional
probabilities, never a duty: $\Phi2$ stands, and \S\ref{gs:sub:fence}
restates it at full strength, including the consequence that no actor is
scored by it.
\end{quote}

\subsection{The Initiative}\label{gs:sub:initiative}

\S\ref{ac:sec} ended with an agent's grounding; this section asks what the grounding computes. Between the two sits one definition, and one sentence relating it to the mathematics.

\begin{definition}[The Golden Initiative; \textbf{[$\Phi$]}]\label{gs:def:initiative}
The Golden Initiative is a biosphere's organized effort toward aligning practical outcomes with the Golden Answer.
\end{definition}

The relation to the Tendency is the subsection: \emph{the Initiative is the conscious-agency component to what the Tendency describes mathematically.} The Tendency is an ensemble fact --- futures reweighted, nothing striving (Remark~\ref{ac:rem:noaim}) --- and the Initiative is not its cause but its agent-side face: what a biosphere's practice looks like when organized around the affirmative rather than the negative. The paper has already consumed this object once. The middle step of Proposition~\ref{ac:prop:forcing} quantifies over exactly it --- \emph{a biosphere whose practices are organized around the negative does less to persist} --- so the Initiative is the affirmative case of the practice-organization that step needs, named rather than left implicit. And the scale discipline of Remark~\ref{ac:rem:bridge} carries over unchanged: the Initiative is a biosphere-scale object, and an individual participates in it without being a lever on it.

One expectation attaches, and its kind must be stated before its confidence. In ensemble form it is the Tendency at the level of practice, and needs nothing new: biospheres found at horizon $t$ have, increasingly in $t$, more of their practice so organized --- Principle~\ref{ac:prin:tendency} read through the middle step. The trajectory form --- one biosphere's Initiative deepening over time, in the far tail plausibly the analog of what self-preservation is to an organism now --- is a different kind of claim, and Remark~\ref{ac:rem:notrajectory} locates rather than forbids it: it becomes available only where technoanatomy carries its own dynamics (\S\ref{gr1:subsub:fluct}, Corollary~\ref{gr1:cor:recurrentanatomy}), and those dynamics are \textbf{[O]}. The forecast is filed where \S\ref{ac:sec} says trajectory claims live, and nothing below consumes it.

\subsection{The rule}\label{gs:sub:rule}

Let an agent hold two mutually exclusive options in disagreeing superposition, $A$ with weight $P(A)$ and $B$ with weight $P(B)$, and let $G$ be the state of the biosphere surviving forever. The 2024 material proposed the comparison
\begin{equation}\label{gs:eq:rule}
    P(A)\,P(G \mid A) \;\ge\; P(B)\,P(G \mid B),
\end{equation}
equivalently
\begin{equation}\label{gs:eq:ratio}
    \frac{P(G \mid A)}{P(G \mid B)} \;\ge\; \frac{P(B)}{P(A)}.
\end{equation}
The feature worth keeping is slightly counterintuitive and real: \emph{the difficulty of adopting an option factors into its evaluation}. An option held at weight $0.2$ against an alternative at $0.8$ must carry a compensating factor of four in conditional survival before the comparison favors it --- the rule prices not only where an option leads but how far from the agent's present superposition it sits.

As written, however, the rule is degenerate exactly where the paper lives. In the critical regime $P(G) = 0$ identically, so \eqref{gs:eq:ratio} is $0/0$ and the comparison says nothing. The degeneracy is not a defect of the idea but of the finite form, and the repair is the limit --- which is the next subsection's derivation, not a further stipulation.

\subsection{The limit form and the harmonic function}\label{gs:sub:limit}

Replace $G$ by survival to a horizon, $S_t$, and take the horizon out:
\begin{equation}\label{gs:eq:limit}
    \lim_{t \to \infty} \frac{P(S_t \mid A)}{P(S_t \mid B)} \;\ge\; \frac{P(B)}{P(A)}.
\end{equation}
In the critical regime this limit is computable in the paper's own machinery --- but the computation must be run at the grade its inputs carry, and they carry two. Claim~\ref{sh:claim:tail} gives the survival tail as $\PP(T > t) \asymp C\,t^{-1/4}$, and the paper's convention fixes $\asymp$ as bounded ratio: $0 < \liminf f/g \le \limsup f/g < \infty$. Two quantities each $\asymp$ the same power have a ratio bounded between positive constants; nothing converges. What delivers convergence is the \emph{refined} form, $\PP(\tau > t) \sim c\,\hm\,t^{-1/4}$, which Theorem~\ref{sh:thm:persist} with Proposition~\ref{sh:prop:harmonic} (general start, \cite{GJW99}) supplies for the hard-wall functional $\tau$, and which Theorem~\ref{sh:thm:Q} already assumes. Ansatz~\ref{sh:ansatz} transfers $T$ to $\tau$ only at $\log \PP(T>t) = \log \PP(\tau>t) + O(1)$ --- $\asymp$-grade, which does not upgrade.

\begin{proposition}[The Golden Score; \textbf{[P]} given \textbf{[A]} and \textbf{[C]} inputs]\label{gs:prop:score}
Identify the options $A$, $B$ with starting states (the identification map pending; see below). For the hard-wall functional $\tau$,
\begin{equation}\label{gs:eq:harmonic}
    \lim_{t \to \infty} \frac{P(\tau > t \mid A)}{P(\tau > t \mid B)} \;=\; \frac{\hm(A)}{\hm(B)},
\end{equation}
the Doob transform density of \eqref{sh:eq:Qdensity}, and \eqref{gs:eq:limit} is equivalent to the single computable functional
\begin{equation}\label{gs:eq:functional}
    P(A)\,\hm(A) \;\ge\; P(B)\,\hm(B).
\end{equation}
For the soft-killing functional $T$ the limit exists but is \emph{not} taken in $\hm$. Splitting $\PP_x(T>t)=\E_x[e^{-\Lambda_t}]$ on $\{\tau>t\}$, the surviving term is $\PP_x(\tau>t)\cdot\E^{(t)}_x[e^{-\Lambda_t}]$, and the conditional factor converges to a \emph{state-dependent} constant $\E^x_Q[e^{-\Lambda_\infty}]$; the sub-wall term contributes a renewal series over excursions. Hence
\begin{equation}\label{gs:eq:hT}
\lim_{t\to\infty}\frac{P(T>t\mid A)}{P(T>t\mid B)}=\frac{\hm_T(A)}{\hm_T(B)},
\qquad
\hm_T(x)=\hm(x)\,\E^x_Q\bigl[e^{-\Lambda_\infty}\bigr]+(\text{excursion series}),
\end{equation}
which differs from $\hm(A)/\hm(B)$ in general, since $\E^x_Q[e^{-\Lambda_\infty}]$ increases with the start. Using $\hm$ in place of $\hm_T$ is the bounded-factor approximation, and $\hm_T=c\,\hm$ exactly in the hard-wall limit.
\end{proposition}

\begin{remark}[What the soft-killing statement costs, and what discharges part of it]
\label{gs:rem:hT}
\eqref{gs:eq:hT} is stronger than the bounded-factor language it replaces --- the rule is
\emph{exact} for $T$ as well, in a different object --- and it is owed three conditions. The
first is that the $Q$-expected residual hazard be finite,
$\sup_t\int_{t_0}^t\E^{(t)}[h^{\mathrm{tot}}]\,ds\to0$; the other two govern the excursion
series (finite crossing mass, and convergence of the renewal sum), and both remain open
\textbf{[O]}.

The first is not open for this model. Under $Q$ the credit grows as $x_s\asymp s^{3/2}$, so a
background decaying only polynomially, $\varsigma_b(x)\sim x^{-p}$, would give residual hazard
$\int s^{-3p/2}ds$ and require $p>2/3$. But \eqref{tc:eq:hazard}'s muffling factor is logistic
and therefore decays \emph{exponentially}: $\varsigma_b(x_s)\sim e^{-s^{3/2}}$, and the trough
contributes $\Gamma e^{-\rho s^{3/2}}$. Both integrals converge absolutely whenever $\rho>0$ ---
which is the regime the Score is stated for anyway, by
Theorem~\ref{gr1:thm:twocond}. So the condition is met with room to spare, and the polynomial
threshold is recorded because it is a real constraint on any \emph{other} choice of soft wall.

The consequence for the identification map is worth naming: what a decision must be mapped to
is a starting state for $\hm_T$, and $\hm_T$ carries the hazard the biosphere will actually
accrue rather than only the geometry of the wall.
\end{remark}

The proof is the cancellation: both tails carry the same $t^{-1/4}$, under $\sim$ the prefactors are all that survive, and the prefactors are harmonic values. So the decision rule stated informally in 2024 is computable in the model's own terms, and the quantity it computes is the same object that defines the Golden Biosphere --- a derivation, not an analogy. The qualification in the proposition is not a retreat: a comparison exact for $\tau$ and bounded-factor for $T$ is precisely what the inputs' grades purchase, and stating it so is what the grading discipline is for.

\emph{Aurutance.} The natural linearization is the logarithm. Define the Aurutance of $A$ over $B$, under the same hypotheses, as
\begin{equation}\label{gs:eq:aurutance}
    \operatorname{Aur}(A, B) \;=\; \lim_{t \to \infty} \ln \frac{P(S_t \mid A)\,P(A)}{P(S_t \mid B)\,P(B)}
    \;=\; \underbrace{\ln \frac{\hm(A)}{\hm(B)}}_{\text{log-Doob density}} \;+\; \underbrace{\ln \frac{P(A)}{P(B)}}_{\text{log-prior}},
\end{equation}
a portmanteau of \emph{aurum} and \emph{importance}. The sign convention falls out of \eqref{gs:eq:functional}: positive Aurutance is the comparison favoring $A$, in the exact sense of the limit and in no other sense --- what further sense an agent attaches to it is \S\ref{gs:sub:fence}'s subject. One reading is fenced now rather than later: Aurutance is the logarithm of a probability ratio, and the additivity of \eqref{gs:eq:aurutance} across independent comparisons is the additivity of logarithms --- it licenses no accounting metaphor beyond that, and in particular no stock of scores to hold, spend, or owe.

\emph{Two gaps, flagged rather than papered over.} \emph{First, the identification map is missing.} $\hm$ is a function of the model's state $(\mu, D)$; $A$ and $B$ are Aethic states --- additions to an agent's Aethus. Nothing in the paper supplies the map from a decision to a point in $(\mu, D)$, and without it \eqref{gs:eq:functional} is exact and unusable. The gap is \textbf{[O]}, and it is the same shape as the reservation already recorded at Remark~\ref{gr1:rem:noproxy} about measuring $\rho$: a premature proxy for $\hm$ would be a scored index, inviting exactly the Goodhart failure that $\rho < 0$ names. The formula waits for the map; it does not license improvising one.
Remark~\ref{sch:rem:initialdata} names the one place a candidate might be
found: the schedule channel delivers a distribution over completion
states, and completion states are initial data. \emph{Second, the computation is regime-specific.} The cancellation is a fact about the critical cell alone --- the next subsection tables what happens elsewhere.

One qualitative consequence is available now, at the class level and no finer. Expenditure that raises technostrength while leaving technoanatomy untouched moves neither $\hm$ nor the comparison --- it generates no Golden Score --- and by Remark~\ref{gr1:rem:againstsing} may generate a negative one. That is a claim about a class of activity, it follows from results already proved, and it is as far as the rule reaches: Theorem~\ref{gr:thm:ceiling} forbids certifying any system's position from finite observation, so no actor, named or unnamed, is scored by it --- a bar \S\ref{gs:sub:fence} states in full.

\subsection{Regime-dependence}\label{gs:sub:regimes}

What the rule computes depends on the cell, in the pattern of \S\ref{gr:subsub:regimes}:

\begin{center}\small
\begin{tabular}{@{}lll@{}}
\toprule
\emph{Regime} & \emph{Survival tail} & \emph{The rule computes} \\
\midrule
Viable & $\PP(T = \infty) > 0$ & an ordinary conditional ratio; no limit needed \\
Critical & $\sim c\,\hm\,t^{-1/4}$ (refined form) & the harmonic ratio, Proposition~\ref{gs:prop:score} \\
Trapped & $\le e^{-h_\ast t}$, frozen $B$ (Theorem~\ref{gr1:thm:trapped}(ii)) & a comparison of floors $h_\ast$, via \eqref{gr1:eq:xstar} \\
\bottomrule
\end{tabular}
\end{center}

The trapped cell deserves finishing rather than dismissal. Across different anatomies the ratio diverges or vanishes, and what the rule collapses to is a comparison of the floors themselves: $h_\ast$, which \eqref{gr1:eq:xstar} gives in closed form as a geometric interpolation between the external and internal scales, weighted by anatomy. The trapped verdict is therefore \emph{lower the floor} --- which is to say \emph{move $\rho$} --- exactly the intervention Remark~\ref{hr:rem:concession} names as the one that changes the indicator rather than the multiplicand. Even where the harmonic structure is gone, the rule points at anatomy.

The asymmetry deserves its sentence: the rule is \emph{most} informative exactly where survival is \emph{not} achievable --- the critical cell, where $P(G) = 0$ and the harmonic structure is richest --- and least informative where survival is either assured in measure or foreclosed. \S\ref{gr:subsub:regimes} records the same inversion about the mutual-singularity claim; this is the second instance of that structural fact, and noticing the repetition costs one sentence and buys the pattern.

% [Placement pass, Aug 21: the rule's decision-theoretic neighbors.]
\begin{remark}[The rule's neighbors: maxipok, ergodicity, ruin]
\label{gs:rem:ruleneighbors}
Three decision-flavored literatures sit near the Score, reached from
premises this work does not share, and the pattern across all three is
the same: convergent conclusion, divergent grounds, with the divergence
carrying the content.

\emph{Maxipok.} Bostrom's rule --- maximize the probability of an
``OK outcome'' --- is the Score's nearest published neighbor
\cite{Bos03}, arrived at from astronomical-waste aggregation, exactly
the expected-value machinery \S\ref{ac:sub:placement} declines. The
deltas: the Score's ranking is by survival-harmonic weight rather than
bare probability, so options are priced by what they make attainable
and not only by what they avoid; the ``OK outcome'' is undefined where
$\nu F$ is a fixed point; and where maxipok invites computation, the
Score forbids it absent the identification map. Same direction of
travel, different vehicle, and the vehicle is the claim.

\emph{Ergodicity economics.} Peters's program replaces
ensemble-expectation with time-average growth, on the ground that a
multiplicative process's expectation is dominated by branches no
history occupies \cite{Peters19}. The kinship is real and structural
--- Aurutance is a log-ratio, additive across independent comparisons,
and the refusal to multiply small probabilities by large payoffs is
common property --- and the delta is the target: there the object
whose growth is averaged is wealth, here the object is the
survival-harmonic weight, and the family-versus-measure distinction
(Remark~\ref{gr:rem:familyvsmeasure}) is what the time-average
critique looks like when the multiplicative quantity is existence.

\emph{Ruin problems.} The precautionary literature's structural core
--- ruin is an absorbing barrier, so cost--benefit reasoning that
averages across it is invalid \cite{TRDNB14, Tal12} --- is the
Base-level shadow of this section's machinery. The delta is the
closure: their ruin is an event, $\mu G$ is the inevitability closure
of the event, strictly larger \eqref{gc:eq:departurefix}, and the
separations of Remark~\ref{ac:rem:losingcondition} live entirely in
the difference. A precautionary rule keyed to observable ruin guards
the wrong boundary: it fires late where departure is silent, and it
would forbid, in the critical regime, exactly the exposure that
constitutes staying alive.
\end{remark}

\subsection{What it does not license}\label{gs:sub:fence}

The fence, at full strength, because this is where a reader most wants the rule to be more than it is. $\Phi2$ stands: the rule licenses a comparison of conditional probabilities, and a comparison is not a duty. Calling one option \emph{better} imports $\Phi6$ \emph{and} the Golden Answer, both of which are the agent's adopted grounding rather than the apparatus's output --- the split of \S\ref{ac:sub:countermeasure}, consumed here exactly as stated there, and the reason the enumeration's fourth step said \emph{comparisons follow}, not \emph{duties do}. An agent who has adopted neither reads \eqref{gs:eq:functional} as arithmetic about an ensemble, and the paper gives them no argument that they are wrong to.

Two further bars. Theorem~\ref{gr:thm:ceiling} forbids certifying any system's position on the interpolation from finite observation, so the rule cannot be used to score an actor --- not a person, not an institution, not a biosphere, one's own included; the class-level consequence of \S\ref{gs:sub:limit} is the most the machinery yields. And the missing identification map means no number anyone produces today \emph{is} an Aurutance: a computed score requires a map the framework does not yet contain, and anything offered in its place is a proxy of exactly the kind Remark~\ref{gr1:rem:noproxy} warns against. Remark~\ref{hr:rem:stakeslimit} fences the neighboring stakes argument at this same scope and in this same tone; the two fences are one commitment stated twice.

\subsection{Golden ethics}\label{gs:sub:ethics}

\S\ref{gs:sub:fence} said what the rule does not license. This subsection
says what it licenses \emph{for an agent who has adopted the grounding of
\S\ref{ac:sec}}, and that qualification is the whole of its content. No
premise is added here: everything below is downstream of $\Phi6$ and the
affirmative answer to the Golden Question, and an agent declining either
is committed to nothing in it.

\subsubsection{The commitment, and why it is not an issued obligation}
\label{gs:ssub:commitment}

The tempting formulation is a duty --- \emph{one ought to prefer the
option with the higher Golden Score}. Stated so it is a second premise,
and $\Phi2$ forbids it: \S\ref{ac:sub:countermeasure} is explicit that
$\Phi6$ is the single doorway through which obligation enters this work,
and a free-standing Golden Obligation would be a second doorway.

The formulation that costs nothing reads the commitment off the stakes
instead of issuing it. \S\ref{hr:subsub:stakes} establishes that
forfeiting is not the same as never having had, and
Remark~\ref{hr:rem:monotone} that the stakes so understood rise
monotonically along a lineage. Those are results, not instructions. What
$\Phi6$ adds is an agent standing in a grounding under which those stakes
are \emph{theirs}. The commitment is then not an imperative anything
issued but the shape those stakes take once occupied --- and it is
revocable in one specific respect and not in others. Three parts, easily
run together and worth separating.

\emph{That there is an obligation} is not conditional at all. Condition~1
of \S\ref{ac:sub:problem} places the agent under forced action, and
Proposition~\ref{ac:prop:nonempty} establishes that they hold a grounding
whether or not anything is said. An agent already acts, and already on
some footing --- so the normative is in place before $\Phi6$ speaks,
which is exactly what \S\ref{ac:sub:countermeasure} means in denying that
any \emph{ought} is derived here from an \emph{is}. None of this issues
from the apparatus, so $\Phi2$ is untouched by it.

\emph{What the obligation is about} is the conditional part, and the only
one. $\Phi6$ selects which frozen question a grounding turns on, and an
agent declining it does not thereby escape obligation --- they hold a
different one, on a footing they have not chosen and cannot inspect.

\emph{And how much rides on it} the stakes supply.
\S\ref{hr:subsub:stakes} and Remark~\ref{hr:rem:monotone} give the
magnitude, and Principle~\ref{gc:prin:continuity} extends it backward:
the stakes are continuous with those our ancestors faced without agency,
so nothing about them begins with us. Stakes supply weight and $\Phi6$
supplies direction, and neither does the other's work. \emph{Golden virtue} names conduct answering to that shape and
claims nothing about virtue in any prior sense.

\subsubsection{Two moralities, and the case that separates them}
\label{gs:ssub:twomoralities}

Write $\mathcal G$ for the Golden Point, $\mathcal N_0$ for the Cosmic
Nexus, and let $\varphi$ range over stated attributes --- dispositions,
policies, practices. Two quantities are available and they come apart.

\begin{definition}[Automorality]\label{gs:def:automorality}
For an Aethus $A$ carrying $\varphi$, with $A'_\varphi$ its complement
under $\varphi$,
\begin{equation}\label{gs:eq:am}
\mathrm{AM}(A,\varphi)=
\ln\frac{\PP(\mathcal G\mid A,\mathcal N_0)}
{\PP(\mathcal G\mid A'_\varphi,\mathcal N_0)} .
\end{equation}
The gain from the attribute's having been carried --- read modally,
\emph{if this agent does this, does the Golden Point become likelier?}
It is act-indexed and indifferent to intent.
\end{definition}

\begin{definition}[Relative morality]\label{gs:def:relmorality}
For a reference Aethus $B$, with $B_{\varphi,A}$ replacing $B$'s
attribute under $\varphi$ by $A$'s,
\begin{equation}\label{gs:eq:rm}
\mathrm{RM}_B(A,\varphi)=
\ln\frac{\PP(\mathcal G\mid B_{\varphi,A},\mathcal N_0)}
{\PP(\mathcal G\mid B,\mathcal N_0)} .
\end{equation}
The gain from the disposition's being \emph{adopted} --- read modally,
\emph{if an agent in general does this, does the Golden Point become
likelier?} It is maxim-indexed rather than act-indexed, and asks what
transplanting the disposition would do rather than what carrying it did.
\end{definition}

They agree in the ordinary case and separate in the instructive one. Take
an attribute exercised with intent to harm which, by the stepping-stone
structure of \S\ref{gc:sec}, happens to open a route that mattered. Its
automorality is positive: it helped. Its relative morality is negative:
installed in a reference agent it would be expected to do the destructive
thing, and the favorable outcome was the accident rather than the
disposition. Table~\ref{gs:tab:moralities} places this at \emph{Biotic evil} --- evil
by disposition, Biotic by outcome --- one cell along from \emph{Departric
evil}, which is the same disposition where the accident did not occur,
and diagonally opposite \emph{Departric walk}, where the disposition was
the good one and the accident went the other way. The pair therefore holds apart what an act \emph{did} from
what a disposition \emph{is}, rather than reducing either to the other.

\begin{table}[htbp]
\centering\small
\begin{tabular}{@{}lccc@{}}
\toprule
& \textbf{Departric} & \textbf{Neffective} & \textbf{Biotic}\\
& \small(autoimmoral) & \small(autoneutral) & \small(automoral)\\
\midrule
\textbf{relatively moral}   & Departric walk    & Neffective walk  & \emph{Golden good}\\
\textbf{relatively neutral} & Departric power   & Neffective       & Biotic power\\
\textbf{relatively immoral} & \emph{Departric evil} & Neffective evil & \emph{Biotic evil}\\
\bottomrule
\end{tabular}
\caption{Dispositions classified by the two moralities. Each name is
built from its coordinates, and the prefix is always a term of art, so no
cell can be mistaken for ordinary usage. The \emph{prefix} reports the
outcome: \emph{Departric} where automorality is negative,
\emph{Neffective} where it is zero, \emph{Biotic} where it is positive
--- so \emph{Bio-} occurs in the automoral column and nowhere else. The
\emph{suffix} reports the disposition: \emph{walk}, \emph{power},
\emph{evil} as relative morality runs $+1$, $0$, $-1$ --- so
\emph{evil} occurs in the bottom row and nowhere else. \emph{Walk} is
chosen for the top row because English supplies no neutral term for a
disposition worth adopting: \emph{merit}, \emph{virtue} and
\emph{goodness} all import an approving verdict on the outcome, which two
of that row's three cells do not have. A walk is a manner of proceeding
rather than an act or a result, so the prefix is left free to report how
the proceeding turned out. \emph{Good} is not
a suffix in the scheme at all, since it asserts a beneficial outcome that
only one column has; it is reserved for the corner where merit meets a
Biotic outcome, which takes \emph{Golden} in place of \emph{Biotic} to
mark the double alignment. The centre cell is bare because both
coordinates vanish there and the ratio of
Remark~\ref{gs:rem:power} is undefined rather than unbounded. So the
bottom row's two evils are one disposition-class distinguished by what it
happened to produce --- the distinction the pair exists to draw and the
one a single coordinate cannot --- while \emph{Departric walk} is its
mirror, a disposition that would have helped meeting an outcome that
harmed. \emph{Neffective evil} imputes no malice, relative immorality
covering negligence and folly and requiring no ill intent. The names are
stipulative in the way \emph{force} is stipulative in mechanics, and are
not offered as analyses of ordinary usage.}\label{gs:tab:moralities}
\end{table}

\noindent Cell by cell. The first clause reports what the disposition
does if adopted generally, the second what carrying it did here; the two
are indexed differently, which is why they may disagree without
contradiction.

\begin{itemize}\setlength\itemsep{1.5pt}
\item \emph{Departric walk.} Generally beneficial; harmful here. A
disposition one would want adopted, meeting a circumstance in which
carrying it moved things toward Departure --- misfortune rather than
fault.

\item \emph{Neffective walk.} Generally beneficial; inert here. A
disposition one would want adopted, which on this occasion found no
purchase, with no verdict on its bearer implied either way.

\item \emph{Golden good.} Generally beneficial; beneficial here. The one
cell with both coordinates positive, which is what earns \emph{Golden}
and what earns \emph{good}.

\item \emph{Departric power.} Generally inert; harmful here. Nothing in
the disposition to transplant, and an outcome owed entirely to where the
attribute happened to fall.

\item \emph{Neffective.} Generally inert; inert here. Neither a
disposition nor an outcome, and nothing for position to stand between ---
the one cell about which the pair has nothing to report.

\item \emph{Biotic power.} Generally inert; beneficial here. The
attribute contributes nothing transplanted, yet in the circumstance it
occupied it raised $\PP(\mathcal G)$: what would be nothing anywhere was
decisive here.

\item \emph{Departric evil.} Generally harmful; harmful here.
Disposition and outcome agree unfavorably --- the cell ordinary usage
means by \emph{evil}, and the only one it fits without loss.

\item \emph{Neffective evil.} Generally harmful; inert here. A
disposition that failed to find its occasion; since relative immorality
covers negligence and folly, no ill intent is implied.

\item \emph{Biotic evil.} Generally harmful; beneficial here. The
disposition is the destructive one and the favorable outcome was the
accident --- the case the two coordinates exist to separate.
\end{itemize}

\noindent The diagonal is agreement and the off-diagonal divergence, and
the divergence is where a single moral verdict has to choose what to
throw away. Ordinary usage collapses \emph{Departric evil} together with
\emph{Biotic evil}, since both are dispositions one would not want
adopted; and it has no settled name at all for \emph{Departric walk},
which is why \S\ref{gs:ssub:twomoralities} needed two coordinates rather
than one.

\begin{remark}[Position against disposition]\label{gs:rem:power}
The ratio $|\mathrm{AM}/\mathrm{RM}|$ measures something neither term
does alone: how far an outcome was owed to the \emph{position} an
attribute was exercised from rather than to the attribute. Large ratios
mark dispositions whose effect exceeds what the same disposition would
achieve elsewhere, which is what is ordinarily meant by leverage. This is
the position-versus-purchase distinction of
Remark~\ref{ac:rem:bridge} appearing on the normative side, and it
inherits that remark's caution: occupying a position of leverage is not a
causal claim about its occupant.

Three notes on the vocabulary this fixes. \emph{Walk} carries its older
moral sense --- a manner of proceeding, as in a walk of life --- and not
the stochastic one the model's integrated random walk uses; the modifier
disambiguates at every occurrence, and the two senses meet in no single
passage. \emph{Neffective} names the
automorality-zero column --- a disposition whose carrying did not move
$\PP(\mathcal G)$ at all, which is the Golden analogue of an attribute of
zero linkage in the companion's sense \cite{Nexus}. And \emph{power} is
the suffix for the relatively-neutral row, which repays being read
literally rather than as a row of weak effects. $\mathrm{RM}_B(A,\varphi)
=0$ says the attribute transplanted into the reference circumstance does
\emph{nothing at all}; $\mathrm{AM}(A,\varphi)\neq0$ says that in the
circumstance it actually occupied it did something. The row is therefore
an \emph{exchange}: multiply the attribute into a standard Aethus and it
is inert, multiply it into this one and it is consequential, and the
whole of the outcome is owed to which Aethus received it. That the ratio
above is unbounded there is not a pathology but the correct reading ---
all position, no disposition.

Which makes the row the disposition-side analogue of a structure the
companion already names. A Miracle is an event of negligible generic
weight whose occurrence on one particular line was decisive
\cite{Nexus}; a power cell is an attribute of null generic effect whose
carrying in one particular circumstance was. The parallel is structural
and not an identity --- those are events and these are dispositions, and
the companion's linkage is indexed to observer existence where these are
indexed to $\mathcal G$ --- but it is the same shape, and it is why
\emph{power} is the right word: the capacity to matter here that would
matter nowhere.

The centre cell is the one place this fails, which is why
Table~\ref{gs:tab:moralities} leaves it unsuffixed. Where automorality
vanishes as well, the ratio is $0/0$ and undefined rather than unbounded:
there is neither a disposition nor an outcome for position to stand
between, so nothing is being exchanged.
\end{remark}

\begin{remark}[The act and the maxim, and what the pair is doing to a
familiar dispute]\label{gs:rem:maxim}
The modal readings above have a shape worth naming. Automorality asks
after an act, in a position, by an agent; relative morality asks what
follows if the disposition is adopted generally. That is the distinction
between an act and a maxim, and it is the distinction the categorical
imperative's universalization test needs in order to be stated at all.

The resemblance should not be inflated into an encoding, and three
differences say why. Kant's test is for \emph{contradiction} --- a maxim
that cannot be conceived or cannot be willed as universal --- whereas
$\mathrm{RM}$ tests for a change in probability, which is
consequentialist content wearing a universalizing form. Kant
universalizes over all rational agents, whereas $\mathrm{RM}$ substitutes
into a stated reference $B$ and carries it essentially, so what is
recovered is a \emph{localized} universalization with the full one as a
limiting case, obtained by ranging $B$ over a population in the manner
\S\ref{gs:ssub:limits} describes. And Kant's test returns a verdict where
$\mathrm{RM}$ returns a magnitude.

What the pair does supply is a measurement of a fault line rather than a
position on it. Automorality is act-consequentialist; relative morality
is the universalizing form with consequentialist content, which is to say
rule-consequentialist. The cells where the two diverge ---
\emph{Biotic evil} and its mirror in
Table~\ref{gs:tab:moralities} --- are exactly the cases over which act-
and rule-readings of consequence have always disagreed, and the two
coordinates \emph{quantify} the disagreement instead of adjudicating it.
A framework that had to choose between them would lose the information
the divergence carries.

One differentia remains, and it is the one \S\ref{gs:ssub:notlongtermism2}
turns on. Every consequentialism needs a quantity to be a consequence
\emph{of}, and the usual choice is welfare aggregated over occupants.
Here it is $\PP(\mathcal G)$ --- probability of reaching the asymptote,
not a sum over anyone. So this is a normative ethics fitted to
infinite-horizon biosphere survivorship rather than to utility, which is
why it ranks as it does and why the agreement with aggregative views is
occasional rather than principled.
\end{remark}

\subsubsection{What this cannot do}\label{gs:ssub:limits}

Three limits, the first severe enough that the construction should be
read as an exercise rather than a proposal.

\emph{It requires a vantage this work proves unavailable.} Both
\eqref{gs:eq:am} and \eqref{gs:eq:rm} consume
$\PP(\mathcal G\mid\cdot)$ for arbitrary counterfactual Aethae.
Theorem~\ref{gr:thm:ceiling} forbids certifying any system's position
from finite observation, and no weaker vantage supplies these either. So
Golden ethics is not awaiting data --- it is a well-defined construction
on a counterfactual the framework itself rules out. That is a stronger
disclaimer than impracticality, and it is the correct one.

\emph{It does not score persons.} Both quantities are defined over
\emph{attributes}, and \S\ref{gs:sub:fence}'s prohibition stands
unamended --- the rule cannot be used to score an actor, ``not a person,
not an institution, not a biosphere, one's own included'' --- so the
apparatus remains articulate about dispositions and classes of conduct
and silent about their bearers. A table of dispositions
is not a table of people, and reading
Table~\ref{gs:tab:moralities} as the latter is the error that remark
exists to prevent.

\emph{And no score is absolute.} $\mathrm{RM}$ carries its reference $B$
essentially rather than incidentally; there is no view from nowhere
against which a disposition scores. The nearest substitute is an extremal
score across a stated population, which is a different object and should
be labeled as one.

\subsubsection{Why this is not longtermism, substantively}
\label{gs:ssub:notlongtermism2}

Remark~\ref{gr:rem:notlongtermism} answers the objection structurally and
concedes that in the viable regime it must be met by $\Phi2$ alone. This
closes that gap, because even granting a normative reading the two
orderings differ.

Golden reasoning grades by probability of reaching the boundary, not by
aggregate life. A thousand persisting far outrank a million persisting
less far, since what is conditioned on is survival to the asymptote and
not a sum over occupants; and a population growing while decaying
exponentially in prospect ranks poorly for the same reason, being no
sustainability at all.

The standard test case separates them cleanly. Asked whether to found a
settlement whose occupants will live and die young, an aggregative
longtermism answers yes --- existence is preferred to non-existence and
the sum increases. Golden reasoning has no such argument available: the
settlement's contribution to $\PP(\mathcal G)$ ranges from uncorrelated
to adverse, and at no point does the framework return a positive verdict
on that ground. Where the two agree, as over the fate of a biosphere,
the agreement is a consequence and not a shared principle.

\subsubsection{Past the Departure Point}\label{gs:ssub:departric}

A stronger version of the following is tempting and unsupportable, so the
weaker one is stated exactly.

The tempting version says that beyond a biosphere's Departure Point all
virtue and meaning cease to exist. Nothing here supports that. What is
supportable is that past $\mu G$ the viability event has probability
zero, so $Q$ is not defined as a conditional measure and every quantity
indexed to it --- Golden Scores, both moralities above, the ordering of
\eqref{gs:eq:functional} --- is undefined rather than small. That is a
claim about an index, not about the world.

The claim has a published cousin that locates it from the outside.
Scheffler's collective-afterlife argument holds that the mattering of
present activity depends, more than we notice, on humanity's continuing
after us --- his doomsday and infertility scenarios drain present
projects of point \cite{Scheffler13}. That is this subsubsection's
content met in the psychological register: his scenarios are
conditionally Departric worlds, and what he argues we would \emph{feel}
--- the deflation of mattering --- is what the adopted ordering
\emph{computes}, indices undefined past $\mu G$. The deltas are two
and they cut in opposite directions. His claim is empirical psychology
and holds for whoever's motivations it describes; this one is
$\Phi6$-conditional and holds for whoever adopts the currency ---
narrower in reach, harder inside its reach. And his argument concerns
valuers who precede the loss, while the striking here concerns the
comparison itself; the two agree exactly where
Remark~\ref{gc:rem:complacency}'s mixed-currency reader disagrees with
both.

It has a consequence for forced choices, and the shape of the argument
needs exact statement, because it resembles a wager and is not one.

Suppose an agent under $\Phi6$ faces two courses: one carrying a small
probability of reaching the Golden Point, the other none. The obvious
objection is that the second may be far safer in the ordinary sense ---
that the agent's own survival is likelier under it. The reply is that
this survival lies entirely in the region past $\mu G$, across which the
adopted ordering is indifferent by the paragraph above: every
continuation there is equivalent under it. So the ordering does not fail
to rank the safe course. It ranks it at \emph{zero}, and the survival
that recommended it falls inside an equivalence class the ordering reads
as one point.

What remains is a comparison of two values of a single well-defined
quantity, one strictly greater than the other --- a percent against
none. Nothing unbounded enters, and no small probability is multiplied
by a large payoff, which is the structure this is most often mistaken
for. It differs in a second way as well: a wager stakes something on a
fact already settled, whereas here the probability is of the agent's own
course succeeding, and which probability the agent faces is exactly what
the choice determines. The quantity is contingent on the action rather
than a credence about the world.

The conditionality is unchanged and does all the limiting work. An agent
who has not adopted $\Phi6$ holds no such ordering, faces no such
comparison, and is owed no such conclusion --- and nothing above
converts the argument into one for them.

\begin{remark}[What this is a formalization of]\label{gs:rem:trope}
The structure above is not novel as an intuition, and saying what it
formalizes is the quickest test of whether the formalization is any
good. It is the familiar narrative case in which an agent accepts a
course they are unlikely to survive because the alternative forecloses
something they take to matter at large scale --- ``for the good of
humanity,'' in the usual phrasing.

Ordinary decision theory reads such a choice as noble and irrational:
near-certain survival traded for near-certain death. What the paragraphs
above supply is the reading on which it is neither. The safe course is
not being outweighed; it is being scored at zero, because the survival
recommending it lies wholly inside a class the adopted ordering reads as
one point. The agent is not overriding their interests but applying an
ordering on which the safe course has no value to override with.

The formalization earns its keep by \emph{discriminating}, since the
trope has a familiar failure mode. The argument requires the risky course
to carry nonzero probability of the thing valued. Where it does not, the
comparison is zero against zero, the ordering is silent, and no gesture
is licensed --- which is Remark~\ref{hr:rem:concession}'s structure
exactly: faced with options of equal viability zero, the response is not
to select among them but to find something that is not zero. Heroism
that accomplishes nothing gets no support here, and that it gets none is
a point in the construction's favor.

One stipulation should be visible rather than smuggled. The scenario
assumes an agent with genuine leverage --- no one better placed, and the
act decisive. Remark~\ref{ac:rem:bridge} denies that in general: an
individual is not a lever on their biosphere's hazard, and holds a
position rather than a purchase. The argument concerns the shape of the
comparison where such leverage obtains, and is not a claim that agents
ordinarily face choices of this kind or should act as though they did.
\end{remark}

% ---------------------------------------------------------------------
% ---------------------------------------------------------------------
% [Instates the author's node dictation (Aug 22): the seminar
%  phenomenology (memoir grade, two undergraduate occasions), the
%  commitment-primitive diagnosis, the testimony collapse, the
%  Golden node-criterion, iterability vs rule-axiomatizations, and
%  the neighbors (MacIntyre, Anscombe, peer disagreement).  A slot is
%  reserved per the author: he may supply the observation freshly when
%  next witnessed; the exhibit upgrades from memoir grade on that day.]
\subsection{Judging the nodes: commitment primitives and the
inference-side crisis}\label{gs:sub:nodes}

The rule above evaluates \emph{options}. This subsection concerns
something the evaluation of options presupposes and moral philosophy
mostly leaves in the dark: the undefended placements from which
practical arguments are launched. It states a phenomenon, a diagnosis,
and what the framework uniquely supplies --- and it ends with the
fences that keep the supply honest.

\emph{The phenomenon, at its provenance grade.} Put a room of capable
people to a live dilemma and ask for conduct. What returns is an array
of stances that mutually render one another absurd --- each participant
certain the others are wrong, and each systematically unable to convey
\emph{why} on demand, however far the case is elaborated; the deeper
the push, the more irresolvable the clash. The author witnessed this
twice as an undergraduate --- once in a philosophy-of-technology
seminar, once in a philosophy-of-science seminar, both times over the
ethics of de-extinction --- and reports the structural content
exactly: worldviews obligating vastly different conduct, not from
anything intrinsic like religion or culture, but from a shallow mutual
throwing-out of reasoned stances on which the whole room then stuck.
Neither occasion was written down while fresh, so the exhibit is
carried at recollection grade per the methodology companion's own
practice \cite{Ben26Method}, with the slot held open for a
contemporaneous record when the dynamic is next observed. Two
distinctions keep the phenomenon exact. It is not the
wisdom-of-crowds setting: aggregation helps where many voices estimate
one shared quantity with independent errors, and here there is no
shared quantity --- the divergence is not noise around a target but
different targets, so no happy medium exists to converge on. And it is
not the peer-disagreement problem of the epistemology literature
\cite{Chr07}, which asks how peers sharing evidence and methods should
update on one another; the seminar deadlock fails the setup's
antecedent, since what differs is neither evidence nor method.

\emph{The diagnosis.} What differs is the \emph{nodes}. Attention in
argument goes to the propagation layer --- the reasoning, and the
equivalence classes of outcomes it defends --- while beneath it sit
commitment primitives: ``it is paramount that we preserve this,''
placed undefended at some spot in the normative landscape, from which
everything else propagates. Nobody in the room knows what anything
really \emph{is}, and it is precisely that non-understanding which
frees each participant to place their primitive anywhere; the
propagations then conflict because the placements do. And the crisis
is one level deeper than the clash, because it attacks testimony
itself: knowing that a speaker's node is sourced in their biographical
and neurological contingency --- knowing they do not know the context
and are guessing at it --- voids the evidential value of everything
downstream, \emph{however valid the propagation}, since there are no
good guesses where guessing has no target, the source is not excusable
from the conclusions, and the source is known to be no good. Call this
the \emph{inference-side crisis}: not that people disagree, but that
each participant is right to discount every other's testimony at the
root, and knows the others are right to discount theirs. Its
first-person phenomenology has a canonical name in the literature ---
the anguish of the value-chooser \cite{Sartre46} --- engaged with its
tradition at \S\ref{gs:sub:meaningplacement}.

\emph{The nearest neighbors, and the delta that matters.} The
diagnosis has a distinguished ancestry and the framework should say so
before saying what it adds. MacIntyre opened \emph{After Virtue} on
exactly this phenomenology --- interminable modern moral disagreement,
rival premises incommensurable, premise-choice arbitrary --- and
Anscombe's diagnosis of a deontic vocabulary surviving the loss of its
root is the same finding one register up \cite{MacIntyre81,
Anscombe58}. The delta is the remedy, and it is total. MacIntyre,
holding a universe-sourced telos to be unavailable to
post-Enlightenment thought, retreats to teleology internal to
traditions and practices; the present framework's entire claim is that
the unavailable object exists: a root sourced in the actual
conditioning structure of the universe --- the Golden Point, with
\S\ref{to:sub:goldenaethus} supplying its constitutive standing under
Strong --- rather than in any tradition, any practice, or any
reasoner's contingency. Where the diagnosis is shared nearly verbatim,
the framework answers the crisis with the object its sharpest
diagnostician declared missing. One provenance fact belongs on the
record here, because the methodology companion makes it evidential
\cite{Ben26Method}: the diagnosis above was arrived at independently
--- from the seminar observations, dated in the collaboration record
--- and located in the literature only afterward, at placement. By the
companion's own standard, independent convergence on a diagnosis is
evidence \emph{for} the diagnosis while moving none of the
contribution's locus, which sits at the remedy; the author has
confirmed the hinge accordingly, and the paper treats it as one: the
crisis this subsection names is the crisis the framework's categorical
and ontological completions (\S\ref{to:sub:goldenaethus}) exist to
answer.

\emph{What the framework supplies: a criterion for the nodes
themselves.} Judge commitment primitives the way
Remark~\ref{gr1:rem:goldendef} judges categories: by their tracing to
the root. A node is non-arbitrary to the degree its content is
individuated by contribution to the Golden Point --- the task-object
machinery of \S\ref{to:sec} run on normative primitives, with the
contribution referent of Remark~\ref{to:rem:goldenroot} pricing what
the commitment's object actually carries. The paradigm case is
already on the books: the frozen answer to the Golden Question is a
commitment primitive whose holding is survivorship-forced rather than
biographically sourced (Proposition~\ref{ac:prop:forcing}) --- the
existence proof that a universe-sourced node is possible at all. And
the de-extinction dilemma that occasioned the observation is, by no
accident, a case the criterion can actually pose: ``what would we be
restoring, and what would its restoration contribute'' is a
task-object identification and a contribution question stated in the
framework's own primitives --- recorded here as a worked-case
candidate \textbf{[O]}, not worked.

\emph{Why it iterates where rule-axiomatizations dissipate.} The
author's second advertised property deserves its structural ground. A
rule-form axiomatization --- Kantian universalization is the standing
example, and Remark~\ref{gs:rem:maxim} locates this framework's
relation to it --- is a finite formula, and finite formulas jam on
complexity: hard cases multiply conditions until the rule
under-determines or self-conflicts. The Golden root is not a formula
but a \emph{conditioning}, and a conditioning is defined wherever its
measure is: it refines under any partition of cases, composes across
scales, and grades rather than jams (the $\varepsilon$-apparatus, the
per-case individuation of task-objects, the Score's comparisons). That
is the exact sense in which the Golden answer ``operatively fits the
shape of the situation rather than dissipating with complexity'': the
criterion's domain is the whole space of situations because it is
carried by a measure on that space, not by a clause-list over it.

\begin{remark}[The deleted middle term, re-derived]\label{gs:rem:middleterm}
The engagement with the interminability diagnosis can be made exact,
because the diagnosis has a precise anatomy \cite{MacIntyre81}:
classical ethics was a three-part scheme --- the agent as they happen
to be, the agent as they would be were their telos realized, and
ethics as the discipline of the transition --- and modernity deleted
the middle term when its warrants (metaphysical biology, divine law)
died, leaving the outer fragments with nothing to connect them. On
that anatomy the Enlightenment justification project failed
necessarily: it inherited the fragments' content while rejecting
every framework that had made content intelligible. The framework's
reply maps term by term, and the mapping belongs on the page rather
than in the implication. The agent as they happen to be: a present
state as a weighted superposition of the poles
(Remark~\ref{gc:rem:pairsuffices}, Remark~\ref{gc:rem:boltzmann}).
The agent realized: $\nu F$ --- with $\mu G$ as its privation, so
the scheme's evaluative asymmetry arrives as the complementarity
theorem rather than as an assumption. The discipline of the
transition: the anatomy condition and the Score's pricing of
movements between the poles (Proposition~\ref{gs:prop:score},
\S\ref{gc:sub:acting}). And the twist that makes the middle term
re-derivable where its classical warrants died: $\nu F$ is
\emph{telos-shaped without final causation} --- a fixed point of a
modal operator over a conditioning, in which nothing acts, pulls, or
intends (Remark~\ref{ra:rem:investreg},
Remark~\ref{gc:rem:teleology}) --- so the modern objection that
killed the classical middle term has no purchase on its successor,
which is not a cause at all.

Two identifications complete the anatomy. The verdict that emotivism,
false as an account of meaning, became true as an account of use is,
in this subsection's vocabulary, the testimony collapse observed from
outside: once every node is known to be biographically sourced, moral
utterance \emph{can} function only as preference plus pressure,
however valid the propagation --- the collapse is the mechanism of
the incommensurability, not a further symptom of it. And the
historical argument that the Kantian project had to fail is answered
at its joint rather than disputed: a finite rule \emph{inherited} its
content and jammed, as above, where the present framework inherits no
content --- it derives the middle term first and prices content
against it --- and the criterionless choice at which the
interminability account locates the final breakdown survives here as
\emph{priced} choice: adoption remains free ($\Phi2$), and what is
deleted is only the criterionlessness. The claim all of this makes is
calibrated at the meta-level, as the fences below require: the hole's
damage was that no dimension of comparison existed among rival
premises, so there was no fact of the matter to be wrong about; the
criterion restores the dimension --- commensurability with respect to
the root, determinate though uncertifiable
(Remark~\ref{gc:rem:determinacygap}) --- and thereby re-poses a
question the diagnosis showed to be ill-posed. Re-posed, not solved;
the distance between those words is priced below. One resonance is
recorded rather than argued: the interminability tradition models its
remedy on the response to Rome's fall --- local communities
preserving practice through a dark age --- while this framework's
backward chain prices Rome's fall as a node in the conditioning
itself (Remark~\ref{gc:rem:backwardchain}). The same event serves one
tradition as the exemplar of retreat and this one as a link in the
route, which is as compact a statement of the delta as the record
supplies.
\end{remark}

% [Instates the author's dictation: the crisis classified as
%  Departric-from-its-vantage; the concept-level reversal; the
%  author's assessment of the formulation's rank; the fall-of-Rome
%  closure.  The 2001 and Tattoo exhibits are reserved for the
%  Introduction per the author's plan.]
\begin{remark}[The crisis, classified: Departric from its vantage]
\label{gs:rem:crisisdepartric}
The framework can now say what the post-Nietzschean condition
\emph{is}, in its own vocabulary, and the identification is exact
rather than evocative. Past the Departure Point the adopted ordering's
indices are undefined rather than small
(\S\ref{gs:ssub:departric}); comparison does not fail by returning
bad values, it fails by returning none, and from inside that failure
every direction is as good as every other --- the principle of
explosion, lived. Nietzsche's report of the condition --- up is down,
the horizon wiped away, straying as through an infinite nothing
\cite{nietzsche1974gay} --- is the first-person phenomenology of
exactly that invalidity, produced without the concept of the Golden
Point, without even the concept of the Departure Point, and the
implicative parallel is nonetheless exact: a meaning-system whose
root is gone \emph{is}, from its own vantage, in the Departric
condition, and the metaphysics a vantage inhabits is, in the epistemic
sense this framework has made precise
(Remark~\ref{to:rem:vantage}, Remark~\ref{gc:rem:determinacygap}),
what the metaphysics from that vantage \emph{is}. The dark age the
interminability tradition forecast is the sociology of that
condition run forward.

Which is why the reversal works at the level it works at, and why the
author ranks this formulation as he does. Departric invalidity is
\emph{conditional}: the indices are undefined given departure, and
what the mere \emph{statement} of the Golden Point does --- the
concept, prior to any question of attainment --- is restore the pole
against which the indices are defined. A vantage that represents
$\nu F$ is, by that fact alone, no longer in the epistemic Departric
condition, whatever else remains uncertain; the explosion is not
argued down but \emph{deprived of its antecedent}. This is the
precise sense of the author's claim, entered here in his voice at
\textbf{[$\Phi$]}: the concept prevents the forecast dark age by
nipping it at the metaphysical root --- not by winning the
trajectory-argument on its own terms, but by changing what the
metaphysics from the present vantage is. And it is why the fallback
the interminability tradition lacked now exists: a
centralization-allowing meaning and moral root that requires neither
tradition, nor practice, nor retreat --- so the response to the
forecast need not be the Benedictine one, because the analogous fall
of Rome is, on this showing, avoidable at the root rather than
survivable in enclaves (Remark~\ref{gs:rem:middleterm}). The claim's
grade is inherited from its parts: the identification is exact, the
reversal is a statement about vantages and is available at every rung
of the ladder, and what remains Strong-conditional is only the
ontological backing, priced where it is claimed.
\end{remark}

\emph{The fences, which this subsection needs more than any other.}
Nothing here solves the seminar. The identification map is missing and
the computation ban of \S\ref{gs:sub:fence} governs every sentence
above: the criterion supplies the \emph{direction of judgment} on
nodes and the standard of their non-arbitrariness, never verdicts on
demand. The criterion inherits the ceiling --- node-verdicts are
disciplined, not certified, and anyone claiming certified ones has
left the framework. The sourcing claim is graded by the Strong/Weak/
Void ladder of \S\ref{to:sub:goldenaethus}, degrading gracefully to
the semantic premise's comparative defense under Void. And $\Phi2$
stands untouched: a non-arbitrary node is not thereby an obligatory
one --- the framework prices roots, and the reader still chooses.

% ---------------------------------------------------------------------
% [Instates the meaning material, author-directed (Aug 22): the third
%  route promoted to the paper's spine; the meaning ledger as
%  scaffolding for the forthcoming Introduction (author's WIP, kept
%  out of this pass by his instruction); citation keys chosen to match
%  the draft Introduction's own (nietzsche1974gay, weinberg1977first)
%  so the eventual merge is collision-free.]
\subsection{The third route: inferred meaning}\label{gs:sub:thirdroute}

The crisis of \S\ref{gs:sub:nodes} has a cultural face, and the
framework's answer to it deserves statement at the same register as
its academic face --- not least because the cultural face is the one
nearly every reader arrives already inhabiting.

\emph{The settlement.} For most of the modern era, educated
reflection on cosmic questions has operated under what may be called
the Nietzschean settlement --- the name is convenient rather than
exact, since Nietzsche marked an inflection in a longer arc rather
than its origin \cite{nietzsche1974gay}. Its content: traditional
meaning-structures have been undermined by the cumulative successes
of natural science in ways no honest agent can reverse by appeal to
faith, and the only remaining honest response is to construct meaning
oneself and commit to it as a deliberate act. Existentialism's
authenticity, absurdism's defiant embrace, secular humanism's
procedural commitments, and the sophisticated modern naturalisms are
strategies \emph{within} the settlement, not alternatives to it
\cite{Camus42, Sartre46, Carroll16}. The settlement's evidential
base is real: an arc running from Copernicus through Darwin to the
present in which scientific accuracy and existential significance
have run negatively correlated --- each deepening of comprehension
shrinking the apparent space for cosmic meaning --- whose most honest
statement is Weinberg's: the more the universe seems comprehensible,
the more it also seems pointless \cite{weinberg1977first}. The
framework's reading of that sentence is charitable in exactly the
sense its own machinery licenses: Weinberg was not committing a
category error; he was accurately extrapolating four centuries of
trajectory from within a vantage at which the resolving object did
not exist. The strengthening the author intends should be
stated exactly, because it is stronger than charity: from those
vantages the pessimistic extrapolation was not merely reasonable but
\emph{uniquely} reasonable --- the alternative, expecting the
trajectory to reverse without evidence after the Enlightenment had
spent centuries systematically dismantling every prior ground of
reversal, would have been the deus-ex-machina hypothesis. Nietzsche
sits near the descent's midpoint --- barely an inflection, since that
leg runs mostly linear --- and his own positioning report says as
much: \emph{I have come too early; my time is not yet}
\cite{nietzsche1974gay}. Weinberg, nearer the present, extrapolates
the leg he stands on, and the honest content of his sentence is that
the descent is probably not bounded from below. Watching the prophecy
arrive on schedule, only a fool would have extrapolated to anything
but a dark age; the interminability tradition's forecast
(\S\ref{gs:sub:nodes}) is that extrapolation done carefully. The same holds one register up for the diagnosis of
\S\ref{gs:sub:nodes}: Nietzsche, Weinberg, and the
MacIntyre--Anscombe finding are one trough with three faces ---
prophetic, physical, academic --- and this framework engages all
three with one object.

% [Instates the author's vantage-operator dictation: the flip
%  condition named (coupling of Miracles), the Tricky-Principle
%  precedent, the passive-influence clause, and the wirings.]
\begin{remark}[The inference operator, and where it flips]
\label{gs:rem:inferenceflip}
``Uniquely reasonable'' can be made mechanical, and the corpus
contains the precedent that shows how. Treat best inference as an
\emph{operator}: for a fixed question, rationality's verdict is a
function of the evidence corpus available at the vantage --- indexed,
per the author's own scoping and at his stated hedge \textbf{[S]}, by
the epistemic state and not by anything as deep as an Aethus, since
what varied across the generations was awareness, with meaning
\emph{uncertain} to the prior vantages rather than their ontology
superposed. On this schema the trough is not a mood but an operator
output: given the corpus running Copernicus through Weinberg, the
downward extrapolation was the unique verdict, and holding today's
conclusion \emph{from that corpus} would have been inferentially out
of line --- which is the exact and final form of the charity this
section has extended throughout.

\emph{The precedent.} The companion's Tricky Principle is the same
event one framework earlier \cite{Aethus}: whether the unobserved
operates by the observed's rules is unfalsifiable --- \emph{until the
discovery of quantum mechanics} --- and the principle's own procedure
records the flip in both directions: pre-quantum, the razor demanded
the uniformity solution, and it was the post-quantum corpus ---
one demonstrated object beyond observation running by different
rules --- that made the razor \emph{demand the reverse}, forcing a
new object tailored to unobserved states, the Aethic blank. Two
disciplines transfer from that instance. The flip is passive: the
empiricism arrives and the razor turns, with the author imposing
nothing --- and the generalization is tailored rather than scaled,
the companion's own caution, so that what the new corpus licenses is
a fitted conclusion and never a blunt one.

\emph{The present instance, and its flip condition.} The question is
the trajectory of meaning, and the evidential arrival that flips the
operator is named exactly: the \emph{coupling of Miracles}
\cite{Nexus}. While a record's improbabilities stand as a tally of
independent surprises, each is severally dismissible --- selection
artifact, coincidence, a multiple-comparison debt incurred per item
--- and the downward verdict stands. The companion's coupling
results change the object, not the mood: committing to one labeled
improbability \emph{forces} others, the record's label-set is the
closure of a root commitment with the statistical debt paid once at
the root, and disagreement between the cosmic and
observer-conditioned weightings is a \emph{detector} of coupling
rather than a suspicion of it. A tally can be shrugged at item by
item; a closure cannot, because there is only one commitment to
audit. That is the speedbump. Past it, the operator's output
inverts --- upward becomes the razor's demand rather than anyone's
hope --- and the inversion matures into its end state when the
coupled cascade is read as what it is: exotic structure proper, the
record's improbable features as a categorical kind rather than a
tally of surprises, at the interpretive grade the companion assigns
that reading. The claim's pedigree is thereby auditable end to end:
the flip inherits the coupling results at their stated grades, and
nothing in it is asserted from taste.

\emph{What the flip buys, in three wirings.} It dates the arc's
inflection at an evidential event rather than a generational
temperament --- the trough's far side begins where the corpus
includes the coupling, not where anyone decides to feel better. It
converts the call to action into an obligation transfer: this paper
imposes no conclusion; it delivers the empiricism, and the reader
whose corpus now contains the coupling finds their own operator
flipped --- running the trough's extrapolation from the new corpus is
what is now out of line. And it identifies the phase-transition
prediction's mechanism (Remark~\ref{gs:rem:valorvacancy}): the
transition is this flip at population scale, propagating exactly as
far and as fast as the evidence does, with the concept-level
reversal of Remark~\ref{gs:rem:crisisdepartric} as the flipped
operator's output stated at the metaphysical root.
\end{remark}

\emph{The taxonomy, which is the load-bearing move.} The settlement's
authority has always rested on an exhaustiveness claim: meaning is
\emph{constructed} (projected by the observer onto a world that
contributes nothing) or \emph{revealed} (delivered by the world
through a non-natural channel, the observer merely receiving), and
since revelation is closed to honest agents, construction is all that
remains. The corpus's claim is that the dichotomy was false. There is
a third route: \emph{inferred} meaning --- discovered by the observer
through study of the natural record under the conditioning structure,
the observer doing the inferring and the world supplying the
structure inferred from, neither pole doing all the work. The Nexic
companion opened the route in the backward direction: the record's
measured improbability cascade as significance present in the rocks,
quantifiable, and specific to the observer it produced \cite{Nexus}.
The present work completes it in the forward and ontological
directions: the Golden root as what things and categories arise from
(\S\ref{to:sub:goldenaethus}), the vantage itself included
(Remark~\ref{to:rem:vantage}), and a universe-sourced criterion for
the commitment primitives on which judgment stands
(\S\ref{gs:sub:nodes}). Backward significance and forward root: the
two chambers of the hole the settlement's sharpest diagnosticians
described, filled by one object at two tenses.

\emph{The claim, at its two tiers.} Stated with the grading the
framework requires of itself. The \emph{unconditional} tier: breaking
the dichotomy does not require the Strong reading, because ``the
universe is meaningful'' becomes, on this framework, a graded,
priced, falsifiable position for the first time --- a reader who ends
by assigning the Strong reading low credence has still lost the
settlement's exhaustiveness claim, since a third route now exists to
assign credence \emph{to}. The \emph{Strong-conditional} tier: under
the Route World's Strong reading the historical correlation itself
reverses --- comprehension and significance become positively
coupled, since every deepening of the record's analysis is a
deepening of the measured structure the conditioning concentrates on
--- and Weinberg's sentence acquires its successor: the more the
universe seems comprehensible, the more it also seems
\emph{significant}. The successor is not a denial of Weinberg's
reading but the next stage of the same arc, exactly as his sentence
was the next stage of Copernicus's.

\emph{The phenomenological signature, recorded because it is a test.}
Constructed meaning carries a tell: what one invented oneself, one
can usually feel that one invented --- it is the kind of thing one
must keep talking oneself into. Inferred meaning carries the opposite
tell, and the companion's register is its exhibit: a level of luck in
the face of improbability \emph{too vast to be conveyed and yet too
personal to be named} \cite{Nexus} --- the vertigo of the unchosen,
arriving with the specific force of what one demonstrably did not
author. The pair is offered as a matched diagnostic, not as proof:
feelings certify nothing (the ceiling governs here as everywhere),
but a route whose phenomenology is the structural opposite of
construction's is at minimum not construction relabeled.

\begin{remark}[The displaced-valor diagnosis, and the phase-transition
prediction]\label{gs:rem:valorvacancy}
One further observation of the author's belongs beside the taxonomy,
because it converts the vacancy from an abstraction into a pressure
reading. The exemplars that reliably produce the aesthetic charge of
valor --- martial music as the drums enter, anthems at their
perilous-fight lines, the scored last stand --- are, at their load
points, \emph{nationalistic}: humans rising to defeat other humans,
with the pride requiring the put-down of the defeated as its
complement. The diagnosis is charitable in the register this section
has used throughout: \emph{nobody is to blame for this}. The drive
being scratched --- the burning attachment of agency and valor to
something that matters --- is real and honorable, and if something
externally meaningful were \emph{known}, it would be loved; nothing
is known, so nationalism is what the drive can reach. That is the
node crisis (\S\ref{gs:sub:nodes}) read affectively: the valor has no
non-arbitrary object, so it attaches to the available arbitrary ones.

From this the author enters a prediction, at \textbf{[S]} and
falsifiable in the only way reception claims are: the vacancy is a
\emph{low-pressure region of vast extent}, and the introduction of a
comprehensible external catalyst into it would kick off a phase
transition rapid enough to surprise effectively every observer in
some respect --- and the mechanism is now named: the transition is
the inference flip of Remark~\ref{gs:rem:inferenceflip} at
population scale, propagating as far and as fast as the coupling
evidence does. The prediction is stated so it can be scored, and it is
stated \emph{with its danger attached}, because the two are one
claim: a catalyst capable of a phase transition is capable of
capture, and the anti-capture architecture of
\S\ref{gs:sub:notreligion} is therefore not an appendix to this
remark but its load-bearing condition --- the six deltas are what
permit a meaning-root to centralize \emph{attachment} without
centralizing enforcement. The same distinction prices the author's
strategic observation: the standing attempts to centralize
progressive coordination have run through the economic base
\cite{MarxEngels48}, and their historical capture record is the
dragon pattern in secular dress --- centralized enforcement seizing
centralized meaning --- where a root carrying the six properties has,
by construction, no certification to sell and no compliance lever to
seize. What is centralized is the object of the drive; what is not
created is any instrument of its policing. The fences hold as
everywhere: the framework outputs no platform and no party, $\Phi2$
stands, ``progress'' here denotes the Golden-directed conduct this
work has already priced (\S\ref{gc:sub:acting}) and nothing else ---
and the long game the seed is planted for is the one the whole paper
is about: the deliberate run to the Golden Point itself.
\end{remark}

\begin{remark}[The arrival's timing, read under the framework's own
guards]\label{gs:rem:arrivaltiming}
The author observes that, under the Route World, it is close to
predetermined that this material be written, and written now of all
times --- at the trough, when morale is lowest and the vacancy
widest. The observation is admissible, and it must be entered wearing
every fence the framework owns, because it is the framework applied
to its author. What the retrocausal reading licenses, at
Strong-conditional grade, is a claim about the \emph{role}: among
$G$-reaching continuations, weight concentrates on those in which the
concept arrives, and the invested reading prices the timing --- the
cheapest arrival is into maximal vacancy, where uptake pressure is
greatest, which is the trough. What no reading licenses is the
promotion of role to occupant: that \emph{this} formulation by
\emph{this} author is the arrival is occupancy until audited
(Remark~\ref{gc:rem:backwardchain}'s discipline, applied reflexively),
necessity is never election
(Remark~\ref{ra:rem:investreg}, Remark~\ref{gc:rem:teleology}), and
the narrative-bias guards of Remark~\ref{gc:rem:historyguards} apply
with special force to an author reading his own moment as pivotal ---
the deflationary alternative, that every author at every trough feels
sent, is recorded beside the claim as its standing rival. What
survives the fences is exactly this much, and it is not nothing: if
the framework is right, then the arrival of \emph{some} such concept,
at \emph{some} such trough, is what the conditioning concentrates on
--- and the reader holding this page is, on that reading, not
witnessing a coincidence. Whether they are witnessing this
sentence's version of it, no one certifies.
\end{remark}

\begin{remark}[The meaning ledger: scaffolding for the Introduction]
\label{gs:rem:meaningledger}
The forthcoming Introduction carries this material in the author's
own voice; this remark is its pointer map, so that every manifesto
sentence has a paying result and the manifesto is never flat. (The
master's own front door now exists at \S\ref{intro:sub:situation},
in the paper's register; the manifesto remains the rhetorical
instrument, and this map serves both.) ``The
meaning by which things take shape in the universe'' is
\S\ref{to:sub:goldenaethus} (constitution) with
Remark~\ref{to:rem:vantage} (the standpoint). ``Break free of the
existential freefall'' is \S\ref{gs:sub:nodes} (the node criterion)
with Proposition~\ref{ac:prop:forcing} (the one universe-sourced
frozen commitment already on the books). ``Not a return'' is the
taxonomy above plus the no-agent toggle
(Remark~\ref{to:rem:toggle}; the conditioning does not act,
Remark~\ref{ra:rem:investreg}). ``Weinberg inverted'' is the
Strong-conditional tier above, on the Accordance reweighting. ``You
cannot separate out the biosphere'' is
Remark~\ref{gr1:rem:goldendef} with Remark~\ref{to:rem:vantage} ---
the severing is not merely unimaginable but ill-posed, since the
individuation runs through the carrier. And the Renaissance the
Introduction will call a Miracle is the same node the backward chain
prices (Remark~\ref{gc:rem:backwardchain}). Each pointer is a
load-bearing wall behind a rhetorical fa\c{c}ade; the Introduction
may be read first, but it is paid here.
\end{remark}

% ---------------------------------------------------------------------
% [Placement pass against the meaning literature, author-directed
%  (Aug 22): existentialism and post-Nietzsche thought engaged at the
%  site of the third route, per house discipline (strongest form
%  first, deltas named).  The author expects possible pruning once the
%  Introduction matures; until then this subsection is the canonical
%  home and the Introduction should point here rather than duplicate.]
\subsection{The meaning literature, placed}\label{gs:sub:meaningplacement}

The settlement was engaged above as a frame; its interior contains
distinct programs, and the discipline this paper applies to its rivals
elsewhere --- strongest form first, deltas named, what survives stated
--- is owed here most of all, since these are the positions nearly
every serious reader arrives holding.

\begin{remark}[The constructed route's interior: the strongest
programs]\label{gs:rem:constructedwing}
\emph{Nietzsche's remedy.} The settlement's founder also wrote its
most ambitious constructive response: revaluation --- values
legislated by the sovereign individual, eternal recurrence as the
affirmation-test, amor fati as the achieved stance
\cite{nietzsche1974gay}. Three deltas. The node diagnosis of
\S\ref{gs:sub:nodes} applies to the remedy directly: self-legislated
values are commitment primitives whose source is the legislator, which
is exactly the arbitrariness the testimony collapse prices --- so the
remedy does not resolve the crisis but universalizes it, each
sovereign legislator a new node, mutually void. The recurrence test is
structurally a fixed-point condition run against the agent's own
capacity to affirm; the framework's test runs against the world's
conditioning --- affirmation against inference, and the difference is
the whole of the third route. And amor fati: where Nietzsche
\emph{commands} love of the realized route, the framework
\emph{derives} a portion of what that command asserts --- the
necessity-content of the realized route under the Golden target
(Remark~\ref{gc:rem:backwardchain}) --- which is the difference
between an imperative and a theorem, with the dark-investment fence
still governing: necessity is never endorsement, so what is derived
is smaller than what was commanded, and honestly so.

\emph{Camus.} The absurd is a two-premise clash: the human demand for
meaning against the universe's silence, with suicide and the leap both
refused and revolt held as the lucid posture \cite{Camus42}. The
framework's engagement concedes the structure and contests the second
premise as what it is --- an \emph{empirical} claim. Silence is a
measurement, and the Accordance reading returns a record that is not
silent but dense: the cascade is in the rocks, measured, specific to
its reader. And the leap-charge is met on Camus's own integrity
terms: inference from a public record with published falsifiers is
not philosophical suicide --- the third route passes his lucidity
test while failing, deliberately, his inventory of options. Nothing
here asks Sisyphus to smile; it disputes the mountain's blankness.
Given the inventory Camus possessed --- construct or leap --- revolt
was the honest posture, and the framework's quarrel is with the
inventory, never the integrity.

\emph{Sartre.} Abandonment --- no signs, no given values --- is the
settlement's exhaustiveness claim stated as ontology; existence
precedes essence; anguish marks the choosing; bad faith names the
flight into pretended givens \cite{Sartre46}. The framework contests
the precedence claim at the scale where it operates: under the Strong
reading the conditioning structure precedes and individuates
(\S\ref{to:sub:goldenaethus}) --- essence-like without an
essence-giver, since nothing acts. Anguish it does not contest but
\emph{locates}: it is the first-person phenomenology of the node
crisis, lived. And bad faith it turns: flight into pretended givens
is bad faith, but givens \emph{argued} --- graded, fenced, falsifier
published --- are not fled to; they are audited, and the third route
is the audit. What survives of Sartre is the word he actually needs:
the choice at adoption is real and remains the reader's
($\Phi2$), so the framework deletes his inventory and keeps his
freedom.

\emph{Kierkegaard.} The revealed route's most sophisticated modern
defense: objective uncertainty embraced as the very medium of the
leap, truth as subjectivity, the knight of faith
\cite{Kie1846}. The delta is clean --- he makes the object's
unreachability load-bearing, and the framework supplies the object
--- but the engagement ends on a consilience that should be recorded
rather than hidden: the ceiling forbids certification
(Theorem~\ref{gr:thm:ceiling}), so the third route, too, is
\emph{held rather than possessed} --- commitment under proven
uncertainty, passion with a referent. The difference is that his
uncertainty was constitutive and welcomed, while ours is theorematic
and priced; but a reader formed by Kierkegaard will recognize the
posture, and is owed the acknowledgment.

\emph{Heidegger, structurally.} Two joints map and are claimed, and
only these. Thrownness --- finding oneself already delivered into a
situation not chosen --- receives a formal referent: one is thrown
into a \emph{conditioned} record (Remark~\ref{gc:rem:carryover},
Remark~\ref{to:rem:vantage}). And being-toward-death as the
individuating horizon is engaged at biosphere scale by the hazard
base itself: mortal moments forever, with the finitude tradition's
descendants answered at Remark~\ref{ac:rem:questionneighbors}
\cite{Hei27}. The engagement is structural, not exegetical, and
claims nothing about the corpus beyond the two mapped joints.
\end{remark}

\begin{remark}[Allies and the analytic wing]\label{gs:rem:meaningallies}
\emph{Frankl is the century's nearest ally, and the nearest miss.}
Logotherapy's founding claim is the third route's slogan stated
clinically: meaning is \emph{detected rather than invented},
discovered in the world \cite{Frankl59}. The deltas locate the
contribution. Frankl asserts the discovered character on
phenomenological and clinical grounds and supplies no ontology on
which detection could be veridical; the framework supplies the
ontology. His scale is biographical where the root here is cosmic,
with individual cases inheriting through task-object individuation
(\S\ref{to:sec}). And his attitudinal wing --- meaning available in
unavoidable suffering --- resonates, at \textbf{[S]} and no
stronger, with where this framework locates hope: the critical
regime, written off by every event metric, is where the entire
survivorship structure lives.

\emph{Nagel, and which external view.} The absurd, on Nagel's
sharper statement, is the clash between engaged seriousness and the
external view from which nothing matters \cite{Nagel71}. The
framework's reply is a question: \emph{which} external view? The one
that deflates is the unconditioned one --- the bare Cosmic Nexus,
whose typical face is the void --- and deflation is exactly what
running that measure over a survivor's record produces. The honest
external view of a conditioned record is the conditioned one, and
run with the right measure the outside view returns significance
rather than deflation: the deflation was an artifact of
conditioning-blindness, not a discovery about the world. One rung
deeper, the view from nowhere is engaged by
Remark~\ref{to:rem:vantage}: a literal nowhere is the abyss, every
determinate view is indexed, and the question was never whether to
view from somewhere but which index the viewing consumes. Nagel's
irony survives as a tone; its ground does not.

\emph{Wolf's placeholder, filled.} The most influential
contemporary account holds meaning to arise when subjective
attraction meets objective attractiveness --- fitting fulfillment
--- with its author explicit that no theory of the objective side is
offered \cite{Wolf10}. The framework slots in as exactly that
missing half: contribution to the root
(Remark~\ref{to:rem:goldenroot}) is a theory of objective
attractiveness, graded and priced, and Wolf's bipolar structure
survives intact --- the subjective side remains genuinely necessary,
since adoption is chosen ($\Phi2$). The third route does not replace
the fitting-fulfillment view; it completes it.

\emph{Metz's taxonomy, and where this files.} The field's standing
survey sorts theories into supernaturalist and naturalist, the
latter into subjectivist and objectivist \cite{Metz13}. For
reviewers who will want the filing: this framework is a naturalist
objectivism of a type the taxonomy does not anticipate ---
\emph{conditioning-based} rather than property-based, external in
source, agent-free, graded by the Strong/Weak/Void ladder, with its
objective side carried by survivorship structure rather than by
metaphysical depth --- adjacent to Metz's fundamentality intuition
in direction, disjoint from it in mechanism. The density of near
neighbors across this subsection is, per the methodology companion,
the proximity corollary in action \cite{Ben26Method}; the deltas
above are the scoring key.
\end{remark}

\subsection{The reduction to religion, refused}\label{gs:sub:notreligion}

The third route will be met, predictably and from both directions,
with one reduction: \emph{this is religion again} --- theism in
survivorship clothing. The reduction deserves the same treatment the
Precipice received: its strongest form stated first, then the deltas,
each checkable on the page, because a defense that consists of
checkable structure is reliable in a way a defense that consists of
indignation is not.

\emph{The objection at its strongest.} A serious critic will not say
the framework mentions God; it does not. They will say the reduction
is \emph{functional}: an external source of meaning, a reverent
register, a vocabulary of Miracles and divinity, a future condition
shaping the present, a first-principles root for norms --- the
structure occupies religion's ecological niche, and whatever occupies
that niche will be captured by the dynamics that captured its
predecessors. The canonical statement of this style of objection is
Midgley's, against evolutionary cosmologies that smuggle escalator
myths and salvation narratives into scientific dress
\cite{Midgley85}, and it is a good objection: frameworks \emph{have}
done this, and the burden of showing this one does not is properly
on this work.

\emph{Six deltas, each verifiable on the page.} First, \emph{no
agent, at any grade}: the conditioning is a measure operation --- the
Doob transform at biosphere scale, the Nexus conditioning at cosmic
scale --- and it does not act, choose, or design
(Remark~\ref{ra:rem:investreg}, Remark~\ref{gc:rem:teleology});
adding an agent would add no explanatory power and would break the
derivations, which is why the reduction is not merely uncharitable
but \emph{ontologically incoherent} --- it attributes to the
formalism the one object the formalism cannot host. Second, \emph{the
epistemic channel}: revelation, scripture, and authority against
derivation from stated postulates with stated falsifiers --- the
Golden-categorical premise carries its falsifier
(Remark~\ref{gr1:rem:goldendef}), the fluctuation thesis its
quantifier check (\S\ref{to:sub:goldenaethus}), the Ansatz its
numerics that could have fired and did not; a religion does not
publish the conditions of its own refutation. Third, and sharpest:
\emph{the ceiling is an anti-dogma theorem}. Dogma requires
certifiable favor --- assurance of standing on the saved side --- and
the framework's central epistemic result forbids exactly that
certification, to every finite observer, forever
(Theorem~\ref{gr:thm:ceiling}, Remark~\ref{gr:rem:ceilingscope}). A
structure that \emph{proves} no one can certify the favorable side
cannot function as a salvation guarantee; the trough channel and the
losing condition are the structural opposite of an escalator myth,
which is precisely the Midgley checklist run and failed --- in the
framework's favor. Fourth, \emph{the fear mechanism is absent by
construction}: no damnation counterfactual (the Departric results
concern undefinedness of comparison, never torment), no in-group
salvation split, no compliance lever --- $\Phi2$ stands, the rule
outputs no duties, the node criterion disciplines without
certifying, and there is no institution whose authority substitutes
for the derivations, which anyone may rerun. Fifth,
\emph{severability}: the claims are graded and independently
adoptable --- a reader at Void keeps the entire mathematics and the
Nexic backward route; no faith is demanded as an entry price, where
the reduction's target is package-deal by nature. Sixth, \emph{the
genealogical delta}: on the record's own showing, dragon-mediated
meaning sources in the fear-inheritance and is captured for control,
while inferred meaning sources in measured improbability structure
--- different provenance, and checkable, which is what
``inferred'' means.

\emph{The vocabulary, priced openly.} The critic will quote
``computational divinity,'' ``Miracle,'' and their kin, and the
paper should pre-quote itself: the corpus borrows sacred register
deliberately, to mark magnitudes for which no secular vocabulary
exists --- and defines every such term mechanically
(``computational divinity'' is a statement about computation time;
a Miracle is an exotic-structure Prerequisite; nothing in either
definition hosts a will) \cite{Nexus}. The borrowing is a cost,
accepted with eyes open, for honesty about the phenomenology: the
awe is real, is shared with religious experience as a matter of
psychology, and is re-sourced --- the delta lives at the source and
the structure, never at the feeling, and a framework that denied the
feeling would be falsifying its own exhibit.

\emph{Why the refusal matters, in both directions.} The reduction is
refused not to win a taxonomy dispute but because accepting it would
be costly in exactly the two ways the record documents. Ontologically
it is incoherent, as above. And practically it would re-license the
capture channel: whatever is filed as religion inherits religion's
captors, and the fear-for-control mechanism that seized the
premodern meaning-structures --- documented in this corpus without
euphemism --- would have precedent to seize this one. The six deltas
are therefore also a \emph{prophylaxis}: each is a structural
property whose presence makes dragon-capture difficult in principle
--- no certification to sell, no duties to enforce, no agent to
speak for, no package to hold hostage --- and the paper commits to
them in advance, in writing, where the commitment can be checked.

\emph{What is not claimed.} That religious persons cannot read,
adopt, or contribute to this framework --- severability cuts that
way too, and the mathematics is indifferent to its reader's
metaphysics. That every religious phenomenon reduces to the capture
mechanism --- the corpus itself holds up universalist strands whose
structure the mechanism displaced, and the critique's target is the
mechanism wherever it occurs, secular instances included. And that
refusing the reduction settles the third route's truth --- it
settles only its \emph{category}, which is what a reduction
disputes. The route's truth is priced where it is claimed, at the
tiers and grades stated above, and nowhere else.


% ============================================================
% DELIVERABLE LISTS — revision 2
% ============================================================
% 1. CLAIMS REGISTER (deltas from rev 1)
%    - gs:prop:score NEW as tagged proposition (B2): [P] given
%      [A]+[C]; hard-wall exact via refined ∼; soft-killing under
%      strengthened Ansatz; else bounded-factor comparison (B1).
%    - Initiative expectation restated: ensemble form free (the
%      Tendency at practice level); trajectory form routed to
%      gr1:subsub:fluct at [O] (B3) — no trajectory claim made.
%    - Trapped cell finished (B4): rule collapses to comparison
%      of h_ast via gr1:eq:xstar; verdict = move ρ
%      (hr:rem:concession).
%    - Currency fence added (C): additivity of logarithms only;
%      no stock to hold, spend, or owe.
%    - A3 record: the discrepancy flag caught an error in the
%      BRIEF (gloss dropped prior terms); 2024 form carried,
%      converged across drafts.
%    (All rev-1 entries otherwise stand.)
% 2. CROSS-REFERENCES — as in body; NEW: sh:ansatz ×3,
%    gr1:eq:xstar ×2, hr:rem:concession, ac:prin:tendency,
%    ac:rem:notrajectory, gr1:subsub:fluct ×2,
%    gr1:cor:recurrentanatomy, ac:sec ×7 (was hardcoded \S2);
%    Claim~ kind word supplied for sh:claim:tail.
% 3. NEW LABELS — as rev 1 plus gs:prop:score.
% ============================================================

The objection also has a canonical \emph{case study}, and running
the checklist against it by name is cheaper than letting a reviewer
do it first. Tipler's Omega Point \cite{Tipler94,BarrowTipler86} is
the future-conditioned cosmology that \emph{was} captured by the
salvation register --- physical eschatology, resurrection mechanics,
guaranteed convergence --- and it fails this subsection's deltas
exactly where the present framework passes them: an agentive terminus
where this conditioning has no agent at any grade; physical
commitments since falsified where this work carries published
falsifiers; certifiable salvation where the ceiling proves
certification unavailable; a package where this work is severable
under Void. ``FAP redux'' is therefore the misfiling this subsection
exists to preempt, and the six deltas are precisely the distance
between the Golden Point and the Omega Point;
Remark~\ref{rw:rem:futurecond} states the placement in full.

% [Instated on the author's query: the chain existed pairwise across
%  ten remarks; this subsection states it once, whole, so it can be
%  audited as a unit. Earns its place as Part IV's synthesis.]
\subsection{The load path, stated once}\label{gs:sub:loadpath}

Part~IV's results are not a collection; they are one chain, and it
should be written out once, link by link, with each link's site and
grade, so that the whole can be audited the way its parts are.

\emph{The coupling of Miracles} --- the record's improbabilities as
the closure of one root commitment, not a tally
(\cite{Nexus}, at the companion's stated grades) --- forces
\emph{the inference flip}: the vantage-indexed operator's output
inverts, passively, on the Tricky-Principle precedent
(Remark~\ref{gs:rem:inferenceflip}). The flipped operator's licensed
object is \emph{the Golden Point stated as a concept} --- the third
route's referent, at the two-tier grade of
\S\ref{gs:sub:thirdroute}. The stated concept \emph{de-conditions
the explosion}: Departric invalidity loses its antecedent, at every
rung of the ladder (Remark~\ref{gs:rem:crisisdepartric}). With the
pole restored, \emph{the middle term returns} --- telos-shaped
without final causation (Remark~\ref{gs:rem:middleterm}) --- and
with the middle term, \emph{the node criterion}: commitment
primitives become judgeable, testimony recoverable in principle
(\S\ref{gs:sub:nodes}). A non-arbitrary object for the valor drive
permits \emph{re-attachment} (Remark~\ref{gs:rem:valorvacancy}),
whose population-scale form is \emph{the phase transition} --- the
flip propagating as far and as fast as the evidence does --- whose
test is \emph{the species-awareness criterion}, with its predicts
and its forbids (Remark~\ref{hr:rem:speciesaware}). And the chain's
terminus is \emph{owned conduct}: the anatomy work of
\S\ref{gc:sub:acting}, now representable as being \emph{about}
something, aimed at $\nu F$ under $\Phi2$ --- the reader still
chooses.

The chain's grading is its honesty. Links two through six run at or
near every rung of the ladder --- the flip is evidential, the
de-conditioning is logic, the criterion is semantics --- and this
near-unconditional prefix is where the paper's meaning-claims carry
their weight. Links seven through nine are reception-grade,
\textbf{[S]} and scoreable, entered with their dangers attached and
the anti-capture architecture as their load-bearing condition
(\S\ref{gs:sub:notreligion}). Link ten is conduct and is nobody's
but the reader's. Each link consumes only its predecessor, so a
reader who stops anywhere keeps everything upstream --- and a critic
who breaks a link inherits the obligation the chain makes explicit:
to say which one, since the sites and grades are now listed in a
single place for exactly that purpose.

\subsection{Ledger}\label{gs:sub:ledger}

\paragraph{Costs.}
\begin{enumerate}
    \item \emph{The identification map} (decision $\to$ $(\mu, D)$), \textbf{[O]} --- the section's principal cost. Without it \eqref{gs:eq:functional} is exact and unusable, and the section says so in the body rather than in a footnote.
    \item \emph{The Golden Score} (Proposition~\ref{gs:prop:score}), \textbf{[P]} given \textbf{[A]} and \textbf{[C]} inputs: exact for the hard-wall functional; for soft killing, under Ansatz~\ref{sh:ansatz} strengthened to the refined form, else a comparison up to bounded factors. Regime-scoped, with the table of \S\ref{gs:sub:regimes} recording the other cells.
    \item \emph{The Initiative's trajectory form}, routed to \S\ref{gr1:subsub:fluct} at the \textbf{[O]} of those dynamics; the ensemble form free (the Tendency at the level of practice); neither consumed below.
    \item \emph{The class-level consequence} (technostrength without technoanatomy scores nothing, possibly negatively), \textbf{[$\Phi$]} over Remark~\ref{gr1:rem:againstsing}; bounded above by Theorem~\ref{gr:thm:ceiling}'s bar on scoring actors.
\end{enumerate}

\paragraph{Buys.} A decision rule stated informally in 2024 turns out to be computable in the model's own terms, with the computed object the Doob density --- a derivation, not an analogy, now stated as a tagged proposition with its hypotheses in the bracket. The $\asymp$/$\sim$ qualification, a shared defect of both drafts caught at referee, makes the result more creditable rather than less. The rule's counterintuitive feature (difficulty prices in) survives the formalization intact. The trapped cell, finished, connects the rule to the paper's one anatomy-moving intervention rather than dying in a table row. And the section closes \S\ref{ac:sec}'s forward commitments: every \texttt{gs:} label pointed at from there now exists.

\paragraph{Costs of \S\ref{gs:sub:ethics}.} \emph{Golden ethics} ---
\textbf{[$\Phi$]} throughout and $\Phi6$-conditional at every step, adding
no premise; \emph{falsifier:} $\Phi6$'s and nothing further. Its two
moralities consume $\PP(\mathcal G\mid\cdot)$ for counterfactual Aethae,
which Theorem~\ref{gr:thm:ceiling} forbids anyone from supplying, so it is
exhibited as an exercise on a counterfactual the framework rules out rather
than offered for use.

\paragraph{Not claimed.} That the rule outputs duties ($\Phi2$ stands, and the fence restates it). That any actor's Golden Score can be computed, today or from any finite observation. That Aurutance can be measured absent the identification map --- no proxy is licensed. That the soft-killing limit converges as the Ansatz stands --- bounded factors are what it purchases. That Aurutance is a currency --- the additivity is that of logarithms and licenses no stock to hold, spend, or owe; Remark~\ref{ra:rem:aurutance} states what the economic vocabulary \emph{does} license, which is an exchange rate within one comparison and not a stock. That the Initiative causes the Tendency, or that any individual's participation moves a biosphere's hazard. That the viable or trapped regimes support the harmonic form.

\paragraph{Inherited dependency.} The section inherits \S\ref{ac:sec}'s exposure whole --- under fixed adverse $\rho$ there is nothing for the rule to compare, since $\hm$ collapses along the single surviving path --- and adds one of its own: everything from Proposition~\ref{gs:prop:score} onward stands on the refined tail asymptotics --- Theorem~\ref{sh:thm:persist} for $\tau$, and for $T$ the strengthened Ansatz --- so a repair there reprices this section before it reprices \S\ref{ac:sec}.

\paragraph{Meaning-wing addendum (structural pass).} Entries added
after this ledger's first writing, folded in at the move that restored
the ledger to the section's end: the node criterion and the
inference-side crisis, with the middle-term mapping and the Departric
classification of the trough (\S\ref{gs:sub:nodes}); the third route
with its two-tier claim and the inference flip that dates its
licence (\S\ref{gs:sub:thirdroute}); the placement of the meaning
literature, deltas as scoring keys (\S\ref{gs:sub:meaningplacement});
and the reduction refused on six checkable deltas
(\S\ref{gs:sub:notreligion}). Dependencies: the wing consumes the
ceiling, the conditioning ladder, and the ontological completion
(\S\ref{to:sub:goldenaethus}); what it adds to the section's own
exposure is confined to the reception-grade predictions, each entered
scoreable. The load path of \S\ref{gs:sub:loadpath} states the
wing's chain whole, with the grade of every link.

\part{The Human Position}

\noindent Part~V locates humanity in the structure the earlier parts
built. Two standing errors about the human role are refuted
(\S\ref{hr:sec}), the species-level self-awareness criterion is
stated with its test and its falsifier, and adversivity supplies the
microfoundation of technoanatomy at the scale where humans actually
act (\S\ref{ad:sec}). What leaves this part is the reader's own
position, which is the paper's point.

\section{The human role: two errors about it}\label{hr:sec}

\begin{quote}\small
Two positions occupy the ends of a spectrum concerning what humans are
for. At one end, techno-optimism: capability is the thing, and more of
it is the answer. At the other, a strain of environmentalism holding
that humanity is a net harm to the biosphere and that its voluntary
disappearance would serve the living world. Golden reasoning finds both
mistaken, and mistaken in ways that are structurally opposite --- each
neglects exactly the factor the other over-weights.

The argument against techno-optimism is already in hand and is
consolidated here: it neglects technoanatomy, and by
Remark~\ref{gr1:rem:againstsing} acceleration without anatomy is not
merely insufficient but actively worse than stasis
(\S\ref{hr:sub:optimism}). The argument against extinctionism is new and
is the section's substance: on the extinctionist's \emph{own} valuation
--- the biosphere's persistence --- removing the biosphere's only
hazard-lowering agent lowers the thing they mean to protect
(\S\ref{hr:sub:extinction}). Both arguments are conditional, and the
conditions are stated rather than assumed.

What emerges is a position on environmental value that neither end
occupies: nature just \emph{is} the biosphere, the Golden Point is the
only route by which it acquires an unbounded expectancy, and human
existence is on that accounting not an accident at odds with the living
world but structurally implicated in its only route to persisting
(\S\ref{hr:sub:environment}). The threat is not thereby diminished ---
$\rho<0$ remains the live danger, and this section is as much an
argument for urgency about anatomy as against despair about presence.
\end{quote}

% ---------------------------------------------------------------------
% [Instates the author's species-self-awareness thesis (dictated with
%  the Tattoo observation, Edinburgh, Aug 2026 --- contemporaneous
%  this time, per the standing slot-reservation practice; the Tattoo
%  and 2001 exhibits themselves are reserved for the Introduction).]
\begin{remark}[Species-level self-awareness: a criterion, and where
it places the discontinuity]\label{hr:rem:speciesaware}
The section's two errors share a presupposition worth naming: both
treat humanity as the kind of thing that already knows what it is
doing. The author's thesis is that it is not --- that
\emph{humanity, as a collective, is not yet self-aware} --- and the
framework can make the thesis exact instead of aphoristic.
Individual self-awareness is, on the companion's account, not merely
ancient but \emph{basal}: the latch result reads the spark as forged
in the shared Mesozoic crucible acting on the whole early-mammalian
stock --- plesiomorphic across the crown group, inherited by descent
--- with three consequences the received picture does not take: the
mirror test \cite{Gallup70} is category-error-shaped as an
instrument, measuring a visual modality where the property is
representational (the companion's blunt reductio: on the received
reading dogs are not self-aware; on the latch's, that verdict is an
artifact of testing a scent-first animal through glass); the latch
\emph{predicts} modality-appropriate self-recognition where the
modality is corrected, with the olfactory-mirror result the worked
instance of independent convergence \cite{Nexus, Horowitz17}; and
the latch \emph{forbids} findings that would count against it,
stated in advance. The property itself is representational: an
entity is self-aware when it can represent its own situation as its
own, including what it would be for that situation to be won or
lost. The species-level criterion below therefore does not introduce
a new concept; it \emph{lifts one that already carries evidential
pedigree} --- the same property, one scale up, with the
proxy-critique transferring wholesale (species-scale behavioral
proxies would be as category-error-shaped as the mirror) and with
the latch's own discipline inherited: the criterion \emph{predicts}
that concept-reception produces a measurable coordination shift ---
which is the phase-transition prediction of
Remark~\ref{gs:rem:valorvacancy}, now recognizable as this
criterion's test --- and it states what would count against it: a
species that received the poles' concept and exhibited no ownership
shift and no coordination consequence would leave the criterion
idle, and the idleness would be evidence. Lift the property to species scale and
the criterion writes itself: \emph{a species is self-aware when it
possesses the concept of its own modal poles} --- when its losing
and winning conditions, $\mu G$ and $\nu F$, are represented within
it as such. By that criterion humanity has, to date, failed the
test: it acts, at scale, with the executive competence of an agent
--- and has no representation of what its actions are \emph{of}.

Two consequences follow, and the second dissolves an apparent
paradox. First, the discontinuity is relocated. On the author's
reading, the evolution of individual agency was smooth --- exotic
structure throughout, no single monolith moment --- so the
turning-point picture, if it applies anywhere, applies not to the
lineage's past but to the species' \emph{conceptual} history: the
one genuine discontinuity available is the reception of the poles'
concept, an event datable in principle and, as of this writing,
ahead. Second, the record's shape before that event is not
paradoxical. That humanity already exhibits agency-shaped conduct
--- coordinated, executive, consequential --- without possessing the
concept its conduct would be about is exactly what the conditioning
predicts: under Past-Alignment the record is Golden-conditioned
(\S\ref{gc:sub:pastalign}, Remark~\ref{gc:rem:carryover}), so
$G$-consistent conduct precedes $G$-representation the way any
conditioned trajectory precedes its description. The species could
be augmented with the executive character of agency before knowing
there was a ``something'' its agency was toward --- and the concept's
arrival, when it comes, changes not the conduct's existence but its
\emph{ownership}. The criterion is the claim entered here; whether
any particular formulation is the one through which reception occurs
is occupancy, priced at Remark~\ref{gs:rem:arrivaltiming}, and the
exhibits the author intends for the Introduction are reserved there.
\end{remark}

\subsection{Against extinctionism}\label{hr:sub:extinction}

\subsubsection{The position, in its own terms}\label{hr:subsub:position}

The target is a specific and explicitly biosphere-valuing position, and
it must be stated in its own terms or the argument does not engage it.
The Voluntary Human Extinction Movement holds that humanity should
abstain from reproduction so as to bring about its own gradual
disappearance, on the ground that this would prevent environmental
degradation \cite{VHEMT}. The reasoning proceeds \emph{from} the value
of the living world: the movement's own framing weighs the greater value
of Earth's entire biosphere against the lesser value of one species, and
holds that once one envisages the biosphere entire and grants the
intrinsic value of every life form, voluntary extinction begins to make
sense.

That framing is what makes the position addressable here. Its premise
--- that the biosphere's persistence and flourishing is what matters ---
is one this framework can evaluate, because persistence of a biosphere
is the quantity \S\ref{gr1:sec} models.

\begin{remark}[What this argument does not touch]\label{hr:rem:scope}
It is essential to distinguish the target from a neighboring position
it is often assimilated to. Benatar's antinatalism holds that coming
into existence is always a serious harm, and concludes that human
extinction would be preferable, on an asymmetry between the badness of
pain and the goodness of pleasure \cite{Benatar06}. That argument does
not proceed from biosphere value at all; its currency is suffering, and
extinction is a consequence rather than the objective.

\emph{Nothing below engages it.} An argument that removing humans lowers
the biosphere's expectancy is no reply to someone who never valued the
expectancy. The suffering-focused position must be met on its own ground
and this section does not attempt it. What follows applies to
extinctionism \emph{in its environmental form} --- and that restriction
is not a weakening, since the environmental form is the one that appeals
to a premise the framework shares.
\end{remark}

\subsubsection{The argument}\label{hr:subsub:extinctionargument}

\begin{proposition}[Removing the muffling agent forfeits viability;
\textbf{[P]} given the model]\label{hr:prop:extinction}
In the model of \S\ref{gr1:sec}, a biosphere with no hazard-lowering
agent is the case $x\equiv0$: no muffling is accumulated, and the total
hazard reduces to the background, $h_t=B_t$. Then
$\Lambda_t=\int_0^t B_u\,du\to\infty$
(Lemma~\ref{tc:lem:infbackground} with $S=[0,\infty)$), so
\[
\PP(\Lambda_\infty<\infty)=0
\]
identically --- not small, but zero, and for every parameter value.

Compare the alternative. By Theorem~\ref{gr1:thm:twocond}, a biosphere
that does accumulate muffling has
$\PP(\Lambda_\infty<\infty)=\one\{\rho>0\}\cdot
\bigl(s(\mu_0)-s(-\infty)\bigr)/\bigl(s(+\infty)-s(-\infty)\bigr)$,
which is strictly positive whenever anatomy is positive and the driver
escapes --- transient toward $+\infty$, or (by the theorem's repaired
recurrent clause) positive recurrent with positive stationary mean,
which \emph{widens} the class of muffling dynamics this comparison can
invoke. The extinctionist proposal is therefore a choice of an option
with viability probability exactly zero over one whose viability
probability may be positive.

Two scopings keep the zero honest at the site. The \textbf{[P]} is
model-internal: it prices the fixed-configuration model, and the
model-to-world step remains this section's ledger cost at
\textbf{[$\Phi$]}. And the zero is configuration-conditional: the
recovery branch --- a successor lineage re-evolving the muffling
capacity --- is precisely the branch a fixed-biosphere model does not
represent, and the corpus prices rather than ignores it: that branch
is the $\varepsilon$ of $\varepsilon$-Departure
(Definition~\ref{gc:def:epsdep}), where the pricing sentence lives.
\end{proposition}

The magnitude of what is forfeited is not left abstract. The companion
treats the external hazard as roughly constant over long timescales,
since the cosmic background does not change appreciably from era to era,
giving a survival probability $e^{-\lambda t}$ and an external life
expectancy of $1/\lambda$ \cite{Nexus}. The framework's own reading of
the Rare Earth and High-Contrast Earth hypotheses is that
$\lambda t_0>1$, where $t_0$ is the length natural history actually
took --- which is to say that \emph{Earth's remaining expectancy under
the external hazard alone is shorter than the span already elapsed.}

\begin{remark}[Undisturbed nature is not preserved nature]
\label{hr:rem:undisturbed}
This is the point on which the position turns, and it is worth stating
without ornament. The extinctionist picture is of a living world
continuing undisturbed once the disturbance is removed. But
``undisturbed'' is not a stable state; it is the case $x\equiv0$, in
which the biosphere has no means of lowering its own hazard and the
background runs unopposed.

The reframing of Definition~\ref{gc:def:departure} sharpens what is
being chosen, and sharpens it against the position. Under a reading on
which the relevant event is biological death, ``leave nature alone''
sounds like survival, and the argument must show that death follows
anyway. Under the reading on which the Departure Point is the closure of
possibility, no such demonstration is needed. At $x\equiv0$,
Proposition~\ref{hr:prop:extinction} gives
$\PP(\Lambda_\infty<\infty)=0$ identically: \emph{no} continuation
available to that biosphere advances. The modal quantifier of
Definition~\ref{gc:def:departure} is therefore satisfied at once, and
the Departure Point is not a fate awaiting the undisturbed world but
its present condition. In the terms of
\S\ref{gc:subsub:cascade}: $x\equiv0$ lies in
$\mu G\setminus\mathrm{Base}$ --- it is a departure at which no hazard
has yet fired, entered by configuration rather than by event. The
extinctionist proposal is a proposal to enter the cascade
deliberately. A biosphere left alone has not been preserved;
it has been placed in the one configuration from which nothing
leads anywhere. The extinctionist is not choosing a world that
survives without us. They are choosing a world that cannot go
anywhere. Destabilization of the star, a sufficiently
close hypernova, orbital perturbation by a passing mass --- these do not
become less likely because humanity has withdrawn. They become the only
things that happen. Remark~\ref{sch:rem:reliance} reaches the same
configuration from the other side: a biosphere that has not completed its
schedule has not yet acquired what the extinctionist proposes to give up,
and the two sit at the same bare background.

And the choice is not between a damaged biosphere and an undamaged one.
It is between a biosphere with a positive probability of persisting
indefinitely and one with none. On the extinctionist's own valuation,
the second is worse.
\end{remark}

\subsubsection{The concession, and why it does not rescue the position}
\label{hr:subsub:concession}

The argument has a real limit, and stating it is what keeps it honest.

\begin{remark}[Adverse anatomy does not vindicate the position]
\label{hr:rem:concession}
The natural concession at this point is that under $\rho<0$ the
extinctionist is closer to right, since $x\equiv0$ leaves hazard merely
at background while adverse anatomy drives it upward. The concession
should not be made, and seeing why changes which options the argument is
really between.

\emph{First, the local comparison goes the other way.} A biosphere with
adverse anatomy that \emph{halts} at the optimum $x_\ast$ of
Proposition~\ref{gr1:prop:trichotomy}(iii) sits at hazard $h_\ast$, and
$h_\ast<h(0)$ for every $\rho<0$. Numerically, at $B=1$, $b=3$, $k=6$:
\begin{center}\small
\begin{tabular}{@{}rrrrr@{}}
\toprule
$\rho$ & $h(0)$ & $x_\ast$ & $h(x_\ast)$ & expectancy gain\\
\midrule
$-0.05$ & $1.00$ & $11.47$ & $0.0046$ & $217\times$\\
$-0.20$ & $1.00$ & $8.88$ & $0.0176$ & $57\times$\\
$-1.00$ & $1.00$ & $4.28$ & $0.407$ & $2.5\times$\\
$-2.00$ & $1.00$ & $1.82$ & $0.898$ & $1.1\times$\\
\bottomrule
\end{tabular}
\end{center}
Even at $\rho=-2$ --- anatomy adverse enough that capability doubles
internal hazard twice over --- halting still beats having no muffling at
all. So $x\equiv0$ is strictly dominated for every $\rho<0$. Adverse
anatomy is worse than absence only if the biosphere \emph{fails to
halt}, which is a failure to stop rather than a property of the anatomy.

\emph{Second, and this is the point the local comparison obscures, both
options forfeit.} By Theorem~\ref{gr1:thm:twocond},
$\PP(\text{viable})=\one\{\rho>0\}\cdot(\text{escape factor})$. At
$x\equiv0$ there is no muffling agent, hazard sits at background, and
$\Lambda_\infty=\infty$. At $\rho<0$ halted at $x_\ast$, hazard sits at
$h_\ast>0$, and $\Lambda_\infty=\infty$. \emph{Both are zero.} The
extinctionist and the biosphere that settles at $x_\ast$ are not
choosing between a worse and a better outcome; they are choosing between
two ways of dying, and the model does not distinguish them in the only
respect that admits an answer other than death.

The response to adverse anatomy is therefore not to select among the
zeros. It is to move $\rho$, which is the one intervention that changes
the indicator rather than the multiplicand --- and which
\S\ref{ad:sec} exists to model. A reader persuaded of the danger and
unpersuaded that $\rho$ can be moved still has a coherent position, and
it remains the strongest form of the view opposed here; but it is a
claim about the tractability of anatomy, not a defense of extinction.
\end{remark}

\subsubsection{Why forfeiting is not the same as never having had}
\label{hr:subsub:stakes}

\noindent Purpose and its grounding are treated once, at
\S\ref{ac:sub:question}; this subsection confines itself to stakes.

The preceding leaves a gap that arithmetic alone will not close. If both
options carry $\PP(\text{viable})=0$, why prefer either to the other,
and why regard the choice as weighty rather than indifferent?

The answer is that the two are not symmetric with respect to what has
already been spent --- and note that under
Definition~\ref{gc:def:departure} what is at stake is the trajectory
rather than the lives on it, which is why the accumulation is the right
object to weigh. A biosphere at $x\equiv0$ has not merely a low
viability probability; it has \emph{surrendered} an accumulation, and
the accumulation is the subject of the companion's Miracle sequence
\cite{Nexus}. Whatever its correct magnitude, the record is a chain of
events each of which had to succeed for the next to be attempted: an
origination, a transfer if panspermia obtained, the eukaryotic
transition, the Dark Era, the end-Cretaceous impact, the
$\alpha$-gal epidemic. The companion's estimate for a single such
step --- the interstellar transfer --- is of order $10^{-21}$
\cite{Nexus}.

\begin{remark}[The stakes rise monotonically]\label{hr:rem:monotone}
One structural feature of that chain is worth extracting, because it is
what makes the argument an argument rather than an appeal to sunk cost.
Each Miracle in the sequence is \emph{more} consequential than the last
--- in the companion's formulation, because the gap to the Golden Point
was by then that much closer \cite{Nexus}. The reason is not sentimental
but combinatorial: a step late in the chain carries all the improbability
of every prior step behind it, so the expected cost of failing at step
$n$ is the whole product $\prod_{i<n}p_i$ rather than $p_n$ alone.

Two consequences follow. The stakes attaching to a decision rise
monotonically along the trajectory, so the present moment carries the
highest stakes the lineage has yet faced --- not by self-importance but
by construction. And sunk cost is not the structure here: the classical
sunk-cost fallacy consists in letting past expenditure govern a choice
whose future payoffs are unchanged by it, whereas here the past
expenditure has \emph{purchased} the option that remains, and forfeiting
is forfeiting that option rather than honoring a debt. The distinction
is that between a bad reason to persist and a live asset one may
discard.
\end{remark}

\begin{remark}[Forfeiture is not observable]\label{hr:rem:silentforfeit}
The cascade of \S\ref{gc:subsub:cascade} adds a feature to the stakes
argument that deserves separate statement, because it is the feature
that makes the argument urgent rather than merely large.

A biosphere enters $\mu G$ at the moment the last advancing route
closes, and by \eqref{gc:eq:departurefix} that moment precedes the
firing of any hazard --- possibly by a great deal. Nothing marks it.
There is no event to observe, no loss to register, and by
\S\ref{gr:sub:ceiling} no finite observation that would distinguish a
biosphere which has just crossed from one which has not. The
accumulation is not forfeited in some future catastrophe one might see
coming and reconsider; it is forfeited silently, at whatever point the
routes ran out, and the catastrophe is only the execution.

This does not make the stakes argument more permissive --- see the
limits below --- but it does relocate what the argument is about. It is
not a claim that a visible disaster would be very bad. It is a claim
that the relevant transition is invisible, which is why deferring
attention to anatomy until the danger is legible is deferring it past
the point where deferral is possible.
\end{remark}

\begin{remark}[What this does and does not license]\label{hr:rem:stakeslimit}
The stakes argument establishes an asymmetry between forfeiting and
never having had, and nothing beyond it. In particular it does not
license: that any present cost is acceptable, since the argument is
silent on rates of exchange; that the Golden Point will be reached,
which \S\ref{gr:sub:ceiling} forbids anyone from knowing; or that
persons alive now bear a duty derived from the accumulation, which
$\Phi2$ excludes and which the framework has no machinery to generate.

What it does license is narrow and, we think, sufficient: that the two
options which forfeit viability are not the neutral defaults they are
often taken for, and that between an option which changes the indicator
and options which do not, the model has a determinate preference.
\end{remark}

% ---------------------------------------------------------------------
\subsection{Against techno-optimism}\label{hr:sub:optimism}

The opposite error requires less new argument, since the result it
founders on is already proved.

Techno-optimism holds that continued growth in capability is the answer
to the risks capability creates. In the model's terms it supplies the
escape factor of Theorem~\ref{gr1:thm:twocond} and is silent on the
anatomy factor --- $\rho$ does not appear in any accounting of energy
capture, capability, or rate of progress, and cannot be inferred from
them.

The verdict is not that this leaves the conclusion unsupported but that
it inverts it. With $\rho<0$, feedback does not fail to secure survival;
it \emph{accelerates the failure}, since $\mu_t\to+\infty$ drives
$x_t\to+\infty$ and hence, by
Proposition~\ref{gr1:prop:trichotomy}(iii), $h\to\infty$
superexponentially. Supplying capability alone is strictly worse than
supplying neither: a poorly-constituted biosphere that stagnates sits
near its hazard floor $h_\ast$, while one that accelerates runs at its
own cliff.

\begin{remark}[The cascade forces anatomy to be a process]
\label{hr:rem:forcesprocess}
Both arguments of this section, taken with
\S\ref{gc:subsub:cascade}, deliver a verdict on the fixed-$\rho$ model
that is worth confronting rather than absorbing.

If $\rho<0$ and $\rho$ is a \emph{constant}, then
Theorem~\ref{gr1:thm:twocond} gives
$\PP(\Lambda_\infty<\infty)=0$ identically --- at every time, including
$t=0$. By the modal definition the biosphere is then in $\mu G$ from the
outset: no continuation available to it ever advances, and it departed
before anything happened. That is a degenerate verdict. It says the
question of what such a biosphere should do was settled before it was
asked, which is not a claim about biospheres but a symptom of the
model.

The degeneracy has exactly one exit, and it is not a matter of taste.
It is avoided if and only if $\rho$ can change --- if adverse anatomy
now leaves open a continuation in which anatomy is not adverse later.
That is the fluctuating-anatomy setting of
\S\ref{gr1:subsub:fluct}, where
Corollary~\ref{gr1:cor:recurrentanatomy} and
Proposition~\ref{gr1:prop:tolerance} take over, and where the question
becomes whether the escape is completed in time rather than whether it
was ever available.

So the cascade converts emergent anatomy from a desirable extension
into a \emph{requirement of coherence}. A framework that treats
technoanatomy as a parameter and departure as modal closure will
pronounce every adverse-anatomy biosphere already lost; only a
framework in which anatomy has its own dynamics can say the thing this
section is trying to say, which is that the matter is open and the
window is closing. The arguments above should be read as holding under
the fluctuating reading, and the fixed-$\rho$ statements as the
limiting case in which the window has already shut.
\end{remark}

\begin{remark}[The two errors are the same error]\label{hr:rem:sameerror}
Placed side by side, the positions are not opposites but a single
mistake taken in two directions. Both treat capability as the only
variable. The techno-optimist maximizes it; the extinctionist minimizes
it; neither carries a parameter for whether capability, once acquired,
is turned inward. Once $\rho$ is in the model the spectrum they occupy
is revealed as one axis of a two-dimensional space, and the interesting
region --- positive anatomy with accumulating capability --- lies on
neither end of it. That is why the disagreement between them has proved
unresolvable by argument along their shared axis: the quantity that
settles it is not on that axis.
\end{remark}

% ---------------------------------------------------------------------
The refutation's arsenal has grown since this subsection was written,
and the additions belong on display here. The invariant of
\S\ref{gr:sub:generalgolden} says the operative threshold does not
live in the space capability measures at all; the grabby engagement
says visibility-weighted selection is not persistence-weighted
selection, so the loud are not thereby the lasting
(Remark~\ref{gr:rem:grabby}); and $\beta<\alpha$ is the one
inequality on which capability-first hope depends and which its
proponents never argue (Remark~\ref{gr:rem:betaalpha}). One further
point is sharper than an arsenal item and is entered with its edge
showing. The escalator-myth checklist that this framework was run
against and passed (\S\ref{gs:sub:notreligion}) \emph{fires} when
run against techno-optimism: guaranteed ascent, salvation through
capability, the future as reward --- the structure Midgley named is
present here in secular dress \cite{Midgley85}, with the two-condition
theorem as the specific fact the myth requires readers not to check.
The reduction-to-religion charge, wherever it is finally owed,
is owed to the position this subsection refutes.

\subsection{What this implies for environmental value}
\label{hr:sub:environment}

Three consequences follow, of decreasing formality and increasing
interest. They are stated as consequences of the foregoing and inherit
its conditions.

\emph{Nature is the biosphere.} The unit \S\ref{gc:sub:vocab} adopts for
methodological reasons --- the totality of living and technical activity
associated with a world, taken as one system --- coincides with what
environmental thought means by nature when it speaks of intrinsic value
in the whole rather than in particular organisms. The framework does not
need to argue for that identification; it arrives at the same unit from
the requirement that a conditioning target apply across the whole record
(\S\ref{gc:sub:continuity}).

\emph{The Golden Point is the only route to an unbounded expectancy.}
Every alternative has a finite one. A biosphere without a
hazard-lowering agent has expectancy set by the external background; a
biosphere with adverse anatomy has a shorter one still. Only the
conjunction of positive anatomy and accumulating capability yields
$\PP(\Lambda_\infty<\infty)>0$. Whatever else environmental value
recommends, if it recommends the living world's persistence then it
recommends passage through that conjunction, since no other route is
available. The modal form of the same claim --- that passage through the
Golden Point is the only route by which the Departure Point is avoided ---
is \S\ref{ac:sub:imperative}'s, where it rests on the complementarity of
\eqref{gc:eq:departurefix} rather than on the expectancy comparison made
here.

\emph{Human existence is not ontologically at odds with the living
world.} This is the claim most likely to be over-read, so its content
and its limits are given precisely.

\begin{remark}[What is and is not established]\label{hr:rem:notmistake}
\emph{Established, conditional on the foregoing:} the outcome in which
the biosphere persists indefinitely runs through the existence of an
agent capable of lowering its hazard. Humanity is, so far as the record
shows, that agent's occurrence in this biosphere. So on this accounting
human existence is not contingent noise afflicting an otherwise
self-sufficient nature; it is structurally implicated in the only
trajectory on which nature persists. Under Golden-conditioning
(\S\ref{gc:sec}) the implication is stronger still, since the record is
then read as produced under a conditioning that the passage occur.

\emph{Not established:} that the harm is small, that any particular
human activity is vindicated, or that the present trajectory is good.
The threat is real, $\rho<0$ is the live danger, and
Remark~\ref{hr:rem:concession} declines the concession that a biosphere
with adverse anatomy is worse than none, and the refusal is an argument
rather than an optimism. The claim concerns humanity's
\emph{structural position}, not its conduct, and the two are
independent: an agent can be necessary to an outcome and be failing at
it.

\emph{Also not established, and worth disclaiming explicitly:} any
conclusion about the standing of human meaning-making at large. It is
tempting to read the above as supplying, from a source outside human
self-interpretation, a verdict that humanity is not a mistake --- and as
answering a broader cultural pessimism thereby. The temptation should be
resisted at this grade. What the framework delivers is a conditional
structural claim about one biosphere's route to persistence, resting on
$\Phi3$ and on the Golden-conditioning argument, both of which are
argued rather than proved. That is a substantial thing to have. It is
not a philosophy of history, and presenting it as one would trade a
defensible result for an indefensible one. \S\ref{ac:ssub:meaning} makes a
claim in this neighborhood and does not disturb the refusal: what it offers
is conditional on an adopted premise ($\Phi6$) rather than issued by the
apparatus, and the difference between a conditional grounding and a verdict
from outside is the distinction this remark is drawing.
\end{remark}

% ---------------------------------------------------------------------
\subsection{Ledger}\label{hr:sub:ledger}

\paragraph{Costs.} \emph{The valuation premise} --- the argument against
extinctionism is conditional on valuing the biosphere's persistence, and
is no argument at all against suffering-focused antinatalism
\cite{Benatar06}; \emph{falsifier:} none required, the restriction is
stated. \emph{The identification of humanity as the hazard-lowering
agent} --- \textbf{[$\Phi$]}; \emph{falsifier:} a mechanism by which a
biosphere lowers its own hazard without anything playing that role.
\emph{$\Phi3$ and Golden-conditioning}, inherited by
\S\ref{hr:sub:environment} and carrying their grades unchanged.

\paragraph{Buys.} An argument against environmental extinctionism that
proceeds \emph{on its own valuation} rather than by denying it. A
diagnosis on which the two ends of the spectrum share one error rather
than opposing each other. And a position on environmental value that is
neither anthropocentric nor misanthropic, obtained without adding a
premise for the purpose.

\paragraph{Not claimed.} That human conduct is vindicated; that the
danger is overstated; that suffering-focused antinatalism is answered;
or that anything here bears on the standing of human meaning-making
beyond the structural claim stated.


\section[Adversivity and the microfoundation of technoanatomy]{Adversivity:
institutional balance and the microfoundation of technoanatomy}\label{ad:sec}

\begin{quote}\small
Technoanatomy entered \S\ref{gr1:sec} as a parameter: the exponent
$\rho$ governing whether capacity acquired is capacity turned inward.
Nothing there says where it comes from. This section builds the
apparatus intended to supply it --- a model of institutions that
inflict hazard on one another and reorganize under a single principle
--- and reports how far that apparatus gets on its own. It does not yet
reach $\rho$; \S\ref{ad:sub:toward} states precisely what remains.

The principle is that every institution seeks \emph{net adversivity
balance}: its weighted alliances should cancel its weighted enmities.
Its two motivating cases are historical, and the section's numerical
work reproduces both. Two literatures border it closely and are placed
in \S\ref{ad:sub:placement}; the principle is not a refinement of
either, and against structural balance in particular it has
\emph{disjoint} ground states.
\end{quote}

% ---------------------------------------------------------------------
\subsection{The principle}\label{ad:sub:principle}

\setcounter{premise}{6} % GUARD --- $\Phi6$ is now taken by minimal arbitrariness
% in \S\ref{ac:sub:countermeasure}; this premise is never cited by number.
\begin{premise}[Net adversivity balance; \emph{governing principle}]
\label{ad:phi:balance}
Each institution assigns every other a signed weight, positive for
alliance and negative for enmity, and seeks to make the sum of those
weights --- taken against the others' sizes --- vanish. It is pressed
away from having only allies exactly as it is pressed away from having
only enemies.\footnote{%
\emph{Provenance.} The principle is recorded in our journal of
August 2025, with both cases below and with the accompanying diagnosis
that political theory has attended to its various doctrinal categories
where the underlying process is one of stable and unstable equilibria:
``if there are all alliances, the weights shift and institutions start
targeting each other, and if there are more adversaries than not,
alliances form simply to counter that weight and for no particular other
reason.'' The formalization, the placement against the two literatures
of \S\ref{ad:sub:placement}, and everything numerical below are later
and are from the drafting collaboration.}
\end{premise}

Two cases motivate it, and both are recovered numerically in
\S\ref{ad:sub:experiments}.

\emph{Removal of a common enemy.} When the Welsh drove back the
Anglo-Saxons, the alliance among Welsh tribes did not survive the
threat that had produced it. On the principle, the alliances were never
about one another: they were the positive weight required to offset a
large negative one, and their justification vanished with it.

\emph{Arrival of a common enemy.} Conversely, the only circumstance
under which Earth's states would plausibly unify is the arrival of
something opposed to all of them --- whereupon, on the principle, they
would unify not by persuasion but by the same arithmetic running the
other way.

\begin{remark}[The load-bearing departure]\label{ad:rem:departure}
The principle's distinctive content is that \emph{all-allied is
unstable}. That is what the removal case requires, and it is precisely
where the principle parts from structural balance, in which all-allied
is a ground state. Everything in \S\ref{ad:sub:placement} turns on this
one disagreement.
\end{remark}

% ---------------------------------------------------------------------
\subsection{Placement: two literatures, and a disjointness}
\label{ad:sub:placement}

\paragraph{Structural balance.}

Signed-graph balance theory begins with Heider \cite{Hei46}. It was
formalized for graphs by Cartwright and Harary \cite{CH56}. A triad is
\emph{balanced} when the product of its three signs is positive, so that
$(+,+,+)$ and $(+,-,-)$ are stable while $(+,+,-)$ and $(-,-,-)$ are
not. The structure theorem states that a complete signed graph is
balanced if and only if it splits into at most two camps, positive
within and negative between. Continuous-time dynamics were given by
Marvel, Kleinberg, Kleinberg and Strogatz \cite{MKKS11} and a
statistical-physics treatment by Antal, Krapivsky and Redner
\cite{AKR05}.

The Welsh case is instructive in that vocabulary. While the
Anglo-Saxons are present, every pair of tribes sits in a triad with two
negative edges, which forces the third positive: the tribes are allied
by triad closure, and nothing in their own relations holds them
together. Remove the shared enemy and those props vanish, leaving
unbalanced triads that must repartition.

\begin{remark}[Disjoint ground states; \textbf{P}]
\label{ad:rem:disjoint}
The two principles are not variants. Writing $\hat E_{\mathrm H}$ for
the normalized Heider energy and $\hat E_{\mathrm A}$ for mean-square
net adversivity, at $n=5$ with unit weights:
\begin{center}\small
\begin{tabular}{@{}lrr@{}}
\toprule
configuration & net adversivity & Heider energy\\
\midrule
all-allied (Heider ground state) & $(4,4,4,4,4)$ & $-10$\\
two blocs $3/2$ (Heider ground state) & $(0,0,0,-2,-2)$ & $-10$\\
pentagon$^{+}$/pentagram$^{-}$ & $(0,0,0,0,0)$ & $0$\\
\bottomrule
\end{tabular}
\end{center}
The adversivity optimum is Heider's \emph{worst} case --- maximally
frustrated --- and both Heider optima fail adversivity balance. Nor is
two-bloc structure available even approximately: with unit weights a
majority bloc of size $p$ and a minority of size $q$ would need
$p-1=q$ and $q-1=p$ simultaneously. The principle's stable states are
frustrated ones, which fits the observation that real international
systems do not settle into two clean camps.
\end{remark}

\paragraph{Balance of power.}

The second neighbor is a different tradition, with a name confusingly
close to the first. Realist theory in international relations holds
that states balance against concentrations of power \cite{Wal79},
refined by Walt into balancing against perceived \emph{threat} rather
than raw capability \cite{Wal87}. Both motivating cases above are textbook for
that tradition, and our stipulation that perceived threat
tracks an adversary's present unity rather than its capacity to
recombine is close to Walt's central move.

The present model sits between the two: the signed-network formalism of
the first, the node-local balancing condition of the second. The nearest
existing constructions we know of are the energy-based structural-balance
dynamics of Marvel, Strogatz and Kleinberg and the continuous-time
balance models descending from Ku{\l}akowski, both of which are
edge-local; a node-local quadratic penalty over signed degree does not
appear in either, and we have not found it elsewhere. That is a negative
search and not a demonstration of novelty, and the claim is held at
\textbf{[O]} accordingly.

% ---------------------------------------------------------------------
\subsection{The energy}\label{ad:sub:energy}

Institutions are nodes; $w_{ij}\in[-1,1]$ is $i$'s weight on $j$;
$s_j>0$ is $j$'s size.

\subsubsection{The model}\label{ad:subsub:constrained}

The model is a two-term energy on a constrained set:
\begin{equation}\label{ad:eq:model}
\min_{W}\ \ a\,E_1(W)+b\,E_2(W)
\qquad\text{subject to}\qquad
W=W^{\!\top},\quad \sum_{j\ne i}|w_{ij}|=\beta ,
\end{equation}
with
\begin{align}
E_1&=\frac1n\sum_i\Bigl(\sum_{j\ne i}w_{ij}s_j\Bigr)^{2}
&&\text{adversivity balance; node-local, quadratic}\label{ad:eq:E1}\\
E_2&=\frac{-1}{2\binom n3}\!\!\sum_{i<j<k}\!\!
\bigl(w_{ij}w_{jk}w_{ki}+w_{ji}w_{kj}w_{ik}\bigr)
&&\text{Heider frustration; triad-local, cubic}\label{ad:eq:E2}
\end{align}
and $\beta=n-1$. Equation \eqref{ad:eq:E2} is computed from
$\operatorname{tr}(W^{3})$, which equals three times the triad sum when
the diagonal vanishes.

There is \emph{one} free parameter, the ratio $b/a$, and the two
constraints are substantive rather than technical: relations are
\emph{mutual}, and attention is \emph{finite}. Both are claims about
what an institution is, and \S\ref{ad:subsub:relaxation} records what
goes wrong if either is dropped.

\subsubsection{The relaxation used for computation}
\label{ad:subsub:relaxation}

Hard constraints are awkward to anneal, so the numerical work below
minimizes a four-channel relaxation in which each constraint is
replaced by a penalty:
\begin{align}
\hat E_3&=\frac{1}{8n(n-1)}\sum_{i\ne j}(w_{ij}-w_{ji})^{2}
&&\text{reciprocity; relaxes } W=W^{\!\top}\label{ad:eq:E3}\\
\hat E_{\mathrm b}&=\frac1n\sum_i\Bigl(\sum_{j\ne i}|w_{ij}|-\beta\Bigr)^{2}
&&\text{engagement budget; relaxes the norm}\label{ad:eq:Eb}
\end{align}
with $E=a\hat E_1+b\hat E_2+c\hat E_3+d\hat E_{\mathrm b}$. The two forms
agree on everything reported here. At $n=9$, $b=4$, over independent
seeds:
\begin{center}\small
\begin{tabular}{@{}lll@{}}
\toprule
& four-term relaxation & two-term constrained\\
\midrule
$a=0$ & lopsided $6/3$, $7/2$ & lopsided $6/3$, $6/3$\\
$a$ small & even $4/5$, $\hat E_1=1.78$ & even $4/5$, $E_1=1.78$\\
$a=1$ & frustrated web, $\hat E_1\approx0$ & frustrated web, $E_1=0.03$\\
\bottomrule
\end{tabular}
\end{center}
The relaxation is therefore a computational convenience and
\eqref{ad:eq:model} is the model. Three facts about the relaxation are
nonetheless worth recording, because each explains why the
corresponding constraint is in \eqref{ad:eq:model} at all.

\begin{remark}[The budget is not optional]\label{ad:rem:apathy}
Drop the norm constraint and the principle has a trivial optimum:
\emph{be indifferent to everyone}. Numerically, pure $\hat E_1$ reaches
exact balance at $n=6$ by leaving eight of fifteen ties weak. Balance by
apathy is not what the principle intends, and finite attention --- a
fixed total spent across one's ties --- is the substantive claim that
rules it out. This is why the norm appears in \eqref{ad:eq:model} as a
constraint rather than as a preference to be traded against others.
\end{remark}

\begin{remark}[Reciprocity is not emergent]\label{ad:rem:recip}
$E_1$ depends only on row sums and is therefore blind to the
antisymmetric part of $W$; $E_2$ constrains cyclic products and does not
force symmetry either. Removing $\hat E_3$ from the relaxation and
re-optimizing produces asymmetries at the maximum,
$|w_{ij}-w_{ji}|=2$: total one-sided devotion. Symmetry must therefore
be \emph{imposed}, which is why it is a constraint in
\eqref{ad:eq:model}. Its near-zero value at relaxed optima indicates
that it is cheap to \emph{satisfy}, not that it is redundant --- an
inference the low value invites and which the test above refuses.
\end{remark}

\begin{remark}[Normalize by the mean, not the maximum]
\label{ad:rem:normalisation}
This concerns the relaxation, and is the one point at which a natural
choice is badly wrong. Dividing $\hat E_1$ by its maximum --- attained
at the all-allied state ---
crushes it, since near any configuration of interest it then varies over
only $O(n^{-2})$, and the Heider term dominates at any comparable
weight. The mean-square form \eqref{ad:eq:E1} instead assigns the
balanced two-bloc state a cost of exactly $1$ \emph{for every $n$},
since there $\sum_j w_{ij}=(p-1)-q=-1$ regardless of size, matching the
Heider floor of $-1$. Only under this normalization is $b/a$ an
interpretable weight, and only then is the model's behavior independent
of $n$. No re-powering of $\hat E_2$ achieves the same effect, and
taking roots of it would destroy the sign structure that makes it
frustration in the first place.
\end{remark}

% ---------------------------------------------------------------------
\subsection{Ground states}\label{ad:sub:ground}

\begin{proposition}[Perfect balance requires $n\equiv1\pmod4$;
\textbf{P}]\label{ad:prop:congruence}
Among symmetric configurations with all ties saturated at $\pm1$ and
equal sizes, exact adversivity balance is attainable if and only if
$n\equiv1\pmod4$. Each node requires $(n-1)/2$ positive ties, forcing
$n$ odd; and the total number of positive edges, $n(n-1)/4$, must be an
integer, forcing $n\equiv1\pmod4$. The admissible sizes are
$n=5,9,13,17,\dots$
\end{proposition}

\begin{proof}
Both conditions are immediate; their conjunction holds exactly on
$n\equiv1\pmod4$. Exhaustive search confirms the minimum of
$\sum_i(\text{net}_i)^2$ is zero at $n=5$ and positive at
$n=3,4,6$, the first case being realized by the pentagon of positive
ties with the complementary pentagram of negative ones.
\end{proof}

\begin{proposition}[Two failure modes, by parity; \textbf{P}]
\label{ad:prop:parity}
Write $P$ and $M$ for the numbers of alliances and enmities, so that
$P+M=\binom n2$. Summing every node's ledger counts each tie twice:
\begin{equation}\label{ad:eq:parity}
\sum_i \mathrm{net}_i \;=\; 2\,(P-M).
\end{equation}
Since $P-M$ and $P+M$ share a parity, the total imbalance is
\emph{forced} away from zero whenever $\binom n2$ is odd:
$\sum_i\mathrm{net}_i\equiv2\pmod 4$. Two regimes follow.

\emph{Even $n$: the imbalance spreads.} Here $n-1$ is odd, so every
$\mathrm{net}_i$ is a sum of an odd number of $\pm1$ terms and is
therefore odd; \emph{no} node can balance, and the cheapest option
distributes the residual --- at $n=4$ the optimum is $(-1,-1,-1,-1)$.

\emph{$n\equiv3\pmod4$: the imbalance concentrates.} Here individual
balance is possible but the total is not, and because the penalty is
\emph{quadratic} it is cheaper to load a mandatory total of $2$ onto one
node ($\sum\mathrm{net}_i^2=4$) than to spread it over three
($-2,-2,+2$ sums correctly but costs $12$). Squares reward
concentration.
\end{proposition}

\begin{remark}[The scapegoat, and what it does not license]
\label{ad:rem:scapegoat}
Proposition~\ref{ad:prop:parity} accounts for every computed residual:
\begin{center}\small
\begin{tabular}{@{}rrrrr@{}}
\toprule
$n$ & $\binom n2$ & forced $\sum\mathrm{net}_i$ & predicted min
$\sum\mathrm{net}_i^2$ & computed\\
\midrule
3 & 3 (odd) & $\equiv2$ & 4 & 4\\
4 & 6 (even) & $\equiv0$ & 4 & 4\\
5 & 10 (even) & $\equiv0$ & \textbf{0} & 0\\
6 & 15 (odd) & $\equiv2$ & 6 & 6\\
7 & 21 (odd) & $\equiv2$ & 4 & 4\\
\bottomrule
\end{tabular}
\end{center}
At $n=7$ the optimum leaves six nodes exactly balanced and one carrying
the whole residual; annealing under continuous weights reproduces it,
six nodes within $0.17$ of zero and one at $-1.64$. So where the
arithmetic forbids universal balance, the cheapest arrangement is not
shared discomfort but a single universally-disliked member --- and this
is forced by counting rather than emergent from dynamics.

Two limits on the reading. The result assumes unit weights and equal
sizes; with continuous weights every $n$ reaches balance by weakening
ties, so the congruence marks where balance is \emph{free} rather than
where it is possible. And the scapegoat mechanism, unlike the
congruence, survives that relaxation: the quadratic penalty rewards
concentrating an unavoidable residual on one member whatever the
weights, which is the part worth carrying forward.
\end{remark}

\begin{remark}[Where factions appear]\label{ad:rem:factions}
Under the normalization of Remark~\ref{ad:rem:normalisation}, the two
channels are commensurate at $b\approx a$, and at $n=9$ the positive-tie
components remain a single connected web up to $b\approx2$, splitting
into two factions between $b=2$ and $b=4$. The interesting regime is
$b\approx4$, where the system exhibits \emph{both} a $5/4$ split
\emph{and} near-perfect balance, $\hat E_1\approx0.005$; only above
$b\approx8$ is balance abandoned, $\hat E_1\to1.6$, in favor of
saturated bipolarity.

At large $b$ the adversivity term takes on a further role, and it is
more than tie-breaking. Heider admits two ground states, one bloc and
two blocs, both at $\hat E_2=-1$; since all-allied costs
$\hat E_1=(n-1)^{2}$ while two-bloc costs $1$, any positive $a$ breaks
that degeneracy in favor of two blocs. But Heider is also indifferent
among two-bloc \emph{partitions}: every bipartition sits at $\hat E_2=
-1$ whatever the bloc sizes. Numerically at $n=9$, $b=4$, the effect of
$a$ across three seeds:
\begin{center}\small
\begin{tabular}{@{}rlrr@{}}
\toprule
$a$ & faction sizes found & $\hat E_1$ & $\max_i|\mathrm{net}_i|$\\
\midrule
$0$ & $6/3,\;3/6,\;7/2$ --- lopsided & $12.1$ & $4.67$\\
$0.001$--$0.1$ & $4/5,\;5/4$ --- even, every seed & $1.78$ & $2.00$\\
$1$ & no clean bipartition & $0.003$ & $0.10$\\
\bottomrule
\end{tabular}
\end{center}
Heider alone produces \emph{lopsided} blocs, since it cannot see their
sizes. A thousandth of $\hat E_1$ makes them \emph{even}, because
near-equal blocs balance far better --- which is a structural prediction
(bipolar systems should have near-equal camps) rather than a
degeneracy-breaking convenience. At $a\sim1$ the bipartition dissolves
into a frustrated balanced web. The adversivity term therefore selects
among Heider's ground states on a criterion Heider does not possess, at
every weight at which it is present.
\end{remark}

% ---------------------------------------------------------------------
\subsection{Two experiments}\label{ad:sub:experiments}

Both motivating cases are simulated by representing the common enemy as
an exogenous field: each institution carries a fixed tie $-h$ to it, so
its ledger reads $\sum_j w_{ij}s_j-h$. A first attempt that instead
\emph{initialized} an enemy node and let it optimize failed, the node
simply joining a faction; a threat must be exogenous rather than a
participant balancing its own books.

\subsubsection{Removal: the Welsh}\label{ad:subsub:welsh}

Seven institutions equilibrated at $h=1$, then relaxed at $h=0$ from
that state.

\begin{center}\small
\begin{tabular}{@{}lrrrr@{}}
\toprule
& mean tie & allied & opposed & $\hat E_1$\\
\midrule
threat present & $+0.146$ & 6 & 3 & $0.728$\\
threat removed, before relaxation & $+0.146$ & 6 & 3 & $1.476$\\
after relaxation & $+0.004$ & 4 & 4 & $0.004$\\
\bottomrule
\end{tabular}
\end{center}

The middle row is the mechanism displayed: nothing has moved, and the
same configuration has doubled in cost. The alliances were balancing the
field; without it they are surplus. Relaxation then sheds them, and the
two factions present under threat become \emph{three plus an isolate}.

\begin{remark}[Fragmentation, not repolarization --- and no flips]
\label{ad:rem:welshreading}
Two features distinguish this from what either neighboring theory would
predict, and the second is a limitation.

Balance theory predicts repolarization into two fresh camps, since two
camps is what a balanced complete graph must be. The present model
predicts \emph{fragmentation}: with no external field, the cheapest
balance is many small mutual enmities rather than two large ones. The
historical record on the Welsh after the Anglo-Saxon reverse ---
increased factional subdivision without descent into general war --- is
consistent with the second and not the first. We record this as a
retrodiction and not a test: the simulation was run before the record
was consulted, but the parameters were not fixed in advance and a single
case cannot discriminate.

A third feature was first recorded here as a limitation and is better
read as a match. \emph{No alliance inverted}: no pair went from ally to
enemy. What happened instead is that two weak alliances lapsed to
neutrality while one previously neutral pair turned hostile, so
alliances fell from six to four and enmities rose from three to four.
That is coalition dissolution accompanied by fresh hostility among
parties who were already distinct --- not the betrayal of close allies.
The historical pattern is the former: the Welsh tribes were separate
before the reverse and became mutually hostile after it, rather than
turning on partners with whom they had been closely bound. A model
producing ally-to-enemy inversions would in fact overshoot the record.

The distinction is worth keeping because the two are easy to conflate
and make different predictions. Betrayal requires a commitment term
penalizing intermediate $|w|$, so that ties cannot fade and must flip;
dissolution-plus-new-hostility requires no such term and is what the
present energy produces. Which of the two a given historical episode
exhibits is checkable, and the model is committed to the second.
\end{remark}

\subsubsection{Arrival: the aliens}\label{ad:subsub:aliens}

Nine institutions equilibrated alone, then a tenth appended with size
exceeding all nine combined and with its own weights \emph{frozen} at
$-1$ toward each --- it takes no view among them.

\begin{center}\small
\begin{tabular}{@{}lrrr@{}}
\toprule
& mean tie among the nine & allied pairs & $\hat E_1$\\
\midrule
alone & $+0.000$ & 10 of 36 & $0.001$\\
threat present & $+0.583$ & \textbf{28 of 36} & $8.1$\\
\bottomrule
\end{tabular}
\end{center}

Eight previously opposed pairs became allied --- inverted, not merely
softened. The reinforcement sustaining those enmities loses its force
once every ledger is dominated by a single enormous term.

\begin{remark}[Incomplete unification, and an asymmetry]
\label{ad:rem:aliensreading}
Two results here were not designed for.

\emph{Unification saturates short of completion.} Only $28$ of $36$
pairs allied, with $\hat E_1$ blown up to $8.1$: even universal alliance
supplies at most $+8$ against a threat term of $-18$, so the system
cannot balance and carries permanent strain. Overwhelming threat
produces near-total but incomplete unification, which is a more
qualified prediction than the verbal argument yields.

\emph{Arrival and removal are not symmetric.} Adding a threat produced
eight sign inversions; removing one produced none. The asymmetry follows
from size-weighting --- an adversary larger than everyone combined
dominates every ledger at once, where a unit-sized field leaves only a
small surplus to shed --- and yields the prediction that \emph{threat
arrival is a sharp transition while threat removal is a slow
relaxation}. Nothing was inserted to produce this.
\end{remark}

% ---------------------------------------------------------------------
\subsection{Toward technoanatomy}\label{ad:sub:toward}

The purpose of this apparatus is to supply what \S\ref{gr1:sec} assumes.
Internal hazard there is $\Gamma e^{-\rho x}$ with both constants
primitive; here internal hazard would be what institutions inflict on
one another, and $\rho$ would be how that scales with capability. In the
notation of \S\ref{gr:subsub:boundary}, where $\rho=-(\alpha-\beta)$
with $\alpha$ the growth of reach and $\beta$ of enforcement, the
adversivity model supplies referents: $\alpha$ is how fast the damage
one institution can do another grows with technology, and $\beta$ how
fast joint-agreement structure grows.

\begin{remark}[What is missing, stated exactly]\label{ad:rem:missing}
The section does not deliver $\rho$, and the obstruction is specific
rather than a matter of unfinished labor. This model contains conflict
dynamics and no coordination scaling: $\hat E_1$ through
$\hat E_{\mathrm b}$ describe how institutions arrange themselves given
their capacities, and nothing in them says how the \emph{cost of
maintaining joint agreement} grows as those capacities grow. That
quantity is $\beta$, and technoanatomy is entirely a claim about it.
Balance dynamics of any kind will supply $\alpha$ and be silent on
$\beta$.

The route from here is therefore twofold, and both parts are open
\textbf{[O]}. A coarse-graining must be specified --- merging
positively-tied clusters into single nodes at the next scope, with sizes
adding and weights averaging --- and the energy must be shown to keep
its form under it, so that every scope obeys the same dynamics and the
biosphere's $\rho$ is the top of a flow rather than a separate
posit. And a coordination cost must be added, whose scaling with
capability is the anatomy the whole construction is for.
\end{remark}

% ---------------------------------------------------------------------
\subsection{Ledger}\label{ad:sub:ledger}

\paragraph{Costs.} \emph{The principle} (\ref{ad:phi:balance}) ---
\textbf{governing premise}; \emph{falsifier:} an institutional system
observed to rest stably in universal alliance without an external
threat, or to be indifferent to the balance of its commitments.
\emph{The engagement budget} --- \textbf{structural premise}, without
which the principle is trivially satisfied by indifference;
\emph{falsifier:} institutions observed to hold arbitrarily many ties
at full strength. \emph{The weighting $b/a$} --- the model's \textbf{sole free
parameter}, controlling the multipolar-to-bipolar transition;
\emph{falsifier:} it is fitted rather than derived, and a system whose
faction structure is inconsistent with any single value across eras
would embarrass the model. \emph{Mutuality and finite attention} ---
\textbf{constraints}, not preferences; \emph{falsifiers:} stable
one-sided regard between institutions, and institutions holding
arbitrarily many ties at full strength, respectively. \emph{Completeness} --- the two-faction
structure theorem assumes a complete graph, and real institutional
networks are sparse; the incomplete case admits more than two factions
and the clean result is unavailable.

\paragraph{Buys.} A single principle reproducing both motivating cases,
with the removal case matching the historical pattern more closely than
either neighboring theory. A counting constraint that produces
scapegoat structure without it being posited. An arrival/removal
asymmetry that was not designed for. A construction with one free
parameter and two interpretable constraints, the four-channel form used
for computation being a relaxation that agrees with it. And a stated,
checkable route to $\rho$, together with an exact account of why this
apparatus cannot reach it alone.

% ============================ 

% ---------------------------------------------------------------------
\appendix

\section{Horizon theory}\label{app:horizon}

% ---------------------------------------------------------------------
\label{sh:sub:horizon}

\subsection{The driftless horizon is a timescale, and it is maximally
soft}\label{sh:subsub:softhorizon}

In the driftless model the naive horizon quantity
$\Pi(\mu,D):=\PP_{(\mu,D)}(D\text{ never returns to }0)$ vanishes
identically (Theorem~\ref{sh:thm:noimm}): null recurrence leaves no escape
probability to be gradual about, and the horizon must be recast in the only
currency the model possesses --- time.

\begin{proposition}[Recovery cost is sublinear in credit; \textbf{P}]
\label{sh:prop:recovery}
Let $\varrho(d)$ denote the median of the return time
$\inf\{t:D_t\le0\}$ from $(\mu,D)=(0,d)$, $d>0$. Then exactly
\begin{equation}\label{sh:eq:recovery}
\varrho(d)=\varrho(1)\,\Bigl(\frac d\sigma\Bigr)^{2/3}\sigma^{2/3}\cdot
\sigma^{-2/3}=\varrho_1\Bigl(\frac d\sigma\Bigr)^{2/3},
\qquad \varrho_1:=\varrho(1)\big|_{\sigma=1}\in(0,\infty),
\end{equation}
and the return time is finite a.s.\ with tail index $\tfrac14$. Doubling
accumulated credit multiplies the recovery timescale by only
$2^{2/3}\approx1.59$: credit is far easier to undo than a linear intuition
suggests.
\end{proposition}

\begin{proof}
Finiteness a.s.\ is contained in the proof of
Theorem~\ref{sh:thm:noimm} (the set $A$ there is unbounded). The scaling
identity is Proposition~\ref{sh:prop:selfsim} with $c=d^{-1/3}$ applied to
the median, a scale-equivariant functional; positivity and finiteness of
$\varrho_1$ follow since the return time is a nondegenerate positive random
variable. The tail index is Theorem~\ref{sh:thm:persist} read for the
returning coordinate.
\end{proof}

\begin{proposition}[Horizon level sets are $\kappa$-rays; \textbf{P}]
\label{sh:prop:rays}
Every scale-invariant function on the open quadrant is constant exactly on
the curves $D\propto\mu^3/\sigma^2$; if a horizon exists in any
scale-invariant sense, it is a ray of constant $\kappa$, and $\kappa$ is
its natural transverse coordinate.
\end{proposition}

\begin{proof}
Proposition~\ref{sh:prop:gauge}(iii).
\end{proof}

\begin{definition}[Deadline-relative horizon and its width]
\label{sh:def:deadline}
Fix a time budget $\mathcal T$. The \emph{deadline horizon} is the level
set $\{t_*(\mu,D)\asymp\mathcal T\}$, located at
$\mu^*\asymp\sigma\mathcal T^{1/2}$, $D^*\asymp\sigma\mathcal T^{3/2}$.
Its \emph{softness} is the logarithmic sensitivity of survival to state.
\end{definition}

\begin{corollary}[Maximal softness; \textbf{A}]\label{sh:cor:soft}
Under Claim~\ref{sh:claim:tail},
$\dfrac{\partial\log\PP(T>\mathcal T)}{\partial\log t_*}=\dfrac14+o(1)$.
Changing survival probability by a factor $e$ therefore requires changing
$t_*$ by $e^4\approx55$ and, since $t_*\propto D^{2/3}$ in the
credit-dominated regime, changing $D$ by $e^6\approx403$. The driftless
horizon is smeared over roughly two orders of magnitude in timescale and
nearly three in credit: there is no regime in which resistance has
arbitrarily little effect, only a $\tfrac14$-power degradation --- about as
gentle a degradation as a power law permits.
\end{corollary}

\subsection{The scale-function criterion: P\'olya's theorem, correctly
translated}\label{sh:subsub:scale}

Let the driver be a regular diffusion on $\R$,
\begin{equation}\label{sh:eq:generaldriver}
d\mu_t=b(\mu_t)\,dt+\varsigma(\mu_t)\,dW_t,\qquad
s(\mu):=\int_{0}^{\mu}\exp\Bigl(-\int_{0}^{v}\frac{2b(w)}{\varsigma^2(w)}
\,dw\Bigr)dv,
\end{equation}
with $\varsigma>0$ locally bounded away from $0$ and $\infty$; $s$ is the
scale function and Feller's test (see \cite[Ch.~15]{KT81} or
\cite[Ch.~VII]{RY99}) classifies the boundary behavior.

\begin{theorem}[Trichotomy; \textbf{P}]\label{sh:thm:trichotomy}
For the model \eqref{sh:eq:model} with driver
\eqref{sh:eq:generaldriver}:
\begin{enumerate}[label=\emph{(\roman*)},itemsep=1pt]
\item If $s(+\infty)<\infty$ and $s(-\infty)=-\infty$, then
$\mu_t\to+\infty$ a.s., $\Lambda_\infty<\infty$ a.s., $\PP(\Imm)=1$, and
$\PP(T=\infty)=\E[e^{-\Lambda_\infty}]\in(0,1)$.
\item If $s(-\infty)>-\infty$ and $s(+\infty)=\infty$, then
$\mu_t\to-\infty$ a.s., $\Lambda_\infty=\infty$ a.s., and $T<\infty$ a.s.
\item If both limits are finite, then
\begin{equation}\label{sh:eq:scaleformula}
\PP(\Imm)\;=\;\frac{s(\mu_0)-s(-\infty)}{s(+\infty)-s(-\infty)},
\end{equation}
with $\Lambda_\infty<\infty$ on $\{\mu_t\to+\infty\}$ and
$\Lambda_\infty=\infty$ on the complementary event
$\{\mu_t\to-\infty\}$.
\item If both limits are infinite, the driver is recurrent and the outcome
is \emph{not} determined by $s$; see
Proposition~\ref{sh:prop:recurrent}.
\end{enumerate}
\end{theorem}

\begin{proof}
(i) $\PP(\lim_t\mu_t=+\infty)=1$ is Feller's test \cite{KT81}. On this
event choose $t_0(\omega)$ with $\mu_t\ge1$ for $t\ge t_0$; then
$D_t\ge D_{t_0}+(t-t_0)$, so
$\Lambda_\infty\le\Lambda_{t_0}+h_0e^{-(D_{t_0}-D_0)}\int_0^\infty
e^{-v}dv<\infty$ a.s. Positivity of $\PP(T=\infty)$ follows since
$e^{-\Lambda_\infty}>0$ a.s.; it is $<1$ since $\Lambda_\infty\ge
\Lambda_1>0$ a.s. (ii) is the mirror image: eventually $\mu_t\le-1$, so
$e^{-(D_t-D_0)}\to\infty$ and $\Lambda_\infty=\infty$; death then occurs in
finite time a.s.\ by Definition~\ref{sh:def:cox}. (iii) The hitting
probabilities of $\pm\infty$ are the classical scale-function formula, and
the two events partition the space up to null sets; apply the pathwise
arguments of (i) and (ii) on each. (iv) Recurrence under unbounded scale is
Feller's test again; Example~\ref{sh:ex:recurrentcounter} shows $s$ alone
cannot decide.
\end{proof}

\begin{proposition}[Recurrent drivers; \textbf{P}]\label{sh:prop:recurrent}
Suppose the driver \eqref{sh:eq:generaldriver} is positive recurrent with
stationary law $\pi$ and $\int|\mu|\,\pi(d\mu)<\infty$; write
$\bar m:=\int\mu\,\pi(d\mu)$.
\begin{enumerate}[label=\emph{(\alph*)},itemsep=1pt]
\item If $\bar m>0$ then $D_t/t\to\bar m$ a.s., $\Lambda_\infty<\infty$
a.s., and $\PP(\Imm)=1$.
\item If $\bar m<0$ then $\Lambda_\infty=\infty$ a.s.\ and $T<\infty$ a.s.
\item If $\bar m=0$ and the cycle integral defined in the proof is
nondegenerate (automatic for nondegenerate diffusions), then
$\Lambda_\infty=\infty$ a.s.: \emph{the critical case falls on the mortal
side.}
\end{enumerate}
For null-recurrent drivers no general statement is offered
\textbf{[O]}; the Brownian driver, the fundamental null-recurrent case, is
settled unconditionally by Theorem~\ref{sh:thm:noimm}.
\end{proposition}

\begin{proof}
(a), (b): the ergodic theorem for positive recurrent one-dimensional
diffusions gives $D_t/t\to\bar m$ a.s.\ (see \cite[Ch.~X]{RY99}); if
$\bar m>0$ then eventually $D_t\ge D_0+\tfrac{\bar m}2t$ and
$\Lambda_\infty<\infty$; if $\bar m<0$ the integrand
$e^{-(D_t-D_0)}\to\infty$.

(c) Fix interior levels $a<b$ and regenerate at up-crossings: let
$R_0:=\inf\{t:\mu_t=b\}$ and inductively
$S_k:=\inf\{t>R_{k-1}:\mu_t=a\}$, $R_k:=\inf\{t>S_k:\mu_t=b\}$. By the
strong Markov property the cycles over $[R_{k-1},R_k)$ are i.i.d.; positive
recurrence gives $\E[L]<\infty$ for the cycle length
$L:=R_1-R_0$, and the cycle occupation identity yields, for the cycle
integral $\xi_k:=\int_{R_{k-1}}^{R_k}\mu_s\,ds$,
\[
\E[\xi_1]=\E[L]\int\mu\,d\pi=0,\qquad
\E|\xi_1|\le\E\Bigl[\int_{\text{cycle}}|\mu_s|ds\Bigr]
=\E[L]\int|\mu|\,d\pi<\infty .
\]
Thus $S_n:=D_{R_n}-D_{R_0}=\sum_{k\le n}\xi_k$ is a nondegenerate mean-zero
i.i.d.\ walk, so $\liminf_nS_n=-\infty$ a.s.\ by the classical
oscillation dichotomy for random walks (see, e.g., \cite{Dur19});
in particular the events $B_n:=\{D_{R_n}\le-1\}$ occur infinitely often
a.s. Choose $\delta>0$ and $M<\infty$ with
\[
q:=\PP_b\Bigl(L\ge\delta,\ \sup_{0\le u\le\delta}
\Bigl|\int_0^u\mu_s\,ds\Bigr|\le M\Bigr)>0,
\]
possible since $L>0$ a.s.\ and the supremum is finite a.s. Let
$G_n$ be the corresponding event for cycle $n{+}1$; by the strong Markov
property $\PP(G_n\mid\Fc_{R_n})=q$, so
$\sum_n\PP(B_n\cap G_n\mid\Fc_{R_n})=q\sum_n\one_{B_n}=\infty$ a.s., and
L\'evy's extension of the Borel--Cantelli lemma (see, e.g.,
\cite{Dur19}) gives
$B_n\cap G_n$ infinitely often a.s. On $B_n\cap G_n$, for
$u\in[0,\delta]$ one has $D_{R_n+u}\le-1+M$, hence
\[
\int_{R_n}^{R_n+\delta}e^{-(D_s-D_0)}\,ds\ \ge\ \delta\,e^{\,1-M+D_0}\,>\,0 ,
\]
and these windows are disjoint. Summing over infinitely many of them,
$\Lambda_\infty=\infty$ a.s.
\end{proof}

\begin{example}[Recurrence alone does not decide]\label{sh:ex:recurrentcounter}
The mean-reverting Ornstein--Uhlenbeck driver
$d\mu=\theta(\bar\mu-\mu)dt+\sigma dW$ with $\theta>0$ has Gaussian-growing
$s'$, hence $s(\pm\infty)=\pm\infty$: it is (positive) recurrent for every
$\bar\mu$. Yet Proposition~\ref{sh:prop:recurrent} gives
$\PP(\Imm)=1$ for $\bar\mu>0$ and certain death for $\bar\mu\le0$. A
horizon exists here --- sharp, at $\bar\mu=0$ --- but it lives in
\emph{parameter} space, not in the initial condition, which is again pure
gauge.
\end{example}

\begin{remark}[Two warnings]\label{sh:rem:warnings}
Two traps in this circle of ideas deserve explicit flags.
\emph{First}, it is tempting to summarize \eqref{sh:eq:scaleformula} as
``recurrent driver $\Rightarrow$ certain death''; Example
\ref{sh:ex:recurrentcounter} shows this is false --- in the recurrent class
the order parameter is the stationary mean $\bar m$, with the critical case
$\bar m=0$ mortal, and only in the null-recurrent Brownian case is
symmetry itself the whole story.
\emph{Second}, under adverse drift it is tempting to read the survival tail
off the typical path, which yields doubly-exponential decay; in fact the
tail is only \emph{exponential}, dominated by an optimal ``fighting-back''
fluctuation --- Proposition~\ref{sh:prop:advdrift} below --- a
Freidlin--Wentzell phenomenon in miniature.
\end{remark}

\begin{center}
\begin{tabular}{@{}ll@{}}
\toprule
\textbf{Random walk (P\'olya)} & \textbf{Hazard dynamics}\\
\midrule
dimension $d$ & scale function $s$ of the driver\\
$d\le2$: recurrent, no escape & $s(\pm\infty)=\pm\infty$: no viability from
data\\
$d\ge3$: escape probability $>0$ & $s(\pm\infty)$ finite: viability
probability in $(0,1)$\\
$d=2$: null recurrent & Brownian driver: certain death, infinite mean
life\\
escape probability from $x$ & normalized scale function at $\mu_0$\\
\bottomrule
\end{tabular}
\end{center}

\begin{remark}[Slogan]\label{sh:rem:slogan}
\emph{Existence} of a horizon in the initial condition
$=$ boundedness of the scale function;
\emph{sharpness} of the horizon $=$ steepness of the scale function.
The Brownian driver has $s(\mu)=\mu$, unbounded and of constant slope: its
horizon degenerates not by becoming infinitely sharp but by becoming
infinitely smeared --- a linear scale function distinguishes no location.
Finally, recall from \eqref{sh:eq:twotier} that
\eqref{sh:eq:scaleformula} computes \emph{viability}; immortality proper is
$\E[e^{-\Lambda_\infty}\one_\Imm]$, strictly smaller, and this is the
unique juncture at which the gauge parameter $h_0$ re-enters
(Lemma~\ref{sh:lem:hgauge}) --- affecting values, never dichotomies.
\end{remark}

\subsection{An exactly solvable gradual horizon}\label{sh:subsub:solvable}

\begin{theorem}[Unstable Ornstein--Uhlenbeck driver; \textbf{P}]
\label{sh:thm:uou}
Let the driver be
\begin{equation}\label{sh:eq:uou}
d\mu_t=\nu\,\mu_t\,dt+\sigma\,dW_t,\qquad \nu>0 .
\end{equation}
Then $e^{-\nu t}\mu_t\to Z$ a.s.\ and in $L^2$, where
$Z\sim N\!\bigl(\mu_0,\ \sigma^2/2\nu\bigr)$, and
\begin{equation}\label{sh:eq:sigmoid}
\PP(\Imm)\;=\;\PP(Z>0)\;=\;\Phi\!\Bigl(\frac{\mu_0\sqrt{2\nu}}{\sigma}\Bigr).
\end{equation}
On $\{Z>0\}$ the hazard decays doubly exponentially,
$\log h_t\sim-(Z/\nu)e^{\nu t}$, and $\Lambda_\infty<\infty$; on
$\{Z<0\}$ the hazard explodes doubly exponentially and $T<\infty$ a.s.
\end{theorem}

\begin{proof}
Variation of constants gives
$\mu_t=e^{\nu t}\bigl(\mu_0+\sigma\int_0^te^{-\nu s}dW_s\bigr)
=:e^{\nu t}M_t$. The martingale $M$ has
$\E\langle M\rangle_\infty=\sigma^2\!\int_0^\infty e^{-2\nu s}ds
=\sigma^2/2\nu<\infty$, hence converges a.s.\ and in $L^2$ to a Gaussian
limit $Z$ with the stated mean and variance. On $\{Z>0\}$ pick
$t_0(\omega)$ with $\mu_t\ge\tfrac Z2e^{\nu t}$ for $t\ge t_0$; then
$D_t\ge D_{t_0}+\tfrac{Z}{2\nu}(e^{\nu t}-e^{\nu t_0})$ and the integrand
$e^{-(D_t-D_0)}$ is dominated by a doubly-exponentially decaying function,
so $\Lambda_\infty<\infty$. On $\{Z<0\}$ symmetrically
$\mu_t\le\tfrac Z2e^{\nu t}$ eventually, $D_t\to-\infty$ doubly
exponentially, the integrand diverges, $\Lambda_\infty=\infty$, and
$T<\infty$ a.s. Finally $\PP(Z=0)=0$, and
$\PP(Z>0)=\Phi(\mu_0/\sqrt{\sigma^2/2\nu})$.
\end{proof}

\begin{remark}[Consistency with the scale-function formula]
\label{sh:rem:crosscheck}
For \eqref{sh:eq:uou}, $s'(v)=\exp(-\nu v^2/\sigma^2)$, so
$s(\pm\infty)$ are finite and the normalized scale function at $\mu_0$ is
the mass below $\mu_0$ of the centered Gaussian with variance
$\sigma^2/2\nu$ --- precisely \eqref{sh:eq:sigmoid}. The exactly solvable
case is thus the generic case of
Theorem~\ref{sh:thm:trichotomy}(iii) in its purest form.
\end{remark}

\begin{definition}[Horizon profile, location, width]\label{sh:def:width}
For a driver with bounded scale, the \emph{horizon profile} is
$\Pi(\mu_0):=\PP(\Imm\mid\mu_0)$, the \emph{location} its median point
$\Pi=\tfrac12$, and the \emph{width} the natural spread of the transition,
e.g.\ the interquartile range of the measure $d\Pi$. For
\eqref{sh:eq:uou}: location $\mu_0=0$; width
\begin{equation}\label{sh:eq:width}
w=\frac{\sigma}{\sqrt{2\nu}},\qquad\text{sharpness}\quad
\Sigma:=w^{-1}=\frac{\sqrt{2\nu}}{\sigma}.
\end{equation}
As $\Sigma\to\infty$ (small noise or strong instability),
$\Pi\to\one\{\mu_0>0\}$: a hard event horizon. As $\Sigma\to0$,
$\Pi\to\tfrac12$ everywhere: no horizon, a fair coin.
\end{definition}

\begin{remark}[Sharpness is a small-noise limit; general principle]
\label{sh:rem:FW}
The width \eqref{sh:eq:width} is exactly the scale of fluctuation of the
terminal charge $Z$ capable of overturning the deterministic verdict
$\operatorname{sign}(\mu_0)$: the sharp horizon is the Freidlin--Wentzell
limit \cite{FW12} of the soft family. We record the principle in the form
answering the question that motivates this part: \emph{a horizon is
never intrinsically sharp; its width is the fluctuation scale of the
driver, and it appears sharp exactly when that scale is small against the
spread of initial conditions.} Operationally, comparing $w$ to the
population spread of $\mu_0$ yields three regimes: $w\ll$ spread --- sharp
horizon, the population splits into the fated; $w\sim$ spread --- genuinely
gradual horizon, initial conditions matter but do not determine;
$w\gg$ spread --- no horizon, the outcome is noise.
\end{remark}

\subsection{Model atlas and the role of driver smoothness}
\label{sh:subsub:atlas}

\begin{proposition}[Adverse linear drift: the price of survival;
\textbf{P/S}]\label{sh:prop:advdrift}
Let $\mu_t=\mu_0+\nu t+\sigma W_t$ with $\nu<0$. Then:
\begin{enumerate}[label=\emph{(\roman*)},itemsep=1pt]
\item \textbf{[P]}
$\displaystyle\liminf_{t\to\infty}\frac1t\log\PP(T>t)\ \ge\
-\frac{\nu^2}{2\sigma^2}$;
\item \textbf{[P]} there exists $c>0$ with $\PP(T>t)\le e^{-ct}$ for all
large $t$;
\item \textbf{[S]} the exact rate is
\begin{equation}\label{sh:eq:rate}
\lim_{t\to\infty}\frac1t\log\PP(T>t)\ =\ -\frac{3\nu^2}{8\sigma^2},
\end{equation}
attained by the optimal survival profile
$\mu_{tv}\approx\tfrac{|\nu|t}{4}\,v\,(2-3v)$,
$D_{tv}\approx\tfrac{|\nu|t^2}{4}\,v^2(1-v)$, $v\in[0,1]$: the conditioned
survivor front-loads credit, lets the rate go negative at $v=\tfrac23$, and
coasts so as to exhaust its credit exactly at the deadline.
\end{enumerate}
\end{proposition}

\begin{proof}
(i) By the Markov property and a finite-time event of positive probability
we may assume $\mu_0\ge1$; constants below depend on the data. Let
$a:=|\nu|/\sigma$, $b_0:=(1-\mu_0)/\sigma\le0$, and
\[
A_t:=\Bigl\{W_u\ge b_0+au\ \ \forall u\in[0,t],\ \ W_t\le b_0+at+1\Bigr\}.
\]
On $A_t$, $\mu_u\ge1$ for all $u\le t$, hence $D_u\ge D_0+u$ and
$\Lambda_t\le h_0$, so $\PP(T>t)\ge e^{-h_0}\PP(A_t)$. Girsanov with drift
$a$ (density $\exp(-a\widetilde W_t-a^2t/2)$, $\widetilde W:=W-a\,\cdot$
a Brownian motion under the tilted measure $Q$) gives
\[
\PP(A_t)\ \ge\ e^{-a^2t/2}\,e^{-a(b_0+1)}\,
Q\bigl(\widetilde W_u\ge b_0\ \forall u\le t,\ \widetilde W_t\in[b_0,b_0+1]
\bigr),
\]
and the $Q$-probability is of order $t^{-3/2}$ by the reflection
principle: polynomial corrections do not affect the rate.
(ii) With $K:=\tfrac{h_0}4t$,
$\PP(T>t)\le e^{-K}+\PP(\Lambda_t\le K)$. On $\{\Lambda_t\le K\}$,
$\Leb\{u\le t:D_u\le D_0\}\le K/(h_0e^{-D_0}\wedge h_0)\lesssim t/4$ (adjust
constants by $D_0$), so there exists
$u_*\in[t/2,\,3t/4+1]$ with $D_{u_*}>D_0$, forcing
$\sigma I_{u_*}\ge|\nu|u_*^2/2-\mu_0u_*\ge|\nu|t^2/16$ for large $t$. By
the Borell--TIS inequality for the Gaussian process $I$ on $[0,t]$ (whose
maximal variance is $t^3/3$ and whose expected supremum is
$O(t^{3/2})=o(t^2)$),
$\PP(\sup_{u\le t}I_u\ge|\nu|t^2/16\sigma)\le
\exp(-c't^4/t^3)=e^{-c't}$; combine.
(iii) Rescale $W_{tv}=\sqrt t\,B_v$, so
$I_{tv}=t^{3/2}\int_0^vB$ and
$D_{tv}=\mu_0tv+\tfrac{\nu t^2v^2}2+\sigma t^{3/2}\!\int_0^vB$. Keeping
$\Lambda_t=O(1)$ requires, at leading order, $D_{tv}\ge0$ for all
$v\in[0,1]$; with $B=\sqrt t\,\gamma$ this reads
$\int_0^v\gamma\ge\beta v^2$, $\beta:=|\nu|/2\sigma$, at
Cameron--Martin cost $\exp(-t\cdot\tfrac12\int_0^1\dot\gamma^2)$. The
variational problem
$\min\tfrac12\int_0^1\dot\gamma^2$ subject to
$\int_0^v\gamma\ge\beta v^2\ \forall v$, $\gamma(0)=0$, is solved by
relaxing to the endpoint constraint
$\int_0^1\gamma=\int_0^1(1-s)\dot\gamma(s)\,ds\ge\beta$: Cauchy--Schwarz
gives $\int\dot\gamma^2\ge3\beta^2$ with equality for
$\dot\gamma(s)=3\beta(1-s)$, i.e.\ $\gamma(v)=3\beta(v-\tfrac{v^2}2)$,
which satisfies the \emph{full} constraint,
$\int_0^v\gamma=\tfrac{\beta v^2}2(3-v)\ge\beta v^2$ on $[0,1]$ with
equality only at $v=1$. Hence the minimum is $\tfrac32\beta^2
=3\nu^2/8\sigma^2$, and unwinding the substitutions yields the profiles
stated: $\mu_{tv}=\mu_0+t(\nu v+\sigma\gamma(v))
\approx\tfrac{|\nu|t}4v(2-3v)$ and
$D_{tv}\approx t^2\bigl(\sigma\int_0^v\gamma-\tfrac{|\nu|v^2}2\bigr)
=\tfrac{|\nu|t^2}4v^2(1-v)$. Elevating this computation to a full
large-deviation proof (upper bound with matching constant, treatment of the
soft constraint) is routine but lengthy and is not carried out here.
\end{proof}

\begin{remark}\label{sh:rem:advdrift}
Three comments. \emph{(1)} The naive lower-bound strategy ``hold
$\mu\ge1$ throughout,'' used in (i), costs $\nu^2/2\sigma^2$ per unit time
and is strictly suboptimal: the true optimizer spends the final third of
the budget with $\mu<0$, spending credit rather than defending the rate ---
survival to a deadline is cheaper than survival in perpetuity, and the
discount is exactly the factor $\tfrac34$. \emph{(2)} Conditioned on the
rare event $\{T>t\}$, the path concentrates on the stated profile: once
again the appearance of purposeful resistance is manufactured by
conditioning. \emph{(3)} The analog for a mean-reverting driver with
$\bar\mu<0$ is an exponential rate given by the corresponding
Freidlin--Wentzell action; we do not compute it here.
\end{remark}

\begin{center}
\small
\begin{tabular}{@{}lllll@{}}
\toprule
\textbf{Driver} $\mu_t$ & $\PP(\Imm)$ & \textbf{tail of }$\PP(T>t)$ &
$\E[T]$ & \textbf{horizon}\\
\midrule
constant $\mu_0>0$ & $1$ & $\to e^{-h_0/\mu_0}>0$ & $\infty$ (atom) &
trivial step\\
constant $\mu_0<0$ & $0$ & $\exp\bigl(-\tfrac{h_0}{|\mu_0|}
e^{|\mu_0|t}\bigr)$ & finite & ---\\
$\mu_0+\sigma W_t$ & $0$ & $t^{-1/4}$ \textbf{[A]} & $\infty$ & none
(null rec.)\\
$\mu_0+\nu t+\sigma W_t$, $\nu>0$ & $1$ & $\to$ const $>0$ & $\infty$ &
dichotomy in $\nu$\\
$\mu_0+\nu t+\sigma W_t$, $\nu<0$ & $0$ &
$e^{-\frac{3\nu^2}{8\sigma^2}t(1+o(1))}$ \textbf{[S]} & finite & ---\\
OU, mean $\bar\mu>0$ & $1$ & $\to$ const $>0$ & $\infty$ & sharp in
$\bar\mu$\\
OU, mean $\bar\mu=0$ & $0$ & $t^{-1/2}$ \textbf{[S]} & $\infty$ & null
rec., new exponent\\
OU, mean $\bar\mu<0$ & $0$ & $e^{-ct}$ & finite & ---\\
unstable OU \eqref{sh:eq:uou} & $\Phi\bigl(\mu_0\sqrt{2\nu}/\sigma\bigr)$
& mixture & $\infty$ & \textbf{gradual, width }$\sigma/\sqrt{2\nu}$\\
general bounded scale & \eqref{sh:eq:scaleformula} & mixture & $\infty$ &
shape $=$ scale fn.\\
\bottomrule
\end{tabular}
\end{center}

\begin{remark}[Exponents measure smoothness, not recurrence class]
\label{sh:rem:smoothness}
The mean-zero mean-reverting driver is mortal
(Proposition~\ref{sh:prop:recurrent}(c)) with expected tail $t^{-1/2}$, not
$t^{-1/4}$: at large scales its credit is Brownian-like (an additive
functional obeying an invariance principle \cite{Bha82}), so the relevant
persistence exponent is that of Brownian motion, $\tfrac12$, not of its
integral, $\tfrac14$. The survival exponent is a \emph{persistence exponent
of the credit process} and is set by the smoothness of the driver rather
than by its recurrence class: smoother driver, smaller exponent, heavier
survival tail. For a fractional Brownian driver of Hurst index $H$ the
persistence exponent of the driver itself is $1-H$ \cite{Mol99}; the
exponent for its integral --- the relevant one here --- is not known
rigorously \textbf{[O]}: the numerical evidence of Molchan--Khokhlov
\cite{MK04} supports the conjecture $\theta(H)=H(1-H)$ for persistence
from the origin, which correctly returns $\tfrac14$ at $H=\tfrac12$;
see also \cite{AS15}.
\end{remark}

% ---------------------------------------------------------------------
\section{The background as a free energy}\label{app:background}

% ---------------------------------------------------------------------
\label{tc:sub:background}

The background must be positive, variable, and never near zero.
Construction \eqref{tc:eq:B} achieves all three: a sum of $n$ positive
terms is small only if \emph{all} are small, an event of probability
$\Phi(-a/\epsilon\sqrt t)^{n}$, exponentially small in $n$. This is the
mechanism we identified. What was not anticipated is that the
construction is a familiar object.

\subsection{Freezing and Gompertz}\label{tc:subsub:freezing}

\begin{remark}[The background is a partition function]
\label{tc:rem:partition}
Writing $\epsilon W^{(i)}_t=\epsilon\sqrt t\,g_i$ with $g_i$ i.i.d.\
standard normal,
\[
B_t=\sum_{i=1}^{n}e^{\beta_t g_i},\qquad \beta_t:=\epsilon\sqrt t,
\]
the partition function of $n$ i.i.d.\ Gaussian energy levels at inverse
temperature $\beta_t$. Since $\beta_t$ \emph{increases} with $t$, the
background is a disordered system that \emph{cools as it ages}: it is
Derrida's random energy model \cite{Der81,Der80} with time as inverse
temperature. The corresponding limit theory for sums of random
exponentials is \cite{BBM05}; the mathematical treatment of the REM is
\cite{Bov06}.
\end{remark}

\begin{theorem}[Freezing; \textbf{C}]\label{tc:thm:freezing}
With $\theta=\epsilon\sqrt{t/\log n}$, as $n\to\infty$ at fixed
$\theta$,
\begin{equation}\label{tc:eq:freezing}
\frac{\log B_t}{\log n}\;\longrightarrow\;
\begin{cases}
1+\theta^{2}/2, & \theta<\sqrt2\quad(\text{delocalized}),\\
\theta\sqrt2, & \theta\ge\sqrt2\quad(\text{frozen}),
\end{cases}
\end{equation}
in probability, with freezing time
$t_{\mathrm{fr}}=2\log n/\epsilon^{2}$.
\end{theorem}

\begin{proof}[Proof and status]
The two candidates are the annealed free energy
$\log\E B_t=\log n+\epsilon^{2}t/2$ and the extremal contribution
$\max_i\epsilon W^{(i)}_t\approx\epsilon\sqrt{2t\log n}$; normalized
these are $1+\theta^{2}/2$ and $\theta\sqrt2$, and
$(1+\theta^{2}/2)-\theta\sqrt2=\tfrac12(\theta-\sqrt2)^{2}\ge0$ with
equality only at $\theta=\sqrt2$ --- the annealed value dominates
everywhere and the curves meet tangentially, the REM's characteristic
contact. The second moment gives
$\E B_t^{2}/(\E B_t)^{2}=e^{\epsilon^{2}t}/n+(n-1)/n\to1$ iff
$\epsilon^{2}t<\log n$, so concentration is immediate for
$t<t_{\mathrm{fr}}/2$; the full delocalized range needs the truncated
second moment, and the frozen phase is the standard extreme-value
computation. Both are classical \cite{Der81,Bov06,BBM05}.
\end{proof}

\begin{remark}[Our two mechanisms are the two phases]
\label{tc:rem:twophases}
The brief describes the background as held up both by ``the high ones
dominating'' and by ``the fuzz of all the intermediate-sized ones adding
together.'' These are not two descriptions of one regime but exact
descriptions of the two: the frozen phase is carried by the single
largest term, the delocalized phase by the bulk.
\end{remark}

\begin{corollary}[Gompertz; \textbf{P}]\label{tc:cor:gompertz}
For $t<t_{\mathrm{fr}}$, $B_t=n\,e^{\epsilon^{2}t/2}(1+o(1))$ in
probability: the background grows exponentially in age at rate
$\epsilon^{2}/2$. This is the Gompertz law \cite{Gom25}, arriving
unbidden --- the construction was chosen for positivity, not to
reproduce it.
\end{corollary}

\subsection{The background is gauge}\label{tc:subsub:gauge}

\begin{lemma}[Channels define moving walls; \textbf{P}]
\label{tc:lem:walls}
Set $D_t:=x_t-x_0$. Each channel becomes $O(1)$ --- i.e.\ contributes
appreciably to $\Lambda$ --- at a \emph{killing wall}:
\begin{equation}\label{tc:eq:killwalls}
x^{\mathrm{crest}}_t=\log B_t+b,
\qquad
x^{\mathrm{trough}}=\frac{\log\Gamma}{\rho},
\end{equation}
the first random and slowly moving, the second a fixed constant. the driftless case's soft-wall reduction applies with its fixed wall replaced by
$\max$ of these.
\end{lemma}

\begin{proof}
Solve $B_te^{\,b-x}=1$ and $\Gamma e^{-\rho x}=1$.
\end{proof}

\begin{theorem}[Exact absorption of a linear background; \textbf{P}]
\label{tc:thm:absorb}
Suppose $\log B_u=a+cu$ exactly --- the Gompertz phase, $a=\log n$,
$c=\epsilon^{2}/2$. Then persistence above the crest wall
\eqref{tc:eq:killwalls} is \emph{identical} to the driftless case's fixed-wall
problem with initial data
\begin{equation}\label{tc:eq:shift}
\mu_0\longmapsto\mu_0-c,\qquad D_0\longmapsto D_0-a-b .
\end{equation}
By that section's gauge proposition the background therefore moves
prefactors and timescales alone, and no asymptotic exponent.
\end{theorem}

\begin{proof}
With $D_u=D_0+\mu_0u+\sigma I_u$, $I_u=\int_0^u W_s\,ds$, the constraint
$D_u\ge a+b+cu-x_0$ reads
$\sigma\bigl(I_u-\tfrac{c-\mu_0}{\sigma}u\bigr)\ge a+b-x_0-D_0$, and
$I_u-\kappa u=\int_0^u(W_s-\kappa)ds$ is the integral of a Brownian
motion started at $-\kappa$: the same Kolmogorov diffusion from shifted
data.
\end{proof}

\begin{remark}[A headwind]\label{tc:rem:headwind}
An exponentially growing background is a \emph{constant headwind} of
size $\epsilon^{2}/2$ on the improvement rate: the system must improve
at that rate merely to hold its hazard fixed, and its initial credit is
debited by $\log n+b$. Aging, here, is not a new mechanism but an
offset.
\end{remark}

\begin{theorem}[Sublinear walls are invisible; \textbf{A}]
\label{tc:thm:sublinear}
If both walls are $o(u^{3/2})$ almost surely --- which holds here, the
crest wall being $O(u)$ in the delocalized phase and $O(\sqrt u)$ in the
frozen phase, and the trough wall constant --- then, granting the
moving-boundary robustness of the driftless case's Ansatz class, the
persistence exponent is $\tfrac14$, unchanged.
\end{theorem}

\begin{proof}[Sketch and status]
After rescaling by $t^{3/2}$ the boundary converges to zero uniformly on
compacts, so the constraint converges to the fixed-wall constraint.
Converting this to an exponent identity is the moving-boundary
robustness discussed in the driftless case and surveyed in \cite{AS15}.
\end{proof}

% ---------------------------------------------------------------------
\section{The waiting room and the gain profile}\label{app:room}

% ---------------------------------------------------------------------
\label{tc:sub:room}

\subsection{Two kinds of wall}\label{tc:subsub:twowalls}

A distinction must be drawn that the informal account elides. There are
two different senses in which a channel ``has a wall'', and both are
needed.

\begin{definition}[Shape and killing walls]\label{tc:def:walls}
The \emph{shape walls} are where a channel rises above the background
floor, i.e.\ where it becomes comparable to $B_t$:
\begin{equation}\label{tc:eq:shapewalls}
x^{\mathrm{crest}}_{\mathrm{shape}}\approx b,
\qquad
x^{\mathrm{trough}}_{\mathrm{shape}}\approx-\frac{k_t}{\rho},
\qquad
k_t:=\log\frac{\E[B_t]}{\Gamma},
\end{equation}
so that the room is the \emph{asymmetric} interval $[-k_t/\rho,\,b]$,
containing the origin in its interior.
The \emph{killing walls} are where a channel makes $\Lambda$ accumulate,
i.e.\ where it becomes $O(1)$ in absolute terms; these are
\eqref{tc:eq:killwalls}. The shape walls govern what the hazard
\emph{looks like} as a function of $x$; the killing walls govern
the driftless case's persistence geometry. They are different objects and are
not interchangeable.
\end{definition}

\begin{remark}[$\Gamma$ primitive, $k$ derived]\label{tc:rem:kderived}
$\Gamma$ is a primitive of Definition~\ref{tc:def:model}; the width
$k_t$ is \emph{not} a parameter but a bookkeeping quantity recording
where $\Gamma$ sits relative to the background. Two consequences follow
immediately and are worth separating, since they run in opposite
directions:
\begin{itemize}[itemsep=1pt]
\item \emph{The room widens.} In the delocalized phase
$\E[B_t]=ne^{\epsilon^{2}t/2}$, so $k_t=\log(n/\Gamma)+\epsilon^{2}t/2$
grows linearly and the shape-wall separation $k_t/\rho$ widens at rate
$\epsilon^{2}/2\rho$: \emph{grace grows as the background rises}.
\item \emph{The killing walls separate faster.} The crest killing wall
$\log(2B_t)$ rises at $\epsilon^{2}/2$ while the trough killing wall
$\rho^{-1}\log\Gamma$ is constant, so their separation grows at
$\epsilon^{2}/2$.
\end{itemize}
Both are harmless asymptotically, being $O(t)=o(t^{3/2})$.
\end{remark}

\subsection{The void, and when it is flat}\label{tc:subsub:void}

\begin{proposition}[The room, and when it is flat; \textbf{P}]
\label{tc:prop:flat}
The crest is a sigmoid of transition width $O(1)$ centered at $x=b$:
$\varsigma_b\approx1$ for $x\lesssim b-2$ and $\varsigma_b\approx
e^{\,b-x}$ for $x\gtrsim b+2$. Hence the total hazard is flat, at
$\approx B_t$, on
\begin{equation}\label{tc:eq:flat}
-\frac{k_t}{\rho}\;\ll\;x\;\lesssim\;b-2 ,
\end{equation}
a genuine plateau when $k_t/\rho+b\gg1$. It contains $x=0$ in its
interior precisely when $b\gtrsim2$, which is the structural role of the
offset.
\end{proposition}

\begin{proof}
$\varsigma_b(x)\in(0.88,1.05)\cdot$ its plateau value for
$x\le b-2$, and the trough exceeds the crest below $-k_t/\rho$ by
Proposition~\ref{tc:prop:dropout}(i).
\end{proof}

\begin{constraint}[Internal onsets before external]\label{tc:C5}
As the amplitude grows, $x$ reaches the trough wall before the crest
wall if and only if
\begin{equation}\label{tc:eq:C5}
\boxed{\;b\;>\;\frac{k_t}{\rho}\;}
\end{equation}
--- the offset must exceed the trough's half-width. Under
\eqref{tc:eq:C5} a system of unfavorable anatomy suffers internal
damage first, and only later does its external hazard begin to fall.
\end{constraint}

\begin{remark}[Why the offset is needed]\label{tc:rem:whyoffset}
At $b=0$ the origin sits at the crest's inflection, where the local gain
is $\tfrac12$: a system with no muffling at all would already be half
effective, contradicting the reading of the void as the regime in which
the system contributes nothing. The offset moves the origin into the
plateau, where by Theorem~\ref{tc:thm:gain}(iii) the gain is $O(e^{-b})$.
The void then contains its own center, and \eqref{tc:eq:C5} orders the
two onsets.
\end{remark}

\begin{remark}[Void as mutedness]\label{tc:rem:mutedness}
The void is where the system contributes nothing. That is stated
qualitatively by Remark~\ref{tc:rem:regimes} and quantitatively by
Proposition~\ref{tc:prop:mingain}: the gain $\Psi'$, which measures how
much the driver moves the hazard, falls to
$\kappa(\rho)e^{-(\rho b+k)/(1+\rho)}$ on the void, exponentially small
in the room's width. A system in the waiting room is not merely at low hazard;
its efforts have almost no purchase, and the wider the room the less.
\end{remark}

\subsection{Activation: time supplies the interpolation}
\label{tc:subsub:activation}

\begin{proposition}[Activation time; \textbf{S}]\label{tc:prop:tact}
In log-time $s=\log t$,
\begin{equation}\label{tc:eq:logamp}
\log\Ac_{e^{s}}=\log\sigma+\tfrac32s+\log R_{e^{s}},
\qquad R\ \text{stationary},
\end{equation}
so the amplitude is exponentially small at early log-times and grows
exponentially, with a stationary fluctuation riding on a deterministic
ramp of slope $\tfrac32$. The void therefore ends, and the channels
engage, at
\begin{equation}\label{tc:eq:tact}
t^{\mathrm{int}}_{\mathrm{act}}\asymp\Bigl(\frac{k}{\rho\,\sigma}
\Bigr)^{2/3},
\qquad
t^{\mathrm{ext}}_{\mathrm{act}}\asymp\Bigl(\frac{b}{\sigma}\Bigr)^{2/3},
\end{equation}
the times at which the typical amplitude reaches the trough and crest
walls respectively. Under Constraint~\ref{tc:C5} these are ordered,
$t^{\mathrm{int}}_{\mathrm{act}}<t^{\mathrm{ext}}_{\mathrm{act}}$: the
internal channel engages first and the external one only later.
\end{proposition}

\begin{proof}[Scaling argument]
$\Ac_t=\sigma t^{3/2}R_t$ with $R_t=O_P(1)$ stationary; setting
$\sigma t^{3/2}\asymp k/\rho$ and $\sigma t^{3/2}\asymp b$ give
\eqref{tc:eq:tact}. Dimensions:
$[(k/\rho\sigma)^{2/3}]=(\mathsf T^{3/2})^{2/3}=\mathsf T$. A sharp
constant would require the first-passage law of the amplitude, which we
do not compute.
\end{proof}

\begin{remark}[Why no capacity parameter is needed]
\label{tc:rem:nocapacity}
Revision~1 installed a capacity (strength) parameter with a leak, whose
purpose was to produce a default regime in which the system contributes
nothing and $\mu$ matters only rarely. Equation \eqref{tc:eq:logamp}
supplies that regime for free: the amplitude of $x$ is exponentially
small in log-time at early times, so the void is not an assumption to be
installed but \emph{what the model already does}, and it is left at
\eqref{tc:eq:tact} --- in two stages, internal then external --- without
any additional primitive. The interpolation
parameter is $k$ --- large $k$, long void, background-dominated; small
$k$, immediate engagement --- and $k$ is itself derived
(Remark~\ref{tc:rem:kderived}).

One distinction is worth drawing here rather than left for a reader to
trip over, since \S\ref{sch:sub:schedule} reinstates a leak. The
parameter dropped above was a \emph{standing} primitive: a permanent
damping of $x$, installed to manufacture a regime the model turned out to
supply unaided. Definition~\ref{sch:def:leak}'s leak is a \emph{transient}
and carries no new primitive of that kind --- it is indexed by the
schedule's stage, it vanishes identically at completion, and the model
above is its $k=n$ boundary case. What was dropped was a permanent
assumption; what is added is a temporary one that removes itself.
\end{remark}

\begin{remark}[The cost of dropping capacity, recorded]
\label{tc:rem:cost}
This is a real loss and should not be glossed.
\begin{center}\small
\begin{tabular}{@{}lll@{}}
\toprule
& capacity (revision~1, dropped) & $k$ (retained)\\
\midrule
what it moves & an \emph{exponent}, $\tfrac14\to\tfrac12$ &
a \emph{timescale}\\
sharpness & genuine threshold at $S_c$ & crossover at
$t_{\mathrm{act}}$\\
\bottomrule
\end{tabular}
\end{center}
Revision~1's sharp-threshold theorem is therefore gone, and its ``sharp
in parameter space'' register weakens to a crossover. What replaces it
is the threshold at $\rho=1$
(Proposition~\ref{tc:prop:dropout}(ii)), beyond which the background
still reaches the survivor's hazard and below which the system closes
off from the external world. That is a threshold in a structural
parameter with a mechanical meaning rather than one installed to produce
a threshold --- but, as Remark~\ref{tc:rem:rhorole}(c) records, it moves
no exponent. The exchange is favorable and honest, not free.
\end{remark}

% ---------------------------------------------------------------------
\subsection{The exact gain}\label{tc:subsub:gain}

\begin{theorem}[Soft-kink gain; \textbf{P}]\label{tc:thm:gain}
Fix $B_t\equiv B$ constant and define the \emph{gain}
\begin{equation}\label{tc:eq:Psi}
\Psi(x):=-\log\frac{h^{\mathrm{tot}}(x)}{h^{\mathrm{tot}}(0)},
\qquad\text{so that}\qquad
h^{\mathrm{tot}}(x)=h^{\mathrm{tot}}(0)\,e^{-\Psi(x)}
\end{equation}
identically, with $h^{\mathrm{tot}}(0)=B+\Gamma$. Then:
\begin{enumerate}[label=\emph{(\roman*)},itemsep=1pt]
\item $\Psi(0)=0$ and $\Psi$ is smooth and strictly increasing;
\item $\Psi(x)=x+O(e^{-x})$ as $x\to+\infty$ and
$\Psi(x)=\rho x+O(1)$ as $x\to-\infty$: a \emph{soft kink} of slopes
$1$ and $\rho$;
\item the local gain is a hazard-weighted average of the two channels'
slopes,
\begin{equation}\label{tc:eq:Psiprime}
\Psi'(x)=\frac{h^{\mathrm{crest}}(x)\,\frac{e^{x}}{1+e^{x}}
+h^{\mathrm{trough}}(x)\,\rho}
{h^{\mathrm{crest}}(x)+h^{\mathrm{trough}}(x)}\;\in\;(0,\rho),
\end{equation}
with $\Psi'(0)\approx e^{-b}+\rho e^{-k}$, $\Psi'\to1$ as
$x\to+\infty$, $\Psi'\to\rho$ as $x\to-\infty$, and an interior minimum
on the void quantified in Proposition~\ref{tc:prop:mingain}.
\end{enumerate}
Consequently, with constant background, the model is exactly the driftless case's with the credit passed through $\Psi$: Constraint~\ref{tc:C4} is
met.
\end{theorem}

\begin{proof}
(i) is immediate from \eqref{tc:eq:Psi} and positivity of
$h^{\mathrm{tot}}$; monotonicity from (iii). (ii): as $x\to+\infty$,
$h^{\mathrm{tot}}=2Be^{-x}(1+O(e^{-(\rho-1)x}))$, so $\Psi=x+O(\cdot)$;
as $x\to-\infty$, $h^{\mathrm{tot}}=\Gamma e^{-\rho x}(1+O(e^{(\rho-1)
x}))$, so $\Psi=\rho x+O(1)$. (iii): for a sum $h=h_1+h_2$ one has
$-h'/h=(h_1s_1+h_2s_2)/(h_1+h_2)$ with $s_i=-h_i'/h_i$; here
$s_2=\rho$ and, differentiating $2B/(1+e^{x})$,
$s_1=e^{x-b}/(1+e^{x-b})\in(0,1)$. The stated limits follow by
evaluating the weights, using $\Gamma/B=e^{-k}$ at $x=0$.
\end{proof}

\begin{proposition}[The void is muted, exponentially in $k$;
\textbf{P}]\label{tc:prop:mingain}
On the void the gain \eqref{tc:eq:Psiprime} attains an interior minimum
at
\begin{equation}\label{tc:eq:xstar}
x^{*}=\frac{b-k+2\log\rho}{1+\rho},
\qquad
\min\Psi'=\kappa(\rho)\,e^{-\frac{\rho b+k}{1+\rho}},
\qquad
\kappa(\rho):=\rho^{\frac{2}{1+\rho}}+\rho^{\frac{1-\rho}{1+\rho}} .
\end{equation}
The gain therefore does not merely dip but is \emph{exponentially small
in $(\rho b+k)/(1+\rho)$}: the wider the room --- on either side --- the
more completely the driver is muted within it. For $\rho=\tfrac32$,
$\kappa=2.045$; for $\rho=\tfrac52$, $\kappa=2.443$.
\end{proposition}

\begin{proof}
On the void the crest is saturated, so its weight is $\approx1$ and its
slope is $s_{1}\approx e^{x-b}$, while the trough weight is
$\approx(\Gamma/B)e^{-\rho x}=e^{-k}e^{-\rho x}$. Hence
$\Psi'\approx f(x):=e^{x-b}+\rho e^{-k}e^{-\rho x}$, and $f'(x)=0$ gives
$e^{(1+\rho)x}=\rho^{2}e^{\,b-k}$, which is \eqref{tc:eq:xstar}.
Substituting back, both terms carry the common factor
$e^{-(\rho b+k)/(1+\rho)}$ with coefficients $\rho^{2/(1+\rho)}$ and
$\rho^{(1-\rho)/(1+\rho)}$, summing to the stated value. (The
approximation is verified numerically to four significant figures over
$\rho\in\{1.5,2.5\}$, $b\in[6,14]$, $k\in[8,20]$.)
\end{proof}

\begin{remark}[Correction to the reference specification]
\label{tc:rem:speccorr}
Two points where the informal specification requires amendment, both
recorded in the ledger.

\emph{(a)} The closed form $\psi_\rho(x)=-\log(e^{-x}+e^{-\rho x})$ is
\emph{not} an identity for \eqref{tc:eq:hazard}: at $x=0$ the model
gives $B+\Gamma$ while $\psi_\rho$ gives $2h_0$. It is the clean
representative of the soft-kink class, not the model's own gain;
\eqref{tc:eq:Psi} is exact by construction and has the same asymptotic
slopes, so nothing structural is lost.

\emph{(b)} The gain range is $(0,\rho)$ and not $[1,\rho]$. The crest
channel's local slope is $e^{x}/(1+e^{x})$, which tends to $0$, not $1$,
as $x\to-\infty$. This is not a blemish but the sharpest available
statement of what the void \emph{is}: on the plateau the gain falls to
$\kappa(\rho)e^{-(\rho b+k)/(1+\rho)}$
(Proposition~\ref{tc:prop:mingain}), so
the driver is muted and its efforts scarcely move the hazard. The gain
profile thus encodes all three regimes --- exponentially small in the
void, $\to1$ under shielding, $\to\rho$ under damage. The mutedness is a
property of a \emph{wide} room: if $\rho b+k$ is $O(1)$ the minimum is
$O(1)$ and there is no muted regime, consistently with
Proposition~\ref{tc:prop:flat}.
\end{remark}

\subsection{What the original driver was concealing}
\label{tc:subsub:concealed}

\begin{corollary}[State-dependent gain; \textbf{P}]
\label{tc:cor:concealed}
Along a path, the effective negative logarithmic derivative of the total
hazard is
\begin{equation}\label{tc:eq:effmu}
\tilde\mu_t=-\frac{d}{dt}\log h^{\mathrm{tot}}_t
=\Psi'(x_t)\,\mu_t,
\qquad \Psi'\in(0,\rho).
\end{equation}
That is: \emph{the driftless case's driver was the present driver times a
state-dependent gain}, and the old Brownian motion was absorbing that
gain as extra noise. The revision therefore extracts structure from the
original noise term rather than adding a primitive to it.
\end{corollary}

\begin{proof}
Chain rule on \eqref{tc:eq:Psi} with $\dot x=\mu$.
\end{proof}

\begin{remark}[Reading of the gain]\label{tc:rem:gainreading}
The new $\mu$ is correspondingly narrower than the old, and the
difference is now explicit rather than hidden. Where the old model
attributed all variation in the log-hazard's slope to the driver, the
present one attributes part of it to \emph{where the system is}: the
same effort buys a full unit of log-hazard reduction under shielding,
$\rho$ units of damage in a trough, and nothing at all in the void.
\end{remark}

% ---------------------------------------------------------------------
\section{Transfer, in detail}\label{app:transfer}

\subsection{Reduction to the driftless case}\label{tc:subsub:reduction}

\begin{proposition}[Reduction; \textbf{P}]\label{tc:prop:reduction}
On the slice $\epsilon=0$ (so $B\equiv n$), $\Gamma\to0$, and $x\gg0$,
the model is the driftless case's with $h_0=2n$ and credit $x-x_0$. More
usefully: for $\epsilon=0$ and any $\Gamma$, the total hazard is a
deterministic function of $x$ alone, and Theorem~\ref{tc:thm:gain}
identifies it as the driftless case's hazard composed with a fixed gain.
\end{proposition}

\begin{proof}
At $\epsilon=0$, $B\equiv n$; for $x\gg0$, $h^{\mathrm{tot}}\approx
2ne^{-x}$ since the trough is $O(e^{-\rho x})=o(e^{-x})$.
\end{proof}

% ---------------------------------------------------------------------
\subsection{Asymptotics: the driftless case survives intact}
\label{tc:sub:asymptotics}

\begin{theorem}[the driftless case is undisturbed; \textbf{A}]
\label{tc:thm:survives}
In the two-channel model, for every $\rho>0$, $\Gamma>0$ and $b\ge0$:
\begin{enumerate}[label=\emph{(\roman*)},itemsep=1pt]
\item the persistence exponent is $\tfrac14$, exactly;
\item $\PP(T>t)\asymp Ct^{-1/4}$;
\item $\E[T]=\infty$ while $T<\infty$ a.s.\ (null recurrence);
\item the survivorship interpolation of the driftless case, its pure Doob
limit, its $s/4t$ linear response and its critical mode, are unchanged.
\end{enumerate}
\emph{No truncation window and no $\gamma=0$ caveat}:
Constraint~\ref{tc:C3} holds as a theorem.
\end{theorem}

\begin{proof}
By Lemma~\ref{tc:lem:walls} the killing walls are $\log(2B_t)$ and
$\rho^{-1}\log\Gamma$. The latter is constant, hence an $O(1)$ shift,
which we proved to be gauge. The former is $O(t)$ in the
delocalized phase and $O(\sqrt t)$ in the frozen phase, hence
$o(t^{3/2})$, and is absorbed exactly in the linear case by
Theorem~\ref{tc:thm:absorb} and invisibly in general by
Theorem~\ref{tc:thm:sublinear}; the offset $b$ enters
\eqref{tc:eq:shift} as a further $O(1)$ shift and is likewise gauge.
Survivors live at large positive $x$, where by
Proposition~\ref{tc:prop:dropout}(ii) the dominant channel is the crest
if $\rho>1$ and the trough if $\rho<1$; in either case the dominant
channel is an exponential in $x$, so the killing wall is a level of the
credit and the persistence problem is the driftless case's up to a constant
rescaling. The conclusions are then the driftless case's, transported --- and,
notably, for every $\rho>0$ rather than only $\rho>1$.
\end{proof}

\begin{remark}[This repairs revision~1's defect]\label{tc:rem:repair}
Revision~1's additive internal hazard contributed $\gamma t$ to
$\Lambda$, destroying the $t^{-1/4}$ tail and forcing a truncation
window with finite life expectancy $\asymp g^{-3/4}$. The trough channel
does not, and the reason is structural rather than fortunate:
$\Gamma e^{-\rho x}$ is small exactly where survivors live, whereas a
constant is not small anywhere. Both walls are $O(1)$-or-gauge shifts,
and we proved such shifts to be gauge.
\end{remark}

\begin{remark}[What was given up]\label{tc:rem:givenup}
Both channels are decreasing in $x$, so the total hazard is monotone and
\emph{more strength is always better}. Revision~1's U-shaped hazard, and
with it the optimal finite strength and the $(S,A)$ optimization, are
gone. This is the price of Constraint~\ref{tc:C3}: an interior optimum
requires a term that grows with $x$, and any such term destroys the
tail. An optimum and an undisturbed the driftless case were never going to
coexist. What survives of the earlier optimization story is the
threshold at $\rho=1$ and the trade recorded in
Remark~\ref{tc:rem:cost}.
\end{remark}

Theorem~\ref{tc:thm:survives} asserted that the driftless case's conclusions
carry over. This subsection makes the transfer quantitative: it
enumerates the four ways the two-channel hazard differs from
$h_{0}e^{-x}$, shows that each lands in gauge or in the prefactor, and
isolates the one place where the two-channel model is genuinely
\emph{better} conditioned than the driftless case.

\subsection{The four differences}\label{tc:subsub:four}

\begin{center}\small
\begin{tabular}{@{}llll@{}}
\toprule
difference & size & where it goes & reference\\
\midrule
wall location $\log B_u+b$ & $\Theta(u)$ & gauge (exact absorption) &
Thm.~\ref{tc:thm:absorb}\\
graded boundary of width $W_u$ & $\Theta(u)$ & gauge (collapses under
rescaling) & Thm.~\ref{tc:thm:collapse}\\
waiting-room passage & $O(1)$ & prefactor & Prop.~\ref{tc:prop:roomcost}\\
background fluctuation & $O(\sqrt u)$ & gauge &
Thm.~\ref{tc:thm:sublinear}\\
\bottomrule
\end{tabular}
\end{center}

\noindent Nothing reaches an exponent. The transfer is therefore exact
at the level at which the driftless case states its results.

\begin{lemma}[Above the room the model is the driftless case's; \textbf{P}]
\label{tc:lem:aboveroom}
For $x\ge b$,
\begin{equation}\label{tc:eq:sandwich}
\tfrac12\bigl(1+e^{-b}\bigr)e^{\,b-x}\;\le\;\varsigma_b(x)\;\le\;
\bigl(1+e^{-b}\bigr)e^{\,b-x},
\end{equation}
so $h^{\mathrm{crest}}\asymp B_ue^{\,b-x}$ with constants in
$[\tfrac12,1]\cdot(1+e^{-b})$. If moreover $\rho>1$ then
$h^{\mathrm{trough}}/h^{\mathrm{crest}}\to0$ as $x\to\infty$, so
$h^{\mathrm{tot}}\asymp B_ue^{\,b-x}$: above the room the hazard is
the driftless case's $h_{0}e^{-x}$ with $h_{0}=B_ue^{b}$, up to bounded
multiplicative constants.
\end{lemma}

\begin{proof}
For $x\ge b$ one has $1\le e^{x-b}$, hence
$e^{x-b}\le1+e^{x-b}\le2e^{x-b}$; invert. The trough comparison is
Proposition~\ref{tc:prop:dropout}(ii).
\end{proof}

\subsection{The graded boundary collapses}\label{tc:subsub:collapse}

The substantive structural difference is that the killing boundary is
neither hard nor soft but \emph{graded}. Above $\log B_u+b$ the hazard
is negligible; across a band below it the hazard is $\asymp B_u$,
bounded; below $(\log\Gamma)/\rho$ it diverges.

\begin{theorem}[The band is asymptotically invisible; \textbf{A}]
\label{tc:thm:collapse}
Let
\begin{equation}\label{tc:eq:band}
W_u:=\underbrace{\log B_u+b}_{\text{crest wall}}
-\underbrace{\frac{\log\Gamma}{\rho}}_{\text{trough wall}}
=\Bigl(1-\frac1\rho\Bigr)\log B_u+b+\frac{k}{\rho}
\end{equation}
be the width of the graded band. In the delocalized phase
$W_u=\Theta(u)$ and in the frozen phase $W_u=\Theta(\sqrt u)$; in both,
\[
\frac{W_u}{u^{3/2}}\longrightarrow0 .
\]
Hence under the driftless case's rescaling the whole band contracts to a single
point, and the graded boundary is asymptotically the hard wall of that
section. Granting the moving-boundary robustness of its Ansatz class,
the persistence exponent is $\tfrac14$.
\end{theorem}

\begin{proof}
Insert $\log B_u=\log n+\epsilon^{2}u/2$ (delocalized,
Corollary~\ref{tc:cor:gompertz}) or
$\log B_u=\epsilon\sqrt{2u\log n}$ (frozen,
Theorem~\ref{tc:thm:freezing}) into \eqref{tc:eq:band}; both are
$o(u^{3/2})$. The rescaling and the robustness step are as in
Theorem~\ref{tc:thm:sublinear}.
\end{proof}

\begin{remark}[Softening and re-hardening do not cancel each other ---
they both cancel against self-similarity]\label{tc:rem:netcancel}
It is tempting to read the situation as a competition: the logistic
bounds the hazard from above and would soften the wall, while the trough
is unbounded and re-hardens it, the two effects offsetting. That is not
what happens. The two effects act at \emph{different places} --- the
logistic at the top of the band, the trough at its bottom --- and
neither undoes the other. What renders both irrelevant is the third
thing: the band they enclose has width $o(u^{3/2})$, while the credit
they constrain fluctuates on scale $u^{3/2}$. A boundary structure of
any shape whatsoever, confined to a band of width $o(u^{3/2})$, is
invisible to the exponent. The net effect of the logistic and the trough
is therefore exactly nil asymptotically --- not by mutual cancellation
but because self-similarity outruns them both.
\end{remark}

\begin{proposition}[The tolerance shrinks: this model is better
conditioned than the driftless case's; \textbf{P}]\label{tc:prop:tolerance}
Survival to time $t$ requires
\begin{equation}\label{tc:eq:tolerance}
\mathrm{Leb}\{u\le t:x_u\lesssim b\}\;=\;O\!\bigl(B_t^{-1}\bigr),
\end{equation}
since the hazard on the band is $\asymp B_u$ and $\Lambda_t=O(1)$ is
required. In the delocalized phase $B_t^{-1}=n^{-1}e^{-\epsilon^{2}t/2}
\to0$, so the admissible occupation below the wall tends to zero: the
problem converges to \emph{strict} persistence.
\end{proposition}

\begin{proof}
$\Lambda_t\ge\bigl(\inf_{u\le t}B_u\bigr)\cdot c\cdot
\mathrm{Leb}\{u\le t:x_u\lesssim b\}$ with $c=\varsigma_b(b)>0$;
survival to $t$ with probability bounded below forces $\Lambda_t=O(1)$.
\end{proof}

\begin{remark}[The Ansatz gap narrows]\label{tc:rem:gapnarrows}
the driftless case's sole load-bearing heuristic was that persistence with
\emph{bounded occupation tolerance} shares the exponent of strict
persistence. The two-channel model does not evade that assumption, but
it weakens what is being assumed: by
Proposition~\ref{tc:prop:tolerance} the tolerance here is not $O(1)$ but
$O(B_t^{-1})\to0$. The richer, more ontologically committed model is
thus \emph{closer} to the rigorous case than the simpler one --- an
unusual direction for added structure to move a proof, and worth
recording as such.
\end{remark}

\subsection{The waiting room is a prefactor, and it is not a
shelter}\label{tc:subsub:roomcost}

\begin{proposition}[Cost of the room passage; \textbf{S}]
\label{tc:prop:roomcost}
The room $[-k/\rho,b]$ lies entirely below the crest killing wall
$\log B_u+b$, since $\log B_u>0$. A trajectory therefore traverses it
only before escaping upward, at time
$t^{\mathrm{ext}}_{\mathrm{act}}\asymp(b/\sigma)^{2/3}$
(Proposition~\ref{tc:prop:tact}), accumulating
\begin{equation}\label{tc:eq:roomcost}
\Lambda_{\mathrm{room}}\;\asymp\;
B_0\,t^{\mathrm{ext}}_{\mathrm{act}}
\;\asymp\;n\Bigl(\frac{b}{\sigma}\Bigr)^{2/3},
\end{equation}
an $O(1)$ quantity in $t$. Its entire effect is to multiply the
prefactor in $\PP(T>t)\asymp Ct^{-1/4}$ by
$e^{-\Theta\left(n(b/\sigma)^{2/3}\right)}$.
\end{proposition}

\begin{remark}[Grace, but not shelter]\label{tc:rem:notshelter}
The room deserves a more careful reading than ``grace interval''
suggests. Inside it the hazard is \emph{exactly the background}: neither
raised by damage nor lowered by shielding. It is therefore grace
relative to the trough --- unfavorable anatomy costs nothing while the
amplitude is small enough to keep $x$ above $-k/\rho$ --- but it is a
\emph{cost} relative to the crest, because survivors must be above the
wall and time spent in the room is time not spent shielding. Formula
\eqref{tc:eq:roomcost} prices that cost. The wider the room and the
larger the background, the more expensive the wait: the offset $b$ buys
the ordering of the two onsets (Constraint~\ref{tc:C5}) and pays for it
in the prefactor, at rate $n(b/\sigma)^{2/3}$ in the exponent.
\end{remark}

\subsection{The conditioning ladder}\label{sh:subsub:ladder}

We quantify how each conditioning enhances survival.

\begin{proposition}[Degrees of the initial data; \textbf{C/A}]
\label{sh:prop:degrees}
\emph{(i) Initial rate.} From $(m,0)$ with $m>0$: the rate stays positive
until $W$ reaches $-m/\sigma$, a time of order $t_\mu=m^2/\sigma^2$ during
which $D$ increases; thereafter, exactly by scaling,
$\PP_{(m,0)}(\tau>t)\sim c_1(t_\mu/t)^{1/4}\propto m^{1/2}$
\emph{(degree $\tfrac12$: a tenfold gain in survival odds costs a
hundredfold increase in $m$)}.
\emph{(ii) Initial credit.} From $(0,d)$ with $d>0$: by scaling with
$c=d^{-1/3}$, $\PP_{(0,d)}(\tau>t)=\PP_{(0,1)}(\sigma^{2/3}\tau>
t\,d^{-2/3}\sigma^{2/3})\asymp(t_D/t)^{1/4}\propto d^{1/6}$
\emph{(degree $\tfrac16$)}.
\emph{(iii) Unification.} All initial data act only through the equivalent
timescale they generate:
\begin{equation}\label{sh:eq:tstar}
\PP(T>t)\ \asymp\ \Bigl(\frac{t_*}{t}\Bigr)^{1/4},\qquad
t_*\ \asymp\ \sigma^{-2/3}\max\bigl(1,\ \hat\mu_0^{\,2},\ \hat
D_0^{\,2/3}\bigr),
\end{equation}
and which term dominates is decided by the horizon coordinate: by
\eqref{sh:eq:kappa}, $t_D/t_\mu=\kappa_0^{2/3}$, so rate dominates for
$\kappa_0\ll1$ and credit for $\kappa_0\gg1$.
\end{proposition}

\begin{proof}
The scaling identities are exact (Proposition~\ref{sh:prop:selfsim});
transfer to $\PP(T>t)$ is Ansatz~\ref{sh:ansatz}. That the exponent from
interior starting states is unchanged is part of
Theorem~\ref{sh:thm:persist}.
\end{proof}

\begin{proposition}[Homogeneity of the harmonic function; \textbf{P/C}]
\label{sh:prop:harmonic}
On the open quadrant define
\[
\hm(\mu,d)\ :=\ \lim_{t\to\infty}t^{1/4}\,\PP_{(\mu,d)}(\tau>t);
\]
the limit exists, is finite and positive \cite{IW94,GJW99}. Then $\hm$ is
invariant-harmonic for the killed Kolmogorov diffusion, and homogeneous
of degree $\tfrac12$ under the scaling of
Proposition~\ref{sh:prop:selfsim}: $\hm(c\mu,c^3d)=c^{1/2}\hm(\mu,d)$.
Consequently
\begin{equation}\label{sh:eq:harm}
\hm(\mu,d)=\mu^{1/2}F\bigl(\kappa\bigr)=d^{1/6}\tilde F(\kappa),
\qquad \kappa=\sigma^2d/\mu^3,
\end{equation}
for profiles $F,\tilde F$ determined by the explicit
confluent-hypergeometric formulas constructed in \cite{GJW99}. (We flag
a provenance point: the special functions arising here are confluent
hypergeometric, not Airy; Airy functions belong to the neighboring
parabolic-drift literature.)
\end{proposition}

\begin{proof}
Given existence of the limit, homogeneity is immediate from
Proposition~\ref{sh:prop:selfsim}:
\[
t^{1/4}\,\PP_{(c\mu,c^3d)}(\tau>t)
= c^{1/2}\,\bigl(t/c^2\bigr)^{1/4}\,\PP_{(\mu,d)}\bigl(\tau>t/c^{2}\bigr)
\ \longrightarrow\ c^{1/2}\,\hm(\mu,d).
\]
Harmonicity follows from the Markov property and the same limit taken one
step later. Representation \eqref{sh:eq:harm} is homogeneity restricted to
the orbit coordinates of Proposition~\ref{sh:prop:gauge}(iii).
\end{proof}

\begin{remark}[Conditioning on survival forever: the Doob transform;
\textbf{C}]\label{sh:rem:doob}
Conditioning on $\{T=\infty\}$ is vacuous (Theorem~\ref{sh:thm:noimm});
its canonical regularization is the $\hm$-transform, the
$t\to\infty$ limit of conditioning on $\{\tau>t\}$. Under the transformed
measure the process never returns to zero credit and, by self-similarity,
\[
\mu_s\ \asymp\ \sigma s^{1/2},\qquad D_s\ \asymp\ \sigma s^{3/2},
\qquad\text{hence}\qquad \log h_s\ \sim\ -\sigma s^{3/2}:
\]
conditioned survivors exhibit superexponentially decaying hazard and
appear, from the inside, unambiguously immortal --- although the
unconditioned process is certainly mortal and the model contains no drift.
The drift is manufactured by selection; this is the strongest true form of
the event-horizon intuition, and it is a statement about ensembles, not
about state space. The accompanying Yaglom-type statement --- convergence
in law of $(\mu_t/\sigma t^{1/2},\,D_t/\sigma t^{3/2})$ conditionally on
$\{\tau>t\}$ to a nondegenerate limit --- is the correct
``stationary profile of a long-lived system,'' stationary only in
self-similar coordinates; see \cite{GJW99,AS15}, and \cite{MV12} for
the general theory of quasi-stationary limits against which this
self-similar variant should be read.
Section~\ref{sh:sub:interp} promotes this remark to the section's central
result: the conditioning that manufactures the drift becomes itself the
object of study, and the transform described here is identified as the
exact limit of the survivorship interpolation.
\end{remark}

% ---------------------------------------------------------------------
\section{Reference sheet: the filter at a glance}\label{app:sheet}

% ---------------------------------------------------------------------
\label{sh:sub:sheet}

This subsection collects, without proof, the quantitative content of the
section in the form most useful for reference. Everything here is
established above; cross-references point to the source.

\subsection{The conditioned dynamics in closed form}
\label{sh:subsub:sheet:dyn}

Under the survivorship limit $Q$ of Theorem~\ref{sh:thm:Q} the system is
\begin{equation}\label{sh:eq:sheet}
d\mu_t=\frac{\sigma^2}{\mu_t}\,G(\kappa_t)\,dt+\sigma\,dW^Q_t,
\qquad
dD_t=\mu_t\,dt,
\qquad
\kappa=\frac{\sigma^2D}{\mu^3},
\end{equation}
with $G(\kappa)=\tfrac12-3\kappa F'(\kappa)/F(\kappa)$ read off
$\hm=\mu^{1/2}F(\kappa)$ as in \eqref{sh:eq:harm}--\eqref{sh:eq:G}.
Dimensionally $[\sigma^2/\mu]=\mathsf T^{-2}=[\mu]/[t]$, as required. The
shape function interpolates between two endpoints
(Theorem~\ref{sh:thm:drift}):
\[
G(0^+)=\tfrac12\ \ (\text{drift }\sigma^2/2\mu,\ \text{maximal}),
\qquad
G(\infty)=0\ \ (\text{drift }0,\ \text{none}),
\]
so that \emph{survivorship pushes hardest near the wall and not at all
once credit is abundant}, the strength along any ray of fixed $\kappa$
scaling as $1/\mu$.

\begin{remark}[The endpoints are ray limits, not coordinate limits]
\label{sh:rem:raylimit}
One may not simply let $\mu\to0$ in $\sigma^2/2\mu$: at fixed $D>0$ that
sends $\kappa\to\infty$ and hence $G\to0$, and the two factors compete.
Only the statement along rays of constant $\kappa$ is unambiguous, which
is precisely why $\kappa$ and not $(\mu,D)$ separately is the coordinate
in which the conditioned dynamics should be read
(Proposition~\ref{sh:prop:rays}).
\end{remark}

\subsection{The three levels, before and after}
\label{sh:subsub:sheet:levels}

\begin{center}\small
\begin{tabular}{@{}p{0.24\textwidth}p{0.34\textwidth}p{0.34\textwidth}@{}}
\toprule
& \textbf{Before} ($\PP$) & \textbf{After} ($Q$)\\
\midrule
\multicolumn{3}{@{}l}{\emph{Level 1: the driver $\mu$}}\\
law & Gaussian, driftless & non-Gaussian, drift
$(\sigma^2/\mu)G(\kappa)\ge0$\\
autonomy & $\mu$ Markov by itself & only the pair $(\mu,D)$ is Markov\\
quadratic variation & $\sigma^2\,dt$ & $\sigma^2\,dt$ \emph{(unchanged)}\\
increments & stationary, independent & neither\\
self-similarity & index $3/2$ & index $3/2$ \emph{(preserved)}\\
\addlinespace
\multicolumn{3}{@{}l}{\emph{Level 2: the hazard $h$}}\\
path regularity & $C^1$ & $C^1$ \emph{(unchanged)}\\
law at time $s$ & lognormal, $\log$-variance
$\sigma^2s^3/3$ & non-Gaussian, one-sided\\
range & exceeds $h_0$ infinitely often,
unbounded above & $h_s<h_0$ for \emph{all} $s$ (if $D_0=0$)\\
credit $D$ & null recurrent,
$\liminf_sD_s=-\infty$ & transient, $D_s\to\infty$,
$D_s\asymp\sigma s^{3/2}$\\
\addlinespace
\multicolumn{3}{@{}l}{\emph{Level 3: survival}}\\
$\Lambda_\infty$ & $=\infty$ a.s. & $<\infty$ a.s.,
$\asymp\hat h_0=h_0\sigma^{-2/3}$\\
quenched curve $e^{-\Lambda_t}$ & $\to0$; collapse
$\exp(-e^{\sigma t^{3/2}})$ & plateau at $e^{-\Lambda_\infty}>0$\\
population curve & $\asymp Ct^{-1/4}$ & $\equiv1$\\
lifetime & $T<\infty$ a.s., $\E[T]=\infty$ & $T=\infty$\\
\bottomrule
\end{tabular}
\end{center}

\noindent Sources: Level 1, Theorems~\ref{sh:thm:Q}
and~\ref{sh:thm:drift}; Level 2, Definition~\ref{sh:def:model} and
Theorem~\ref{sh:thm:singular}; Level 3, Theorem~\ref{sh:thm:noimm},
Claim~\ref{sh:claim:tail}, Corollary~\ref{sh:cor:mean}.

\begin{remark}[Notable structures at Level 2]\label{sh:rem:sheet:L2}
Three features are worth isolating. \emph{(i)} The hazard is one
derivative \emph{smoother} than the noise driving it, a direct
consequence of the log-derivative placement
(Remark~\ref{sh:rem:iwp}); conditioning does not change this, since
Girsanov alters drift and never path regularity. \emph{(ii)} Before
conditioning $\log h_s$ is exactly Gaussian with variance
$\sigma^2s^3/3$ --- \emph{cubic} in time, the integrated-Brownian
fingerprint --- so $h_s$ is lognormal. \emph{(iii)} After conditioning,
the survivor's hazard never once exceeds its initial value, whereas
before it does so infinitely often and by unbounded amounts: this is the
wall of \S\ref{sh:subsub:reduction} in its most concrete form. Finer
still, Groeneboom, Jongbloed and Wellner \cite{GJW99} identify
$t\mapsto t^{9/10}$ as a critical curve for the conditioned process, the
expected time spent below $t^\alpha$ being finite for $\alpha<9/10$ and
infinite for $\alpha\ge9/10$: survivors do lag their typical
$s^{3/2}$ growth, and $9/10$ is exactly how much.
\end{remark}

\begin{proposition}[The quenched--annealed gap; \textbf{P}]
\label{sh:prop:gap}
In the driftless model, almost surely
\begin{equation}\label{sh:eq:gap}
\log\Lambda_t=\sup_{u\le t}\bigl(D_0-D_u\bigr)+O(\log t),
\end{equation}
and $t^{-3/2}\sup_{u\le t}(D_0-D_u)\to\sigma\sup_{v\le1}(-I_v)$ in law, a
nondegenerate positive limit. Hence the quenched survival curve
collapses \emph{doubly exponentially}: writing
$c_t:=t^{-3/2}\log\Lambda_t$,
\[
e^{-\Lambda_t}=\exp\bigl(-e^{\,c_t\,t^{3/2}}\bigr),
\qquad
c_t\ \xrightarrow{\;d\;}\ \sigma\sup_{v\le1}(-I_v)\;>\;0,
\]
while by Claim~\ref{sh:claim:tail} the population curve decays only as
$Ct^{-1/4}$. The entire population tail is therefore carried by a set of
paths of vanishing probability.
\end{proposition}

\begin{proof}
Write $M_t=\sup_{u\le t}(D_0-D_u)$ and
$\Lambda_t=h_0\int_0^te^{D_0-D_u}du$. The upper bound
$\Lambda_t\le h_0te^{M_t}$ is immediate. For the lower bound, $D$ is
$C^1$ with $\dot D=\mu$ and $\sup_{u\le t}|\mu_u|=O(t^{1/2})$ a.s.\ up to
logarithmic factors, so on an interval of length $\asymp t^{-1/2}$ around
a near-maximizer the integrand exceeds $e^{M_t-1}$, giving
$\Lambda_t\gtrsim h_0t^{-1/2}e^{M_t}$. Taking logarithms yields
\eqref{sh:eq:gap}. The scaling limit is
Proposition~\ref{sh:prop:selfsim} applied to $I$, with
$\sup_{v\le1}(-I_v)>0$ a.s.
\end{proof}

\begin{remark}[The interpolation as a one-parameter family]
\label{sh:rem:sheet:family}
Between the two columns, the survival curve of a $t$-survivor is
\[
\PP(T>t+s\mid T>t)\;\approx\;\Bigl(1+\frac st\Bigr)^{-1/4},
\]
a scale-free family flattening from a power law at $t=0$ to the constant
$1$ as $t\to\infty$; equivalently the $t$-survivor carries effective
hazard $1/4t$ (Corollaries~\ref{sh:cor:residual}
and~\ref{sh:cor:lindy}). The filter does not deform the survival curve
arbitrarily: it slides it along a single one-parameter orbit.
\end{remark}

\subsection{What conditioning can and cannot do}
\label{sh:subsub:sheet:rigid}

\begin{proposition}[Rigidity of the survivorship transform; \textbf{P}]
\label{sh:prop:rigid}
Let $\widetilde\PP$ be any measure locally absolutely continuous with
respect to $\PP$ on each $\Fc_s$ --- in particular $Q$. Then:
\begin{enumerate}[label=\emph{(\roman*)},itemsep=1pt]
\item \emph{The kinematics survive exactly.} $dD_t=\mu_t\,dt$ holds
$\widetilde\PP$-a.s.\ unchanged: credit remains literally the integral of
rate. Conditioning cannot alter it.
\item \emph{Only the driver may be tilted.} Since the noise enters solely
the $\mu$ coordinate in \eqref{sh:eq:kolm}, Girsanov can add drift only
to $\mu$; the transform has exactly one functional degree of freedom.
\item \emph{Self-similarity fixes its form.} If in addition
$\widetilde\PP$ respects the scaling of
Proposition~\ref{sh:prop:selfsim}, that degree of freedom is forced into
the shape $(\sigma^2/\mu)G(\kappa)$ of \eqref{sh:eq:sheet}, a single
function of one dimensionless variable.
\item \emph{The cascade is destroyed.} Under $\PP$ the system is
triangular: $\mu$ is autonomous and $D$ is slaved to it. Under $Q$ the
credit re-enters the drift of $\mu$ through $\kappa$, so $\mu$ alone is
no longer Markov although the pair remains so.
\end{enumerate}
\end{proposition}

\begin{proof}
(i) and (ii) are Girsanov's theorem for \eqref{sh:eq:kolm}, whose
diffusion matrix is degenerate with range spanned by $\partial_\mu$; the
finite-variation coordinate admits no equivalent change. (iii) is
Proposition~\ref{sh:prop:harmonic} together with
Theorem~\ref{sh:thm:drift}. (iv) is the $\kappa$-dependence in
\eqref{sh:eq:G}, non-constant by Theorem~\ref{sh:thm:drift}(ii)--(iii).
\end{proof}

\begin{remark}[Survivorship coupling is collider bias in continuous time]
\label{sh:rem:collider}
The feedback of Proposition~\ref{sh:prop:rigid}(iv) is informational, not
mechanical: nothing pushes $\mu$. Under $\PP$ the credit accumulated by
time $s$ and the future increments of the driver are independent, but
both feed a common outcome, survival; conditioning on that outcome
renders them dependent. This is exactly collider bias --- Berkson's
paradox, or ``explaining away'' --- realized in continuous time, and
\eqref{sh:eq:sheet} is its quantitative form: little credit banked
implies the remaining path must work harder, and the drift says by how
much. The same object appears in large-deviation theory as the driven or
effective process associated with a rare event. The moral of
Proposition~\ref{sh:prop:rigid} is then sharp: survivorship cannot bend
the kinematics of a system, only the statistics of what drives it ---
and, when the system is self-similar, only along a one-parameter family
of shapes.
\end{remark}

% ---------------------------------------------------------------------
\section{Specification of numerical tests}\label{app:numerics}

% ---------------------------------------------------------------------
\label{sh:sub:numerics}

The one load-bearing heuristic, Ansatz~\ref{sh:ansatz}, admits a direct
test, which we specify for completeness; this subsection may be read as an
appendix.

\begin{lemma}[Exact simulation of the pair; \textbf{P}]\label{sh:lem:exact}
For a step $\Delta>0$ let $\xi:=W_{t+\Delta}-W_t$ and
$\eta:=\int_t^{t+\Delta}(W_s-W_t)\,ds$. Then $(\xi,\eta)$ is independent of
$\Fc_t$, centered Gaussian with covariance
$\bigl(\begin{smallmatrix}\Delta&\Delta^2/2\\
\Delta^2/2&\Delta^3/3\end{smallmatrix}\bigr)$, and the update
\[
\mu_{t+\Delta}=\mu_t+\sigma\xi,\qquad
D_{t+\Delta}=D_t+\mu_t\Delta+\sigma\eta
\]
reproduces the law of \eqref{sh:eq:kolm} at the grid points with zero
discretization error.
\end{lemma}

\begin{proof}
$(\xi,\eta)$ is a measurable functional of the post-$t$ increments and is a
copy of $(W_\Delta,I_\Delta)$; apply Lemma~\ref{sh:lem:secondorder}. The
update identity is
$I_{t+\Delta}=I_t+W_t\Delta+\eta$ multiplied through by $\sigma$.
\end{proof}

The remaining error is the quadrature of
$\Lambda_t=h_0e^{-D_0}\int_0^te^{-(D_u-D_0)}du$, whose integrand varies on
the scale $|\Delta D|\lesssim1$, i.e.\ $\Delta\lesssim1/|\mu|\sim
t^{-1/2}$: the constraint tightens with time. Accumulate in logarithmic
space; the integrand ranges over hundreds of orders of magnitude. To reach
the tail --- probabilities of $10^{-3}$ to $10^{-4}$ require
$t/t_*$ up to $10^{12}$ and beyond --- use multilevel splitting along the
thresholds $D_k=\sigma(\theta^kt_*)^{3/2}$, or importance sampling tilted
by the harmonic function of Proposition~\ref{sh:prop:harmonic}, the
asymptotically optimal tilt. Validate the code first on the hard-wall
functional, where \eqref{sh:eq:persist} is a theorem; then run the
soft-killing functional and fit $\log\PP$ against $\log t$ over at least
three decades, discarding $t\lesssim10\,t_*$.
\emph{Falsifiers:} a fitted slope robustly different from $-\tfrac14$
(Ansatz fails: the killing term carries its own scale); any dependence of
the slope on $h_0$ (gauge violation, contradicting
Proposition~\ref{sh:prop:gauge}); non-convergence of the conditioned
rescaled credit $D_s/\sigma s^{3/2}$ (failure of the self-similar Yaglom
picture of Remark~\ref{sh:rem:doob}).

\emph{Results (executed August 2026; \cite{Ben26Extract}, \S6).} The
falsifiers did not fire. With exact pair simulation,
$N=1.5\times10^{5}$ paths, and a hybrid uniform--geometric grid to
$t=3\times10^{4}$ (control run at half the step sizes agreeing within
intervals): hard-wall slope $-0.2499$ [$-0.254,-0.243$] against the
theorem's $-\tfrac14$; soft-killing slopes $-0.2487$, $-0.2480$,
$-0.2483$ at $h_0\in\{0.3,1,3\}$, with the $h_0=30$ transient
relaxing toward $-\tfrac14$ over late windows exactly as
Lemma~\ref{sh:lem:hgauge} predicts; the ratio
$\PP(T>t)/\PP(\tau>t)$ constant to three decimals over two fitted
decades ($[0.1768,0.1788]$ at $h_0=1$) --- sharper than the Ansatz
asserts; the kinematic factor matching $(1-s/t)^{-1/4}$ at grid
spacing ($0.2874\pm0.0133$ measured, $0.2859$ exact); Yaglom
rescaled-credit distributions coinciding to KS $0.010$ between
$t=2\times10^{3}$ and $3\times10^{4}$; and the survivorship drift
positive at the predicted order for small $\kappa$ and vanishing at
large $\kappa$. The Ansatz is accordingly \textbf{[A]} with direct
numerical support; the deep-tail computation at
$t/t_*\gtrsim10^{12}$ via splitting or harmonic-function importance
sampling remains open. The survivorship results of
\S\ref{sh:sub:interp} add three sharper tests: estimate the window bias
$\beta_t(s)$ under splitting and regress it against $\hm(X_s)$, whose
proportionality tests Theorem~\ref{sh:thm:Q}; regress conditioned
increments $\Delta\mu$ on $1/\mu$ at small $\kappa$, whose slope
$\sigma^2/2$ tests the wall asymptotics of
Corollary~\ref{sh:cor:bessel}; and verify the linear response
$\beta_t\approx1+s/4t$ at large $t/s$, whose coefficient $\tfrac14$ tests
Proposition~\ref{sh:prop:rate}.

% ---------------------------------------------------------------------
\label{tc:sub:numerics}

Four tests specific to this part, additional to the driftless case's.

\emph{(i) Freezing.} Simulate $B_t$ for $n\in\{10^{3},10^{4},10^{5}\}$
and plot $\log B_t/\log n$ against $\theta$; expect collapse onto
\eqref{tc:eq:freezing} with tangential contact at $\theta=\sqrt2$, and
sample-to-sample fluctuations $o(1)$ before the transition and $O(1)$
after, the frozen phase being non-self-averaging.

\emph{(ii) The two coordinates under conditioning.} Verify the sign
truncation and the anatomy selection separately, since they are the
section's two measured claims and they differ in kind. For the first,
estimate $\PP(x_T<0\mid T_{\text{death}}>T)$ by the particle filter and
compare with the unconditioned $\tfrac12$ of
Theorem~\ref{tc:thm:half}; the prediction is an abrupt collapse, to
order $10^{-3}$ within a few multiples of the intrinsic timescale. For
the second, place a prior on the phase --- $\rho\sim
\mathrm{Unif}(-1,3)$ serves --- and track the conditioned
$\PP(\rho>0)$ against horizon; the prediction is a \emph{gradual} rise
from the prior together with steady truncation of the negative tail. A
selection on $\rho$ as abrupt as the selection on $\operatorname{sign}
x$, or an absence of selection on $\rho$ altogether, would each
contradict \S\ref{gr1:sub:anatomy}.

\emph{(iii) Exponent robustness.} Estimate the survival tail for a grid
of $(\epsilon,n,\Gamma,\rho)$ and fit the local slope. Prediction:
$-\tfrac14$ throughout, with no dependence on any of the four ---
this is Theorem~\ref{tc:thm:survives} and the section's central
falsifiable claim.

\emph{(iv) The gain profile.} Compute $\Psi'$ from
\eqref{tc:eq:Psiprime} and check the limiting values of
Theorem~\ref{tc:thm:gain}(iii); then test
Proposition~\ref{tc:prop:mingain} by regressing $\log\min\Psi'$ on
$\rho b+k$ across a grid of $\rho$. The predicted slope is
$-1/(1+\rho)$, the
predicted intercept $\log\kappa(\rho)$, and the minimizer is at
\eqref{tc:eq:xstar}; all three are sharp, making this the most stringent
test in the section.

\emph{Falsifiers.} A fitted exponent differing from $-\tfrac14$, or
depending on $\Gamma$ or $\rho$ (contradicting
Theorem~\ref{tc:thm:survives}); a gain reaching a floor of $1$ rather
than $0$ on the plateau (contradicting
Remark~\ref{tc:rem:speccorr}(b)); an absence of anatomy selection under
conditioning (contradicting Theorem~\ref{gr1:thm:trapped}, which makes
adverse anatomy eventually fatal and hence selectable).


% =====================================================================
\PrintRefs

\end{document}
